\documentclass[10pt,a4paper]{article}
\usepackage[margin=0.75in]{geometry}
\usepackage{amsmath,amssymb,amsthm,mathtools}
\usepackage{array,booktabs,tabularx,graphicx,microtype}
\usepackage{needspace,xurl,longtable}
\usepackage[section]{placeins}
\usepackage{tikz,pgfplots}
\usepgfplotslibrary{groupplots}
\pgfplotsset{compat=1.18}
\usetikzlibrary{arrows.meta}
\usepackage[scaled=0.92]{helvet}
\usepackage{sansmath}
\tikzset{every picture/.append style={execute at begin picture={\sansmath}}}
\usepackage[colorlinks=true,allcolors=blue]{hyperref}
\usepackage[numbers]{natbib}
\usepackage{titlesec}
\titleformat{\section}{\normalfont\large\bfseries}{Supplementary Note~\thesection.}{0.5em}{}
\renewcommand{\figurename}{Supplementary Fig.}
\renewcommand{\tablename}{Supplementary Table}
\providecommand{\bibcommenthead}{}
\newtheorem{proposition}{Proposition}
\renewcommand{\theproposition}{S\arabic{proposition}}
\newcommand{\KL}{\operatorname{KL}}
\newcolumntype{Y}{>{\raggedright\arraybackslash}X}
% Generated by extension/identified_set/make_manuscript_assets.py. Do not edit.
\newcommand{\IdRecipientQueries}{8{,}586}
\newcommand{\IdUnidentifiedCount}{8{,}586}
\newcommand{\IdUnidentifiedPct}{100.0}
\newcommand{\IdUnidCognatePct}{100.0}
\newcommand{\IdCertifiedEndpoints}{17{,}172}
\newcommand{\IdCertMaxDiff}{4.1\times10^{-11}}
\newcommand{\IdScalingFailures}{0}
\newcommand{\IdUnidCam}{100.0}
\newcommand{\IdUnidCognateCam}{100.0}
\newcommand{\IdWidthCam}{14.8}
\newcommand{\IdWidthLoCam}{8.9}
\newcommand{\IdWidthHiCam}{22.9}
\newcommand{\IdWidthMinCam}{5.0}
\newcommand{\IdRiskCam}{0.266}
\newcommand{\IdHeldFreqCam}{0.0223}
\newcommand{\IdHeldFullCam}{0.0186}
\newcommand{\IdSpearRiskCam}{0.03}
\newcommand{\IdSpearProdCam}{0.03}
\newcommand{\IdRevCam}{374}
\newcommand{\IdRevPctCam}{19.2}
\newcommand{\IdRevTiesCam}{8}
\newcommand{\IdAgreeCam}{59.6}
\newcommand{\IdAgreeNCam}{218/366}
\newcommand{\IdAgreeLoCam}{55.0}
\newcommand{\IdAgreeHiCam}{64.2}
\newcommand{\IdAgreeAdjLoCam}{53.8}
\newcommand{\IdAgreeAdjHiCam}{65.5}
\newcommand{\IdNonemptyCam}{59.5}
\newcommand{\IdNonemptyAdjLoCam}{54.1}
\newcommand{\IdRevGainShareCam}{18.8}
\newcommand{\IdUnidNew}{100.0}
\newcommand{\IdUnidCognateNew}{100.0}
\newcommand{\IdWidthNew}{15.1}
\newcommand{\IdWidthLoNew}{11.1}
\newcommand{\IdWidthHiNew}{19.7}
\newcommand{\IdWidthMinNew}{5.3}
\newcommand{\IdRiskNew}{0.319}
\newcommand{\IdHeldFreqNew}{0.0323}
\newcommand{\IdHeldFullNew}{0.0229}
\newcommand{\IdSpearRiskNew}{0.08}
\newcommand{\IdSpearProdNew}{0.09}
\newcommand{\IdRevNew}{1{,}029}
\newcommand{\IdRevPctNew}{22.7}
\newcommand{\IdRevTiesNew}{65}
\newcommand{\IdAgreeNew}{75.3}
\newcommand{\IdAgreeNNew}{726/964}
\newcommand{\IdAgreeLoNew}{71.8}
\newcommand{\IdAgreeHiNew}{78.7}
\newcommand{\IdAgreeAdjLoNew}{70.9}
\newcommand{\IdAgreeAdjHiNew}{79.7}
\newcommand{\IdNonemptyNew}{75.9}
\newcommand{\IdNonemptyAdjLoNew}{71.6}
\newcommand{\IdRevGainShareNew}{37.3}
\newcommand{\IdUnidTall}{100.0}
\newcommand{\IdUnidCognateTall}{100.0}
\newcommand{\IdWidthTall}{16.4}
\newcommand{\IdWidthLoTall}{12.2}
\newcommand{\IdWidthHiTall}{22.2}
\newcommand{\IdWidthMinTall}{5.6}
\newcommand{\IdRiskTall}{0.370}
\newcommand{\IdHeldFreqTall}{0.0585}
\newcommand{\IdHeldFullTall}{0.0260}
\newcommand{\IdSpearRiskTall}{0.04}
\newcommand{\IdSpearProdTall}{0.03}
\newcommand{\IdRevTall}{344}
\newcommand{\IdRevPctTall}{28.3}
\newcommand{\IdRevTiesTall}{3}
\newcommand{\IdAgreeTall}{84.5}
\newcommand{\IdAgreeNTall}{288/341}
\newcommand{\IdAgreeLoTall}{80.6}
\newcommand{\IdAgreeHiTall}{88.1}
\newcommand{\IdAgreeAdjLoTall}{79.5}
\newcommand{\IdAgreeAdjHiTall}{89.0}
\newcommand{\IdNonemptyTall}{83.3}
\newcommand{\IdNonemptyAdjLoTall}{77.7}
\newcommand{\IdRevGainShareTall}{37.5}
\newcommand{\IdUnidColon}{100.0}
\newcommand{\IdUnidCognateColon}{100.0}
\newcommand{\IdWidthColon}{12.6}
\newcommand{\IdWidthLoColon}{9.6}
\newcommand{\IdWidthHiColon}{19.5}
\newcommand{\IdWidthMinColon}{5.6}
\newcommand{\IdRiskColon}{0.265}
\newcommand{\IdHeldFreqColon}{0.0259}
\newcommand{\IdHeldFullColon}{0.0218}
\newcommand{\IdSpearRiskColon}{0.11}
\newcommand{\IdSpearProdColon}{0.10}
\newcommand{\IdRevColon}{201}
\newcommand{\IdRevPctColon}{22.6}
\newcommand{\IdRevTiesColon}{4}
\newcommand{\IdAgreeColon}{68.0}
\newcommand{\IdAgreeNColon}{134/197}
\newcommand{\IdAgreeLoColon}{59.1}
\newcommand{\IdAgreeHiColon}{76.3}
\newcommand{\IdAgreeAdjLoColon}{56.7}
\newcommand{\IdAgreeAdjHiColon}{78.7}
\newcommand{\IdNonemptyColon}{69.2}
\newcommand{\IdNonemptyAdjLoColon}{57.8}
\newcommand{\IdRevGainShareColon}{67.3}
\newcommand{\IdAgreeRetroMin}{60}
\newcommand{\IdAgreeRetroMax}{84}
\newcommand{\IdRevPctMin}{19}
\newcommand{\IdRevPctMax}{28}
\newcommand{\IdAgreeColonRound}{68}
\newcommand{\IdFreeExtLowMedian}{0.89}
\newcommand{\IdFreeExtHighMedian}{1.16}
\newcommand{\IdFreeExtMax}{5.40}
\newcommand{\IdFreeStarWidth}{10.6}
\newcommand{\IdFreeWidth}{12.5}

% Generated by extension/spectral_transfer/make_scale_assets.py. Do not edit.
\newcommand{\SoPairCam}{91.1}
\newcommand{\SoPairLoCam}{88.9}
\newcommand{\SoPairHiCam}{93.0}
\newcommand{\SoClosedCam}{75.3}
\newcommand{\SoClosedLoCam}{59.7}
\newcommand{\SoClosedHiCam}{86.4}
\newcommand{\SoFirstCam}{\ensuremath{-}202}
\newcommand{\SoPairNew}{98.0}
\newcommand{\SoPairLoNew}{97.1}
\newcommand{\SoPairHiNew}{98.8}
\newcommand{\SoClosedNew}{72.2}
\newcommand{\SoClosedLoNew}{62.9}
\newcommand{\SoClosedHiNew}{80.0}
\newcommand{\SoFirstNew}{\ensuremath{-}314}
\newcommand{\SoPairTall}{96.2}
\newcommand{\SoPairLoTall}{94.6}
\newcommand{\SoPairHiTall}{97.4}
\newcommand{\SoClosedTall}{88.8}
\newcommand{\SoClosedLoTall}{83.6}
\newcommand{\SoClosedHiTall}{92.8}
\newcommand{\SoFirstTall}{\ensuremath{-}246}
\newcommand{\SoPairColon}{98.0}
\newcommand{\SoPairLoColon}{96.2}
\newcommand{\SoPairHiColon}{99.8}
\newcommand{\SoClosedColon}{66.4}
\newcommand{\SoClosedLoColon}{47.2}
\newcommand{\SoClosedHiColon}{83.2}
\newcommand{\SoFirstColon}{\ensuremath{-}15}
\newcommand{\SoClipped}{116}
\newcommand{\SoPairMin}{91}
\newcommand{\SoPairMax}{98}
\newcommand{\SoClosedMin}{66}
\newcommand{\SoClosedMax}{89}
\newcommand{\ScDonors}{12}
\newcommand{\ScGenes}{209}
\newcommand{\ScProteins}{177}
\newcommand{\ScPairs}{36{,}993}
\newcommand{\ScErrIndep}{1.000}
\newcommand{\ScCorrIndep}{0.000}
\newcommand{\ScRsqIndep}{0.000}
\newcommand{\ScErrVar}{0.936}
\newcommand{\ScCorrVar}{0.438}
\newcommand{\ScRsqVar}{\ensuremath{-}0.393}
\newcommand{\ScErrRefCorr}{0.477}
\newcommand{\ScCorrRefCorr}{0.882}
\newcommand{\ScRsqRefCorr}{\ensuremath{-}6.523}
\newcommand{\ScErrRefCov}{0.669}
\newcommand{\ScCorrRefCov}{0.822}
\newcommand{\ScRsqRefCov}{\ensuremath{-}180}
\newcommand{\ScErrReg}{0.380}
\newcommand{\ScCorrReg}{0.922}
\newcommand{\ScRsqReg}{0.283}
\newcommand{\ScErrCells}{0.371}
\newcommand{\ScCorrCells}{0.927}
\newcommand{\ScRsqCells}{0.160}
\newcommand{\ScErrClosed}{0.341}
\newcommand{\ScCorrClosed}{0.934}
\newcommand{\ScRsqClosed}{0.285}
\newcommand{\ScErrFirst}{678}
\newcommand{\ScCorrFirst}{0.708}
\newcommand{\ScRsqFirst}{\ensuremath{-}125{,}313}
\newcommand{\ScErrOracle}{0.367}
\newcommand{\ScCorrOracle}{0.924}
\newcommand{\ScRsqOracle}{0.018}
\newcommand{\ScErrPfMeans}{0.536}
\newcommand{\ScCorrPfMeans}{0.840}
\newcommand{\ScRsqPfMeans}{0.089}
\newcommand{\ScErrPfLow}{0.460}
\newcommand{\ScCorrPfLow}{0.881}
\newcommand{\ScRsqPfLow}{0.115}
\newcommand{\ScErrPfMeansReg}{0.740}
\newcommand{\ScCorrPfMeansReg}{0.731}
\newcommand{\ScRsqPfMeansReg}{0.014}
\newcommand{\ScErrPfLowReg}{0.515}
\newcommand{\ScCorrPfLowReg}{0.855}
\newcommand{\ScRsqPfLowReg}{0.116}
\newcommand{\ScDiffVar}{\ensuremath{-}0.595}
\newcommand{\ScDiffLoVar}{\ensuremath{-}0.670}
\newcommand{\ScDiffHiVar}{\ensuremath{-}0.497}
\newcommand{\ScWinsVar}{12/12}
\newcommand{\ScDiffReg}{\ensuremath{-}0.039}
\newcommand{\ScDiffLoReg}{\ensuremath{-}0.056}
\newcommand{\ScDiffHiReg}{\ensuremath{-}0.024}
\newcommand{\ScWinsReg}{12/12}
\newcommand{\ScDiffRefCorr}{\ensuremath{-}0.136}
\newcommand{\ScDiffLoRefCorr}{\ensuremath{-}0.186}
\newcommand{\ScDiffHiRefCorr}{\ensuremath{-}0.090}
\newcommand{\ScWinsRefCorr}{12/12}
\newcommand{\ScDiffRefCov}{\ensuremath{-}0.328}
\newcommand{\ScDiffLoRefCov}{\ensuremath{-}0.401}
\newcommand{\ScDiffHiRefCov}{\ensuremath{-}0.256}
\newcommand{\ScWinsRefCov}{12/12}
\newcommand{\ScDiffCells}{0.030}
\newcommand{\ScDiffLoCells}{0.020}
\newcommand{\ScDiffHiCells}{0.040}
\newcommand{\PfDiffClosed}{0.119}
\newcommand{\PfDiffLoClosed}{0.091}
\newcommand{\PfDiffHiClosed}{0.155}
\newcommand{\PfDiffReg}{0.080}
\newcommand{\PfDiffLoReg}{0.042}
\newcommand{\PfDiffHiReg}{0.115}
\newcommand{\PfDiffVar}{\ensuremath{-}0.476}
\newcommand{\PfDiffLoVar}{\ensuremath{-}0.561}
\newcommand{\PfDiffHiVar}{\ensuremath{-}0.381}
\newcommand{\PfDiffIndep}{\ensuremath{-}0.540}
\newcommand{\PfDiffLoIndep}{\ensuremath{-}0.618}
\newcommand{\PfDiffHiIndep}{\ensuremath{-}0.448}
\newcommand{\PfShareOfPaired}{82}
\newcommand{\PfShareOfPairedMeans}{70}
\newcommand{\HcErrIndep}{1.000}
\newcommand{\HcRsqIndep}{0.000}
\newcommand{\HcErrVar}{0.943}
\newcommand{\HcRsqVar}{\ensuremath{-}0.799}
\newcommand{\HcErrRefCorr}{0.510}
\newcommand{\HcRsqRefCorr}{\ensuremath{-}9.618}
\newcommand{\HcErrRefCov}{0.705}
\newcommand{\HcRsqRefCov}{\ensuremath{-}258}
\newcommand{\HcErrReg}{0.424}
\newcommand{\HcRsqReg}{0.259}
\newcommand{\HcErrCells}{0.397}
\newcommand{\HcRsqCells}{0.102}
\newcommand{\HcErrClosed}{0.367}
\newcommand{\HcRsqClosed}{0.263}
\newcommand{\HcErrFirst}{1{,}041}
\newcommand{\HcRsqFirst}{\ensuremath{-}351{,}681}
\newcommand{\HcErrOracle}{0.384}
\newcommand{\HcRsqOracle}{\ensuremath{-}0.028}
\newcommand{\HcDiffVar}{\ensuremath{-}0.576}
\newcommand{\HcDiffLoVar}{\ensuremath{-}0.670}
\newcommand{\HcDiffHiVar}{\ensuremath{-}0.444}
\newcommand{\HcDiffReg}{\ensuremath{-}0.057}
\newcommand{\HcDiffLoReg}{\ensuremath{-}0.082}
\newcommand{\HcDiffHiReg}{\ensuremath{-}0.031}
\newcommand{\HcRecipients}{8}
\newcommand{\PcObs}{0.536}
\newcommand{\PcNull}{1.255}
\newcommand{\PcPerms}{200}
\newcommand{\PcP}{1/201}
\newcommand{\PcDonors}{12}
\newcommand{\LcErrFive}{0.723}
\newcommand{\LcErrTen}{0.610}
\newcommand{\LcErrFifteen}{0.553}
\newcommand{\LcErrAll}{0.536}
\newcommand{\PfRidgeEdge}{11 of 12}
\newcommand{\PfSamples}{19--20}

% Generated by extension/cross_study/make_assets.py. Do not edit.
\newcommand{\XsPatients}{118}
\newcommand{\XsTrainCells}{223{,}640}
\newcommand{\XsHalfCells}{111{,}816}
\newcommand{\XsGenes}{206}
\newcommand{\XsProteins}{122}
\newcommand{\XsPfScale}{0.01}
\newcommand{\XsPfRank}{32}
\newcommand{\XsPfMomentIter}{9{,}858}
\newcommand{\XsPfMomentGrad}{2.4\times10^{-4}}
\newcommand{\XsPairedRidge}{0.01}
\newcommand{\XsRegLambda}{0.1}
\newcommand{\XsHaoCells}{161{,}764}
\newcommand{\XsHaoScoring}{80{,}884}
\newcommand{\XsColonScoring}{7{,}979}
\newcommand{\XsHaoDonors}{8}
\newcommand{\XsHaoErrPfMeans}{0.292}
\newcommand{\XsHaoRhoPfMeans}{0.964}
\newcommand{\XsHaoErrPfMoment}{0.389}
\newcommand{\XsHaoRhoPfMoment}{0.957}
\newcommand{\XsHaoErrPaired}{0.134}
\newcommand{\XsHaoRhoPaired}{0.995}
\newcommand{\XsHaoErrReg}{0.522}
\newcommand{\XsHaoRhoReg}{0.901}
\newcommand{\XsHaoErrCorr}{0.630}
\newcommand{\XsHaoRhoCorr}{0.846}
\newcommand{\XsHaoErrCov}{0.951}
\newcommand{\XsHaoRhoCov}{0.570}
\newcommand{\XsHaoErrVar}{0.976}
\newcommand{\XsHaoRhoVar}{0.394}
\newcommand{\XsHaoErrBench}{0.120}
\newcommand{\XsHaoRhoBench}{0.998}
\newcommand{\XsHaoErrIndep}{1.000}
\newcommand{\XsHaoRhoIndep}{0.000}
\newcommand{\XsHaoNullErr}{0.949}
\newcommand{\XsHaoNullMin}{0.639}
\newcommand{\XsHaoPermP}{0.005}
\newcommand{\XsHaoPermMaxDonorP}{0.005}
\newcommand{\XsHaoShare}{82}
\newcommand{\XsHaoShareLo}{80}
\newcommand{\XsHaoShareHi}{84}
\newcommand{\XsHaoShareMoment}{71}
\newcommand{\XsHaoShareMomentLo}{69}
\newcommand{\XsHaoShareMomentHi}{72}
\newcommand{\XsHaoWErrPfMeans}{1.082}
\newcommand{\XsHaoWErrPfMeansLo}{1.014}
\newcommand{\XsHaoWErrPfMeansHi}{1.162}
\newcommand{\XsHaoWBelowPfMeans}{2}
\newcommand{\XsHaoWErrPfMoment}{1.038}
\newcommand{\XsHaoWErrPfMomentLo}{0.981}
\newcommand{\XsHaoWErrPfMomentHi}{1.104}
\newcommand{\XsHaoWBelowPfMoment}{4}
\newcommand{\XsHaoWErrPaired}{0.528}
\newcommand{\XsHaoWErrPairedLo}{0.502}
\newcommand{\XsHaoWErrPairedHi}{0.565}
\newcommand{\XsHaoWBelowPaired}{8}
\newcommand{\XsHaoWErrBench}{0.348}
\newcommand{\XsHaoWErrBenchLo}{0.323}
\newcommand{\XsHaoWErrBenchHi}{0.378}
\newcommand{\XsHaoWBelowBench}{8}
\newcommand{\XsHaoWShare}{\ensuremath{-}20}
\newcommand{\XsHaoWShareLo}{\ensuremath{-}41}
\newcommand{\XsHaoWShareHi}{\ensuremath{-}3}
\newcommand{\XsHaoDiffReg}{0.388}
\newcommand{\XsHaoDiffRegLo}{0.364}
\newcommand{\XsHaoDiffRegHi}{0.415}
\newcommand{\XsHaoFavReg}{8}
\newcommand{\XsHaoDiffCorr}{0.495}
\newcommand{\XsHaoDiffCorrLo}{0.466}
\newcommand{\XsHaoDiffCorrHi}{0.531}
\newcommand{\XsHaoFavCorr}{8}
\newcommand{\XsHaoDiffCov}{0.817}
\newcommand{\XsHaoDiffCovLo}{0.615}
\newcommand{\XsHaoDiffCovHi}{1.078}
\newcommand{\XsHaoFavCov}{8}
\newcommand{\XsHaoDiffVar}{0.842}
\newcommand{\XsHaoDiffVarLo}{0.834}
\newcommand{\XsHaoDiffVarHi}{0.852}
\newcommand{\XsHaoFavVar}{8}
\newcommand{\XsHaoMomentMinusMeans}{0.097}
\newcommand{\XsHaoMomentMinusMeansLo}{0.078}
\newcommand{\XsHaoMomentMinusMeansHi}{0.119}
\newcommand{\XsHaoMomentBetter}{0}
\newcommand{\XsHaoPfMeansBetterThanPaired}{0}
\newcommand{\XsHaoPfMomentBetterThanPaired}{0}
\newcommand{\XsHaoPfMeansBetterThanBench}{0}
\newcommand{\XsHaoPairedBetterThanBench}{1}
\newcommand{\XsColonDonors}{12}
\newcommand{\XsColonErrPfMeans}{0.988}
\newcommand{\XsColonRhoPfMeans}{0.397}
\newcommand{\XsColonErrPfMoment}{0.907}
\newcommand{\XsColonRhoPfMoment}{0.463}
\newcommand{\XsColonErrPaired}{0.719}
\newcommand{\XsColonRhoPaired}{0.689}
\newcommand{\XsColonErrReg}{0.859}
\newcommand{\XsColonRhoReg}{0.516}
\newcommand{\XsColonErrCorr}{1.132}
\newcommand{\XsColonRhoCorr}{0.416}
\newcommand{\XsColonErrCov}{4.651}
\newcommand{\XsColonRhoCov}{0.230}
\newcommand{\XsColonErrVar}{0.971}
\newcommand{\XsColonRhoVar}{0.278}
\newcommand{\XsColonErrBench}{0.539}
\newcommand{\XsColonRhoBench}{0.842}
\newcommand{\XsColonErrIndep}{1.000}
\newcommand{\XsColonRhoIndep}{0.000}
\newcommand{\XsColonNullErr}{1.168}
\newcommand{\XsColonNullMin}{1.019}
\newcommand{\XsColonPermP}{0.005}
\newcommand{\XsColonPermMaxDonorP}{0.090}
\newcommand{\XsColonShare}{\ensuremath{-}12}
\newcommand{\XsColonShareLo}{\ensuremath{-}49}
\newcommand{\XsColonShareHi}{15}
\newcommand{\XsColonShareMoment}{25}
\newcommand{\XsColonShareMomentLo}{5}
\newcommand{\XsColonShareMomentHi}{42}
\newcommand{\XsColonWErrPfMeans}{1.274}
\newcommand{\XsColonWErrPfMeansLo}{1.239}
\newcommand{\XsColonWErrPfMeansHi}{1.307}
\newcommand{\XsColonWBelowPfMeans}{0}
\newcommand{\XsColonWErrPfMoment}{1.227}
\newcommand{\XsColonWErrPfMomentLo}{1.200}
\newcommand{\XsColonWErrPfMomentHi}{1.253}
\newcommand{\XsColonWBelowPfMoment}{0}
\newcommand{\XsColonWErrPaired}{1.000}
\newcommand{\XsColonWErrPairedLo}{0.987}
\newcommand{\XsColonWErrPairedHi}{1.013}
\newcommand{\XsColonWBelowPaired}{7}
\newcommand{\XsColonWErrBench}{0.936}
\newcommand{\XsColonWErrBenchLo}{0.922}
\newcommand{\XsColonWErrBenchHi}{0.953}
\newcommand{\XsColonWBelowBench}{11}
\newcommand{\XsColonWShare}{\ensuremath{-}4126}
\newcommand{\XsColonWShareLo}{\ensuremath{-}12228}
\newcommand{\XsColonWShareHi}{891}
\newcommand{\XsColonDiffReg}{0.140}
\newcommand{\XsColonDiffRegLo}{0.099}
\newcommand{\XsColonDiffRegHi}{0.180}
\newcommand{\XsColonFavReg}{12}
\newcommand{\XsColonDiffCorr}{0.413}
\newcommand{\XsColonDiffCorrLo}{0.348}
\newcommand{\XsColonDiffCorrHi}{0.500}
\newcommand{\XsColonFavCorr}{12}
\newcommand{\XsColonDiffCov}{3.932}
\newcommand{\XsColonDiffCovLo}{2.510}
\newcommand{\XsColonDiffCovHi}{5.443}
\newcommand{\XsColonFavCov}{12}
\newcommand{\XsColonDiffVar}{0.252}
\newcommand{\XsColonDiffVarLo}{0.182}
\newcommand{\XsColonDiffVarHi}{0.324}
\newcommand{\XsColonFavVar}{12}
\newcommand{\XsColonMomentMinusMeans}{\ensuremath{-}0.081}
\newcommand{\XsColonMomentMinusMeansLo}{\ensuremath{-}0.106}
\newcommand{\XsColonMomentMinusMeansHi}{\ensuremath{-}0.053}
\newcommand{\XsColonMomentBetter}{11}
\newcommand{\XsColonPfMeansBetterThanPaired}{0}
\newcommand{\XsColonPfMomentBetterThanPaired}{0}
\newcommand{\XsColonPfMeansBetterThanBench}{0}
\newcommand{\XsColonPairedBetterThanBench}{0}
\newcommand{\XsHaoWRhoPfMeans}{0.43}
\newcommand{\XsHaoWRhoPaired}{0.85}
\newcommand{\XsHaoWRhoBench}{0.96}
\newcommand{\XsHaoLTwoErrPfMeans}{1.484}
\newcommand{\XsHaoLTwoErrPfMoment}{1.435}
\newcommand{\XsHaoLTwoErrPaired}{0.868}
\newcommand{\XsHaoLTwoErrBench}{0.541}
\newcommand{\XsDimM}{11}
\newcommand{\XsEffDim}{16.9}
\newcommand{\XsHaoIdRho}{38}
\newcommand{\XsHaoIdRhoMin}{35}
\newcommand{\XsHaoIdRhoMax}{41}
\newcommand{\XsHaoIdFloor}{0.16}
\newcommand{\XsHaoIdFloorMin}{0.13}
\newcommand{\XsHaoIdFloorMax}{0.20}
\newcommand{\XsHaoIdRadius}{0.88}
\newcommand{\XsHaoIdRadiusMin}{0.86}
\newcommand{\XsHaoIdRadiusMax}{0.90}
\newcommand{\XsHaoWIdRho}{12}
\newcommand{\XsHaoWIdRhoMin}{11}
\newcommand{\XsHaoWIdRhoMax}{13}
\newcommand{\XsHaoWIdFloor}{0.72}
\newcommand{\XsHaoWIdFloorMin}{0.61}
\newcommand{\XsHaoWIdFloorMax}{0.82}
\newcommand{\XsHaoWIdRadius}{1.84}
\newcommand{\XsHaoWIdRadiusMin}{1.70}
\newcommand{\XsHaoWIdRadiusMax}{2.23}
\newcommand{\XsColonIdRho}{19}
\newcommand{\XsColonIdRhoMin}{15}
\newcommand{\XsColonIdRhoMax}{24}
\newcommand{\XsColonIdFloor}{0.71}
\newcommand{\XsColonIdFloorMin}{0.59}
\newcommand{\XsColonIdFloorMax}{0.87}
\newcommand{\XsColonIdRadius}{1.32}
\newcommand{\XsColonIdRadiusMin}{1.05}
\newcommand{\XsColonIdRadiusMax}{1.51}
\newcommand{\XsColonWIdRho}{9}
\newcommand{\XsColonWIdRhoMin}{8}
\newcommand{\XsColonWIdRhoMax}{10}
\newcommand{\XsColonWIdFloor}{0.92}
\newcommand{\XsColonWIdFloorMin}{0.87}
\newcommand{\XsColonWIdFloorMax}{0.98}
\newcommand{\XsColonWIdRadius}{1.89}
\newcommand{\XsColonWIdRadiusMin}{1.52}
\newcommand{\XsColonWIdRadiusMax}{2.26}
\newcommand{\XsIdUnits}{48}
\newcommand{\XsFloorErrCorr}{0.89}
\newcommand{\XsRhoErrCorr}{\ensuremath{-}0.94}
\newcommand{\XsMaxFMin}{0.44}
\newcommand{\XsMaxFMax}{4.29}
\newcommand{\XsMaxFAboveOne}{46}
\newcommand{\XsMaxFTotalMin}{1.02}
\newcommand{\XsGeneSpearmanUnidHao}{0.86}
\newcommand{\XsGeneSpearmanAxisHao}{0.17}
\newcommand{\XsGeneSpearmanUnidColon}{0.72}
\newcommand{\XsGeneSpearmanAxisColon}{\ensuremath{-}0.01}
\newcommand{\XsRestartBChange}{9\times10^{-7}}
\newcommand{\XsRestartEChange}{2\times10^{-6}}
\newcommand{\XsLcHaoEight}{0.637}
\newcommand{\XsLcColonEight}{1.024}
\newcommand{\XsLcDimEight}{3.7}
\newcommand{\XsLcHaoSixteen}{0.493}
\newcommand{\XsLcColonSixteen}{0.992}
\newcommand{\XsLcDimSixteen}{5.1}
\newcommand{\XsLcHaoThirtyTwo}{0.395}
\newcommand{\XsLcColonThirtyTwo}{0.931}
\newcommand{\XsLcDimThirtyTwo}{4.2}
\newcommand{\XsLcHaoSixtyFour}{0.336}
\newcommand{\XsLcColonSixtyFour}{0.996}
\newcommand{\XsLcDimSixtyFour}{11.4}
\newcommand{\XsLcHaoAll}{0.292}
\newcommand{\XsLcColonAll}{0.988}
\newcommand{\XsLcDimAll}{11}
\newcommand{\XsRigErrMeans}{0.566}
\newcommand{\XsRigErrMoment}{0.680}
\newcommand{\XsRigNullMeans}{1.086}
\newcommand{\XsRigNullMoment}{1.048}
\newcommand{\XsRigMaxPMeans}{0.005}
\newcommand{\XsRigMaxPMoment}{0.143}
\newcommand{\XsRigDonorsSigMeans}{12}
\newcommand{\XsRigDonorsSigMoment}{10}
\newcommand{\XsRigEdge}{0}
\newcommand{\XsRigExtended}{0}
\newcommand{\XsRigRankMin}{8}
\newcommand{\XsRigRankMax}{8}
\newcommand{\XsRigConverged}{9}
\newcommand{\XsRigDimMin}{11}
\newcommand{\XsRigDimMax}{12}
\newcommand{\XsRigBiopsiesMin}{19}
\newcommand{\XsRigBiopsiesMax}{20}
\newcommand{\XsRigIdRho}{37}
\newcommand{\XsRigIdFloor}{0.59}
\newcommand{\XsRigIdRadius}{1.13}
\newcommand{\XsRigMaxF}{1.46}
\newcommand{\XsRigShareMeans}{65}
\newcommand{\XsRigShareMoment}{50}
\newcommand{\XsPoneUnits}{512}
\newcommand{\XsPoneDim}{86}
\newcommand{\XsPoneErr}{1.287}
\newcommand{\XsPoneErrLTwo}{2.163}
\newcommand{\XsPoneErrTotal}{0.660}
\newcommand{\XsPoneRho}{50}
\newcommand{\XsPoneFloor}{0.28}
\newcommand{\XsMixHaoBetweenShare}{97}
\newcommand{\XsMixHaoBetweenShareMin}{96}
\newcommand{\XsMixHaoBetweenShareMax}{98}
\newcommand{\XsMixHaoWithinShare}{11}
\newcommand{\XsMixHaoWithinShareMin}{10}
\newcommand{\XsMixHaoWithinShareMax}{14}
\newcommand{\XsMixHaoErrBetween}{0.71}
\newcommand{\XsMixHaoErrBetweenMin}{0.65}
\newcommand{\XsMixHaoErrBetweenMax}{0.76}
\newcommand{\XsMixHaoErrBetweenBelow}{8}
\newcommand{\XsMixHaoErrWithin}{1.92}
\newcommand{\XsMixHaoErrWithinMin}{1.55}
\newcommand{\XsMixHaoErrWithinMax}{2.56}
\newcommand{\XsMixHaoErrWithinBelow}{0}
\newcommand{\XsMixHaoErrDirect}{0.75}
\newcommand{\XsMixHaoErrDirectMin}{0.69}
\newcommand{\XsMixHaoErrDirectMax}{0.81}
\newcommand{\XsMixHaoErrDirectBelow}{8}
\newcommand{\XsMixHaoRBetween}{0.75}
\newcommand{\XsMixHaoRWithin}{0.15}
\newcommand{\XsMixHaoNormWithin}{1.8}
\newcommand{\XsMixHaoNormBetween}{0.99}
\newcommand{\XsMixHaoRhoBetween}{88}
\newcommand{\XsMixHaoRhoBetweenMin}{87}
\newcommand{\XsMixHaoRhoBetweenMax}{89}
\newcommand{\XsMixHaoRhoWithin}{12}
\newcommand{\XsMixHaoRhoWithinMin}{10}
\newcommand{\XsMixHaoRhoWithinMax}{13}
\newcommand{\XsMixHaoVarBetween}{39}
\newcommand{\XsMixHaoVarBetweenMin}{34}
\newcommand{\XsMixHaoVarBetweenMax}{42}
\newcommand{\XsMixColonBetweenShare}{81}
\newcommand{\XsMixColonBetweenShareMin}{45}
\newcommand{\XsMixColonBetweenShareMax}{94}
\newcommand{\XsMixColonWithinShare}{43}
\newcommand{\XsMixColonWithinShareMin}{27}
\newcommand{\XsMixColonWithinShareMax}{84}
\newcommand{\XsMixColonErrBetween}{1.14}
\newcommand{\XsMixColonErrBetweenMin}{0.85}
\newcommand{\XsMixColonErrBetweenMax}{1.29}
\newcommand{\XsMixColonErrBetweenBelow}{1}
\newcommand{\XsMixColonErrWithin}{1.38}
\newcommand{\XsMixColonErrWithinMin}{1.21}
\newcommand{\XsMixColonErrWithinMax}{1.53}
\newcommand{\XsMixColonErrWithinBelow}{0}
\newcommand{\XsMixColonErrDirect}{1.20}
\newcommand{\XsMixColonErrDirectMin}{0.91}
\newcommand{\XsMixColonErrDirectMax}{1.34}
\newcommand{\XsMixColonErrDirectBelow}{1}
\newcommand{\XsMixColonRBetween}{0.35}
\newcommand{\XsMixColonRWithin}{0.07}
\newcommand{\XsMixColonNormWithin}{1.0}
\newcommand{\XsMixColonNormBetween}{1.00}
\newcommand{\XsMixColonRhoBetween}{57}
\newcommand{\XsMixColonRhoBetweenMin}{50}
\newcommand{\XsMixColonRhoBetweenMax}{66}
\newcommand{\XsMixColonRhoWithin}{9}
\newcommand{\XsMixColonRhoWithinMin}{8}
\newcommand{\XsMixColonRhoWithinMax}{10}
\newcommand{\XsMixColonVarBetween}{20}
\newcommand{\XsMixColonVarBetweenMin}{9}
\newcommand{\XsMixColonVarBetweenMax}{28}
\newcommand{\XsRigContinued}{3}
\newcommand{\XsRigContinueConverged}{3}
\newcommand{\XsRigContinueMaxChange}{0.0015}
\newcommand{\XsRigContinueBound}{0.002}
\newcommand{\XsRigContinueMaxWChange}{0.031}
\newcommand{\XsPairFcgraTotTruth}{0.75}
\newcommand{\XsPairFcgraTotPfMeans}{0.29}
\newcommand{\XsPairFcgraTotPaired}{0.61}
\newcommand{\XsPairFcgraTotBench}{0.67}
\newcommand{\XsPairFcgraWTruth}{0.63}
\newcommand{\XsPairFcgraWPfMeans}{0.12}
\newcommand{\XsPairFcgraWPaired}{0.52}
\newcommand{\XsPairFcgraWBench}{0.46}
\newcommand{\XsPairKlrbTotTruth}{0.69}
\newcommand{\XsPairKlrbTotPfMeans}{0.26}
\newcommand{\XsPairKlrbTotPaired}{0.44}
\newcommand{\XsPairKlrbTotBench}{0.60}
\newcommand{\XsPairKlrbWTruth}{0.51}
\newcommand{\XsPairKlrbWPfMeans}{0.07}
\newcommand{\XsPairKlrbWPaired}{0.30}
\newcommand{\XsPairKlrbWBench}{0.28}
\newcommand{\XsPairIghdTotTruth}{0.76}
\newcommand{\XsPairIghdTotPfMeans}{0.42}
\newcommand{\XsPairIghdTotPaired}{0.60}
\newcommand{\XsPairIghdTotBench}{0.67}
\newcommand{\XsPairIghdWTruth}{0.51}
\newcommand{\XsPairIghdWPfMeans}{0.10}
\newcommand{\XsPairIghdWPaired}{0.36}
\newcommand{\XsPairIghdWBench}{0.35}
\newcommand{\XsPairIlrTotTruth}{0.82}
\newcommand{\XsPairIlrTotPfMeans}{0.57}
\newcommand{\XsPairIlrTotPaired}{0.68}
\newcommand{\XsPairIlrTotBench}{0.73}
\newcommand{\XsPairIlrWTruth}{0.48}
\newcommand{\XsPairIlrWPfMeans}{0.20}
\newcommand{\XsPairIlrWPaired}{0.39}
\newcommand{\XsPairIlrWBench}{0.29}
\newcommand{\XsPairLyzTotTruth}{0.89}
\newcommand{\XsPairLyzTotPfMeans}{0.77}
\newcommand{\XsPairLyzTotPaired}{0.81}
\newcommand{\XsPairLyzTotBench}{0.80}
\newcommand{\XsPairLyzWTruth}{0.21}
\newcommand{\XsPairLyzWPfMeans}{0.22}
\newcommand{\XsPairLyzWPaired}{0.32}
\newcommand{\XsPairLyzWBench}{0.17}
\newcommand{\XsPairCdaTotTruth}{0.87}
\newcommand{\XsPairCdaTotPfMeans}{0.61}
\newcommand{\XsPairCdaTotPaired}{0.77}
\newcommand{\XsPairCdaTotBench}{0.77}
\newcommand{\XsPairCdaWTruth}{0.05}
\newcommand{\XsPairCdaWPfMeans}{0.09}
\newcommand{\XsPairCdaWPaired}{0.21}
\newcommand{\XsPairCdaWBench}{0.04}

% Generated by extension/recoverability/make_assets.py. Do not edit.
\newcommand{\RcDevFReported}{2.78}
\newcommand{\RcDevFClr}{2.69}
\newcommand{\RcDevFRaw}{2.75}
\newcommand{\RcDevTotalAnalyses}{20}
\newcommand{\RcDevFMeasTotal}{2.78}
\newcommand{\RcDevFLatTotal}{0.28}
\newcommand{\RcDevFMeasAbove}{20}
\newcommand{\RcDevFMeasLowerAbove}{18}
\newcommand{\RcDevFLatBloodMin}{0.85}
\newcommand{\RcDevFLatBloodMax}{1.90}
\newcommand{\RcDevFLatBloodAbove}{6}
\newcommand{\RcDevFLatBloodLowerAbove}{6}
\newcommand{\RcDevFLatColonMax}{0.31}
\newcommand{\RcDevWithinAnalyses}{28}
\newcommand{\RcDevFMeasWithin}{3.23}
\newcommand{\RcDevFLatWithin}{0.51}
\newcommand{\RcDevFLatWithinLowerAbove}{4}
\newcommand{\RcDevPOneAnalyses}{48}
\newcommand{\RcDevPOneFLat}{3.02}
\newcommand{\RcDevPOneLowerAbove}{45}
\newcommand{\RcDevDeparture}{0.82}
\newcommand{\RcDevDepartureMin}{0.70}
\newcommand{\RcDevDepartureMax}{1.24}
\newcommand{\RcDevBoundsNonempty}{37}
\newcommand{\RcDevBoundsAnalyses}{48}
\newcommand{\RcDevBoundsSigned}{0}
\newcommand{\RcDevBoundsInsideMin}{0.928}
\newcommand{\RcDevLabelBlood}{0.12}
\newcommand{\RcDevLabelBloodMin}{0.10}
\newcommand{\RcDevLabelBloodMax}{0.16}
\newcommand{\RcDevPfBlood}{0.29}
\newcommand{\RcDevPairedBlood}{0.13}
\newcommand{\RcDevLabelColon}{0.57}
\newcommand{\RcDevLabelColonMin}{0.38}
\newcommand{\RcDevLabelColonMax}{0.95}
\newcommand{\RcDevPfColon}{0.99}
\newcommand{\RcDevPairedColon}{0.72}
\newcommand{\RcDevLabelBeatsPaired}{16}
\newcommand{\RcDevLabelUnits}{20}
\newcommand{\RcDevPairs}{1,776}
\newcommand{\RcBmSamples}{12}
\newcommand{\RcBmDonors}{9}
\newcommand{\RcBmCells}{90,261}
\newcommand{\RcBmGenes}{201}
\newcommand{\RcBmProteins}{88}
\newcommand{\RcBmFineMin}{24}
\newcommand{\RcBmFineMax}{36}
\newcommand{\RcBmFrozenTot}{0.52}
\newcommand{\RcBmFrozenMeasTot}{0.52}
\newcommand{\RcBmFrozenBelowTot}{12}
\newcommand{\RcBmPOneTot}{0.68}
\newcommand{\RcBmColonTot}{0.49}
\newcommand{\RcBmColonBelowTot}{12}
\newcommand{\RcBmBenchTot}{0.17}
\newcommand{\RcBmFrozenCoarse}{0.99}
\newcommand{\RcBmFrozenMeasCoarse}{1.02}
\newcommand{\RcBmFrozenBelowCoarse}{6}
\newcommand{\RcBmPOneCoarse}{1.01}
\newcommand{\RcBmColonCoarse}{0.93}
\newcommand{\RcBmColonBelowCoarse}{10}
\newcommand{\RcBmBenchCoarse}{0.41}
\newcommand{\RcBmFrozenFine}{1.30}
\newcommand{\RcBmFrozenMeasFine}{1.43}
\newcommand{\RcBmFrozenBelowFine}{0}
\newcommand{\RcBmPOneFine}{1.44}
\newcommand{\RcBmColonFine}{1.10}
\newcommand{\RcBmColonBelowFine}{2}
\newcommand{\RcBmBenchFine}{0.66}
\newcommand{\RcBmLabelTot}{0.21}
\newcommand{\RcBmFMeasTot}{2.80}
\newcommand{\RcBmFLatTot}{0.52}
\newcommand{\RcBmFLatMax}{0.94}
\newcommand{\RcBmFMeasMin}{1.85}
\newcommand{\RcBmFMeasLowerAbove}{36}
\newcommand{\RcBmColonFLatMax}{0.73}
\newcommand{\RcBmPOneFLatMin}{1.00}
\newcommand{\RcBmPOneFLatMax}{3.65}
\newcommand{\RcBmPOneLowerAbove}{35}
\newcommand{\RcBmAnalyses}{36}
\newcommand{\RcBmPairs}{1,476}
\newcommand{\RcBmFailures}{838}
\newcommand{\RcBmPermAbstain}{95}
\newcommand{\RcBmLcAbstain}{7.5}
\newcommand{\RcBmAucPrimary}{0.87}
\newcommand{\RcBmAucLabelFree}{0.84}
\newcommand{\RcBmAucCoverage}{0.84}
\newcommand{\RcBmAucComposition}{0.84}
\newcommand{\RcBmAucCells}{0.46}
\newcommand{\RcBmAucNTrain}{0.54}
\newcommand{\RcBmAucShift}{0.83}
\newcommand{\RcBmAucCompat}{0.29}
\newcommand{\RcBmAucTypeDis}{0.59}
\newcommand{\RcBmDiffPrimaryCoverage}{0.030}
\newcommand{\RcBmDiffPrimaryCoverageLo}{\ensuremath{-}0.014}
\newcommand{\RcBmDiffPrimaryCoverageHi}{0.059}
\newcommand{\RcBmDiffPrimaryComposition}{0.034}
\newcommand{\RcBmDiffPrimaryCompositionLo}{0.024}
\newcommand{\RcBmDiffPrimaryCompositionHi}{0.048}
\newcommand{\RcBmDiffPrimaryCells}{0.417}
\newcommand{\RcBmDiffPrimaryCellsLo}{0.257}
\newcommand{\RcBmDiffPrimaryCellsHi}{0.451}
\newcommand{\RcBmDiffPrimaryNTrain}{0.334}
\newcommand{\RcBmDiffPrimaryNTrainLo}{0.261}
\newcommand{\RcBmDiffPrimaryNTrainHi}{0.392}
\newcommand{\RcBmDiffPrimaryShift}{0.045}
\newcommand{\RcBmDiffPrimaryShiftLo}{\ensuremath{-}0.012}
\newcommand{\RcBmDiffPrimaryShiftHi}{0.086}
\newcommand{\RcBmDiffLabelFreeCoverage}{0.001}
\newcommand{\RcBmDiffLabelFreeCoverageLo}{\ensuremath{-}0.014}
\newcommand{\RcBmDiffLabelFreeCoverageHi}{0.020}
\newcommand{\RcBmHiPairs}{251}
\newcommand{\RcBmHiFailures}{19}
\newcommand{\RcBmHiPOne}{36}
\newcommand{\RcBmHiFailuresPOne}{18}
\newcommand{\RcBmHiAucPrimary}{0.99}
\newcommand{\RcBmHiAucPrimaryLo}{0.98}
\newcommand{\RcBmHiAucPrimaryHi}{1.00}
\newcommand{\RcBmHiAucTypeDis}{0.93}
\newcommand{\RcBmHiAucTypeDisLo}{0.90}
\newcommand{\RcBmHiAucTypeDisHi}{0.95}
\newcommand{\RcBmHiAucCompat}{0.89}
\newcommand{\RcBmHiAucCompatLo}{0.85}
\newcommand{\RcBmHiAucCompatHi}{0.92}
\newcommand{\RcBmHiAucLabelFree}{0.22}
\newcommand{\RcBmHiAucLabelFreeLo}{0.14}
\newcommand{\RcBmHiAucLabelFreeHi}{0.37}
\newcommand{\RcBmRankPrimary}{0.69}
\newcommand{\RcBmRankCoverage}{0.48}
\newcommand{\RcBmRankShift}{0.46}
\newcommand{\RcBmRefRandomFifty}{1.01}
\newcommand{\RcBmRefBalancedFifty}{1.04}
\newcommand{\RcBmRefGuidedFifty}{0.96}
\newcommand{\RcBmRefHybridFifty}{0.21}
\newcommand{\RcBmRefPairedFifty}{0.64}
\newcommand{\RcBmRefRandomHundred}{0.91}
\newcommand{\RcBmRefBalancedHundred}{0.93}
\newcommand{\RcBmRefGuidedHundred}{0.87}
\newcommand{\RcBmRefHybridHundred}{0.19}
\newcommand{\RcBmRefPairedHundred}{0.51}
\newcommand{\RcBmRefRandomTwoHundred}{0.80}
\newcommand{\RcBmRefBalancedTwoHundred}{0.84}
\newcommand{\RcBmRefGuidedTwoHundred}{0.77}
\newcommand{\RcBmRefHybridTwoHundred}{0.18}
\newcommand{\RcBmRefPairedTwoHundred}{0.38}
\newcommand{\RcBmRefRandomFourHundred}{0.69}
\newcommand{\RcBmRefBalancedFourHundred}{0.74}
\newcommand{\RcBmRefGuidedFourHundred}{0.67}
\newcommand{\RcBmRefHybridFourHundred}{0.16}
\newcommand{\RcBmRefPairedFourHundred}{0.29}
\newcommand{\RcBmRefRandomEightHundred}{0.58}
\newcommand{\RcBmRefBalancedEightHundred}{0.62}
\newcommand{\RcBmRefGuidedEightHundred}{0.56}
\newcommand{\RcBmRefHybridEightHundred}{0.14}
\newcommand{\RcBmRefPairedEightHundred}{0.21}
\newcommand{\RcBmRefBetweenOnly}{0.21}
\newcommand{\RcBmRefGainRandom}{0.034}
\newcommand{\RcBmRefGainRandomLo}{0.026}
\newcommand{\RcBmRefGainRandomHi}{0.041}
\newcommand{\RcBmRefGainBalanced}{0.068}
\newcommand{\RcBmRefGainBalancedLo}{0.056}
\newcommand{\RcBmRefGainBalancedHi}{0.087}
\newcommand{\RcBmRefAllSamples}{all}
\newcommand{\RcBmRefEquivTwoHundred}{243}
\newcommand{\RcBmRefEquivEightHundred}{above 800}
\newcommand{\RcBmBoundsSigned}{0}
\newcommand{\RcBmBoundsAnalyses}{36}
\newcommand{\RcBmBoundsHalfwidth}{0.74}
\newcommand{\RcBmBoundsTarget}{0.14}
\newcommand{\RcBmHOne}{not met}
\newcommand{\RcBmHThree}{met}
\newcommand{\RcBmHSix}{met}
\newcommand{\RcFmPairs}{1,776}
\newcommand{\RcFmFailures}{1,235}
\newcommand{\RcFmPrimaryAuc}{0.95}
\newcommand{\RcFmPrimaryCoverage}{\ensuremath{-}0.12}
\newcommand{\RcFmPrimaryLogF}{0.57}
\newcommand{\RcFmPrimaryAgreement}{\ensuremath{-}1.25}
\newcommand{\RcFmPrimaryComposition}{\ensuremath{-}2.97}
\newcommand{\RcFmLabelFreeAuc}{0.89}
\newcommand{\RcFmLabelFreeCoverage}{\ensuremath{-}2.23}
\newcommand{\RcFmLabelFreeLogF}{0.86}
\newcommand{\RcPhMeansOnly}{0.21}
\newcommand{\RcPhMeansOnlyLo}{0.19}
\newcommand{\RcPhMeansOnlyHi}{0.24}
\newcommand{\RcPhTotFiftyRandom}{0.22}
\newcommand{\RcPhTotFiftyRandomDiff}{0.009}
\newcommand{\RcPhTotFiftyRandomDiffLo}{0.000}
\newcommand{\RcPhTotFiftyRandomDiffHi}{0.014}
\newcommand{\RcPhTotFiftyGuided}{0.21}
\newcommand{\RcPhTotFiftyGuidedDiff}{\ensuremath{-}0.002}
\newcommand{\RcPhTotFiftyGuidedDiffLo}{\ensuremath{-}0.009}
\newcommand{\RcPhTotFiftyGuidedDiffHi}{0.003}
\newcommand{\RcPhTotFiftyAlone}{0.64}
\newcommand{\RcPhTotFiftyGuidedMinusRandom}{\ensuremath{-}0.010}
\newcommand{\RcPhTotFiftyGuidedMinusRandomLo}{\ensuremath{-}0.014}
\newcommand{\RcPhTotFiftyGuidedMinusRandomHi}{\ensuremath{-}0.006}
\newcommand{\RcPhTotHundredRandom}{0.20}
\newcommand{\RcPhTotHundredRandomDiff}{\ensuremath{-}0.011}
\newcommand{\RcPhTotHundredRandomDiffLo}{\ensuremath{-}0.022}
\newcommand{\RcPhTotHundredRandomDiffHi}{\ensuremath{-}0.004}
\newcommand{\RcPhTotHundredGuided}{0.19}
\newcommand{\RcPhTotHundredGuidedDiff}{\ensuremath{-}0.016}
\newcommand{\RcPhTotHundredGuidedDiffLo}{\ensuremath{-}0.029}
\newcommand{\RcPhTotHundredGuidedDiffHi}{\ensuremath{-}0.008}
\newcommand{\RcPhTotHundredAlone}{0.51}
\newcommand{\RcPhTotHundredGuidedMinusRandom}{\ensuremath{-}0.006}
\newcommand{\RcPhTotHundredGuidedMinusRandomLo}{\ensuremath{-}0.008}
\newcommand{\RcPhTotHundredGuidedMinusRandomHi}{\ensuremath{-}0.004}
\newcommand{\RcPhTotTwoHundredRandom}{0.18}
\newcommand{\RcPhTotTwoHundredRandomDiff}{\ensuremath{-}0.028}
\newcommand{\RcPhTotTwoHundredRandomDiffLo}{\ensuremath{-}0.045}
\newcommand{\RcPhTotTwoHundredRandomDiffHi}{\ensuremath{-}0.017}
\newcommand{\RcPhTotTwoHundredGuided}{0.18}
\newcommand{\RcPhTotTwoHundredGuidedDiff}{\ensuremath{-}0.032}
\newcommand{\RcPhTotTwoHundredGuidedDiffLo}{\ensuremath{-}0.051}
\newcommand{\RcPhTotTwoHundredGuidedDiffHi}{\ensuremath{-}0.020}
\newcommand{\RcPhTotTwoHundredAlone}{0.38}
\newcommand{\RcPhTotTwoHundredGuidedMinusRandom}{\ensuremath{-}0.004}
\newcommand{\RcPhTotTwoHundredGuidedMinusRandomLo}{\ensuremath{-}0.006}
\newcommand{\RcPhTotTwoHundredGuidedMinusRandomHi}{\ensuremath{-}0.003}
\newcommand{\RcPhTotFourHundredRandom}{0.16}
\newcommand{\RcPhTotFourHundredRandomDiff}{\ensuremath{-}0.046}
\newcommand{\RcPhTotFourHundredRandomDiffLo}{\ensuremath{-}0.068}
\newcommand{\RcPhTotFourHundredRandomDiffHi}{\ensuremath{-}0.032}
\newcommand{\RcPhTotFourHundredGuided}{0.16}
\newcommand{\RcPhTotFourHundredGuidedDiff}{\ensuremath{-}0.049}
\newcommand{\RcPhTotFourHundredGuidedDiffLo}{\ensuremath{-}0.072}
\newcommand{\RcPhTotFourHundredGuidedDiffHi}{\ensuremath{-}0.035}
\newcommand{\RcPhTotFourHundredAlone}{0.29}
\newcommand{\RcPhTotFourHundredGuidedMinusRandom}{\ensuremath{-}0.003}
\newcommand{\RcPhTotFourHundredGuidedMinusRandomLo}{\ensuremath{-}0.005}
\newcommand{\RcPhTotFourHundredGuidedMinusRandomHi}{\ensuremath{-}0.002}
\newcommand{\RcPhTotEightHundredRandom}{0.15}
\newcommand{\RcPhTotEightHundredRandomDiff}{\ensuremath{-}0.064}
\newcommand{\RcPhTotEightHundredRandomDiffLo}{\ensuremath{-}0.090}
\newcommand{\RcPhTotEightHundredRandomDiffHi}{\ensuremath{-}0.048}
\newcommand{\RcPhTotEightHundredGuided}{0.14}
\newcommand{\RcPhTotEightHundredGuidedDiff}{\ensuremath{-}0.067}
\newcommand{\RcPhTotEightHundredGuidedDiffLo}{\ensuremath{-}0.093}
\newcommand{\RcPhTotEightHundredGuidedDiffHi}{\ensuremath{-}0.051}
\newcommand{\RcPhTotEightHundredAlone}{0.21}
\newcommand{\RcPhTotEightHundredGuidedMinusRandom}{\ensuremath{-}0.003}
\newcommand{\RcPhTotEightHundredGuidedMinusRandomLo}{\ensuremath{-}0.004}
\newcommand{\RcPhTotEightHundredGuidedMinusRandomHi}{\ensuremath{-}0.002}
\newcommand{\RcPhAloneEightHundredMinusMeans}{\ensuremath{-}0.001}
\newcommand{\RcPhAloneEightHundredMinusMeansLo}{\ensuremath{-}0.033}
\newcommand{\RcPhAloneEightHundredMinusMeansHi}{0.020}
\newcommand{\RcPhSavingFifty}{31}
\newcommand{\RcPhSavingFiftyLo}{25}
\newcommand{\RcPhSavingFiftyHi}{34}
\newcommand{\RcPhRandomCellsFifty}{73}
\newcommand{\RcPhRandomCellsFiftyLo}{66}
\newcommand{\RcPhRandomCellsFiftyHi}{76}
\newcommand{\RcPhSavingHundred}{22}
\newcommand{\RcPhSavingHundredLo}{18}
\newcommand{\RcPhSavingHundredHi}{25}
\newcommand{\RcPhRandomCellsHundred}{128}
\newcommand{\RcPhRandomCellsHundredLo}{122}
\newcommand{\RcPhRandomCellsHundredHi}{133}
\newcommand{\RcPhSavingTwoHundred}{18}
\newcommand{\RcPhSavingTwoHundredLo}{14}
\newcommand{\RcPhSavingTwoHundredHi}{20}
\newcommand{\RcPhRandomCellsTwoHundred}{243}
\newcommand{\RcPhRandomCellsTwoHundredLo}{233}
\newcommand{\RcPhRandomCellsTwoHundredHi}{251}
\newcommand{\RcPhSavingFourHundred}{13}
\newcommand{\RcPhSavingFourHundredLo}{10}
\newcommand{\RcPhSavingFourHundredHi}{17}
\newcommand{\RcPhRandomCellsFourHundred}{462}
\newcommand{\RcPhRandomCellsFourHundredLo}{443}
\newcommand{\RcPhRandomCellsFourHundredHi}{480}
\newcommand{\RcPhIdAll}{0.84}
\newcommand{\RcPhIdAllLo}{0.78}
\newcommand{\RcPhIdAllHi}{0.90}
\newcommand{\RcPhIdInAll}{0.85}
\newcommand{\RcPhIdInAllLo}{0.80}
\newcommand{\RcPhIdInAllHi}{0.90}
\newcommand{\RcPhIdFamAll}{0.52}
\newcommand{\RcPhIdFamAllLo}{0.51}
\newcommand{\RcPhIdFamAllHi}{0.53}
\newcommand{\RcPhIdDiffAll}{0.030}
\newcommand{\RcPhIdDiffAllLo}{0.017}
\newcommand{\RcPhIdDiffAllHi}{0.045}
\newcommand{\RcPhPrimaryAll}{0.87}
\newcommand{\RcPhPrimaryAllLo}{0.81}
\newcommand{\RcPhPrimaryAllHi}{0.92}
\newcommand{\RcPhIdHi}{0.96}
\newcommand{\RcPhIdHiLo}{0.90}
\newcommand{\RcPhIdHiHi}{0.99}
\newcommand{\RcPhIdInHi}{0.96}
\newcommand{\RcPhIdInHiLo}{0.91}
\newcommand{\RcPhIdInHiHi}{0.99}
\newcommand{\RcPhIdFamHi}{0.94}
\newcommand{\RcPhIdFamHiLo}{0.88}
\newcommand{\RcPhIdFamHiHi}{0.96}
\newcommand{\RcPhIdDiffHi}{0.029}
\newcommand{\RcPhIdDiffHiLo}{0.002}
\newcommand{\RcPhIdDiffHiHi}{0.082}
\newcommand{\RcPhPrimaryHi}{0.99}
\newcommand{\RcPhPrimaryHiLo}{0.98}
\newcommand{\RcPhPrimaryHiHi}{1.00}
\newcommand{\RcPhLofoPOneHi}{0.99}
\newcommand{\RcPhLofoPOneHiLo}{0.97}
\newcommand{\RcPhLofoPOneHiHi}{0.99}
\newcommand{\RcPhLofoPOneAll}{0.87}
\newcommand{\RcPhLofoPOneAllLo}{0.80}
\newcommand{\RcPhLofoPOneAllHi}{0.92}
\newcommand{\RcPhLofoPOneFamily}{0.89}
\newcommand{\RcPhLofoPOneFamilyLo}{0.69}
\newcommand{\RcPhLofoPOneFamilyHi}{0.97}
\newcommand{\RcPhPOneWithinPairs}{24}
\newcommand{\RcPhPOneWithinFailures}{18}
\newcommand{\RcPhPOneWithinPrimary}{0.77}
\newcommand{\RcPhPOneWithinPrimaryLo}{0.51}
\newcommand{\RcPhPOneWithinPrimaryHi}{0.88}
\newcommand{\RcPhPOneWithinCoverage}{0.94}
\newcommand{\RcPhPOneWithinCoverageLo}{0.85}
\newcommand{\RcPhPOneWithinCoverageHi}{1.00}
\newcommand{\RcPhFamLcPrimary}{0.88}
\newcommand{\RcPhFamLcCoverage}{0.85}
\newcommand{\RcPhFamFrozenPrimary}{0.84}
\newcommand{\RcPhFamFrozenCoverage}{0.96}
\newcommand{\RcPhFamPOnePrimary}{0.92}
\newcommand{\RcPhFamPOneCoverage}{0.98}
\newcommand{\RcPhFamColonPrimary}{0.80}
\newcommand{\RcPhFamColonCoverage}{0.94}
\newcommand{\RcPhLabRna}{93}
\newcommand{\RcPhLabProtein}{94}
\newcommand{\RcPhLabAgree}{92}
\newcommand{\RcPhLabErr}{0.21}
\newcommand{\RcPhLabErrJoint}{0.21}
\newcommand{\RcPhLabDiff}{\ensuremath{-}0.002}
\newcommand{\RcPhLabDiffLo}{\ensuremath{-}0.009}
\newcommand{\RcPhLabDiffHi}{0.002}
\newcommand{\RcPhLabBetter}{12}
\newcommand{\RcPhLabAucAll}{0.87}
\newcommand{\RcPhLabAucHi}{0.99}
\newcommand{\RcPhLabAucDiff}{0.000}
\newcommand{\RcSimRows}{144}
\newcommand{\RcSimNuBiasMax}{0.02}
\newcommand{\RcSimNuMaeMax}{0.07}
\newcommand{\RcSimNuMaeMin}{0.03}
\newcommand{\RcSimCouplingMax}{6}
\newcommand{\RcSimErrMeasMin}{0.12}
\newcommand{\RcSimErrMeasMax}{1.01}
\newcommand{\RcSimErrCorrMin}{0.06}
\newcommand{\RcSimErrCorrMax}{0.38}
\newcommand{\RcSimCorrRatioMin}{0.22}
\newcommand{\RcSimCorrRatioMax}{0.93}
\newcommand{\RcSimMeasRatioMin}{1.33}
\newcommand{\RcSimMeasRatioMax}{3.03}
\newcommand{\RcSimFrozenRatioMin}{0.22}
\newcommand{\RcSimFrozenRatioMax}{0.63}
\newcommand{\RcSimBracket}{431}
\newcommand{\RcSimBracketTotal}{432}
\newcommand{\RcSimFrozenTrue}{1.51}
\newcommand{\RcSimFrozenMeas}{2.62}
\newcommand{\RcSimFrozenCorr}{0.43}
\newcommand{\RcSimFrozenTrueAbove}{24}
\newcommand{\RcSimFrozenCorrAbove}{2}
\newcommand{\RcSimDOne}{24}
\newcommand{\RcSimOracleMin}{0.96}
\newcommand{\RcSimOracleMax}{0.97}
\newcommand{\RcSimPOneCorrAbove}{141}
\newcommand{\RcSimPOneTrueAbove}{144}

% Generated by extension/perturbation/make_assets.py. Do not edit.
\newcommand{\PtCells}{218,331}
\newcommand{\PtTargets}{248}
\newcommand{\PtHeldTargets}{62}
\newcommand{\PtTrainTargets}{186}
\newcommand{\PtSingleCells}{111,453}
\newcommand{\PtNTCells}{16,319}
\newcommand{\PtGenes}{200}
\newcommand{\PtEncodingGenes}{19}
\newcommand{\PtProteins}{20}
\newcommand{\PtTrainPops}{505}
\newcommand{\PtGroups}{141}
\newcommand{\PtHeldEligible}{51}
\newcommand{\PtDimS}{1}
\newcommand{\PtDimSPooled}{2}
\newcommand{\PtRfPrimary}{6.8}
\newcommand{\PtRfPrimaryLo}{4.1}
\newcommand{\PtRfPrimaryHi}{9.4}
\newcommand{\PtRfPrimaryLatent}{7.0}
\newcommand{\PtRfPrimaryLatentLo}{4.4}
\newcommand{\PtRfPrimaryLatentHi}{9.5}
\newcommand{\PtRfRegression}{9.2}
\newcommand{\PtRfRegressionLo}{7.9}
\newcommand{\PtRfRegressionHi}{10.5}
\newcommand{\PtRfCondMarg}{15}
\newcommand{\PtRfCondMargLo}{13}
\newcommand{\PtRfCondMargHi}{17}
\newcommand{\PtRfPooled}{\ensuremath{-}36}
\newcommand{\PtRfPooledLo}{\ensuremath{-}42}
\newcommand{\PtRfPooledHi}{\ensuremath{-}32}
\newcommand{\PtRfPaired}{37}
\newcommand{\PtRfPairedLo}{32}
\newcommand{\PtRfPairedHi}{42}
\newcommand{\PtRfPairedLatent}{40}
\newcommand{\PtRfPairedLatentLo}{34}
\newcommand{\PtRfPairedLatentHi}{44}
\newcommand{\PtRfRefReg}{53}
\newcommand{\PtRfRefRegLo}{49}
\newcommand{\PtRfRefRegHi}{57}
\newcommand{\PtRfTransCorr}{97}
\newcommand{\PtRfTransCorrLo}{93}
\newcommand{\PtRfTransCorrHi}{101}
\newcommand{\PtRfBench}{24}
\newcommand{\PtRfBenchLo}{20}
\newcommand{\PtRfBenchHi}{28}
\newcommand{\PtRfPseudo}{\ensuremath{-}0.1}
\newcommand{\PtRfPseudoLo}{\ensuremath{-}0.1}
\newcommand{\PtRfPseudoHi}{\ensuremath{-}0.1}
\newcommand{\PtRfDerange}{\ensuremath{-}0.4}
\newcommand{\PtRfDerangeLo}{\ensuremath{-}0.6}
\newcommand{\PtRfDerangeHi}{\ensuremath{-}0.3}
\newcommand{\PtRfPooledPseudo}{\ensuremath{-}46}
\newcommand{\PtRfPooledPseudoLo}{\ensuremath{-}52}
\newcommand{\PtRfPooledPseudoHi}{\ensuremath{-}41}
\newcommand{\PtRfPooledDerange}{\ensuremath{-}75}
\newcommand{\PtRfPooledDerangeLo}{\ensuremath{-}82}
\newcommand{\PtRfPooledDerangeHi}{\ensuremath{-}70}
\newcommand{\PtRfDiffPseudo}{6.9}
\newcommand{\PtRfDiffPseudoLo}{4.2}
\newcommand{\PtRfDiffPseudoHi}{9.5}
\newcommand{\PtRfDiffDerange}{7.2}
\newcommand{\PtRfDiffDerangeLo}{4.6}
\newcommand{\PtRfDiffDerangeHi}{9.7}
\newcommand{\PtRfShare}{18}
\newcommand{\PtRfShareLo}{12}
\newcommand{\PtRfShareHi}{25}
\newcommand{\PtRfDiffPooled}{43}
\newcommand{\PtRfDiffPooledLo}{40}
\newcommand{\PtRfDiffPooledHi}{47}
\newcommand{\PtCondAlonePrimary}{11}
\newcommand{\PtNtAlonePrimary}{15}
\newcommand{\PtCondAlonePaired}{24}
\newcommand{\PtNtAlonePaired}{93}
\newcommand{\PtCondAloneTransCorr}{98}
\newcommand{\PtNtAloneTransCorr}{98}
\newcommand{\PtCondIfngPrimary}{16}
\newcommand{\PtNtIfngPrimary}{28}
\newcommand{\PtCondIfngPaired}{43}
\newcommand{\PtNtIfngPaired}{97}
\newcommand{\PtCondIfngTransCorr}{97}
\newcommand{\PtNtIfngTransCorr}{99}
\newcommand{\PtCondCocultPrimary}{\ensuremath{-}5.4}
\newcommand{\PtNtCocultPrimary}{13}
\newcommand{\PtCondCocultPaired}{43}
\newcommand{\PtNtCocultPaired}{94}
\newcommand{\PtCondCocultTransCorr}{96}
\newcommand{\PtNtCocultTransCorr}{98}
\newcommand{\PtNtPrimary}{19}
\newcommand{\PtNtPrimaryLo}{16}
\newcommand{\PtNtPrimaryHi}{21}
\newcommand{\PtNtPaired}{94}
\newcommand{\PtNtPairedLo}{93}
\newcommand{\PtNtPairedHi}{96}
\newcommand{\PtNtPooled}{\ensuremath{-}12}
\newcommand{\PtNtPooledLo}{\ensuremath{-}16}
\newcommand{\PtNtPooledHi}{\ensuremath{-}8}
\newcommand{\PtNtTransCorr}{98}
\newcommand{\PtNtTransCorrLo}{97}
\newcommand{\PtNtTransCorrHi}{100}
\newcommand{\PtNtRefReg}{91}
\newcommand{\PtNtRefRegLo}{90}
\newcommand{\PtNtRefRegHi}{93}
\newcommand{\PtNtBench}{90}
\newcommand{\PtNtBenchLo}{88}
\newcommand{\PtNtBenchHi}{91}
\newcommand{\PtNtShare}{20}
\newcommand{\PtNtShareLo}{17}
\newcommand{\PtNtShareHi}{22}
\newcommand{\PtCoverage}{2.4}
\newcommand{\PtCoverageNt}{3.6}
\newcommand{\PtAdaptMedian}{88}
\newcommand{\PtAdaptNtMin}{1,989}
\newcommand{\PtAdaptNtMax}{3,222}
\newcommand{\PtDevPrimary}{10}
\newcommand{\PtDevPooled}{\ensuremath{-}33}
\newcommand{\PtDevPaired}{38}
\newcommand{\PtHOne}{met}
\newcommand{\PtHThree}{met}
\newcommand{\PtHFour}{not met}
\newcommand{\PtHTwo}{met}
\newcommand{\PgIfnGenes}{34}
\newcommand{\PgIfnProteinCount}{3}
\newcommand{\PgIfnPf}{43}
\newcommand{\PgIfnPfLo}{32}
\newcommand{\PgIfnPfHi}{51}
\newcommand{\PgIfnPaired}{64}
\newcommand{\PgIfnPairedLo}{58}
\newcommand{\PgIfnPairedHi}{70}
\newcommand{\PgIfnTrans}{94}
\newcommand{\PgIfnTransLo}{88}
\newcommand{\PgIfnTransHi}{102}
\newcommand{\PgIfnRefReg}{69}
\newcommand{\PgIfnRefRegLo}{63}
\newcommand{\PgIfnRefRegHi}{74}
\newcommand{\PgIfnBench}{48}
\newcommand{\PgIfnBenchLo}{44}
\newcommand{\PgIfnBenchHi}{53}
\newcommand{\PgIfnRepl}{\ensuremath{-}24}
\newcommand{\PgIfnReplLo}{\ensuremath{-}40}
\newcommand{\PgIfnReplHi}{\ensuremath{-}10}
\newcommand{\PgIfnShare}{66}
\newcommand{\PgIfnShareLo}{54}
\newcommand{\PgIfnShareHi}{77}
\newcommand{\PgIfnShareTrans}{45}
\newcommand{\PgIfnShareTransLo}{35}
\newcommand{\PgIfnShareTransHi}{54}
\newcommand{\PgAlongPf}{27}
\newcommand{\PgAlongPfLo}{\ensuremath{-}1}
\newcommand{\PgAlongPfHi}{50}
\newcommand{\PgAlongPaired}{69}
\newcommand{\PgAlongPairedLo}{56}
\newcommand{\PgAlongPairedHi}{84}
\newcommand{\PgAlongShare}{39}
\newcommand{\PgAlongShareLo}{\ensuremath{-}1}
\newcommand{\PgAlongShareHi}{65}
\newcommand{\PgCompPf}{5.7}
\newcommand{\PgCompPfLo}{3.9}
\newcommand{\PgCompPfHi}{7.3}
\newcommand{\PgCompPaired}{36}
\newcommand{\PgCompPairedLo}{29}
\newcommand{\PgCompPairedHi}{43}
\newcommand{\PgCompTrans}{98}
\newcommand{\PgCompTransLo}{92}
\newcommand{\PgCompTransHi}{104}
\newcommand{\PgDiff}{37}
\newcommand{\PgDiffLo}{26}
\newcommand{\PgDiffHi}{46}
\newcommand{\PgIfnGenesPf}{16}
\newcommand{\PgIfnGenesPfLo}{6}
\newcommand{\PgIfnGenesPfHi}{26}
\newcommand{\PgIfnGenesPaired}{33}
\newcommand{\PgIfnGenesPairedLo}{26}
\newcommand{\PgIfnGenesPairedHi}{41}
\newcommand{\PgIfnRel}{0.29}
\newcommand{\PgCompRel}{0.11}
\newcommand{\PgIfnAlonePf}{19}
\newcommand{\PgIfnAlonePaired}{16}
\newcommand{\PgIfnIfngPf}{65}
\newcommand{\PgIfnIfngPaired}{79}
\newcommand{\PgIfnCocultPf}{19}
\newcommand{\PgIfnCocultPaired}{62}
\newcommand{\PgNtIfnPf}{61}
\newcommand{\PgNtCompPf}{12}
\newcommand{\PgNtIfnPaired}{93}
\newcommand{\PgNtCompPaired}{95}
\newcommand{\PgNtIfnTrans}{98}
\newcommand{\PgNtCompTrans}{99}
\newcommand{\PgNtIfnRepl}{94}
\newcommand{\PgNtCompRepl}{88}
\newcommand{\PgNtIfnRel}{0.95}
\newcommand{\PgOtherMax}{7.8}
\newcommand{\PgOtherCount}{5}
\newcommand{\PgSupportIfn}{0.094}
\newcommand{\PgSupportOtherMax}{0.021}
\newcommand{\PgPOne}{met}
\newcommand{\PgPTwo}{not met}
\newcommand{\PgPThree}{met}
\newcommand{\DsNoneAll}{6.8}
\newcommand{\DsNoneIfn}{43}
\newcommand{\DsFiveAllExp}{17}
\newcommand{\DsFiveAllStr}{3.7}
\newcommand{\DsFiveAllRand}{1.3}
\newcommand{\DsFiveAllRandLow}{\ensuremath{-}24}
\newcommand{\DsFiveAllDiff}{\ensuremath{-}15}
\newcommand{\DsFiveAllDiffLo}{\ensuremath{-}17}
\newcommand{\DsFiveAllDiffHi}{\ensuremath{-}13}
\newcommand{\DsFiveAllBelow}{50}
\newcommand{\DsFiveAllDiffStr}{\ensuremath{-}2.4}
\newcommand{\DsFiveAllDiffStrLo}{\ensuremath{-}5.3}
\newcommand{\DsFiveAllDiffStrHi}{0.1}
\newcommand{\DsFiveIfnExp}{41}
\newcommand{\DsFiveIfnStr}{7.9}
\newcommand{\DsFiveIfnRand}{36}
\newcommand{\DsFiveIfnRandLow}{6.0}
\newcommand{\DsFiveIfnDiff}{\ensuremath{-}4.1}
\newcommand{\DsFiveIfnDiffLo}{\ensuremath{-}12.8}
\newcommand{\DsFiveIfnDiffHi}{3.4}
\newcommand{\DsFiveIfnBelow}{9}
\newcommand{\DsFiveIfnDiffStr}{29}
\newcommand{\DsFiveIfnDiffStrLo}{18}
\newcommand{\DsFiveIfnDiffStrHi}{38}
\newcommand{\DsFiveDimExp}{0}
\newcommand{\DsFiveDimRand}{2.5}
\newcommand{\DsFiveZeroRand}{0}
\newcommand{\DsTenAllExp}{3.5}
\newcommand{\DsTenAllStr}{3.2}
\newcommand{\DsTenAllRand}{2.3}
\newcommand{\DsTenAllRandLow}{\ensuremath{-}24}
\newcommand{\DsTenAllDiff}{\ensuremath{-}1.2}
\newcommand{\DsTenAllDiffLo}{\ensuremath{-}4.2}
\newcommand{\DsTenAllDiffHi}{1.3}
\newcommand{\DsTenAllBelow}{7}
\newcommand{\DsTenAllDiffStr}{\ensuremath{-}0.9}
\newcommand{\DsTenAllDiffStrLo}{\ensuremath{-}4.0}
\newcommand{\DsTenAllDiffStrHi}{1.6}
\newcommand{\DsTenIfnExp}{8.0}
\newcommand{\DsTenIfnStr}{4.7}
\newcommand{\DsTenIfnRand}{38}
\newcommand{\DsTenIfnRandLow}{7.6}
\newcommand{\DsTenIfnDiff}{30}
\newcommand{\DsTenIfnDiffLo}{19}
\newcommand{\DsTenIfnDiffHi}{39}
\newcommand{\DsTenIfnBelow}{3}
\newcommand{\DsTenIfnDiffStr}{33}
\newcommand{\DsTenIfnDiffStrLo}{22}
\newcommand{\DsTenIfnDiffStrHi}{42}
\newcommand{\DsTenDimExp}{0}
\newcommand{\DsTenDimRand}{2.1}
\newcommand{\DsTenZeroRand}{0}
\newcommand{\DsTwentyAllExp}{3.5}
\newcommand{\DsTwentyAllStr}{3.4}
\newcommand{\DsTwentyAllRand}{2.5}
\newcommand{\DsTwentyAllRandLow}{\ensuremath{-}26}
\newcommand{\DsTwentyAllDiff}{\ensuremath{-}1.0}
\newcommand{\DsTwentyAllDiffLo}{\ensuremath{-}3.8}
\newcommand{\DsTwentyAllDiffHi}{1.5}
\newcommand{\DsTwentyAllBelow}{7}
\newcommand{\DsTwentyAllDiffStr}{\ensuremath{-}0.9}
\newcommand{\DsTwentyAllDiffStrLo}{\ensuremath{-}3.7}
\newcommand{\DsTwentyAllDiffStrHi}{1.7}
\newcommand{\DsTwentyIfnExp}{8.4}
\newcommand{\DsTwentyIfnStr}{5.1}
\newcommand{\DsTwentyIfnRand}{38}
\newcommand{\DsTwentyIfnRandLow}{5.2}
\newcommand{\DsTwentyIfnDiff}{29}
\newcommand{\DsTwentyIfnDiffLo}{19}
\newcommand{\DsTwentyIfnDiffHi}{39}
\newcommand{\DsTwentyIfnBelow}{5}
\newcommand{\DsTwentyIfnDiffStr}{33}
\newcommand{\DsTwentyIfnDiffStrLo}{22}
\newcommand{\DsTwentyIfnDiffStrHi}{42}
\newcommand{\DsTwentyDimExp}{0}
\newcommand{\DsTwentyDimRand}{2.0}
\newcommand{\DsTwentyZeroRand}{0}
\newcommand{\DsTopFour}{JAK2, IFNGR1, JAK1 and IFNGR2}
\newcommand{\DsTopFive}{JAK2, IFNGR1, JAK1, IFNGR2 and TMED10}
\newcommand{\DsTopStrengthFive}{JAK2, JAK1, IFNGR1, IFNGR2 and CTPS1}
\newcommand{\DsExposureGap}{6}
\newcommand{\DsLooSpearman}{\ensuremath{-}0.09}
\newcommand{\DsLooP}{0.87}
\newcommand{\DsLooTop}{JAK2 and JAK1}
\newcommand{\DsDraws}{50}
\newcommand{\DsDOne}{not met}
\newcommand{\PtTopGenes}{B2M, HLA-B, CD74, GBP1, HLA-DRA, HLA-A}

% Generated by extension/perturbation_replication/make_assets.py. Do not edit.
\newcommand{\RpCells}{20,652}
\newcommand{\RpTargets}{25}
\newcommand{\RpGroups}{55}
\newcommand{\RpTestTargets}{21}
\newcommand{\RpReplicates}{3}
\newcommand{\RpGenes}{200}
\newcommand{\RpProteins}{4}
\newcommand{\RpIfnGenes}{16}
\newcommand{\RpDimSMedian}{18}
\newcommand{\RpDimSMin}{18}
\newcommand{\RpDimSMax}{61}
\newcommand{\RpPrimary}{\ensuremath{-}212}
\newcommand{\RpPrimaryLo}{\ensuremath{-}287}
\newcommand{\RpPrimaryHi}{\ensuremath{-}157}
\newcommand{\RpDiffPseudo}{\ensuremath{-}202}
\newcommand{\RpDiffPseudoLo}{\ensuremath{-}275}
\newcommand{\RpDiffPseudoHi}{\ensuremath{-}151}
\newcommand{\RpDiffDerange}{\ensuremath{-}167}
\newcommand{\RpDiffDerangeLo}{\ensuremath{-}234}
\newcommand{\RpDiffDerangeHi}{\ensuremath{-}119}
\newcommand{\RpIfnShare}{\ensuremath{-}3.5}
\newcommand{\RpIfnShareLo}{\ensuremath{-}41.3}
\newcommand{\RpIfnShareHi}{24.3}
\newcommand{\RpRemDiff}{\ensuremath{-}44}
\newcommand{\RpRemDiffLo}{\ensuremath{-}91}
\newcommand{\RpRemDiffHi}{\ensuremath{-}3}
\newcommand{\RpPaired}{32}
\newcommand{\RpPairedLo}{21}
\newcommand{\RpPairedHi}{43}
\newcommand{\RpIfnPrimary}{\ensuremath{-}3.0}
\newcommand{\RpIfnPrimaryLo}{\ensuremath{-}35.6}
\newcommand{\RpIfnPrimaryHi}{19.5}
\newcommand{\RpIfnPaired}{84}
\newcommand{\RpIfnPairedLo}{73}
\newcommand{\RpIfnPairedHi}{100}
\newcommand{\RpShare}{\ensuremath{-}656}
\newcommand{\RpShareLo}{\ensuremath{-}1195}
\newcommand{\RpShareHi}{\ensuremath{-}429}
\newcommand{\RpIfnRemDiff}{1.9}
\newcommand{\RpIfnRemDiffLo}{\ensuremath{-}21.1}
\newcommand{\RpIfnRemDiffHi}{20.2}
\newcommand{\RpAlongPrimary}{\ensuremath{-}124}
\newcommand{\RpAlongPrimaryLo}{\ensuremath{-}188}
\newcommand{\RpAlongPrimaryHi}{\ensuremath{-}75}
\newcommand{\RpAlongShare}{\ensuremath{-}171}
\newcommand{\RpAlongShareLo}{\ensuremath{-}266}
\newcommand{\RpAlongShareHi}{\ensuremath{-}107}
\newcommand{\RpIfnTrans}{94}
\newcommand{\RpIfnTransLo}{80}
\newcommand{\RpIfnTransHi}{115}
\newcommand{\RpIfnReplicate}{19}
\newcommand{\RpIfnReplicateLo}{\ensuremath{-}2}
\newcommand{\RpIfnReplicateHi}{36}
\newcommand{\RpAllPf}{\ensuremath{-}212}
\newcommand{\RpAllPaired}{32}
\newcommand{\RpAllTrans}{81}
\newcommand{\RpAllRefReg}{37}
\newcommand{\RpAllBench}{22}
\newcommand{\RpAllRepl}{\ensuremath{-}205}
\newcommand{\RpAllPseudo}{\ensuremath{-}10}
\newcommand{\RpAllDerange}{\ensuremath{-}46}
\newcommand{\RpAllRemExp}{\ensuremath{-}172}
\newcommand{\RpAllRemStr}{\ensuremath{-}219}
\newcommand{\RpAllRemRand}{\ensuremath{-}216}
\newcommand{\RpAllRel}{0.14}
\newcommand{\RpIfnPf}{\ensuremath{-}3.0}
\newcommand{\RpIfnRefReg}{84}
\newcommand{\RpIfnBench}{61}
\newcommand{\RpIfnRepl}{19}
\newcommand{\RpIfnPseudo}{\ensuremath{-}3.2}
\newcommand{\RpIfnDerange}{\ensuremath{-}12}
\newcommand{\RpIfnRemExp}{\ensuremath{-}5.6}
\newcommand{\RpIfnRemStr}{\ensuremath{-}23}
\newcommand{\RpIfnRemRand}{\ensuremath{-}3.7}
\newcommand{\RpIfnRel}{0.34}
\newcommand{\RpCompPf}{\ensuremath{-}201}
\newcommand{\RpCompPaired}{25}
\newcommand{\RpCompTrans}{84}
\newcommand{\RpCompRefReg}{31}
\newcommand{\RpCompBench}{3.1}
\newcommand{\RpCompRepl}{\ensuremath{-}390}
\newcommand{\RpCompPseudo}{\ensuremath{-}16}
\newcommand{\RpCompDerange}{\ensuremath{-}54}
\newcommand{\RpCompRemExp}{\ensuremath{-}164}
\newcommand{\RpCompRemStr}{\ensuremath{-}173}
\newcommand{\RpCompRemRand}{\ensuremath{-}227}
\newcommand{\RpCompRel}{0.09}
\newcommand{\RpAlongPf}{\ensuremath{-}124}
\newcommand{\RpAlongPaired}{72}
\newcommand{\RpAlongTrans}{87}
\newcommand{\RpAlongRefReg}{73}
\newcommand{\RpAlongBench}{50}
\newcommand{\RpAlongRepl}{\ensuremath{-}5.0}
\newcommand{\RpAlongPseudo}{\ensuremath{-}8.1}
\newcommand{\RpAlongDerange}{\ensuremath{-}37}
\newcommand{\RpAlongRemExp}{\ensuremath{-}114}
\newcommand{\RpAlongRemStr}{\ensuremath{-}155}
\newcommand{\RpAlongRemRand}{\ensuremath{-}123}
\newcommand{\RpAlongRel}{0.30}
\newcommand{\RpRepOnePf}{\ensuremath{-}98}
\newcommand{\RpRepOneIfnPf}{25}
\newcommand{\RpRepOneIfnPaired}{87}
\newcommand{\RpRepThreePf}{\ensuremath{-}298}
\newcommand{\RpRepThreeIfnPf}{8.6}
\newcommand{\RpRepThreeIfnPaired}{74}
\newcommand{\RpRepFourPf}{\ensuremath{-}393}
\newcommand{\RpRepFourIfnPf}{\ensuremath{-}129}
\newcommand{\RpRepFourIfnPaired}{93}
\newcommand{\RpNtPrimary}{\ensuremath{-}90}
\newcommand{\RpNtPrimaryLo}{\ensuremath{-}140}
\newcommand{\RpNtPrimaryHi}{\ensuremath{-}52}
\newcommand{\RpNtPaired}{82}
\newcommand{\RpNtIfnPf}{21}
\newcommand{\RpNtIfnPaired}{89}
\newcommand{\RpNtIfnRepl}{88}
\newcommand{\RpTopRemoved}{CUL3, STAT1, SMAD4, IFNGR2, JAK2}
\newcommand{\RpFixedPointOneAll}{\ensuremath{-}76}
\newcommand{\RpFixedPointOneIfn}{32}
\newcommand{\RpFixedPointOneDim}{1}
\newcommand{\RpFixedOneAll}{4.6}
\newcommand{\RpFixedOneIfn}{45}
\newcommand{\RpFixedOneDim}{0}
\newcommand{\RpROne}{not met}
\newcommand{\RpRThree}{not met}
\newcommand{\RpRFour}{not met}
\newcommand{\RpRFive}{not met}
\newcommand{\RpRTwo}{not met}

% Generated by extension/semipaired/make_assets.py. Do not edit.
\newcommand{\HyFrBlockHundred}{38}
\newcommand{\HyFrCurrent}{3.6}
\newcommand{\HyFrHybridHundred}{35}
\newcommand{\HyFrHybridPcHundred}{36}
\newcommand{\HyFrNoiseAware}{\ensuremath{-}12}
\newcommand{\HyFrPairedJsHundred}{16}
\newcommand{\HyPaBlockHundred}{40}
\newcommand{\HyPaCurrent}{\ensuremath{-}154}
\newcommand{\HyPaHybridHundred}{25}
\newcommand{\HyPaHybridPcHundred}{38}
\newcommand{\HyPaNoiseAware}{\ensuremath{-}314}
\newcommand{\HyPaPairedJsHundred}{22}
\newcommand{\SpBoundBelowIdealized}{yes}
\newcommand{\SpBoundBudgetMaxFr}{800}
\newcommand{\SpBoundBudgetMaxOc}{400}
\newcommand{\SpBoundBudgetMaxPa}{800}
\newcommand{\SpBoundCellsFr}{21,171}
\newcommand{\SpBoundCellsOc}{457}
\newcommand{\SpBoundCellsPa}{2,256}
\newcommand{\SpBoundConditionFr}{yes}
\newcommand{\SpBoundConditionOc}{yes}
\newcommand{\SpBoundConditionPa}{yes}
\newcommand{\SpBoundCoveredFr}{yes}
\newcommand{\SpBoundCoveredOc}{yes}
\newcommand{\SpBoundCoveredPa}{yes}
\newcommand{\SpBoundDraws}{200}
\newcommand{\SpBoundEffDimMin}{4.4}
\newcommand{\SpBoundRatioMax}{3.3}
\newcommand{\SpBoundRatioMin}{1.3}
\newcommand{\SpDevChampText}{In Frangieh, Champollion recovered 40\% at 100 paired cells against 48\% for the proposed estimator, and the proposed estimator needed 1.4 to 1.7 times fewer paired cells to reach recovered fractions of 0.3 to 0.6 (lowest 95\% bound, 1.4). In Papalexi (7 held-out targets, 21 groups), Champollion recovered 41\% at 100 paired cells against 29\% for the proposed estimator, so it reached recovered fractions of 0.3 and 0.4 with fewer paired cells (ratios 0.55 and 0.69; 95\% CI, 0.35--0.79 and 0.35--1.00), and the two were similar at 0.5 and 0.6 (0.91 and 0.97). Champollion's lasso weight was chosen per budget on the evaluation groups, which favours it.}
\newcommand{\SpEffDimBudgetFr}{1,600}
\newcommand{\SpEffDimBudgetPa}{1,600}
\newcommand{\SpEffDimMin}{4.2}
\newcommand{\SpEffDimSmallMin}{3.3}
\newcommand{\SpExBestDiffHundred}{8.1}
\newcommand{\SpExBestDiffHundredHi}{13}
\newcommand{\SpExBestDiffHundredLo}{2.4}
\newcommand{\SpExBestHundred}{ridge\_p/mapped}
\newcommand{\SpExCells}{908}
\newcommand{\SpExCellsAll}{111}
\newcommand{\SpExCellsAllFour}{186}
\newcommand{\SpExCellsAllTwo}{58}
\newcommand{\SpExCellsChamp}{131}
\newcommand{\SpExCellsChampFour}{190}
\newcommand{\SpExCellsChampTwo}{87}
\newcommand{\SpExCellsOriginal}{115}
\newcommand{\SpExCellsOriginalFour}{187}
\newcommand{\SpExCellsOriginalTwo}{61}
\newcommand{\SpExCellsPairedOnly}{357}
\newcommand{\SpExCellsPairedOnlyFour}{400}
\newcommand{\SpExCellsPairedOnlyTwo}{196}
\newcommand{\SpExCellsProposed}{73}
\newcommand{\SpExCellsProposedRel}{}
\newcommand{\SpExCellsSemi}{122}
\newcommand{\SpExCellsSemiFour}{236}
\newcommand{\SpExCellsSemiTwo}{64}
\newcommand{\SpExChampFifty}{8.6}
\newcommand{\SpExChampFourHundred}{54}
\newcommand{\SpExChampHundred}{23}
\newcommand{\SpExChampTwentyFive}{0.9}
\newcommand{\SpExChampTwoHundred}{41}
\newcommand{\SpExCrossPairedOnlyHiBudget}{400}
\newcommand{\SpExCrossPairedOnlyHiRf}{32}
\newcommand{\SpExCrossPairedOnlyLoBudget}{200}
\newcommand{\SpExCrossPairedOnlyLoRf}{20}
\newcommand{\SpExCrossProposedHiBudget}{100}
\newcommand{\SpExCrossProposedHiRf}{36}
\newcommand{\SpExCrossProposedLoBudget}{50}
\newcommand{\SpExCrossProposedLoRf}{23}
\newcommand{\SpExFrozenArmAllMet}{yes}
\newcommand{\SpExFrozenArmCells}{101}
\newcommand{\SpExFrozenArmChamp}{1.3}
\newcommand{\SpExFrozenArmChampHi}{1.5}
\newcommand{\SpExFrozenArmChampLo}{0.96}
\newcommand{\SpExFrozenArmPairedOnly}{3.5}
\newcommand{\SpExFrozenArmPairedOnlyHi}{4.1}
\newcommand{\SpExFrozenArmPairedOnlyLo}{2.6}
\newcommand{\SpExFrozenArmSemi}{1.2}
\newcommand{\SpExFrozenArmSemiHi}{1.3}
\newcommand{\SpExFrozenArmSemiLo}{1.09}
\newcommand{\SpExFrozenAt}{14:49 EDT on 27 September 2026}
\newcommand{\SpExHOne}{met}
\newcommand{\SpExHThree}{met}
\newcommand{\SpExHTwo}{met}
\newcommand{\SpExLossBestOtherHundred}{0.85}
\newcommand{\SpExLossChampHundred}{0.87}
\newcommand{\SpExLossFifty}{0.87}
\newcommand{\SpExLossFourHundred}{0.68}
\newcommand{\SpExLossHundred}{0.80}
\newcommand{\SpExLossIndependence}{1.00}
\newcommand{\SpExLossPairedOnlyHundred}{0.93}
\newcommand{\SpExLossSemiHundred}{0.85}
\newcommand{\SpExLossTwentyFive}{0.97}
\newcommand{\SpExLossTwoHundred}{0.74}
\newcommand{\SpExOrfs}{11}
\newcommand{\SpExPairedAll}{64}
\newcommand{\SpExPairedOnlyFifty}{5.7}
\newcommand{\SpExPairedOnlyFourHundred}{32}
\newcommand{\SpExPairedOnlyHundred}{12}
\newcommand{\SpExPairedOnlyTwentyFive}{2.6}
\newcommand{\SpExPairedOnlyTwoHundred}{20}
\newcommand{\SpExPoolPops}{28}
\newcommand{\SpExReservoir}{457}
\newcommand{\SpExRfFifty}{23}
\newcommand{\SpExRfFiftyHi}{28}
\newcommand{\SpExRfFiftyLo}{15}
\newcommand{\SpExRfFourHundred}{57}
\newcommand{\SpExRfFourHundredHi}{68}
\newcommand{\SpExRfFourHundredLo}{47}
\newcommand{\SpExRfHundred}{36}
\newcommand{\SpExRfHundredHi}{44}
\newcommand{\SpExRfHundredLo}{27}
\newcommand{\SpExRfTwentyFive}{5.2}
\newcommand{\SpExRfTwentyFiveHi}{11}
\newcommand{\SpExRfTwentyFiveLo}{\ensuremath{-}6.6}
\newcommand{\SpExRfTwoHundred}{47}
\newcommand{\SpExRfTwoHundredHi}{56}
\newcommand{\SpExRfTwoHundredLo}{38}
\newcommand{\SpExSaveAll}{1.5}
\newcommand{\SpExSaveAllFour}{1.4}
\newcommand{\SpExSaveAllFourHi}{1.9}
\newcommand{\SpExSaveAllFourLo}{1.02}
\newcommand{\SpExSaveAllFourRel}{}
\newcommand{\SpExSaveAllHi}{1.8}
\newcommand{\SpExSaveAllLo}{1.11}
\newcommand{\SpExSaveAllRel}{}
\newcommand{\SpExSaveAllTwo}{1.3}
\newcommand{\SpExSaveAllTwoHi}{1.5}
\newcommand{\SpExSaveAllTwoLo}{1.02}
\newcommand{\SpExSaveAllTwoRel}{}
\newcommand{\SpExSaveChamp}{1.8}
\newcommand{\SpExSaveChampFour}{1.5}
\newcommand{\SpExSaveChampFourHi}{2.0}
\newcommand{\SpExSaveChampFourLo}{1.02}
\newcommand{\SpExSaveChampFourRel}{}
\newcommand{\SpExSaveChampHi}{2.2}
\newcommand{\SpExSaveChampLo}{1.2}
\newcommand{\SpExSaveChampRel}{}
\newcommand{\SpExSaveChampTwo}{1.9}
\newcommand{\SpExSaveChampTwoHi}{2.4}
\newcommand{\SpExSaveChampTwoLo}{1.4}
\newcommand{\SpExSaveChampTwoRel}{}
\newcommand{\SpExSaveOriginal}{1.6}
\newcommand{\SpExSaveOriginalFour}{1.5}
\newcommand{\SpExSaveOriginalFourHi}{1.9}
\newcommand{\SpExSaveOriginalFourLo}{1.02}
\newcommand{\SpExSaveOriginalFourRel}{}
\newcommand{\SpExSaveOriginalHi}{1.9}
\newcommand{\SpExSaveOriginalLo}{1.12}
\newcommand{\SpExSaveOriginalRel}{}
\newcommand{\SpExSaveOriginalTwo}{1.4}
\newcommand{\SpExSaveOriginalTwoHi}{1.6}
\newcommand{\SpExSaveOriginalTwoLo}{1.03}
\newcommand{\SpExSaveOriginalTwoRel}{}
\newcommand{\SpExSavePairedOnly}{4.9}
\newcommand{\SpExSavePairedOnlyFour}{3.1}
\newcommand{\SpExSavePairedOnlyFourHi}{4.7}
\newcommand{\SpExSavePairedOnlyFourLo}{1.7}
\newcommand{\SpExSavePairedOnlyFourRel}{more than }
\newcommand{\SpExSavePairedOnlyHi}{6.0}
\newcommand{\SpExSavePairedOnlyLo}{3.2}
\newcommand{\SpExSavePairedOnlyRel}{}
\newcommand{\SpExSavePairedOnlyTwo}{4.4}
\newcommand{\SpExSavePairedOnlyTwoHi}{5.2}
\newcommand{\SpExSavePairedOnlyTwoLo}{3.2}
\newcommand{\SpExSavePairedOnlyTwoRel}{}
\newcommand{\SpExSaveSemi}{1.7}
\newcommand{\SpExSaveSemiFour}{1.8}
\newcommand{\SpExSaveSemiFourHi}{2.5}
\newcommand{\SpExSaveSemiFourLo}{1.5}
\newcommand{\SpExSaveSemiFourRel}{}
\newcommand{\SpExSaveSemiHi}{2.0}
\newcommand{\SpExSaveSemiLo}{1.3}
\newcommand{\SpExSaveSemiRel}{}
\newcommand{\SpExSaveSemiTwo}{1.4}
\newcommand{\SpExSaveSemiTwoHi}{1.7}
\newcommand{\SpExSaveSemiTwoLo}{1.2}
\newcommand{\SpExSaveSemiTwoRel}{}
\newcommand{\SpExSecondsChamp}{8.7}
\newcommand{\SpExSecondsChampFit}{5.5}
\newcommand{\SpExSecondsProposed}{14}
\newcommand{\SpExSemiFifty}{16}
\newcommand{\SpExSemiFourHundred}{47}
\newcommand{\SpExSemiHundred}{27}
\newcommand{\SpExSemiTwentyFive}{12}
\newcommand{\SpExSemiTwoHundred}{38}
\newcommand{\SpExTarget}{0.3}
\newcommand{\SpFixedCriterion}{60.8}
\newcommand{\SpFixedFrFifty}{37}
\newcommand{\SpFixedPaFifty}{16}
\newcommand{\SpFrBestDiffHundred}{16}
\newcommand{\SpFrBestDiffHundredHi}{17}
\newcommand{\SpFrBestDiffHundredLo}{15}
\newcommand{\SpFrBestWins}{7}
\newcommand{\SpFrBudgets}{7}
\newcommand{\SpFrCellsFour}{59}
\newcommand{\SpFrChampBudgets}{4}
\newcommand{\SpFrChampFolds}{1}
\newcommand{\SpFrChampGroups}{141}
\newcommand{\SpFrChampHundred}{40}
\newcommand{\SpFrChampPropHundred}{48}
\newcommand{\SpFrChampSaveFour}{1.7}
\newcommand{\SpFrChampSaveFourHi}{1.9}
\newcommand{\SpFrChampSaveFourLo}{1.5}
\newcommand{\SpFrChampSaveMax}{1.7}
\newcommand{\SpFrChampSaveMin}{1.4}
\newcommand{\SpFrChampWins}{4}
\newcommand{\SpFrDoseEightThousandHundred}{47}
\newcommand{\SpFrDoseFiveHundredHundred}{42}
\newcommand{\SpFrDoseTwoThousandHundred}{45}
\newcommand{\SpFrGroups}{141}
\newcommand{\SpFrLossLowestEverywhere}{yes}
\newcommand{\SpFrMapGain}{3.5}
\newcommand{\SpFrMapMeasuredHundred}{43}
\newcommand{\SpFrMapPairedHundred}{\ensuremath{-}398}
\newcommand{\SpFrMarkerHundred}{15}
\newcommand{\SpFrOneSidedHundred}{39}
\newcommand{\SpFrPairBasisHundred}{11}
\newcommand{\SpFrPairedAll}{97}
\newcommand{\SpFrPairedOnlyCellsFour}{347}
\newcommand{\SpFrPairedOnlyFourHundred}{43}
\newcommand{\SpFrPairedOnlyHundred}{16}
\newcommand{\SpFrPercondHundred}{26}
\newcommand{\SpFrPooledHundred}{45}
\newcommand{\SpFrRfFifty}{37}
\newcommand{\SpFrRfFourHundred}{72}
\newcommand{\SpFrRfHundred}{48}
\newcommand{\SpFrRidgeHundred}{27}
\newcommand{\SpFrSaveBestCiMin}{1.7}
\newcommand{\SpFrSaveBestFour}{2.4}
\newcommand{\SpFrSaveBestFourHi}{2.6}
\newcommand{\SpFrSaveBestFourLo}{2.2}
\newcommand{\SpFrSaveBestMax}{2.4}
\newcommand{\SpFrSaveBestMin}{1.8}
\newcommand{\SpFrSavePairedOnlyCiMin}{4.1}
\newcommand{\SpFrSavePairedOnlyFour}{5.9}
\newcommand{\SpFrSavePairedOnlyFourHi}{6.5}
\newcommand{\SpFrSavePairedOnlyFourLo}{5.3}
\newcommand{\SpFrSavePairedOnlyMax}{5.9}
\newcommand{\SpFrSavePairedOnlyMin}{4.4}
\newcommand{\SpFrSavePublishedCiMin}{1.8}
\newcommand{\SpFrSavePublishedFour}{2.7}
\newcommand{\SpFrSavePublishedFourHi}{2.9}
\newcommand{\SpFrSavePublishedFourLo}{2.4}
\newcommand{\SpFrSavePublishedMax}{2.9}
\newcommand{\SpFrSavePublishedMin}{1.9}
\newcommand{\SpFrSemiccaHundred}{31}
\newcommand{\SpFrTargets}{51}
\newcommand{\SpPaBestDiffHundred}{4.6}
\newcommand{\SpPaBestDiffHundredHi}{11}
\newcommand{\SpPaBestDiffHundredLo}{\ensuremath{-}2.2}
\newcommand{\SpPaBestWins}{0}
\newcommand{\SpPaBudgets}{6}
\newcommand{\SpPaCellsFour}{112}
\newcommand{\SpPaChampBudgets}{4}
\newcommand{\SpPaChampFolds}{7}
\newcommand{\SpPaChampGroups}{21}
\newcommand{\SpPaChampHundred}{41}
\newcommand{\SpPaChampPropHundred}{29}
\newcommand{\SpPaChampSaveFour}{0.69}
\newcommand{\SpPaChampSaveFourHi}{1.00}
\newcommand{\SpPaChampSaveFourLo}{0.35}
\newcommand{\SpPaChampSaveMax}{0.97}
\newcommand{\SpPaChampSaveMin}{0.55}
\newcommand{\SpPaChampWins}{0}
\newcommand{\SpPaDoseEightThousandHundred}{37}
\newcommand{\SpPaDoseFiveHundredHundred}{35}
\newcommand{\SpPaDoseTwoThousandHundred}{37}
\newcommand{\SpPaGroups}{55}
\newcommand{\SpPaLossLowestEverywhere}{no}
\newcommand{\SpPaMapGain}{1.4}
\newcommand{\SpPaMapMeasuredHundred}{34}
\newcommand{\SpPaMapPairedHundred}{\ensuremath{-}306}
\newcommand{\SpPaMarkerHundred}{19}
\newcommand{\SpPaOneSidedHundred}{37}
\newcommand{\SpPaPairBasisHundred}{16}
\newcommand{\SpPaPairedAll}{81}
\newcommand{\SpPaPairedOnlyCellsFour}{236}
\newcommand{\SpPaPairedOnlyFourHundred}{51}
\newcommand{\SpPaPairedOnlyHundred}{21}
\newcommand{\SpPaPercondHundred}{14}
\newcommand{\SpPaPooledHundred}{36}
\newcommand{\SpPaRfFifty}{16}
\newcommand{\SpPaRfFourHundred}{67}
\newcommand{\SpPaRfHundred}{37}
\newcommand{\SpPaRidgeHundred}{31}
\newcommand{\SpPaSaveBestCiMin}{0.65}
\newcommand{\SpPaSaveBestFour}{1.12}
\newcommand{\SpPaSaveBestFourHi}{1.3}
\newcommand{\SpPaSaveBestFourLo}{0.88}
\newcommand{\SpPaSaveBestMax}{1.12}
\newcommand{\SpPaSaveBestMin}{0.99}
\newcommand{\SpPaSavePairedOnlyCiMin}{1.5}
\newcommand{\SpPaSavePairedOnlyFour}{2.1}
\newcommand{\SpPaSavePairedOnlyFourHi}{2.6}
\newcommand{\SpPaSavePairedOnlyFourLo}{1.7}
\newcommand{\SpPaSavePairedOnlyMax}{2.5}
\newcommand{\SpPaSavePairedOnlyMin}{1.9}
\newcommand{\SpPaSavePublishedCiMin}{0.82}
\newcommand{\SpPaSavePublishedFour}{1.2}
\newcommand{\SpPaSavePublishedFourHi}{1.5}
\newcommand{\SpPaSavePublishedFourLo}{0.94}
\newcommand{\SpPaSavePublishedMax}{1.3}
\newcommand{\SpPaSavePublishedMin}{1.14}
\newcommand{\SpPaSemiccaHundred}{28}
\newcommand{\SpPaTargets}{21}
\newcommand{\SpProvArrays}{200}
\newcommand{\SpProvManifestAt}{15:06:29 EDT}
\newcommand{\SpProvMaxDiff}{0}
\newcommand{\SpSelfCriterion}{61.9}
\newcommand{\SpSelfFrFifty}{41}
\newcommand{\SpSelfFullCriterion}{57.3}
\newcommand{\SpSelfPaFifty}{19}
\newcommand{\SpSelfPartitionCriterion}{57.4}

% Generated by extension/generality/make_assets.py. Do not edit.
\newcommand{\GenCellsBanc}{2,487}
\newcommand{\GenCellsBmCite}{36,000}
\newcommand{\GenCellsBmMulti}{36,450}
\newcommand{\GenCellsColon}{15,953}
\newcommand{\GenCellsGouwens}{3,395}
\newcommand{\GenCellsHao}{72,000}
\newcommand{\GenCellsMaleCns}{3,009}
\newcommand{\GenCellsPairedOnlyBanc}{90}
\newcommand{\GenCellsPairedOnlyBmCite}{$>$800}
\newcommand{\GenCellsPairedOnlyBmMulti}{$>$800}
\newcommand{\GenCellsPairedOnlyColon}{$>$800}
\newcommand{\GenCellsPairedOnlyGouwens}{219}
\newcommand{\GenCellsPairedOnlyHao}{$>$800}
\newcommand{\GenCellsPairedOnlyMaleCns}{73}
\newcommand{\GenCellsPairedOnlyScala}{$\leq$25}
\newcommand{\GenCellsPairedOnlyStephenson}{$>$800}
\newcommand{\GenCellsProposedBanc}{25}
\newcommand{\GenCellsProposedBmCite}{168}
\newcommand{\GenCellsProposedBmMulti}{285}
\newcommand{\GenCellsProposedColon}{126}
\newcommand{\GenCellsProposedGouwens}{157}
\newcommand{\GenCellsProposedHao}{408}
\newcommand{\GenCellsProposedMaleCns}{$\leq$25}
\newcommand{\GenCellsProposedScala}{$>$100}
\newcommand{\GenCellsProposedStephenson}{390}
\newcommand{\GenCellsScala}{1,203}
\newcommand{\GenCellsStephenson}{70,800}
\newcommand{\GenFrozenGOne}{1.5}
\newcommand{\GenFrozenGOneHi}{1.6}
\newcommand{\GenFrozenGOneLo}{1.4}
\newcommand{\GenFrozenGOneMet}{met}
\newcommand{\GenFrozenGTwo}{1.3}
\newcommand{\GenFrozenGTwoHi}{1.4}
\newcommand{\GenFrozenGTwoLo}{1.05}
\newcommand{\GenFrozenGTwoMet}{met}
\newcommand{\GenFrozenNTrivial}{five}
\newcommand{\GenFrozenSaveBanc}{3.7}
\newcommand{\GenFrozenSaveBmCite}{2.5}
\newcommand{\GenFrozenSaveBmMulti}{1.00}
\newcommand{\GenFrozenSaveColon}{1.00}
\newcommand{\GenFrozenSaveGouwens}{1.4}
\newcommand{\GenFrozenSaveHao}{1.00}
\newcommand{\GenFrozenSaveMaleCns}{3.1}
\newcommand{\GenFrozenSaveScala}{1.00}
\newcommand{\GenFrozenSaveStephenson}{1.00}
\newcommand{\GenGOne}{2.9}
\newcommand{\GenGOneAll}{2.2}
\newcommand{\GenGOneAllHi}{2.8}
\newcommand{\GenGOneAllLo}{1.8}
\newcommand{\GenGOneAllMet}{met}
\newcommand{\GenGOneHi}{3.0}
\newcommand{\GenGOneLo}{2.3}
\newcommand{\GenGOneMet}{met}
\newcommand{\GenGOneQuarter}{2.9}
\newcommand{\GenGOneQuarterAll}{2.2}
\newcommand{\GenGOneQuarterAllHi}{2.7}
\newcommand{\GenGOneQuarterAllLo}{1.7}
\newcommand{\GenGOneQuarterHi}{3.1}
\newcommand{\GenGOneQuarterLo}{2.2}
\newcommand{\GenGOneThreeQuarters}{1.9}
\newcommand{\GenGOneThreeQuartersAll}{1.5}
\newcommand{\GenGOneThreeQuartersAllHi}{1.9}
\newcommand{\GenGOneThreeQuartersAllLo}{1.4}
\newcommand{\GenGOneThreeQuartersHi}{2.0}
\newcommand{\GenGOneThreeQuartersLo}{1.7}
\newcommand{\GenGOriginal}{1.4}
\newcommand{\GenGOriginalAll}{1.13}
\newcommand{\GenGOriginalAllHi}{1.3}
\newcommand{\GenGOriginalAllLo}{0.94}
\newcommand{\GenGOriginalHi}{1.4}
\newcommand{\GenGOriginalLo}{1.10}
\newcommand{\GenGThree}{0.58}
\newcommand{\GenGThreeAll}{0.61}
\newcommand{\GenGThreeAllHi}{0.71}
\newcommand{\GenGThreeAllLo}{0.59}
\newcommand{\GenGThreeHi}{0.69}
\newcommand{\GenGThreeLo}{0.55}
\newcommand{\GenGTwo}{2.3}
\newcommand{\GenGTwoAll}{2.1}
\newcommand{\GenGTwoAllHi}{2.2}
\newcommand{\GenGTwoAllLo}{1.6}
\newcommand{\GenGTwoAllMet}{met}
\newcommand{\GenGTwoHi}{2.4}
\newcommand{\GenGTwoLo}{1.8}
\newcommand{\GenGTwoMet}{met}
\newcommand{\GenLawC}{0.92}
\newcommand{\GenLawCP}{0.008}
\newcommand{\GenLawCPairs}{24}
\newcommand{\GenLawFrozenMet}{not met}
\newcommand{\GenLawFrozenN}{12}
\newcommand{\GenLawFrozenP}{0.627}
\newcommand{\GenLawFrozenRho}{-0.09}
\newcommand{\GenLawMet}{not met}
\newcommand{\GenLawN}{11}
\newcommand{\GenLawNCensored}{six}
\newcommand{\GenLawNConsistent}{five}
\newcommand{\GenLawNUncensored}{five}
\newcommand{\GenLawP}{0.104}
\newcommand{\GenLawRatioMax}{1.3}
\newcommand{\GenLawRatioMin}{0.6}
\newcommand{\GenLawRho}{0.42}
\newcommand{\GenLawUncensoredRho}{0.90}
\newcommand{\GenLevelBanc}{68}
\newcommand{\GenLevelBancX}{49}
\newcommand{\GenLevelBmCite}{46}
\newcommand{\GenLevelBmMulti}{25}
\newcommand{\GenLevelColon}{25}
\newcommand{\GenLevelGouwens}{42}
\newcommand{\GenLevelHao}{35}
\newcommand{\GenLevelMaleCns}{69}
\newcommand{\GenLevelScala}{0.0}
\newcommand{\GenLevelStephenson}{13}
\newcommand{\GenMaxBudgetBanc}{400}
\newcommand{\GenMaxBudgetBmCite}{800}
\newcommand{\GenMaxBudgetBmMulti}{800}
\newcommand{\GenMaxBudgetColon}{800}
\newcommand{\GenMaxBudgetGouwens}{400}
\newcommand{\GenMaxBudgetHao}{800}
\newcommand{\GenMaxBudgetMaleCns}{400}
\newcommand{\GenMaxBudgetScala}{100}
\newcommand{\GenMaxBudgetStephenson}{800}
\newcommand{\GenNInformative}{eight}
\newcommand{\GenNLbAboveOne}{five}
\newcommand{\GenNLbAboveOneSemi}{five}
\newcommand{\GenNotInformative}{motor cortex}
\newcommand{\GenPairedOnlyHundredBanc}{36}
\newcommand{\GenPairedOnlyHundredBancX}{\ensuremath{-}52}
\newcommand{\GenPairedOnlyHundredBmCite}{1.1}
\newcommand{\GenPairedOnlyHundredBmMulti}{0.0}
\newcommand{\GenPairedOnlyHundredColon}{1.2}
\newcommand{\GenPairedOnlyHundredGouwens}{12}
\newcommand{\GenPairedOnlyHundredHao}{0.7}
\newcommand{\GenPairedOnlyHundredMaleCns}{40}
\newcommand{\GenPairedOnlyHundredScala}{0.0}
\newcommand{\GenPairedOnlyHundredStephenson}{0.4}
\newcommand{\GenPairedOnlyMaxBanc}{55}
\newcommand{\GenPairedOnlyMaxBmCite}{5.9}
\newcommand{\GenPairedOnlyMaxBmMulti}{0.6}
\newcommand{\GenPairedOnlyMaxColon}{7.7}
\newcommand{\GenPairedOnlyMaxGouwens}{29}
\newcommand{\GenPairedOnlyMaxHao}{2.1}
\newcommand{\GenPairedOnlyMaxMaleCns}{56}
\newcommand{\GenPairedOnlyMaxScala}{0.0}
\newcommand{\GenPairedOnlyMaxStephenson}{3.7}
\newcommand{\GenPropMaxBancX}{\ensuremath{-}30}
\newcommand{\GenPropPooledMaxBancX}{\ensuremath{-}9.1}
\newcommand{\GenRefBanc}{76}
\newcommand{\GenRefBmCite}{66}
\newcommand{\GenRefBmMulti}{\ensuremath{-}148}
\newcommand{\GenRefColon}{\ensuremath{-}237}
\newcommand{\GenRefGouwens}{43}
\newcommand{\GenRefHao}{70}
\newcommand{\GenRefMaleCns}{77}
\newcommand{\GenRefMax}{77}
\newcommand{\GenRefScala}{\ensuremath{-}274}
\newcommand{\GenRefStephenson}{\ensuremath{-}1042}
\newcommand{\GenRfHundredBanc}{56}
\newcommand{\GenRfHundredBancX}{\ensuremath{-}78}
\newcommand{\GenRfHundredBmCite}{16}
\newcommand{\GenRfHundredBmMulti}{\ensuremath{-}3.7}
\newcommand{\GenRfHundredColon}{11}
\newcommand{\GenRfHundredGouwens}{15}
\newcommand{\GenRfHundredHao}{4.3}
\newcommand{\GenRfHundredMaleCns}{54}
\newcommand{\GenRfHundredScala}{\ensuremath{-}11}
\newcommand{\GenRfHundredStephenson}{\ensuremath{-}0.7}
\newcommand{\GenRfMaxBanc}{68}
\newcommand{\GenRfMaxBmCite}{46}
\newcommand{\GenRfMaxBmMulti}{25}
\newcommand{\GenRfMaxColon}{25}
\newcommand{\GenRfMaxGouwens}{32}
\newcommand{\GenRfMaxHao}{30}
\newcommand{\GenRfMaxMaleCns}{67}
\newcommand{\GenRfMaxScala}{\ensuremath{-}11}
\newcommand{\GenRfMaxStephenson}{9.3}
\newcommand{\GenRidgeHundredBancX}{31}
\newcommand{\GenRidgeMinBancX}{11}
\newcommand{\GenSaveMax}{6.4}
\newcommand{\GenSaveMin}{1.4}
\newcommand{\GenSavePairedOnlyBanc}{3.6}
\newcommand{\GenSavePairedOnlyBancHi}{3.8}
\newcommand{\GenSavePairedOnlyBancLo}{3.1}
\newcommand{\GenSavePairedOnlyBancX}{$\geq$1.00}
\newcommand{\GenSavePairedOnlyBancXHi}{1.00}
\newcommand{\GenSavePairedOnlyBancXLo}{1.00}
\newcommand{\GenSavePairedOnlyBmCite}{$\geq$4.7}
\newcommand{\GenSavePairedOnlyBmCiteHi}{6.0}
\newcommand{\GenSavePairedOnlyBmCiteLo}{1.00}
\newcommand{\GenSavePairedOnlyBmMulti}{$\geq$2.8}
\newcommand{\GenSavePairedOnlyBmMultiHi}{3.0}
\newcommand{\GenSavePairedOnlyBmMultiLo}{1.7}
\newcommand{\GenSavePairedOnlyColon}{$\geq$6.4}
\newcommand{\GenSavePairedOnlyColonHi}{7.7}
\newcommand{\GenSavePairedOnlyColonLo}{3.3}
\newcommand{\GenSavePairedOnlyGouwens}{1.4}
\newcommand{\GenSavePairedOnlyGouwensHi}{1.6}
\newcommand{\GenSavePairedOnlyGouwensLo}{1.2}
\newcommand{\GenSavePairedOnlyHao}{$\geq$2.0}
\newcommand{\GenSavePairedOnlyHaoHi}{2.1}
\newcommand{\GenSavePairedOnlyHaoLo}{1.00}
\newcommand{\GenSavePairedOnlyMaleCns}{2.9}
\newcommand{\GenSavePairedOnlyMaleCnsHi}{3.1}
\newcommand{\GenSavePairedOnlyMaleCnsLo}{2.6}
\newcommand{\GenSavePairedOnlyScala}{$\leq$0.25}
\newcommand{\GenSavePairedOnlyScalaHi}{2.7}
\newcommand{\GenSavePairedOnlyScalaLo}{0.25}
\newcommand{\GenSavePairedOnlyStephenson}{$\geq$2.0}
\newcommand{\GenSavePairedOnlyStephensonHi}{2.6}
\newcommand{\GenSavePairedOnlyStephensonLo}{1.00}
\newcommand{\GenSaveSemiBanc}{2.3}
\newcommand{\GenSaveSemiBancHi}{2.4}
\newcommand{\GenSaveSemiBancLo}{2.0}
\newcommand{\GenSaveSemiBmCite}{3.7}
\newcommand{\GenSaveSemiBmCiteHi}{4.6}
\newcommand{\GenSaveSemiBmCiteLo}{1.00}
\newcommand{\GenSaveSemiBmMulti}{$\geq$2.8}
\newcommand{\GenSaveSemiBmMultiHi}{3.0}
\newcommand{\GenSaveSemiBmMultiLo}{1.7}
\newcommand{\GenSaveSemiColon}{$\geq$6.4}
\newcommand{\GenSaveSemiColonHi}{7.7}
\newcommand{\GenSaveSemiColonLo}{1.9}
\newcommand{\GenSaveSemiGouwens}{1.3}
\newcommand{\GenSaveSemiGouwensHi}{1.5}
\newcommand{\GenSaveSemiGouwensLo}{1.10}
\newcommand{\GenSaveSemiHao}{$\geq$2.0}
\newcommand{\GenSaveSemiHaoHi}{2.1}
\newcommand{\GenSaveSemiHaoLo}{1.00}
\newcommand{\GenSaveSemiMaleCns}{1.9}
\newcommand{\GenSaveSemiMaleCnsHi}{2.1}
\newcommand{\GenSaveSemiMaleCnsLo}{1.7}
\newcommand{\GenSaveSemiScala}{$\geq$1.00}
\newcommand{\GenSaveSemiScalaHi}{1.3}
\newcommand{\GenSaveSemiScalaLo}{0.25}
\newcommand{\GenSaveSemiStephenson}{1.02}
\newcommand{\GenSaveSemiStephensonHi}{1.3}
\newcommand{\GenSaveSemiStephensonLo}{0.53}
\newcommand{\GenSelfBanc}{0.62}
\newcommand{\GenSelfBetter}{none}
\newcommand{\GenSelfBmCite}{0.39}
\newcommand{\GenSelfBmMulti}{0.81}
\newcommand{\GenSelfColon}{0.59}
\newcommand{\GenSelfGouwens}{0.59}
\newcommand{\GenSelfHao}{0.54}
\newcommand{\GenSelfMaleCns}{0.69}
\newcommand{\GenSelfPooledHundredBancX}{7.0}
\newcommand{\GenSelfScala}{1.00}
\newcommand{\GenSelfStephenson}{0.49}
\newcommand{\GenSemiHundredBanc}{44}
\newcommand{\GenSemiHundredBancX}{28}
\newcommand{\GenSemiHundredBmCite}{7.7}
\newcommand{\GenSemiHundredBmMulti}{\ensuremath{-}1.9}
\newcommand{\GenSemiHundredColon}{1.6}
\newcommand{\GenSemiHundredGouwens}{11}
\newcommand{\GenSemiHundredHao}{6.5}
\newcommand{\GenSemiHundredMaleCns}{47}
\newcommand{\GenSemiHundredScala}{\ensuremath{-}3.6}
\newcommand{\GenSemiHundredStephenson}{1.6}
\newcommand{\GenUnitsBanc}{913}
\newcommand{\GenUnitsBmCite}{9}
\newcommand{\GenUnitsBmMulti}{10}
\newcommand{\GenUnitsColon}{12}
\newcommand{\GenUnitsGouwens}{362}
\newcommand{\GenUnitsHao}{8}
\newcommand{\GenUnitsMaleCns}{1,010}
\newcommand{\GenUnitsScala}{263}
\newcommand{\GenUnitsStephenson}{118}

% Generated by extension/predictability/make_assets.py. Do not edit.
\newcommand{\PredBootErr}{6}
\newcommand{\PredBootErrSeven}{8}
\newcommand{\PredBootErrThree}{33}
\newcommand{\PredBootRatio}{1.02}
\newcommand{\PredCellsBlocksRatio}{1.23}
\newcommand{\PredCellsJsRatio}{1.04}
\newcommand{\PredCensAgree}{34}
\newcommand{\PredCensTotal}{37}
\newcommand{\PredCheckCases}{72}
\newcommand{\PredCheckMaxRatio}{0.10}
\newcommand{\PredCurveErrBlocks}{4.2}
\newcommand{\PredCurveErrJs}{0.8}
\newcommand{\PredDataAt}{21:56 EDT}
\newcommand{\PredDevLawErr}{7}
\newcommand{\PredDevLawErrMax}{22}
\newcommand{\PredDevN}{14}
\newcommand{\PredDevSaveMax}{14.7}
\newcommand{\PredDevSaveMin}{1.4}
\newcommand{\PredFrozenAt}{21:39 EDT}
\newcommand{\PredHFive}{met}
\newcommand{\PredHFour}{met}
\newcommand{\PredHOne}{met}
\newcommand{\PredHSix}{met}
\newcommand{\PredHThree}{met}
\newcommand{\PredHTwo}{met}
\newcommand{\PredLawErr}{21}
\newcommand{\PredLawErrHi}{29}
\newcommand{\PredLawErrLo}{11}
\newcommand{\PredLawErrSeven}{4}
\newcommand{\PredLawErrThree}{74}
\newcommand{\PredLawRatio}{0.83}
\newcommand{\PredLawRatioSeven}{0.96}
\newcommand{\PredLawRatioThree}{0.57}
\newcommand{\PredLawWithin}{60}
\newcommand{\PredNCensored}{19}
\newcommand{\PredNData}{ten}
\newcommand{\PredNEasy}{three}
\newcommand{\PredNEstimable}{10}
\newcommand{\PredNEstimableData}{four}
\newcommand{\PredNProblems}{29}
\newcommand{\PredNSeven}{five}
\newcommand{\PredNThree}{eight}
\newcommand{\PredPilotErr}{12}
\newcommand{\PredPilotErrHi}{147}
\newcommand{\PredPilotErrLo}{4}
\newcommand{\PredPilotErrSeven}{2}
\newcommand{\PredPilotErrThree}{96}
\newcommand{\PredPilotWithin}{70}
\newcommand{\PredProfileN}{eight}
\newcommand{\PredProfileTop}{73}
\newcommand{\PredProfileTopTwenty}{94}
\newcommand{\PredRandomLaw}{1.01}
\newcommand{\PredRandomN}{11}
\newcommand{\PredRandomObs}{1.00}
\newcommand{\PredRho}{0.98}
\newcommand{\PredRhoRel}{<0.001}
\newcommand{\PredSaveMax}{11.6}
\newcommand{\PredSaveMin}{1.9}
\newcommand{\PredSnareLawMax}{408}
\newcommand{\PredSnareLawMin}{396}
\newcommand{\PredSnareObsMax}{1,357}
\newcommand{\PredSnareObsMin}{783}
\newcommand{\PredTransferErr}{40}
\newcommand{\PredTransferErrHi}{45}
\newcommand{\PredTransferErrLo}{31}
\newcommand{\PredTransferWithin}{10}

% Generated by extension/synthesis/make_assets.py. Do not edit.
\newcommand{\AlNDonors}{18}
\newcommand{\AlNPools}{seven}
\newcommand{\AlNSamples}{34}
\newcommand{\AlNTwoTimepoints}{16}
\newcommand{\AltCellsAtlas}{147}
\newcommand{\AltCellsAtlasHi}{185}
\newcommand{\AltCellsAtlasLo}{125}
\newcommand{\AltCellsChamp}{180}
\newcommand{\AltCellsChampAtlas}{147}
\newcommand{\AltCellsChampAtlasHi}{189}
\newcommand{\AltCellsChampAtlasLo}{123}
\newcommand{\AltCellsChampHi}{198}
\newcommand{\AltCellsChampLo}{165}
\newcommand{\AltCellsOther}{70}
\newcommand{\AltCellsOtherHi}{85}
\newcommand{\AltCellsOtherLo}{58}
\newcommand{\AltCellsOwn}{110}
\newcommand{\AltCellsOwnHi}{165}
\newcommand{\AltCellsOwnLo}{83}
\newcommand{\AltCellsPo}{$>$800}
\newcommand{\AltCellsPoHi}{800}
\newcommand{\AltCellsPoLo}{800}
\newcommand{\AltCellsRidge}{169}
\newcommand{\AltCellsRidgeHi}{236}
\newcommand{\AltCellsRidgeLo}{131}
\newcommand{\AltCellsRidgePMapped}{235}
\newcommand{\AltCellsRidgePMappedHi}{275}
\newcommand{\AltCellsRidgePMappedLo}{203}
\newcommand{\AltCellsRidgePPercond}{197}
\newcommand{\AltCellsRidgePPercondHi}{241}
\newcommand{\AltCellsRidgePPercondLo}{158}
\newcommand{\AltCellsRidgePPooled}{$>$800}
\newcommand{\AltCellsRidgePPooledHi}{800}
\newcommand{\AltCellsRidgePPooledLo}{800}
\newcommand{\AltCellsRidgeUMapped}{214}
\newcommand{\AltCellsRidgeUMappedHi}{248}
\newcommand{\AltCellsRidgeUMappedLo}{186}
\newcommand{\AltCellsRidgeUPercond}{175}
\newcommand{\AltCellsRidgeUPercondHi}{414}
\newcommand{\AltCellsRidgeUPercondLo}{136}
\newcommand{\AltCellsRidgeUPooled}{$>$800}
\newcommand{\AltCellsRidgeUPooledHi}{800}
\newcommand{\AltCellsRidgeUPooledLo}{800}
\newcommand{\AltCellsSemi}{496}
\newcommand{\AltCellsSemiHi}{800}
\newcommand{\AltCellsSemiLo}{308}
\newcommand{\AltCellsSemiMapped}{579}
\newcommand{\AltCellsSemiMappedHi}{800}
\newcommand{\AltCellsSemiMappedLo}{417}
\newcommand{\AltCellsSemiPercond}{$>$800}
\newcommand{\AltCellsSemiPercondHi}{800}
\newcommand{\AltCellsSemiPercondLo}{308}
\newcommand{\AltCellsSemiPooled}{$>$800}
\newcommand{\AltCellsSemiPooledHi}{800}
\newcommand{\AltCellsSemiPooledLo}{800}
\newcommand{\AltChampFit}{50}
\newcommand{\AltChampFolds}{34}
\newcommand{\AltCompared}{11,220}
\newcommand{\AltDiffOneChampEightHundred}{0.7}
\newcommand{\AltDiffOneChampEightHundredHi}{1.1}
\newcommand{\AltDiffOneChampEightHundredLo}{0.4}
\newcommand{\AltDiffOneChampFifty}{\ensuremath{-}1.3}
\newcommand{\AltDiffOneChampFiftyHi}{0.4}
\newcommand{\AltDiffOneChampFiftyLo}{\ensuremath{-}3.0}
\newcommand{\AltDiffOneChampFourHundred}{0.2}
\newcommand{\AltDiffOneChampFourHundredHi}{0.7}
\newcommand{\AltDiffOneChampFourHundredLo}{\ensuremath{-}0.3}
\newcommand{\AltDiffOneChampHundred}{\ensuremath{-}1.2}
\newcommand{\AltDiffOneChampHundredHi}{\ensuremath{-}0.0}
\newcommand{\AltDiffOneChampHundredLo}{\ensuremath{-}2.4}
\newcommand{\AltDiffOneChampTwentyFive}{4.0}
\newcommand{\AltDiffOneChampTwentyFiveHi}{6.9}
\newcommand{\AltDiffOneChampTwentyFiveLo}{1.0}
\newcommand{\AltDiffOneChampTwoHundred}{\ensuremath{-}0.3}
\newcommand{\AltDiffOneChampTwoHundredHi}{0.3}
\newcommand{\AltDiffOneChampTwoHundredLo}{\ensuremath{-}0.9}
\newcommand{\AltDiffOnePoEightHundred}{1.5}
\newcommand{\AltDiffOnePoEightHundredHi}{1.8}
\newcommand{\AltDiffOnePoEightHundredLo}{1.0}
\newcommand{\AltDiffOnePoFifty}{12.9}
\newcommand{\AltDiffOnePoFiftyHi}{13.8}
\newcommand{\AltDiffOnePoFiftyLo}{12.0}
\newcommand{\AltDiffOnePoFourHundred}{4.1}
\newcommand{\AltDiffOnePoFourHundredHi}{4.7}
\newcommand{\AltDiffOnePoFourHundredLo}{3.4}
\newcommand{\AltDiffOnePoHundred}{11.2}
\newcommand{\AltDiffOnePoHundredHi}{12.0}
\newcommand{\AltDiffOnePoHundredLo}{10.3}
\newcommand{\AltDiffOnePoOneFifty}{9.1}
\newcommand{\AltDiffOnePoOneFiftyHi}{9.9}
\newcommand{\AltDiffOnePoOneFiftyLo}{8.3}
\newcommand{\AltDiffOnePoTwentyFive}{9.2}
\newcommand{\AltDiffOnePoTwentyFiveHi}{9.9}
\newcommand{\AltDiffOnePoTwentyFiveLo}{8.5}
\newcommand{\AltDiffOnePoTwoHundred}{7.6}
\newcommand{\AltDiffOnePoTwoHundredHi}{8.2}
\newcommand{\AltDiffOnePoTwoHundredLo}{6.8}
\newcommand{\AltDiffOneRidgeEightHundred}{\ensuremath{-}1.5}
\newcommand{\AltDiffOneRidgeEightHundredHi}{\ensuremath{-}1.3}
\newcommand{\AltDiffOneRidgeEightHundredLo}{\ensuremath{-}1.7}
\newcommand{\AltDiffOneRidgeFifty}{0.0}
\newcommand{\AltDiffOneRidgeFiftyHi}{0.3}
\newcommand{\AltDiffOneRidgeFiftyLo}{\ensuremath{-}0.2}
\newcommand{\AltDiffOneRidgeFourHundred}{\ensuremath{-}1.3}
\newcommand{\AltDiffOneRidgeFourHundredHi}{\ensuremath{-}1.0}
\newcommand{\AltDiffOneRidgeFourHundredLo}{\ensuremath{-}1.6}
\newcommand{\AltDiffOneRidgeHundred}{\ensuremath{-}0.1}
\newcommand{\AltDiffOneRidgeHundredHi}{0.3}
\newcommand{\AltDiffOneRidgeHundredLo}{\ensuremath{-}0.5}
\newcommand{\AltDiffOneRidgeOneFifty}{\ensuremath{-}0.6}
\newcommand{\AltDiffOneRidgeOneFiftyHi}{\ensuremath{-}0.2}
\newcommand{\AltDiffOneRidgeOneFiftyLo}{\ensuremath{-}1.0}
\newcommand{\AltDiffOneRidgeTwentyFive}{\ensuremath{-}1.7}
\newcommand{\AltDiffOneRidgeTwentyFiveHi}{\ensuremath{-}1.1}
\newcommand{\AltDiffOneRidgeTwentyFiveLo}{\ensuremath{-}2.5}
\newcommand{\AltDiffOneRidgeTwoHundred}{\ensuremath{-}0.9}
\newcommand{\AltDiffOneRidgeTwoHundredHi}{\ensuremath{-}0.6}
\newcommand{\AltDiffOneRidgeTwoHundredLo}{\ensuremath{-}1.2}
\newcommand{\AltDiffOneSemiEightHundred}{2.4}
\newcommand{\AltDiffOneSemiEightHundredHi}{2.9}
\newcommand{\AltDiffOneSemiEightHundredLo}{1.9}
\newcommand{\AltDiffOneSemiFifty}{11.1}
\newcommand{\AltDiffOneSemiFiftyHi}{12.5}
\newcommand{\AltDiffOneSemiFiftyLo}{9.7}
\newcommand{\AltDiffOneSemiFourHundred}{5.1}
\newcommand{\AltDiffOneSemiFourHundredHi}{5.6}
\newcommand{\AltDiffOneSemiFourHundredLo}{4.6}
\newcommand{\AltDiffOneSemiHundred}{9.6}
\newcommand{\AltDiffOneSemiHundredHi}{11.1}
\newcommand{\AltDiffOneSemiHundredLo}{8.3}
\newcommand{\AltDiffOneSemiOneFifty}{9.0}
\newcommand{\AltDiffOneSemiOneFiftyHi}{9.7}
\newcommand{\AltDiffOneSemiOneFiftyLo}{7.9}
\newcommand{\AltDiffOneSemiTwentyFive}{11.2}
\newcommand{\AltDiffOneSemiTwentyFiveHi}{12.7}
\newcommand{\AltDiffOneSemiTwentyFiveLo}{9.8}
\newcommand{\AltDiffOneSemiTwoHundred}{7.2}
\newcommand{\AltDiffOneSemiTwoHundredHi}{8.2}
\newcommand{\AltDiffOneSemiTwoHundredLo}{6.2}
\newcommand{\AltDiffThreePoEightHundred}{8.8}
\newcommand{\AltDiffThreePoEightHundredHi}{10.5}
\newcommand{\AltDiffThreePoEightHundredLo}{7.0}
\newcommand{\AltDiffThreePoFifty}{13.0}
\newcommand{\AltDiffThreePoFiftyHi}{14.1}
\newcommand{\AltDiffThreePoFiftyLo}{11.5}
\newcommand{\AltDiffThreePoFourHundred}{10.2}
\newcommand{\AltDiffThreePoFourHundredHi}{11.6}
\newcommand{\AltDiffThreePoFourHundredLo}{8.5}
\newcommand{\AltDiffThreePoHundred}{14.2}
\newcommand{\AltDiffThreePoHundredHi}{15.7}
\newcommand{\AltDiffThreePoHundredLo}{12.5}
\newcommand{\AltDiffThreePoOneFifty}{14.4}
\newcommand{\AltDiffThreePoOneFiftyHi}{15.4}
\newcommand{\AltDiffThreePoOneFiftyLo}{12.9}
\newcommand{\AltDiffThreePoTwentyFive}{7.9}
\newcommand{\AltDiffThreePoTwentyFiveHi}{9.0}
\newcommand{\AltDiffThreePoTwentyFiveLo}{7.0}
\newcommand{\AltDiffThreePoTwoHundred}{14.1}
\newcommand{\AltDiffThreePoTwoHundredHi}{15.2}
\newcommand{\AltDiffThreePoTwoHundredLo}{12.7}
\newcommand{\AltDiffThreeRidgeEightHundred}{17.5}
\newcommand{\AltDiffThreeRidgeEightHundredHi}{18.3}
\newcommand{\AltDiffThreeRidgeEightHundredLo}{15.7}
\newcommand{\AltDiffThreeRidgeFifty}{6.9}
\newcommand{\AltDiffThreeRidgeFiftyHi}{7.7}
\newcommand{\AltDiffThreeRidgeFiftyLo}{6.1}
\newcommand{\AltDiffThreeRidgeFourHundred}{10.9}
\newcommand{\AltDiffThreeRidgeFourHundredHi}{12.0}
\newcommand{\AltDiffThreeRidgeFourHundredLo}{10.1}
\newcommand{\AltDiffThreeRidgeHundred}{11.3}
\newcommand{\AltDiffThreeRidgeHundredHi}{12.0}
\newcommand{\AltDiffThreeRidgeHundredLo}{10.5}
\newcommand{\AltDiffThreeRidgeOneFifty}{13.6}
\newcommand{\AltDiffThreeRidgeOneFiftyHi}{14.5}
\newcommand{\AltDiffThreeRidgeOneFiftyLo}{12.2}
\newcommand{\AltDiffThreeRidgeTwentyFive}{2.4}
\newcommand{\AltDiffThreeRidgeTwentyFiveHi}{3.1}
\newcommand{\AltDiffThreeRidgeTwentyFiveLo}{1.5}
\newcommand{\AltDiffThreeRidgeTwoHundred}{12.8}
\newcommand{\AltDiffThreeRidgeTwoHundredHi}{14.3}
\newcommand{\AltDiffThreeRidgeTwoHundredLo}{11.3}
\newcommand{\AltDiffThreeSemiEightHundred}{27.8}
\newcommand{\AltDiffThreeSemiEightHundredHi}{29.5}
\newcommand{\AltDiffThreeSemiEightHundredLo}{26.1}
\newcommand{\AltDiffThreeSemiFifty}{19.3}
\newcommand{\AltDiffThreeSemiFiftyHi}{21.1}
\newcommand{\AltDiffThreeSemiFiftyLo}{17.4}
\newcommand{\AltDiffThreeSemiFourHundred}{30.2}
\newcommand{\AltDiffThreeSemiFourHundredHi}{31.5}
\newcommand{\AltDiffThreeSemiFourHundredLo}{29.0}
\newcommand{\AltDiffThreeSemiHundred}{26.0}
\newcommand{\AltDiffThreeSemiHundredHi}{27.8}
\newcommand{\AltDiffThreeSemiHundredLo}{24.1}
\newcommand{\AltDiffThreeSemiOneFifty}{26.0}
\newcommand{\AltDiffThreeSemiOneFiftyHi}{27.7}
\newcommand{\AltDiffThreeSemiOneFiftyLo}{24.4}
\newcommand{\AltDiffThreeSemiTwentyFive}{11.4}
\newcommand{\AltDiffThreeSemiTwentyFiveHi}{13.1}
\newcommand{\AltDiffThreeSemiTwentyFiveLo}{9.7}
\newcommand{\AltDiffThreeSemiTwoHundred}{26.7}
\newcommand{\AltDiffThreeSemiTwoHundredHi}{28.4}
\newcommand{\AltDiffThreeSemiTwoHundredLo}{25.2}
\newcommand{\AltDiffTwoChampEightHundred}{2.3}
\newcommand{\AltDiffTwoChampEightHundredHi}{3.5}
\newcommand{\AltDiffTwoChampEightHundredLo}{1.3}
\newcommand{\AltDiffTwoChampFifty}{7.8}
\newcommand{\AltDiffTwoChampFiftyHi}{9.7}
\newcommand{\AltDiffTwoChampFiftyLo}{6.0}
\newcommand{\AltDiffTwoChampFourHundred}{3.7}
\newcommand{\AltDiffTwoChampFourHundredHi}{5.0}
\newcommand{\AltDiffTwoChampFourHundredLo}{2.5}
\newcommand{\AltDiffTwoChampHundred}{4.7}
\newcommand{\AltDiffTwoChampHundredHi}{6.8}
\newcommand{\AltDiffTwoChampHundredLo}{2.6}
\newcommand{\AltDiffTwoChampTwentyFive}{11.4}
\newcommand{\AltDiffTwoChampTwentyFiveHi}{13.8}
\newcommand{\AltDiffTwoChampTwentyFiveLo}{9.2}
\newcommand{\AltDiffTwoChampTwoHundred}{6.2}
\newcommand{\AltDiffTwoChampTwoHundredHi}{7.6}
\newcommand{\AltDiffTwoChampTwoHundredLo}{4.9}
\newcommand{\AltDiffTwoPoEightHundred}{5.6}
\newcommand{\AltDiffTwoPoEightHundredHi}{6.4}
\newcommand{\AltDiffTwoPoEightHundredLo}{4.8}
\newcommand{\AltDiffTwoPoFifty}{20.1}
\newcommand{\AltDiffTwoPoFiftyHi}{21.3}
\newcommand{\AltDiffTwoPoFiftyLo}{18.9}
\newcommand{\AltDiffTwoPoFourHundred}{9.9}
\newcommand{\AltDiffTwoPoFourHundredHi}{10.9}
\newcommand{\AltDiffTwoPoFourHundredLo}{9.0}
\newcommand{\AltDiffTwoPoHundred}{19.3}
\newcommand{\AltDiffTwoPoHundredHi}{20.4}
\newcommand{\AltDiffTwoPoHundredLo}{18.0}
\newcommand{\AltDiffTwoPoOneFifty}{16.8}
\newcommand{\AltDiffTwoPoOneFiftyHi}{17.9}
\newcommand{\AltDiffTwoPoOneFiftyLo}{15.6}
\newcommand{\AltDiffTwoPoTwentyFive}{15.7}
\newcommand{\AltDiffTwoPoTwentyFiveHi}{16.6}
\newcommand{\AltDiffTwoPoTwentyFiveLo}{14.8}
\newcommand{\AltDiffTwoPoTwoHundred}{15.2}
\newcommand{\AltDiffTwoPoTwoHundredHi}{16.0}
\newcommand{\AltDiffTwoPoTwoHundredLo}{14.2}
\newcommand{\AltDiffTwoRidgeEightHundred}{\ensuremath{-}4.1}
\newcommand{\AltDiffTwoRidgeEightHundredHi}{\ensuremath{-}3.1}
\newcommand{\AltDiffTwoRidgeEightHundredLo}{\ensuremath{-}5.0}
\newcommand{\AltDiffTwoRidgeFifty}{3.0}
\newcommand{\AltDiffTwoRidgeFiftyHi}{3.9}
\newcommand{\AltDiffTwoRidgeFiftyLo}{1.8}
\newcommand{\AltDiffTwoRidgeFourHundred}{\ensuremath{-}5.1}
\newcommand{\AltDiffTwoRidgeFourHundredHi}{\ensuremath{-}4.2}
\newcommand{\AltDiffTwoRidgeFourHundredLo}{\ensuremath{-}6.0}
\newcommand{\AltDiffTwoRidgeHundred}{\ensuremath{-}2.8}
\newcommand{\AltDiffTwoRidgeHundredHi}{\ensuremath{-}1.9}
\newcommand{\AltDiffTwoRidgeHundredLo}{\ensuremath{-}3.6}
\newcommand{\AltDiffTwoRidgeOneFifty}{\ensuremath{-}4.3}
\newcommand{\AltDiffTwoRidgeOneFiftyHi}{\ensuremath{-}3.4}
\newcommand{\AltDiffTwoRidgeOneFiftyLo}{\ensuremath{-}5.0}
\newcommand{\AltDiffTwoRidgeTwentyFive}{1.7}
\newcommand{\AltDiffTwoRidgeTwentyFiveHi}{2.8}
\newcommand{\AltDiffTwoRidgeTwentyFiveLo}{0.5}
\newcommand{\AltDiffTwoRidgeTwoHundred}{\ensuremath{-}4.5}
\newcommand{\AltDiffTwoRidgeTwoHundredHi}{\ensuremath{-}3.6}
\newcommand{\AltDiffTwoRidgeTwoHundredLo}{\ensuremath{-}5.5}
\newcommand{\AltDiffTwoSemiEightHundred}{3.1}
\newcommand{\AltDiffTwoSemiEightHundredHi}{4.2}
\newcommand{\AltDiffTwoSemiEightHundredLo}{2.2}
\newcommand{\AltDiffTwoSemiFifty}{12.5}
\newcommand{\AltDiffTwoSemiFiftyHi}{13.9}
\newcommand{\AltDiffTwoSemiFiftyLo}{11.2}
\newcommand{\AltDiffTwoSemiFourHundred}{10.0}
\newcommand{\AltDiffTwoSemiFourHundredHi}{11.5}
\newcommand{\AltDiffTwoSemiFourHundredLo}{8.6}
\newcommand{\AltDiffTwoSemiHundred}{15.4}
\newcommand{\AltDiffTwoSemiHundredHi}{17.1}
\newcommand{\AltDiffTwoSemiHundredLo}{13.8}
\newcommand{\AltDiffTwoSemiOneFifty}{14.8}
\newcommand{\AltDiffTwoSemiOneFiftyHi}{15.9}
\newcommand{\AltDiffTwoSemiOneFiftyLo}{13.8}
\newcommand{\AltDiffTwoSemiTwentyFive}{9.8}
\newcommand{\AltDiffTwoSemiTwentyFiveHi}{11.5}
\newcommand{\AltDiffTwoSemiTwentyFiveLo}{8.1}
\newcommand{\AltDiffTwoSemiTwoHundred}{14.5}
\newcommand{\AltDiffTwoSemiTwoHundredHi}{15.5}
\newcommand{\AltDiffTwoSemiTwoHundredLo}{13.7}
\newcommand{\AltFrozenAt}{16:31 EDT on 29 September 2026}
\newcommand{\AltIntegrity}{$1.2\times10^{-7}$}
\newcommand{\AltKOne}{met}
\newcommand{\AltKTwo}{met}
\newcommand{\AltOneAtlasEightHundred}{94}
\newcommand{\AltOneAtlasFifty}{80}
\newcommand{\AltOneAtlasFourHundred}{92}
\newcommand{\AltOneAtlasHundred}{86}
\newcommand{\AltOneAtlasOneFifty}{89}
\newcommand{\AltOneAtlasTwentyFive}{71}
\newcommand{\AltOneAtlasTwoHundred}{90}
\newcommand{\AltOneChampAtlasEightHundred}{94}
\newcommand{\AltOneChampAtlasFifty}{80}
\newcommand{\AltOneChampAtlasFourHundred}{92}
\newcommand{\AltOneChampAtlasHundred}{87}
\newcommand{\AltOneChampAtlasTwentyFive}{69}
\newcommand{\AltOneChampAtlasTwoHundred}{90}
\newcommand{\AltOneChampEightHundred}{93}
\newcommand{\AltOneChampFifty}{82}
\newcommand{\AltOneChampFourHundred}{92}
\newcommand{\AltOneChampHundred}{88}
\newcommand{\AltOneChampTwentyFive}{65}
\newcommand{\AltOneChampTwoHundred}{90}
\newcommand{\AltOneOtherEightHundred}{96}
\newcommand{\AltOneOtherFifty}{85}
\newcommand{\AltOneOtherFourHundred}{94}
\newcommand{\AltOneOtherHundred}{89}
\newcommand{\AltOneOtherOneFifty}{91}
\newcommand{\AltOneOtherTwentyFive}{77}
\newcommand{\AltOneOtherTwoHundred}{92}
\newcommand{\AltOneOwnEightHundred}{94}
\newcommand{\AltOneOwnFifty}{84}
\newcommand{\AltOneOwnFourHundred}{92}
\newcommand{\AltOneOwnHundred}{88}
\newcommand{\AltOneOwnOneFifty}{89}
\newcommand{\AltOneOwnTwentyFive}{76}
\newcommand{\AltOneOwnTwoHundred}{90}
\newcommand{\AltOnePoEightHundred}{92}
\newcommand{\AltOnePoFifty}{67}
\newcommand{\AltOnePoFourHundred}{88}
\newcommand{\AltOnePoHundred}{75}
\newcommand{\AltOnePoOneFifty}{80}
\newcommand{\AltOnePoTwentyFive}{61}
\newcommand{\AltOnePoTwoHundred}{82}
\newcommand{\AltOneRidgeEightHundred}{95}
\newcommand{\AltOneRidgeFifty}{80}
\newcommand{\AltOneRidgeFourHundred}{94}
\newcommand{\AltOneRidgeHundred}{86}
\newcommand{\AltOneRidgeOneFifty}{89}
\newcommand{\AltOneRidgePMappedEightHundred}{93}
\newcommand{\AltOneRidgePMappedFifty}{80}
\newcommand{\AltOneRidgePMappedFourHundred}{91}
\newcommand{\AltOneRidgePMappedHundred}{85}
\newcommand{\AltOneRidgePMappedOneFifty}{88}
\newcommand{\AltOneRidgePMappedTwentyFive}{72}
\newcommand{\AltOneRidgePMappedTwoHundred}{89}
\newcommand{\AltOneRidgePPercondEightHundred}{95}
\newcommand{\AltOneRidgePPercondFifty}{76}
\newcommand{\AltOneRidgePPercondFourHundred}{93}
\newcommand{\AltOneRidgePPercondHundred}{85}
\newcommand{\AltOneRidgePPercondOneFifty}{88}
\newcommand{\AltOneRidgePPercondTwentyFive}{57}
\newcommand{\AltOneRidgePPercondTwoHundred}{89}
\newcommand{\AltOneRidgePPooledEightHundred}{88}
\newcommand{\AltOneRidgePPooledFifty}{76}
\newcommand{\AltOneRidgePPooledFourHundred}{86}
\newcommand{\AltOneRidgePPooledHundred}{80}
\newcommand{\AltOneRidgePPooledOneFifty}{83}
\newcommand{\AltOneRidgePPooledTwentyFive}{69}
\newcommand{\AltOneRidgePPooledTwoHundred}{84}
\newcommand{\AltOneRidgeTwentyFive}{72}
\newcommand{\AltOneRidgeTwoHundred}{91}
\newcommand{\AltOneRidgeUMappedEightHundred}{93}
\newcommand{\AltOneRidgeUMappedFifty}{80}
\newcommand{\AltOneRidgeUMappedFourHundred}{91}
\newcommand{\AltOneRidgeUMappedHundred}{85}
\newcommand{\AltOneRidgeUMappedOneFifty}{88}
\newcommand{\AltOneRidgeUMappedTwentyFive}{72}
\newcommand{\AltOneRidgeUMappedTwoHundred}{89}
\newcommand{\AltOneRidgeUPercondEightHundred}{95}
\newcommand{\AltOneRidgeUPercondFifty}{76}
\newcommand{\AltOneRidgeUPercondFourHundred}{94}
\newcommand{\AltOneRidgeUPercondHundred}{86}
\newcommand{\AltOneRidgeUPercondOneFifty}{89}
\newcommand{\AltOneRidgeUPercondTwentyFive}{57}
\newcommand{\AltOneRidgeUPercondTwoHundred}{91}
\newcommand{\AltOneRidgeUPooledEightHundred}{88}
\newcommand{\AltOneRidgeUPooledFifty}{76}
\newcommand{\AltOneRidgeUPooledFourHundred}{87}
\newcommand{\AltOneRidgeUPooledHundred}{81}
\newcommand{\AltOneRidgeUPooledOneFifty}{83}
\newcommand{\AltOneRidgeUPooledTwentyFive}{70}
\newcommand{\AltOneRidgeUPooledTwoHundred}{84}
\newcommand{\AltOneSemiEightHundred}{92}
\newcommand{\AltOneSemiFifty}{69}
\newcommand{\AltOneSemiFourHundred}{87}
\newcommand{\AltOneSemiHundred}{76}
\newcommand{\AltOneSemiMappedEightHundred}{89}
\newcommand{\AltOneSemiMappedFifty}{67}
\newcommand{\AltOneSemiMappedFourHundred}{87}
\newcommand{\AltOneSemiMappedHundred}{73}
\newcommand{\AltOneSemiMappedOneFifty}{79}
\newcommand{\AltOneSemiMappedTwentyFive}{60}
\newcommand{\AltOneSemiMappedTwoHundred}{83}
\newcommand{\AltOneSemiOneFifty}{80}
\newcommand{\AltOneSemiPercondEightHundred}{92}
\newcommand{\AltOneSemiPercondFifty}{69}
\newcommand{\AltOneSemiPercondFourHundred}{87}
\newcommand{\AltOneSemiPercondHundred}{76}
\newcommand{\AltOneSemiPercondOneFifty}{80}
\newcommand{\AltOneSemiPercondTwentyFive}{56}
\newcommand{\AltOneSemiPercondTwoHundred}{81}
\newcommand{\AltOneSemiPooledEightHundred}{85}
\newcommand{\AltOneSemiPooledFifty}{65}
\newcommand{\AltOneSemiPooledFourHundred}{83}
\newcommand{\AltOneSemiPooledHundred}{72}
\newcommand{\AltOneSemiPooledOneFifty}{76}
\newcommand{\AltOneSemiPooledTwentyFive}{59}
\newcommand{\AltOneSemiPooledTwoHundred}{79}
\newcommand{\AltOneSemiTwentyFive}{60}
\newcommand{\AltOneSemiTwoHundred}{83}
\newcommand{\AltRatioChamp}{1.23}
\newcommand{\AltRatioChampHi}{1.44}
\newcommand{\AltRatioChampLo}{0.98}
\newcommand{\AltRatioPoOverAtlas}{at least 5.4}
\newcommand{\AltRatioPoOverAtlasHi}{6.4}
\newcommand{\AltRatioPoOverAtlasLo}{4.3}
\newcommand{\AltRatioPoOverRidge}{at least 4.7}
\newcommand{\AltRatioPoOverRidgeHi}{6.1}
\newcommand{\AltRatioPoOverRidgeLo}{3.4}
\newcommand{\AltRatioPoOverSemi}{at least 1.61}
\newcommand{\AltRatioPoOverSemiHi}{2.6}
\newcommand{\AltRatioPoOverSemiLo}{1.00}
\newcommand{\AltRatioRidge}{1.15}
\newcommand{\AltRatioRidgeHi}{1.36}
\newcommand{\AltRatioRidgeLo}{1.003}
\newcommand{\AltRatioRidgeOverOwn}{1.54}
\newcommand{\AltRatioRidgeOverOwnHi}{1.87}
\newcommand{\AltRatioRidgeOverOwnLo}{1.27}
\newcommand{\AltRatioSemi}{3.4}
\newcommand{\AltRatioSemiHi}{4.6}
\newcommand{\AltRatioSemiLo}{2.4}
\newcommand{\AltRatioSemiOverOwn}{4.5}
\newcommand{\AltRatioSemiOverOwnHi}{5.8}
\newcommand{\AltRatioSemiOverOwnLo}{3.4}
\newcommand{\AltReproduces}{$3.3\times10^{-6}$}
\newcommand{\AltRfAtlasEightHundred}{48}
\newcommand{\AltRfAtlasFifty}{11}
\newcommand{\AltRfAtlasFourHundred}{42}
\newcommand{\AltRfAtlasHundred}{23}
\newcommand{\AltRfAtlasOneFifty}{29}
\newcommand{\AltRfAtlasTwentyFive}{0.5}
\newcommand{\AltRfAtlasTwoHundred}{33}
\newcommand{\AltRfChampAtlasEightHundred}{48}
\newcommand{\AltRfChampAtlasFifty}{11}
\newcommand{\AltRfChampAtlasFourHundred}{42}
\newcommand{\AltRfChampAtlasHundred}{22}
\newcommand{\AltRfChampAtlasTwentyFive}{\ensuremath{-}1.3}
\newcommand{\AltRfChampAtlasTwoHundred}{33}
\newcommand{\AltRfChampEightHundred}{40}
\newcommand{\AltRfChampFifty}{11}
\newcommand{\AltRfChampFourHundred}{35}
\newcommand{\AltRfChampHundred}{20}
\newcommand{\AltRfChampTwentyFive}{0.0}
\newcommand{\AltRfChampTwoHundred}{30}
\newcommand{\AltRfMaxAtlas}{48}
\newcommand{\AltRfMaxChamp}{40}
\newcommand{\AltRfMaxChampAtlas}{48}
\newcommand{\AltRfMaxOther}{57}
\newcommand{\AltRfMaxOwn}{45}
\newcommand{\AltRfMaxPo}{14}
\newcommand{\AltRfMaxRidge}{45}
\newcommand{\AltRfMaxRidgePMapped}{41}
\newcommand{\AltRfMaxRidgePPercond}{43}
\newcommand{\AltRfMaxRidgePPooled}{18}
\newcommand{\AltRfMaxRidgeUMapped}{45}
\newcommand{\AltRfMaxRidgeUPercond}{42}
\newcommand{\AltRfMaxRidgeUPooled}{18}
\newcommand{\AltRfMaxSemi}{32}
\newcommand{\AltRfMaxSemiMapped}{32}
\newcommand{\AltRfMaxSemiPercond}{28}
\newcommand{\AltRfMaxSemiPooled}{11}
\newcommand{\AltRfOtherEightHundred}{57}
\newcommand{\AltRfOtherFifty}{23}
\newcommand{\AltRfOtherFourHundred}{50}
\newcommand{\AltRfOtherHundred}{34}
\newcommand{\AltRfOtherOneFifty}{39}
\newcommand{\AltRfOtherTwentyFive}{8.8}
\newcommand{\AltRfOtherTwoHundred}{43}
\newcommand{\AltRfOwnEightHundred}{45}
\newcommand{\AltRfOwnFifty}{19}
\newcommand{\AltRfOwnFourHundred}{40}
\newcommand{\AltRfOwnHundred}{27}
\newcommand{\AltRfOwnOneFifty}{31}
\newcommand{\AltRfOwnTwentyFive}{5.6}
\newcommand{\AltRfOwnTwoHundred}{34}
\newcommand{\AltRfPoEightHundred}{13}
\newcommand{\AltRfPoFifty}{2.4}
\newcommand{\AltRfPoFourHundred}{14}
\newcommand{\AltRfPoHundred}{3.5}
\newcommand{\AltRfPoOneFifty}{5.1}
\newcommand{\AltRfPoTwentyFive}{0.5}
\newcommand{\AltRfPoTwoHundred}{7.7}
\newcommand{\AltRfRidgeEightHundred}{45}
\newcommand{\AltRfRidgeFifty}{14}
\newcommand{\AltRfRidgeFourHundred}{38}
\newcommand{\AltRfRidgeHundred}{22}
\newcommand{\AltRfRidgeOneFifty}{28}
\newcommand{\AltRfRidgePMappedEightHundred}{41}
\newcommand{\AltRfRidgePMappedFifty}{11}
\newcommand{\AltRfRidgePMappedFourHundred}{36}
\newcommand{\AltRfRidgePMappedHundred}{19}
\newcommand{\AltRfRidgePMappedOneFifty}{23}
\newcommand{\AltRfRidgePMappedTwentyFive}{2.1}
\newcommand{\AltRfRidgePMappedTwoHundred}{26}
\newcommand{\AltRfRidgePPercondEightHundred}{43}
\newcommand{\AltRfRidgePPercondFifty}{14}
\newcommand{\AltRfRidgePPercondFourHundred}{37}
\newcommand{\AltRfRidgePPercondHundred}{22}
\newcommand{\AltRfRidgePPercondOneFifty}{26}
\newcommand{\AltRfRidgePPercondTwentyFive}{3.2}
\newcommand{\AltRfRidgePPercondTwoHundred}{28}
\newcommand{\AltRfRidgePPooledEightHundred}{18}
\newcommand{\AltRfRidgePPooledFifty}{5.2}
\newcommand{\AltRfRidgePPooledFourHundred}{16}
\newcommand{\AltRfRidgePPooledHundred}{8.6}
\newcommand{\AltRfRidgePPooledOneFifty}{10}
\newcommand{\AltRfRidgePPooledTwentyFive}{0.7}
\newcommand{\AltRfRidgePPooledTwoHundred}{12}
\newcommand{\AltRfRidgeTwentyFive}{5.6}
\newcommand{\AltRfRidgeTwoHundred}{30}
\newcommand{\AltRfRidgeUMappedEightHundred}{45}
\newcommand{\AltRfRidgeUMappedFifty}{11}
\newcommand{\AltRfRidgeUMappedFourHundred}{38}
\newcommand{\AltRfRidgeUMappedHundred}{20}
\newcommand{\AltRfRidgeUMappedOneFifty}{25}
\newcommand{\AltRfRidgeUMappedTwentyFive}{5.6}
\newcommand{\AltRfRidgeUMappedTwoHundred}{27}
\newcommand{\AltRfRidgeUPercondEightHundred}{42}
\newcommand{\AltRfRidgeUPercondFifty}{11}
\newcommand{\AltRfRidgeUPercondFourHundred}{35}
\newcommand{\AltRfRidgeUPercondHundred}{20}
\newcommand{\AltRfRidgeUPercondOneFifty}{28}
\newcommand{\AltRfRidgeUPercondTwentyFive}{2.6}
\newcommand{\AltRfRidgeUPercondTwoHundred}{30}
\newcommand{\AltRfRidgeUPooledEightHundred}{18}
\newcommand{\AltRfRidgeUPooledFifty}{5.2}
\newcommand{\AltRfRidgeUPooledFourHundred}{15}
\newcommand{\AltRfRidgeUPooledHundred}{9.2}
\newcommand{\AltRfRidgeUPooledOneFifty}{11}
\newcommand{\AltRfRidgeUPooledTwentyFive}{2.6}
\newcommand{\AltRfRidgeUPooledTwoHundred}{12}
\newcommand{\AltRfSemiEightHundred}{32}
\newcommand{\AltRfSemiFifty}{11}
\newcommand{\AltRfSemiFourHundred}{27}
\newcommand{\AltRfSemiHundred}{17}
\newcommand{\AltRfSemiMappedEightHundred}{32}
\newcommand{\AltRfSemiMappedFifty}{\ensuremath{-}0.2}
\newcommand{\AltRfSemiMappedFourHundred}{24}
\newcommand{\AltRfSemiMappedHundred}{3.6}
\newcommand{\AltRfSemiMappedOneFifty}{7.4}
\newcommand{\AltRfSemiMappedTwentyFive}{1.1}
\newcommand{\AltRfSemiMappedTwoHundred}{13}
\newcommand{\AltRfSemiOneFifty}{20}
\newcommand{\AltRfSemiPercondEightHundred}{28}
\newcommand{\AltRfSemiPercondFifty}{11}
\newcommand{\AltRfSemiPercondFourHundred}{27}
\newcommand{\AltRfSemiPercondHundred}{17}
\newcommand{\AltRfSemiPercondOneFifty}{20}
\newcommand{\AltRfSemiPercondTwentyFive}{3.1}
\newcommand{\AltRfSemiPercondTwoHundred}{20}
\newcommand{\AltRfSemiPooledEightHundred}{11}
\newcommand{\AltRfSemiPooledFifty}{0.3}
\newcommand{\AltRfSemiPooledFourHundred}{8.4}
\newcommand{\AltRfSemiPooledHundred}{1.8}
\newcommand{\AltRfSemiPooledOneFifty}{2.9}
\newcommand{\AltRfSemiPooledTwentyFive}{0.6}
\newcommand{\AltRfSemiPooledTwoHundred}{4.5}
\newcommand{\AltRfSemiTwentyFive}{3.1}
\newcommand{\AltRfSemiTwoHundred}{20}
\newcommand{\AltTarget}{28}
\newcommand{\AltThreeAtlasEightHundred}{73}
\newcommand{\AltThreeAtlasFifty}{24}
\newcommand{\AltThreeAtlasFourHundred}{64}
\newcommand{\AltThreeAtlasHundred}{39}
\newcommand{\AltThreeAtlasOneFifty}{47}
\newcommand{\AltThreeAtlasTwentyFive}{13}
\newcommand{\AltThreeAtlasTwoHundred}{53}
\newcommand{\AltThreeOtherEightHundred}{74}
\newcommand{\AltThreeOtherFifty}{24}
\newcommand{\AltThreeOtherFourHundred}{63}
\newcommand{\AltThreeOtherHundred}{36}
\newcommand{\AltThreeOtherOneFifty}{44}
\newcommand{\AltThreeOtherTwentyFive}{14}
\newcommand{\AltThreeOtherTwoHundred}{50}
\newcommand{\AltThreeOwnEightHundred}{67}
\newcommand{\AltThreeOwnFifty}{23}
\newcommand{\AltThreeOwnFourHundred}{57}
\newcommand{\AltThreeOwnHundred}{34}
\newcommand{\AltThreeOwnOneFifty}{40}
\newcommand{\AltThreeOwnTwentyFive}{14}
\newcommand{\AltThreeOwnTwoHundred}{45}
\newcommand{\AltThreePoEightHundred}{65}
\newcommand{\AltThreePoFifty}{11}
\newcommand{\AltThreePoFourHundred}{54}
\newcommand{\AltThreePoHundred}{25}
\newcommand{\AltThreePoOneFifty}{33}
\newcommand{\AltThreePoTwentyFive}{4.7}
\newcommand{\AltThreePoTwoHundred}{39}
\newcommand{\AltThreeRidgeEightHundred}{56}
\newcommand{\AltThreeRidgeFifty}{17}
\newcommand{\AltThreeRidgeFourHundred}{53}
\newcommand{\AltThreeRidgeHundred}{28}
\newcommand{\AltThreeRidgeOneFifty}{33}
\newcommand{\AltThreeRidgePMappedEightHundred}{52}
\newcommand{\AltThreeRidgePMappedFifty}{17}
\newcommand{\AltThreeRidgePMappedFourHundred}{49}
\newcommand{\AltThreeRidgePMappedHundred}{28}
\newcommand{\AltThreeRidgePMappedOneFifty}{33}
\newcommand{\AltThreeRidgePMappedTwentyFive}{10}
\newcommand{\AltThreeRidgePMappedTwoHundred}{40}
\newcommand{\AltThreeRidgePPercondEightHundred}{51}
\newcommand{\AltThreeRidgePPercondFifty}{14}
\newcommand{\AltThreeRidgePPercondFourHundred}{45}
\newcommand{\AltThreeRidgePPercondHundred}{26}
\newcommand{\AltThreeRidgePPercondOneFifty}{32}
\newcommand{\AltThreeRidgePPercondTwentyFive}{3.3}
\newcommand{\AltThreeRidgePPercondTwoHundred}{36}
\newcommand{\AltThreeRidgePPooledEightHundred}{46}
\newcommand{\AltThreeRidgePPooledFifty}{15}
\newcommand{\AltThreeRidgePPooledFourHundred}{43}
\newcommand{\AltThreeRidgePPooledHundred}{24}
\newcommand{\AltThreeRidgePPooledOneFifty}{29}
\newcommand{\AltThreeRidgePPooledTwentyFive}{9.0}
\newcommand{\AltThreeRidgePPooledTwoHundred}{35}
\newcommand{\AltThreeRidgeTwentyFive}{10}
\newcommand{\AltThreeRidgeTwoHundred}{40}
\newcommand{\AltThreeRidgeUMappedEightHundred}{56}
\newcommand{\AltThreeRidgeUMappedFifty}{13}
\newcommand{\AltThreeRidgeUMappedFourHundred}{53}
\newcommand{\AltThreeRidgeUMappedHundred}{24}
\newcommand{\AltThreeRidgeUMappedOneFifty}{30}
\newcommand{\AltThreeRidgeUMappedTwentyFive}{7.0}
\newcommand{\AltThreeRidgeUMappedTwoHundred}{34}
\newcommand{\AltThreeRidgeUPercondEightHundred}{56}
\newcommand{\AltThreeRidgeUPercondFifty}{13}
\newcommand{\AltThreeRidgeUPercondFourHundred}{49}
\newcommand{\AltThreeRidgeUPercondHundred}{26}
\newcommand{\AltThreeRidgeUPercondOneFifty}{32}
\newcommand{\AltThreeRidgeUPercondTwentyFive}{2.7}
\newcommand{\AltThreeRidgeUPercondTwoHundred}{35}
\newcommand{\AltThreeRidgeUPooledEightHundred}{48}
\newcommand{\AltThreeRidgeUPooledFifty}{11}
\newcommand{\AltThreeRidgeUPooledFourHundred}{47}
\newcommand{\AltThreeRidgeUPooledHundred}{21}
\newcommand{\AltThreeRidgeUPooledOneFifty}{26}
\newcommand{\AltThreeRidgeUPooledTwentyFive}{6.0}
\newcommand{\AltThreeRidgeUPooledTwoHundred}{29}
\newcommand{\AltThreeSemiEightHundred}{46}
\newcommand{\AltThreeSemiFifty}{5.0}
\newcommand{\AltThreeSemiFourHundred}{34}
\newcommand{\AltThreeSemiHundred}{13}
\newcommand{\AltThreeSemiMappedEightHundred}{38}
\newcommand{\AltThreeSemiMappedFifty}{4.6}
\newcommand{\AltThreeSemiMappedFourHundred}{34}
\newcommand{\AltThreeSemiMappedHundred}{13}
\newcommand{\AltThreeSemiMappedOneFifty}{21}
\newcommand{\AltThreeSemiMappedTwentyFive}{1.1}
\newcommand{\AltThreeSemiMappedTwoHundred}{26}
\newcommand{\AltThreeSemiOneFifty}{21}
\newcommand{\AltThreeSemiPercondEightHundred}{46}
\newcommand{\AltThreeSemiPercondFifty}{5.0}
\newcommand{\AltThreeSemiPercondFourHundred}{32}
\newcommand{\AltThreeSemiPercondHundred}{9.8}
\newcommand{\AltThreeSemiPercondOneFifty}{14}
\newcommand{\AltThreeSemiPercondTwentyFive}{1.2}
\newcommand{\AltThreeSemiPercondTwoHundred}{19}
\newcommand{\AltThreeSemiPooledEightHundred}{33}
\newcommand{\AltThreeSemiPooledFifty}{3.7}
\newcommand{\AltThreeSemiPooledFourHundred}{31}
\newcommand{\AltThreeSemiPooledHundred}{11}
\newcommand{\AltThreeSemiPooledOneFifty}{18}
\newcommand{\AltThreeSemiPooledTwentyFive}{0.8}
\newcommand{\AltThreeSemiPooledTwoHundred}{23}
\newcommand{\AltThreeSemiTwentyFive}{1.2}
\newcommand{\AltThreeSemiTwoHundred}{26}
\newcommand{\AltTwoAtlasEightHundred}{67}
\newcommand{\AltTwoAtlasFifty}{35}
\newcommand{\AltTwoAtlasFourHundred}{62}
\newcommand{\AltTwoAtlasHundred}{47}
\newcommand{\AltTwoAtlasOneFifty}{52}
\newcommand{\AltTwoAtlasTwentyFive}{23}
\newcommand{\AltTwoAtlasTwoHundred}{55}
\newcommand{\AltTwoChampAtlasEightHundred}{67}
\newcommand{\AltTwoChampAtlasFifty}{36}
\newcommand{\AltTwoChampAtlasFourHundred}{62}
\newcommand{\AltTwoChampAtlasHundred}{47}
\newcommand{\AltTwoChampAtlasTwentyFive}{23}
\newcommand{\AltTwoChampAtlasTwoHundred}{55}
\newcommand{\AltTwoChampEightHundred}{64}
\newcommand{\AltTwoChampFifty}{28}
\newcommand{\AltTwoChampFourHundred}{58}
\newcommand{\AltTwoChampHundred}{42}
\newcommand{\AltTwoChampTwentyFive}{11}
\newcommand{\AltTwoChampTwoHundred}{49}
\newcommand{\AltTwoOtherEightHundred}{67}
\newcommand{\AltTwoOtherFifty}{38}
\newcommand{\AltTwoOtherFourHundred}{61}
\newcommand{\AltTwoOtherHundred}{46}
\newcommand{\AltTwoOtherOneFifty}{50}
\newcommand{\AltTwoOtherTwentyFive}{28}
\newcommand{\AltTwoOtherTwoHundred}{54}
\newcommand{\AltTwoOwnEightHundred}{57}
\newcommand{\AltTwoOwnFifty}{36}
\newcommand{\AltTwoOwnFourHundred}{52}
\newcommand{\AltTwoOwnHundred}{41}
\newcommand{\AltTwoOwnOneFifty}{44}
\newcommand{\AltTwoOwnTwentyFive}{27}
\newcommand{\AltTwoOwnTwoHundred}{46}
\newcommand{\AltTwoPoEightHundred}{61}
\newcommand{\AltTwoPoFifty}{15}
\newcommand{\AltTwoPoFourHundred}{52}
\newcommand{\AltTwoPoHundred}{28}
\newcommand{\AltTwoPoOneFifty}{35}
\newcommand{\AltTwoPoTwentyFive}{6.9}
\newcommand{\AltTwoPoTwoHundred}{40}
\newcommand{\AltTwoRidgeEightHundred}{71}
\newcommand{\AltTwoRidgeFifty}{32}
\newcommand{\AltTwoRidgeFourHundred}{67}
\newcommand{\AltTwoRidgeHundred}{50}
\newcommand{\AltTwoRidgeOneFifty}{57}
\newcommand{\AltTwoRidgePMappedEightHundred}{63}
\newcommand{\AltTwoRidgePMappedFifty}{32}
\newcommand{\AltTwoRidgePMappedFourHundred}{58}
\newcommand{\AltTwoRidgePMappedHundred}{43}
\newcommand{\AltTwoRidgePMappedOneFifty}{48}
\newcommand{\AltTwoRidgePMappedTwentyFive}{21}
\newcommand{\AltTwoRidgePMappedTwoHundred}{51}
\newcommand{\AltTwoRidgePPercondEightHundred}{69}
\newcommand{\AltTwoRidgePPercondFifty}{32}
\newcommand{\AltTwoRidgePPercondFourHundred}{65}
\newcommand{\AltTwoRidgePPercondHundred}{48}
\newcommand{\AltTwoRidgePPercondOneFifty}{55}
\newcommand{\AltTwoRidgePPercondTwentyFive}{9.3}
\newcommand{\AltTwoRidgePPercondTwoHundred}{58}
\newcommand{\AltTwoRidgePPooledEightHundred}{33}
\newcommand{\AltTwoRidgePPooledFifty}{23}
\newcommand{\AltTwoRidgePPooledFourHundred}{32}
\newcommand{\AltTwoRidgePPooledHundred}{28}
\newcommand{\AltTwoRidgePPooledOneFifty}{30}
\newcommand{\AltTwoRidgePPooledTwentyFive}{17}
\newcommand{\AltTwoRidgePPooledTwoHundred}{31}
\newcommand{\AltTwoRidgeTwentyFive}{21}
\newcommand{\AltTwoRidgeTwoHundred}{60}
\newcommand{\AltTwoRidgeUMappedEightHundred}{64}
\newcommand{\AltTwoRidgeUMappedFifty}{32}
\newcommand{\AltTwoRidgeUMappedFourHundred}{60}
\newcommand{\AltTwoRidgeUMappedHundred}{44}
\newcommand{\AltTwoRidgeUMappedOneFifty}{49}
\newcommand{\AltTwoRidgeUMappedTwentyFive}{21}
\newcommand{\AltTwoRidgeUMappedTwoHundred}{52}
\newcommand{\AltTwoRidgeUPercondEightHundred}{71}
\newcommand{\AltTwoRidgeUPercondFifty}{32}
\newcommand{\AltTwoRidgeUPercondFourHundred}{67}
\newcommand{\AltTwoRidgeUPercondHundred}{50}
\newcommand{\AltTwoRidgeUPercondOneFifty}{57}
\newcommand{\AltTwoRidgeUPercondTwentyFive}{9.5}
\newcommand{\AltTwoRidgeUPercondTwoHundred}{60}
\newcommand{\AltTwoRidgeUPooledEightHundred}{32}
\newcommand{\AltTwoRidgeUPooledFifty}{23}
\newcommand{\AltTwoRidgeUPooledFourHundred}{32}
\newcommand{\AltTwoRidgeUPooledHundred}{28}
\newcommand{\AltTwoRidgeUPooledOneFifty}{29}
\newcommand{\AltTwoRidgeUPooledTwentyFive}{17}
\newcommand{\AltTwoRidgeUPooledTwoHundred}{30}
\newcommand{\AltTwoSemiEightHundred}{63}
\newcommand{\AltTwoSemiFifty}{23}
\newcommand{\AltTwoSemiFourHundred}{52}
\newcommand{\AltTwoSemiHundred}{32}
\newcommand{\AltTwoSemiMappedEightHundred}{49}
\newcommand{\AltTwoSemiMappedFifty}{16}
\newcommand{\AltTwoSemiMappedFourHundred}{45}
\newcommand{\AltTwoSemiMappedHundred}{22}
\newcommand{\AltTwoSemiMappedOneFifty}{29}
\newcommand{\AltTwoSemiMappedTwentyFive}{13}
\newcommand{\AltTwoSemiMappedTwoHundred}{36}
\newcommand{\AltTwoSemiOneFifty}{37}
\newcommand{\AltTwoSemiPercondEightHundred}{63}
\newcommand{\AltTwoSemiPercondFifty}{23}
\newcommand{\AltTwoSemiPercondFourHundred}{52}
\newcommand{\AltTwoSemiPercondHundred}{32}
\newcommand{\AltTwoSemiPercondOneFifty}{37}
\newcommand{\AltTwoSemiPercondTwentyFive}{8.2}
\newcommand{\AltTwoSemiPercondTwoHundred}{41}
\newcommand{\AltTwoSemiPooledEightHundred}{33}
\newcommand{\AltTwoSemiPooledFifty}{14}
\newcommand{\AltTwoSemiPooledFourHundred}{32}
\newcommand{\AltTwoSemiPooledHundred}{20}
\newcommand{\AltTwoSemiPooledOneFifty}{26}
\newcommand{\AltTwoSemiPooledTwentyFive}{10}
\newcommand{\AltTwoSemiPooledTwoHundred}{29}
\newcommand{\AltTwoSemiTwentyFive}{13}
\newcommand{\AltTwoSemiTwoHundred}{41}
\newcommand{\BoBOne}{met}
\newcommand{\BoBThree}{met}
\newcommand{\BoBTwo}{met}
\newcommand{\BoCellsOneAtlas}{29}
\newcommand{\BoCellsOneOther}{$\leq$25}
\newcommand{\BoCellsOneOwn}{$\leq$25}
\newcommand{\BoCellsOnePo}{84}
\newcommand{\BoCellsTwoAtlas}{45}
\newcommand{\BoCellsTwoOther}{37}
\newcommand{\BoCellsTwoOwn}{42}
\newcommand{\BoCellsTwoPo}{136}
\newcommand{\BoCogAtlasEightHundred}{99}
\newcommand{\BoCogAtlasEightHundredHi}{100}
\newcommand{\BoCogAtlasEightHundredLo}{99}
\newcommand{\BoCogAtlasFifty}{85}
\newcommand{\BoCogAtlasFiftyHi}{87}
\newcommand{\BoCogAtlasFiftyLo}{83}
\newcommand{\BoCogAtlasFourHundred}{99}
\newcommand{\BoCogAtlasFourHundredHi}{99}
\newcommand{\BoCogAtlasFourHundredLo}{98}
\newcommand{\BoCogAtlasHundred}{93}
\newcommand{\BoCogAtlasHundredHi}{94}
\newcommand{\BoCogAtlasHundredLo}{92}
\newcommand{\BoCogAtlasOneFifty}{96}
\newcommand{\BoCogAtlasOneFiftyHi}{96}
\newcommand{\BoCogAtlasOneFiftyLo}{95}
\newcommand{\BoCogAtlasTwentyFive}{76}
\newcommand{\BoCogAtlasTwentyFiveHi}{78}
\newcommand{\BoCogAtlasTwentyFiveLo}{75}
\newcommand{\BoCogAtlasTwoHundred}{97}
\newcommand{\BoCogAtlasTwoHundredHi}{97}
\newcommand{\BoCogAtlasTwoHundredLo}{96}
\newcommand{\BoCogOtherEightHundred}{99}
\newcommand{\BoCogOtherEightHundredHi}{100}
\newcommand{\BoCogOtherEightHundredLo}{99}
\newcommand{\BoCogOtherFifty}{87}
\newcommand{\BoCogOtherFiftyHi}{88}
\newcommand{\BoCogOtherFiftyLo}{86}
\newcommand{\BoCogOtherFourHundred}{99}
\newcommand{\BoCogOtherFourHundredHi}{99}
\newcommand{\BoCogOtherFourHundredLo}{98}
\newcommand{\BoCogOtherHundred}{93}
\newcommand{\BoCogOtherHundredHi}{94}
\newcommand{\BoCogOtherHundredLo}{92}
\newcommand{\BoCogOtherOneFifty}{96}
\newcommand{\BoCogOtherOneFiftyHi}{96}
\newcommand{\BoCogOtherOneFiftyLo}{95}
\newcommand{\BoCogOtherTwentyFive}{78}
\newcommand{\BoCogOtherTwentyFiveHi}{80}
\newcommand{\BoCogOtherTwentyFiveLo}{77}
\newcommand{\BoCogOtherTwoHundred}{97}
\newcommand{\BoCogOtherTwoHundredHi}{98}
\newcommand{\BoCogOtherTwoHundredLo}{96}
\newcommand{\BoCogOwnEightHundred}{99}
\newcommand{\BoCogOwnEightHundredHi}{100}
\newcommand{\BoCogOwnEightHundredLo}{99}
\newcommand{\BoCogOwnFifty}{85}
\newcommand{\BoCogOwnFiftyHi}{87}
\newcommand{\BoCogOwnFiftyLo}{84}
\newcommand{\BoCogOwnFourHundred}{98}
\newcommand{\BoCogOwnFourHundredHi}{99}
\newcommand{\BoCogOwnFourHundredLo}{98}
\newcommand{\BoCogOwnHundred}{92}
\newcommand{\BoCogOwnHundredHi}{93}
\newcommand{\BoCogOwnHundredLo}{90}
\newcommand{\BoCogOwnOneFifty}{94}
\newcommand{\BoCogOwnOneFiftyHi}{95}
\newcommand{\BoCogOwnOneFiftyLo}{93}
\newcommand{\BoCogOwnTwentyFive}{78}
\newcommand{\BoCogOwnTwentyFiveHi}{80}
\newcommand{\BoCogOwnTwentyFiveLo}{76}
\newcommand{\BoCogOwnTwoHundred}{96}
\newcommand{\BoCogOwnTwoHundredHi}{97}
\newcommand{\BoCogOwnTwoHundredLo}{95}
\newcommand{\BoCogPoEightHundred}{99}
\newcommand{\BoCogPoEightHundredHi}{100}
\newcommand{\BoCogPoEightHundredLo}{99}
\newcommand{\BoCogPoFifty}{90}
\newcommand{\BoCogPoFiftyHi}{92}
\newcommand{\BoCogPoFiftyLo}{87}
\newcommand{\BoCogPoFourHundred}{99}
\newcommand{\BoCogPoFourHundredHi}{99}
\newcommand{\BoCogPoFourHundredLo}{99}
\newcommand{\BoCogPoHundred}{95}
\newcommand{\BoCogPoHundredHi}{96}
\newcommand{\BoCogPoHundredLo}{94}
\newcommand{\BoCogPoOneFifty}{97}
\newcommand{\BoCogPoOneFiftyHi}{98}
\newcommand{\BoCogPoOneFiftyLo}{96}
\newcommand{\BoCogPoTwentyFive}{82}
\newcommand{\BoCogPoTwentyFiveHi}{84}
\newcommand{\BoCogPoTwentyFiveLo}{80}
\newcommand{\BoCogPoTwoHundred}{98}
\newcommand{\BoCogPoTwoHundredHi}{98}
\newcommand{\BoCogPoTwoHundredLo}{97}
\newcommand{\BoCompared}{11,220}
\newcommand{\BoDiffOneFifty}{12.9}
\newcommand{\BoDiffOneFiftyHi}{13.7}
\newcommand{\BoDiffOneFiftyLo}{12.1}
\newcommand{\BoDiffOneHundred}{11.2}
\newcommand{\BoDiffOneHundredHi}{12.0}
\newcommand{\BoDiffOneHundredLo}{10.3}
\newcommand{\BoDiffOneOneFifty}{9.1}
\newcommand{\BoDiffOneOneFiftyHi}{9.8}
\newcommand{\BoDiffOneOneFiftyLo}{8.4}
\newcommand{\BoDiffThreeHundred}{14.2}
\newcommand{\BoDiffThreeHundredHi}{15.8}
\newcommand{\BoDiffThreeHundredLo}{12.5}
\newcommand{\BoDiffThreeOneFifty}{14.4}
\newcommand{\BoDiffThreeOneFiftyHi}{15.3}
\newcommand{\BoDiffThreeOneFiftyLo}{12.9}
\newcommand{\BoDiffTwoFifty}{20.1}
\newcommand{\BoDiffTwoFiftyHi}{21.3}
\newcommand{\BoDiffTwoFiftyLo}{18.9}
\newcommand{\BoDiffTwoHundred}{19.3}
\newcommand{\BoDiffTwoHundredHi}{20.5}
\newcommand{\BoDiffTwoHundredLo}{18.0}
\newcommand{\BoDiffTwoOneFifty}{16.8}
\newcommand{\BoDiffTwoOneFiftyHi}{18.0}
\newcommand{\BoDiffTwoOneFiftyLo}{15.6}
\newcommand{\BoFrozenAt}{07:06 EDT on 29 September 2026}
\newcommand{\BoIntegrity}{$1.2\times10^{-7}$}
\newcommand{\BoNCognate}{16}
\newcommand{\BoNPairs}{22,200}
\newcommand{\BoNRep}{120,649}
\newcommand{\BoNUnits}{170}
\newcommand{\BoOneAtlasEightHundred}{94}
\newcommand{\BoOneAtlasEightHundredHi}{95}
\newcommand{\BoOneAtlasEightHundredLo}{93}
\newcommand{\BoOneAtlasFifty}{80}
\newcommand{\BoOneAtlasFiftyHi}{82}
\newcommand{\BoOneAtlasFiftyLo}{79}
\newcommand{\BoOneAtlasFourHundred}{92}
\newcommand{\BoOneAtlasFourHundredHi}{93}
\newcommand{\BoOneAtlasFourHundredLo}{91}
\newcommand{\BoOneAtlasHundred}{86}
\newcommand{\BoOneAtlasHundredHi}{87}
\newcommand{\BoOneAtlasHundredLo}{85}
\newcommand{\BoOneAtlasOneFifty}{89}
\newcommand{\BoOneAtlasOneFiftyHi}{90}
\newcommand{\BoOneAtlasOneFiftyLo}{87}
\newcommand{\BoOneAtlasTwentyFive}{71}
\newcommand{\BoOneAtlasTwentyFiveHi}{72}
\newcommand{\BoOneAtlasTwentyFiveLo}{70}
\newcommand{\BoOneAtlasTwoHundred}{90}
\newcommand{\BoOneAtlasTwoHundredHi}{91}
\newcommand{\BoOneAtlasTwoHundredLo}{89}
\newcommand{\BoOneOtherEightHundred}{96}
\newcommand{\BoOneOtherEightHundredHi}{96}
\newcommand{\BoOneOtherEightHundredLo}{95}
\newcommand{\BoOneOtherFifty}{85}
\newcommand{\BoOneOtherFiftyHi}{86}
\newcommand{\BoOneOtherFiftyLo}{84}
\newcommand{\BoOneOtherFourHundred}{94}
\newcommand{\BoOneOtherFourHundredHi}{95}
\newcommand{\BoOneOtherFourHundredLo}{93}
\newcommand{\BoOneOtherHundred}{89}
\newcommand{\BoOneOtherHundredHi}{90}
\newcommand{\BoOneOtherHundredLo}{88}
\newcommand{\BoOneOtherOneFifty}{91}
\newcommand{\BoOneOtherOneFiftyHi}{92}
\newcommand{\BoOneOtherOneFiftyLo}{90}
\newcommand{\BoOneOtherTwentyFive}{77}
\newcommand{\BoOneOtherTwentyFiveHi}{78}
\newcommand{\BoOneOtherTwentyFiveLo}{75}
\newcommand{\BoOneOtherTwoHundred}{92}
\newcommand{\BoOneOtherTwoHundredHi}{93}
\newcommand{\BoOneOtherTwoHundredLo}{91}
\newcommand{\BoOneOwnEightHundred}{94}
\newcommand{\BoOneOwnEightHundredHi}{95}
\newcommand{\BoOneOwnEightHundredLo}{93}
\newcommand{\BoOneOwnFifty}{84}
\newcommand{\BoOneOwnFiftyHi}{85}
\newcommand{\BoOneOwnFiftyLo}{83}
\newcommand{\BoOneOwnFourHundred}{92}
\newcommand{\BoOneOwnFourHundredHi}{93}
\newcommand{\BoOneOwnFourHundredLo}{92}
\newcommand{\BoOneOwnHundred}{88}
\newcommand{\BoOneOwnHundredHi}{89}
\newcommand{\BoOneOwnHundredLo}{87}
\newcommand{\BoOneOwnOneFifty}{89}
\newcommand{\BoOneOwnOneFiftyHi}{90}
\newcommand{\BoOneOwnOneFiftyLo}{88}
\newcommand{\BoOneOwnTwentyFive}{76}
\newcommand{\BoOneOwnTwentyFiveHi}{77}
\newcommand{\BoOneOwnTwentyFiveLo}{74}
\newcommand{\BoOneOwnTwoHundred}{90}
\newcommand{\BoOneOwnTwoHundredHi}{91}
\newcommand{\BoOneOwnTwoHundredLo}{89}
\newcommand{\BoOnePoEightHundred}{92}
\newcommand{\BoOnePoEightHundredHi}{93}
\newcommand{\BoOnePoEightHundredLo}{92}
\newcommand{\BoOnePoFifty}{67}
\newcommand{\BoOnePoFiftyHi}{69}
\newcommand{\BoOnePoFiftyLo}{66}
\newcommand{\BoOnePoFourHundred}{88}
\newcommand{\BoOnePoFourHundredHi}{89}
\newcommand{\BoOnePoFourHundredLo}{87}
\newcommand{\BoOnePoHundred}{75}
\newcommand{\BoOnePoHundredHi}{76}
\newcommand{\BoOnePoHundredLo}{74}
\newcommand{\BoOnePoOneFifty}{80}
\newcommand{\BoOnePoOneFiftyHi}{81}
\newcommand{\BoOnePoOneFiftyLo}{78}
\newcommand{\BoOnePoTwentyFive}{61}
\newcommand{\BoOnePoTwentyFiveHi}{63}
\newcommand{\BoOnePoTwentyFiveLo}{61}
\newcommand{\BoOnePoTwoHundred}{82}
\newcommand{\BoOnePoTwoHundredHi}{83}
\newcommand{\BoOnePoTwoHundredLo}{81}
\newcommand{\BoPerfectTwo}{97}
\newcommand{\BoRandomTwo}{1.6}
\newcommand{\BoRepB}{5,077}
\newcommand{\BoRepCDEight}{56,637}
\newcommand{\BoRepCDFour}{27,388}
\newcommand{\BoRepMono}{21,021}
\newcommand{\BoRepNK}{10,526}
\newcommand{\BoRepShare}{3.2}
\newcommand{\BoSaveOne}{2.9}
\newcommand{\BoSaveOneHi}{3.1}
\newcommand{\BoSaveOneLo}{2.6}
\newcommand{\BoSaveTwo}{3.0}
\newcommand{\BoSaveTwoHi}{3.2}
\newcommand{\BoSaveTwoLo}{2.8}
\newcommand{\BoTargetOne}{73}
\newcommand{\BoTargetTwo}{34}
\newcommand{\BoThreeAtlasEightHundred}{73}
\newcommand{\BoThreeAtlasEightHundredHi}{76}
\newcommand{\BoThreeAtlasEightHundredLo}{71}
\newcommand{\BoThreeAtlasFifty}{24}
\newcommand{\BoThreeAtlasFiftyHi}{27}
\newcommand{\BoThreeAtlasFiftyLo}{22}
\newcommand{\BoThreeAtlasFourHundred}{64}
\newcommand{\BoThreeAtlasFourHundredHi}{67}
\newcommand{\BoThreeAtlasFourHundredLo}{62}
\newcommand{\BoThreeAtlasHundred}{39}
\newcommand{\BoThreeAtlasHundredHi}{41}
\newcommand{\BoThreeAtlasHundredLo}{36}
\newcommand{\BoThreeAtlasOneFifty}{47}
\newcommand{\BoThreeAtlasOneFiftyHi}{50}
\newcommand{\BoThreeAtlasOneFiftyLo}{44}
\newcommand{\BoThreeAtlasTwentyFive}{13}
\newcommand{\BoThreeAtlasTwentyFiveHi}{14}
\newcommand{\BoThreeAtlasTwentyFiveLo}{11}
\newcommand{\BoThreeAtlasTwoHundred}{53}
\newcommand{\BoThreeAtlasTwoHundredHi}{56}
\newcommand{\BoThreeAtlasTwoHundredLo}{50}
\newcommand{\BoThreeOtherEightHundred}{74}
\newcommand{\BoThreeOtherEightHundredHi}{77}
\newcommand{\BoThreeOtherEightHundredLo}{70}
\newcommand{\BoThreeOtherFifty}{24}
\newcommand{\BoThreeOtherFiftyHi}{26}
\newcommand{\BoThreeOtherFiftyLo}{22}
\newcommand{\BoThreeOtherFourHundred}{63}
\newcommand{\BoThreeOtherFourHundredHi}{66}
\newcommand{\BoThreeOtherFourHundredLo}{60}
\newcommand{\BoThreeOtherHundred}{36}
\newcommand{\BoThreeOtherHundredHi}{39}
\newcommand{\BoThreeOtherHundredLo}{33}
\newcommand{\BoThreeOtherOneFifty}{44}
\newcommand{\BoThreeOtherOneFiftyHi}{47}
\newcommand{\BoThreeOtherOneFiftyLo}{40}
\newcommand{\BoThreeOtherTwentyFive}{14}
\newcommand{\BoThreeOtherTwentyFiveHi}{16}
\newcommand{\BoThreeOtherTwentyFiveLo}{13}
\newcommand{\BoThreeOtherTwoHundred}{50}
\newcommand{\BoThreeOtherTwoHundredHi}{53}
\newcommand{\BoThreeOtherTwoHundredLo}{47}
\newcommand{\BoThreeOwnEightHundred}{67}
\newcommand{\BoThreeOwnEightHundredHi}{70}
\newcommand{\BoThreeOwnEightHundredLo}{63}
\newcommand{\BoThreeOwnFifty}{23}
\newcommand{\BoThreeOwnFiftyHi}{25}
\newcommand{\BoThreeOwnFiftyLo}{21}
\newcommand{\BoThreeOwnFourHundred}{57}
\newcommand{\BoThreeOwnFourHundredHi}{60}
\newcommand{\BoThreeOwnFourHundredLo}{53}
\newcommand{\BoThreeOwnHundred}{34}
\newcommand{\BoThreeOwnHundredHi}{36}
\newcommand{\BoThreeOwnHundredLo}{31}
\newcommand{\BoThreeOwnOneFifty}{40}
\newcommand{\BoThreeOwnOneFiftyHi}{43}
\newcommand{\BoThreeOwnOneFiftyLo}{37}
\newcommand{\BoThreeOwnTwentyFive}{14}
\newcommand{\BoThreeOwnTwentyFiveHi}{16}
\newcommand{\BoThreeOwnTwentyFiveLo}{13}
\newcommand{\BoThreeOwnTwoHundred}{45}
\newcommand{\BoThreeOwnTwoHundredHi}{48}
\newcommand{\BoThreeOwnTwoHundredLo}{42}
\newcommand{\BoThreePoEightHundred}{65}
\newcommand{\BoThreePoEightHundredHi}{67}
\newcommand{\BoThreePoEightHundredLo}{61}
\newcommand{\BoThreePoFifty}{11}
\newcommand{\BoThreePoFiftyHi}{13}
\newcommand{\BoThreePoFiftyLo}{9.9}
\newcommand{\BoThreePoFourHundred}{54}
\newcommand{\BoThreePoFourHundredHi}{58}
\newcommand{\BoThreePoFourHundredLo}{50}
\newcommand{\BoThreePoHundred}{25}
\newcommand{\BoThreePoHundredHi}{27}
\newcommand{\BoThreePoHundredLo}{23}
\newcommand{\BoThreePoOneFifty}{33}
\newcommand{\BoThreePoOneFiftyHi}{36}
\newcommand{\BoThreePoOneFiftyLo}{30}
\newcommand{\BoThreePoTwentyFive}{4.7}
\newcommand{\BoThreePoTwentyFiveHi}{5.6}
\newcommand{\BoThreePoTwentyFiveLo}{3.8}
\newcommand{\BoThreePoTwoHundred}{39}
\newcommand{\BoThreePoTwoHundredHi}{42}
\newcommand{\BoThreePoTwoHundredLo}{36}
\newcommand{\BoTwoAtlasEightHundred}{67}
\newcommand{\BoTwoAtlasEightHundredHi}{68}
\newcommand{\BoTwoAtlasEightHundredLo}{65}
\newcommand{\BoTwoAtlasFifty}{35}
\newcommand{\BoTwoAtlasFiftyHi}{37}
\newcommand{\BoTwoAtlasFiftyLo}{34}
\newcommand{\BoTwoAtlasFourHundred}{62}
\newcommand{\BoTwoAtlasFourHundredHi}{63}
\newcommand{\BoTwoAtlasFourHundredLo}{60}
\newcommand{\BoTwoAtlasHundred}{47}
\newcommand{\BoTwoAtlasHundredHi}{49}
\newcommand{\BoTwoAtlasHundredLo}{46}
\newcommand{\BoTwoAtlasOneFifty}{52}
\newcommand{\BoTwoAtlasOneFiftyHi}{54}
\newcommand{\BoTwoAtlasOneFiftyLo}{51}
\newcommand{\BoTwoAtlasTwentyFive}{23}
\newcommand{\BoTwoAtlasTwentyFiveHi}{24}
\newcommand{\BoTwoAtlasTwentyFiveLo}{21}
\newcommand{\BoTwoAtlasTwoHundred}{55}
\newcommand{\BoTwoAtlasTwoHundredHi}{56}
\newcommand{\BoTwoAtlasTwoHundredLo}{54}
\newcommand{\BoTwoOtherEightHundred}{67}
\newcommand{\BoTwoOtherEightHundredHi}{69}
\newcommand{\BoTwoOtherEightHundredLo}{66}
\newcommand{\BoTwoOtherFifty}{38}
\newcommand{\BoTwoOtherFiftyHi}{40}
\newcommand{\BoTwoOtherFiftyLo}{36}
\newcommand{\BoTwoOtherFourHundred}{61}
\newcommand{\BoTwoOtherFourHundredHi}{63}
\newcommand{\BoTwoOtherFourHundredLo}{60}
\newcommand{\BoTwoOtherHundred}{46}
\newcommand{\BoTwoOtherHundredHi}{48}
\newcommand{\BoTwoOtherHundredLo}{44}
\newcommand{\BoTwoOtherOneFifty}{50}
\newcommand{\BoTwoOtherOneFiftyHi}{52}
\newcommand{\BoTwoOtherOneFiftyLo}{49}
\newcommand{\BoTwoOtherTwentyFive}{28}
\newcommand{\BoTwoOtherTwentyFiveHi}{30}
\newcommand{\BoTwoOtherTwentyFiveLo}{27}
\newcommand{\BoTwoOtherTwoHundred}{54}
\newcommand{\BoTwoOtherTwoHundredHi}{55}
\newcommand{\BoTwoOtherTwoHundredLo}{52}
\newcommand{\BoTwoOwnEightHundred}{57}
\newcommand{\BoTwoOwnEightHundredHi}{59}
\newcommand{\BoTwoOwnEightHundredLo}{54}
\newcommand{\BoTwoOwnFifty}{36}
\newcommand{\BoTwoOwnFiftyHi}{37}
\newcommand{\BoTwoOwnFiftyLo}{35}
\newcommand{\BoTwoOwnFourHundred}{52}
\newcommand{\BoTwoOwnFourHundredHi}{54}
\newcommand{\BoTwoOwnFourHundredLo}{50}
\newcommand{\BoTwoOwnHundred}{41}
\newcommand{\BoTwoOwnHundredHi}{43}
\newcommand{\BoTwoOwnHundredLo}{40}
\newcommand{\BoTwoOwnOneFifty}{44}
\newcommand{\BoTwoOwnOneFiftyHi}{46}
\newcommand{\BoTwoOwnOneFiftyLo}{43}
\newcommand{\BoTwoOwnTwentyFive}{27}
\newcommand{\BoTwoOwnTwentyFiveHi}{28}
\newcommand{\BoTwoOwnTwentyFiveLo}{25}
\newcommand{\BoTwoOwnTwoHundred}{46}
\newcommand{\BoTwoOwnTwoHundredHi}{48}
\newcommand{\BoTwoOwnTwoHundredLo}{45}
\newcommand{\BoTwoPoEightHundred}{61}
\newcommand{\BoTwoPoEightHundredHi}{62}
\newcommand{\BoTwoPoEightHundredLo}{60}
\newcommand{\BoTwoPoFifty}{15}
\newcommand{\BoTwoPoFiftyHi}{16}
\newcommand{\BoTwoPoFiftyLo}{14}
\newcommand{\BoTwoPoFourHundred}{52}
\newcommand{\BoTwoPoFourHundredHi}{53}
\newcommand{\BoTwoPoFourHundredLo}{51}
\newcommand{\BoTwoPoHundred}{28}
\newcommand{\BoTwoPoHundredHi}{29}
\newcommand{\BoTwoPoHundredLo}{27}
\newcommand{\BoTwoPoOneFifty}{35}
\newcommand{\BoTwoPoOneFiftyHi}{36}
\newcommand{\BoTwoPoOneFiftyLo}{34}
\newcommand{\BoTwoPoTwentyFive}{6.9}
\newcommand{\BoTwoPoTwentyFiveHi}{7.8}
\newcommand{\BoTwoPoTwentyFiveLo}{6.0}
\newcommand{\BoTwoPoTwoHundred}{40}
\newcommand{\BoTwoPoTwoHundredHi}{41}
\newcommand{\BoTwoPoTwoHundredLo}{39}
\newcommand{\BoTypeBAtlasFifty}{63}
\newcommand{\BoTypeBAtlasHundred}{67}
\newcommand{\BoTypeBAtlasOneFifty}{70}
\newcommand{\BoTypeBPoFifty}{57}
\newcommand{\BoTypeBPoHundred}{61}
\newcommand{\BoTypeBPoOneFifty}{63}
\newcommand{\BoTypeCDEightAtlasFifty}{85}
\newcommand{\BoTypeCDEightAtlasHundred}{91}
\newcommand{\BoTypeCDEightAtlasOneFifty}{94}
\newcommand{\BoTypeCDEightPoFifty}{70}
\newcommand{\BoTypeCDEightPoHundred}{79}
\newcommand{\BoTypeCDEightPoOneFifty}{83}
\newcommand{\BoTypeCDFourAtlasFifty}{87}
\newcommand{\BoTypeCDFourAtlasHundred}{92}
\newcommand{\BoTypeCDFourAtlasOneFifty}{93}
\newcommand{\BoTypeCDFourPoFifty}{71}
\newcommand{\BoTypeCDFourPoHundred}{77}
\newcommand{\BoTypeCDFourPoOneFifty}{81}
\newcommand{\BoTypeMonoAtlasFifty}{66}
\newcommand{\BoTypeMonoAtlasHundred}{73}
\newcommand{\BoTypeMonoAtlasOneFifty}{76}
\newcommand{\BoTypeMonoPoFifty}{64}
\newcommand{\BoTypeMonoPoHundred}{72}
\newcommand{\BoTypeMonoPoOneFifty}{77}
\newcommand{\BoTypeNKAtlasFifty}{71}
\newcommand{\BoTypeNKAtlasHundred}{78}
\newcommand{\BoTypeNKAtlasOneFifty}{82}
\newcommand{\BoTypeNKPoFifty}{64}
\newcommand{\BoTypeNKPoHundred}{69}
\newcommand{\BoTypeNKPoOneFifty}{73}
\newcommand{\BoUnitsB}{34}
\newcommand{\BoUnitsCDEight}{34}
\newcommand{\BoUnitsCDFour}{34}
\newcommand{\BoUnitsMono}{34}
\newcommand{\BoUnitsNK}{34}
\newcommand{\CalCon}{1.00}
\newcommand{\CalDfMax}{0.95}
\newcommand{\CalDfMin}{0.64}
\newcommand{\CbBoot}{20,000}
\newcommand{\CbDiffOneLoMin}{7.1}
\newcommand{\CbDiffOneMax}{11.8}
\newcommand{\CbDiffOneMin}{8.3}
\newcommand{\CbDiffTwoLoMin}{15.2}
\newcommand{\CbDiffTwoMax}{21.0}
\newcommand{\CbDiffTwoMin}{16.5}
\newcommand{\CbGlobalThreshold}{$1.1\times10^{-3}$}
\newcommand{\CbHolmMax}{$\leq2\times10^{-4}$}
\newcommand{\CbHolmOne}{$\leq2\times10^{-4}$}
\newcommand{\CbHolmThree}{$\leq2\times10^{-4}$}
\newcommand{\CbHolmTwo}{$\leq2\times10^{-4}$}
\newcommand{\CbMajFdrAllPooled}{58}
\newcommand{\CbMajFdrAllSample}{58}
\newcommand{\CbMajFdrAllType}{60}
\newcommand{\CbMajFdrAllUnit}{61}
\newcommand{\CbMajFdrUnitPooled}{57}
\newcommand{\CbMajFdrUnitSample}{58}
\newcommand{\CbMajFdrUnitType}{60}
\newcommand{\CbMajFdrUnitUnit}{62}
\newcommand{\CbMajPrimaryPooled}{57}
\newcommand{\CbMajPrimarySample}{57}
\newcommand{\CbMajPrimaryType}{58}
\newcommand{\CbMajPrimaryUnit}{58}
\newcommand{\CbMajThousandthPooled}{58}
\newcommand{\CbMajThousandthSample}{58}
\newcommand{\CbMajThousandthType}{60}
\newcommand{\CbMajThousandthUnit}{61}
\newcommand{\CbMajTwentiethPooled}{55}
\newcommand{\CbMajTwentiethSample}{56}
\newcommand{\CbMajTwentiethType}{56}
\newcommand{\CbMajTwentiethUnit}{56}
\newcommand{\CbNFdrAll}{80,156}
\newcommand{\CbNFdrUnit}{83,863}
\newcommand{\CbNPrimary}{120,649}
\newcommand{\CbNThousandth}{79,414}
\newcommand{\CbNTwentieth}{189,472}
\newcommand{\CbOneFdrAllPooledAtlasFifty}{83}
\newcommand{\CbOneFdrAllPooledAtlasHundred}{89}
\newcommand{\CbOneFdrAllPooledAtlasOneFifty}{91}
\newcommand{\CbOneFdrAllPooledDiffFifty}{12.7}
\newcommand{\CbOneFdrAllPooledDiffFiftyHi}{13.6}
\newcommand{\CbOneFdrAllPooledDiffFiftyLo}{11.7}
\newcommand{\CbOneFdrAllPooledDiffHundred}{10.3}
\newcommand{\CbOneFdrAllPooledDiffHundredHi}{11.1}
\newcommand{\CbOneFdrAllPooledDiffHundredLo}{9.3}
\newcommand{\CbOneFdrAllPooledDiffOneFifty}{7.8}
\newcommand{\CbOneFdrAllPooledDiffOneFiftyHi}{8.5}
\newcommand{\CbOneFdrAllPooledDiffOneFiftyLo}{7.0}
\newcommand{\CbOneFdrAllPooledPoFifty}{70}
\newcommand{\CbOneFdrAllPooledPoHundred}{79}
\newcommand{\CbOneFdrAllPooledPoOneFifty}{84}
\newcommand{\CbOneFdrAllSampleAtlasFifty}{83}
\newcommand{\CbOneFdrAllSampleAtlasHundred}{89}
\newcommand{\CbOneFdrAllSampleAtlasOneFifty}{91}
\newcommand{\CbOneFdrAllSampleDiffFifty}{12.5}
\newcommand{\CbOneFdrAllSampleDiffFiftyHi}{13.5}
\newcommand{\CbOneFdrAllSampleDiffFiftyLo}{11.5}
\newcommand{\CbOneFdrAllSampleDiffHundred}{10.0}
\newcommand{\CbOneFdrAllSampleDiffHundredHi}{10.9}
\newcommand{\CbOneFdrAllSampleDiffHundredLo}{9.0}
\newcommand{\CbOneFdrAllSampleDiffOneFifty}{7.6}
\newcommand{\CbOneFdrAllSampleDiffOneFiftyHi}{8.4}
\newcommand{\CbOneFdrAllSampleDiffOneFiftyLo}{6.8}
\newcommand{\CbOneFdrAllSamplePoFifty}{71}
\newcommand{\CbOneFdrAllSamplePoHundred}{79}
\newcommand{\CbOneFdrAllSamplePoOneFifty}{84}
\newcommand{\CbOneFdrAllTypeAtlasFifty}{76}
\newcommand{\CbOneFdrAllTypeAtlasHundred}{82}
\newcommand{\CbOneFdrAllTypeAtlasOneFifty}{85}
\newcommand{\CbOneFdrAllTypeDiffFifty}{9.8}
\newcommand{\CbOneFdrAllTypeDiffFiftyHi}{10.4}
\newcommand{\CbOneFdrAllTypeDiffFiftyLo}{9.2}
\newcommand{\CbOneFdrAllTypeDiffHundred}{8.9}
\newcommand{\CbOneFdrAllTypeDiffHundredHi}{9.6}
\newcommand{\CbOneFdrAllTypeDiffHundredLo}{8.1}
\newcommand{\CbOneFdrAllTypeDiffOneFifty}{6.9}
\newcommand{\CbOneFdrAllTypeDiffOneFiftyHi}{7.5}
\newcommand{\CbOneFdrAllTypeDiffOneFiftyLo}{6.2}
\newcommand{\CbOneFdrAllTypePoFifty}{67}
\newcommand{\CbOneFdrAllTypePoHundred}{73}
\newcommand{\CbOneFdrAllTypePoOneFifty}{78}
\newcommand{\CbOneFdrAllUnitAtlasFifty}{77}
\newcommand{\CbOneFdrAllUnitAtlasHundred}{83}
\newcommand{\CbOneFdrAllUnitAtlasOneFifty}{86}
\newcommand{\CbOneFdrAllUnitDiffFifty}{9.9}
\newcommand{\CbOneFdrAllUnitDiffFiftyHi}{10.7}
\newcommand{\CbOneFdrAllUnitDiffFiftyLo}{9.2}
\newcommand{\CbOneFdrAllUnitDiffHundred}{8.7}
\newcommand{\CbOneFdrAllUnitDiffHundredHi}{9.5}
\newcommand{\CbOneFdrAllUnitDiffHundredLo}{8.0}
\newcommand{\CbOneFdrAllUnitDiffOneFifty}{6.5}
\newcommand{\CbOneFdrAllUnitDiffOneFiftyHi}{7.1}
\newcommand{\CbOneFdrAllUnitDiffOneFiftyLo}{6.0}
\newcommand{\CbOneFdrAllUnitPoFifty}{67}
\newcommand{\CbOneFdrAllUnitPoHundred}{74}
\newcommand{\CbOneFdrAllUnitPoOneFifty}{79}
\newcommand{\CbOneFdrUnitPooledAtlasFifty}{83}
\newcommand{\CbOneFdrUnitPooledAtlasHundred}{89}
\newcommand{\CbOneFdrUnitPooledAtlasOneFifty}{92}
\newcommand{\CbOneFdrUnitPooledDiffFifty}{13.4}
\newcommand{\CbOneFdrUnitPooledDiffFiftyHi}{14.3}
\newcommand{\CbOneFdrUnitPooledDiffFiftyLo}{12.4}
\newcommand{\CbOneFdrUnitPooledDiffHundred}{10.9}
\newcommand{\CbOneFdrUnitPooledDiffHundredHi}{11.7}
\newcommand{\CbOneFdrUnitPooledDiffHundredLo}{9.9}
\newcommand{\CbOneFdrUnitPooledDiffOneFifty}{8.5}
\newcommand{\CbOneFdrUnitPooledDiffOneFiftyHi}{9.3}
\newcommand{\CbOneFdrUnitPooledDiffOneFiftyLo}{7.6}
\newcommand{\CbOneFdrUnitPooledPoFifty}{70}
\newcommand{\CbOneFdrUnitPooledPoHundred}{78}
\newcommand{\CbOneFdrUnitPooledPoOneFifty}{83}
\newcommand{\CbOneFdrUnitSampleAtlasFifty}{84}
\newcommand{\CbOneFdrUnitSampleAtlasHundred}{89}
\newcommand{\CbOneFdrUnitSampleAtlasOneFifty}{92}
\newcommand{\CbOneFdrUnitSampleDiffFifty}{13.1}
\newcommand{\CbOneFdrUnitSampleDiffFiftyHi}{14.1}
\newcommand{\CbOneFdrUnitSampleDiffFiftyLo}{12.1}
\newcommand{\CbOneFdrUnitSampleDiffHundred}{10.6}
\newcommand{\CbOneFdrUnitSampleDiffHundredHi}{11.6}
\newcommand{\CbOneFdrUnitSampleDiffHundredLo}{9.5}
\newcommand{\CbOneFdrUnitSampleDiffOneFifty}{8.2}
\newcommand{\CbOneFdrUnitSampleDiffOneFiftyHi}{9.1}
\newcommand{\CbOneFdrUnitSampleDiffOneFiftyLo}{7.2}
\newcommand{\CbOneFdrUnitSamplePoFifty}{71}
\newcommand{\CbOneFdrUnitSamplePoHundred}{79}
\newcommand{\CbOneFdrUnitSamplePoOneFifty}{84}
\newcommand{\CbOneFdrUnitTypeAtlasFifty}{76}
\newcommand{\CbOneFdrUnitTypeAtlasHundred}{82}
\newcommand{\CbOneFdrUnitTypeAtlasOneFifty}{85}
\newcommand{\CbOneFdrUnitTypeDiffFifty}{9.8}
\newcommand{\CbOneFdrUnitTypeDiffFiftyHi}{10.6}
\newcommand{\CbOneFdrUnitTypeDiffFiftyLo}{9.1}
\newcommand{\CbOneFdrUnitTypeDiffHundred}{8.7}
\newcommand{\CbOneFdrUnitTypeDiffHundredHi}{9.6}
\newcommand{\CbOneFdrUnitTypeDiffHundredLo}{7.7}
\newcommand{\CbOneFdrUnitTypeDiffOneFifty}{6.8}
\newcommand{\CbOneFdrUnitTypeDiffOneFiftyHi}{7.6}
\newcommand{\CbOneFdrUnitTypeDiffOneFiftyLo}{5.9}
\newcommand{\CbOneFdrUnitTypePoFifty}{67}
\newcommand{\CbOneFdrUnitTypePoHundred}{74}
\newcommand{\CbOneFdrUnitTypePoOneFifty}{78}
\newcommand{\CbOneFdrUnitUnitAtlasFifty}{78}
\newcommand{\CbOneFdrUnitUnitAtlasHundred}{84}
\newcommand{\CbOneFdrUnitUnitAtlasOneFifty}{86}
\newcommand{\CbOneFdrUnitUnitDiffFifty}{10.4}
\newcommand{\CbOneFdrUnitUnitDiffFiftyHi}{11.9}
\newcommand{\CbOneFdrUnitUnitDiffFiftyLo}{9.0}
\newcommand{\CbOneFdrUnitUnitDiffHundred}{8.3}
\newcommand{\CbOneFdrUnitUnitDiffHundredHi}{9.6}
\newcommand{\CbOneFdrUnitUnitDiffHundredLo}{7.1}
\newcommand{\CbOneFdrUnitUnitDiffOneFifty}{6.0}
\newcommand{\CbOneFdrUnitUnitDiffOneFiftyHi}{6.9}
\newcommand{\CbOneFdrUnitUnitDiffOneFiftyLo}{5.0}
\newcommand{\CbOneFdrUnitUnitPoFifty}{68}
\newcommand{\CbOneFdrUnitUnitPoHundred}{75}
\newcommand{\CbOneFdrUnitUnitPoOneFifty}{80}
\newcommand{\CbOnePrimaryPooledAtlasFifty}{80}
\newcommand{\CbOnePrimaryPooledAtlasHundred}{86}
\newcommand{\CbOnePrimaryPooledAtlasOneFifty}{89}
\newcommand{\CbOnePrimaryPooledDiffFifty}{12.9}
\newcommand{\CbOnePrimaryPooledDiffFiftyHi}{13.7}
\newcommand{\CbOnePrimaryPooledDiffFiftyLo}{12.1}
\newcommand{\CbOnePrimaryPooledDiffHundred}{11.2}
\newcommand{\CbOnePrimaryPooledDiffHundredHi}{12.0}
\newcommand{\CbOnePrimaryPooledDiffHundredLo}{10.3}
\newcommand{\CbOnePrimaryPooledDiffOneFifty}{9.1}
\newcommand{\CbOnePrimaryPooledDiffOneFiftyHi}{9.8}
\newcommand{\CbOnePrimaryPooledDiffOneFiftyLo}{8.4}
\newcommand{\CbOnePrimaryPooledPoFifty}{67}
\newcommand{\CbOnePrimaryPooledPoHundred}{75}
\newcommand{\CbOnePrimaryPooledPoOneFifty}{80}
\newcommand{\CbOnePrimarySampleAtlasFifty}{80}
\newcommand{\CbOnePrimarySampleAtlasHundred}{86}
\newcommand{\CbOnePrimarySampleAtlasOneFifty}{89}
\newcommand{\CbOnePrimarySampleDiffFifty}{12.8}
\newcommand{\CbOnePrimarySampleDiffFiftyHi}{13.7}
\newcommand{\CbOnePrimarySampleDiffFiftyLo}{11.9}
\newcommand{\CbOnePrimarySampleDiffHundred}{11.1}
\newcommand{\CbOnePrimarySampleDiffHundredHi}{11.9}
\newcommand{\CbOnePrimarySampleDiffHundredLo}{10.1}
\newcommand{\CbOnePrimarySampleDiffOneFifty}{9.0}
\newcommand{\CbOnePrimarySampleDiffOneFiftyHi}{9.7}
\newcommand{\CbOnePrimarySampleDiffOneFiftyLo}{8.1}
\newcommand{\CbOnePrimarySamplePoFifty}{68}
\newcommand{\CbOnePrimarySamplePoHundred}{75}
\newcommand{\CbOnePrimarySamplePoOneFifty}{80}
\newcommand{\CbOnePrimaryTypeAtlasFifty}{75}
\newcommand{\CbOnePrimaryTypeAtlasHundred}{80}
\newcommand{\CbOnePrimaryTypeAtlasOneFifty}{83}
\newcommand{\CbOnePrimaryTypeDiffFifty}{10.3}
\newcommand{\CbOnePrimaryTypeDiffFiftyHi}{10.8}
\newcommand{\CbOnePrimaryTypeDiffFiftyLo}{9.7}
\newcommand{\CbOnePrimaryTypeDiffHundred}{9.9}
\newcommand{\CbOnePrimaryTypeDiffHundredHi}{10.4}
\newcommand{\CbOnePrimaryTypeDiffHundredLo}{9.3}
\newcommand{\CbOnePrimaryTypeDiffOneFifty}{8.2}
\newcommand{\CbOnePrimaryTypeDiffOneFiftyHi}{8.7}
\newcommand{\CbOnePrimaryTypeDiffOneFiftyLo}{7.6}
\newcommand{\CbOnePrimaryTypePoFifty}{64}
\newcommand{\CbOnePrimaryTypePoHundred}{70}
\newcommand{\CbOnePrimaryTypePoOneFifty}{75}
\newcommand{\CbOnePrimaryUnitAtlasFifty}{75}
\newcommand{\CbOnePrimaryUnitAtlasHundred}{81}
\newcommand{\CbOnePrimaryUnitAtlasOneFifty}{84}
\newcommand{\CbOnePrimaryUnitDiffFifty}{10.3}
\newcommand{\CbOnePrimaryUnitDiffFiftyHi}{10.9}
\newcommand{\CbOnePrimaryUnitDiffFiftyLo}{9.7}
\newcommand{\CbOnePrimaryUnitDiffHundred}{9.9}
\newcommand{\CbOnePrimaryUnitDiffHundredHi}{10.4}
\newcommand{\CbOnePrimaryUnitDiffHundredLo}{9.3}
\newcommand{\CbOnePrimaryUnitDiffOneFifty}{7.9}
\newcommand{\CbOnePrimaryUnitDiffOneFiftyHi}{8.3}
\newcommand{\CbOnePrimaryUnitDiffOneFiftyLo}{7.4}
\newcommand{\CbOnePrimaryUnitPoFifty}{65}
\newcommand{\CbOnePrimaryUnitPoHundred}{71}
\newcommand{\CbOnePrimaryUnitPoOneFifty}{76}
\newcommand{\CbOneThousandthPooledAtlasFifty}{83}
\newcommand{\CbOneThousandthPooledAtlasHundred}{89}
\newcommand{\CbOneThousandthPooledAtlasOneFifty}{91}
\newcommand{\CbOneThousandthPooledDiffFifty}{12.7}
\newcommand{\CbOneThousandthPooledDiffFiftyHi}{13.6}
\newcommand{\CbOneThousandthPooledDiffFiftyLo}{11.7}
\newcommand{\CbOneThousandthPooledDiffHundred}{10.2}
\newcommand{\CbOneThousandthPooledDiffHundredHi}{11.0}
\newcommand{\CbOneThousandthPooledDiffHundredLo}{9.3}
\newcommand{\CbOneThousandthPooledDiffOneFifty}{7.8}
\newcommand{\CbOneThousandthPooledDiffOneFiftyHi}{8.5}
\newcommand{\CbOneThousandthPooledDiffOneFiftyLo}{7.0}
\newcommand{\CbOneThousandthPooledPoFifty}{70}
\newcommand{\CbOneThousandthPooledPoHundred}{79}
\newcommand{\CbOneThousandthPooledPoOneFifty}{84}
\newcommand{\CbOneThousandthSampleAtlasFifty}{83}
\newcommand{\CbOneThousandthSampleAtlasHundred}{89}
\newcommand{\CbOneThousandthSampleAtlasOneFifty}{92}
\newcommand{\CbOneThousandthSampleDiffFifty}{12.5}
\newcommand{\CbOneThousandthSampleDiffFiftyHi}{13.5}
\newcommand{\CbOneThousandthSampleDiffFiftyLo}{11.5}
\newcommand{\CbOneThousandthSampleDiffHundred}{10.0}
\newcommand{\CbOneThousandthSampleDiffHundredHi}{10.9}
\newcommand{\CbOneThousandthSampleDiffHundredLo}{9.0}
\newcommand{\CbOneThousandthSampleDiffOneFifty}{7.6}
\newcommand{\CbOneThousandthSampleDiffOneFiftyHi}{8.4}
\newcommand{\CbOneThousandthSampleDiffOneFiftyLo}{6.8}
\newcommand{\CbOneThousandthSamplePoFifty}{71}
\newcommand{\CbOneThousandthSamplePoHundred}{79}
\newcommand{\CbOneThousandthSamplePoOneFifty}{84}
\newcommand{\CbOneThousandthTypeAtlasFifty}{76}
\newcommand{\CbOneThousandthTypeAtlasHundred}{82}
\newcommand{\CbOneThousandthTypeAtlasOneFifty}{85}
\newcommand{\CbOneThousandthTypeDiffFifty}{9.7}
\newcommand{\CbOneThousandthTypeDiffFiftyHi}{10.3}
\newcommand{\CbOneThousandthTypeDiffFiftyLo}{9.1}
\newcommand{\CbOneThousandthTypeDiffHundred}{8.8}
\newcommand{\CbOneThousandthTypeDiffHundredHi}{9.5}
\newcommand{\CbOneThousandthTypeDiffHundredLo}{8.0}
\newcommand{\CbOneThousandthTypeDiffOneFifty}{6.8}
\newcommand{\CbOneThousandthTypeDiffOneFiftyHi}{7.4}
\newcommand{\CbOneThousandthTypeDiffOneFiftyLo}{6.1}
\newcommand{\CbOneThousandthTypePoFifty}{67}
\newcommand{\CbOneThousandthTypePoHundred}{73}
\newcommand{\CbOneThousandthTypePoOneFifty}{78}
\newcommand{\CbOneThousandthUnitAtlasFifty}{77}
\newcommand{\CbOneThousandthUnitAtlasHundred}{83}
\newcommand{\CbOneThousandthUnitAtlasOneFifty}{86}
\newcommand{\CbOneThousandthUnitDiffFifty}{9.9}
\newcommand{\CbOneThousandthUnitDiffFiftyHi}{10.6}
\newcommand{\CbOneThousandthUnitDiffFiftyLo}{9.2}
\newcommand{\CbOneThousandthUnitDiffHundred}{8.6}
\newcommand{\CbOneThousandthUnitDiffHundredHi}{9.4}
\newcommand{\CbOneThousandthUnitDiffHundredLo}{7.9}
\newcommand{\CbOneThousandthUnitDiffOneFifty}{6.5}
\newcommand{\CbOneThousandthUnitDiffOneFiftyHi}{7.1}
\newcommand{\CbOneThousandthUnitDiffOneFiftyLo}{5.9}
\newcommand{\CbOneThousandthUnitPoFifty}{67}
\newcommand{\CbOneThousandthUnitPoHundred}{75}
\newcommand{\CbOneThousandthUnitPoOneFifty}{79}
\newcommand{\CbOneTwentiethPooledAtlasFifty}{76}
\newcommand{\CbOneTwentiethPooledAtlasHundred}{82}
\newcommand{\CbOneTwentiethPooledAtlasOneFifty}{85}
\newcommand{\CbOneTwentiethPooledDiffFifty}{12.5}
\newcommand{\CbOneTwentiethPooledDiffFiftyHi}{13.2}
\newcommand{\CbOneTwentiethPooledDiffFiftyLo}{11.7}
\newcommand{\CbOneTwentiethPooledDiffHundred}{11.8}
\newcommand{\CbOneTwentiethPooledDiffHundredHi}{12.4}
\newcommand{\CbOneTwentiethPooledDiffHundredLo}{11.0}
\newcommand{\CbOneTwentiethPooledDiffOneFifty}{10.2}
\newcommand{\CbOneTwentiethPooledDiffOneFiftyHi}{10.9}
\newcommand{\CbOneTwentiethPooledDiffOneFiftyLo}{9.5}
\newcommand{\CbOneTwentiethPooledPoFifty}{64}
\newcommand{\CbOneTwentiethPooledPoHundred}{70}
\newcommand{\CbOneTwentiethPooledPoOneFifty}{74}
\newcommand{\CbOneTwentiethSampleAtlasFifty}{77}
\newcommand{\CbOneTwentiethSampleAtlasHundred}{82}
\newcommand{\CbOneTwentiethSampleAtlasOneFifty}{85}
\newcommand{\CbOneTwentiethSampleDiffFifty}{12.4}
\newcommand{\CbOneTwentiethSampleDiffFiftyHi}{13.2}
\newcommand{\CbOneTwentiethSampleDiffFiftyLo}{11.6}
\newcommand{\CbOneTwentiethSampleDiffHundred}{11.7}
\newcommand{\CbOneTwentiethSampleDiffHundredHi}{12.4}
\newcommand{\CbOneTwentiethSampleDiffHundredLo}{10.9}
\newcommand{\CbOneTwentiethSampleDiffOneFifty}{10.1}
\newcommand{\CbOneTwentiethSampleDiffOneFiftyHi}{10.8}
\newcommand{\CbOneTwentiethSampleDiffOneFiftyLo}{9.4}
\newcommand{\CbOneTwentiethSamplePoFifty}{64}
\newcommand{\CbOneTwentiethSamplePoHundred}{70}
\newcommand{\CbOneTwentiethSamplePoOneFifty}{75}
\newcommand{\CbOneTwentiethTypeAtlasFifty}{72}
\newcommand{\CbOneTwentiethTypeAtlasHundred}{77}
\newcommand{\CbOneTwentiethTypeAtlasOneFifty}{80}
\newcommand{\CbOneTwentiethTypeDiffFifty}{10.2}
\newcommand{\CbOneTwentiethTypeDiffFiftyHi}{10.7}
\newcommand{\CbOneTwentiethTypeDiffFiftyLo}{9.7}
\newcommand{\CbOneTwentiethTypeDiffHundred}{10.6}
\newcommand{\CbOneTwentiethTypeDiffHundredHi}{11.1}
\newcommand{\CbOneTwentiethTypeDiffHundredLo}{10.0}
\newcommand{\CbOneTwentiethTypeDiffOneFifty}{9.3}
\newcommand{\CbOneTwentiethTypeDiffOneFiftyHi}{9.9}
\newcommand{\CbOneTwentiethTypeDiffOneFiftyLo}{8.8}
\newcommand{\CbOneTwentiethTypePoFifty}{62}
\newcommand{\CbOneTwentiethTypePoHundred}{67}
\newcommand{\CbOneTwentiethTypePoOneFifty}{71}
\newcommand{\CbOneTwentiethUnitAtlasFifty}{72}
\newcommand{\CbOneTwentiethUnitAtlasHundred}{78}
\newcommand{\CbOneTwentiethUnitAtlasOneFifty}{80}
\newcommand{\CbOneTwentiethUnitDiffFifty}{10.2}
\newcommand{\CbOneTwentiethUnitDiffFiftyHi}{10.7}
\newcommand{\CbOneTwentiethUnitDiffFiftyLo}{9.7}
\newcommand{\CbOneTwentiethUnitDiffHundred}{10.5}
\newcommand{\CbOneTwentiethUnitDiffHundredHi}{11.0}
\newcommand{\CbOneTwentiethUnitDiffHundredLo}{10.0}
\newcommand{\CbOneTwentiethUnitDiffOneFifty}{9.1}
\newcommand{\CbOneTwentiethUnitDiffOneFiftyHi}{9.6}
\newcommand{\CbOneTwentiethUnitDiffOneFiftyLo}{8.6}
\newcommand{\CbOneTwentiethUnitPoFifty}{62}
\newcommand{\CbOneTwentiethUnitPoHundred}{67}
\newcommand{\CbOneTwentiethUnitPoOneFifty}{71}
\newcommand{\CbOverallFdrAll}{58}
\newcommand{\CbOverallFdrUnit}{57}
\newcommand{\CbOverallPrimary}{57}
\newcommand{\CbOverallThousandth}{58}
\newcommand{\CbOverallTwentieth}{55}
\newcommand{\CbPResolution}{$5.0\times10^{-5}$}
\newcommand{\CbPlanMaxP}{$\leq5\times10^{-5}$}
\newcommand{\CbPlanOneFiftyP}{$\leq5\times10^{-5}$}
\newcommand{\CbPlanOneFiftySimHi}{14.1}
\newcommand{\CbPlanOneFiftySimLo}{11.7}
\newcommand{\CbPlanOneHundredP}{$\leq5\times10^{-5}$}
\newcommand{\CbPlanOneHundredSimHi}{12.2}
\newcommand{\CbPlanOneHundredSimLo}{10.1}
\newcommand{\CbPlanOneOneFiftyP}{$\leq5\times10^{-5}$}
\newcommand{\CbPlanOneOneFiftySimHi}{10.2}
\newcommand{\CbPlanOneOneFiftySimLo}{8.0}
\newcommand{\CbPlanThreeHundredP}{$\leq5\times10^{-5}$}
\newcommand{\CbPlanThreeHundredSimHi}{16.2}
\newcommand{\CbPlanThreeHundredSimLo}{11.9}
\newcommand{\CbPlanThreeOneFiftyP}{$\leq5\times10^{-5}$}
\newcommand{\CbPlanThreeOneFiftySimHi}{15.8}
\newcommand{\CbPlanThreeOneFiftySimLo}{12.4}
\newcommand{\CbPlanTwoFiftyP}{$\leq5\times10^{-5}$}
\newcommand{\CbPlanTwoFiftySimHi}{21.8}
\newcommand{\CbPlanTwoFiftySimLo}{18.5}
\newcommand{\CbPlanTwoHundredP}{$\leq5\times10^{-5}$}
\newcommand{\CbPlanTwoHundredSimHi}{20.9}
\newcommand{\CbPlanTwoHundredSimLo}{17.4}
\newcommand{\CbPlanTwoOneFiftyP}{$\leq5\times10^{-5}$}
\newcommand{\CbPlanTwoOneFiftySimHi}{18.3}
\newcommand{\CbPlanTwoOneFiftySimLo}{15.1}
\newcommand{\CbPosFdrAll}{58}
\newcommand{\CbPosFdrUnit}{57}
\newcommand{\CbPosPrimary}{57}
\newcommand{\CbPosThousandth}{58}
\newcommand{\CbPosTwentieth}{55}
\newcommand{\CbRandomTwoFdrAll}{1.1}
\newcommand{\CbRandomTwoFdrUnit}{1.1}
\newcommand{\CbRandomTwoPrimary}{1.6}
\newcommand{\CbRandomTwoThousandth}{1.1}
\newcommand{\CbRandomTwoTwentieth}{2.5}
\newcommand{\CbShareCDEightFdrAll}{52}
\newcommand{\CbShareCDEightFdrUnit}{57}
\newcommand{\CbShareCDEightPrimary}{47}
\newcommand{\CbShareCDEightThousandth}{53}
\newcommand{\CbShareCDEightTwentieth}{41}
\newcommand{\CbTwoFdrAllPooledAtlasFifty}{31}
\newcommand{\CbTwoFdrAllPooledAtlasHundred}{42}
\newcommand{\CbTwoFdrAllPooledAtlasOneFifty}{47}
\newcommand{\CbTwoFdrAllPooledDiffFifty}{17.7}
\newcommand{\CbTwoFdrAllPooledDiffFiftyHi}{18.9}
\newcommand{\CbTwoFdrAllPooledDiffFiftyLo}{16.4}
\newcommand{\CbTwoFdrAllPooledDiffHundred}{17.2}
\newcommand{\CbTwoFdrAllPooledDiffHundredHi}{18.5}
\newcommand{\CbTwoFdrAllPooledDiffHundredLo}{15.9}
\newcommand{\CbTwoFdrAllPooledDiffOneFifty}{14.8}
\newcommand{\CbTwoFdrAllPooledDiffOneFiftyHi}{16.0}
\newcommand{\CbTwoFdrAllPooledDiffOneFiftyLo}{13.5}
\newcommand{\CbTwoFdrAllPooledPoFifty}{14}
\newcommand{\CbTwoFdrAllPooledPoHundred}{25}
\newcommand{\CbTwoFdrAllPooledPoOneFifty}{32}
\newcommand{\CbTwoFdrAllSampleAtlasFifty}{31}
\newcommand{\CbTwoFdrAllSampleAtlasHundred}{42}
\newcommand{\CbTwoFdrAllSampleAtlasOneFifty}{47}
\newcommand{\CbTwoFdrAllSampleDiffFifty}{17.7}
\newcommand{\CbTwoFdrAllSampleDiffFiftyHi}{18.9}
\newcommand{\CbTwoFdrAllSampleDiffFiftyLo}{16.4}
\newcommand{\CbTwoFdrAllSampleDiffHundred}{17.2}
\newcommand{\CbTwoFdrAllSampleDiffHundredHi}{18.5}
\newcommand{\CbTwoFdrAllSampleDiffHundredLo}{15.9}
\newcommand{\CbTwoFdrAllSampleDiffOneFifty}{14.8}
\newcommand{\CbTwoFdrAllSampleDiffOneFiftyHi}{16.0}
\newcommand{\CbTwoFdrAllSampleDiffOneFiftyLo}{13.5}
\newcommand{\CbTwoFdrAllSamplePoFifty}{14}
\newcommand{\CbTwoFdrAllSamplePoHundred}{25}
\newcommand{\CbTwoFdrAllSamplePoOneFifty}{32}
\newcommand{\CbTwoFdrAllTypeAtlasFifty}{31}
\newcommand{\CbTwoFdrAllTypeAtlasHundred}{42}
\newcommand{\CbTwoFdrAllTypeAtlasOneFifty}{47}
\newcommand{\CbTwoFdrAllTypeDiffFifty}{17.7}
\newcommand{\CbTwoFdrAllTypeDiffFiftyHi}{18.9}
\newcommand{\CbTwoFdrAllTypeDiffFiftyLo}{16.4}
\newcommand{\CbTwoFdrAllTypeDiffHundred}{17.2}
\newcommand{\CbTwoFdrAllTypeDiffHundredHi}{18.5}
\newcommand{\CbTwoFdrAllTypeDiffHundredLo}{15.9}
\newcommand{\CbTwoFdrAllTypeDiffOneFifty}{14.8}
\newcommand{\CbTwoFdrAllTypeDiffOneFiftyHi}{16.0}
\newcommand{\CbTwoFdrAllTypeDiffOneFiftyLo}{13.5}
\newcommand{\CbTwoFdrAllTypePoFifty}{14}
\newcommand{\CbTwoFdrAllTypePoHundred}{25}
\newcommand{\CbTwoFdrAllTypePoOneFifty}{32}
\newcommand{\CbTwoFdrAllUnitAtlasFifty}{31}
\newcommand{\CbTwoFdrAllUnitAtlasHundred}{42}
\newcommand{\CbTwoFdrAllUnitAtlasOneFifty}{47}
\newcommand{\CbTwoFdrAllUnitDiffFifty}{17.7}
\newcommand{\CbTwoFdrAllUnitDiffFiftyHi}{18.9}
\newcommand{\CbTwoFdrAllUnitDiffFiftyLo}{16.4}
\newcommand{\CbTwoFdrAllUnitDiffHundred}{17.2}
\newcommand{\CbTwoFdrAllUnitDiffHundredHi}{18.5}
\newcommand{\CbTwoFdrAllUnitDiffHundredLo}{15.9}
\newcommand{\CbTwoFdrAllUnitDiffOneFifty}{14.8}
\newcommand{\CbTwoFdrAllUnitDiffOneFiftyHi}{16.0}
\newcommand{\CbTwoFdrAllUnitDiffOneFiftyLo}{13.5}
\newcommand{\CbTwoFdrAllUnitPoFifty}{14}
\newcommand{\CbTwoFdrAllUnitPoHundred}{25}
\newcommand{\CbTwoFdrAllUnitPoOneFifty}{32}
\newcommand{\CbTwoFdrUnitPooledAtlasFifty}{31}
\newcommand{\CbTwoFdrUnitPooledAtlasHundred}{41}
\newcommand{\CbTwoFdrUnitPooledAtlasOneFifty}{46}
\newcommand{\CbTwoFdrUnitPooledDiffFifty}{17.2}
\newcommand{\CbTwoFdrUnitPooledDiffFiftyHi}{18.6}
\newcommand{\CbTwoFdrUnitPooledDiffFiftyLo}{15.9}
\newcommand{\CbTwoFdrUnitPooledDiffHundred}{16.5}
\newcommand{\CbTwoFdrUnitPooledDiffHundredHi}{17.8}
\newcommand{\CbTwoFdrUnitPooledDiffHundredLo}{15.2}
\newcommand{\CbTwoFdrUnitPooledDiffOneFifty}{14.2}
\newcommand{\CbTwoFdrUnitPooledDiffOneFiftyHi}{15.5}
\newcommand{\CbTwoFdrUnitPooledDiffOneFiftyLo}{12.8}
\newcommand{\CbTwoFdrUnitPooledPoFifty}{14}
\newcommand{\CbTwoFdrUnitPooledPoHundred}{25}
\newcommand{\CbTwoFdrUnitPooledPoOneFifty}{32}
\newcommand{\CbTwoFdrUnitSampleAtlasFifty}{31}
\newcommand{\CbTwoFdrUnitSampleAtlasHundred}{41}
\newcommand{\CbTwoFdrUnitSampleAtlasOneFifty}{46}
\newcommand{\CbTwoFdrUnitSampleDiffFifty}{17.2}
\newcommand{\CbTwoFdrUnitSampleDiffFiftyHi}{18.6}
\newcommand{\CbTwoFdrUnitSampleDiffFiftyLo}{15.9}
\newcommand{\CbTwoFdrUnitSampleDiffHundred}{16.5}
\newcommand{\CbTwoFdrUnitSampleDiffHundredHi}{17.8}
\newcommand{\CbTwoFdrUnitSampleDiffHundredLo}{15.2}
\newcommand{\CbTwoFdrUnitSampleDiffOneFifty}{14.2}
\newcommand{\CbTwoFdrUnitSampleDiffOneFiftyHi}{15.5}
\newcommand{\CbTwoFdrUnitSampleDiffOneFiftyLo}{12.8}
\newcommand{\CbTwoFdrUnitSamplePoFifty}{14}
\newcommand{\CbTwoFdrUnitSamplePoHundred}{25}
\newcommand{\CbTwoFdrUnitSamplePoOneFifty}{32}
\newcommand{\CbTwoFdrUnitTypeAtlasFifty}{31}
\newcommand{\CbTwoFdrUnitTypeAtlasHundred}{41}
\newcommand{\CbTwoFdrUnitTypeAtlasOneFifty}{46}
\newcommand{\CbTwoFdrUnitTypeDiffFifty}{17.2}
\newcommand{\CbTwoFdrUnitTypeDiffFiftyHi}{18.6}
\newcommand{\CbTwoFdrUnitTypeDiffFiftyLo}{15.9}
\newcommand{\CbTwoFdrUnitTypeDiffHundred}{16.5}
\newcommand{\CbTwoFdrUnitTypeDiffHundredHi}{17.8}
\newcommand{\CbTwoFdrUnitTypeDiffHundredLo}{15.2}
\newcommand{\CbTwoFdrUnitTypeDiffOneFifty}{14.2}
\newcommand{\CbTwoFdrUnitTypeDiffOneFiftyHi}{15.5}
\newcommand{\CbTwoFdrUnitTypeDiffOneFiftyLo}{12.8}
\newcommand{\CbTwoFdrUnitTypePoFifty}{14}
\newcommand{\CbTwoFdrUnitTypePoHundred}{25}
\newcommand{\CbTwoFdrUnitTypePoOneFifty}{32}
\newcommand{\CbTwoFdrUnitUnitAtlasFifty}{31}
\newcommand{\CbTwoFdrUnitUnitAtlasHundred}{41}
\newcommand{\CbTwoFdrUnitUnitAtlasOneFifty}{46}
\newcommand{\CbTwoFdrUnitUnitDiffFifty}{17.2}
\newcommand{\CbTwoFdrUnitUnitDiffFiftyHi}{18.6}
\newcommand{\CbTwoFdrUnitUnitDiffFiftyLo}{15.9}
\newcommand{\CbTwoFdrUnitUnitDiffHundred}{16.5}
\newcommand{\CbTwoFdrUnitUnitDiffHundredHi}{17.8}
\newcommand{\CbTwoFdrUnitUnitDiffHundredLo}{15.2}
\newcommand{\CbTwoFdrUnitUnitDiffOneFifty}{14.2}
\newcommand{\CbTwoFdrUnitUnitDiffOneFiftyHi}{15.5}
\newcommand{\CbTwoFdrUnitUnitDiffOneFiftyLo}{12.8}
\newcommand{\CbTwoFdrUnitUnitPoFifty}{14}
\newcommand{\CbTwoFdrUnitUnitPoHundred}{25}
\newcommand{\CbTwoFdrUnitUnitPoOneFifty}{32}
\newcommand{\CbTwoPrimaryPooledAtlasFifty}{35}
\newcommand{\CbTwoPrimaryPooledAtlasHundred}{47}
\newcommand{\CbTwoPrimaryPooledAtlasOneFifty}{52}
\newcommand{\CbTwoPrimaryPooledDiffFifty}{20.1}
\newcommand{\CbTwoPrimaryPooledDiffFiftyHi}{21.3}
\newcommand{\CbTwoPrimaryPooledDiffFiftyLo}{18.9}
\newcommand{\CbTwoPrimaryPooledDiffHundred}{19.3}
\newcommand{\CbTwoPrimaryPooledDiffHundredHi}{20.5}
\newcommand{\CbTwoPrimaryPooledDiffHundredLo}{18.0}
\newcommand{\CbTwoPrimaryPooledDiffOneFifty}{16.8}
\newcommand{\CbTwoPrimaryPooledDiffOneFiftyHi}{18.0}
\newcommand{\CbTwoPrimaryPooledDiffOneFiftyLo}{15.6}
\newcommand{\CbTwoPrimaryPooledPoFifty}{15}
\newcommand{\CbTwoPrimaryPooledPoHundred}{28}
\newcommand{\CbTwoPrimaryPooledPoOneFifty}{35}
\newcommand{\CbTwoPrimarySampleAtlasFifty}{35}
\newcommand{\CbTwoPrimarySampleAtlasHundred}{47}
\newcommand{\CbTwoPrimarySampleAtlasOneFifty}{52}
\newcommand{\CbTwoPrimarySampleDiffFifty}{20.1}
\newcommand{\CbTwoPrimarySampleDiffFiftyHi}{21.3}
\newcommand{\CbTwoPrimarySampleDiffFiftyLo}{18.9}
\newcommand{\CbTwoPrimarySampleDiffHundred}{19.3}
\newcommand{\CbTwoPrimarySampleDiffHundredHi}{20.5}
\newcommand{\CbTwoPrimarySampleDiffHundredLo}{18.0}
\newcommand{\CbTwoPrimarySampleDiffOneFifty}{16.8}
\newcommand{\CbTwoPrimarySampleDiffOneFiftyHi}{18.0}
\newcommand{\CbTwoPrimarySampleDiffOneFiftyLo}{15.6}
\newcommand{\CbTwoPrimarySamplePoFifty}{15}
\newcommand{\CbTwoPrimarySamplePoHundred}{28}
\newcommand{\CbTwoPrimarySamplePoOneFifty}{35}
\newcommand{\CbTwoPrimaryTypeAtlasFifty}{35}
\newcommand{\CbTwoPrimaryTypeAtlasHundred}{47}
\newcommand{\CbTwoPrimaryTypeAtlasOneFifty}{52}
\newcommand{\CbTwoPrimaryTypeDiffFifty}{20.1}
\newcommand{\CbTwoPrimaryTypeDiffFiftyHi}{21.3}
\newcommand{\CbTwoPrimaryTypeDiffFiftyLo}{18.9}
\newcommand{\CbTwoPrimaryTypeDiffHundred}{19.3}
\newcommand{\CbTwoPrimaryTypeDiffHundredHi}{20.5}
\newcommand{\CbTwoPrimaryTypeDiffHundredLo}{18.0}
\newcommand{\CbTwoPrimaryTypeDiffOneFifty}{16.8}
\newcommand{\CbTwoPrimaryTypeDiffOneFiftyHi}{18.0}
\newcommand{\CbTwoPrimaryTypeDiffOneFiftyLo}{15.6}
\newcommand{\CbTwoPrimaryTypePoFifty}{15}
\newcommand{\CbTwoPrimaryTypePoHundred}{28}
\newcommand{\CbTwoPrimaryTypePoOneFifty}{35}
\newcommand{\CbTwoPrimaryUnitAtlasFifty}{35}
\newcommand{\CbTwoPrimaryUnitAtlasHundred}{47}
\newcommand{\CbTwoPrimaryUnitAtlasOneFifty}{52}
\newcommand{\CbTwoPrimaryUnitDiffFifty}{20.1}
\newcommand{\CbTwoPrimaryUnitDiffFiftyHi}{21.3}
\newcommand{\CbTwoPrimaryUnitDiffFiftyLo}{18.9}
\newcommand{\CbTwoPrimaryUnitDiffHundred}{19.3}
\newcommand{\CbTwoPrimaryUnitDiffHundredHi}{20.5}
\newcommand{\CbTwoPrimaryUnitDiffHundredLo}{18.0}
\newcommand{\CbTwoPrimaryUnitDiffOneFifty}{16.8}
\newcommand{\CbTwoPrimaryUnitDiffOneFiftyHi}{18.0}
\newcommand{\CbTwoPrimaryUnitDiffOneFiftyLo}{15.6}
\newcommand{\CbTwoPrimaryUnitPoFifty}{15}
\newcommand{\CbTwoPrimaryUnitPoHundred}{28}
\newcommand{\CbTwoPrimaryUnitPoOneFifty}{35}
\newcommand{\CbTwoThousandthPooledAtlasFifty}{31}
\newcommand{\CbTwoThousandthPooledAtlasHundred}{42}
\newcommand{\CbTwoThousandthPooledAtlasOneFifty}{47}
\newcommand{\CbTwoThousandthPooledDiffFifty}{17.6}
\newcommand{\CbTwoThousandthPooledDiffFiftyHi}{18.9}
\newcommand{\CbTwoThousandthPooledDiffFiftyLo}{16.4}
\newcommand{\CbTwoThousandthPooledDiffHundred}{17.2}
\newcommand{\CbTwoThousandthPooledDiffHundredHi}{18.4}
\newcommand{\CbTwoThousandthPooledDiffHundredLo}{15.8}
\newcommand{\CbTwoThousandthPooledDiffOneFifty}{14.8}
\newcommand{\CbTwoThousandthPooledDiffOneFiftyHi}{16.0}
\newcommand{\CbTwoThousandthPooledDiffOneFiftyLo}{13.4}
\newcommand{\CbTwoThousandthPooledPoFifty}{13}
\newcommand{\CbTwoThousandthPooledPoHundred}{25}
\newcommand{\CbTwoThousandthPooledPoOneFifty}{32}
\newcommand{\CbTwoThousandthSampleAtlasFifty}{31}
\newcommand{\CbTwoThousandthSampleAtlasHundred}{42}
\newcommand{\CbTwoThousandthSampleAtlasOneFifty}{47}
\newcommand{\CbTwoThousandthSampleDiffFifty}{17.6}
\newcommand{\CbTwoThousandthSampleDiffFiftyHi}{18.9}
\newcommand{\CbTwoThousandthSampleDiffFiftyLo}{16.4}
\newcommand{\CbTwoThousandthSampleDiffHundred}{17.2}
\newcommand{\CbTwoThousandthSampleDiffHundredHi}{18.4}
\newcommand{\CbTwoThousandthSampleDiffHundredLo}{15.8}
\newcommand{\CbTwoThousandthSampleDiffOneFifty}{14.8}
\newcommand{\CbTwoThousandthSampleDiffOneFiftyHi}{16.0}
\newcommand{\CbTwoThousandthSampleDiffOneFiftyLo}{13.4}
\newcommand{\CbTwoThousandthSamplePoFifty}{13}
\newcommand{\CbTwoThousandthSamplePoHundred}{25}
\newcommand{\CbTwoThousandthSamplePoOneFifty}{32}
\newcommand{\CbTwoThousandthTypeAtlasFifty}{31}
\newcommand{\CbTwoThousandthTypeAtlasHundred}{42}
\newcommand{\CbTwoThousandthTypeAtlasOneFifty}{47}
\newcommand{\CbTwoThousandthTypeDiffFifty}{17.6}
\newcommand{\CbTwoThousandthTypeDiffFiftyHi}{18.9}
\newcommand{\CbTwoThousandthTypeDiffFiftyLo}{16.4}
\newcommand{\CbTwoThousandthTypeDiffHundred}{17.2}
\newcommand{\CbTwoThousandthTypeDiffHundredHi}{18.4}
\newcommand{\CbTwoThousandthTypeDiffHundredLo}{15.8}
\newcommand{\CbTwoThousandthTypeDiffOneFifty}{14.8}
\newcommand{\CbTwoThousandthTypeDiffOneFiftyHi}{16.0}
\newcommand{\CbTwoThousandthTypeDiffOneFiftyLo}{13.4}
\newcommand{\CbTwoThousandthTypePoFifty}{13}
\newcommand{\CbTwoThousandthTypePoHundred}{25}
\newcommand{\CbTwoThousandthTypePoOneFifty}{32}
\newcommand{\CbTwoThousandthUnitAtlasFifty}{31}
\newcommand{\CbTwoThousandthUnitAtlasHundred}{42}
\newcommand{\CbTwoThousandthUnitAtlasOneFifty}{47}
\newcommand{\CbTwoThousandthUnitDiffFifty}{17.6}
\newcommand{\CbTwoThousandthUnitDiffFiftyHi}{18.9}
\newcommand{\CbTwoThousandthUnitDiffFiftyLo}{16.4}
\newcommand{\CbTwoThousandthUnitDiffHundred}{17.2}
\newcommand{\CbTwoThousandthUnitDiffHundredHi}{18.4}
\newcommand{\CbTwoThousandthUnitDiffHundredLo}{15.8}
\newcommand{\CbTwoThousandthUnitDiffOneFifty}{14.8}
\newcommand{\CbTwoThousandthUnitDiffOneFiftyHi}{16.0}
\newcommand{\CbTwoThousandthUnitDiffOneFiftyLo}{13.4}
\newcommand{\CbTwoThousandthUnitPoFifty}{13}
\newcommand{\CbTwoThousandthUnitPoHundred}{25}
\newcommand{\CbTwoThousandthUnitPoOneFifty}{32}
\newcommand{\CbTwoTwentiethPooledAtlasFifty}{40}
\newcommand{\CbTwoTwentiethPooledAtlasHundred}{52}
\newcommand{\CbTwoTwentiethPooledAtlasOneFifty}{57}
\newcommand{\CbTwoTwentiethPooledDiffFifty}{22.2}
\newcommand{\CbTwoTwentiethPooledDiffFiftyHi}{23.3}
\newcommand{\CbTwoTwentiethPooledDiffFiftyLo}{21.1}
\newcommand{\CbTwoTwentiethPooledDiffHundred}{21.0}
\newcommand{\CbTwoTwentiethPooledDiffHundredHi}{22.0}
\newcommand{\CbTwoTwentiethPooledDiffHundredLo}{19.9}
\newcommand{\CbTwoTwentiethPooledDiffOneFifty}{18.1}
\newcommand{\CbTwoTwentiethPooledDiffOneFiftyHi}{19.1}
\newcommand{\CbTwoTwentiethPooledDiffOneFiftyLo}{17.1}
\newcommand{\CbTwoTwentiethPooledPoFifty}{18}
\newcommand{\CbTwoTwentiethPooledPoHundred}{31}
\newcommand{\CbTwoTwentiethPooledPoOneFifty}{39}
\newcommand{\CbTwoTwentiethSampleAtlasFifty}{40}
\newcommand{\CbTwoTwentiethSampleAtlasHundred}{52}
\newcommand{\CbTwoTwentiethSampleAtlasOneFifty}{57}
\newcommand{\CbTwoTwentiethSampleDiffFifty}{22.2}
\newcommand{\CbTwoTwentiethSampleDiffFiftyHi}{23.3}
\newcommand{\CbTwoTwentiethSampleDiffFiftyLo}{21.1}
\newcommand{\CbTwoTwentiethSampleDiffHundred}{21.0}
\newcommand{\CbTwoTwentiethSampleDiffHundredHi}{22.0}
\newcommand{\CbTwoTwentiethSampleDiffHundredLo}{19.9}
\newcommand{\CbTwoTwentiethSampleDiffOneFifty}{18.1}
\newcommand{\CbTwoTwentiethSampleDiffOneFiftyHi}{19.1}
\newcommand{\CbTwoTwentiethSampleDiffOneFiftyLo}{17.1}
\newcommand{\CbTwoTwentiethSamplePoFifty}{18}
\newcommand{\CbTwoTwentiethSamplePoHundred}{31}
\newcommand{\CbTwoTwentiethSamplePoOneFifty}{39}
\newcommand{\CbTwoTwentiethTypeAtlasFifty}{40}
\newcommand{\CbTwoTwentiethTypeAtlasHundred}{52}
\newcommand{\CbTwoTwentiethTypeAtlasOneFifty}{57}
\newcommand{\CbTwoTwentiethTypeDiffFifty}{22.2}
\newcommand{\CbTwoTwentiethTypeDiffFiftyHi}{23.3}
\newcommand{\CbTwoTwentiethTypeDiffFiftyLo}{21.1}
\newcommand{\CbTwoTwentiethTypeDiffHundred}{21.0}
\newcommand{\CbTwoTwentiethTypeDiffHundredHi}{22.0}
\newcommand{\CbTwoTwentiethTypeDiffHundredLo}{19.9}
\newcommand{\CbTwoTwentiethTypeDiffOneFifty}{18.1}
\newcommand{\CbTwoTwentiethTypeDiffOneFiftyHi}{19.1}
\newcommand{\CbTwoTwentiethTypeDiffOneFiftyLo}{17.1}
\newcommand{\CbTwoTwentiethTypePoFifty}{18}
\newcommand{\CbTwoTwentiethTypePoHundred}{31}
\newcommand{\CbTwoTwentiethTypePoOneFifty}{39}
\newcommand{\CbTwoTwentiethUnitAtlasFifty}{40}
\newcommand{\CbTwoTwentiethUnitAtlasHundred}{52}
\newcommand{\CbTwoTwentiethUnitAtlasOneFifty}{57}
\newcommand{\CbTwoTwentiethUnitDiffFifty}{22.2}
\newcommand{\CbTwoTwentiethUnitDiffFiftyHi}{23.3}
\newcommand{\CbTwoTwentiethUnitDiffFiftyLo}{21.1}
\newcommand{\CbTwoTwentiethUnitDiffHundred}{21.0}
\newcommand{\CbTwoTwentiethUnitDiffHundredHi}{22.0}
\newcommand{\CbTwoTwentiethUnitDiffHundredLo}{19.9}
\newcommand{\CbTwoTwentiethUnitDiffOneFifty}{18.1}
\newcommand{\CbTwoTwentiethUnitDiffOneFiftyHi}{19.1}
\newcommand{\CbTwoTwentiethUnitDiffOneFiftyLo}{17.1}
\newcommand{\CbTwoTwentiethUnitPoFifty}{18}
\newcommand{\CbTwoTwentiethUnitPoHundred}{31}
\newcommand{\CbTwoTwentiethUnitPoOneFifty}{39}
\newcommand{\CbTypeOneBDiffFifty}{6.1}
\newcommand{\CbTypeOneBDiffFiftyHi}{7.3}
\newcommand{\CbTypeOneBDiffFiftyLo}{5.1}
\newcommand{\CbTypeOneBDiffHundred}{6.6}
\newcommand{\CbTypeOneBDiffHundredHi}{7.7}
\newcommand{\CbTypeOneBDiffHundredLo}{5.5}
\newcommand{\CbTypeOneBDiffOneFifty}{7.1}
\newcommand{\CbTypeOneBDiffOneFiftyHi}{8.3}
\newcommand{\CbTypeOneBDiffOneFiftyLo}{5.2}
\newcommand{\CbTypeOneCDEightDiffFifty}{15.2}
\newcommand{\CbTypeOneCDEightDiffFiftyHi}{16.1}
\newcommand{\CbTypeOneCDEightDiffFiftyLo}{14.1}
\newcommand{\CbTypeOneCDEightDiffHundred}{12.8}
\newcommand{\CbTypeOneCDEightDiffHundredHi}{13.9}
\newcommand{\CbTypeOneCDEightDiffHundredLo}{11.6}
\newcommand{\CbTypeOneCDEightDiffOneFifty}{10.4}
\newcommand{\CbTypeOneCDEightDiffOneFiftyHi}{11.4}
\newcommand{\CbTypeOneCDEightDiffOneFiftyLo}{9.4}
\newcommand{\CbTypeOneCDFourDiffFifty}{16.5}
\newcommand{\CbTypeOneCDFourDiffFiftyHi}{18.0}
\newcommand{\CbTypeOneCDFourDiffFiftyLo}{15.0}
\newcommand{\CbTypeOneCDFourDiffHundred}{15.2}
\newcommand{\CbTypeOneCDFourDiffHundredHi}{16.1}
\newcommand{\CbTypeOneCDFourDiffHundredLo}{13.5}
\newcommand{\CbTypeOneCDFourDiffOneFifty}{12.8}
\newcommand{\CbTypeOneCDFourDiffOneFiftyHi}{13.9}
\newcommand{\CbTypeOneCDFourDiffOneFiftyLo}{11.6}
\newcommand{\CbTypeOneMonoDiffFifty}{1.6}
\newcommand{\CbTypeOneMonoDiffFiftyHi}{2.6}
\newcommand{\CbTypeOneMonoDiffFiftyLo}{0.6}
\newcommand{\CbTypeOneMonoDiffHundred}{1.3}
\newcommand{\CbTypeOneMonoDiffHundredHi}{1.9}
\newcommand{\CbTypeOneMonoDiffHundredLo}{0.5}
\newcommand{\CbTypeOneMonoDiffOneFifty}{\ensuremath{-}0.4}
\newcommand{\CbTypeOneMonoDiffOneFiftyHi}{0.5}
\newcommand{\CbTypeOneMonoDiffOneFiftyLo}{\ensuremath{-}1.2}
\newcommand{\CbTypeOneNKDiffFifty}{7.3}
\newcommand{\CbTypeOneNKDiffFiftyHi}{8.6}
\newcommand{\CbTypeOneNKDiffFiftyLo}{6.2}
\newcommand{\CbTypeOneNKDiffHundred}{9.5}
\newcommand{\CbTypeOneNKDiffHundredHi}{10.8}
\newcommand{\CbTypeOneNKDiffHundredLo}{8.4}
\newcommand{\CbTypeOneNKDiffOneFifty}{9.5}
\newcommand{\CbTypeOneNKDiffOneFiftyHi}{10.8}
\newcommand{\CbTypeOneNKDiffOneFiftyLo}{7.9}
\newcommand{\CbTypeTwoBDiffFifty}{0.4}
\newcommand{\CbTypeTwoBDiffFiftyHi}{1.0}
\newcommand{\CbTypeTwoBDiffFiftyLo}{\ensuremath{-}0.2}
\newcommand{\CbTypeTwoBDiffHundred}{1.1}
\newcommand{\CbTypeTwoBDiffHundredHi}{2.1}
\newcommand{\CbTypeTwoBDiffHundredLo}{0.2}
\newcommand{\CbTypeTwoBDiffOneFifty}{0.7}
\newcommand{\CbTypeTwoBDiffOneFiftyHi}{1.8}
\newcommand{\CbTypeTwoBDiffOneFiftyLo}{\ensuremath{-}0.4}
\newcommand{\CbTypeTwoCDEightDiffFifty}{24.0}
\newcommand{\CbTypeTwoCDEightDiffFiftyHi}{27.3}
\newcommand{\CbTypeTwoCDEightDiffFiftyLo}{20.8}
\newcommand{\CbTypeTwoCDEightDiffHundred}{10.7}
\newcommand{\CbTypeTwoCDEightDiffHundredHi}{13.4}
\newcommand{\CbTypeTwoCDEightDiffHundredLo}{8.2}
\newcommand{\CbTypeTwoCDEightDiffOneFifty}{5.4}
\newcommand{\CbTypeTwoCDEightDiffOneFiftyHi}{7.0}
\newcommand{\CbTypeTwoCDEightDiffOneFiftyLo}{4.0}
\newcommand{\CbTypeTwoCDFourDiffFifty}{43.8}
\newcommand{\CbTypeTwoCDFourDiffFiftyHi}{47.1}
\newcommand{\CbTypeTwoCDFourDiffFiftyLo}{40.0}
\newcommand{\CbTypeTwoCDFourDiffHundred}{42.3}
\newcommand{\CbTypeTwoCDFourDiffHundredHi}{45.3}
\newcommand{\CbTypeTwoCDFourDiffHundredLo}{39.0}
\newcommand{\CbTypeTwoCDFourDiffOneFifty}{33.5}
\newcommand{\CbTypeTwoCDFourDiffOneFiftyHi}{36.2}
\newcommand{\CbTypeTwoCDFourDiffOneFiftyLo}{30.7}
\newcommand{\CbTypeTwoMonoDiffFifty}{16.2}
\newcommand{\CbTypeTwoMonoDiffFiftyHi}{18.8}
\newcommand{\CbTypeTwoMonoDiffFiftyLo}{13.3}
\newcommand{\CbTypeTwoMonoDiffHundred}{23.8}
\newcommand{\CbTypeTwoMonoDiffHundredHi}{26.0}
\newcommand{\CbTypeTwoMonoDiffHundredLo}{21.7}
\newcommand{\CbTypeTwoMonoDiffOneFifty}{25.7}
\newcommand{\CbTypeTwoMonoDiffOneFiftyHi}{27.7}
\newcommand{\CbTypeTwoMonoDiffOneFiftyLo}{23.4}
\newcommand{\CbTypeTwoNKDiffFifty}{9.2}
\newcommand{\CbTypeTwoNKDiffFiftyHi}{10.3}
\newcommand{\CbTypeTwoNKDiffFiftyLo}{8.1}
\newcommand{\CbTypeTwoNKDiffHundred}{12.2}
\newcommand{\CbTypeTwoNKDiffHundredHi}{14.6}
\newcommand{\CbTypeTwoNKDiffHundredLo}{9.8}
\newcommand{\CbTypeTwoNKDiffOneFifty}{14.0}
\newcommand{\CbTypeTwoNKDiffOneFiftyHi}{16.2}
\newcommand{\CbTypeTwoNKDiffOneFiftyLo}{11.8}
\newcommand{\CtAssayMax}{7.3}
\newcommand{\CtCellsOther}{125}
\newcommand{\CtCellsSame}{185}
\newcommand{\CtFrozenOtherFirst}{\ensuremath{-}535}
\newcommand{\CtFrozenOtherHundred}{\ensuremath{-}3.9}
\newcommand{\CtOtherFifty}{7.4}
\newcommand{\CtOtherFirst}{0.6}
\newcommand{\CtOtherHundred}{17}
\newcommand{\CtOtherMax}{35}
\newcommand{\CtOverSame}{0.68}
\newcommand{\CtOverSameHi}{0.93}
\newcommand{\CtOverSameLo}{0.61}
\newcommand{\CtPairedFirst}{0.3}
\newcommand{\CtPairedMax}{5.9}
\newcommand{\CtSaveOther}{6.4}
\newcommand{\CtSaveOtherHi}{8.3}
\newcommand{\CtSaveOtherLo}{3.7}
\newcommand{\CtVsFrozen}{1.6}
\newcommand{\CtVsFrozenHi}{1.9}
\newcommand{\CtVsFrozenLo}{1.2}
\newcommand{\DpCalConMax}{1.00}
\newcommand{\DpCalConMin}{1.00}
\newcommand{\DpCalDfMax}{0.91}
\newcommand{\DpCalDfMin}{0.64}
\newcommand{\DpCalDraws}{20}
\newcommand{\DpCellsOther}{205}
\newcommand{\DpCellsSame}{235}
\newcommand{\DpCellsSamePool}{169}
\newcommand{\DpDOne}{met}
\newcommand{\DpDThree}{met}
\newcommand{\DpDTwo}{not met}
\newcommand{\DpFirstBudget}{25}
\newcommand{\DpFrozenAt}{19:35 EDT on 28 September 2026}
\newcommand{\DpLevel}{36}
\newcommand{\DpMaxBudget}{800}
\newcommand{\DpNFolds}{twelve}
\newcommand{\DpNSites}{four}
\newcommand{\DpNonInf}{0.87}
\newcommand{\DpNonInfHi}{0.96}
\newcommand{\DpNonInfLo}{0.81}
\newcommand{\DpNonInfQuarter}{0.93}
\newcommand{\DpNonInfQuarterHi}{0.99}
\newcommand{\DpNonInfQuarterLo}{0.89}
\newcommand{\DpNonInfThreeQ}{0.83}
\newcommand{\DpNonInfThreeQHi}{0.88}
\newcommand{\DpNonInfThreeQLo}{0.75}
\newcommand{\DpOverlapAssay}{0.53}
\newcommand{\DpOverlapAssayTwenty}{0.56}
\newcommand{\DpOverlapOther}{0.74}
\newcommand{\DpOverlapOtherTwenty}{0.71}
\newcommand{\DpOwnNegativeBelow}{200}
\newcommand{\DpOwnNegativeMax}{100}
\newcommand{\DpPoConFifty}{0.9}
\newcommand{\DpPoConHundred}{2.1}
\newcommand{\DpPoConTwentyFive}{0.3}
\newcommand{\DpPoDfFifty}{0.0}
\newcommand{\DpPoDfHundred}{0.9}
\newcommand{\DpPoDfTwentyFive}{\ensuremath{-}16}
\newcommand{\DpPoolOtherMax}{17,621}
\newcommand{\DpPoolOtherMin}{17,369}
\newcommand{\DpPoolSameMax}{4,514}
\newcommand{\DpPoolSameMin}{4,207}
\newcommand{\DpRAssay}{0.59}
\newcommand{\DpROther}{0.80}
\newcommand{\DpReservoir}{1,471}
\newcommand{\DpRfAssayMax}{7.3}
\newcommand{\DpRfMeansBest}{\ensuremath{-}443}
\newcommand{\DpRfOtherMax}{36}
\newcommand{\DpRfPairedMax}{5.9}
\newcommand{\DpRfSameMax}{34}
\newcommand{\DpRfSamePoolFirst}{1.9}
\newcommand{\DpSaveOther}{3.9}
\newcommand{\DpSaveOtherHi}{4.5}
\newcommand{\DpSaveOtherLo}{3.3}
\newcommand{\DpSaveOtherQuarter}{5.3}
\newcommand{\DpSaveOtherQuarterHi}{5.9}
\newcommand{\DpSaveOtherQuarterLo}{4.7}
\newcommand{\DpSaveOtherThreeQ}{2.3}
\newcommand{\DpSaveOtherThreeQHi}{2.6}
\newcommand{\DpSaveOtherThreeQLo}{2.1}
\newcommand{\DpSaveSame}{3.4}
\newcommand{\DpSaveSamePool}{4.7}
\newcommand{\DpSiteMax}{4.4}
\newcommand{\DpSiteMin}{2.0}
\newcommand{\DpTarget}{18}
\newcommand{\FxDiffPaired}{1.4}
\newcommand{\FxDiffProposed}{0.5}
\newcommand{\FxOwnHundred}{60}
\newcommand{\FxOwnMax}{74}
\newcommand{\FxPairedOwnHundred}{36}
\newcommand{\FxPairedPoolHundred}{\ensuremath{-}52}
\newcommand{\FxPoolHundred}{\ensuremath{-}78}
\newcommand{\FxSameHundred}{56}
\newcommand{\LwAssay}{3.0}
\newcommand{\LwAssayHi}{4.3}
\newcommand{\LwAssayLo}{2.0}
\newcommand{\LwOther}{4.3}
\newcommand{\LwOtherHi}{6.2}
\newcommand{\LwOtherLo}{3.3}
\newcommand{\LwPiAssay}{11}
\newcommand{\LwPiOther}{14}
\newcommand{\LwPiSame}{15}
\newcommand{\LwSame}{6.9}
\newcommand{\LwSameHi}{10.5}
\newcommand{\LwSameLo}{4.0}
\newcommand{\LwWAssay}{45}
\newcommand{\LwWOther}{68}
\newcommand{\LwWSame}{77}
\newcommand{\MpAssayMapped}{7.3}
\newcommand{\MpAssayOtherMap}{26}
\newcommand{\MpAssayUnmapped}{6.9}
\newcommand{\MpOtherMapped}{35}
\newcommand{\MpOtherUnmapped}{7.0}
\newcommand{\MpSameUnmapped}{6.0}
\newcommand{\PrNAbove}{10}
\newcommand{\PrNBelow}{9}
\newcommand{\PrNNeither}{6}
\newcommand{\PsSmallFifty}{29}
\newcommand{\PsSmallTwentyFive}{60}
\newcommand{\RbCellCountFewBMono}{47,683}
\newcommand{\RbCellCountFewT}{34,942}
\newcommand{\RbCellCountFull}{63,796}
\newcommand{\RbCellCountPatientsFifteen}{8,231}
\newcommand{\RbCellCountPatientsFiftyNine}{32,350}
\newcommand{\RbCellCountPatientsFour}{2,132}
\newcommand{\RbCellCountPatientsSeven}{3,823}
\newcommand{\RbCellCountPatientsThirty}{16,510}
\newcommand{\RbCellCountWithoutB}{55,371}
\newcommand{\RbCellCountWithoutCDEightT}{52,250}
\newcommand{\RbCellCountWithoutCDFourT}{42,339}
\newcommand{\RbCellCountWithoutMono}{53,494}
\newcommand{\RbCellCountWithoutNK}{51,730}
\newcommand{\RbCellsFewBMono}{$>$800}
\newcommand{\RbCellsFewBMonoHi}{800}
\newcommand{\RbCellsFewBMonoLo}{800}
\newcommand{\RbCellsFewT}{$>$800}
\newcommand{\RbCellsFewTHi}{800}
\newcommand{\RbCellsFewTLo}{800}
\newcommand{\RbCellsFull}{147}
\newcommand{\RbCellsFullHi}{182}
\newcommand{\RbCellsFullLo}{126}
\newcommand{\RbCellsPatientsFifteen}{185}
\newcommand{\RbCellsPatientsFifteenHi}{246}
\newcommand{\RbCellsPatientsFifteenLo}{155}
\newcommand{\RbCellsPatientsFiftyNine}{174}
\newcommand{\RbCellsPatientsFiftyNineHi}{217}
\newcommand{\RbCellsPatientsFiftyNineLo}{147}
\newcommand{\RbCellsPatientsFour}{$>$800}
\newcommand{\RbCellsPatientsFourHi}{800}
\newcommand{\RbCellsPatientsFourLo}{800}
\newcommand{\RbCellsPatientsSeven}{385}
\newcommand{\RbCellsPatientsSevenHi}{800}
\newcommand{\RbCellsPatientsSevenLo}{281}
\newcommand{\RbCellsPatientsThirty}{177}
\newcommand{\RbCellsPatientsThirtyHi}{219}
\newcommand{\RbCellsPatientsThirtyLo}{151}
\newcommand{\RbCellsWithoutB}{174}
\newcommand{\RbCellsWithoutBHi}{222}
\newcommand{\RbCellsWithoutBLo}{147}
\newcommand{\RbCellsWithoutCDEightT}{$>$800}
\newcommand{\RbCellsWithoutCDEightTHi}{800}
\newcommand{\RbCellsWithoutCDEightTLo}{800}
\newcommand{\RbCellsWithoutCDFourT}{159}
\newcommand{\RbCellsWithoutCDFourTHi}{201}
\newcommand{\RbCellsWithoutCDFourTLo}{132}
\newcommand{\RbCellsWithoutMono}{209}
\newcommand{\RbCellsWithoutMonoHi}{316}
\newcommand{\RbCellsWithoutMonoLo}{162}
\newcommand{\RbCellsWithoutNK}{184}
\newcommand{\RbCellsWithoutNKHi}{241}
\newcommand{\RbCellsWithoutNKLo}{152}
\newcommand{\RbDonorAtlasMax}{741}
\newcommand{\RbDonorMax}{8.1}
\newcommand{\RbDonorMin}{1.08}
\newcommand{\RbDonorsCensored}{15}
\newcommand{\RbDonorsPoReached}{three}
\newcommand{\RbLopoMax}{5.6}
\newcommand{\RbLopoMin}{5.2}
\newcommand{\RbLopoOwnMax}{1.42}
\newcommand{\RbLopoOwnMin}{1.24}
\newcommand{\RbMappedFewBMono}{20}
\newcommand{\RbMappedFewT}{\ensuremath{-}2.7}
\newcommand{\RbMappedFull}{48}
\newcommand{\RbMappedMaxFewBMono}{20}
\newcommand{\RbMappedMaxFewT}{14}
\newcommand{\RbMappedMaxFull}{48}
\newcommand{\RbNDonors}{18}
\newcommand{\RbPoolMax}{8.4}
\newcommand{\RbPoolMin}{3.4}
\newcommand{\RbPopsFewBMono}{395}
\newcommand{\RbPopsFewT}{351}
\newcommand{\RbPopsFull}{566}
\newcommand{\RbPopsPatientsFifteen}{73}
\newcommand{\RbPopsPatientsFiftyNine}{281}
\newcommand{\RbPopsPatientsFour}{17}
\newcommand{\RbPopsPatientsSeven}{34}
\newcommand{\RbPopsPatientsThirty}{145}
\newcommand{\RbPopsWithoutB}{457}
\newcommand{\RbPopsWithoutCDEightT}{444}
\newcommand{\RbPopsWithoutCDFourT}{440}
\newcommand{\RbPopsWithoutMono}{480}
\newcommand{\RbPopsWithoutNK}{443}
\newcommand{\RbReproduces}{0}
\newcommand{\RbRfHundredFewBMono}{2.8}
\newcommand{\RbRfHundredFewT}{14}
\newcommand{\RbRfHundredFull}{23}
\newcommand{\RbRfHundredPatientsFifteen}{20}
\newcommand{\RbRfHundredPatientsFiftyNine}{20}
\newcommand{\RbRfHundredPatientsFour}{9.5}
\newcommand{\RbRfHundredPatientsSeven}{16}
\newcommand{\RbRfHundredPatientsThirty}{20}
\newcommand{\RbRfHundredWithoutB}{21}
\newcommand{\RbRfHundredWithoutCDEightT}{9.8}
\newcommand{\RbRfHundredWithoutCDFourT}{23}
\newcommand{\RbRfHundredWithoutMono}{20}
\newcommand{\RbRfHundredWithoutNK}{21}
\newcommand{\RbRfMaxFewBMono}{20}
\newcommand{\RbRfMaxFewT}{14}
\newcommand{\RbRfMaxFull}{48}
\newcommand{\RbRfMaxPatientsFifteen}{41}
\newcommand{\RbRfMaxPatientsFiftyNine}{46}
\newcommand{\RbRfMaxPatientsFour}{15}
\newcommand{\RbRfMaxPatientsSeven}{31}
\newcommand{\RbRfMaxPatientsThirty}{44}
\newcommand{\RbRfMaxWithoutB}{43}
\newcommand{\RbRfMaxWithoutCDEightT}{25}
\newcommand{\RbRfMaxWithoutCDFourT}{43}
\newcommand{\RbRfMaxWithoutMono}{38}
\newcommand{\RbRfMaxWithoutNK}{41}
\newcommand{\RbSaveFewBMono}{none}
\newcommand{\RbSaveFewBMonoHi}{1.00}
\newcommand{\RbSaveFewBMonoLo}{1.00}
\newcommand{\RbSaveFewT}{none}
\newcommand{\RbSaveFewTHi}{1.00}
\newcommand{\RbSaveFewTLo}{1.00}
\newcommand{\RbSaveFull}{at least 5.4}
\newcommand{\RbSaveFullHi}{6.3}
\newcommand{\RbSaveFullLo}{4.4}
\newcommand{\RbSavePatientsFifteen}{at least 4.3}
\newcommand{\RbSavePatientsFifteenHi}{5.2}
\newcommand{\RbSavePatientsFifteenLo}{3.3}
\newcommand{\RbSavePatientsFiftyNine}{at least 4.6}
\newcommand{\RbSavePatientsFiftyNineHi}{5.4}
\newcommand{\RbSavePatientsFiftyNineLo}{3.7}
\newcommand{\RbSavePatientsFour}{none}
\newcommand{\RbSavePatientsFourHi}{1.00}
\newcommand{\RbSavePatientsFourLo}{1.00}
\newcommand{\RbSavePatientsSeven}{at least 2.1}
\newcommand{\RbSavePatientsSevenHi}{2.8}
\newcommand{\RbSavePatientsSevenLo}{1.00}
\newcommand{\RbSavePatientsThirty}{at least 4.5}
\newcommand{\RbSavePatientsThirtyHi}{5.3}
\newcommand{\RbSavePatientsThirtyLo}{3.6}
\newcommand{\RbSaveWithoutB}{at least 4.6}
\newcommand{\RbSaveWithoutBHi}{5.5}
\newcommand{\RbSaveWithoutBLo}{3.6}
\newcommand{\RbSaveWithoutCDEightT}{none}
\newcommand{\RbSaveWithoutCDEightTHi}{1.00}
\newcommand{\RbSaveWithoutCDEightTLo}{1.00}
\newcommand{\RbSaveWithoutCDFourT}{at least 5.0}
\newcommand{\RbSaveWithoutCDFourTHi}{6.1}
\newcommand{\RbSaveWithoutCDFourTLo}{4.0}
\newcommand{\RbSaveWithoutMono}{at least 3.8}
\newcommand{\RbSaveWithoutMonoHi}{4.9}
\newcommand{\RbSaveWithoutMonoLo}{2.5}
\newcommand{\RbSaveWithoutNK}{at least 4.4}
\newcommand{\RbSaveWithoutNKHi}{5.3}
\newcommand{\RbSaveWithoutNKLo}{3.3}
\newcommand{\RbUnmapFewBMono}{14}
\newcommand{\RbUnmapFewT}{13}
\newcommand{\RbUnmapFull}{14}
\newcommand{\RvAtlasBaseCells}{147}
\newcommand{\RvAtlasBaseKTwo}{1.13}
\newcommand{\RvAtlasCellsMax}{170}
\newcommand{\RvAtlasCellsMin}{141}
\newcommand{\RvAtlasKTwoAbove}{8}
\newcommand{\RvAtlasKTwoHi}{1.40}
\newcommand{\RvAtlasKTwoLo}{0.82}
\newcommand{\RvAtlasKTwoMax}{1.27}
\newcommand{\RvAtlasKTwoMed}{1.07}
\newcommand{\RvAtlasKTwoMin}{0.99}
\newcommand{\RvAtlasKTwoShare}{61}
\newcommand{\RvAtlasReps}{10}
\newcommand{\RvAtlasVOneHi}{6.3}
\newcommand{\RvAtlasVOneLo}{3.9}
\newcommand{\RvAtlasVOneMax}{5.7}
\newcommand{\RvAtlasVOneMin}{4.7}
\newcommand{\RvAxesCells}{more than 800}
\newcommand{\RvAxesCellsHi}{800}
\newcommand{\RvAxesCellsLo}{800}
\newcommand{\RvAxesEight}{14}
\newcommand{\RvAxesEightHi}{16}
\newcommand{\RvAxesEightLo}{12}
\newcommand{\RvAxesFifty}{4.0}
\newcommand{\RvAxesFiftyHi}{4.8}
\newcommand{\RvAxesFiftyLo}{3.2}
\newcommand{\RvAxesFourHundred}{13}
\newcommand{\RvAxesFourHundredHi}{14}
\newcommand{\RvAxesFourHundredLo}{11}
\newcommand{\RvAxesHundred}{7.8}
\newcommand{\RvAxesHundredHi}{8.9}
\newcommand{\RvAxesHundredLo}{6.6}
\newcommand{\RvAxesOneFifty}{9.3}
\newcommand{\RvAxesOneFiftyHi}{11}
\newcommand{\RvAxesOneFiftyLo}{7.6}
\newcommand{\RvAxesSave}{at least 1.0}
\newcommand{\RvAxesSaveHi}{1.0}
\newcommand{\RvAxesSaveLo}{1.0}
\newcommand{\RvAxesTwentyFive}{\ensuremath{-}0.7}
\newcommand{\RvAxesTwentyFiveHi}{0.5}
\newcommand{\RvAxesTwentyFiveLo}{\ensuremath{-}2.5}
\newcommand{\RvAxesTwoHundred}{10}
\newcommand{\RvAxesTwoHundredHi}{12}
\newcommand{\RvAxesTwoHundredLo}{8.8}
\newcommand{\RvBothCells}{145}
\newcommand{\RvBothCellsHi}{183}
\newcommand{\RvBothCellsLo}{124}
\newcommand{\RvBothEight}{48}
\newcommand{\RvBothEightHi}{52}
\newcommand{\RvBothEightLo}{44}
\newcommand{\RvBothFifty}{11}
\newcommand{\RvBothFiftyHi}{14}
\newcommand{\RvBothFiftyLo}{8.8}
\newcommand{\RvBothFourHundred}{42}
\newcommand{\RvBothFourHundredHi}{45}
\newcommand{\RvBothFourHundredLo}{38}
\newcommand{\RvBothHundred}{23}
\newcommand{\RvBothHundredHi}{25}
\newcommand{\RvBothHundredLo}{20}
\newcommand{\RvBothOneFifty}{29}
\newcommand{\RvBothOneFiftyHi}{32}
\newcommand{\RvBothOneFiftyLo}{25}
\newcommand{\RvBothSave}{at least 5.5}
\newcommand{\RvBothSaveHi}{6.5}
\newcommand{\RvBothSaveLo}{4.4}
\newcommand{\RvBothTwentyFive}{0.5}
\newcommand{\RvBothTwentyFiveHi}{2.3}
\newcommand{\RvBothTwentyFiveLo}{\ensuremath{-}2.0}
\newcommand{\RvBothTwoHundred}{33}
\newcommand{\RvBothTwoHundredHi}{36}
\newcommand{\RvBothTwoHundredLo}{30}
\newcommand{\RvCalBmBlockMax}{1.31}
\newcommand{\RvCalBmBlockMin}{1.00}
\newcommand{\RvCalBmJackGainMax}{6.0}
\newcommand{\RvCalBmJackGainMin}{0.9}
\newcommand{\RvCalBmJackMax}{1.54}
\newcommand{\RvCalBmJackMin}{1.05}
\newcommand{\RvCalBmMax}{1.029}
\newcommand{\RvCalBmMin}{1.005}
\newcommand{\RvCalBmOrderCvEight}{4.4}
\newcommand{\RvCalBmOrderCvHundred}{9.7}
\newcommand{\RvCalBmOrderCvMax}{12}
\newcommand{\RvCalBmOrderFactorEight}{0.015}
\newcommand{\RvCalBmOrderFactorHundred}{0.11}
\newcommand{\RvCalBmOrderFactorMax}{0.22}
\newcommand{\RvCalBmOrthoMax}{1.034}
\newcommand{\RvCalBmOrthoMin}{1.005}
\newcommand{\RvCalBmOwnMax}{1.01}
\newcommand{\RvCalBmOwnMin}{0.65}
\newcommand{\RvCalValBlockMax}{1.07}
\newcommand{\RvCalValBlockMin}{1.00}
\newcommand{\RvCalValJackGainMax}{2.9}
\newcommand{\RvCalValJackGainMin}{0.6}
\newcommand{\RvCalValJackMax}{1.48}
\newcommand{\RvCalValJackMin}{1.03}
\newcommand{\RvCalValMax}{1.016}
\newcommand{\RvCalValMin}{1.005}
\newcommand{\RvCalValOrderCvEight}{2.8}
\newcommand{\RvCalValOrderCvHundred}{8.5}
\newcommand{\RvCalValOrderCvMax}{12}
\newcommand{\RvCalValOrderFactorEight}{0.016}
\newcommand{\RvCalValOrderFactorHundred}{0.10}
\newcommand{\RvCalValOrderFactorMax}{0.22}
\newcommand{\RvCalValOrthoMax}{1.010}
\newcommand{\RvCalValOrthoMin}{1.003}
\newcommand{\RvCalValOwnMax}{1.00}
\newcommand{\RvCalValOwnMin}{0.67}
\newcommand{\RvDonorBetterPo}{18}
\newcommand{\RvDonorBetterReg}{10}
\newcommand{\RvDonorMinusPoMax}{16}
\newcommand{\RvDonorMinusPoMin}{6.8}
\newcommand{\RvDonorMinusRegMax}{1.8}
\newcommand{\RvDonorMinusRegMin}{\ensuremath{-}2.0}
\newcommand{\RvDonorN}{18}
\newcommand{\RvDrawCellsMax}{154}
\newcommand{\RvDrawCellsMin}{133}
\newcommand{\RvDrawFileIdentical}{identical}
\newcommand{\RvDrawInflMax}{537}
\newcommand{\RvDrawInflMedianMax}{29}
\newcommand{\RvDrawInflMedianMin}{2}
\newcommand{\RvDrawKTwoHi}{1.45}
\newcommand{\RvDrawKTwoLo}{0.85}
\newcommand{\RvDrawKTwoMax}{1.17}
\newcommand{\RvDrawKTwoMed}{1.04}
\newcommand{\RvDrawKTwoMin}{0.95}
\newcommand{\RvDrawKTwoShare}{63}
\newcommand{\RvDrawMaxDiff}{$4\times10^{-9}$}
\newcommand{\RvDrawTests}{23}
\newcommand{\RvDrawVOneHi}{6.7}
\newcommand{\RvDrawVOneLo}{4.2}
\newcommand{\RvDrawVOneMax}{6.0}
\newcommand{\RvDrawVOneMed}{5.5}
\newcommand{\RvDrawVOneMin}{5.2}
\newcommand{\RvDrawVOneShare}{100}
\newcommand{\RvFewBMonoAxesCells}{more than 800}
\newcommand{\RvFewBMonoAxesCellsHi}{800}
\newcommand{\RvFewBMonoAxesCellsLo}{800}
\newcommand{\RvFewBMonoAxesEight}{14}
\newcommand{\RvFewBMonoAxesEightHi}{16}
\newcommand{\RvFewBMonoAxesEightLo}{12}
\newcommand{\RvFewBMonoAxesFifty}{4.2}
\newcommand{\RvFewBMonoAxesFiftyHi}{5.1}
\newcommand{\RvFewBMonoAxesFiftyLo}{3.4}
\newcommand{\RvFewBMonoAxesFourHundred}{12}
\newcommand{\RvFewBMonoAxesFourHundredHi}{14}
\newcommand{\RvFewBMonoAxesFourHundredLo}{11}
\newcommand{\RvFewBMonoAxesHundred}{8.0}
\newcommand{\RvFewBMonoAxesHundredHi}{9.1}
\newcommand{\RvFewBMonoAxesHundredLo}{6.7}
\newcommand{\RvFewBMonoAxesOneFifty}{9.3}
\newcommand{\RvFewBMonoAxesOneFiftyHi}{11}
\newcommand{\RvFewBMonoAxesOneFiftyLo}{7.5}
\newcommand{\RvFewBMonoAxesSave}{at least 1.0}
\newcommand{\RvFewBMonoAxesSaveHi}{1.0}
\newcommand{\RvFewBMonoAxesSaveLo}{1.0}
\newcommand{\RvFewBMonoAxesTwentyFive}{\ensuremath{-}0.5}
\newcommand{\RvFewBMonoAxesTwentyFiveHi}{0.8}
\newcommand{\RvFewBMonoAxesTwentyFiveLo}{\ensuremath{-}2.3}
\newcommand{\RvFewBMonoAxesTwoHundred}{10}
\newcommand{\RvFewBMonoAxesTwoHundredHi}{12}
\newcommand{\RvFewBMonoAxesTwoHundredLo}{8.6}
\newcommand{\RvFewBMonoMappedCells}{more than 800}
\newcommand{\RvFewBMonoMappedCellsHi}{800}
\newcommand{\RvFewBMonoMappedCellsLo}{800}
\newcommand{\RvFewBMonoMappedEight}{20}
\newcommand{\RvFewBMonoMappedEightHi}{26}
\newcommand{\RvFewBMonoMappedEightLo}{13}
\newcommand{\RvFewBMonoMappedFifty}{\ensuremath{-}6.1}
\newcommand{\RvFewBMonoMappedFiftyHi}{\ensuremath{-}2.4}
\newcommand{\RvFewBMonoMappedFiftyLo}{\ensuremath{-}11}
\newcommand{\RvFewBMonoMappedFourHundred}{17}
\newcommand{\RvFewBMonoMappedFourHundredHi}{22}
\newcommand{\RvFewBMonoMappedFourHundredLo}{12}
\newcommand{\RvFewBMonoMappedHundred}{2.8}
\newcommand{\RvFewBMonoMappedHundredHi}{5.9}
\newcommand{\RvFewBMonoMappedHundredLo}{\ensuremath{-}0.6}
\newcommand{\RvFewBMonoMappedOneFifty}{5.7}
\newcommand{\RvFewBMonoMappedOneFiftyHi}{11}
\newcommand{\RvFewBMonoMappedOneFiftyLo}{\ensuremath{-}0.5}
\newcommand{\RvFewBMonoMappedSave}{at least 1.0}
\newcommand{\RvFewBMonoMappedSaveHi}{1.0}
\newcommand{\RvFewBMonoMappedSaveLo}{1.0}
\newcommand{\RvFewBMonoMappedTwentyFive}{\ensuremath{-}13}
\newcommand{\RvFewBMonoMappedTwentyFiveHi}{\ensuremath{-}9.9}
\newcommand{\RvFewBMonoMappedTwentyFiveLo}{\ensuremath{-}17}
\newcommand{\RvFewBMonoMappedTwoHundred}{10.0}
\newcommand{\RvFewBMonoMappedTwoHundredHi}{15}
\newcommand{\RvFewBMonoMappedTwoHundredLo}{5.1}
\newcommand{\RvFewBMonoRwCells}{148}
\newcommand{\RvFewBMonoRwCellsHi}{192}
\newcommand{\RvFewBMonoRwCellsLo}{122}
\newcommand{\RvFewBMonoRwEight}{46}
\newcommand{\RvFewBMonoRwEightHi}{49}
\newcommand{\RvFewBMonoRwEightLo}{42}
\newcommand{\RvFewBMonoRwFifty}{10.0}
\newcommand{\RvFewBMonoRwFiftyHi}{13}
\newcommand{\RvFewBMonoRwFiftyLo}{6.9}
\newcommand{\RvFewBMonoRwFourHundred}{40}
\newcommand{\RvFewBMonoRwFourHundredHi}{43}
\newcommand{\RvFewBMonoRwFourHundredLo}{37}
\newcommand{\RvFewBMonoRwHundred}{23}
\newcommand{\RvFewBMonoRwHundredHi}{25}
\newcommand{\RvFewBMonoRwHundredLo}{20}
\newcommand{\RvFewBMonoRwOneFifty}{29}
\newcommand{\RvFewBMonoRwOneFiftyHi}{32}
\newcommand{\RvFewBMonoRwOneFiftyLo}{25}
\newcommand{\RvFewBMonoRwSave}{at least 5.4}
\newcommand{\RvFewBMonoRwSaveHi}{6.5}
\newcommand{\RvFewBMonoRwSaveLo}{4.2}
\newcommand{\RvFewBMonoRwTwentyFive}{\ensuremath{-}25}
\newcommand{\RvFewBMonoRwTwentyFiveHi}{\ensuremath{-}4.3}
\newcommand{\RvFewBMonoRwTwentyFiveLo}{\ensuremath{-}64}
\newcommand{\RvFewBMonoRwTwoHundred}{32}
\newcommand{\RvFewBMonoRwTwoHundredHi}{35}
\newcommand{\RvFewBMonoRwTwoHundredLo}{29}
\newcommand{\RvFewTAxesCells}{more than 800}
\newcommand{\RvFewTAxesCellsHi}{800}
\newcommand{\RvFewTAxesCellsLo}{800}
\newcommand{\RvFewTAxesEight}{13}
\newcommand{\RvFewTAxesEightHi}{15}
\newcommand{\RvFewTAxesEightLo}{11}
\newcommand{\RvFewTAxesFifty}{2.5}
\newcommand{\RvFewTAxesFiftyHi}{3.2}
\newcommand{\RvFewTAxesFiftyLo}{1.8}
\newcommand{\RvFewTAxesFourHundred}{12}
\newcommand{\RvFewTAxesFourHundredHi}{13}
\newcommand{\RvFewTAxesFourHundredLo}{10.0}
\newcommand{\RvFewTAxesHundred}{5.9}
\newcommand{\RvFewTAxesHundredHi}{6.9}
\newcommand{\RvFewTAxesHundredLo}{4.9}
\newcommand{\RvFewTAxesOneFifty}{7.6}
\newcommand{\RvFewTAxesOneFiftyHi}{8.8}
\newcommand{\RvFewTAxesOneFiftyLo}{6.1}
\newcommand{\RvFewTAxesSave}{at least 1.0}
\newcommand{\RvFewTAxesSaveHi}{1.0}
\newcommand{\RvFewTAxesSaveLo}{1.0}
\newcommand{\RvFewTAxesTwentyFive}{\ensuremath{-}1.3}
\newcommand{\RvFewTAxesTwentyFiveHi}{\ensuremath{-}0.3}
\newcommand{\RvFewTAxesTwentyFiveLo}{\ensuremath{-}2.7}
\newcommand{\RvFewTAxesTwoHundred}{8.9}
\newcommand{\RvFewTAxesTwoHundredHi}{10}
\newcommand{\RvFewTAxesTwoHundredLo}{7.5}
\newcommand{\RvFewTMappedCells}{more than 800}
\newcommand{\RvFewTMappedCellsHi}{800}
\newcommand{\RvFewTMappedCellsLo}{800}
\newcommand{\RvFewTMappedEight}{\ensuremath{-}2.7}
\newcommand{\RvFewTMappedEightHi}{12}
\newcommand{\RvFewTMappedEightLo}{\ensuremath{-}19}
\newcommand{\RvFewTMappedFifty}{7.2}
\newcommand{\RvFewTMappedFiftyHi}{10.0}
\newcommand{\RvFewTMappedFiftyLo}{4.0}
\newcommand{\RvFewTMappedFourHundred}{8.8}
\newcommand{\RvFewTMappedFourHundredHi}{20}
\newcommand{\RvFewTMappedFourHundredLo}{\ensuremath{-}3.9}
\newcommand{\RvFewTMappedHundred}{14}
\newcommand{\RvFewTMappedHundredHi}{18}
\newcommand{\RvFewTMappedHundredLo}{9.3}
\newcommand{\RvFewTMappedOneFifty}{15}
\newcommand{\RvFewTMappedOneFiftyHi}{21}
\newcommand{\RvFewTMappedOneFiftyLo}{7.0}
\newcommand{\RvFewTMappedSave}{at least 1.0}
\newcommand{\RvFewTMappedSaveHi}{1.0}
\newcommand{\RvFewTMappedSaveLo}{1.0}
\newcommand{\RvFewTMappedTwentyFive}{\ensuremath{-}3.3}
\newcommand{\RvFewTMappedTwentyFiveHi}{0.3}
\newcommand{\RvFewTMappedTwentyFiveLo}{\ensuremath{-}8.1}
\newcommand{\RvFewTMappedTwoHundred}{14}
\newcommand{\RvFewTMappedTwoHundredHi}{22}
\newcommand{\RvFewTMappedTwoHundredLo}{4.8}
\newcommand{\RvFewTRwCells}{311}
\newcommand{\RvFewTRwCellsHi}{378}
\newcommand{\RvFewTRwCellsLo}{266}
\newcommand{\RvFewTRwEight}{41}
\newcommand{\RvFewTRwEightHi}{43}
\newcommand{\RvFewTRwEightLo}{37}
\newcommand{\RvFewTRwFifty}{0.6}
\newcommand{\RvFewTRwFiftyHi}{3.6}
\newcommand{\RvFewTRwFiftyLo}{\ensuremath{-}3.2}
\newcommand{\RvFewTRwFourHundred}{32}
\newcommand{\RvFewTRwFourHundredHi}{35}
\newcommand{\RvFewTRwFourHundredLo}{29}
\newcommand{\RvFewTRwHundred}{12}
\newcommand{\RvFewTRwHundredHi}{14}
\newcommand{\RvFewTRwHundredLo}{9.6}
\newcommand{\RvFewTRwOneFifty}{18}
\newcommand{\RvFewTRwOneFiftyHi}{20}
\newcommand{\RvFewTRwOneFiftyLo}{14}
\newcommand{\RvFewTRwSave}{at least 2.6}
\newcommand{\RvFewTRwSaveHi}{3.0}
\newcommand{\RvFewTRwSaveLo}{2.1}
\newcommand{\RvFewTRwTwentyFive}{\ensuremath{-}35}
\newcommand{\RvFewTRwTwentyFiveHi}{\ensuremath{-}11}
\newcommand{\RvFewTRwTwentyFiveLo}{\ensuremath{-}72}
\newcommand{\RvFewTRwTwoHundred}{22}
\newcommand{\RvFewTRwTwoHundredHi}{24}
\newcommand{\RvFewTRwTwoHundredLo}{19}
\newcommand{\RvFineOneAtlas}{74}
\newcommand{\RvFineOneMinusPo}{2.8}
\newcommand{\RvFineOneMinusPoHi}{3.4}
\newcommand{\RvFineOneMinusPoLo}{2.1}
\newcommand{\RvFineOneMinusReg}{\ensuremath{-}2.0}
\newcommand{\RvFineOneMinusRegHi}{\ensuremath{-}1.3}
\newcommand{\RvFineOneMinusRegLo}{\ensuremath{-}2.7}
\newcommand{\RvFineOnePo}{71}
\newcommand{\RvFineOneReg}{76}
\newcommand{\RvFineReps}{31,509}
\newcommand{\RvFineRfAtlas}{\ensuremath{-}78}
\newcommand{\RvFineRfMinusPo}{\ensuremath{-}71}
\newcommand{\RvFineRfMinusPoHi}{\ensuremath{-}58}
\newcommand{\RvFineRfMinusPoLo}{\ensuremath{-}87}
\newcommand{\RvFineRfMinusReg}{\ensuremath{-}70}
\newcommand{\RvFineRfMinusRegHi}{\ensuremath{-}57}
\newcommand{\RvFineRfMinusRegLo}{\ensuremath{-}85}
\newcommand{\RvFineRfPo}{\ensuremath{-}6.6}
\newcommand{\RvFineRfReg}{\ensuremath{-}7.6}
\newcommand{\RvFineTwoAtlas}{17}
\newcommand{\RvFineTwoMinusPo}{6.5}
\newcommand{\RvFineTwoMinusPoHi}{7.0}
\newcommand{\RvFineTwoMinusPoLo}{6.1}
\newcommand{\RvFineTwoMinusReg}{\ensuremath{-}1.9}
\newcommand{\RvFineTwoMinusRegHi}{\ensuremath{-}1.6}
\newcommand{\RvFineTwoMinusRegLo}{\ensuremath{-}2.2}
\newcommand{\RvFineTwoPo}{10}
\newcommand{\RvFineTwoReg}{19}
\newcommand{\RvFineUnits}{170}
\newcommand{\RvFisherFive}{4.90}
\newcommand{\RvFisherOne}{1.07}
\newcommand{\RvFisherOneMax}{1.10}
\newcommand{\RvFisherOneMin}{1.06}
\newcommand{\RvFisherPerm}{20}
\newcommand{\RvFisherSmallFive}{4.54}
\newcommand{\RvFisherSmallLo}{40}
\newcommand{\RvFisherSmallOne}{1.07}
\newcommand{\RvFisherSmallRange}{40 to 79}
\newcommand{\RvFisherSmallTenth}{0.18}
\newcommand{\RvFisherTenth}{0.149}
\newcommand{\RvFisherUnits}{170}
\newcommand{\RvFixedSameRunningMax}{none}
\newcommand{\RvFixedSets}{10}
\newcommand{\RvFixedTenAbove}{6}
\newcommand{\RvFixedTenBelow}{2}
\newcommand{\RvFixedTenBothAbove}{two}
\newcommand{\RvFixedTenEstBelow}{no}
\newcommand{\RvFixedTenGm}{1.7}
\newcommand{\RvFixedTenMax}{2.6}
\newcommand{\RvFixedTenMin}{1.2}
\newcommand{\RvFixedTenN}{2}
\newcommand{\RvFixedTenPoAbove}{four}
\newcommand{\RvFixedThirtyAbove}{7}
\newcommand{\RvFixedThirtyBelow}{2}
\newcommand{\RvFixedThirtyBothAbove}{four}
\newcommand{\RvFixedThirtyEstBelow}{two}
\newcommand{\RvFixedThirtyGm}{4.9}
\newcommand{\RvFixedThirtyMax}{4.9}
\newcommand{\RvFixedThirtyMin}{4.9}
\newcommand{\RvFixedThirtyN}{1}
\newcommand{\RvFixedThirtyPoAbove}{three}
\newcommand{\RvFixedTwentyAbove}{6}
\newcommand{\RvFixedTwentyBelow}{2}
\newcommand{\RvFixedTwentyBothAbove}{two}
\newcommand{\RvFixedTwentyEstBelow}{two}
\newcommand{\RvFixedTwentyGm}{2.5}
\newcommand{\RvFixedTwentyMax}{4.4}
\newcommand{\RvFixedTwentyMin}{1.4}
\newcommand{\RvFixedTwentyN}{2}
\newcommand{\RvFixedTwentyPoAbove}{four}
\newcommand{\RvFrozenAt}{04:30 EDT on 30 September 2026}
\newcommand{\RvFullRwCells}{131}
\newcommand{\RvFullRwCellsHi}{162}
\newcommand{\RvFullRwCellsLo}{114}
\newcommand{\RvFullRwEight}{50}
\newcommand{\RvFullRwEightHi}{53}
\newcommand{\RvFullRwEightLo}{47}
\newcommand{\RvFullRwFifty}{8.3}
\newcommand{\RvFullRwFiftyHi}{12}
\newcommand{\RvFullRwFiftyLo}{2.9}
\newcommand{\RvFullRwFourHundred}{44}
\newcommand{\RvFullRwFourHundredHi}{47}
\newcommand{\RvFullRwFourHundredLo}{40}
\newcommand{\RvFullRwHundred}{24}
\newcommand{\RvFullRwHundredHi}{26}
\newcommand{\RvFullRwHundredLo}{21}
\newcommand{\RvFullRwOneFifty}{31}
\newcommand{\RvFullRwOneFiftyHi}{34}
\newcommand{\RvFullRwOneFiftyLo}{27}
\newcommand{\RvFullRwSave}{at least 6.1}
\newcommand{\RvFullRwSaveHi}{7.0}
\newcommand{\RvFullRwSaveLo}{4.9}
\newcommand{\RvFullRwTwentyFive}{\ensuremath{-}37}
\newcommand{\RvFullRwTwentyFiveHi}{\ensuremath{-}8.2}
\newcommand{\RvFullRwTwentyFiveLo}{\ensuremath{-}84}
\newcommand{\RvFullRwTwoHundred}{35}
\newcommand{\RvFullRwTwoHundredHi}{38}
\newcommand{\RvFullRwTwoHundredLo}{32}
\newcommand{\RvLawHalfBelow}{8}
\newcommand{\RvLawHalfErr}{21}
\newcommand{\RvLawHalfErrHi}{29}
\newcommand{\RvLawHalfErrLo}{11}
\newcommand{\RvLawHalfErrStHi}{26}
\newcommand{\RvLawHalfErrStLo}{13}
\newcommand{\RvLawHalfLooMax}{26}
\newcommand{\RvLawHalfLooMin}{17}
\newcommand{\RvLawHalfMeansP}{0.042}
\newcommand{\RvLawHalfMeansRho}{1.00}
\newcommand{\RvLawHalfN}{10}
\newcommand{\RvLawHalfNonN}{19}
\newcommand{\RvLawHalfNotReached}{10}
\newcommand{\RvLawHalfOneAbove}{1}
\newcommand{\RvLawHalfOver}{4}
\newcommand{\RvLawHalfSets}{4}
\newcommand{\RvLawHalfSlope}{1.08}
\newcommand{\RvLawHalfSlopeHi}{1.26}
\newcommand{\RvLawHalfSlopeLo}{0.99}
\newcommand{\RvLawHalfWithinP}{0.007}
\newcommand{\RvLawHalfWithinPMin}{0.007}
\newcommand{\RvLawHalfWithinPerm}{144}
\newcommand{\RvLawSevenBelow}{0}
\newcommand{\RvLawSevenErr}{4}
\newcommand{\RvLawSevenErrHi}{5}
\newcommand{\RvLawSevenErrLo}{2}
\newcommand{\RvLawSevenErrStHi}{4}
\newcommand{\RvLawSevenErrStLo}{4}
\newcommand{\RvLawSevenLooMax}{5}
\newcommand{\RvLawSevenLooMin}{2}
\newcommand{\RvLawSevenN}{5}
\newcommand{\RvLawSevenNonN}{24}
\newcommand{\RvLawSevenNotReached}{15}
\newcommand{\RvLawSevenOneAbove}{9}
\newcommand{\RvLawSevenOver}{1}
\newcommand{\RvLawSevenSets}{2}
\newcommand{\RvLawSevenSlope}{1.02}
\newcommand{\RvLawSevenSlopeHi}{1.15}
\newcommand{\RvLawSevenSlopeLo}{0.99}
\newcommand{\RvLawSevenWithinP}{0.083}
\newcommand{\RvLawSevenWithinPMin}{0.083}
\newcommand{\RvLawSevenWithinPerm}{12}
\newcommand{\RvLawThreeBelow}{9}
\newcommand{\RvLawThreeErr}{74}
\newcommand{\RvLawThreeErrHi}{144}
\newcommand{\RvLawThreeErrLo}{54}
\newcommand{\RvLawThreeErrStHi}{144}
\newcommand{\RvLawThreeErrStLo}{61}
\newcommand{\RvLawThreeLooMax}{99}
\newcommand{\RvLawThreeLooMin}{61}
\newcommand{\RvLawThreeMeansP}{0.167}
\newcommand{\RvLawThreeMeansRho}{1.00}
\newcommand{\RvLawThreeN}{8}
\newcommand{\RvLawThreeNonN}{21}
\newcommand{\RvLawThreeNotReached}{7}
\newcommand{\RvLawThreeOneAbove}{5}
\newcommand{\RvLawThreeOver}{3}
\newcommand{\RvLawThreeSets}{3}
\newcommand{\RvLawThreeSlope}{0.98}
\newcommand{\RvLawThreeSlopeHi}{1.22}
\newcommand{\RvLawThreeSlopeLo}{-2.39}
\newcommand{\RvLawThreeWithinP}{0.028}
\newcommand{\RvLawThreeWithinPMin}{0.014}
\newcommand{\RvLawThreeWithinPerm}{72}
\newcommand{\RvMapOnlyCells}{more than 800}
\newcommand{\RvMapOnlyCellsHi}{800}
\newcommand{\RvMapOnlyCellsLo}{683}
\newcommand{\RvMapOnlyEight}{28}
\newcommand{\RvMapOnlyEightHi}{31}
\newcommand{\RvMapOnlyEightLo}{25}
\newcommand{\RvMapOnlyFifty}{1.8}
\newcommand{\RvMapOnlyFiftyHi}{2.3}
\newcommand{\RvMapOnlyFiftyLo}{1.3}
\newcommand{\RvMapOnlyFourHundred}{18}
\newcommand{\RvMapOnlyFourHundredHi}{20}
\newcommand{\RvMapOnlyFourHundredLo}{15}
\newcommand{\RvMapOnlyHundred}{4.7}
\newcommand{\RvMapOnlyHundredHi}{5.5}
\newcommand{\RvMapOnlyHundredLo}{3.9}
\newcommand{\RvMapOnlyOneFifty}{7.1}
\newcommand{\RvMapOnlyOneFiftyHi}{8.5}
\newcommand{\RvMapOnlyOneFiftyLo}{5.5}
\newcommand{\RvMapOnlySave}{at least 1.0}
\newcommand{\RvMapOnlySaveHi}{1.2}
\newcommand{\RvMapOnlySaveLo}{1.0}
\newcommand{\RvMapOnlyTwentyFive}{\ensuremath{-}0.2}
\newcommand{\RvMapOnlyTwentyFiveHi}{0.3}
\newcommand{\RvMapOnlyTwentyFiveLo}{\ensuremath{-}0.8}
\newcommand{\RvMapOnlyTwoHundred}{9.5}
\newcommand{\RvMapOnlyTwoHundredHi}{11}
\newcommand{\RvMapOnlyTwoHundredLo}{8.0}
\newcommand{\RvNeitherCells}{more than 800}
\newcommand{\RvNeitherCellsHi}{800}
\newcommand{\RvNeitherCellsLo}{800}
\newcommand{\RvNeitherEight}{9.9}
\newcommand{\RvNeitherEightHi}{11}
\newcommand{\RvNeitherEightLo}{8.5}
\newcommand{\RvNeitherFifty}{1.1}
\newcommand{\RvNeitherFiftyHi}{1.2}
\newcommand{\RvNeitherFiftyLo}{0.9}
\newcommand{\RvNeitherFourHundred}{7.1}
\newcommand{\RvNeitherFourHundredHi}{7.9}
\newcommand{\RvNeitherFourHundredLo}{6.2}
\newcommand{\RvNeitherHundred}{2.4}
\newcommand{\RvNeitherHundredHi}{2.7}
\newcommand{\RvNeitherHundredLo}{2.0}
\newcommand{\RvNeitherOneFifty}{3.4}
\newcommand{\RvNeitherOneFiftyHi}{3.9}
\newcommand{\RvNeitherOneFiftyLo}{2.8}
\newcommand{\RvNeitherTwentyFive}{0.0}
\newcommand{\RvNeitherTwentyFiveHi}{0.3}
\newcommand{\RvNeitherTwentyFiveLo}{\ensuremath{-}0.3}
\newcommand{\RvNeitherTwoHundred}{4.4}
\newcommand{\RvNeitherTwoHundredHi}{5.0}
\newcommand{\RvNeitherTwoHundredLo}{3.8}
\newcommand{\RvOtherAxesCells}{more than 800}
\newcommand{\RvOtherAxesCellsHi}{800}
\newcommand{\RvOtherAxesCellsLo}{800}
\newcommand{\RvOtherAxesEight}{14}
\newcommand{\RvOtherAxesEightHi}{17}
\newcommand{\RvOtherAxesEightLo}{12}
\newcommand{\RvOtherAxesFifty}{7.5}
\newcommand{\RvOtherAxesFiftyHi}{9.1}
\newcommand{\RvOtherAxesFiftyLo}{5.8}
\newcommand{\RvOtherAxesFourHundred}{13}
\newcommand{\RvOtherAxesFourHundredHi}{15}
\newcommand{\RvOtherAxesFourHundredLo}{11}
\newcommand{\RvOtherAxesHundred}{10}
\newcommand{\RvOtherAxesHundredHi}{12}
\newcommand{\RvOtherAxesHundredLo}{8.7}
\newcommand{\RvOtherAxesOneFifty}{11}
\newcommand{\RvOtherAxesOneFiftyHi}{13}
\newcommand{\RvOtherAxesOneFiftyLo}{8.9}
\newcommand{\RvOtherAxesSave}{at least 1.0}
\newcommand{\RvOtherAxesSaveHi}{1.0}
\newcommand{\RvOtherAxesSaveLo}{1.0}
\newcommand{\RvOtherAxesTwentyFive}{1.6}
\newcommand{\RvOtherAxesTwentyFiveHi}{3.5}
\newcommand{\RvOtherAxesTwentyFiveLo}{\ensuremath{-}1.0}
\newcommand{\RvOtherAxesTwoHundred}{12}
\newcommand{\RvOtherAxesTwoHundredHi}{14}
\newcommand{\RvOtherAxesTwoHundredLo}{9.8}
\newcommand{\RvOwnAxesCells}{more than 800}
\newcommand{\RvOwnAxesCellsHi}{800}
\newcommand{\RvOwnAxesCellsLo}{800}
\newcommand{\RvOwnAxesEight}{13}
\newcommand{\RvOwnAxesEightHi}{15}
\newcommand{\RvOwnAxesEightLo}{10}
\newcommand{\RvOwnAxesFifty}{6.5}
\newcommand{\RvOwnAxesFiftyHi}{8.3}
\newcommand{\RvOwnAxesFiftyLo}{4.9}
\newcommand{\RvOwnAxesFourHundred}{12}
\newcommand{\RvOwnAxesFourHundredHi}{14}
\newcommand{\RvOwnAxesFourHundredLo}{9.6}
\newcommand{\RvOwnAxesHundred}{9.2}
\newcommand{\RvOwnAxesHundredHi}{11}
\newcommand{\RvOwnAxesHundredLo}{7.3}
\newcommand{\RvOwnAxesOneFifty}{9.7}
\newcommand{\RvOwnAxesOneFiftyHi}{12}
\newcommand{\RvOwnAxesOneFiftyLo}{7.2}
\newcommand{\RvOwnAxesSave}{at least 1.0}
\newcommand{\RvOwnAxesSaveHi}{1.0}
\newcommand{\RvOwnAxesSaveLo}{1.0}
\newcommand{\RvOwnAxesTwentyFive}{0.5}
\newcommand{\RvOwnAxesTwentyFiveHi}{3.0}
\newcommand{\RvOwnAxesTwentyFiveLo}{\ensuremath{-}3.2}
\newcommand{\RvOwnAxesTwoHundred}{10}
\newcommand{\RvOwnAxesTwoHundredHi}{13}
\newcommand{\RvOwnAxesTwoHundredLo}{8.1}
\newcommand{\RvPilotHalfFiftyAchieved}{54}
\newcommand{\RvPilotHalfFiftyBudget}{84}
\newcommand{\RvPilotHalfFiftyFlagged}{5}
\newcommand{\RvPilotHalfFiftyProblems}{29}
\newcommand{\RvPilotHalfFiftyRatio}{1.16}
\newcommand{\RvPilotHalfFiftyReached}{64}
\newcommand{\RvPilotHalfFiftyWithin}{72}
\newcommand{\RvPilotHalfHundredAchieved}{53}
\newcommand{\RvPilotHalfHundredBudget}{100}
\newcommand{\RvPilotHalfHundredFlagged}{8}
\newcommand{\RvPilotHalfHundredProblems}{29}
\newcommand{\RvPilotHalfHundredRatio}{1.12}
\newcommand{\RvPilotHalfHundredReached}{65}
\newcommand{\RvPilotHalfHundredWithin}{75}
\newcommand{\RvPilotHalfResAchieved}{46}
\newcommand{\RvPilotHalfResBudget}{86}
\newcommand{\RvPilotHalfResFlagged}{0}
\newcommand{\RvPilotHalfResProblems}{29}
\newcommand{\RvPilotHalfResRatio}{0.77}
\newcommand{\RvPilotHalfResReached}{26}
\newcommand{\RvPilotHalfResWithin}{50}
\newcommand{\RvPilotHalfTwoHundredAchieved}{53}
\newcommand{\RvPilotHalfTwoHundredBudget}{200}
\newcommand{\RvPilotHalfTwoHundredFlagged}{10}
\newcommand{\RvPilotHalfTwoHundredProblems}{29}
\newcommand{\RvPilotHalfTwoHundredRatio}{0.87}
\newcommand{\RvPilotHalfTwoHundredReached}{65}
\newcommand{\RvPilotHalfTwoHundredWithin}{77}
\newcommand{\RvPilotSevenFiftyAchieved}{68}
\newcommand{\RvPilotSevenFiftyBudget}{181}
\newcommand{\RvPilotSevenFiftyFlagged}{22}
\newcommand{\RvPilotSevenFiftyProblems}{29}
\newcommand{\RvPilotSevenFiftyRatio}{0.82}
\newcommand{\RvPilotSevenFiftyReached}{46}
\newcommand{\RvPilotSevenFiftyWithin}{56}
\newcommand{\RvPilotSevenHundredAchieved}{69}
\newcommand{\RvPilotSevenHundredBudget}{181}
\newcommand{\RvPilotSevenHundredFlagged}{24}
\newcommand{\RvPilotSevenHundredProblems}{29}
\newcommand{\RvPilotSevenHundredRatio}{0.85}
\newcommand{\RvPilotSevenHundredReached}{48}
\newcommand{\RvPilotSevenHundredWithin}{59}
\newcommand{\RvPilotSevenResAchieved}{67}
\newcommand{\RvPilotSevenResBudget}{181}
\newcommand{\RvPilotSevenResFlagged}{0}
\newcommand{\RvPilotSevenResProblems}{29}
\newcommand{\RvPilotSevenResRatio}{0.82}
\newcommand{\RvPilotSevenResReached}{23}
\newcommand{\RvPilotSevenResWithin}{49}
\newcommand{\RvPilotSevenTwoHundredAchieved}{70}
\newcommand{\RvPilotSevenTwoHundredBudget}{200}
\newcommand{\RvPilotSevenTwoHundredFlagged}{25}
\newcommand{\RvPilotSevenTwoHundredProblems}{29}
\newcommand{\RvPilotSevenTwoHundredRatio}{0.99}
\newcommand{\RvPilotSevenTwoHundredReached}{56}
\newcommand{\RvPilotSevenTwoHundredWithin}{61}
\newcommand{\RvPilotThreeFiftyAchieved}{41}
\newcommand{\RvPilotThreeFiftyBudget}{50}
\newcommand{\RvPilotThreeFiftyFlagged}{0}
\newcommand{\RvPilotThreeFiftyProblems}{29}
\newcommand{\RvPilotThreeFiftyRatio}{0.99}
\newcommand{\RvPilotThreeFiftyReached}{73}
\newcommand{\RvPilotThreeFiftyWithin}{78}
\newcommand{\RvPilotThreeHundredAchieved}{41}
\newcommand{\RvPilotThreeHundredBudget}{100}
\newcommand{\RvPilotThreeHundredFlagged}{1}
\newcommand{\RvPilotThreeHundredProblems}{29}
\newcommand{\RvPilotThreeHundredRatio}{1.02}
\newcommand{\RvPilotThreeHundredReached}{77}
\newcommand{\RvPilotThreeHundredWithin}{82}
\newcommand{\RvPilotThreeResAchieved}{23}
\newcommand{\RvPilotThreeResBudget}{32}
\newcommand{\RvPilotThreeResFlagged}{0}
\newcommand{\RvPilotThreeResProblems}{29}
\newcommand{\RvPilotThreeResRatio}{0.57}
\newcommand{\RvPilotThreeResReached}{34}
\newcommand{\RvPilotThreeResWithin}{45}
\newcommand{\RvPilotThreeTwoHundredAchieved}{48}
\newcommand{\RvPilotThreeTwoHundredBudget}{200}
\newcommand{\RvPilotThreeTwoHundredFlagged}{1}
\newcommand{\RvPilotThreeTwoHundredProblems}{29}
\newcommand{\RvPilotThreeTwoHundredRatio}{1.16}
\newcommand{\RvPilotThreeTwoHundredReached}{83}
\newcommand{\RvPilotThreeTwoHundredWithin}{87}
\newcommand{\RvPoCells}{more than 800}
\newcommand{\RvPoCellsHi}{800}
\newcommand{\RvPoCellsLo}{800}
\newcommand{\RvPoEight}{13}
\newcommand{\RvPoEightHi}{20}
\newcommand{\RvPoEightLo}{11}
\newcommand{\RvPoFifty}{2.4}
\newcommand{\RvPoFiftyHi}{3.8}
\newcommand{\RvPoFiftyLo}{1.2}
\newcommand{\RvPoFourHundred}{14}
\newcommand{\RvPoFourHundredHi}{19}
\newcommand{\RvPoFourHundredLo}{7.9}
\newcommand{\RvPoHundred}{3.5}
\newcommand{\RvPoHundredHi}{6.1}
\newcommand{\RvPoHundredLo}{2.9}
\newcommand{\RvPoOneFifty}{5.1}
\newcommand{\RvPoOneFiftyHi}{9.3}
\newcommand{\RvPoOneFiftyLo}{3.5}
\newcommand{\RvPoTwentyFive}{0.5}
\newcommand{\RvPoTwentyFiveHi}{0.9}
\newcommand{\RvPoTwentyFiveLo}{0.4}
\newcommand{\RvPoTwoHundred}{7.7}
\newcommand{\RvPoTwoHundredHi}{12}
\newcommand{\RvPoTwoHundredLo}{4.5}
\newcommand{\RvPrimaryOneAtlas}{86}
\newcommand{\RvPrimaryOneMinusPo}{11}
\newcommand{\RvPrimaryOneMinusPoHi}{12}
\newcommand{\RvPrimaryOneMinusPoLo}{10}
\newcommand{\RvPrimaryOneMinusReg}{\ensuremath{-}0.1}
\newcommand{\RvPrimaryOneMinusRegHi}{0.2}
\newcommand{\RvPrimaryOneMinusRegLo}{\ensuremath{-}0.5}
\newcommand{\RvPrimaryOnePo}{75}
\newcommand{\RvPrimaryOneReg}{86}
\newcommand{\RvPrimaryReps}{120,649}
\newcommand{\RvPrimaryRfAtlas}{23}
\newcommand{\RvPrimaryRfMinusPo}{19}
\newcommand{\RvPrimaryRfMinusPoHi}{20}
\newcommand{\RvPrimaryRfMinusPoLo}{17}
\newcommand{\RvPrimaryRfMinusReg}{0.1}
\newcommand{\RvPrimaryRfMinusRegHi}{1.6}
\newcommand{\RvPrimaryRfMinusRegLo}{\ensuremath{-}1.3}
\newcommand{\RvPrimaryRfPo}{3.5}
\newcommand{\RvPrimaryRfReg}{22}
\newcommand{\RvPrimaryTwoAtlas}{47}
\newcommand{\RvPrimaryTwoMinusPo}{19}
\newcommand{\RvPrimaryTwoMinusPoHi}{20}
\newcommand{\RvPrimaryTwoMinusPoLo}{18}
\newcommand{\RvPrimaryTwoMinusReg}{\ensuremath{-}2.8}
\newcommand{\RvPrimaryTwoMinusRegHi}{\ensuremath{-}1.9}
\newcommand{\RvPrimaryTwoMinusRegLo}{\ensuremath{-}3.6}
\newcommand{\RvPrimaryTwoPo}{28}
\newcommand{\RvPrimaryTwoReg}{50}
\newcommand{\RvPrimaryUnits}{170}
\newcommand{\RvRegCells}{152}
\newcommand{\RvRegCellsHi}{231}
\newcommand{\RvRegCellsLo}{123}
\newcommand{\RvRegEight}{45}
\newcommand{\RvRegEightHi}{48}
\newcommand{\RvRegEightLo}{42}
\newcommand{\RvRegFifty}{14}
\newcommand{\RvRegFiftyHi}{16}
\newcommand{\RvRegFiftyLo}{11}
\newcommand{\RvRegFourHundred}{38}
\newcommand{\RvRegFourHundredHi}{41}
\newcommand{\RvRegFourHundredLo}{35}
\newcommand{\RvRegHundred}{22}
\newcommand{\RvRegHundredHi}{24}
\newcommand{\RvRegHundredLo}{20}
\newcommand{\RvRegOneFifty}{28}
\newcommand{\RvRegOneFiftyHi}{32}
\newcommand{\RvRegOneFiftyLo}{24}
\newcommand{\RvRegSave}{at least 5.3}
\newcommand{\RvRegSaveHi}{6.5}
\newcommand{\RvRegSaveLo}{3.5}
\newcommand{\RvRegTwentyFive}{5.6}
\newcommand{\RvRegTwentyFiveHi}{6.5}
\newcommand{\RvRegTwentyFiveLo}{4.7}
\newcommand{\RvRegTwoHundred}{30}
\newcommand{\RvRegTwoHundredHi}{35}
\newcommand{\RvRegTwoHundredLo}{26}
\newcommand{\RvRelBetterPo}{7}
\newcommand{\RvRelBetterReg}{7}
\newcommand{\RvRelDonorWinPo}{81}
\newcommand{\RvRelDonorWinReg}{77}
\newcommand{\RvRelEightAtlasPred}{0.11}
\newcommand{\RvRelEightAtlasRf}{55}
\newcommand{\RvRelEightAtlasSign}{100}
\newcommand{\RvRelEightDonorsPo}{18 of 18}
\newcommand{\RvRelEightDonorsReg}{18 of 18}
\newcommand{\RvRelEightPoPred}{0.06}
\newcommand{\RvRelEightPoRf}{29}
\newcommand{\RvRelEightPoSign}{96}
\newcommand{\RvRelEightRegPred}{0.08}
\newcommand{\RvRelEightRegRf}{38}
\newcommand{\RvRelEightRegSign}{100}
\newcommand{\RvRelEightTruth}{0.33}
\newcommand{\RvRelEightTruthShare}{100}
\newcommand{\RvRelEightUnits}{34}
\newcommand{\RvRelFiveAtlasPred}{\ensuremath{-}0.28}
\newcommand{\RvRelFiveAtlasRf}{69}
\newcommand{\RvRelFiveAtlasSign}{99}
\newcommand{\RvRelFiveDonorsPo}{16 of 18}
\newcommand{\RvRelFiveDonorsReg}{10 of 18}
\newcommand{\RvRelFivePoPred}{\ensuremath{-}0.22}
\newcommand{\RvRelFivePoRf}{56}
\newcommand{\RvRelFivePoSign}{99}
\newcommand{\RvRelFiveRegPred}{\ensuremath{-}0.32}
\newcommand{\RvRelFiveRegRf}{66}
\newcommand{\RvRelFiveRegSign}{100}
\newcommand{\RvRelFiveTruth}{\ensuremath{-}0.49}
\newcommand{\RvRelFiveTruthShare}{100}
\newcommand{\RvRelFiveUnits}{34}
\newcommand{\RvRelFourAtlasPred}{0.28}
\newcommand{\RvRelFourAtlasRf}{65}
\newcommand{\RvRelFourAtlasSign}{99}
\newcommand{\RvRelFourDonorsPo}{12 of 18}
\newcommand{\RvRelFourDonorsReg}{12 of 18}
\newcommand{\RvRelFourPoPred}{0.28}
\newcommand{\RvRelFourPoRf}{57}
\newcommand{\RvRelFourPoSign}{94}
\newcommand{\RvRelFourRegPred}{0.25}
\newcommand{\RvRelFourRegRf}{63}
\newcommand{\RvRelFourRegSign}{93}
\newcommand{\RvRelFourTruth}{0.50}
\newcommand{\RvRelFourTruthShare}{100}
\newcommand{\RvRelFourUnits}{34}
\newcommand{\RvRelN}{8}
\newcommand{\RvRelOneAtlasPred}{0.02}
\newcommand{\RvRelOneAtlasRf}{27}
\newcommand{\RvRelOneAtlasSign}{72}
\newcommand{\RvRelOneDonorsPo}{15 of 18}
\newcommand{\RvRelOneDonorsReg}{14 of 18}
\newcommand{\RvRelOnePoPred}{0.01}
\newcommand{\RvRelOnePoRf}{15}
\newcommand{\RvRelOnePoSign}{76}
\newcommand{\RvRelOneRegPred}{0.01}
\newcommand{\RvRelOneRegRf}{21}
\newcommand{\RvRelOneRegSign}{79}
\newcommand{\RvRelOneTruth}{0.07}
\newcommand{\RvRelOneTruthShare}{100}
\newcommand{\RvRelOneUnits}{34}
\newcommand{\RvRelSevenAtlasPred}{0.01}
\newcommand{\RvRelSevenAtlasRf}{\ensuremath{-}81}
\newcommand{\RvRelSevenAtlasSign}{36}
\newcommand{\RvRelSevenDonorsPo}{3 of 18}
\newcommand{\RvRelSevenDonorsReg}{4 of 18}
\newcommand{\RvRelSevenPoPred}{0.00}
\newcommand{\RvRelSevenPoRf}{0.0}
\newcommand{\RvRelSevenPoSign}{0}
\newcommand{\RvRelSevenRegPred}{\ensuremath{-}0.03}
\newcommand{\RvRelSevenRegRf}{\ensuremath{-}20}
\newcommand{\RvRelSevenRegSign}{46}
\newcommand{\RvRelSevenTruth}{\ensuremath{-}0.05}
\newcommand{\RvRelSevenTruthShare}{68}
\newcommand{\RvRelSevenUnits}{34}
\newcommand{\RvRelSixAtlasPred}{0.13}
\newcommand{\RvRelSixAtlasRf}{53}
\newcommand{\RvRelSixAtlasSign}{98}
\newcommand{\RvRelSixDonorsPo}{18 of 18}
\newcommand{\RvRelSixDonorsReg}{18 of 18}
\newcommand{\RvRelSixPoPred}{0.04}
\newcommand{\RvRelSixPoRf}{18}
\newcommand{\RvRelSixPoSign}{81}
\newcommand{\RvRelSixRegPred}{0.09}
\newcommand{\RvRelSixRegRf}{37}
\newcommand{\RvRelSixRegSign}{95}
\newcommand{\RvRelSixTruth}{0.38}
\newcommand{\RvRelSixTruthShare}{100}
\newcommand{\RvRelSixUnits}{34}
\newcommand{\RvRelThreeAtlasPred}{0.10}
\newcommand{\RvRelThreeAtlasRf}{81}
\newcommand{\RvRelThreeAtlasSign}{100}
\newcommand{\RvRelThreeDonorsPo}{17 of 18}
\newcommand{\RvRelThreeDonorsReg}{17 of 18}
\newcommand{\RvRelThreePoPred}{0.04}
\newcommand{\RvRelThreePoRf}{39}
\newcommand{\RvRelThreePoSign}{100}
\newcommand{\RvRelThreeRegPred}{0.05}
\newcommand{\RvRelThreeRegRf}{54}
\newcommand{\RvRelThreeRegSign}{100}
\newcommand{\RvRelThreeTruth}{0.15}
\newcommand{\RvRelThreeTruthShare}{100}
\newcommand{\RvRelThreeUnits}{34}
\newcommand{\RvRelTwoAtlasPred}{0.31}
\newcommand{\RvRelTwoAtlasRf}{80}
\newcommand{\RvRelTwoAtlasSign}{100}
\newcommand{\RvRelTwoDonorsPo}{18 of 18}
\newcommand{\RvRelTwoDonorsReg}{18 of 18}
\newcommand{\RvRelTwoPoPred}{0.05}
\newcommand{\RvRelTwoPoRf}{19}
\newcommand{\RvRelTwoPoSign}{93}
\newcommand{\RvRelTwoRegPred}{0.13}
\newcommand{\RvRelTwoRegRf}{46}
\newcommand{\RvRelTwoRegSign}{100}
\newcommand{\RvRelTwoTruth}{0.47}
\newcommand{\RvRelTwoTruthShare}{100}
\newcommand{\RvRelTwoUnits}{34}
\newcommand{\RvRelUnits}{272}
\newcommand{\RvReplFalseShare}{0.21}
\newcommand{\RvReplNominal}{5.0}
\newcommand{\RvReplNull}{6.8}
\newcommand{\RvReplObserved}{3.2}
\newcommand{\RvTechFineOneAtlas}{73}
\newcommand{\RvTechFineOneMinusPo}{3.2}
\newcommand{\RvTechFineOneMinusPoHi}{3.6}
\newcommand{\RvTechFineOneMinusPoLo}{2.6}
\newcommand{\RvTechFineOneMinusReg}{\ensuremath{-}0.6}
\newcommand{\RvTechFineOneMinusRegHi}{\ensuremath{-}0.0}
\newcommand{\RvTechFineOneMinusRegLo}{\ensuremath{-}1.2}
\newcommand{\RvTechFineOnePo}{70}
\newcommand{\RvTechFineOneReg}{74}
\newcommand{\RvTechFineReps}{26,750}
\newcommand{\RvTechFineRfAtlas}{\ensuremath{-}97}
\newcommand{\RvTechFineRfMinusPo}{\ensuremath{-}89}
\newcommand{\RvTechFineRfMinusPoHi}{\ensuremath{-}73}
\newcommand{\RvTechFineRfMinusPoLo}{\ensuremath{-}107}
\newcommand{\RvTechFineRfMinusReg}{\ensuremath{-}87}
\newcommand{\RvTechFineRfMinusRegHi}{\ensuremath{-}71}
\newcommand{\RvTechFineRfMinusRegLo}{\ensuremath{-}104}
\newcommand{\RvTechFineRfPo}{\ensuremath{-}8.2}
\newcommand{\RvTechFineRfReg}{\ensuremath{-}10}
\newcommand{\RvTechFineTwoAtlas}{15}
\newcommand{\RvTechFineTwoMinusPo}{5.6}
\newcommand{\RvTechFineTwoMinusPoHi}{6.1}
\newcommand{\RvTechFineTwoMinusPoLo}{5.1}
\newcommand{\RvTechFineTwoMinusReg}{\ensuremath{-}1.0}
\newcommand{\RvTechFineTwoMinusRegHi}{\ensuremath{-}0.7}
\newcommand{\RvTechFineTwoMinusRegLo}{\ensuremath{-}1.2}
\newcommand{\RvTechFineTwoPo}{9.6}
\newcommand{\RvTechFineTwoReg}{16}
\newcommand{\RvTechFineUnits}{170}
\newcommand{\RvTechOneAtlas}{86}
\newcommand{\RvTechOneMinusPo}{10}
\newcommand{\RvTechOneMinusPoHi}{11}
\newcommand{\RvTechOneMinusPoLo}{9.6}
\newcommand{\RvTechOneMinusReg}{0.1}
\newcommand{\RvTechOneMinusRegHi}{0.5}
\newcommand{\RvTechOneMinusRegLo}{\ensuremath{-}0.3}
\newcommand{\RvTechOnePo}{75}
\newcommand{\RvTechOneReg}{86}
\newcommand{\RvTechReps}{98,739}
\newcommand{\RvTechRfAtlas}{17}
\newcommand{\RvTechRfMinusPo}{13}
\newcommand{\RvTechRfMinusPoHi}{16}
\newcommand{\RvTechRfMinusPoLo}{11}
\newcommand{\RvTechRfMinusReg}{\ensuremath{-}3.9}
\newcommand{\RvTechRfMinusRegHi}{\ensuremath{-}2.0}
\newcommand{\RvTechRfMinusRegLo}{\ensuremath{-}5.8}
\newcommand{\RvTechRfPo}{3.2}
\newcommand{\RvTechRfReg}{20}
\newcommand{\RvTechTwoAtlas}{44}
\newcommand{\RvTechTwoMinusPo}{18}
\newcommand{\RvTechTwoMinusPoHi}{19}
\newcommand{\RvTechTwoMinusPoLo}{17}
\newcommand{\RvTechTwoMinusReg}{\ensuremath{-}1.7}
\newcommand{\RvTechTwoMinusRegHi}{\ensuremath{-}1.1}
\newcommand{\RvTechTwoMinusRegLo}{\ensuremath{-}2.3}
\newcommand{\RvTechTwoPo}{26}
\newcommand{\RvTechTwoReg}{46}
\newcommand{\RvTechUnits}{170}
\newcommand{\RvTruthMinusPoMax}{18}
\newcommand{\RvTruthMinusPoMin}{2.8}
\newcommand{\RvValMeanDiff}{$10^{-7}$}
\newcommand{\RvValMsqDiff}{machine precision}
\newcommand{\SyExt}{4.9}
\newcommand{\SyGeo}{3.1}
\newcommand{\SyGeoHi}{3.2}
\newcommand{\SyGeoLo}{2.4}
\newcommand{\SyMax}{6.4}
\newcommand{\SyMin}{1.4}
\newcommand{\SyNInf}{nine}
\newcommand{\SyNLb}{six}
\newcommand{\SyNTest}{ten}
\newcommand{\VlAtlasOverOther}{2.10}
\newcommand{\VlAtlasOverOtherHi}{2.26}
\newcommand{\VlAtlasOverOtherLo}{1.97}
\newcommand{\VlAtlasOverOwn}{1.34}
\newcommand{\VlAtlasOverOwnHi}{1.56}
\newcommand{\VlAtlasOverOwnLo}{1.15}
\newcommand{\VlCalConMax}{1.00}
\newcommand{\VlCalConMin}{1.00}
\newcommand{\VlCalDfMax}{0.95}
\newcommand{\VlCalDfMin}{0.67}
\newcommand{\VlCellsAtlas}{147}
\newcommand{\VlCellsDf}{178}
\newcommand{\VlCellsOther}{70}
\newcommand{\VlCellsOwn}{110}
\newcommand{\VlCellsOwnPool}{114}
\newcommand{\VlCellsPaired}{800}
\newcommand{\VlFrozenAt}{21:54 EDT on 28 September 2026}
\newcommand{\VlHashedAt}{22:16 EDT}
\newcommand{\VlLevel}{57}
\newcommand{\VlNAntibodies}{111}
\newcommand{\VlNDonors}{18}
\newcommand{\VlNFolds}{34}
\newcommand{\VlNPools}{seven}
\newcommand{\VlPoolMax}{8.2}
\newcommand{\VlPoolMin}{4.5}
\newcommand{\VlPopsFifty}{10}
\newcommand{\VlPopsHundred}{10}
\newcommand{\VlPopsTwentyFive}{9}
\newcommand{\VlRfAtlasEightHundred}{48}
\newcommand{\VlRfAtlasEightHundredHi}{52}
\newcommand{\VlRfAtlasEightHundredLo}{44}
\newcommand{\VlRfAtlasFifty}{11}
\newcommand{\VlRfAtlasFiftyHi}{14}
\newcommand{\VlRfAtlasFiftyLo}{8.8}
\newcommand{\VlRfAtlasHundred}{23}
\newcommand{\VlRfAtlasHundredHi}{25}
\newcommand{\VlRfAtlasHundredLo}{20}
\newcommand{\VlRfAtlasTwentyFive}{0.5}
\newcommand{\VlRfAtlasTwentyFiveHi}{2.3}
\newcommand{\VlRfAtlasTwentyFiveLo}{\ensuremath{-}2.0}
\newcommand{\VlRfDfEightHundred}{48}
\newcommand{\VlRfDfEightHundredHi}{52}
\newcommand{\VlRfDfEightHundredLo}{44}
\newcommand{\VlRfDfFifty}{\ensuremath{-}60}
\newcommand{\VlRfDfFiftyHi}{\ensuremath{-}52}
\newcommand{\VlRfDfFiftyLo}{\ensuremath{-}69}
\newcommand{\VlRfDfHundred}{13}
\newcommand{\VlRfDfHundredHi}{16}
\newcommand{\VlRfDfHundredLo}{9.3}
\newcommand{\VlRfDfTwentyFive}{\ensuremath{-}408}
\newcommand{\VlRfDfTwentyFiveHi}{\ensuremath{-}376}
\newcommand{\VlRfDfTwentyFiveLo}{\ensuremath{-}447}
\newcommand{\VlRfOtherEightHundred}{57}
\newcommand{\VlRfOtherEightHundredHi}{60}
\newcommand{\VlRfOtherEightHundredLo}{53}
\newcommand{\VlRfOtherFifty}{23}
\newcommand{\VlRfOtherFiftyHi}{26}
\newcommand{\VlRfOtherFiftyLo}{20}
\newcommand{\VlRfOtherHundred}{34}
\newcommand{\VlRfOtherHundredHi}{37}
\newcommand{\VlRfOtherHundredLo}{31}
\newcommand{\VlRfOtherTwentyFive}{8.8}
\newcommand{\VlRfOtherTwentyFiveHi}{12}
\newcommand{\VlRfOtherTwentyFiveLo}{5.6}
\newcommand{\VlRfOwnEightHundred}{45}
\newcommand{\VlRfOwnEightHundredHi}{49}
\newcommand{\VlRfOwnEightHundredLo}{41}
\newcommand{\VlRfOwnFifty}{19}
\newcommand{\VlRfOwnFiftyHi}{23}
\newcommand{\VlRfOwnFiftyLo}{15}
\newcommand{\VlRfOwnHundred}{27}
\newcommand{\VlRfOwnHundredHi}{31}
\newcommand{\VlRfOwnHundredLo}{24}
\newcommand{\VlRfOwnPoolEightHundred}{44}
\newcommand{\VlRfOwnPoolEightHundredHi}{48}
\newcommand{\VlRfOwnPoolEightHundredLo}{40}
\newcommand{\VlRfOwnPoolFifty}{19}
\newcommand{\VlRfOwnPoolFiftyHi}{23}
\newcommand{\VlRfOwnPoolFiftyLo}{14}
\newcommand{\VlRfOwnPoolHundred}{27}
\newcommand{\VlRfOwnPoolHundredHi}{31}
\newcommand{\VlRfOwnPoolHundredLo}{23}
\newcommand{\VlRfOwnPoolTwentyFive}{8.2}
\newcommand{\VlRfOwnPoolTwentyFiveHi}{13}
\newcommand{\VlRfOwnPoolTwentyFiveLo}{2.4}
\newcommand{\VlRfOwnTwentyFive}{5.6}
\newcommand{\VlRfOwnTwentyFiveHi}{9.1}
\newcommand{\VlRfOwnTwentyFiveLo}{1.5}
\newcommand{\VlRfPairedEightHundred}{13}
\newcommand{\VlRfPairedEightHundredHi}{20}
\newcommand{\VlRfPairedEightHundredLo}{11}
\newcommand{\VlRfPairedFifty}{2.4}
\newcommand{\VlRfPairedFiftyHi}{3.8}
\newcommand{\VlRfPairedFiftyLo}{1.1}
\newcommand{\VlRfPairedHundred}{3.5}
\newcommand{\VlRfPairedHundredHi}{6.2}
\newcommand{\VlRfPairedHundredLo}{2.9}
\newcommand{\VlRfPairedTwentyFive}{0.5}
\newcommand{\VlRfPairedTwentyFiveHi}{0.9}
\newcommand{\VlRfPairedTwentyFiveLo}{0.4}
\newcommand{\VlRowsAtlas}{137}
\newcommand{\VlRowsFifty}{40}
\newcommand{\VlRowsHundred}{90}
\newcommand{\VlRowsMax}{790}
\newcommand{\VlRowsTwentyFive}{16}
\newcommand{\VlSaveAtlas}{5.4}
\newcommand{\VlSaveAtlasHi}{6.0}
\newcommand{\VlSaveAtlasLo}{4.8}
\newcommand{\VlSaveAtlasQuarter}{13.4}
\newcommand{\VlSaveAtlasQuarterHi}{13.7}
\newcommand{\VlSaveAtlasQuarterLo}{4.5}
\newcommand{\VlSaveAtlasRows}{5.8}
\newcommand{\VlSaveAtlasThreeQ}{1.9}
\newcommand{\VlSaveAtlasThreeQHi}{2.1}
\newcommand{\VlSaveAtlasThreeQLo}{1.6}
\newcommand{\VlSaveDf}{4.5}
\newcommand{\VlSaveDfHi}{4.8}
\newcommand{\VlSaveDfLo}{4.1}
\newcommand{\VlSaveOther}{11.4}
\newcommand{\VlSaveOtherHi}{12.8}
\newcommand{\VlSaveOtherLo}{10.3}
\newcommand{\VlSaveOwn}{7.3}
\newcommand{\VlSaveOwnHi}{8.8}
\newcommand{\VlSaveOwnLo}{5.8}
\newcommand{\VlSaveOwnPool}{7.0}
\newcommand{\VlSaveOwnPoolHi}{8.9}
\newcommand{\VlSaveOwnPoolLo}{5.4}
\newcommand{\VlSinglesFifty}{0.9}
\newcommand{\VlSinglesHundred}{0.2}
\newcommand{\VlSinglesTwentyFive}{2.2}
\newcommand{\VlTarget}{28}
\newcommand{\VlVFour}{not met}
\newcommand{\VlVOne}{met}
\newcommand{\VlVThree}{met}
\newcommand{\VlVTwo}{met}

\title{Supplementary Information\\[4pt]\large Estimating Cross-Modal Dependence with Fewer Paired Cells}
\author{}
\date{}
\hypersetup{pdftitle={Supplementary Information: Estimating Cross-Modal Dependence with Fewer Paired Cells},pdfauthor={}}
\begin{document}
\maketitle
This Supplementary Information follows the order of the Results. Supplementary Table~\ref{tab:registry} lists every prespecified hypothesis of the manuscript with its criterion, estimate and result, and every analysis done after scoring. Supplementary Note~\ref{sec:ident} states what unpaired data identify (Propositions~\ref{prop:pairing}--\ref{prop:gaussian}) and Supplementary Note~\ref{sec:noise} how the latent RNA correlation is estimated (Proposition~\ref{prop:noise}). Supplementary Note~\ref{sec:deploy} gives the bone-marrow test, in which the unpaired cells came from other sites and assays, its analyses after scoring and the analyses with unpaired neurons of other animals; Supplementary Note~\ref{sec:val} gives the validation test of the revised workflow in a blood study not analysed before, with the allocation of its donors, its biological readout and the readout's calibration, the comparison with other methods given the same inputs and the robustness of its result. Supplementary Note~\ref{sec:semi} defines the estimator and its guarantees (Propositions~\ref{prop:bjs}--\ref{prop:minimax}) and reports the development comparisons and the external test; Supplementary Note~\ref{sec:gen} gives the benchmark and the saving in every test data set; Supplementary Note~\ref{sec:pred} gives the theory of the saving (Propositions~\ref{prop:oracle}--\ref{prop:reveal}) with its development and prospective test. Supplementary Notes~\ref{sec:cross}--\ref{sec:rep} give the analyses of learning without pairs: the blood test with its identification analyses, the perturbation test and its replication in a second screen. Supplementary Note~\ref{sec:review} gives the follow-up checks and analyses, registered before any of their quantities was computed: the scoring of repeated draws, the noise estimate after contrasts (Proposition~\ref{prop:contrast}), the reference's axes and condition maps with reweighting, the validation's saving for one draw and one atlas, named relationships, the association test and finer cell states, the law with data sets as units, a pilot that sets the budget, and savings at fixed accuracy. Supplementary Figs.~\ref{sfig:ident} and~\ref{sfig:transfer} accompany the Results of the main text, and Supplementary Fig.~\ref{sfig:entries} accompanies Supplementary Note~\ref{sec:cross}.

The notes use the terms of the main text together with a few technical equivalents. The dependence between two modalities within a population is their cross-correlation matrix (cross-covariance where stated), and predicting no dependence is called independence. The axes are the eigenvectors of the unpaired correlation matrices, also called the unpaired bases, and a block is a set of coordinates in these bases. The estimator, also called the proposed estimator, is block James--Stein shrinkage along the unpaired axes followed by the condition maps; single-block (one-block) James--Stein shrinks all coordinates by one factor. For a block $b$, $r_b^2$ is the sum of squares of its true dependence (its signal) and $\tau_b$ the total variance of one cell's products in it; the dependence and noise shares are $w_b=r_b^2/\sum_{b'}r_{b'}^2$ and $\pi_b=\tau_b/\sum_{b'}\tau_{b'}$. Contrasts are the Helmert contrasts of a population's drawn paired cells (Methods). An arm is one version of an analysis within a test. A test was registered by recording a cryptographic hash of its plan, code and data, so that any later change would be detectable; a prespecified test was registered before any held-out result was computed. Learning without pairs is also called pairing-free learning.
\section*{Registry of hypotheses and analyses after scoring}
Every estimate and result below is read directly from the tests' result records by the code that writes this table.
\newcommand{\RegistryCaption}{\textbf{Every prespecified hypothesis of the manuscript and every analysis done after scoring.} For each prespecified test, each hypothesis with its criterion, the estimate with its 95\% interval and the result. Savings are ratios of paired cells needed (comparator over the estimator); ``at least'' marks a comparator that did not reach the target within the largest budget. The readout of the validation test, and later the comparison with other methods given the same atlas (K1, K2), were specified after the validation test had been scored and before their own quantities were computed. Below, analyses done after scoring, amendments and deviations, with where they are reported.}
{\footnotesize\setlength{\LTleft}{0pt}\setlength{\LTright}{0pt}\setlength{\LTcapwidth}{\textwidth}
% Generated by extension/synthesis/registry.py. Do not edit.
\begin{longtable}{@{}>{\raggedright\arraybackslash}p{0.12\textwidth}>{\raggedright\arraybackslash}p{0.32\textwidth}>{\raggedright\arraybackslash}p{0.16\textwidth}>{\raggedright\arraybackslash}p{0.22\textwidth}>{\raggedright\arraybackslash}p{0.07\textwidth}@{}}
\caption{\RegistryCaption}\label{tab:registry}\\
\toprule
Hypothesis & Statement & Criterion & Estimate (95\% CI) & Result \\
\midrule
\endfirsthead
\multicolumn{5}{@{}l}{\textit{\tablename~\thetable, continued}} \\
\toprule
Hypothesis & Statement & Criterion & Estimate (95\% CI) & Result \\
\midrule
\endhead
\bottomrule
\endfoot
\multicolumn{5}{@{}l}{\textit{External test (OverCITE-seq ORF screen)}} \\
H1 (primary) & Fewer paired cells than paired-only estimation at 30\% recovery & 95\% lower bound of the ratio $>$ 1 & 4.9 (3.2 to 6.0) & met \\
H2 & Fewer paired cells than SemiCCA & 95\% lower bound $>$ 1 & 1.7 (1.3 to 2.0) & met \\
H3 & No more than 1.25 times the paired cells of Champollion & 95\% lower bound $>$ 0.8 & 1.8 (1.2 to 2.2) & met \\
\multicolumn{5}{@{}l}{\textit{Benchmark (nine data sets)}} \\
G1 (primary) & Geometric-mean saving over paired-only estimation at the plan's target (half of the unshrunk estimator's recovery with every training cell paired) & 95\% lower bound $>$ 1 & 1.50 (1.40 to 1.55) & met \\
G2 & The same against SemiCCA & 95\% lower bound $>$ 1 & 1.34 (1.05 to 1.39) & met \\
G3 & Self-tuning against the fixed estimator (secondary) & two-sided; no threshold & 0.77 (0.75 to 0.82) & reported \\
G4 & Predicted and observed savings rank the data sets alike (structural law) & Spearman $>$ 0, one-sided P $<$ 0.05 & $\rho$ \ensuremath{-}0.09, P = 0.63 & not met \\
G1 (amended) & As G1, at half of the best recovery of any method, over the 8 data sets with recovered dependence (amended after three data sets had been scored) & 95\% lower bound $>$ 1 & 2.90 (2.26 to 3.03) & met \\
G2 (amended) & As G2 under the amended rule & 95\% lower bound $>$ 1 & 2.29 (1.82 to 2.38) & met \\
G4 (amended) & As G4 under the amended rule & Spearman $>$ 0, one-sided P $<$ 0.05 & $\rho$ 0.42, P = 0.10 & not met \\
\multicolumn{5}{@{}l}{\textit{Prospective test of the law (ten data sets, 29 problems)}} \\
H1 (primary) & The law is calibrated (10 of 29 problems with an estimable saving) & median error $<$ 25\% & 21\% (11\% to 29\%) & met \\
H2 & The law ranks the problems & Spearman $>$ 0, one-sided P $<$ 0.05 & $\rho$ 0.98, P $<$ 0.001 & met \\
H3 & No saving in random bases & median saving in [0.8, 1.25] & 1.00 & met \\
H4 & A pilot of up to 200 paired cells predicts the saving & median error $<$ 50\% & 12\% (4\% to 147\%) & met \\
H5 & Unpaired noise with the development dependence profile predicts it & median error $<$ 50\% & 40\% (31\% to 45\%) & met \\
H6 & Every saving lies within the unpaired bounds & all in [0.8, 1.25 / min $\pi$] & 0 outside & met \\
\multicolumn{5}{@{}l}{\textit{Bone-marrow test (four sites)}} \\
D1 (primary) & Saving over paired-only estimation with unpaired cells of other sites & 95\% lower bound $>$ 1 & at least 3.91 (3.29 to 4.45) & met \\
D2 & The same with nuclear RNA of other sites & 95\% lower bound $>$ 1 & neither curve reached the target & not met \\
D3 & Other sites need no more paired cells than the own site & 95\% upper bound $<$ 1.25 & 0.87 (0.81 to 0.96) & met \\
\multicolumn{5}{@{}l}{\textit{Validation test (blood study not analysed before; atlas of another study)}} \\
V1 (primary) & Saving over paired-only estimation with the atlas & 95\% lower bound $>$ 1 & at least 5.43 (4.82 to 5.96) & met \\
V2 & The atlas arm recovers dependence with 50 and with 100 paired cells & 95\% lower bounds $>$ 0 & 11\% (8.8\% to 14\%); 23\% (20\% to 25\%) & met \\
V3 & Saving with the study's other pools and donors & 95\% lower bound $>$ 1 & at least 11.4 (10.3 to 12.8) & met \\
V4 & The atlas needs no more paired cells than the paired samples' own unpaired cells & 95\% upper bound $<$ 1.25 & 1.34 (1.15 to 1.56) & not met \\
\multicolumn{5}{@{}l}{\textit{Readout of the validation test (specified after V1-V4 were scored)}} \\
B1 & Direction of reproducible within-cell-type associations: atlas minus paired-only & 95\% lower bound $>$ 0 at 50, 100, 150 paired cells & 50: 12.9 points (12.1 to 13.7); 100: 11.2 points (10.3 to 12.0); 150: 9.1 points (8.4 to 9.8) & met \\
B2 & Top-100 nominations that are reproducible associations: atlas minus paired-only & 95\% lower bound $>$ 0 at 50, 100, 150 paired cells & 50: 20.1 points (18.9 to 21.3); 100: 19.3 points (18.0 to 20.5); 150: 16.8 points (15.6 to 18.0) & met \\
B3 & Recovered fraction of cognate RNA-protein correlations: atlas minus paired-only & 95\% lower bound $>$ 0 at 100, 150 paired cells & 100: 14.2 points (12.5 to 15.8); 150: 14.4 points (12.9 to 15.3) & met \\
\multicolumn{5}{@{}l}{\textit{Other methods given the validation's atlas (specified after V1-V4 and B1-B3 were scored)}} \\
K1 & SemiCCA (best of three forms), given the estimator's paired cells, atlas and tuning information, needs more paired cells than the estimator & 95\% lower bound of the ratio of paired cells $>$ 1 & 3.36 (2.41 to 4.63) & met \\
K2 & Reference regression (ridge; best of six forms) likewise & 95\% lower bound of the ratio of paired cells $>$ 1 & 1.15 (1.00 to 1.36) & met \\
\multicolumn{5}{@{}l}{\textit{Learning without pairs, blood: confirmatory cohort}} \\
H1 (primary) & Beats a donor-aware permutation null & P $<$ 0.05 & error 0.29 against 0.95; P = 0.005 & met \\
H2 (primary) & Recovers most of the paired operator's improvement over independence & 95\% lower bound of the share $>$ 0.5 & 82\% (80\% to 84\%) & met \\
H3 & Recovers dependence within cell types & 95\% upper bound of the error $<$ 1 & 1.08 (1.01 to 1.16) & not met \\
H4 & The paired closed form beats four matched baselines & Bonferroni intervals exclude zero & 4 of 4 baselines & met \\
\multicolumn{5}{@{}l}{\textit{Learning without pairs, colon cohort (analysed before)}} \\
H1 (primary) & Beats a donor-aware permutation null & P $<$ 0.05 & error 0.99 against 1.17; P = 0.005 & met \\
H2 (primary) & Recovers most of the paired operator's improvement over independence & 95\% lower bound of the share $>$ 0.5 & \ensuremath{-}12\% (\ensuremath{-}49\% to 15\%) & not met \\
H3 & Recovers dependence within cell types & 95\% upper bound of the error $<$ 1 & 1.27 (1.24 to 1.31) & not met \\
H4 & The paired closed form beats four matched baselines & Bonferroni intervals exclude zero & 4 of 4 baselines & met \\
\multicolumn{5}{@{}l}{\textit{Learning without pairs, melanoma screen}} \\
H1 & Recovery of within-condition dependence in held-out targets & 95\% lower bound $>$ 0 & 6.8\% (4.1\% to 9.4\%) & met \\
H2a & Above the pseudo-population control & difference, lower bound $>$ 0 & 6.9\% (4.2\% to 9.5\%) & met \\
H2b & Above the derangement control & difference, lower bound $>$ 0 & 7.2\% (4.6\% to 9.7\%) & met \\
H3 & Recovery in non-targeting cells & 95\% lower bound $>$ 0 & 19\% (16\% to 21\%) & met \\
H4 & Most of what paired cells recover & share, lower bound $>$ 50\% & 18\% (12\% to 25\%) & not met \\
\multicolumn{5}{@{}l}{\textit{Learning without pairs, monocyte screen (replication)}} \\
R1 & Recovery over all test groups & 95\% lower bound $>$ 0 & \ensuremath{-}212\% (\ensuremath{-}287\% to \ensuremath{-}157\%) & not met \\
R2a & Above the pseudo-population control & lower bound $>$ 0 & \ensuremath{-}202\% (\ensuremath{-}275\% to \ensuremath{-}151\%) & not met \\
R2b & Above the derangement control & lower bound $>$ 0 & \ensuremath{-}167\% (\ensuremath{-}234\% to \ensuremath{-}119\%) & not met \\
R3 & Interferon block: at least half of the paired closed form & lower bound $>$ 50\% & \ensuremath{-}3.5\% (\ensuremath{-}41.3\% to 24.3\%) & not met \\
R4 & Recovery in non-targeting cells & lower bound $>$ 0 & \ensuremath{-}90\% (\ensuremath{-}140\% to \ensuremath{-}52\%) & not met \\
R5 & Removing the most exposed targets costs more than random ones & lower bound $>$ 0 & \ensuremath{-}44\% (\ensuremath{-}91\% to \ensuremath{-}3\%) & not met \\
\midrule
\multicolumn{5}{@{}l}{\textit{Analyses after scoring, amendments and deviations}} \\
\multicolumn{3}{@{}>{\raggedright\arraybackslash}p{0.63\textwidth}}{Development comparisons and ablations (melanoma and monocyte screens)} & \multicolumn{2}{>{\raggedright\arraybackslash}p{0.3\textwidth}@{}}{Supplementary Note on the estimator} \\
\multicolumn{3}{@{}>{\raggedright\arraybackslash}p{0.63\textwidth}}{Benchmark amendments: primary target half of the best recovery of any method; data sets without recovered dependence set aside (recorded after three and four of nine data sets were scored)} & \multicolumn{2}{>{\raggedright\arraybackslash}p{0.3\textwidth}@{}}{Methods; savings table} \\
\multicolumn{3}{@{}>{\raggedright\arraybackslash}p{0.63\textwidth}}{Summary across test data sets (geometric mean 3.1)} & \multicolumn{2}{>{\raggedright\arraybackslash}p{0.3\textwidth}@{}}{Results; savings table} \\
\multicolumn{3}{@{}>{\raggedright\arraybackslash}p{0.63\textwidth}}{Development of the law on the benchmark data sets} & \multicolumn{2}{>{\raggedright\arraybackslash}p{0.3\textwidth}@{}}{Supplementary Note on the law} \\
\multicolumn{3}{@{}>{\raggedright\arraybackslash}p{0.63\textwidth}}{External test: the registered scoring code scored the one-sided variant of the estimator, not the two-sided variant named in the plan (deviation; both variants met H1-H3)} & \multicolumn{2}{>{\raggedright\arraybackslash}p{0.3\textwidth}@{}}{Methods; Supplementary Note on the estimator} \\
\multicolumn{3}{@{}>{\raggedright\arraybackslash}p{0.63\textwidth}}{Bone-marrow test: centring by contrasts; RNA from nuclei (maps, axes, the law along each pool's axes, principal angles); population sizes; calibration of the noise estimate; paired-only control under both centrings} & \multicolumn{2}{>{\raggedright\arraybackslash}p{0.3\textwidth}@{}}{Results; Supplementary Note on the bone-marrow test} \\
\multicolumn{3}{@{}>{\raggedright\arraybackslash}p{0.63\textwidth}}{Fly connectome: unpaired neurons of other animals with contrasts} & \multicolumn{2}{>{\raggedright\arraybackslash}p{0.3\textwidth}@{}}{Results; Supplementary Note on the bone-marrow test} \\
\multicolumn{3}{@{}>{\raggedright\arraybackslash}p{0.63\textwidth}}{Validation test: budget axis counted in contrasts; results per donor and pool; leaving out each pool; atlas size; skewed or missing cell types in the atlas, and the skewed atlases scored without condition maps} & \multicolumn{2}{>{\raggedright\arraybackslash}p{0.3\textwidth}@{}}{Results; Supplementary Note on the validation test} \\
\multicolumn{3}{@{}>{\raggedright\arraybackslash}p{0.63\textwidth}}{Validation test: sample, donor and pool allocation of every fold, checked from sample labels (every fold donor-disjoint)} & \multicolumn{2}{>{\raggedright\arraybackslash}p{0.3\textwidth}@{}}{Methods; Supplementary Note on the validation test} \\
\multicolumn{3}{@{}>{\raggedright\arraybackslash}p{0.63\textwidth}}{Validation test: Champollion given the same paired cells and atlas, on the first draw of every fold and budget (specified with K1-K2, descriptive)} & \multicolumn{2}{>{\raggedright\arraybackslash}p{0.3\textwidth}@{}}{Results; Supplementary Note on the validation test} \\
\multicolumn{3}{@{}>{\raggedright\arraybackslash}p{0.63\textwidth}}{Readout of the validation test, calibrated: significance thresholds, false-discovery-rate truth sets, majority-sign baselines, weighting of samples and cell types, bootstrap P-values, simultaneous intervals and Holm adjustment (specified with K1-K2)} & \multicolumn{2}{>{\raggedright\arraybackslash}p{0.3\textwidth}@{}}{Results; Supplementary Note on the validation test} \\
\multicolumn{3}{@{}>{\raggedright\arraybackslash}p{0.63\textwidth}}{Blood test of learning without pairs: dependence split into parts between and within cell types; identification analyses} & \multicolumn{2}{>{\raggedright\arraybackslash}p{0.3\textwidth}@{}}{Supplementary Note on learning without pairs in blood} \\
\multicolumn{3}{@{}>{\raggedright\arraybackslash}p{0.63\textwidth}}{Melanoma screen: analyses of the interferon response and of removing perturbations (specified before computation, after the scoring cells were opened)} & \multicolumn{2}{>{\raggedright\arraybackslash}p{0.3\textwidth}@{}}{Results; Supplementary Note on the perturbation screen} \\
\multicolumn{3}{@{}>{\raggedright\arraybackslash}p{0.63\textwidth}}{Monocyte screen: a penalty fixed after the scoring files had been opened} & \multicolumn{2}{>{\raggedright\arraybackslash}p{0.3\textwidth}@{}}{Supplementary Note on the replication} \\
\multicolumn{3}{@{}>{\raggedright\arraybackslash}p{0.63\textwidth}}{Check that every recovery curve averages the scores of single draws (every test rerun with its registered code)} & \multicolumn{2}{>{\raggedright\arraybackslash}p{0.3\textwidth}@{}}{Methods; Supplementary Note on follow-up analyses} \\
\multicolumn{3}{@{}>{\raggedright\arraybackslash}p{0.63\textwidth}}{Follow-up analyses, registered together before any was computed: calibration of the noise estimate after contrasts; the atlas's axes and condition maps at every budget; maps reweighted to the paired cells' mix of cell types; one draw of paired cells and resampled atlas patients; named within-type relationships; calibration of the association test; finer cell states and technical covariates; the law with data sets as units; a pilot that sets the total budget; savings at fixed accuracy} & \multicolumn{2}{>{\raggedright\arraybackslash}p{0.3\textwidth}@{}}{Results; Supplementary Note on follow-up analyses} \\
\multicolumn{3}{@{}>{\raggedright\arraybackslash}p{0.63\textwidth}}{Follow-up analyses: two runs departed from the planned order of execution (deviations; no code, setting or value changed)} & \multicolumn{2}{>{\raggedright\arraybackslash}p{0.3\textwidth}@{}}{Supplementary Note on follow-up analyses} \\
\end{longtable}
}

\clearpage
\section*{Supplementary Figures accompanying the main text}

\begin{figure}[p]
\centering
% Generated by extension/cross_study/make_assets.py. Do not edit.
\definecolor{donorcol}{HTML}{5D6670}
\definecolor{haocol}{HTML}{217D91}
\definecolor{coloncol}{HTML}{A64442}
\definecolor{betweencol}{HTML}{217D91}
\definecolor{withincol}{HTML}{C9803A}
\begin{tikzpicture}[font=\sffamily]
\begin{axis}[name=pa,width=.28\textwidth,height=4.6cm,xmin=0,xmax=1.05,ymin=.4,ymax=4.6,
 xtick={0,.25,.5,.75,1},ytick={1,2,3,4},yticklabels={{Colon, within},{Colon, between},{Blood, within},{Blood, between}},xlabel={Norm relative to total},title={\textbf{a}\enspace Size of each part},
 axis x line*=bottom,axis y line*=left,tick align=outside,y tick style={draw=none},xmajorgrids=true,grid style={gray!18},axis line style={gray!55},tick label style={font=\sffamily\scriptsize},label style={font=\sffamily\small},clip=false,title style={font=\sffamily\small,at={(0,1)},anchor=base west,yshift=5pt}]
\draw[black!25] (axis cs:0,2.5) -- (axis cs:1.05,2.5);
\addplot[only marks,mark=*,mark size=.9pt,color=withincol!70] coordinates { (0.4233,0.850) (0.8419,0.877) (0.4147,0.905) (0.3821,0.932) (0.5345,0.959) (0.3194,0.986) (0.3447,1.014) (0.3670,1.041) (0.4749,1.068) (0.3810,1.095) (0.2721,1.123) (0.4058,1.150) };
\addplot[only marks,mark=|,mark size=4pt,line width=1.2pt,color=black] coordinates { (0.4301,1) };
\addplot[only marks,mark=*,mark size=.9pt,color=betweencol!70] coordinates { (0.8211,1.850) (0.4498,1.877) (0.8296,1.905) (0.8458,1.932) (0.7491,1.959) (0.8763,1.986) (0.8686,2.014) (0.8598,2.041) (0.8419,2.068) (0.8234,2.095) (0.9382,2.123) (0.8584,2.150) };
\addplot[only marks,mark=|,mark size=4pt,line width=1.2pt,color=black] coordinates { (0.8135,2) };
\addplot[only marks,mark=*,mark size=.9pt,color=withincol!70] coordinates { (0.1436,2.850) (0.1240,2.893) (0.1052,2.936) (0.1011,2.979) (0.1192,3.021) (0.0952,3.064) (0.1069,3.107) (0.1119,3.150) };
\addplot[only marks,mark=|,mark size=4pt,line width=1.2pt,color=black] coordinates { (0.1134,3) };
\addplot[only marks,mark=*,mark size=.9pt,color=betweencol!70] coordinates { (0.9596,3.850) (0.9627,3.893) (0.9728,3.936) (0.9790,3.979) (0.9666,4.021) (0.9772,4.064) (0.9747,4.107) (0.9762,4.150) };
\addplot[only marks,mark=|,mark size=4pt,line width=1.2pt,color=black] coordinates { (0.9711,4) };
\end{axis}
\begin{axis}[at={(pa.east)},anchor=west,xshift=1.0cm,name=pb,width=.28\textwidth,height=4.6cm,xmin=0,xmax=2.7,ymin=.4,ymax=4.6,
 xtick={0,.5,1,1.5,2,2.5},ytick={1,2,3,4},yticklabels={},xlabel={Relative error},title={\textbf{b}\enspace Prediction error},
 axis x line*=bottom,axis y line*=left,tick align=outside,y tick style={draw=none},xmajorgrids=true,grid style={gray!18},axis line style={gray!55},tick label style={font=\sffamily\scriptsize},label style={font=\sffamily\small},clip=false,title style={font=\sffamily\small,at={(0,1)},anchor=base west,yshift=5pt}]
\draw[black!25] (axis cs:0,2.5) -- (axis cs:2.7,2.5);
\draw[dashed,black!45] (axis cs:1,.4) -- (axis cs:1,4.6);
\addplot[only marks,mark=*,mark size=.9pt,color=withincol!70] coordinates { (1.4633,0.850) (1.2107,0.877) (1.4279,0.905) (1.5323,0.932) (1.3258,0.959) (1.5077,0.986) (1.3997,1.014) (1.3129,1.041) (1.2660,1.068) (1.4275,1.095) (1.3212,1.123) (1.3428,1.150) };
\addplot[only marks,mark=|,mark size=4pt,line width=1.2pt,color=black] coordinates { (1.3782,1) };
\addplot[only marks,mark=*,mark size=.9pt,color=betweencol!70] coordinates { (1.2597,1.850) (1.2915,1.877) (1.0163,1.905) (1.2113,1.932) (1.1092,1.959) (1.0779,1.986) (1.2769,2.014) (1.2425,2.041) (1.1182,2.068) (1.1504,2.095) (0.8499,2.123) (1.0999,2.150) };
\addplot[only marks,mark=|,mark size=4pt,line width=1.2pt,color=black] coordinates { (1.1420,2) };
\addplot[only marks,mark=*,mark size=.9pt,color=withincol!70] coordinates { (1.6535,2.850) (1.8537,2.893) (2.0413,2.936) (2.5572,2.979) (1.5532,3.021) (1.7068,3.064) (1.7664,3.107) (2.2263,3.150) };
\addplot[only marks,mark=|,mark size=4pt,line width=1.2pt,color=black] coordinates { (1.9198,3) };
\addplot[only marks,mark=*,mark size=.9pt,color=betweencol!70] coordinates { (0.7587,3.850) (0.6902,3.893) (0.6533,3.936) (0.7263,3.979) (0.6587,4.021) (0.7264,4.064) (0.7272,4.107) (0.7414,4.150) };
\addplot[only marks,mark=|,mark size=4pt,line width=1.2pt,color=black] coordinates { (0.7103,4) };
\end{axis}
\begin{axis}[at={(pb.east)},anchor=west,xshift=1.0cm,name=pc,width=.28\textwidth,height=4.6cm,xmin=0,xmax=1.05,ymin=.4,ymax=4.6,
 xtick={0,.25,.5,.75,1},ytick={1,2,3,4},yticklabels={},xlabel={Share of RNA covariance},title={\textbf{c}\enspace Share inside \textit{S}},
 axis x line*=bottom,axis y line*=left,tick align=outside,y tick style={draw=none},xmajorgrids=true,grid style={gray!18},axis line style={gray!55},tick label style={font=\sffamily\scriptsize},label style={font=\sffamily\small},clip=false,title style={font=\sffamily\small,at={(0,1)},anchor=base west,yshift=5pt}]
\draw[black!25] (axis cs:0,2.5) -- (axis cs:1.05,2.5);
\addplot[only marks,mark=*,mark size=.9pt,color=withincol!70] coordinates { (0.0914,0.850) (0.0816,0.877) (0.0859,0.905) (0.0785,0.932) (0.0750,0.959) (0.0943,0.986) (0.0855,1.014) (0.0961,1.041) (0.0800,1.068) (0.0953,1.095) (0.1011,1.123) (0.0814,1.150) };
\addplot[only marks,mark=|,mark size=4pt,line width=1.2pt,color=black] coordinates { (0.0872,1) };
\addplot[only marks,mark=*,mark size=.9pt,color=betweencol!70] coordinates { (0.5861,1.850) (0.5182,1.877) (0.5630,1.905) (0.5029,1.932) (0.5232,1.959) (0.5633,1.986) (0.6049,2.014) (0.5842,2.041) (0.5282,2.068) (0.6475,2.095) (0.6644,2.123) (0.4950,2.150) };
\addplot[only marks,mark=|,mark size=4pt,line width=1.2pt,color=black] coordinates { (0.5651,2) };
\addplot[only marks,mark=*,mark size=.9pt,color=withincol!70] coordinates { (0.1265,2.850) (0.1254,2.893) (0.1047,2.936) (0.1132,2.979) (0.1273,3.021) (0.1046,3.064) (0.1332,3.107) (0.1331,3.150) };
\addplot[only marks,mark=|,mark size=4pt,line width=1.2pt,color=black] coordinates { (0.1210,3) };
\addplot[only marks,mark=*,mark size=.9pt,color=betweencol!70] coordinates { (0.8893,3.850) (0.8923,3.893) (0.8886,3.936) (0.8801,3.979) (0.8700,4.021) (0.8897,4.064) (0.8801,4.107) (0.8826,4.150) };
\addplot[only marks,mark=|,mark size=4pt,line width=1.2pt,color=black] coordinates { (0.8841,4) };
\end{axis}
\begin{axis}[at={(pa.south west)},anchor=north west,yshift=-2.45cm,name=pd,
 width=.40\textwidth,height=5.0cm,xmode=log,ymode=log,xmin=.4,xmax=8,ymin=.1,ymax=8,
 xtick={.5,1,2,4},xticklabels={0.5,1,2,4},ytick={.1,.3,1,3},yticklabels={0.1,0.3,1,3},
 xlabel={$\|F\|$, true latent RNA covariance},ylabel={$\|F\|$, estimated},
 title={\textbf{d}\enspace Compatibility in simulations},
 axis x line*=bottom,axis y line*=left,tick align=outside,xmajorgrids=true,grid style={gray!18},axis line style={gray!55},tick label style={font=\sffamily\scriptsize},label style={font=\sffamily\small},clip=false,title style={font=\sffamily\small,at={(0,1)},anchor=base west,yshift=5pt},ymajorgrids=true,
 legend style={draw=none,font=\sffamily\scriptsize,at={(0.5,-0.36)},anchor=north,legend columns=2,column sep=5pt,cells={anchor=west}}]
\draw[dashed,black!45] (axis cs:1,.1) -- (axis cs:1,8);
\draw[dashed,black!45] (axis cs:.4,1) -- (axis cs:8,1);
\draw[black!30] (axis cs:.4,.4) -- (axis cs:8,8);
\addplot[only marks,mark=*,mark size=1.1pt,color=haocol,mark options={fill=white},forget plot] coordinates { (1.6380,2.6064) (1.1506,2.5665) (1.7078,2.7831) (1.4065,3.0374) (1.5154,2.5685) (1.1926,2.5968) (1.4678,2.9238) (1.1491,2.6169) (1.3573,2.2635) (1.0703,2.3409) (1.5034,2.8600) (1.1911,2.7323) (1.0634,2.4075) (0.8041,1.9401) (1.3253,2.4248) (1.1412,2.4816) (1.3500,2.6345) (1.1619,2.6481) (1.4989,2.5157) (1.2137,2.4963) (1.4908,2.5153) (1.1780,2.5331) (1.2471,2.2828) (1.0763,2.3185) (2.2305,3.0809) (1.7245,2.6276) (2.2885,3.0462) (2.0287,2.9261) (1.9995,2.8089) (1.6469,2.5100) (1.7230,3.0443) (1.4265,2.5920) (1.7234,2.5602) (1.5625,2.4347) (1.8226,2.9459) (1.5184,2.6935) (1.1263,2.4245) (0.8701,1.8791) (1.4956,2.5230) (1.3083,2.4676) (1.5632,2.7406) (1.3946,2.6457) (1.9966,2.7724) (1.6496,2.4909) (1.9331,2.8806) (1.6154,2.4431) (1.4711,2.5077) (1.3298,2.4162) };
\addplot[only marks,mark=*,mark size=1.1pt,color=haocol] coordinates { (1.6380,0.3303) (1.1506,0.2023) (1.7078,0.5636) (1.4065,0.3337) (1.5154,0.4366) (1.1926,0.2714) (1.4678,0.3254) (1.1491,0.2607) (1.3573,0.3419) (1.0703,0.2154) (1.5034,0.4534) (1.1911,0.2796) (1.0634,0.2718) (0.8041,0.2257) (1.3253,0.3809) (1.1412,0.2892) (1.3500,0.4355) (1.1619,0.3256) (1.4989,0.3876) (1.2137,0.2447) (1.4908,0.3633) (1.1780,0.2314) (1.2471,0.3918) (1.0763,0.2143) (2.2305,1.0252) (1.7245,0.5514) (2.2885,1.2374) (2.0287,0.6123) (1.9995,0.9016) (1.6469,0.4441) (1.7230,0.4798) (1.4265,0.3132) (1.7234,0.4175) (1.5625,0.3597) (1.8226,0.5512) (1.5184,0.3622) (1.1263,0.2799) (0.8701,0.2258) (1.4956,0.4219) (1.3083,0.3967) (1.5632,0.3742) (1.3946,0.3906) (1.9966,0.7762) (1.6496,0.5532) (1.9331,0.5449) (1.6154,0.4263) (1.4711,0.5220) (1.3298,0.3671) };
\addlegendentry{Channel from donors}
\addplot[only marks,mark=*,mark size=1.1pt,color=withincol,mark options={fill=white},forget plot] coordinates { (3.1244,5.5307) (2.3693,5.3944) (3.1419,5.4734) (2.6837,5.6422) (3.0481,5.3390) (2.4623,5.1231) (2.7522,5.4794) (2.4048,5.5432) (2.4524,4.6844) (2.2732,4.9906) (2.8399,5.5688) (2.3404,5.6081) (2.4570,5.1592) (2.1900,4.9439) (2.3392,4.6427) (2.2058,5.0129) (2.4658,5.1845) (2.3441,5.4966) (2.8088,4.9408) (2.3462,4.8506) (2.8430,5.2463) (2.3930,5.2694) (2.6257,4.8455) (2.3377,4.9592) (4.3282,6.2126) (3.7985,5.7286) (4.2058,5.9755) (3.9516,5.9057) (4.0093,5.7686) (3.6110,5.3154) (3.1665,5.6804) (2.9163,5.6201) (3.5205,5.0301) (3.5447,5.2217) (3.4173,5.9182) (3.0983,5.7902) (2.6410,5.2272) (2.3940,4.9632) (2.6884,4.9120) (2.7018,5.2184) (2.8899,5.3518) (2.9047,5.5469) (3.8739,5.2710) (3.5268,4.9844) (4.0196,5.9274) (3.7431,5.6698) (3.3859,5.1526) (3.1469,5.1112) };
\addplot[only marks,mark=*,mark size=1.1pt,color=withincol] coordinates { (3.1244,2.0429) (2.3693,1.4177) (3.1419,2.4147) (2.6837,1.7531) (3.0481,2.4429) (2.4623,1.7568) (2.7522,1.8118) (2.4048,1.6092) (2.4524,2.0241) (2.2732,1.6242) (2.8399,2.1139) (2.3404,1.7132) (2.4570,1.5058) (2.1900,1.3081) (2.3392,1.9196) (2.2058,1.9606) (2.4658,2.0992) (2.3441,1.9558) (2.8088,2.0193) (2.3462,1.4569) (2.8430,2.0712) (2.3930,1.7354) (2.6257,1.8713) (2.3377,1.3587) (4.3282,3.8511) (3.7985,3.5940) (4.2058,3.9028) (3.9516,3.3904) (4.0093,3.5725) (3.6110,3.1347) (3.1665,2.7621) (2.9163,2.7318) (3.5205,2.8373) (3.5447,3.0486) (3.4173,2.8763) (3.0983,2.6498) (2.6410,1.7579) (2.3940,1.6589) (2.6884,2.1258) (2.7018,2.1484) (2.8899,2.4899) (2.9047,2.6367) (3.8739,3.1068) (3.5268,2.8184) (4.0196,3.4933) (3.7431,3.2197) (3.3859,3.0379) (3.1469,2.8528) };
\addlegendentry{Channel from donor-type groups}
\addplot[only marks,mark=*,mark size=1.1pt,color=coloncol,mark options={fill=white},forget plot] coordinates { (1.1700,1.9941) (0.9010,1.9764) (1.1101,1.9394) (0.9054,2.0520) (1.0687,1.8537) (0.8827,2.0205) (0.9316,1.8408) (0.8364,2.0130) (0.8515,1.5195) (0.7703,1.7602) (0.9556,1.9725) (0.7991,2.0903) (0.8193,1.8077) (0.7134,1.7556) (0.8376,1.7129) (0.7818,1.8807) (0.8675,1.8420) (0.7729,1.9244) (0.9654,1.6752) (0.8099,1.7360) (0.9911,1.7334) (0.7550,1.7730) (0.8892,1.5931) (0.7929,1.7611) (1.5903,2.2306) (1.4043,2.1815) (1.6314,2.2007) (1.5471,2.1982) (1.4786,2.1054) (1.3792,2.1119) (1.1318,1.9790) (1.0733,2.0633) (1.0967,1.6658) (1.0755,1.8375) (1.1655,2.1615) (1.0599,2.1765) (0.8515,1.8498) (0.7464,1.7920) (0.9465,1.7473) (0.9123,1.8669) (1.0045,1.9873) (0.9342,2.0145) (1.3073,1.9106) (1.1862,1.8973) (1.3453,2.0770) (1.1505,1.9754) (1.1511,1.7755) (1.0817,1.8621) };
\addplot[only marks,mark=*,mark size=1.1pt,color=coloncol] coordinates { (1.1700,0.2871) (0.9010,0.1588) (1.1101,0.2568) (0.9054,0.1625) (1.0687,0.3303) (0.8827,0.1997) (0.9316,0.2628) (0.8364,0.2226) (0.8515,0.1863) (0.7703,0.1449) (0.9556,0.2802) (0.7991,0.2261) (0.8193,0.1840) (0.7134,0.1556) (0.8376,0.2923) (0.7818,0.2455) (0.8675,0.2207) (0.7729,0.1867) (0.9654,0.3392) (0.8099,0.1548) (0.9911,0.2198) (0.7550,0.1453) (0.8892,0.2703) (0.7929,0.2054) (1.5903,0.8301) (1.4043,0.5609) (1.6314,0.9368) (1.5471,0.7179) (1.4786,0.6332) (1.3792,0.4947) (1.1318,0.3579) (1.0733,0.2886) (1.0967,0.2845) (1.0755,0.2154) (1.1655,0.4450) (1.0599,0.3790) (0.8515,0.1643) (0.7464,0.1474) (0.9465,0.2465) (0.9123,0.2081) (1.0045,0.3764) (0.9342,0.3060) (1.3073,0.7527) (1.1862,0.5562) (1.3453,0.3457) (1.1505,0.2109) (1.1511,0.5859) (1.0817,0.3506) };
\addlegendentry{Channel from colon}
\addlegendimage{only marks,mark=*,mark size=1.1pt,color=black!55,mark options={fill=white}}
\addlegendentry{Measured covariance}
\addlegendimage{only marks,mark=*,mark size=1.1pt,color=black!55}
\addlegendentry{Noise-corrected}
\end{axis}
\begin{axis}[at={(pd.east)},anchor=west,xshift=2.0cm,name=pe,
 width=.33\textwidth,height=5.0cm,xmode=log,log basis x=2,
 xtick={8,16,32,64,118},xticklabels={8,16,32,64,118},
 ymin=0.2,ymax=1.2,xlabel={Training donors},ylabel={No pairs: mean error},
 title={\textbf{e}\enspace More training populations},
 title style={font=\sffamily\small,at={(0,1)},anchor=base west,yshift=5pt},
 axis x line*=bottom,axis y line*=left,tick align=outside,
 ymajorgrids=true,grid style={gray!18},axis line style={gray!55},
 tick label style={font=\sffamily\scriptsize},label style={font=\sffamily\small},
 legend style={draw=none,font=\sffamily\scriptsize,at={(0.5,-0.36)},anchor=north,legend columns=1,cells={anchor=west}}]
\draw[dashed,black!45] (axis cs:6,1) -- (axis cs:150,1);
\addplot[color=haocol,mark=*,mark size=1.5pt,line width=.8pt,error bars/.cd,y dir=both,y explicit] coordinates { (8,0.6369) += (0,0.2951) -= (0,0.1975) (16,0.4927) += (0,0.2730) -= (0,0.1401) (32,0.3954) += (0,0.1964) -= (0,0.1213) (64,0.3365) += (0,0.0323) -= (0,0.0505) (118,0.2916) += (0,0.0000) -= (0,0.0000) };
\addlegendentry{Confirmatory blood}
\addplot[color=coloncol,mark=*,mark size=1.5pt,line width=.8pt,error bars/.cd,y dir=both,y explicit] coordinates { (8,1.0243) += (0,0.1363) -= (0,0.0964) (16,0.9917) += (0,0.1029) -= (0,0.1710) (32,0.9314) += (0,0.2152) -= (0,0.1100) (64,0.9960) += (0,0.0918) -= (0,0.0644) (118,0.9878) += (0,0.0000) -= (0,0.0000) };
\addlegendentry{Colon}
\end{axis}
\end{tikzpicture}

\caption{\textbf{Population variation determines which dependence is recovered.}
\textbf{a}--\textbf{c}, Each recipient's scoring-cell cross-covariance over level-1 (blood) or coarse (colon) cell types, split into a between-type and a within-type part, both scaled by the total standard deviations. \textbf{a}, Frobenius norm of each part relative to that of the total. \textbf{b}, Relative error of the prespecified pairing-free channel's prediction of each part, applied without saturation; the dashed line marks a prediction of zero. \textbf{c}, Share of the recipient's corresponding RNA covariance that lies in the effective supported subspace \textit{S}. Points are donors ($n=\XsHaoDonors$ blood and \XsColonDonors{} colon) and black marks are means; exploratory analyses added after scoring.
\textbf{d}, Compatibility statistic $\|F\|$ estimated from measured (open) and noise-corrected (filled) RNA covariance against its value for the true latent covariance, in simulations built from the \RcBmSamples{} bone-marrow samples at their observed depths (all cells and within coarse types; two constructions of latent expression). Values above one mean that no valid joint distribution agrees with the channel; the grey line marks equality.
\textbf{e}, Mean pairing-free error against the number of training donors. Points are means over 10 random subsets of donors and bars span them; all \XsPatients{} donors give a single fit.}
\label{sfig:ident}
\end{figure}

\begin{figure}[p]
\centering
% Generated by extension/synthesis/make_assets.py. Do not edit.
\definecolor{propcol}{HTML}{217D91}
\definecolor{pairedcol}{HTML}{5D6670}
\definecolor{semicol}{HTML}{C9803A}
\definecolor{champcol}{HTML}{8E4B8C}
\definecolor{ridgecol}{HTML}{9AA3AB}
\definecolor{meancol}{HTML}{A64442}
\definecolor{sitecol}{HTML}{6A8FB3}
\begin{tikzpicture}[font=\sffamily]
\begin{axis}[name=a,scale only axis,width=.3\textwidth,height=3.5cm,xmode=log,log basis x=2,
 xmin=20,xmax=1000,ymin=-20,ymax=40,xtick={25,50,100,200,400,800},xticklabels={25,50,100,200,400,800},ytick={-20,0,20,40},
 xlabel={Paired cells},ylabel={Dependence recovered (\%)},
 title={\textbf{a}\enspace Bone marrow, four sites},
 axis x line*=bottom,axis y line*=left,tick align=outside,ymajorgrids=true,grid style={gray!18},axis line style={gray!55},tick label style={font=\sffamily\scriptsize},label style={font=\sffamily\small},clip=false,title style={font=\sffamily\small,at={(0,1)},anchor=base west,yshift=5pt}]
\draw[dashed,black!45] (axis cs:20,17.93) -- (axis cs:1000,17.93);
\addplot[color=black,line width=.9pt,dashed,mark=*,mark size=1.3pt] coordinates { (25,1.87) (50,7.25) (100,14.24) (200,19.13) (400,26.57) (800,32.61) };
\addplot[color=sitecol,line width=.9pt,mark=diamond*,mark size=1.7pt] coordinates { (25,-20.00) (50,-20.00) (100,-6.41) (200,15.42) (400,26.29) (800,33.52) };
\addplot[color=propcol,line width=1.2pt,mark=*,mark size=1.6pt] coordinates { (25,-20.00) (50,-20.00) (100,-3.89) (200,17.54) (400,29.15) (800,35.86) };
\addplot[color=propcol,line width=1.1pt,densely dashed,mark=*,mark size=1.4pt,mark options={solid,fill=white}] coordinates { (25,0.57) (50,7.41) (100,16.60) (200,20.75) (400,29.10) (800,35.40) };
\addplot[color=semicol,line width=.9pt,mark=triangle*,mark size=1.6pt] coordinates { (25,-20.00) (50,-20.00) (100,-10.47) (200,0.48) (400,5.61) (800,7.28) };
\addplot[color=meancol,line width=.9pt,mark=x,mark size=1.8pt] coordinates { (25,-20.00) (50,-20.00) (100,-20.00) (200,-20.00) (400,-20.00) (800,-20.00) };
\addplot[color=pairedcol,line width=.8pt,mark=square*,mark size=1.3pt] coordinates { (25,-16.39) (50,0.00) (100,0.86) (200,1.60) (400,4.37) (800,5.92) };
\end{axis}
\begin{axis}[at={(a.outer north east)},anchor=outer north west,xshift=.4cm,name=b,scale only axis,width=.2\textwidth,height=3.5cm,xmode=log,log basis x=2,
 xmin=20,xmax=1600,ymin=.4,ymax=7.6,xtick={25,100,400,1600},xticklabels={25,100,400,$>$800},
 ytick={1,2,3,4,5,6,7},
 yticklabels={{Paired cells only},{Other sites, their averages},{Other sites, RNA from nuclei},{Other sites, contrasts (after scoring)},{Other sites},{Own site},{Same samples, their averages}},
 xlabel={Paired cells to the target},
 title={\textbf{b}\enspace Paired cells needed},
 axis x line*=bottom,axis y line*=left,tick align=outside,xmajorgrids=true,grid style={gray!18},axis line style={gray!55},tick label style={font=\sffamily\scriptsize},label style={font=\sffamily\small},clip=false,title style={font=\sffamily\small,at={(0,1)},anchor=base west,yshift=5pt},y tick style={draw=none}]
\draw[black!25] (axis cs:800,.4) -- (axis cs:800,7.6);
\draw[->,color=pairedcol,line width=.8pt] (axis cs:800,1) -- (axis cs:1400,1);
\addplot[only marks,mark=o,mark size=1.8pt,color=pairedcol] coordinates { (800,1) };
\draw[->,color=meancol,line width=.8pt] (axis cs:800,2) -- (axis cs:1400,2);
\addplot[only marks,mark=o,mark size=1.8pt,color=meancol] coordinates { (800,2) };
\draw[->,color=semicol,line width=.8pt] (axis cs:800,3) -- (axis cs:1400,3);
\addplot[only marks,mark=o,mark size=1.8pt,color=semicol] coordinates { (800,3) };
\draw[color=propcol,line width=1.6pt,densely dashed] (axis cs:20,4) -- (axis cs:124.9,4);
\node[font=\sffamily\scriptsize,anchor=west] at (axis cs:134.9,4) {125};
\draw[color=propcol,line width=2.4pt] (axis cs:20,5) -- (axis cs:204.7,5);
\node[font=\sffamily\scriptsize,anchor=west] at (axis cs:221.1,5) {205};
\draw[color=sitecol,line width=2.4pt] (axis cs:20,6) -- (axis cs:234.7,6);
\node[font=\sffamily\scriptsize,anchor=west] at (axis cs:253.5,6) {235};
\draw[color=black,line width=1.6pt,densely dashed] (axis cs:20,7) -- (axis cs:168.7,7);
\node[font=\sffamily\scriptsize,anchor=west] at (axis cs:182.2,7) {169};
\end{axis}
\path (a.outer south west) -- (b.outer south east) coordinate[pos=.5] (legab);
\begin{axis}[hide axis,scale only axis,width=1pt,height=1pt,xmin=0,xmax=1,ymin=0,ymax=1,at={(legab)},anchor=north,yshift=-2pt,legend style={draw=none,font=\sffamily\scriptsize,at={(0.5,0.5)},anchor=north,legend columns=2,column sep=6pt,cells={anchor=west}}]
\addlegendimage{color=black,line width=.9pt,dashed,mark=*,mark size=1.3pt}
\addlegendentry{Same samples, their averages}
\addlegendimage{color=sitecol,line width=.9pt,mark=diamond*,mark size=1.7pt}
\addlegendentry{Own site (to \ensuremath{-}642\% at 25)}
\addlegendimage{color=propcol,line width=1.2pt,mark=*,mark size=1.6pt}
\addlegendentry{Other sites (to \ensuremath{-}535\% at 25)}
\addlegendimage{color=propcol,line width=1.1pt,densely dashed,mark=*,mark size=1.4pt,mark options={solid,fill=white}}
\addlegendentry{Other sites, contrasts (after scoring)}
\addlegendimage{color=semicol,line width=.9pt,mark=triangle*,mark size=1.6pt}
\addlegendentry{Other sites, RNA from nuclei (to \ensuremath{-}295\% at 25)}
\addlegendimage{color=meancol,line width=.9pt,mark=x,mark size=1.8pt}
\addlegendentry{Other sites, their averages (to \ensuremath{-}843\% at 800)}
\addlegendimage{color=pairedcol,line width=.8pt,mark=square*,mark size=1.3pt}
\addlegendentry{Paired cells only}
\end{axis}
\begin{axis}[at={(a.outer south west)},anchor=outer north west,yshift=-1.95cm,name=c,scale only axis,width=.3\textwidth,height=3.5cm,xmode=log,log basis x=2,
 xmin=20,xmax=480,ymin=-20,ymax=80,xtick={25,50,100,200,400},xticklabels={25,50,100,200,400},ytick={-20,0,20,40,60,80},
 xlabel={Paired neurons},ylabel={Dependence recovered (\%)},
 title={\textbf{c}\enspace Fly wiring, other animals},
 axis x line*=bottom,axis y line*=left,tick align=outside,ymajorgrids=true,grid style={gray!18},axis line style={gray!55},tick label style={font=\sffamily\scriptsize},label style={font=\sffamily\small},clip=false,title style={font=\sffamily\small,at={(0,1)},anchor=base west,yshift=5pt}]
\addplot[color=black,line width=.9pt,dashed,mark=*,mark size=1.3pt] coordinates { (25,33.97) (50,45.97) (100,56.24) (200,63.95) (400,68.14) };
\addplot[color=propcol,line width=1.2pt,mark=*,mark size=1.6pt] coordinates { (25,36.57) (50,51.33) (100,59.99) (200,68.35) (400,74.12) };
\addplot[color=meancol,line width=.9pt,mark=x,mark size=1.8pt] coordinates { (25,-20.00) (50,-20.00) (100,-20.00) (200,-20.00) (400,-20.00) };
\addplot[color=pairedcol,line width=.8pt,mark=square*,mark size=1.3pt] coordinates { (25,13.53) (50,24.70) (100,36.49) (200,48.09) (400,58.97) };
\end{axis}
\begin{axis}[hide axis,scale only axis,width=1pt,height=1pt,xmin=0,xmax=1,ymin=0,ymax=1,at={(c.outer south west)},anchor=north west,yshift=-2pt,legend style={draw=none,font=\sffamily\scriptsize,at={(0,0.5)},anchor=north west,legend columns=1,column sep=5pt,cells={anchor=west}}]
\addlegendimage{color=black,line width=.9pt,dashed,mark=*,mark size=1.3pt}
\addlegendentry{Same animal, its averages}
\addlegendimage{color=propcol,line width=1.2pt,mark=*,mark size=1.6pt}
\addlegendentry{Other animals, contrasts}
\addlegendimage{color=meancol,line width=.9pt,mark=x,mark size=1.8pt}
\addlegendentry{Other animals, their averages (to \ensuremath{-}113\% at 50)}
\addlegendimage{color=pairedcol,line width=.8pt,mark=square*,mark size=1.3pt}
\addlegendentry{Paired neurons only, contrasts}
\end{axis}
\begin{axis}[at={(c.outer north east)},anchor=outer north west,xshift=.4cm,name=d,scale only axis,width=.3\textwidth,height=3.2cm,xmode=log,log basis x=2,
 xmin=20,xmax=1000,ymin=.6,ymax=1.7,xtick={25,50,100,200,400,800},xticklabels={25,50,100,200,400,800},
 ytick={.6,.8,1,1.2,1.4,1.6},xlabel={Paired cells},ylabel={Estimated / actual noise},
 title={\textbf{d}\enspace Noise estimate after contrasts},
 axis x line*=bottom,axis y line*=left,tick align=outside,ymajorgrids=true,grid style={gray!18},axis line style={gray!55},tick label style={font=\sffamily\scriptsize},label style={font=\sffamily\small},clip=false,title style={font=\sffamily\small,at={(0,1)},anchor=base west,yshift=5pt}]
\draw[dashed,black!45] (axis cs:20,1) -- (axis cs:1000,1);
\addplot[color=propcol,line width=1.1pt,mark=*,mark size=1.4pt] coordinates { (25,1.0049) (50,1.0080) (100,1.0109) (200,1.0136) (400,1.0155) (800,1.0164) };
\addplot[color=semicol,line width=.8pt,mark=square*,mark size=1.2pt] coordinates { (25,1.4784) (50,1.2576) (100,1.1339) (200,1.0732) (400,1.0446) (800,1.0307) };
\addplot[color=meancol,line width=.8pt,mark=x,mark size=1.8pt] coordinates { (25,0.6737) (50,0.7985) (100,0.8984) (200,0.9579) (400,0.9884) (800,1.0030) };
\addplot[color=propcol,line width=1.1pt,mark=*,mark size=1.4pt,densely dashed] coordinates { (25,1.0050) (50,1.0116) (100,1.0161) (200,1.0237) (400,1.0286) (800,1.0288) };
\addplot[color=semicol,line width=.8pt,mark=square*,mark size=1.2pt,densely dashed] coordinates { (25,1.5377) (50,1.3275) (100,1.1929) (200,1.1174) (400,1.0751) (800,1.0507) };
\addplot[color=meancol,line width=.8pt,mark=x,mark size=1.8pt,densely dashed] coordinates { (25,0.6501) (50,0.7585) (100,0.8555) (200,0.9340) (400,0.9825) (800,1.0058) };
\end{axis}
\begin{axis}[hide axis,scale only axis,width=1pt,height=1pt,xmin=0,xmax=1,ymin=0,ymax=1,at={(d.outer south east)},anchor=north east,yshift=-2pt,legend style={draw=none,font=\sffamily\scriptsize,at={(1,0.5)},anchor=north east,legend columns=1,column sep=5pt,cells={anchor=west}}]
\addlegendimage{color=propcol,line width=1.1pt,mark=*,mark size=1.4pt}
\addlegendentry{Implemented (contrasts)}
\addlegendimage{color=semicol,line width=.8pt,mark=square*,mark size=1.2pt}
\addlegendentry{Jackknife}
\addlegendimage{color=meancol,line width=.8pt,mark=x,mark size=1.8pt}
\addlegendentry{Own-mean centring}
\addlegendimage{color=black,line width=.8pt}
\addlegendentry{Validation}
\addlegendimage{color=black,line width=.8pt,densely dashed}
\addlegendentry{Bone marrow}
\end{axis}
\end{tikzpicture}

\caption{\textbf{Correlation structure carries over between sites and animals; averages do not.}
\textbf{a}, Bone-marrow test (prespecified; CITE-seq from four sites): fraction of the dependence within cell types recovered in a held-out donor against the number of paired cells from the other two donors of its site, pooled over the twelve held-out site-by-donor batches. Unpaired cells came from the same samples as the paired cells, whose averages centred them (as in development and the benchmark); from the paired cells' own site; from other sites and donors; from those donors with RNA measured in single nuclei; or from other sites whose averages centred the paired cells. Otherwise, as prespecified, the paired cells were centred on their own population averages and multiplied by $\sqrt{n/(n-1)}$ to restore their spread, or, dashed, centred by contrasts (other sites; exploratory). Dashed line, the accuracy target (half the best recovery), used in \textbf{b}.
\textbf{b}, Paired cells needed to reach the target; arrows, not reached within 800 paired cells.
\textbf{c}, Female fly connectome: fraction of the dependence between a neuron's brain and nerve-cord partners recovered in held-out neuron types with 25 to 400 paired neurons. Unpaired neurons came from the same animal, whose averages centred the paired neurons (the benchmark's design), or from two other animals, whose averages centred them (prespecified), or with the paired neurons centred by contrasts (exploratory). In \textbf{a} and \textbf{c}, values below $-20\%$ are drawn at $-20\%$, and the legend gives the lowest value of each such curve.
\textbf{d}, The noise estimate that sets each block's shrinkage, averaged over 500 independent resamples of paired cells with contrasts and scaling rebuilt each time, divided by the variance of the block averages across the resamples, summed over blocks and pooled over held-out units (follow-up): the implemented estimate from contrasts, the jackknife over cells, and centring on the drawn cells' own averages as in the bone-marrow test; validation test (solid) and bone-marrow test (dashed). Dashed line, a calibrated estimate.}
\label{sfig:transfer}
\end{figure}

\clearpage

\section{What unpaired data identify}\label{sec:ident}
Let RNA $x$ and protein $y$ of a cell population have joint law $P$. Separately measured assays reveal only the marginal laws $P_x$ and $P_y$, and paired cells reveal $P$. On finite supports with positive marginal probabilities and finite interaction $S(x,y)$, the law with these marginals that is closest in Kullback--Leibler divergence to the reference $e^{S(x,y)}P_x(x)P_y(y)$ has the form
\begin{equation}
 q(x,y)=\exp\{S(x,y)+a(x)+b(y)\},
 \label{eq:information-projection}
\end{equation}
with adjusting functions $a$ and $b$ that enforce the marginals; on finite supports it is computed by matrix scaling\cite{sinkhorn1964}. Proposition~\ref{prop:pairing} states that unpaired data carry no information about the interaction and what a channel shared by populations changes, Proposition~\ref{prop:means} what population means identify, Proposition~\ref{prop:mixture} what they identify when populations differ in cell-type composition, and Proposition~\ref{prop:gaussian} how a channel is transferred to a recipient with its own marginals.

\subsection*{Unpaired data alone}
\begin{proposition}[Identification without pairs]\label{prop:pairing}
\begin{enumerate}
\item On fixed finite supports, under Eq.~\eqref{eq:information-projection} with population-specific adjusting
functions, every finite interaction is compatible with every collection of strictly positive
within-assay probability mass functions. Unpaired data from any number of populations
therefore carry no information about the interaction.
\item Suppose instead that, across samples $s$, $y\mid x\sim N(Wx+b,\Psi)$
with $W,b,\Psi$ shared, $\Psi$ positive definite and the RNA law sample-specific. Then
$\mu_{y,s}=W\mu_{x,s}+b$ and $\Sigma_{y,s}=W\Sigma_{x,s}W^\top+\Psi$. If the
centred RNA means span $\mathbb R^p$, the separately measured assay moments
identify $W$, $b$ and then $\Psi$. This model is the special case of the
Gaussian reconstruction with $B=W^\top\Psi^{-1}$ in which the protein-side
adjusting function is shared.
\end{enumerate}
\end{proposition}

\begin{proof}
(1) Matrix scaling reaches any positive marginals for any finite interaction.
(2) The moment equations follow from the law of total expectation and
variance. With spanning centred means, the mean equations form a full-rank
linear system for $(W,b)$, and $\Psi$ follows from any covariance equation.
The joint density of the shared-channel model is proportional to
$r_s(x)\exp\{-\frac12(y-Wx-b)^\top\Psi^{-1}(y-Wx-b)\}$, whose cross term is
$x^\top W^\top\Psi^{-1}y$.
\end{proof}

With fewer samples than RNA features, the mean equations identify $W$ only on
the span of the sample means (Proposition~\ref{prop:means}). If $\Psi$ is
also shared, the covariance equations
$W(\Sigma_{x,s}-\Sigma_{x,t})W^\top=\Sigma_{y,s}-\Sigma_{y,t}$ constrain $W$
off that span as well. They are quadratic, so they determine $W$ at most up to
transformations that preserve every between-sample difference of RNA
covariances. The analyses therefore regularize $W$ and evaluate it on held-out
paired cells.

\subsection*{What population variation identifies}
Proposition~\ref{prop:means} states what population means identify, what the
ridge estimator of the main text recovers, and how far the unidentified
remainder can move a recipient's cross-covariance. It requires only that the
conditional mean of protein given RNA is shared; $\Psi$ may differ between
populations.

\begin{proposition}[Identification from population means]\label{prop:means}
Let populations $s=1,\dots,S$ share $E(y\mid x,s)=Wx+b$. Let their RNA means
$\mu_{x,s}$ have weights $w_s>0$ summing to one, weighted mean $\bar\mu_x$, and
between-population matrix
$G=\sum_s w_s(\mu_{x,s}-\bar\mu_x)(\mu_{x,s}-\bar\mu_x)^\top$. Let $M$ be the
column space of $G$, the span of the centred RNA means ($\dim M\le S-1$), with
orthogonal projector $P_M$.
\begin{enumerate}
\item A channel $W'$ reproduces every population's protein mean, with some
intercept, if and only if $W'P_M=WP_M$.
\item At the population moments, the weighted ridge regression of protein
means on RNA means with penalty $\lambda>0$ is
$W_\lambda=WG(G+\lambda I)^{-1}$. It keeps each eigendirection of $G$ with
eigenvalue $g$ in proportion $g/(g+\lambda)$, and $W_\lambda\to WP_M$ as
$\lambda\to0$.
\item For a recipient with within-assay covariances $\Sigma_x,\Sigma_y$ whose
channel agrees with $W$ on $M$, the cross-covariance $C=\Sigma_xW^\top$ is the
sum of $\Sigma_xP_MW^\top$, which the population means fix, and
$\Sigma_x(I-P_M)W^\top$, about which they carry no information.
\item Suppose $\Sigma_x$ and $\Sigma_y$ are positive definite. Let $Q_N$ and $Q_\perp$ be orthonormal bases of $N=\Sigma_x^{-1/2}M$ and
$N^\perp$, and $F=Q_N^\top\Sigma_x^{1/2}W^\top\Sigma_y^{-1/2}$, which depends
on $W$ only through $WP_M$. The cross-covariances compatible with $WP_M$ and
with a valid joint law of the recipient's two assays exist if and only if
$\|F\|_{\mathrm{op}}\le1$, and they are then exactly
\[
 \mathcal C=\bigl\{\Sigma_x^{1/2}(Q_NF+Q_\perp\Gamma D)\Sigma_y^{1/2}:
 \|\Gamma\|_{\mathrm{op}}\le1\bigr\},\qquad D=(I-F^\top F)^{1/2}.
\]
$\mathcal C$ is symmetric about $C_0=\Sigma_x^{1/2}Q_NF\Sigma_y^{1/2}$. Its
Chebyshev radius in Frobenius norm, the smallest worst-case error of any
estimate that uses only this information, is
\[
 r=\Bigl\{\sum_i\sigma_i\bigl(\Sigma_x^{1/2}Q_\perp\bigr)^2
 \sigma_i\bigl(D\Sigma_y^{1/2}\bigr)^2\Bigr\}^{1/2},
\]
with both sets of singular values in decreasing order.
\end{enumerate}
\end{proposition}

\begin{proof}
(1) The centred protein means are $W(\mu_{x,s}-\bar\mu_x)$, so two channels
give the same means exactly when they agree on every centred RNA mean, that is
on $M$; the intercept absorbs the rest. (2) The weighted cross-moment of the
centred protein and RNA means is $WG$, so the ridge solution is
$WG(G+\lambda I)^{-1}$; diagonalize $G$. (3) Write
$W^\top=P_MW^\top+(I-P_M)W^\top$ and apply (1). (4) The joint covariance is
positive semidefinite if and only if $K=\Sigma_x^{-1/2}C\Sigma_y^{-1/2}$ is a
contraction. For $m\in M$,
$(\Sigma_x^{-1/2}m)^\top K=(Wm)^\top\Sigma_y^{-1/2}$, so the component of $K$
along $N$ is $F$. Since $N^\perp=\Sigma_x^{1/2}M^\perp$, the component along
$N^\perp$ equals a term fixed by $WP_M$ plus
$Q_\perp^\top\Sigma_x^{1/2}(I-P_M)W^\top\Sigma_y^{-1/2}$, which takes every
value as $W$ varies off $M$. The stacked matrix with blocks $F$ and $X$ is a
contraction if and only if $\|F\|_{\mathrm{op}}\le1$ and
$X^\top X\preceq I-F^\top F$, that is $X=\Gamma D$
with $\|\Gamma\|_{\mathrm{op}}\le1$ (Douglas' lemma \cite{douglas1966}). Replacing $\Gamma$ by
$-\Gamma$ gives the symmetry. With $L=\Sigma_x^{1/2}Q_\perp$ and
$R=D\Sigma_y^{1/2}$, $\|L\Gamma R\|_F^2=\operatorname{tr}(\Gamma^\top
L^\top L\Gamma RR^\top)\le\sum_i\sigma_i(L)^2\sigma_i(R)^2$ by von Neumann's
trace inequality and $\sigma_i(\Gamma^\top L^\top L\Gamma)\le\sigma_i(L)^2$.
Equality holds for the partial isometry that aligns the singular vectors of
$L$ and $R$. For a set symmetric about $C_0$, every estimate $\hat C$ and
member $C_0+Z$ satisfy
$\max(\|\hat C-C_0-Z\|,\|\hat C-C_0+Z\|)\ge\|Z\|$, so the worst-case error is
smallest at $C_0$ and equals $r$.
\end{proof}

Two quantities follow. The share of a recipient's RNA variance inside $M$,
$\operatorname{tr}(P_M\Sigma_xP_M)/\operatorname{tr}\Sigma_x$, needs no paired
data. After scoring, the relative size of the unidentified term in (3),
computed with the recipient's own regression on its scoring cells, is the
error of an oracle that knows the recipient's channel on $M$ and nothing else.
In the analyses, $M$ is the span of the directions that the fitted ridge
operator keeps by at least one half ($g\ge\lambda$), in within-population
standard-deviation units.

\subsection*{Populations that differ in composition}
Proposition~\ref{prop:means} assumes a channel shared by every cell. Blood
samples differ from one another mainly in the proportions of cell types, and
the next result describes what their means identify when nothing is shared
except the laws of the types themselves.

\begin{proposition}[Populations that differ in composition]\label{prop:mixture}
Let cells of type $k=1,\dots,K$ have the same joint law of $(x,y)$ in every
population, with means $\mu_{x,k},\mu_{y,k}$, within-type covariances
$\Sigma_{xx,k},\Sigma_{yy,k}$ and within-type cross-covariance
$\Sigma_{xy,k}$. Let population $s$ contain the types in proportions
$\pi_s$, let the centred proportions $\pi_s-\bar\pi$ span the vectors whose
entries sum to zero, and let the RNA type means be affinely independent. For
proportions $\pi$ write $\bar\mu_x(\pi)=\sum_k\pi_k\mu_{x,k}$, and similarly
for protein, and let
$\Sigma^{\mathrm b}_x(\pi)=\sum_k\pi_k\{\mu_{x,k}-\bar\mu_x(\pi)\}\{\mu_{x,k}-\bar\mu_x(\pi)\}^\top$
and $\Sigma^{\mathrm w}_x(\pi)=\sum_k\pi_k\Sigma_{xx,k}$ be the between-type
and within-type RNA covariances.
\begin{enumerate}
\item Let $M$ be the span of the differences $\mu_{x,k}-\mu_{x,l}$. There is
a unique linear map $W_{\mathrm b}$ on $M$ with
$W_{\mathrm b}(\mu_{x,k}-\mu_{x,l})=\mu_{y,k}-\mu_{y,l}$ for all $k,l$. $M$ is
the span of the centred RNA means of the populations, their centred protein
means satisfy $\mu_{y,s}-\bar\mu_y=W_{\mathrm b}(\mu_{x,s}-\bar\mu_x)$ exactly,
and the population means determine $W_{\mathrm b}$.
\item For any proportions $\pi$, observed or not, the between-type part of the
cross-covariance is
\[
 \sum_k\pi_k\{\mu_{x,k}-\bar\mu_x(\pi)\}\{\mu_{y,k}-\bar\mu_y(\pi)\}^\top
 =\Sigma^{\mathrm b}_x(\pi)W_{\mathrm b}^\top .
\]
\item Changing the within-type cross-covariances $\Sigma_{xy,k}$, while keeping
every type's two marginal laws, changes neither assay's distribution in any
population. Unpaired data from any number of such populations are therefore
compatible with every coupling of the types' marginal laws, and carry no
information about the within-type part $\sum_k\pi_k\Sigma_{xy,k}$ beyond the
constraints that those marginals impose; for Gaussian types, every
$\Sigma_{xy,k}$ that keeps the type's joint covariance positive semidefinite
remains possible.
\item Let $W$ agree with $W_{\mathrm b}$ on $M$. The channel predicts the
within-type part as $\Sigma^{\mathrm w}_x(\pi)W^\top$, and this prediction is
correct for every $\pi$ if and only if $\Sigma_{xy,k}=\Sigma_{xx,k}W^\top$ for
every type, that is, if every type's regression of protein on RNA has the
slope $W$.
\end{enumerate}
\end{proposition}

\begin{proof}
(1) Because the centred proportions sum to zero,
$\mu_{x,s}-\bar\mu_x=\sum_k(\pi_{s,k}-\bar\pi_k)(\mu_{x,k}-\mu_{x,1})$, so the
centred RNA means lie in $M$, and they span it because the centred proportions
span the vectors that sum to zero. Affine independence makes
$v\mapsto\sum_kv_k\mu_{x,k}$ injective on those vectors, so
$W_{\mathrm b}(\sum_kv_k\mu_{x,k})=\sum_kv_k\mu_{y,k}$ is well defined on $M$;
it gives $\mu_{y,s}-\bar\mu_y=W_{\mathrm b}(\mu_{x,s}-\bar\mu_x)$.
Any linear map that reproduces the centred protein means agrees with
$W_{\mathrm b}$ on their span (Proposition~\ref{prop:means}(1)). (2) Each
$\mu_{x,k}-\bar\mu_x(\pi)=\sum_l\pi_l(\mu_{x,k}-\mu_{x,l})$ lies in $M$, and its
image under $W_{\mathrm b}$ is $\mu_{y,k}-\bar\mu_y(\pi)$; substitute. (3) A
population's RNA law is the mixture $\sum_k\pi_{s,k}L_{x,k}$ of the types' RNA
laws, and likewise for protein; neither involves how $x$ and $y$ are coupled
within a type, and any coupling of the type marginals is a valid law. (4) Both
sides are linear in $\pi$; compare them at the vertices of the simplex.
\end{proof}

Blood populations fit the shared-channel mean equation of
Proposition~\ref{prop:means} exactly under this model, but the map they
identify is the between-type slope, which need not describe any within-cell
channel. Applied to a
recipient's total RNA covariance
$\Sigma_x=\Sigma^{\mathrm b}_x+\Sigma^{\mathrm w}_x$, it predicts the
between-type part correctly and the within-type part as
$\Sigma^{\mathrm w}_xW^\top$. When the within-type slopes along $M$ are
shallower than the between-type slope, this overstates the protein covariance
along $W_{\mathrm b}M$, and $\|F\|_{\mathrm{op}}$ in
Proposition~\ref{prop:means}(4) can exceed one although the population means
fit the channel exactly. This is the ecological fallacy
\cite{robinson1950,simpson1951} in the form that matters here.

\subsection*{Transfer to a recipient}
\begin{proposition}[Closed-form transfer]\label{prop:gaussian}
Let $S_x\succ0$ and $S_y\succ0$, and write
$S_x^{1/2}BS_y^{1/2}=U\operatorname{diag}(\beta)V^\top$ for a singular value
decomposition. The reconstruction with interaction $x^\top By$ and marginals
$N(\mu_x,S_x)$ and $N(\mu_y,S_y)$ is the Gaussian law with cross-covariance
\[
 C=S_x^{1/2}U\operatorname{diag}\{g(\beta)\}V^\top S_y^{1/2},\qquad
 g(\beta)=\frac{\sqrt{1+4\beta^2}-1}{2\beta},\ g(0)=0,
\]
and its precision has off-diagonal block $-B$. Moreover $0\le g(\beta)<1$ and
$C=S_xBS_y-S_xBS_yB^\top S_xBS_y+O(\|B\|^5)$.
\end{proposition}

\begin{proof}
A Gaussian law whose precision has off-diagonal block $-B$ has density
$\exp\{x^\top By+a(x)+b(y)\}$ with quadratic $a,b$, which is the stationarity
condition for the strictly convex problem in
Eq.~\eqref{eq:information-projection}; so if such a law has the required
marginals, it is the reconstruction. Block inversion shows that the precision
of $\bigl(\begin{smallmatrix}S_x&C\\C^\top&S_y\end{smallmatrix}\bigr)$ has
off-diagonal block $-S_x^{-1}C(S_y-C^\top S_x^{-1}C)^{-1}$. Setting it to $-B$
and writing $\hat C=S_x^{-1/2}CS_y^{-1/2}$ gives
$\hat C=\hat B(I-\hat C^\top\hat C)$ with $\hat B=S_x^{1/2}BS_y^{1/2}$. The
choice $\hat C=U\operatorname{diag}(\gamma)V^\top$ solves it when
$\gamma=\beta(1-\gamma^2)$, whose root in $[0,1)$ is $g(\beta)$; then
$\hat C^\top\hat C\prec I$, so the joint covariance is positive definite.
Expanding $g(\beta)=\beta-\beta^3+O(\beta^5)$ gives the series.
\end{proof}

Its linear term, $S_xBS_y$, is the response of the reconstruction to a weak interaction; the second-order term vanishes because Gaussian third central moments are zero, and the cubic term shows how the transferred dependence saturates. The same
closed form describes entropic optimal transport between Gaussian measures
\cite{janati2020,mallasto2022}; here it transfers a channel learned without paired cells to a recipient with its own RNA and protein covariances (Supplementary Note~\ref{sec:cross}).

\section{Measurement noise and the latent RNA correlation}\label{sec:noise}
Population means average over cells, and hence over the sampling noise of
single-cell RNA counts; a recipient's single-cell covariance does not. The
next result states what this does to Proposition~\ref{prop:means}.

\begin{proposition}[Latent and measured RNA]\label{prop:noise}
In every population let each cell have a latent RNA state $z$ and measured RNA
$x=z+e$, with $E(e\mid z,y)=0$ and $\operatorname{Cov}(e\mid z)=N(z)$
diagonal, and let $E(y\mid z,s)=Az+b$ with $A$ shared by all populations.
Write $\Sigma_z$ and $\Sigma_x$ for a recipient's latent and measured RNA
covariances and $\bar N=E\,N(z)$.
\begin{enumerate}
\item Population means satisfy $E(y\mid s)=A\,E(x\mid s)+b$, so they identify
$A$, the channel acting on the latent state, on $M$
(Proposition~\ref{prop:means}(1)).
\item $\Sigma_x=\Sigma_z+\bar N$ and $\operatorname{Cov}(x,y)=\Sigma_zA^\top$;
the cell-level regression slope of $y$ on $x$ is
$A\Sigma_z\Sigma_x^{-1}$, which differs from $A$ whenever $\bar N\neq0$.
\item Let $F(\Sigma)=Q_N^\top\Sigma^{1/2}A^\top\Sigma_y^{-1/2}$ with $Q_N$ an
orthonormal basis of $\Sigma^{-1/2}M$. The cross-covariances compatible with
$AP_M$ and with a valid joint law of $(z,y)$ exist if and only if
$\|F(\Sigma_z)\|_{\mathrm{op}}\le1$, and
$\|F(\Sigma_x)\|_{\mathrm{op}}\ge\|F(\Sigma_z)\|_{\mathrm{op}}$. Measurement
noise can therefore empty the set of Proposition~\ref{prop:means}(4) computed
with $\Sigma_x$ when the latent set is not empty.
\item If $x$ is a function of counts $c\sim\mathrm{Poisson}(\lambda)$ given
$\lambda$, the binomial halves $c_1\mid c\sim\mathrm{Bin}(c,1/2)$ and
$c_2=c-c_1$ are independent given $\lambda$, so any functions $f(c_1)$ and
$f(c_2)$ have covariance $\operatorname{Var}\{E(f(c_1)\mid\lambda)\}$, the
latent variance at half depth, which requires no paired cell.
\end{enumerate}
\end{proposition}

\begin{proof}
(1) $E(y\mid s)=E\{E(y\mid z,s)\mid s\}=A\,E(z\mid s)+b$ and
$E(x\mid s)=E(z\mid s)$. (2) By the law of total covariance
$\Sigma_x=\Sigma_z+E\,N(z)$; $\operatorname{Cov}(e,y)=0$ because
$E(e\mid z,y)=0$, so $\operatorname{Cov}(x,y)=\operatorname{Cov}(z,y)
=\Sigma_zA^\top$. (3) The first statement is Proposition~\ref{prop:means}(4)
applied to $(z,y)$. With $V$ an orthonormal basis of $M$ and $A_M=AV$,
$F(\Sigma)^\top F(\Sigma)=\Sigma_y^{-1/2}A_M(V^\top\Sigma^{-1}V)^{-1}A_M^\top\Sigma_y^{-1/2}$,
because $\Sigma^{1/2}Q_NQ_N^\top\Sigma^{1/2}=V(V^\top\Sigma^{-1}V)^{-1}V^\top$.
$\Sigma_x\succeq\Sigma_z$ gives $V^\top\Sigma_x^{-1}V\preceq V^\top\Sigma_z^{-1}V$
and, inverting, $F(\Sigma_x)^\top F(\Sigma_x)\succeq F(\Sigma_z)^\top F(\Sigma_z)$.
(4) Binomial thinning of a Poisson count gives independent Poisson counts with
half the mean; functions of independent variables are independent, and the
law of total covariance gives the stated covariance.
\end{proof}

In the analyses, $x$ is log-normalized expression. For each gene, the
correlation across cells of the log-normalized values of the two halves is its
split-half reliability; the Spearman--Brown formula $2r/(1+r)$ gives the
reliability at full depth, and one minus this is the share $\nu_j$ of the
gene's standardized variance that is sampling noise. The latent correlation
matrix is estimated as $R_x-\operatorname{diag}(\nu_jR_{x,jj})$, after the off-diagonal entries of $R_x$
have been multiplied by 0.9, with eigenvalues floored at 0.001 times the mean diagonal. The Spearman--Brown step treats the two
half-depth values as parallel measurements, which is exact on the count scale
and approximate after the logarithm, and library normalization couples genes;
in simulations with known latent covariance both effects were small
(below).

The proof of (3) shows that $\|F(\Sigma)\|^2_{\mathrm{op}}$ is the largest
eigenvalue of the channel applied to $(V^\top\Sigma^{-1}V)^{-1}$, the variance
of the $M$-coordinates of $x$ that the orthogonal coordinates leave unexplained.
Noise raises this conditional variance, so the statistic computed with measured
covariance bounds the latent one from above. After the noise diagonal is removed
from a covariance estimated from $n$ cells, fluctuations of the noise covariance
of spectral norm of order $\nu(p/n)^{1/2}$ remain off the diagonal\cite{marchenko1967}.
They create near-null directions that shrink the conditional variance, and in
simulations the statistic computed with the corrected covariance understated
the latent one, so that the two computed statistics bracket it
(below).

\subsection*{Validation of the noise correction}
Each of \RcBmSamples{} bone-marrow samples of the NeurIPS 2021 CITE-seq data\cite{luecken2021sandbox} (adaptation cells) defined a
latent truth. Each cell's expected proportions of the panel genes were the
average of the observed proportions (count divided by library) over the cell and
its 14 nearest neighbours in 30 principal components of standardized log
expression, and the rest of the transcriptome kept the remaining share. In a
second construction, independent gene-specific lognormal variation was added,
sized so that each gene's latent within-type variance matched the real data's
within-type variance times one minus its count-splitting noise share. Expected
libraries were the observed ones times 0.5, 1 or 2. Counts of the panel genes and
of the rest of the transcriptome were drawn from Poisson distributions, and the
library was their sum, so normalization coupled genes as in real data. Forty
independent replicate draws per cell gave the noise covariance
$N=E\{\operatorname{Cov}(x\mid\text{truth})\}$, a full matrix, and the latent
covariance $\operatorname{Cov}\{E(x\mid\text{truth})\}$, estimated as the
covariance of replicate means minus $N/40$, for
$x=\log(1+10^4\,\text{count}/\text{library})$. The first replicate was analysed
by the pipeline of the pairing-free analyses, and the truth was expressed on the
same scale (standard deviations of the analysed replicate, same shrinkage), with
the pairing-free channels of Supplementary Note~\ref{sec:cross} and the sample's real protein correlation.

Estimated noise shares had a median absolute error of
\RcSimNuMaeMin--\RcSimNuMaeMax{} and a median bias of at most \RcSimNuBiasMax{} in
every setting (Supplementary Table~\ref{tab:sim}). Off-diagonal noise covariance,
which normalization creates, was at most \RcSimCouplingMax\% of the off-diagonal
latent covariance in norm. Correcting the diagonal reduced the relative error of
the RNA correlation matrix from \RcSimErrMeasMin--\RcSimErrMeasMax{} to
\RcSimErrCorrMin--\RcSimErrCorrMax, as much as correction with the true noise
variances did. The compatibility statistic was nevertheless biased: across
settings and channels, its median ratio to the latent value was
\RcSimCorrRatioMin--\RcSimCorrRatioMax{} with corrected covariance
(\RcSimFrozenRatioMin--\RcSimFrozenRatioMax{} for the prespecified channel) and
\RcSimMeasRatioMin--\RcSimMeasRatioMax{} with measured covariance, and the latent
value lay between the two in \RcSimBracket{} of \RcSimBracketTotal{} cases. Correction
with the true noise variances gave nearly the same statistics (median ratio to
the estimated correction, \RcSimOracleMin--\RcSimOracleMax), so the bias arises
from finite-cell fluctuations of the noise covariance rather than from the
estimated shares (Proposition~\ref{prop:noise} and the remark after it). At the observed
depth, the prespecified blood channel's latent statistic exceeded one in all \RcSimDOne{}
simulated total analyses (median \RcSimFrozenTrue), while its corrected statistic
exceeded one in \RcSimFrozenCorrAbove.

\begin{table}[h]
\caption{\textbf{Validation of the noise correction.} Medians over
\RcBmSamples{} bone-marrow samples. Noise share, true median share of sampling
noise in gene variance; share error, mean absolute error of the count-splitting
estimate; matrix error, relative Frobenius error of the measured and corrected RNA
correlation matrices against the latent one (all with the same shrinkage);
$\|F\|$ / true, ratio of the compatibility statistic of the prespecified blood channel
computed with measured or corrected covariance to its latent value. Depth, factor
applied to the observed library sizes.}
\label{tab:sim}
\centering\footnotesize\setlength{\tabcolsep}{4pt}
% Generated by extension/recoverability/make_assets.py. Do not edit.
\begin{tabular}{llrcccccc}
\toprule
& & & & & \multicolumn{2}{c}{Matrix error} & \multicolumn{2}{c}{$\|F\|$ / true, blood channel} \\
\cmidrule(lr){6-7}\cmidrule(lr){8-9}
Latent & Cells & Depth & Noise share & Share error & Measured & Corrected & Measured & Corrected \\
\midrule
Smoothed & All & 0.5 & 0.56 & 0.055 & 0.29 & 0.11 & 2.25 & 0.23 \\
 & All & 1 & 0.39 & 0.041 & 0.18 & 0.08 & 1.76 & 0.27 \\
 & All & 2 & 0.25 & 0.028 & 0.12 & 0.06 & 1.46 & 0.47 \\
 & Within types & 0.5 & 0.78 & 0.067 & 0.95 & 0.31 & 2.79 & 0.22 \\
 & Within types & 1 & 0.64 & 0.057 & 0.61 & 0.23 & 2.18 & 0.23 \\
 & Within types & 2 & 0.49 & 0.041 & 0.38 & 0.16 & 1.79 & 0.26 \\
Plus gene-specific & All & 0.5 & 0.48 & 0.055 & 0.32 & 0.13 & 1.92 & 0.27 \\
 & All & 1 & 0.33 & 0.039 & 0.20 & 0.10 & 1.55 & 0.29 \\
 & All & 2 & 0.21 & 0.027 & 0.13 & 0.07 & 1.37 & 0.63 \\
 & Within types & 0.5 & 0.66 & 0.070 & 1.01 & 0.38 & 2.09 & 0.25 \\
 & Within types & 1 & 0.52 & 0.055 & 0.67 & 0.28 & 1.67 & 0.27 \\
 & Within types & 2 & 0.36 & 0.038 & 0.44 & 0.20 & 1.43 & 0.41 \\
\bottomrule
\end{tabular}

\end{table}

\section{The bone-marrow test: unpaired cells from other sites, assays and animals}\label{sec:deploy}
In every other test, the unpaired cells came from the same samples as the paired cells, with one modality withheld. This note describes the bone-marrow test, a prespecified test in which they came from other sites and donors or, for RNA, from another assay; two analyses of that test added after scoring, one of how the paired cells are centred and one of where RNA from nuclei fails; and the analyses of the fly connectome in which the unpaired neurons came from other animals.

\subsection*{Protocol}
The plan, its code and the two data files were registered before any estimator of this test was run on the data under its design. Before that, the pipeline had been run end to end on synthetic values with the real design, and the unpaired summaries of two folds had been computed on the real pools to check that the condition maps were stable (pool sizes, reliabilities, map norms and the latent eigenvalue spectrum); no paired-cell estimate and no statistic of held-out cells had been computed. Predictions of every fold were recorded before any held-out cell was read.

The data were the NeurIPS 2021 bone-marrow CITE-seq and multiome profiles\cite{luecken2021sandbox} as extracted for the benchmark (3,000 cells per batch), in the eight coarse cell types present in both assays; the benchmark had analysed the CITE-seq and multiome data sets with unpaired cells of the same samples. Every CITE-seq batch (site by donor, 12 in all) was held out once. The paired reservoir was the 25\% of the cells of the other two batches of its site selected by the development's pseudo-random rule (about \DpReservoir{} cells). The unpaired pools were the remaining cells of those two batches (own site, \DpPoolSameMin--\DpPoolSameMax{} cells), every CITE-seq cell of the other sites' batches whose donor did not appear at the held-out site (other sites, \DpPoolOtherMin--\DpPoolOtherMax), and, for the nuclear-RNA arm, every multiome nucleus of those sites and donors, with protein from the other-site pool. Pool cells contributed RNA or protein, never both. Panels were chosen from the other-site pool: 200 genes among those measured in both assays by the benchmark's dispersion rule and the antibodies detected in at least 1\% of cells.

The estimator, its blocks, map and comparators were those of the benchmark, unchanged. In every draw, paired cells were centred on the means of their own populations of cell type and donor and scaled by their own pooled standard deviations, with centred values multiplied by $\sqrt{n/(n-1)}$ for a population of $n\ge2$ drawn cells and single cells dropped, so that the unpaired cells supplied only correlation structure: the eigenbases, the blocks' noise and the condition map. Two secondary arms instead centred and scaled the paired cells by a pool's statistics: the own-site pool's population means, the design of every earlier test, and the other-site pool's cell-type means. Paired-only estimation used the same own-centred cells. Budgets were 25 to 800 paired cells with five draws, and all arms of a draw used the same paired cells. The endpoint was the noise-unbiased recovered fraction of within-population gene--protein cross-correlation of held-out cells, pooled over cell types and folds. The primary target was half of the best recovery that any arm reached; D1 (primary) and D2 required 95\% lower bounds above one for the saving over paired-only estimation with other-site and with nuclear unpaired profiles, and D3 an upper bound below 1.25 for the ratio of paired cells needed with other-site to own-site unpaired profiles. Intervals came from 2,000 resamples of held-out batches stratified by site.

\subsection*{Results}
The best recovery was \DpLevel\%, so the primary target was \DpTarget\% (Supplementary Table~\ref{tab:deploy}). With other-site unpaired cells, the estimator reached it with \DpCellsOther{} paired cells, whereas paired-only estimation did not reach it within \DpMaxBudget{} (saving at least \DpSaveOther; 95\% CI, \DpSaveOtherLo--\DpSaveOtherHi; D1 \DpDOne), and it needed \DpNonInf{} times the paired cells needed with own-site unpaired cells (\DpNonInfLo--\DpNonInfHi; D3 \DpDThree). At a quarter and three quarters of the best recovery, the savings were \DpSaveOtherQuarter{} (\DpSaveOtherQuarterLo--\DpSaveOtherQuarterHi) and \DpSaveOtherThreeQ{} (\DpSaveOtherThreeQLo--\DpSaveOtherThreeQHi), and the ratios to own-site cells \DpNonInfQuarter{} and \DpNonInfThreeQ. Every site gave a saving (\DpSiteMin{} to \DpSiteMax; Supplementary Table~\ref{tab:deploysite}). With nuclear RNA, the estimator recovered at most \DpRfAssayMax\% and never reached the target, so D2 was \DpDTwo. Centred on the other sites' means, every estimate was worse than independence (at best \DpRfMeansBest\%), and with means from unpaired cells of the same samples, the design of the earlier tests, the estimator was informative from \DpFirstBudget{} paired cells and reached the target with \DpCellsSamePool. Below \DpOwnNegativeBelow{} paired cells, the own-centred estimates were worse than independence.

\begin{table}[h]
\caption{\textbf{Bone-marrow test: recovered fraction of within-cell-type dependence by source of unpaired cells.} Pooled over the twelve held-out batches; recovered fraction (\%) at each budget of paired cells, averaged over five draws, and paired cells needed to reach the primary target (\DpTarget\%, half of the best recovery), interpolated on the running maximum of each curve. Paired cells only, the best at each budget of James--Stein shrinkage, SCOSE, FCOSE and cross-validated low rank, pooled or per cell type. Same samples, their averages: paired cells centred on the population averages of unpaired cells of the same batches; other sites, their averages: centred on the other sites' cell-type averages; RNA from nuclei: unpaired RNA from multiome nuclei of the other sites; all other rows: paired cells centred and scaled on their own statistics, as prespecified.}
\label{tab:deploy}
\centering\footnotesize
% Generated by extension/synthesis/make_assets.py. Do not edit.
\begin{tabular}{@{}lrrrrrrr@{}}
\toprule
 & \multicolumn{6}{c}{Recovered (\%) with paired cells} & Paired cells \\
\cmidrule(lr){2-7}
Unpaired cells & 25 & 50 & 100 & 200 & 400 & 800 & to the target \\
\midrule
Paired cells only & \ensuremath{-}16 & 0.0 & 0.9 & 1.6 & 4.4 & 5.9 & $>$800 \\
Same samples, their averages & 1.9 & 7.2 & 14 & 19 & 27 & 33 & 169 \\
Own site & \ensuremath{-}642 & \ensuremath{-}116 & \ensuremath{-}6.4 & 15 & 26 & 34 & 235 \\
Other sites & \ensuremath{-}535 & \ensuremath{-}93 & \ensuremath{-}3.9 & 18 & 29 & 36 & 205 \\
Other sites, RNA from nuclei & \ensuremath{-}295 & \ensuremath{-}58 & \ensuremath{-}10 & 0.5 & 5.6 & 7.3 & $>$800 \\
Other sites, their averages & \ensuremath{-}443 & \ensuremath{-}619 & \ensuremath{-}660 & \ensuremath{-}765 & \ensuremath{-}827 & \ensuremath{-}843 & $>$800 \\
\bottomrule
\end{tabular}

\end{table}

\begin{table}[h]
\caption{\textbf{Bone-marrow test by site.} Savings over paired-only estimation at half of each site's best recovery, pooled over the site's three held-out batches. $\geq$, paired-only estimation did not reach the target within 800 paired cells, so the saving is a lower bound; n.r., neither curve reached it.}
\label{tab:deploysite}
\centering\footnotesize
% Generated by extension/synthesis/make_assets.py. Do not edit.
\begin{tabular}{@{}lrrr@{}}
\toprule
Site & Best recovery (\%) & Other sites & Other sites, RNA from nuclei \\
\midrule
s1 & 41 & $\geq$4.1 & n.r. \\
s2 & 32 & $\geq$3.3 & n.r. \\
s3 & 37 & $\geq$2.0 & n.r. \\
s4 & 36 & $\geq$4.4 & n.r. \\
\bottomrule
\end{tabular}

\end{table}

\subsection*{Centring within populations}
Centring a population's drawn cells on their own mean makes them dependent. For $n$ cells with $\tilde x_i=\sqrt{n/(n-1)}\,(x_i-\bar x)$ and $\tilde y_i$ likewise, the products $z_i=\tilde x_i\tilde y_i^\top$ have mean equal to the population's cross-covariance, but their sum is $n$ times the unbiased sample cross-covariance, whose entries have variance $(\sigma_{x}^2\sigma_{y}^2+\sigma_{xy}^2)/(n-1)$ for Gaussian cells, whereas each product has variance $\sigma_{x}^2\sigma_{y}^2+\sigma_{xy}^2$. The spread of the products therefore understates the variance of their sum by the factor $n/(n-1)$; for two cells the two products are equal. Every estimator that takes its noise from the spread of per-cell products (the James--Stein factors of the proposed estimator and of the paired-only comparators, SCOSE and FCOSE) then shrinks too little, most in the large blocks dominated by noise, and cross-validation over cells (low rank) sees the cells of a population on both sides of a split. At 25 paired cells, \PsSmallTwentyFive\% of the drawn cells lay in populations of at most three cells, and at 50, \PsSmallFifty\%.

The $n-1$ orthonormal Helmert contrasts of a population's cells, $h_j=(u_1+\dots+u_j-j\,u_{j+1})/\sqrt{j(j+1)}$ for $j=1,\dots,n-1$, applied to the RNA vectors ($u_i=x_i$, giving $h^x_j$) and to the protein vectors ($u_i=y_i$, giving $h^y_j$), have mean zero and the population's covariance, are independent for Gaussian cells, and their cross-products sum to those of the centred cells, $\sum_jh^x_jh^{y\top}_j=\sum_i(x_i-\bar x)(y_i-\bar y)^\top$. Replacing each population's cells by its contrasts therefore keeps the estimand and makes the rows uncorrelated, at the cost of one row per population; their products are independent only for Gaussian cells, and Proposition~\ref{prop:contrast} gives the resulting bias of the noise estimate (Supplementary Note~\ref{sec:review}). After the test had been scored, every arm was recomputed with the paired cells so replaced and scaled by the contrasts' pooled standard deviations, with every other setting unchanged (the prespecified own-centred arm of the first fold was reproduced exactly). With other-site unpaired cells, the estimator then recovered \CtOtherFirst\% of the dependence with 25 paired cells, \CtOtherFifty\% with 50 and \CtOtherHundred\% with 100, against \CtFrozenOtherFirst\% and \CtFrozenOtherHundred\% with the prespecified centring, and it matched the prespecified centring from 400 paired cells on (Supplementary Table~\ref{tab:deploycon}). It reached the primary target with \CtCellsOther{} paired cells, \CtVsFrozen{} times fewer than with the prespecified centring (95\% CI, \CtVsFrozenLo--\CtVsFrozenHi), while paired-only estimation with contrasts did not reach it within 800 (saving at least \CtSaveOther; \CtSaveOtherLo--\CtSaveOtherHi); with own-site unpaired cells it needed \CtCellsSame{} (ratio of other-site to own-site cells, \CtOverSame; \CtOverSameLo--\CtOverSameHi). The other-site pools held about four times as many cells as the own-site pools. With 20 fresh draws per budget from the same reservoirs, the noise term of the estimator's blocks matched the variance of the block means across draws under contrasts (ratio \DpCalConMin{} to \DpCalConMax{} at 25 to 200 paired cells) and was too small under the test's centring (\DpCalDfMin{} to \DpCalDfMax; Supplementary Table~\ref{tab:audit}); 500 independent resamples passed through the whole pipeline later showed the estimate under contrasts slightly conservative (Supplementary Table~\ref{tab:noisecal}).

\begin{table}[h]
\caption{\textbf{Bone-marrow test with paired cells centred by contrasts (after scoring).} As Supplementary Table~\ref{tab:deploy}, with every population's drawn paired cells replaced by their Helmert contrasts and scaled by the contrasts' pooled standard deviations; target as there.}
\label{tab:deploycon}
\centering\footnotesize
% Generated by extension/synthesis/make_assets.py. Do not edit.
\begin{tabular}{@{}lrrrrrrr@{}}
\toprule
 & \multicolumn{6}{c}{Recovered (\%) with paired cells} & Paired cells \\
\cmidrule(lr){2-7}
Unpaired cells & 25 & 50 & 100 & 200 & 400 & 800 & to the target \\
\midrule
Paired cells only & 0.3 & 0.9 & 2.1 & 2.4 & 4.4 & 5.9 & $>$800 \\
Own site & \ensuremath{-}1.5 & 4.5 & 14 & 18 & 26 & 33 & 185 \\
Other sites & 0.6 & 7.4 & 17 & 21 & 29 & 35 & 125 \\
Other sites, RNA from nuclei & \ensuremath{-}2.6 & \ensuremath{-}0.2 & 4.6 & 4.0 & 6.3 & 7.3 & $>$800 \\
\bottomrule
\end{tabular}

\end{table}

\subsection*{Where RNA from nuclei fails}
Three analyses after scoring located the failure of nuclear RNA; all used pool or reservoir cells, and the second also held-out cells. First, the pools' latent RNA correlation matrices were compared in every fold. Relative to the own-site pool, the leading five eigenvectors of the other-site pool overlapped by \DpOverlapOther{} (mean squared cosine of the principal angles) and the leading 20 by \DpOverlapOtherTwenty, and the off-diagonal entries correlated at \DpROther; for the nuclear pool, the same quantities were \DpOverlapAssay, \DpOverlapAssayTwenty{} and \DpRAssay. Second, with paired cells centred by contrasts, the estimate of each pool was scored with and without its condition map, and the nuclear bases were combined with the other-site maps. With 800 paired cells, the estimates without maps recovered \MpOtherUnmapped\% (other sites) and \MpAssayUnmapped\% (nuclear RNA) of the within-cell-type dependence, against \MpOtherMapped\% and \MpAssayMapped\% with their maps; the nuclear bases with the other-site maps recovered \MpAssayOtherMap\%. Most of the within-cell-type dependence thus came through the condition maps, which are built from each cell type's RNA covariance, and nuclear RNA failed mainly through them. Third, the law of Supplementary Note~\ref{sec:pred} was applied to each pool's bases with block moments of all reservoir cells centred by contrasts: the leading five RNA eigenvectors carried \LwWSame\%, \LwWOther\% and \LwWAssay\% of the reservoir's dependence in the own-site, other-site and nuclear bases, and the predicted savings of block over single-block shrinkage were \LwSame, \LwOther{} and \LwAssay{} (geometric means over folds). The nuclear bases still concentrated the dependence, less than whole-cell bases did.

\subsection*{Unpaired neurons from other animals}
A prespecified secondary analysis of the benchmark kept the paired BANC neurons, the folds and the held-out cell types of the female connectome but took the unpaired pool from two other animals: brain-side wiring from the FAFB brain connectome\cite{dorkenwald2024,schlegel2024} and nerve-cord-side wiring from the MANC nerve-cord connectome\cite{takemura2024}, with the feature panels chosen from these pools. As in every earlier analysis, the paired neurons were centred on the pool's means and scaled by its standard deviations, those of other animals, so the product of the offsets between the animals' means entered the mean of their cross-products. With 100 paired neurons, the estimator recovered \GenRfHundredBancX\% and paired-only estimation with the same centring \GenPairedOnlyHundredBancX\%; SemiCCA recovered \GenSemiHundredBancX\% and reference regression \GenRidgeHundredBancX\%, and the best recovery of any method was \GenLevelBancX\% (reference regression refitted per condition, 400 paired neurons), so no saving could be estimated.

After scoring, the analysis was repeated with every setting unchanged (folds, reservoir, budgets, draws, seeds, unpaired pool and held-out cell types) except that the paired neurons were centred within their populations by Helmert contrasts and scaled by the contrasts' pooled standard deviations. The other animals then contributed only correlation structure, and the estimator recovered \FxOwnHundred\% with 100 paired neurons, against \FxPairedOwnHundred\% for paired-only estimation and \FxSameHundred\% with the same animal's unpaired neurons (Supplementary Table~\ref{tab:fly}). With two conditions and many neurons in each, centring on the own means with the factor $\sqrt{n/(n-1)}$, as in the bone-marrow test, changed the estimator's recovered fractions by at most \FxDiffProposed{} and those of paired-only estimation by at most \FxDiffPaired{} percentage points.

\begin{table}[h]
\caption{\textbf{Female fly connectome with unpaired neurons of other animals.} Recovered fraction (\%) of the dependence between brain and nerve-cord wiring in held-out neuron types, by budget of paired neurons (averaged over five draws). Centred on its or their averages: paired neurons centred and scaled by the unpaired pool's statistics (for the same animal, the benchmark; for other animals, the prespecified secondary analysis); centred by contrasts: by Helmert contrasts (after scoring). Paired neurons only, the best of the paired-only estimators at each budget.}
\label{tab:fly}
\centering\footnotesize
% Generated by extension/synthesis/make_assets.py. Do not edit.
\begin{tabular}{@{}lrrrrr@{}}
\toprule
 & \multicolumn{5}{c}{Recovered (\%) with paired neurons} \\
\cmidrule(lr){2-6}
Estimator & 25 & 50 & 100 & 200 & 400 \\
\midrule
Same animal, centred on its averages (benchmark) & 34 & 46 & 56 & 64 & 68 \\
Other animals, centred on their averages (prespecified) & \ensuremath{-}30 & \ensuremath{-}113 & \ensuremath{-}78 & \ensuremath{-}101 & \ensuremath{-}109 \\
Other animals, centred by contrasts (after scoring) & 37 & 51 & 60 & 68 & 74 \\
Paired neurons only, same animal's averages & 14 & 23 & 36 & 46 & 55 \\
Paired neurons only, other animals' averages & 0.0 & \ensuremath{-}69 & \ensuremath{-}52 & \ensuremath{-}66 & \ensuremath{-}69 \\
Paired neurons only, centred by contrasts & 14 & 25 & 36 & 48 & 59 \\
\bottomrule
\end{tabular}

\end{table}

\section{The validation test: a blood study not analysed before}\label{sec:val}
This note describes the prespecified test of the workflow that the analyses of Supplementary Note~\ref{sec:deploy} arrived at after the bone-marrow test had been scored: a few paired cells of a small study, centred within their populations by Helmert contrasts, with unpaired correlation structure from an external reference, at small budgets, in a study that this work had not analysed. The centring was thus defined after one test had been scored, and it was tested here under a protocol registered before any paired-cell estimate or held-out statistic of the new study had been computed. The note then gives the allocation of donors to the folds, the biological readout of the test, specified after it had been scored, the readout's calibration and a comparison with other methods given the same inputs, both specified after scoring, and the robustness of the test's result.

\subsection*{Protocol}
The plan, its code and the data files were registered once the data had been extracted. Before that, the data had been extracted (counts and metadata; cell counts, library sizes and the antibody name match were the only statistics computed), the pipeline had been run end to end on synthetic values with the real design, and the unpaired summaries of the atlas and of two other-pool references had been computed to check their stability (cells, populations, map reliabilities, map norms and leading latent eigenvalues); no paired-cell estimate and no statistic of held-out cells had been computed. Predictions of every fold were recorded before any held-out cell was read.

\paragraph{Data.} GEO GSE314416\cite{babdor2026} profiles blood of healthy adults by 10x 5$'$ CITE-seq with the TotalSeq-C panel of 140 antibodies. Only its pools with antibody counts were used: \VlNFolds{} samples of \VlNDonors{} donors, 16 of them at two time points in different pools, in \VlNPools{} processing pools of two to four GEM wells. Cell types were the authors' labels, mapped by name to B, CD4 T (including regulatory T cells), CD8 T, NK and CD14 monocytes, the five types that the atlas holds with at least 10 cells per patient in each assay half (it averages 15 CD16 monocytes and 7 dendritic cells per patient). Those labels were transferred from a reference annotated using RNA and protein. The atlas was the extract of the Stephenson et al. COVID-19 blood study\cite{stephenson2021} used in the benchmark (118 patients from three UK sites; 600 cells per patient; 192 TotalSeq-C antibodies), with its labels mapped to the same five types. Its annotations used RNA clusters supported by protein markers. These supplied annotations were shared by all methods.

\paragraph{Features.} The gene panel was 200 genes chosen by the benchmark's binned-dispersion rule from the atlas's RNA-half cells, among the atlas extract's genes, all of which the validation data measure. Validation antibodies were matched to atlas antibodies by antibody name or, when a target gene identified one unmatched antibody on each side, by target gene (isotype controls never matched); the \VlNAntibodies{} matched antibodies detected in at least 1\% of the atlas's protein-half cells formed the protein panel. RNA was log-normalized to $10^4$ counts per cell and protein transformed by the centred log-ratio over the panel. The panel was fixed for every fold and used no statistic of the validation data.

\paragraph{Folds and references.} Every sample was held out once. Within its pool, samples were put in a fixed pseudo-random order, and the held-out sample's paired samples were the next two in that order, so that every fold is a small paired study of two samples of one pool. The paired reservoir was 25\% of the paired samples' cells, selected by the development's pseudo-random rule (about 1,500--2,200 cells), and held-out cells were split into two halves pseudo-randomly. The unpaired references were the atlas (primary); every cell of the samples in other pools whose donor does not appear in the held-out sample's pool (other pools); and the non-reservoir cells of the paired samples (own samples). RNA and protein summaries were computed from disjoint sets of reference cells, conditional on the supplied annotations.

\paragraph{Donor separation.} Every fold is donor-disjoint (Supplementary Table~\ref{tab:allocation}; checked from the sample labels after scoring). Each pool holds at most one sample of a donor, and the \AlNTwoTimepoints{} donors sampled twice have their two samples in different pools, so the two paired samples come from two donors other than the held-out donor. The other-pool reference excludes every donor of the held-out sample's pool, and with it the held-out donor's other sample; the own-sample reference consists of the paired samples' remaining cells; and the atlas is another study. The held-out cells that define every endpoint, including both halves of the readout below, therefore come from a donor that contributed no paired and no reference cell to its fold. Across folds, each sample serves as a paired sample in two other folds of its pool, and the bootstrap resamples donors with all their held-out samples.

\begin{table}[h]
\caption{\textbf{Sample, donor and pool allocation of the validation test.} Each row is a fold: the held-out sample's donor and time point (b, baseline; f, follow-up), the two paired samples of its pool (the next two in the fixed order), the other-pool reference (every sample of the other pools whose donor does not appear in the held-out sample's pool; the same for every fold of a pool) and the own-sample reference (the paired samples' non-reservoir cells). Cells of the five cell types analysed. No donor appears on both sides of any fold.}
\label{tab:allocation}
\centering\footnotesize
% Generated by extension/synthesis/make_assets.py. Do not edit.
\begin{tabular}{@{}lllrrr@{}}
\toprule
 & & & \multicolumn{2}{c}{Other-pool reference} & Own-sample \\
\cmidrule(lr){4-5}
Pool & Held-out sample & Paired samples & Samples (donors) & Cells & reference cells \\
\midrule
DB8 & HS6 (b) & HS32 (b), HS94 (f) & 28 (15) & 105,238 & 6,262 \\
 & HS32 (b) & HS94 (f), HS6 (b) &  &  & 6,505 \\
 & HS94 (f) & HS6 (b), HS32 (b) &  &  & 6,679 \\
\midrule
DB9 & HS86 (f) & HS24 (f), HS46 (b) & 22 (12) & 83,929 & 5,294 \\
 & HS24 (f) & HS46 (b), HS88 (b) &  &  & 5,749 \\
 & HS46 (b) & HS88 (b), HS16 (f) &  &  & 5,557 \\
 & HS88 (b) & HS16 (f), HS10 (b) &  &  & 5,268 \\
 & HS16 (f) & HS10 (b), HS86 (f) &  &  & 5,467 \\
 & HS10 (b) & HS86 (f), HS24 (f) &  &  & 5,301 \\
\midrule
DB10 & HS68 (b) & HS24 (b), HS47 (f) & 26 (14) & 99,582 & 5,021 \\
 & HS24 (b) & HS47 (f), HS93 (f) &  &  & 4,770 \\
 & HS47 (f) & HS93 (f), HS68 (b) &  &  & 5,526 \\
 & HS93 (f) & HS68 (b), HS24 (b) &  &  & 5,777 \\
\midrule
DB11 & HS46 (f) & HS22 (f), HS68 (f) & 22 (12) & 81,772 & 6,346 \\
 & HS22 (f) & HS68 (f), HS75 (b) &  &  & 5,771 \\
 & HS68 (f) & HS75 (b), HS88 (f) &  &  & 5,914 \\
 & HS75 (b) & HS88 (f), HS28 (b) &  &  & 5,983 \\
 & HS88 (f) & HS28 (b), HS46 (f) &  &  & 5,273 \\
 & HS28 (b) & HS46 (f), HS22 (f) &  &  & 5,779 \\
\midrule
DB12 & HS10 (f) & HS93 (b), HS47 (b) & 24 (13) & 92,020 & 5,719 \\
 & HS93 (b) & HS47 (b), HS16 (b) &  &  & 5,780 \\
 & HS47 (b) & HS16 (b), HS75 (f) &  &  & 5,386 \\
 & HS16 (b) & HS75 (f), HS10 (f) &  &  & 5,910 \\
 & HS75 (f) & HS10 (f), HS93 (b) &  &  & 6,161 \\
\midrule
DB13 & HS22 (b) & HS6 (f), HS48 (b) & 27 (14) & 100,636 & 5,759 \\
 & HS6 (f) & HS48 (b), HS30 (f) &  &  & 5,669 \\
 & HS48 (b) & HS30 (f), HS22 (b) &  &  & 5,546 \\
 & HS30 (f) & HS22 (b), HS6 (f) &  &  & 5,636 \\
\midrule
DB14 & HS12 (f) & HS30 (b), HS94 (b) & 23 (12) & 86,518 & 4,861 \\
 & HS30 (b) & HS94 (b), HS86 (b) &  &  & 5,029 \\
 & HS94 (b) & HS86 (b), HS28 (f) &  &  & 6,041 \\
 & HS86 (b) & HS28 (f), HS32 (f) &  &  & 5,775 \\
 & HS28 (f) & HS32 (f), HS12 (f) &  &  & 5,468 \\
 & HS32 (f) & HS12 (f), HS30 (b) &  &  & 5,566 \\
\bottomrule
\end{tabular}

\end{table}

\paragraph{Estimation and arms.} Populations were cell type by sample. In every draw, each population's drawn cells were replaced by their $n-1$ Helmert contrasts, scaled by the contrasts' pooled standard deviations, and the estimator, blocks, map and comparators were those of the benchmark with the reference's unpaired summaries. Arms: paired-only estimation (the best at each budget of James--Stein shrinkage, SCOSE, FCOSE and cross-validated low rank, pooled or per cell type, on the same contrasts); the estimator with each reference; and two secondary arms with the own-sample reference's population means, the design of the earlier tests, or with the atlas and the centring of the bone-marrow test ($\sqrt{n/(n-1)}$ times the deviations from the paired cells' own means). Budgets were 25 to 800 paired cells with five draws.

\paragraph{Endpoint and hypotheses.} The endpoint was the noise-unbiased recovered fraction of within-cell-type gene--protein cross-correlation of the held-out sample, pooled over cell types and folds; paired cells needed were read from the running maximum of each curve, and a curve that did not reach the target counted as needing 800. The primary target was half of the best recovery of any arm. V1 (primary): the saving over paired-only estimation with the atlas, 95\% lower bound above one. V2: the atlas arm's recovered fraction, 95\% lower bounds above zero with 50 and with 100 paired cells. V3: as V1 with the other-pool reference. V4 (non-inferiority): paired cells needed with the atlas over those needed with the own-sample reference, 95\% upper bound below 1.25. Intervals came from 2,000 resamples of donors, each carrying its held-out samples.

\subsection*{Results}
The best recovery was \VlLevel\%, so the primary target was \VlTarget\% (Supplementary Table~\ref{tab:val}). With the atlas as the only reference, the estimator reached it with \VlCellsAtlas{} paired cells, whereas paired-only estimation did not reach it within 800 (saving at least \VlSaveAtlas; 95\% CI, \VlSaveAtlasLo--\VlSaveAtlasHi; V1 \VlVOne). Its recovered fraction was \VlRfAtlasTwentyFive\% with 25 paired cells (\VlRfAtlasTwentyFiveLo{} to \VlRfAtlasTwentyFiveHi\%), \VlRfAtlasFifty\% with 50 (\VlRfAtlasFiftyLo--\VlRfAtlasFiftyHi\%) and \VlRfAtlasHundred\% with 100 (\VlRfAtlasHundredLo--\VlRfAtlasHundredHi\%; V2 \VlVTwo), against \VlRfPairedFifty\% and \VlRfPairedHundred\% for paired-only estimation with 50 and 100. With the other-pool reference, it reached the target with \VlCellsOther{} paired cells (saving at least \VlSaveOther; \VlSaveOtherLo--\VlSaveOtherHi; V3 \VlVThree), and with the own-sample reference with \VlCellsOwn; the atlas needed \VlAtlasOverOwn{} times as many paired cells as the own samples (\VlAtlasOverOwnLo--\VlAtlasOverOwnHi; V4 \VlVFour) and \VlAtlasOverOther{} times as many as the other pools (\VlAtlasOverOtherLo--\VlAtlasOverOtherHi). At a quarter and at three quarters of the best recovery, the atlas's savings over paired-only estimation were at least \VlSaveAtlasQuarter{} (\VlSaveAtlasQuarterLo--\VlSaveAtlasQuarterHi) and \VlSaveAtlasThreeQ{} (\VlSaveAtlasThreeQLo--\VlSaveAtlasThreeQHi). Every pool gave a saving with the atlas: at least \VlPoolMin{} to \VlPoolMax{} at each pool's own target, half of the best recovery in that pool, and at least \RbPoolMin{} to \RbPoolMax{} at the test's primary target (Supplementary Table~\ref{tab:pools}). The secondary arm with the own samples' means reached the target with \VlCellsOwnPool{} paired cells; the atlas with the centring of the bone-marrow test recovered \VlRfDfTwentyFive\% and \VlRfDfFifty\% with 25 and 50 paired cells and reached the target with \VlCellsDf.

\begin{table}[h]
\caption{\textbf{Validation test: recovered fraction of within-cell-type dependence by reference.} Pooled over the \VlNFolds{} held-out samples; recovered fraction (\%) at each budget of paired cells, averaged over five draws, and paired cells needed to reach the primary target (\VlTarget\%, half of the best recovery), interpolated on the running maximum of each curve. Paired cells were centred within their populations by contrasts in every arm except the last two: own samples, their averages, centred on the population averages of the own-sample reference; atlas, centred as in bone marrow, centred on their own population averages with the factor $\sqrt{n/(n-1)}$ of the bone-marrow test.}
\label{tab:val}
\centering\footnotesize
% Generated by extension/synthesis/make_assets.py. Do not edit.
\begin{tabular}{@{}lrrrrrrr@{}}
\toprule
 & \multicolumn{6}{c}{Recovered (\%) with paired cells} & Paired cells \\
\cmidrule(lr){2-7}
Arm & 25 & 50 & 100 & 200 & 400 & 800 & to the target \\
\midrule
Paired cells only & 0.5 & 2.4 & 3.5 & 7.7 & 14 & 13 & $>$800 \\
Atlas (other study) & 0.5 & 11 & 23 & 33 & 42 & 48 & 147 \\
Other pools, other donors & 8.8 & 23 & 34 & 43 & 50 & 57 & 70 \\
Own samples & 5.6 & 19 & 27 & 34 & 40 & 45 & 110 \\
Own samples, their averages & 8.2 & 19 & 27 & 34 & 40 & 44 & 114 \\
Atlas, centred as in bone marrow & \ensuremath{-}408 & \ensuremath{-}60 & 13 & 31 & 42 & 48 & 178 \\
\bottomrule
\end{tabular}

\end{table}

\subsection*{Degrees of freedom and the noise term}
Contrasts spend one cell of every population with a drawn cell and all of a single cell, and the budget counts every drawn cell. Supplementary Table~\ref{tab:audit} gives, for each budget, the populations with a drawn cell, the single cells and the contrast rows; paired-only estimation used the same contrasts, so every comparison is made at equal degrees of freedom. With 25 paired cells over ten populations, \VlRowsTwentyFive{} contrast rows remained on average, with 50, \VlRowsFifty, and with 100, \VlRowsHundred.

The table also compares, for the estimator's blocks along the atlas's axes, the noise term that sets each block's shrinkage (the trace of the covariance of the block mean estimated from the spread of per-row products) with the variance of the block mean across the five draws, corrected for the finite reservoir (prespecified). The same comparison was made in the bone-marrow test, along the other sites' axes, with \DpCalDraws{} fresh draws per budget from its reservoirs (after scoring). Under contrasts the two agreed at every budget from 25 to 200 paired cells in both data sets (ratio \CalCon; with 500 independent resamples through the whole pipeline, the estimate was slightly conservative, Supplementary Table~\ref{tab:noisecal}); under the centring of the bone-marrow test the noise term was too small (\VlCalDfMin{} to \VlCalDfMax{} in the validation and \DpCalDfMin{} to \DpCalDfMax{} in bone marrow), most at the smallest budgets, where that centring gave estimates worse than independence.

Two further checks, made after scoring, concern the budget axis and the paired-only control. Counted in contrast rows rather than drawn cells, with the mean rows at each budget as the axis and the same interpolation, the atlas arm reached the primary target with \VlRowsAtlas{} rows and paired-only estimation did not reach it within \VlRowsMax, a saving of at least \VlSaveAtlasRows{} against \VlSaveAtlas{} counted in cells; because populations cost a larger share of small budgets, counting drawn cells favours paired-only estimation. In the bone-marrow test, paired-only estimation under contrasts recovered \DpPoConTwentyFive\%, \DpPoConFifty\% and \DpPoConHundred\% with 25, 50 and 100 paired cells, against \DpPoDfTwentyFive\%, \DpPoDfFifty\% and \DpPoDfHundred\% under the test's own centring, and the two agreed to within 0.1 percentage points at 400 and 800 (Supplementary Tables~\ref{tab:deploy} and~\ref{tab:deploycon}), so centring by contrasts does not weaken the paired-only control.

\begin{table}[h]
\caption{\textbf{Degrees of freedom of the paired draws and calibration of the noise term.} Means over draws and folds. Contrast rows: drawn cells minus populations with a drawn cell. Noise term / variance across draws: over the estimator's blocks along the reference's axes (validation, atlas; bone marrow, other sites) and over folds, the sum of the estimated noise traces of the block means divided by the sum of their variances across draws, each variance divided by $1-B/N$ for $B$ paired cells drawn from a reservoir of $N$; one means calibrated, below one means the noise term is too small. Validation: five draws per budget (prespecified); bone marrow: \DpCalDraws{} fresh draws per budget (after scoring).}
\label{tab:audit}
\centering\footnotesize
% Generated by extension/synthesis/make_assets.py. Do not edit.
\begin{tabular}{@{}lrrrr@{}}
\toprule
Paired cells drawn & 25 & 50 & 100 & 200 \\
\midrule
\multicolumn{5}{@{}l}{\textit{Validation (prespecified)}} \\
Populations with a drawn cell & 8.7 & 9.7 & 9.9 & 10.0 \\
Single cells (no contrast) & 2.2 & 0.9 & 0.2 & 0.0 \\
Contrast rows & 16.3 & 40.3 & 90.1 & 190.0 \\
Noise term / variance across draws, contrasts & 1.00 & 1.00 & 1.00 & 1.00 \\
Noise term / variance across draws, own averages, $\sqrt{n/(n-1)}$ & 0.67 & 0.79 & 0.89 & 0.95 \\
\multicolumn{5}{@{}l}{\textit{Bone marrow (after scoring)}} \\
Populations with a drawn cell & 10.8 & 13.1 & 14.9 & 15.6 \\
Single cells (no contrast) & 4.1 & 2.7 & 1.6 & 0.8 \\
Contrast rows & 14.2 & 36.9 & 85.1 & 184.4 \\
Noise term / variance across draws, contrasts & 1.00 & 1.00 & 1.00 & 1.00 \\
Noise term / variance across draws, own averages, $\sqrt{n/(n-1)}$ & 0.64 & 0.75 & 0.84 & 0.91 \\
\bottomrule
\end{tabular}

\end{table}

\subsection*{Biological readout (specified after scoring)}
The test measured the fraction of within-cell-type cross-correlation recovered in a held-out sample. A laboratory would use such an estimate to say which genes and proteins co-vary within a cell type, and in which direction. After the test had been scored, a second analysis was written and registered with its code, the data and the test's recorded predictions before any of its endpoints had been computed; what was known then were the test's pooled recovered fractions, paired cells needed and verdicts.

\paragraph{Predictions.} The test stored the mean of the five draws' predictions for every fold, arm and budget, whereas a single study obtains one draw. The predictions of every draw were therefore regenerated with the test's registered code, folds, panels, references and seeds; at the test's six budgets their means reproduced the recorded ones (largest absolute difference \BoIntegrity{} over \BoCompared{} predictions, squared norms identical). A budget of 150 paired cells, just above the 147 with which the atlas arm had reached the test's primary target, was added with new draws under the same seed rule.

\paragraph{Reproducible associations.} In each held-out sample and cell type with at least 20 held-out cells in each of its two halves (\BoNUnits{} units), the correlation of every gene with every protein was computed in each half, and a pair was a reproducible association if both correlations had the same sign with two-sided $P<0.01$ by Fisher's $z$. A pair without association passes with probability below $5\times10^{-5}$. There were \BoNRep{} reproducible associations, \BoRepShare\% of the \BoNPairs{} pairs of each unit (B cells \BoRepB, CD4 T \BoRepCDFour, CD8 T \BoRepCDEight, NK \BoRepNK, CD14 monocytes \BoRepMono).

\paragraph{Endpoints and hypotheses.} Direction: the fraction of reproducible associations whose sign the estimate gave (an estimate of zero counting one half; chance 50\%). Nomination: the fraction of the 100 pairs with the largest absolute estimate in each unit that were reproducible associations of the estimated sign (ties at the boundary by their expectation; random pairs \BoRandomTwo\%, a perfect ranking \BoPerfectTwo\%). Cognate pairs: the recovered fraction, by the test's formula, over the correlations between the \BoNCognate{} panel genes and the panel proteins they encode (CD11b and \textit{ITGAM}, CD8 and \textit{CD8A}, CD14, CD16 and \textit{FCGR3A}, CD31 and \textit{PECAM1}, CD161 and \textit{KLRB1}, KLRG1, HLA-DR and \textit{HLA-DRA}, CD314 and \textit{KLRK1}, CD79b, CD122 and \textit{IL2RB}, CD83, CD124 and \textit{IL4R}, CD127 and \textit{IL7R}, CD94 and \textit{KLRD1}, CD85j and \textit{LILRB1}). Each endpoint was computed per draw, averaged over draws and pooled over units; paired-only estimation was, for each endpoint and budget, the best of the eight paired-only methods. B1 and B2 required 95\% lower bounds above zero for the difference between the atlas arm and paired-only estimation in direction and in nomination at 50, 100 and 150 paired cells, and B3 in the cognate recovered fraction at 100 and 150; intervals came from 2,000 resamples of donors.

\paragraph{Results.} All three hypotheses were met (B1 \BoBOne, B2 \BoBTwo, B3 \BoBThree; Supplementary Table~\ref{tab:readout}). Direction: with the atlas, \BoOneAtlasFifty\%, \BoOneAtlasHundred\% and \BoOneAtlasOneFifty\% at 50, 100 and 150 paired cells, against \BoOnePoFifty\%, \BoOnePoHundred\% and \BoOnePoOneFifty\% for paired-only estimation (differences \BoDiffOneFifty, \BoDiffOneHundred{} and \BoDiffOneOneFifty{} percentage points; 95\% CI, \BoDiffOneFiftyLo--\BoDiffOneFiftyHi, \BoDiffOneHundredLo--\BoDiffOneHundredHi{} and \BoDiffOneOneFiftyLo--\BoDiffOneOneFiftyHi). Nomination: \BoTwoAtlasFifty\%, \BoTwoAtlasHundred\% and \BoTwoAtlasOneFifty\%, against \BoTwoPoFifty\%, \BoTwoPoHundred\% and \BoTwoPoOneFifty\% (differences \BoDiffTwoFifty, \BoDiffTwoHundred{} and \BoDiffTwoOneFifty{} points; \BoDiffTwoFiftyLo--\BoDiffTwoFiftyHi, \BoDiffTwoHundredLo--\BoDiffTwoHundredHi{} and \BoDiffTwoOneFiftyLo--\BoDiffTwoOneFiftyHi). Cognate pairs: \BoThreeAtlasHundred\% and \BoThreeAtlasOneFifty\% at 100 and 150 paired cells, against \BoThreePoHundred\% and \BoThreePoOneFifty\% (differences \BoDiffThreeHundred{} and \BoDiffThreeOneFifty{} points; \BoDiffThreeHundredLo--\BoDiffThreeHundredHi{} and \BoDiffThreeOneFiftyLo--\BoDiffThreeOneFiftyHi). Per donor, with 100 paired cells, the atlas arm gave the direction of more associations than paired-only estimation in \RvDonorBetterPo{} of \RvDonorN{} donors and than the best form of reference regression in \RvDonorBetterReg{} (Fig.~2e of the main text). The other-pool and own-sample references gave the direction of more associations than the atlas at every budget, whereas from 100 paired cells on the atlas nominated more reproducible pairs than the own-sample reference and about as many as the other-pool reference.

To give the direction of \BoTargetOne\% of the associations (the midpoint between chance and the best arm), the atlas arm needed \BoCellsOneAtlas{} paired cells, the other-pool reference \BoCellsOneOther{}, the own-sample reference \BoCellsOneOwn{} and paired-only estimation \BoCellsOnePo{} (saving with the atlas \BoSaveOne; 95\% CI, \BoSaveOneLo--\BoSaveOneHi). To make \BoTargetTwo\% of the top pairs reproducible (half of the best arm), they needed \BoCellsTwoAtlas, \BoCellsTwoOther, \BoCellsTwoOwn{} and \BoCellsTwoPo{} (saving \BoSaveTwo; \BoSaveTwoLo--\BoSaveTwoHi).

By cell type, the atlas gave the direction of more associations than paired-only estimation at every budget in CD4 T, CD8 T and NK cells, which held \BoRepCDFour, \BoRepCDEight{} and \BoRepNK{} of the associations. In B cells it gave slightly more up to 400 paired cells (\BoTypeBAtlasHundred\% against \BoTypeBPoHundred\% with 100), and in CD14 monocytes about as many below 400 (\BoTypeMonoAtlasHundred\% against \BoTypeMonoPoHundred\% with 100) and fewer from 400 on. Among cognate pairs that were reproducible associations, paired-only estimation gave the direction slightly more often than the atlas at every budget (\BoCogPoHundred\% against \BoCogAtlasHundred\% with 100 paired cells); these associations are strong, and the paired cells alone recover their sign.

\begin{table}[h]
\caption{\textbf{Biological readout of the validation test.} Pooled over the scored units of the \VlNFolds{} held-out samples; each endpoint computed per draw and averaged over five draws. Paired-only estimation is the best of the eight paired-only methods for each endpoint and budget; every arm centred the paired cells by contrasts. Direction: chance 50\%. Top 100 pairs: random pairs \BoRandomTwo\%, a perfect ranking \BoPerfectTwo\%. By cell type, the number of reproducible associations in parentheses.}
\label{tab:readout}
\centering\footnotesize
% Generated by extension/synthesis/make_assets.py. Do not edit.
\begin{tabular}{@{}lrrrrrrr@{}}
\toprule
 & \multicolumn{7}{c}{Paired cells} \\
\cmidrule(lr){2-8}
Arm & 25 & 50 & 100 & 150 & 200 & 400 & 800 \\
\midrule
\multicolumn{8}{@{}l}{\textit{Direction of reproducible associations (\%)}} \\
Paired cells only & 61 & 67 & 75 & 80 & 82 & 88 & 92 \\
Atlas (other study) & 71 & 80 & 86 & 89 & 90 & 92 & 94 \\
Other pools, other donors & 77 & 85 & 89 & 91 & 92 & 94 & 96 \\
Own samples & 76 & 84 & 88 & 89 & 90 & 92 & 94 \\
\midrule
\multicolumn{8}{@{}l}{\textit{Top 100 pairs reproducible with the estimated sign (\%)}} \\
Paired cells only & 6.9 & 15 & 28 & 35 & 40 & 52 & 61 \\
Atlas (other study) & 23 & 35 & 47 & 52 & 55 & 62 & 67 \\
Other pools, other donors & 28 & 38 & 46 & 50 & 54 & 61 & 67 \\
Own samples & 27 & 36 & 41 & 44 & 46 & 52 & 57 \\
\midrule
\multicolumn{8}{@{}l}{\textit{Cognate pairs: recovered fraction (\%)}} \\
Paired cells only & 4.7 & 11 & 25 & 33 & 39 & 54 & 65 \\
Atlas (other study) & 13 & 24 & 39 & 47 & 53 & 64 & 73 \\
Other pools, other donors & 14 & 24 & 36 & 44 & 50 & 63 & 74 \\
Own samples & 14 & 23 & 34 & 40 & 45 & 57 & 67 \\
\midrule
\multicolumn{8}{@{}l}{\textit{Cognate pairs: direction (\%)}} \\
Paired cells only & 82 & 90 & 95 & 97 & 98 & 99 & 99 \\
Atlas (other study) & 76 & 85 & 93 & 96 & 97 & 99 & 99 \\
Other pools, other donors & 78 & 87 & 93 & 96 & 97 & 99 & 99 \\
Own samples & 78 & 85 & 92 & 94 & 96 & 98 & 99 \\
\midrule
\multicolumn{8}{@{}l}{\textit{Direction by cell type, atlas / paired only (\%)}} \\
B (5,077) & 57 / 54 & 63 / 57 & 67 / 61 & 70 / 63 & 72 / 66 & 77 / 74 & 81 / 82 \\
CD4 T (27,388) & 77 / 63 & 87 / 71 & 92 / 77 & 93 / 81 & 94 / 84 & 96 / 89 & 97 / 94 \\
CD8 T (56,637) & 74 / 64 & 85 / 70 & 91 / 79 & 94 / 83 & 95 / 86 & 96 / 93 & 97 / 96 \\
NK (10,526) & 63 / 60 & 71 / 64 & 78 / 69 & 82 / 73 & 82 / 74 & 87 / 81 & 90 / 87 \\
CD14 Mono (21,021) & 60 / 58 & 66 / 64 & 73 / 72 & 76 / 77 & 79 / 79 & 84 / 86 & 88 / 92 \\
\bottomrule
\end{tabular}

\end{table}

\subsection*{Calibration of the readout (specified after scoring)}
After the test and its readout had been scored, a further plan was registered with its code, the data, the test's recorded predictions and the readout's results before any quantity in it had been computed. It specified the calibration of the readout described here and the comparison with other methods in the next section. What was known then were the test's and the readout's results and the donor allocation; no count or endpoint under the truth sets below, no value per held-out sample and cell type and no share of positive associations had been computed. Its code regenerated the per-draw predictions of every arm as the readout did (largest difference from the recorded means \AltIntegrity{} over \AltCompared{} predictions) and reproduced every value of the readout under the readout's own definitions (largest difference \AltReproduces, in one paired-only method at one budget, from the number of computing threads).

\paragraph{Truth sets.} Besides the readout's definition ($P<0.01$ in both halves with the same sign), reproducible associations were defined at $P<0.05$ and at $P<0.001$ in both halves, and with control of the false discovery rate. For the latter, the replicability $P$-value of a pair, $2\{1-\Phi(\min(|z_A|,|z_B|))\}$ when the two halves agree in sign and 1 otherwise, is at most a level exactly when both halves pass that level with the same sign; it was controlled at 5\% by the Benjamini--Hochberg procedure\cite{benjamini1995} within each held-out sample and cell type (a unit; 22,200 pairs each), or over all \BoNUnits{} units together (threshold $P\leq{}$\CbGlobalThreshold).

\paragraph{Baselines and weighting.} Besides chance (50\%), a majority-sign baseline gave every reproducible association of a unit the sign held by most of that unit's associations; it is an oracle, since it knows each unit's majority from the held-out cells. The readout pooled units, weighting each by its number of reproducible associations, so that CD8 T cells, which held \CbShareCDEightPrimary\% of them, carried the most weight. The direction was therefore also averaged with every unit weighted alike, with every cell type weighted alike (units pooled within a type) and with every held-out sample weighted alike (units pooled within a sample); the nomination endpoint, which counts 100 pairs in every unit, weights units alike when pooled.

\paragraph{Intervals and multiplicity.} Differences between the atlas arm and paired-only estimation, whose envelope over the eight methods was recomputed in every resample, have 95\% intervals from the readout's 2,000 resamples of donors. Each of B1--B3 required a lower bound above zero at every one of its budgets, an intersection--union test\cite{berger1982}, which needs no adjustment within a hypothesis. For the eight planned comparisons (B1 and B2 at 50, 100 and 150 paired cells, B3 at 100 and 150), one-sided bootstrap $P$-values, one plus the number of resampled differences at most zero divided by one plus the number of resamples, and Bonferroni-simultaneous intervals at level $1-0.05/8$ were computed from \CbBoot{} resamples of donors; the $P$-value of each hypothesis, the largest over its budgets, was adjusted over B1--B3 by Holm's procedure\cite{holm1979}.

\paragraph{Results.} Positive associations outnumbered negative ones: \CbPosPrimary\% of the readout's reproducible associations were positive, and the majority-sign baseline gave the direction of \CbMajPrimaryPooled\% of them pooled and \CbMajPrimaryType\% with every cell type weighted alike, against \BoOneAtlasHundred\% for the atlas arm and \BoOnePoHundred\% for paired-only estimation with 100 paired cells (Supplementary Table~\ref{tab:calibration}). The other truth sets held \CbNTwentieth{} ($P<0.05$), \CbNThousandth{} ($P<0.001$), \CbNFdrUnit{} (5\% false discovery rate within units) and \CbNFdrAll{} (over all units) associations. Under each truth set and each weighting, the atlas arm gave the direction of more associations than paired-only estimation with 100 paired cells, by \CbDiffOneMin{} to \CbDiffOneMax{} percentage points with 95\% lower bounds of at least \CbDiffOneLoMin, and made more of its top pairs reproducible, by \CbDiffTwoMin{} to \CbDiffTwoMax{} points (lower bounds at least \CbDiffTwoLoMin). With every cell type weighted alike, the atlas arm gave the direction of \CbOnePrimaryTypeAtlasHundred\% of the associations against \CbOnePrimaryTypePoHundred\%. By cell type, the difference with 100 paired cells was \CbTypeOneCDEightDiffHundred{} points in CD8 T cells (95\% CI, \CbTypeOneCDEightDiffHundredLo{} to \CbTypeOneCDEightDiffHundredHi), \CbTypeOneCDFourDiffHundred{} in CD4 T cells (\CbTypeOneCDFourDiffHundredLo{} to \CbTypeOneCDFourDiffHundredHi), \CbTypeOneNKDiffHundred{} in NK cells (\CbTypeOneNKDiffHundredLo{} to \CbTypeOneNKDiffHundredHi), \CbTypeOneBDiffHundred{} in B cells (\CbTypeOneBDiffHundredLo{} to \CbTypeOneBDiffHundredHi) and \CbTypeOneMonoDiffHundred{} in CD14 monocytes (\CbTypeOneMonoDiffHundredLo{} to \CbTypeOneMonoDiffHundredHi). All eight planned comparisons had Bonferroni-simultaneous intervals above zero (Supplementary Table~\ref{tab:multiplicity}); no resampled difference was at most zero, so every one-sided bootstrap $P$-value was at the resolution of the resamples (largest \CbPlanMaxP), and the Holm-adjusted $P$-values of B1--B3 were all \CbHolmMax.

\begin{table}[h]
\caption{\textbf{Calibration of the readout with 100 paired cells.} Values in percent, differences in percentage points with 95\% intervals from the readout's 2,000 resamples of donors. Associations: reproducible associations under each truth set, pooled over the \BoNUnits{} units (held-out sample by cell type). Majority: the oracle majority-sign baseline for the direction, each unit's associations given its majority sign, aggregated with the same weighting. Atlas and paired: the atlas arm and paired-only estimation (the best of the eight paired-only methods for each endpoint, truth set and weighting). Truth sets: reproducible at the stated $P$ in both halves with the same sign, or at a false discovery rate (FDR) of 5\% on the replicability $P$-value, within each unit or over all units. Weighting: pooled over units (the readout's), or averaged with every unit, cell type or held-out sample weighted alike. Cell types: pooled within the type; associations of the type in the readout.}
\label{tab:calibration}
\centering\footnotesize
% Generated by extension/synthesis/make_assets.py. Do not edit.
\begin{tabular}{@{}lrrrrrrrr@{}}
\toprule
 & & & \multicolumn{3}{c}{Direction (E1, \%)} & \multicolumn{3}{c}{Top 100 reproducible (E2, \%)} \\
\cmidrule(lr){4-6}\cmidrule(l){7-9}
 & Associations & Majority & Atlas & Paired & Difference (95\% CI) & Atlas & Paired & Difference (95\% CI) \\
\midrule
\multicolumn{9}{@{}l}{\textit{Truth set (pooled)}} \\
$P<0.01$ (readout) & 120,649 & 57 & 86 & 75 & 11.2 (10.3--12.0) & 47 & 28 & 19.3 (18.0--20.5) \\
$P<0.05$ & 189,472 & 55 & 82 & 70 & 11.8 (11.0--12.4) & 52 & 31 & 21.0 (19.9--22.0) \\
$P<0.001$ & 79,414 & 58 & 89 & 79 & 10.2 (9.3--11.0) & 42 & 25 & 17.2 (15.8--18.4) \\
FDR 5\%, per unit & 83,863 & 57 & 89 & 78 & 10.9 (9.9--11.7) & 41 & 25 & 16.5 (15.2--17.8) \\
FDR 5\%, all units & 80,156 & 58 & 89 & 79 & 10.3 (9.3--11.1) & 42 & 25 & 17.2 (15.9--18.5) \\
\midrule
\multicolumn{9}{@{}l}{\textit{Weighting (readout's truth set)}} \\
Pooled (readout) & & 57 & 86 & 75 & 11.2 (10.3--12.0) & 47 & 28 & 19.3 (18.0--20.5) \\
Units alike & & 58 & 81 & 71 & 9.9 (9.3--10.4) & 47 & 28 & 19.3 (18.0--20.5) \\
Cell types alike & & 58 & 80 & 70 & 9.9 (9.3--10.4) & 47 & 28 & 19.3 (18.0--20.5) \\
Samples alike & & 57 & 86 & 75 & 11.1 (10.1--11.9) & 47 & 28 & 19.3 (18.0--20.5) \\
\midrule
\multicolumn{9}{@{}l}{\textit{Cell type (readout's truth set)}} \\
B & 5,077 & & 67 & 61 & 6.6 (5.5--7.7) & 3.5 & 2.3 & 1.1 (0.2--2.1) \\
CD4 T & 27,388 & & 92 & 77 & 15.2 (13.5--16.1) & 76 & 34 & 42.3 (39.0--45.3) \\
CD8 T & 56,637 & & 91 & 79 & 12.8 (11.6--13.9) & 85 & 74 & 10.7 (8.2--13.4) \\
NK & 10,526 & & 78 & 69 & 9.5 (8.4--10.8) & 30 & 18 & 12.2 (9.8--14.6) \\
CD14 Mono & 21,021 & & 73 & 72 & 1.3 (0.5--1.9) & 41 & 17 & 23.8 (21.7--26.0) \\
\bottomrule
\end{tabular}

\end{table}

\begin{table}[h]
\caption{\textbf{Planned comparisons of the readout: intervals and multiplicity.} Atlas arm minus paired-only estimation, in percentage points (direction and top 100 pairs) or points of recovered fraction (cognate pairs). 95\% CI: the readout's 2,000 resamples of donors. Simultaneous CI: Bonferroni over the eight comparisons (level $1-0.05/8$), \CbBoot{} resamples of donors. One-sided $P$: one plus the number of resampled differences at most zero, over one plus the number of resamples ($\leq$, no resampled difference at most zero). Each hypothesis's $P$-value is the largest over its budgets (intersection--union), adjusted over the three hypotheses by Holm's procedure.}
\label{tab:multiplicity}
\centering\footnotesize
% Generated by extension/synthesis/make_assets.py. Do not edit.
\begin{tabular}{@{}llrrrrr@{}}
\toprule
Hypothesis & Endpoint & Paired cells & Difference & 95\% CI & Simultaneous CI & One-sided $P$ \\
\midrule
B1 & Direction & 50 & 12.9 & 12.1 to 13.7 & 11.7 to 14.1 & $\leq5\times10^{-5}$ \\
B1 & Direction & 100 & 11.2 & 10.3 to 12.0 & 10.1 to 12.2 & $\leq5\times10^{-5}$ \\
B1 & Direction & 150 & 9.1 & 8.4 to 9.8 & 8.0 to 10.2 & $\leq5\times10^{-5}$ \\
B2 & Top 100 reproducible & 50 & 20.1 & 18.9 to 21.3 & 18.5 to 21.8 & $\leq5\times10^{-5}$ \\
B2 & Top 100 reproducible & 100 & 19.3 & 18.0 to 20.5 & 17.4 to 20.9 & $\leq5\times10^{-5}$ \\
B2 & Top 100 reproducible & 150 & 16.8 & 15.6 to 18.0 & 15.1 to 18.3 & $\leq5\times10^{-5}$ \\
B3 & Cognate pairs, recovered & 100 & 14.2 & 12.5 to 15.8 & 11.9 to 16.2 & $\leq5\times10^{-5}$ \\
B3 & Cognate pairs, recovered & 150 & 14.4 & 12.9 to 15.3 & 12.4 to 15.8 & $\leq5\times10^{-5}$ \\
\midrule
B1 & \multicolumn{5}{l}{Intersection--union $P$ (largest over its budgets); Holm-adjusted over B1--B3} & $\leq2\times10^{-4}$ \\
B2 & \multicolumn{5}{l}{Intersection--union $P$ (largest over its budgets); Holm-adjusted over B1--B3} & $\leq2\times10^{-4}$ \\
B3 & \multicolumn{5}{l}{Intersection--union $P$ (largest over its budgets); Holm-adjusted over B1--B3} & $\leq2\times10^{-4}$ \\
\bottomrule
\end{tabular}

\end{table}

\subsection*{Other methods given the same atlas (specified after scoring)}
The test showed that the atlas reduced the paired cells needed by this estimator. Whether established methods that also use unpaired data would exploit the atlas as well was specified in the same plan. SemiCCA and reference regression received the test's paired draws (folds, budgets, five draws and seeds), centred by the same contrasts, on the same genes and proteins, with the same atlas summaries; tuning used only the paired cells and the atlas. SemiCCA\cite{kimura2013} projected the paired cross-covariance on its canonical subspaces computed with the atlas's covariances, with its weight (0.2, 0.5 or 0.8) and rank (1, 2, 4, 8 or all proteins) chosen by two-fold cross-validation within the paired contrasts. Reference regression regressed protein on RNA by ridge regression in the paired contrasts, with the Gram matrix of the paired contrasts or of the atlas and the penalty chosen likewise, and predicted the cross-correlation as the atlas's RNA correlation times the coefficients. Each ran in the benchmark's three forms: pooled over cell types, mapped to each cell type through the atlas's condition maps (the estimator's maps), and refitted in each cell type with the atlas's summaries of that type. At each budget the best form, chosen on the held-out endpoint, represented each method, which favours the alternatives. K1 and K2 required 95\% lower bounds above one for the ratio of the paired cells that SemiCCA and reference regression needed to reach the test's primary target to those the estimator needed (2,000 resamples of donors, target fixed).

Champollion\cite{champollion} learns a cost between the two modalities from paired cells by inverse optimal transport and transports unpaired cells with it. It applies here as it did in development: it was fitted on the paired contrasts with the settings that its pilot had fixed for the external test on another data set (epsilon 1, 2,000 iterations, lasso weight by budget), and its transport plan between 2,000 RNA and 2,000 protein cells of each cell type of the atlas, centred on their population means, gave the prediction. This cap made its atlas sample smaller than that used for the covariance summaries of the other methods. Its fits took from seconds to minutes each on two cores, growing with the number of paired cells, so it was run on the first of the five draws of every fold and budget, and the estimator was scored on the same draws.

\paragraph{Results.} Every alternative gained from the atlas. The best forms of SemiCCA and reference regression reached the test's primary target with \AltCellsSemi{} and \AltCellsRidge{} paired cells, whereas paired-only estimation did not reach it within 800 (savings \AltRatioPoOverSemi{} and \AltRatioPoOverRidge; Supplementary Table~\ref{tab:alternatives}). The estimator needed \AltCellsAtlas, so SemiCCA needed \AltRatioSemi{} times as many paired cells (95\% CI, \AltRatioSemiLo--\AltRatioSemiHi; K1 \AltKOne) and reference regression \AltRatioRidge{} times as many (\AltRatioRidgeLo--\AltRatioRidgeHi; K2 \AltKTwo). Pooled over cell types, no alternative reached the target; their gains came through the per-cell-type and mapped forms, as the estimator's come through its condition maps. With 100 paired cells, reference regression gave the direction of \AltOneRidgeHundred\% of the readout's reproducible associations and made \AltTwoRidgeHundred\% of its top pairs reproducible, against \AltOneAtlasHundred\% and \AltTwoAtlasHundred\% for the estimator, and recovered \AltThreeRidgeHundred\% of the cognate correlations against \AltThreeAtlasHundred\%; SemiCCA gave \AltOneSemiHundred\%, \AltTwoSemiHundred\% and \AltThreeSemiHundred\%. On the first draws, Champollion recovered \AltRfChampHundred\% of the dependence with 100 paired cells and at most \AltRfMaxChamp\% with 800, against \AltRfChampAtlasHundred\% and \AltRfMaxChampAtlas\% for the estimator on the same draws, and it needed \AltCellsChamp{} paired cells to reach the target against \AltCellsChampAtlas{} (ratio \AltRatioChamp; \AltRatioChampLo--\AltRatioChampHi). With 100 paired cells it gave the direction of \AltOneChampHundred\% of the reproducible associations, against \AltOneChampAtlasHundred\% for the estimator on the same draws, and made \AltTwoChampHundred\% of its top pairs reproducible, against \AltTwoChampAtlasHundred\%.

\begin{table}[h]
\caption{\textbf{Other methods using the same paired draws and source atlas.} Validation test, pooled over the \VlNFolds{} held-out samples, five draws (last block: the first draw of every fold and budget). Champollion used 2,000 cells per assay and cell type from the atlas; the other methods used its full covariance summaries. Recovered: fraction of within-cell-type dependence recovered (\%). Paired cells: needed to reach the test's primary target (\VlTarget\%), with 95\% intervals from 2,000 resamples of donors; $>$800, not reached. Ratio: paired cells needed over those the estimator with the atlas needed (K1, SemiCCA; K2, reference regression; $\geq$, a lower bound). Readout, 100: direction (E1) and top 100 pairs (E2) with 100 paired cells. Ridge: reference regression with the Gram matrix of the paired contrasts or of the atlas. Forms: pooled over cell types; map, the pooled estimate mapped to each cell type through the atlas's condition maps; per cell type, refitted on each type's contrasts with the atlas's summaries of that type. Best form: the best at each budget, chosen on the held-out endpoint.}
\label{tab:alternatives}
\centering\footnotesize
% Generated by extension/synthesis/make_assets.py. Do not edit.
\begin{tabular}{@{}lrrrrrrrr@{}}
\toprule
 & \multicolumn{4}{c}{Recovered (\%), paired cells} & & & \multicolumn{2}{c}{Readout, 100 (\%)} \\
\cmidrule(lr){2-5}\cmidrule(l){8-9}
Arm & 50 & 100 & 200 & 800 & Paired cells (95\% CI) & Ratio (95\% CI) & E1 & E2 \\
\midrule
Estimator & 11 & 23 & 33 & 48 & 147 (125--185) & 1 & 86 & 47 \\
SemiCCA, pooled & 0.3 & 1.8 & 4.5 & 11 & $>$800 &  & 72 & 20 \\
SemiCCA, map & \ensuremath{-}0.2 & 3.6 & 13 & 32 & 579 (417--$>$800) &  & 73 & 22 \\
SemiCCA, per cell type & 11 & 17 & 20 & 28 & $>$800 &  & 76 & 32 \\
SemiCCA, best form & 11 & 17 & 20 & 32 & 496 (308--$>$800) & 3.4 (2.4--4.6) & 76 & 32 \\
Ridge, paired Gram, pooled & 5.2 & 8.6 & 12 & 18 & $>$800 &  & 80 & 28 \\
Ridge, paired Gram, map & 11 & 19 & 26 & 41 & 235 (203--275) &  & 85 & 43 \\
Ridge, paired Gram, per type & 14 & 22 & 28 & 43 & 197 (158--241) &  & 85 & 48 \\
Ridge, atlas Gram, pooled & 5.2 & 9.2 & 12 & 18 & $>$800 &  & 81 & 28 \\
Ridge, atlas Gram, map & 11 & 20 & 27 & 45 & 214 (186--248) &  & 85 & 44 \\
Ridge, atlas Gram, per type & 11 & 20 & 30 & 42 & 175 (136--414) &  & 86 & 50 \\
Reference regression, best form & 14 & 22 & 30 & 45 & 169 (131--236) & 1.15 (1.003--1.36) & 86 & 50 \\
Paired cells only & 2.4 & 3.5 & 7.7 & 13 & $>$800 & $\geq$5.4 (4.3--6.4) & 75 & 28 \\
\midrule
\multicolumn{9}{@{}l}{\textit{First draw of every fold only}} \\
Champollion & 11 & 20 & 30 & 40 & 180 (165--198) & 1.23 (0.98--1.44) & 88 & 42 \\
Estimator, same draws & 11 & 22 & 33 & 48 & 147 (123--189) & 1 & 87 & 47 \\
\bottomrule
\end{tabular}

\end{table}

\subsection*{Robustness of the atlas result (after scoring)}
\paragraph{Pools and donors.} From the recorded predictions, every processing pool gave a saving (Supplementary Table~\ref{tab:pools}), and with any one pool left out and the primary target recomputed, the saving was at least \RbLopoMin{} to \RbLopoMax; without any one pool, the atlas needed \RbLopoOwnMin{} to \RbLopoOwnMax{} times the paired cells needed with the own samples' unpaired cells, so the result of V4 does not rest on one pool. Every donor's held-out samples gave a saving, from at least \RbDonorMin{} to \RbDonorMax{} (Supplementary Table~\ref{tab:donors}); paired-only estimation reached the target in the samples of \RbDonorsPoReached{} donors, and the atlas arm needed at most \RbDonorAtlasMax{} paired cells in any donor's samples.

\paragraph{Size and composition of the atlas.} The atlas arm was recomputed with the reference changed and every other setting of the test unchanged (panel, folds, draws and seeds), and scored against the prespecified paired-only arm at the test's primary target; recomputed with the unchanged atlas, it reproduced the recorded atlas predictions exactly (largest difference \RbReproduces). Patients were taken in a fixed pseudo-random order. Supplementary Table~\ref{tab:robust} gives every variant. With the cells of 59, 30 and 15 patients the atlas arm needed \RbCellsPatientsFiftyNine, \RbCellsPatientsThirty{} and \RbCellsPatientsFifteen{} paired cells (95\% CI, \RbCellsPatientsFiftyNineLo--\RbCellsPatientsFiftyNineHi, \RbCellsPatientsThirtyLo--\RbCellsPatientsThirtyHi{} and \RbCellsPatientsFifteenLo--\RbCellsPatientsFifteenHi), with 7 patients \RbCellsPatientsSeven, and with 4 patients (\RbCellCountPatientsFour{} cells, \RbPopsPatientsFour{} populations) it recovered at most \RbRfMaxPatientsFour\% and did not reach the target. With one cell type left out, it needed \RbCellsWithoutB{} paired cells without B cells, \RbCellsWithoutCDFourT{} without CD4 T cells, \RbCellsWithoutNK{} without NK cells and \RbCellsWithoutMono{} without CD14 monocytes, whereas without CD8 T cells, which held \BoRepCDEight{} of the \BoNRep{} reproducible associations of the readout, it did not reach the target. With B cells and CD14 monocytes from 15 patients and the other types from all 118 (\RbCellCountFewBMono{} cells) it recovered at most \RbMappedMaxFewBMono\%, and with CD4 and CD8 T cells from 15 patients (\RbCellCountFewT{} cells) at most \RbMappedMaxFewT\%, falling to \RbMappedFewT\% with 800 paired cells.

To locate this failure, the full atlas and the two skewed atlases were scored without condition maps, every cell type receiving the pooled estimate along the atlas's axes. With 800 paired cells, these estimates recovered \RbUnmapFull\%, \RbUnmapFewBMono\% and \RbUnmapFewT\%, whereas with maps the full atlas recovered \RbMappedFull\%. The axes therefore carried over under the skewed compositions, and the condition maps, which convert the pooled estimate into each cell type's through the atlas's pooled covariance, failed when the atlas's composition departed from that of the paired cells; a cell type absent from the atlas escapes this, because its cells keep the pooled estimate.

\begin{table}[h]
\caption{\textbf{The atlas arm with the reference changed, and by processing pool.} Atlas arm recomputed with every setting of the test unchanged (panel, folds, draws, seeds) except the reference; patients chosen in a fixed pseudo-random order. A cell type absent from the reference receives the pooled estimate without a condition map. With 100, recovered fraction with 100 paired cells. Cells to target: paired cells needed to reach the test's primary target (\VlTarget\%), with 95\% intervals from 2,000 resamples of donors; saving over the prespecified paired-only arm, which did not reach the target within 800 paired cells, so every saving is a lower bound. With the target held fixed, the interval of the unchanged atlas is wider than that of hypothesis V1, whose bootstrap recomputed the target in every resample.}
\label{tab:robust}
\centering\footnotesize
% Generated by extension/synthesis/make_assets.py. Do not edit.
\begin{tabular}{@{}lrrrrr@{}}
\toprule
Atlas reference & Patients & Cells & With 100 (\%) & Cells to target (95\% CI) & Saving (95\% CI) \\
\midrule
All 118 patients & 118 & 63,796 & 23 & 147 (126--182) & $\geq$5.4 (4.4--6.3) \\
59 patients & 59 & 32,350 & 20 & 174 (147--217) & $\geq$4.6 (3.7--5.4) \\
30 patients & 30 & 16,510 & 20 & 177 (151--219) & $\geq$4.5 (3.6--5.3) \\
15 patients & 15 & 8,231 & 20 & 185 (155--246) & $\geq$4.3 (3.3--5.2) \\
7 patients & 7 & 3,823 & 16 & 385 (281--$>$800) & $\geq$2.1 (1.00--2.8) \\
4 patients & 4 & 2,132 & 9.5 & $>$800 & neither reached \\
B, CD14 Mono from 15 patients & 118 & 47,683 & 2.8 & $>$800 & neither reached \\
CD4, CD8 T from 15 patients & 118 & 34,942 & 14 & $>$800 & neither reached \\
Without B cells & 118 & 55,371 & 21 & 174 (147--222) & $\geq$4.6 (3.6--5.5) \\
Without CD4 T & 118 & 42,339 & 23 & 159 (132--201) & $\geq$5.0 (4.0--6.1) \\
Without CD8 T & 118 & 52,250 & 9.8 & $>$800 & neither reached \\
Without NK & 118 & 51,730 & 21 & 184 (152--241) & $\geq$4.4 (3.3--5.3) \\
Without CD14 Mono & 118 & 53,494 & 20 & 209 (162--316) & $\geq$3.8 (2.5--4.9) \\
\bottomrule
\end{tabular}

\end{table}

\begin{table}[h]
\caption{\textbf{The validation test by processing pool.} From the recorded predictions. In each pool: paired cells the atlas arm needed to reach the test's primary target, and its saving over paired-only estimation ($\geq$, paired-only estimation did not reach the target within 800 paired cells). Without each pool in turn: the primary target recomputed as half of the best recovery, the saving, and the paired cells the atlas needed divided by those needed with the own samples' unpaired cells (the ratio of hypothesis V4).}
\label{tab:pools}
\centering\footnotesize
% Generated by extension/synthesis/make_assets.py. Do not edit.
\begin{tabular}{@{}lrrrrrr@{}}
\toprule
 & & \multicolumn{2}{c}{In the pool} & \multicolumn{3}{c}{Without the pool} \\
\cmidrule(lr){3-4}\cmidrule(l){5-7}
Pool & Samples & Atlas cells & Saving & Target (\%) & Saving & Atlas / own \\
\midrule
DB8 & 3 & 95 & $\geq$8.4 & 28 & $\geq$5.2 & 1.30 \\
DB9 & 6 & 129 & $\geq$6.2 & 28 & $\geq$5.4 & 1.24 \\
DB10 & 4 & 238 & $\geq$3.4 & 29 & $\geq$5.5 & 1.36 \\
DB11 & 6 & 184 & $\geq$4.3 & 29 & $\geq$5.6 & 1.42 \\
DB12 & 5 & 186 & $\geq$4.3 & 28 & $\geq$5.6 & 1.36 \\
DB13 & 4 & 129 & $\geq$6.2 & 28 & $\geq$5.4 & 1.35 \\
DB14 & 6 & 132 & $\geq$6.0 & 28 & $\geq$5.3 & 1.39 \\
\bottomrule
\end{tabular}

\end{table}

\begin{table}[h]
\caption{\textbf{The validation test by donor.} From the recorded predictions, each donor's held-out samples (two for donors sampled at two time points): recovered fraction with 100 paired cells with the atlas and with paired cells only, paired cells the atlas arm needed to reach the test's primary target, and its saving over paired-only estimation ($\geq$, paired-only estimation did not reach the target within 800 paired cells).}
\label{tab:donors}
\centering\footnotesize
% Generated by extension/synthesis/make_assets.py. Do not edit.
\begin{tabular}{@{}lrrrrr@{}}
\toprule
 & & \multicolumn{2}{c}{Recovered with 100 (\%)} & & \\
\cmidrule(lr){3-4}
Donor & Samples & Atlas & Paired only & Atlas cells & Saving \\
\midrule
HS6 & 2 & 30 & 12 & 89 & 3.1 \\
HS10 & 2 & 26 & 4.0 & 117 & $\geq$6.8 \\
HS12 & 1 & 29 & 9.0 & 98 & $\geq$8.1 \\
HS16 & 2 & 24 & 6.1 & 123 & $\geq$6.5 \\
HS22 & 2 & 22 & 2.8 & 164 & $\geq$4.9 \\
HS24 & 2 & 16 & 2.5 & 741 & $\geq$1.08 \\
HS28 & 2 & 22 & 5.6 & 149 & $\geq$5.4 \\
HS30 & 2 & 24 & 3.9 & 126 & $\geq$6.4 \\
HS32 & 2 & 27 & 10 & 106 & 5.0 \\
HS46 & 2 & 23 & 3.2 & 129 & $\geq$6.2 \\
HS47 & 2 & 13 & 3.5 & 292 & $\geq$2.7 \\
HS48 & 1 & 23 & 5.7 & 200 & $\geq$4.0 \\
HS68 & 2 & 10 & 0.7 & 436 & $\geq$1.8 \\
HS75 & 2 & 18 & 2.4 & 236 & $\geq$3.4 \\
HS86 & 2 & 20 & 6.6 & 161 & $\geq$5.0 \\
HS88 & 2 & 27 & 8.4 & 108 & $\geq$7.4 \\
HS93 & 2 & 23 & 5.2 & 155 & 4.7 \\
HS94 & 2 & 22 & 6.4 & 250 & $\geq$3.2 \\
\bottomrule
\end{tabular}

\end{table}

\section{The estimator: theory, development and external test}\label{sec:semi}

This note defines the estimator, proves a risk bound for an idealized version of it, the condition map under a shared channel and a matching lower bound in an idealized Gaussian-sequence model (Propositions~\ref{prop:bjs}--\ref{prop:minimax}), and gives the development comparisons and the protocol and complete results of the prespecified external test. Supplementary Note~\ref{sec:pred} bounds the implemented rule itself and derives the saving.

\subsection*{Setting and estimator}
Training cells of eligible populations are divided by a fixed pseudo-random rule into a paired calibration reservoir (25\%) and an unpaired pool (75\%). Pool cells contribute RNA from one assay half and protein from the other, never both, and give three unpaired summaries: the pooled within-population RNA correlation, its latent version after binomial count splitting (Supplementary Note~\ref{sec:noise}), and the pooled within-population protein correlation, each also computed within every condition. From the reservoir, $B$ paired cells are drawn; each is centred on the unpaired means of its population and scaled by pooled within-population standard deviations, giving RNA vectors $x_i\in\mathbb{R}^p$ and protein vectors $y_i\in\mathbb{R}^q$ and per-cell products $z_i=x_iy_i^\top$. Their mean $\hat\theta$ estimates the pooled within-population cross-covariance $\theta$.

Let $U$ be the eigenvectors of the latent RNA correlation and $V$ those of the protein correlation, both ordered by eigenvalue and both from unpaired cells. The proposed estimator writes $\hat\theta$ in these bases, $A=U^\top\hat\theta V$, partitions $A$ into blocks formed by RNA eigen-coordinates 1--5, 6--20, 21--60 and the rest and protein eigen-coordinates 1--3 and the rest, and multiplies each block $A_b$ by the positive-part James--Stein factor\cite{james1961,baranchik1970} $(1-\hat t_b/\|A_b\|^2)_+$, where $\hat t_b$ is the trace of the covariance of the block's mean, estimated from the per-cell products. The shrunk coefficients are rotated back, $\tilde\theta=U\tilde AV^\top$. For condition $c$, the estimate is mapped by $G_c=\Sigma_{c}U_k\Lambda_k^{-1}U_k^\top+(I-U_kU_k^\top)$ on the leading $k=120$ latent eigendirections, where $U_k$ ($p\times k$) and $\Lambda_k$ ($k\times k$) are the leading eigenvectors and eigenvalues of the pooled latent RNA correlation and $\Sigma_c=D_{x,c}R_{z,c}D_{x,c}$ combines the condition's latent RNA correlation $R_{z,c}$ with the ratios $D_{x,c}$ of its measured RNA standard deviations to the pooled ones (Proposition~\ref{prop:map} motivates the leading term). The map is shrunk towards the identity by the split-half reliability $s_c$ of the condition-minus-pooled covariance (unpaired cells only), $\hat G_c=I+s_c(G_c-I)$, and rows and columns are rescaled by the likewise shrunk ratios, $\hat D_{x,c}=I+s_c(D_{x,c}-I)$ and $\hat D_{y,c}$, so that the mapped estimate is $\hat D_{x,c}^{-1}\hat G_c\tilde\theta\hat D_{y,c}^{-1}$. The estimator needs one eigendecomposition per assay; its blocks, shrinkage rule and map rank were fixed in development, and nothing is tuned to a data set.

\subsection*{Risk bound}
\begin{proposition}[Block James--Stein with correlated, unequal noise]\label{prop:bjs}
Let $U$ and $V$ be fixed orthonormal bases and let the coefficients of block $b$ satisfy $\hat a_b\sim N(a_b,S_b)$ for any positive semidefinite $S_b$. Write $t_b=\operatorname{tr}S_b$, $\ell_b=\lambda_{\max}(S_b)$ and $\delta_b=(1-t_b/\|\hat a_b\|^2)\hat a_b$, with $\delta_b=0$ at $\hat a_b=0$. If $t_b\ge4\ell_b$, then
\[
\mathbb{E}\|\delta_b-a_b\|^2\le\frac{\|a_b\|^2t_b}{\|a_b\|^2+t_b}+4\ell_b\le\min(\|a_b\|^2,t_b)+4\ell_b .
\]
The fraction is defined as zero when $\|a_b\|^2=t_b=0$.
\end{proposition}
\noindent\textit{Proof.} If $t=0$, then $X=a$ almost surely and the bound is immediate. Otherwise, with $X=\hat a_b$ and $g(X)=X/\|X\|^2$, Stein's identity for a Gaussian vector with covariance $S$\cite{stein1981} gives $\mathbb{E}\langle X-a,g(X)\rangle=\mathbb{E}\operatorname{tr}(S\,Dg(X))$ and $\operatorname{tr}(S\,Dg(X))=\operatorname{tr}S/\|X\|^2-2X^\top SX/\|X\|^4$. Hence, with $c=t$,
\[
\mathbb{E}\|\delta-a\|^2=t-2t\,\mathbb{E}\Big[\frac{t}{\|X\|^2}-\frac{2X^\top SX}{\|X\|^4}\Big]+t^2\,\mathbb{E}\frac{1}{\|X\|^2}\le t-t(t-4\ell)\,\mathbb{E}\frac{1}{\|X\|^2},
\]
using $X^\top SX\le\ell\|X\|^2$. For $t\ge4\ell$ the bracket is nonnegative, and Jensen's inequality, $\mathbb{E}\|X\|^{-2}\ge(\|a\|^2+t)^{-1}$, gives $\mathbb{E}\|\delta-a\|^2\le t(\|a\|^2+4\ell)/(\|a\|^2+t)$, which is at most the stated bound. $\square$

Since $M\mapsto U^\top MV$ is an isometry in the Frobenius norm, the risk of $\tilde\theta$ is the sum of the block risks. With $S_b=\Sigma_b/B$, $\tau_b=\operatorname{tr}\Sigma_b$ and $\rho_b=\lambda_{\max}(\Sigma_b)$, for any set $\mathcal{A}$ of blocks,
\[
R(B)\le\sum_{b\notin\mathcal{A}}\|a_b\|^2+\frac{1}{B}\Big(\sum_{b\in\mathcal{A}}\tau_b+4\sum_b\rho_b\Big),
\]
so, in the idealized setting and for known block norms $\|a_b\|^2$, a budget $B\ge(\sum_{b\in\mathcal{A}}\tau_b+4\sum_b\rho_b)/(e-\sum_{b\notin\mathcal{A}}\|a_b\|^2)$ guarantees risk at most $e$ when its denominator is positive. For $\|\theta\|>0$, the expected recovered fraction of the pooled target is then at least $1-e/\|\theta\|^2$. The block norms are properties of the dependence that only paired cells reveal (Supplementary Note~\ref{sec:pred}). The leading budget term depends on the noise of the signal-bearing blocks, with the remainder $4\sum_b\rho_b$ accounting for all blocks. A single block of all $pq$ coefficients instead retains the noise of every coordinate. Bases estimated from the paired cells are not independent of their noise, so the bound does not apply to them, and random blocks do not concentrate the dependence. The condition $\tau_b\ge4\rho_b$ does not depend on $B$; the smallest median effective dimension $\hat t_b\hat\ell_b^{-1}$ over blocks, budgets and screens was \SpEffDimMin. With spherical noise, truncating the factor at zero always lowers the risk\cite{baranchik1970}. Related results on shrinkage under general quadratic loss and non-spherical noise are refs.~\citenum{bock1975,matsuda2024}.

\paragraph{Scope and numerical check.} The implemented bases and blocks meet the fixed-basis conditions of Proposition~\ref{prop:bjs}: $U$ and $V$ come from pool cells disjoint from the reservoir. Its Gaussian and known-noise assumptions do not describe the implementation. The block means average non-Gaussian per-cell products, the trace $t_b$ is estimated from those products and the factor is truncated at zero. Proposition~\ref{prop:oracle} (Supplementary Note~\ref{sec:pred}) bounds the implemented rule for independent product vectors without Gaussian assumptions. Centring on estimated unpaired means adds the outer product of the RNA and protein mean errors to the target cross-covariance; each mean error is of order $n^{-1/2}$ for $n$ independent unpaired cells from a population with finite covariance. In a numerical comparison with the idealized bound, each reservoir was treated as the population (OverCITE-seq, \SpBoundCellsOc{} cells; Frangieh, \SpBoundCellsFr; Papalexi, \SpBoundCellsPa), so that $\theta$, $\tau_b$ and $\rho_b$ were its exact moments, with bases from the pool as in the estimator and \SpBoundDraws{} draws with replacement at each budget from 25 to \SpBoundBudgetMaxFr{} cells (OverCITE-seq, to \SpBoundBudgetMaxOc). At every budget in all three data sets, the mean squared error of the implemented estimator was below the idealized bound, by factors of \SpBoundRatioMin{} to \SpBoundRatioMax, and below that of the idealized rule, and every block had $\tau_b\ge4\rho_b$ (smallest ratio, \SpBoundEffDimMin), apart from the Papalexi blocks formed with the protein direction that centred log-ratios leave without variance, whose signal and noise are zero. This numerical comparison does not extend Proposition~\ref{prop:bjs} to the implementation. The savings over other methods are measured empirically; Proposition~\ref{prop:minimax} below gives a matching lower bound in a Gaussian-sequence model.


\subsection*{Condition map}
\begin{proposition}[Condition map under a shared channel]\label{prop:map}
Let latent RNA $z$, measured RNA $x=z+\varepsilon$ and protein $y=Wz+e$ have $W$ shared by all conditions. Within each condition, let $\varepsilon$ be independent of $(z,y)$ and $e$ independent of $z$. Write $\operatorname{Cov}_c(z)=\Sigma_c$ and let $\Sigma=\sum_c w_c\Sigma_c$ be the positive-definite pooled within-condition covariance, with fixed weights $w_c>0$ summing to one. Then $\operatorname{Cov}_c(x,y)=\Sigma_cW^\top$, the corresponding pooled cross-covariance is $C=\sum_c w_c\operatorname{Cov}_c(x,y)=\Sigma W^\top$, and $\operatorname{Cov}_c(x,y)=\Sigma_c\Sigma^{-1}C$.
\end{proposition}
\noindent\textit{Proof.} $\operatorname{Cov}_c(x,y)=\operatorname{Cov}_c(z,Wz+e)=\Sigma_cW^\top$, and likewise for the pooled covariance; eliminating $W^\top=\Sigma^{-1}C$ gives the map. $\square$

Measured covariances would give $\Sigma_{x,c}\Sigma_x^{-1}$, with measurement noise on both diagonals, which is why the map uses latent covariances.

\paragraph{Error of the mapped estimate.} Given the unpaired cells, the implemented map is a fixed linear operator on $p\times q$ matrices, $\mathcal{G}_c(M)=\hat D_{x,c}^{-1}\hat G_cM\hat D_{y,c}^{-1}$, with $\hat D_{x,c}^{-1}\hat G_c$ of size $p\times p$ and $\hat D_{y,c}^{-1}$ of size $q\times q$, and it does not depend on the paired cells. For the target $\theta_c$ of condition $c$, Minkowski's inequality and $\|AMB\|\le\|A\|_{\mathrm{op}}\|M\|\|B\|_{\mathrm{op}}$ give
\[
\big(\mathbb{E}\|\mathcal{G}_c(\tilde\theta)-\theta_c\|^2\big)^{1/2}\le\|\hat D_{x,c}^{-1}\hat G_c\|_{\mathrm{op}}\,\|\hat D_{y,c}^{-1}\|_{\mathrm{op}}\,R(B)^{1/2}+\|\mathcal{G}_c(\theta)-\theta_c\|,
\]
where $R(B)=\mathbb{E}\|\tilde\theta-\theta\|^2$ is the risk of the pooled estimate. The first term carries whatever bound holds for $R(B)$. The second is the error of the map itself; it does not depend on $B$, and the inequality does not control it. It would vanish for the exact map of Proposition~\ref{prop:map}, $\Sigma_{z,c}\Sigma_z^{-1}$ with exact latent covariances and exact scaling, under that proposition's assumptions. The implemented map is not that map: it inverts only $k=120$ leading directions and is the identity on the rest, it is shrunk towards the identity by $s_c$, and it rescales latent correlations by ratios of measured standard deviations. Even with its inputs known exactly ($s_c=1$) and unit scaling, under Proposition~\ref{prop:map} its error is $(G_c-\Sigma_{z,c}\Sigma_z^{-1})\theta=(\Sigma_z-\Sigma_{z,c})P_{-k}W^\top$, with $P_{-k}$ the projection on the pooled eigendirections beyond the leading $k$. It is nonzero whenever the condition changes the latent covariance on directions beyond the leading $k$ that the channel $W$ uses, and further error arises when latent and measured standard deviations change differently between conditions. The development comparisons measure the net effect of the map (it added \SpFrMapGain{} percentage points at 100 paired cells in Frangieh).

\subsection*{Near-optimality in an idealized model}
\begin{proposition}[Near-optimality over dependence with given block norms, Gaussian-sequence model]\label{prop:minimax}
In the setting of Proposition~\ref{prop:bjs}, let the noise of the block means be spherical and independent across blocks, $S_b=(\tau_b/(d_bB))I_{d_b}$ for a block of $d_b\ge4$ coefficients, and let $\mathcal{P}(r)$ be the priors on $\theta$ with $\mathbb{E}\|a_b\|^2\le r_b^2$ in every block. Then
\[
\inf_{\hat\theta}\,\sup_{\pi\in\mathcal{P}(r)}\mathbb{E}_\pi\mathbb{E}\|\hat\theta-\theta\|^2=\sum_b\frac{r_b^2\,\tau_b/B}{r_b^2+\tau_b/B}=:R^*(B),
\]
with a summand defined as zero when $r_b^2=\tau_b=0$, and the infimum taken over all estimators in this Gaussian-sequence experiment. Block James--Stein, which does not know $r$, attains $\sup_{\pi\in\mathcal{P}(r)}\mathbb{E}_\pi\mathbb{E}\|\tilde\theta-\theta\|^2\le R^*(B)+4\sum_b\tau_b/(d_bB)$.
\end{proposition}
\noindent\textit{Proof.} Blocks with $r_b^2=\tau_b=0$ have zero signal and noise almost surely; set $c_b=0$ and omit them from the following calculation. Lower bound: under the prior $\pi^*$ with independent $a_b\sim N(0,(r_b^2/d_b)I)$, which belongs to $\mathcal{P}(r)$, no estimator has smaller Bayes risk than the posterior mean, whose risk is $\sum_bd_b\,v_b\sigma_b^2/(v_b+\sigma_b^2)=R^*(B)$ with $v_b=r_b^2/d_b$ and $\sigma_b^2=\tau_b/(d_bB)$. Attainment: the linear rule $c_bX_b$ with $c_b=r_b^2/(r_b^2+\tau_b/B)$ has risk $(1-c_b)^2\mathbb{E}\|a_b\|^2+c_b^2\tau_b/B\le R^*$ block by block under every prior in $\mathcal{P}(r)$, so the infimum equals $R^*(B)$. James--Stein: here $t_b=\tau_b/B$ and $\ell_b=\tau_b/(d_bB)$, so $t_b\ge4\ell_b$, and Proposition~\ref{prop:bjs} bounds its risk at every $\theta$ by $\sum_b\|a_b\|^2t_b/(\|a_b\|^2+t_b)+4\ell_b$; the first term is concave and increasing in $\|a_b\|^2$, so by Jensen's inequality its expectation under any prior in $\mathcal{P}(r)$ is at most $R^*(B)$. $\square$

Two consequences follow in this idealized model. First, over dependence matrices with given block norms, the paired cells that any estimator needs to guarantee a risk $e$ are at least the $B$ solving $R^*(B)=e$, and block James--Stein needs at most the $B$ solving $R^*(B)+4\sum_b\tau_b/(d_bB)=e$. Given the block norms, the block structure of the unpaired bases therefore fixes the paired budget up to a term that shrinks with block size; the block norms themselves are not revealed by unpaired cells (Supplementary Note~\ref{sec:pred}). Second, if the same dependence is described by a single block, the minimal budget to halve the risk of estimating zero is $\tau/A$, with $A=\sum_b\|a_b\|^2$ and $\tau=\sum_b\tau_b$. With the block structure, it is the $B$ solving $R^*(B)=A/2$. Their ratio is the saving predicted by the structural law tested in Supplementary Note~\ref{sec:gen}. The model has spherical noise within blocks, which per-cell products approximate, and the comparison is among estimators given the unpaired bases and the block norms; estimators that exploit other structure, such as low rank across blocks, lie outside it.

\subsection*{Development comparison}
Both screens had been opened by the analyses of Supplementary Notes~\ref{sec:pert} and~\ref{sec:rep} and served as development data. Evaluation groups and endpoint were those of those notes: the noise-unbiased recovered fraction of within-group cross-correlation in held-out target groups (Frangieh, \SpFrGroups{} groups from \SpFrTargets{} targets; Papalexi, \SpPaGroups{} groups, each target held out in turn), and the uncorrected relative squared loss. Budgets were 50 to 3,200 paired cells (Frangieh, 10 draws) and 50 to 1,600 (Papalexi, 4 draws per fold); every arm used the same paired cells in a draw and the same unpaired pool. Comparators, each pooled, refitted per condition and with the proposed map, were James--Stein shrinkage in marker coordinates, SCOSE and FCOSE\cite{muandet2014,ramdas2015}, cross-validated low rank, SemiCCA\cite{kimura2013} (the paired cross-covariance projected on SemiCCA subspaces fitted to paired and unpaired cells, trade-off and rank by cross-validation) and ridge regression of protein on RNA with a paired or unpaired Gram matrix. Champollion\cite{champollion} was fitted on the paired cells with its default entropic regularization and 2,000 iterations, with lasso weight $c/\sqrt{B}$ and $c$ chosen for each budget from $\{1, 1.5, 2.25\}$ on the evaluation groups themselves, a choice that favours it; it transported between the unpaired RNA and protein cells of each condition's pool (up to 2,000 of each), and the cross-correlation of the transport plan was its estimate (budgets 50 to 800; Frangieh, 3 draws; Papalexi, \SpPaChampFolds{} folds, 2 draws). The paired cells needed to reach a target were read from the running maximum of each budget curve by log-linear interpolation between the two budgets that bracket the target, so they are interpolated values rather than budgets that were run. Intervals came from 1,000 bootstrap resamples of held-out targets (Frangieh, \SpFrTargets; Papalexi, \SpPaTargets), the independent units of these intervals. Draws of paired cells were averaged within each budget and are not replicates. Each screen comes from one cell source, patient-derived melanoma cells in Frangieh and the THP-1 monocytic line in Papalexi, so neither replicates across donors; Supplementary Note~\ref{sec:gen} holds out donors, mice and neuron types.

The two-sided eigenbasis was chosen during development: relative to blocks holding all protein coordinates, splitting the protein coordinates by the unpaired protein eigenbasis raised recovery at 100 paired cells from \SpFrOneSidedHundred\% to \SpFrRfHundred\% in Frangieh and left Papalexi unchanged by construction, because the centred log-ratios of its four proteins span three dimensions, which form a single protein block. Other variants examined in development and not adopted were one block per protein, a lasso regression with the unpaired Gram matrix, a Stein-unbiased choice of a smooth spectral profile, and an average of block shrinkage and ridge regression.

\subsection*{External test}
The protocol was registered before any statistic of held-out cells other than cell counts had been computed. OverCITE-seq\cite{legut2022} profiles primary human CD8$^+$ T cells, each transduced with one open reading frame (ORF) and cultured with IL-2 alone or with CD3/CD28 activation, for RNA and 14 surface proteins. Every third ORF in a fixed pseudo-random order was held out with all its cells (\SpExOrfs{} ORFs, \SpExCells{} cells). Features were the genes encoding the proteins targeted by the antibodies and the most variable genes of the training cells (200 genes, ORF transgenes excluded). Arms, budgets (25 to 400 paired cells, 20 draws) and comparators were those of the development comparison, with Champollion's lasso weight chosen per budget on the training ORFs alone. The endpoint was the recovered fraction of within-condition cross-correlation among held-out ORFs, each cell centred on its ORF-by-condition mean and pooled within condition, because held-out groups were small. The accuracy target, a recovered fraction of 0.3, was fixed from a pilot on training ORFs. Hypotheses compared the paired cells the proposed estimator needed with those of paired-only estimation (H1, primary) and SemiCCA (H2), each requiring a 95\% bootstrap lower bound of the ratio above one, and required non-inferiority to Champollion (H3: lower bound above 0.8, that is, at most 1.25 times Champollion's paired cells), with 2,000 resamples of held-out ORFs; comparisons with every comparator in its original form or with the map were reported without a claim, and a comparison was decided only if the proposed curve reached the target within 400 cells. As in development, the paired cells a method needed were interpolated log-linearly on the running maximum of its budget curve between the two budgets bracketing the target. The \SpExOrfs{} held-out ORFs are the independent units of the bootstrap; the 20 draws of paired cells at each budget were averaged and are not replicates. The deposited data do not identify donors, so this test replicates across perturbations, not across donors.

\subsection*{Development results}
Supplementary Table~\ref{tab:semidev} gives recovered fractions at every budget. In Frangieh, the proposed estimator exceeded the best comparator at every budget, including comparators with the proposed map (difference at 100 paired cells, \SpFrBestDiffHundred{} percentage points; 95\% CI, \SpFrBestDiffHundredLo{} to \SpFrBestDiffHundredHi), and reached recovered fractions of 0.3 to 0.6 with \SpFrSavePairedOnlyMin{} to \SpFrSavePairedOnlyMax{} times fewer paired cells than paired-only estimation (to 0.4: \SpFrCellsFour{} against \SpFrPairedOnlyCellsFour{} cells; factor \SpFrSavePairedOnlyFour, \SpFrSavePairedOnlyFourLo--\SpFrSavePairedOnlyFourHi), \SpFrSavePublishedMin{} to \SpFrSavePublishedMax{} times fewer than the other comparators in their original forms and \SpFrSaveBestMin{} to \SpFrSaveBestMax{} times fewer than any comparator with the map. In Papalexi, it needed \SpPaSavePairedOnlyMin{} to \SpPaSavePairedOnlyMax{} times fewer paired cells than paired-only estimation (to 0.4: factor \SpPaSavePairedOnlyFour, \SpPaSavePairedOnlyFourLo--\SpPaSavePairedOnlyFourHi) and \SpPaSavePublishedMin{} to \SpPaSavePublishedMax{} times fewer than the other comparators in their original forms, with intervals that included one at some targets; reference regression given the proposed map matched the proposed estimator from 200 to 800 paired cells and exceeded it at 1,600, and at 100 paired cells the proposed estimator was ahead by \SpPaBestDiffHundred{} percentage points (95\% CI, \SpPaBestDiffHundredLo{} to \SpPaBestDiffHundredHi). \SpDevChampText{} Its uncorrected relative loss was the lowest of all arms at every budget in Frangieh. All training cells, used as paired cells, recovered \SpFrPairedAll\% (Frangieh) and \SpPaPairedAll\% (Papalexi); the remainder is dependence specific to the held-out targets.

\begin{table}[h]
\caption{\textbf{Recovered fraction (\%) of within-group cross-correlation in development, by number of paired cells.} Proposed, two-sided block shrinkage with the condition map; ablations change one component; comparator rows give the best of the listed fits at each budget (paired only: James--Stein, SCOSE, FCOSE and low rank, pooled or per condition; SemiCCA and reference regression: pooled or per condition); $^a$Lasso weight chosen per budget on the evaluation groups; Frangieh, 2 draws; Papalexi, \SpPaChampFolds{} of the held-out targets (\SpPaChampGroups{} groups), 1 draw. Other rows: Frangieh, \SpFrGroups{} groups, 10 draws; Papalexi, \SpPaGroups{} groups, 4 draws.}
\label{tab:semidev}
\centering\footnotesize
% Generated by extension/semipaired/make_assets.py. Do not edit.
\begin{tabular}{@{}lrrrrrrr@{}}
\toprule
Paired cells & 50 & 100 & 200 & 400 & 800 & 1,600 & 3,200 \\
\midrule
\multicolumn{8}{@{}l}{\textit{Frangieh (melanoma)}} \\
Proposed (two-sided, mapped) & 37 & 48 & 63 & 72 & 79 & 84 & 87 \\
Two-sided, pooled & 35 & 45 & 57 & 65 & 71 & 75 & 77 \\
Two-sided, per condition & \ensuremath{-}9.0 & 26 & 45 & 62 & 73 & 80 & 85 \\
One-sided (RNA eigenbasis), mapped & 27 & 39 & 57 & 69 & 77 & 83 & 87 \\
Bases from the paired cells & \ensuremath{-}8.1 & 11 & 35 & 53 & 65 & 72 & 76 \\
Marker coordinates, random blocks & 7.1 & 15 & 28 & 43 & 56 & 66 & 73 \\
Map from measured RNA covariance & 30 & 43 & 61 & 70 & 76 & 80 & 83 \\
500 unpaired cells per assay & 32 & 42 & 54 & 64 & 70 & 75 & 78 \\
2,000 unpaired cells per assay & 34 & 45 & 59 & 68 & 74 & 79 & 82 \\
Paired only (best of four) & 9.1 & 16 & 28 & 43 & 58 & 67 & 74 \\
SemiCCA & 16 & 31 & 45 & 53 & 66 & 70 & 79 \\
Reference regression & 14 & 27 & 43 & 57 & 66 & 72 & 82 \\
Any comparator with the map & 17 & 32 & 47 & 62 & 71 & 79 & 84 \\
Champollion$^a$ & 19 & 40 & 54 & 70 & -- & -- & -- \\
\midrule
\multicolumn{8}{@{}l}{\textit{Papalexi (monocyte)}} \\
Proposed (two-sided, mapped) & 16 & 37 & 55 & 67 & 73 & 78 & -- \\
Two-sided, pooled & 18 & 36 & 53 & 63 & 69 & 73 & -- \\
Two-sided, per condition & \ensuremath{-}17 & 14 & 34 & 52 & 65 & 74 & -- \\
One-sided (RNA eigenbasis), mapped & 16 & 37 & 55 & 67 & 73 & 78 & -- \\
Bases from the paired cells & \ensuremath{-}4.9 & 16 & 42 & 57 & 66 & 72 & -- \\
Marker coordinates, random blocks & 4.8 & 19 & 36 & 51 & 62 & 70 & -- \\
Map from measured RNA covariance & 13 & 34 & 52 & 65 & 70 & 74 & -- \\
500 unpaired cells per assay & 15 & 35 & 52 & 63 & 70 & 75 & -- \\
2,000 unpaired cells per assay & 17 & 37 & 55 & 67 & 73 & 77 & -- \\
Paired only (best of four) & 9.3 & 21 & 37 & 51 & 64 & 71 & -- \\
SemiCCA & 20 & 28 & 43 & 53 & 56 & 64 & -- \\
Reference regression & 21 & 31 & 52 & 62 & 69 & 76 & -- \\
Any comparator with the map & 21 & 32 & 55 & 67 & 73 & 79 & -- \\
Champollion$^a$ & 28 & 41 & 53 & 63 & -- & -- & -- \\
\bottomrule
\end{tabular}

\end{table}

\subsection*{External results}
The plan and code were registered, predictions of every arm were computed from training cells and recorded before the held-out files were opened, and scoring followed. The reservoir held \SpExReservoir{} paired cells from \SpExPoolPops{} training populations; the maps' split-half reliabilities were 0.49 in both conditions. Supplementary Table~\ref{tab:semiext} gives the results. The running maximum of the proposed estimator's curve crossed a recovered fraction of 0.3 between \SpExCrossProposedLoBudget{} paired cells (\SpExCrossProposedLoRf\%) and \SpExCrossProposedHiBudget{} (\SpExCrossProposedHiRf\%), at an interpolated \SpExCellsProposed{} cells; that of paired-only estimation crossed it between \SpExCrossPairedOnlyLoBudget{} (\SpExCrossPairedOnlyLoRf\%) and \SpExCrossPairedOnlyHiBudget{} (\SpExCrossPairedOnlyHiRf\%), at an interpolated \SpExCellsPairedOnly. Paired-only estimation thus needed \SpExCellsPairedOnly{} (ratio \SpExSavePairedOnly; 95\% CI, \SpExSavePairedOnlyLo--\SpExSavePairedOnlyHi; H1 met), SemiCCA \SpExCellsSemi{} (\SpExSaveSemi; \SpExSaveSemiLo--\SpExSaveSemiHi; H2 met) and Champollion \SpExCellsChamp{} (\SpExSaveChamp; \SpExSaveChampLo--\SpExSaveChampHi; H3, non-inferiority, met). At the secondary targets, the ratios against paired-only estimation, SemiCCA and Champollion were \SpExSavePairedOnlyTwo, \SpExSaveSemiTwo{} and \SpExSaveChampTwo{} at 0.2 and \SpExSaveChampFour{} for Champollion at 0.4, where paired-only estimation did not reach the target within 400 cells (ratio at least \SpExSavePairedOnlyFour). The estimator's uncorrected relative loss was the lowest of all arms from 50 paired cells on.

\paragraph{Scoring correction.} The registered plan names the two-sided estimator as the proposed estimator, and no line of it names the one-sided variant. Its predictions for every budget, draw and condition (\SpProvArrays{} arrays) were recorded before any held-out file was opened, and recomputing them from the training data with the registered code reproduced them exactly (largest absolute difference, \SpProvMaxDiff). The registered scoring code, however, set the proposed arm to the one-sided variant, a line left from its draft. This was noticed after the held-out files had been opened, and the unchanged scoring function was then applied to the plan's arm; both outputs are kept. The one-sided arm needed \SpExFrozenArmCells{} paired cells, with ratios \SpExFrozenArmPairedOnly{} (\SpExFrozenArmPairedOnlyLo--\SpExFrozenArmPairedOnlyHi), \SpExFrozenArmSemi{} (\SpExFrozenArmSemiLo--\SpExFrozenArmSemiHi) and \SpExFrozenArmChamp{} (\SpExFrozenArmChampLo--\SpExFrozenArmChampHi), so it also met all three hypotheses (Supplementary Table~\ref{tab:extarms}). The correction rests on the provenance above, not on this agreement.

\begin{table}[h]
\caption{\textbf{External test: recovered fraction (\%) by number of paired cells, uncorrected relative loss, paired cells needed to reach a recovered fraction of 0.3 and prespecified comparisons.} Recovered fraction (\%) of within-condition gene--protein cross-correlation among \SpExOrfs{} held-out ORFs (\SpExCells{} cells), averaged over 20 draws of paired cells from a reservoir of \SpExReservoir, and uncorrected relative squared loss (independence, 1). Below, the paired cells each method needed to reach the prespecified target, interpolated log-linearly between the bracketing budgets on the running maximum of its curve, and their ratio to the proposed estimator's, with 95\% intervals from 2,000 bootstrap resamples of the held-out ORFs; H1 and H2 required a lower bound above one and H3, a non-inferiority test, above 0.8. $^a$Best at each budget of James--Stein shrinkage, SCOSE, FCOSE and cross-validated low rank, pooled or per condition. $^b$Better of pooled and per-condition fits (reference regression: paired or unpaired Gram matrix).}
\label{tab:semiext}
\centering\footnotesize
% Generated by extension/semipaired/make_assets.py. Do not edit.
\begin{tabular}{@{}lrrrrrrr@{}}
\toprule
& \multicolumn{5}{c}{Recovered (\%) at paired cells} & \multicolumn{2}{c}{Loss} \\
\cmidrule(lr){2-6}\cmidrule(l){7-8}
Estimator & 25 & 50 & 100 & 200 & 400 & 100 & 400 \\
\midrule
Proposed & 5.2 & 23 & 36 & 47 & 57 & 0.800 & 0.685 \\
Paired cells only$^a$ & 2.6 & 5.7 & 12 & 20 & 32 & 0.933 & 0.823 \\
SemiCCA$^b$ & 12 & 16 & 27 & 38 & 47 & 0.851 & 0.741 \\
Champollion & 0.9 & 8.6 & 23 & 41 & 54 & 0.874 & 0.698 \\
Reference regression$^b$ & 12 & 17 & 27 & 35 & 47 & 0.850 & 0.738 \\
Any comparator with the map & 12 & 18 & 28 & 39 & 50 & 0.845 & 0.724 \\
\midrule
\multicolumn{8}{@{}l}{Paired cells to reach a recovered fraction of 0.3 (proposed: 73)} \\
Comparison & \multicolumn{2}{r}{Cells} & \multicolumn{3}{r}{Ratio (95\% CI)} & \multicolumn{2}{r}{Verdict} \\
H1: paired cells only & \multicolumn{2}{r}{357} & \multicolumn{3}{r}{4.9 (3.2--6.0)} & \multicolumn{2}{r}{met} \\
H2: SemiCCA & \multicolumn{2}{r}{122} & \multicolumn{3}{r}{1.7 (1.3--2.0)} & \multicolumn{2}{r}{met} \\
H3: Champollion (non-inferiority) & \multicolumn{2}{r}{131} & \multicolumn{3}{r}{1.8 (1.2--2.2)} & \multicolumn{2}{r}{met} \\
Every comparator, original form & \multicolumn{2}{r}{115} & \multicolumn{3}{r}{1.6 (1.12--1.9)} & \multicolumn{2}{r}{--} \\
Every comparator, with or without map & \multicolumn{2}{r}{111} & \multicolumn{3}{r}{1.5 (1.11--1.8)} & \multicolumn{2}{r}{--} \\
\bottomrule
\end{tabular}

\end{table}

\begin{table}[h]
\caption{\textbf{External test: the arm named in the registered plan and the arm named in the registered scoring code.} Paired cells needed to reach a recovered fraction of 0.3, interpolated as in the main analysis, and ratios of each comparator's paired cells to the arm's, with 95\% intervals from 2,000 resamples of held-out ORFs. H1 (against paired-only estimation) and H2 (against SemiCCA) required lower bounds above one, and H3 (against Champollion) a lower bound above 0.8. The main text reports the plan's arm.}
\label{tab:extarms}
\centering\footnotesize
\begin{tabular}{@{}lrlllc@{}}
\toprule
Arm (named by) & Paired cells & Paired only & SemiCCA & Champollion & H1--H3 \\
\midrule
Two-sided blocks (registered plan) & \SpExCellsProposed & \SpExSavePairedOnly{} (\SpExSavePairedOnlyLo--\SpExSavePairedOnlyHi) & \SpExSaveSemi{} (\SpExSaveSemiLo--\SpExSaveSemiHi) & \SpExSaveChamp{} (\SpExSaveChampLo--\SpExSaveChampHi) & met \\
One-sided blocks (registered scoring code) & \SpExFrozenArmCells & \SpExFrozenArmPairedOnly{} (\SpExFrozenArmPairedOnlyLo--\SpExFrozenArmPairedOnlyHi) & \SpExFrozenArmSemi{} (\SpExFrozenArmSemiLo--\SpExFrozenArmSemiHi) & \SpExFrozenArmChamp{} (\SpExFrozenArmChampLo--\SpExFrozenArmChampHi) & met \\
\bottomrule
\end{tabular}
\end{table}

\section{The benchmark: tissues, pairs of modalities and nervous systems}\label{sec:gen}

This note gives the protocol and complete results of the benchmark, which tested whether the saving of the estimator holds when independent replication units are held out, across tissues, species, modality pairs and nervous systems, and whether its size follows from how the dependence is concentrated in the blocks along the unpaired axes. The plan, the code and the extracted data files were registered before the estimators were run on any of these data sets and before any statistic of held-out cells had been computed.

\subsection*{Data sets}
Nine data sets were analysed (Supplementary Table~\ref{tab:gen}). Four pair RNA with surface proteins: peripheral blood from 8 donors\cite{hao2021}, peripheral blood from 118 patients with COVID-19 or controls\cite{stephenson2021}, bone marrow from 9 donors and colon from 12 patients\cite{mennillo2024}. One pairs RNA with chromatin accessibility (bone-marrow multiome, 10 donors). Two pair the transcriptome of a neuron with its electrophysiological features, measured in the same cell by Patch-seq: mouse motor cortex\cite{scala2021} (263 mice) and mouse visual-cortex interneurons\cite{gouwens2020,gala2021} (362 recording days; expression was available only as log counts per million, so the latent correlation equalled the measured one). Two pair the wiring of the same neurons in two parts of the fly central nervous system: for every descending and ascending neuron, the synapse counts with partner groups in the brain (X) and in the nerve cord (Y), in a female brain-and-nerve-cord connectome\cite{bates2026} (BANC; 913 cell types) and a male central-nervous-system connectome\cite{berg2025} (1,010 cell types). Partner groups were the super class with the hemilineage, or the cell class where no hemilineage was annotated, with inputs and outputs counted separately. Brain-side partners had their soma in the brain and nerve-cord-side partners in the nerve cord; partners that cross the neck and non-neuronal profiles were excluded, and neurons needed at least 50 synapses on each side. Droplet data were capped at 600 (COVID-19) or 3,000 cells per sample. Extraction used cell counts and detection rates of one modality at a time and no statistic of dependence between modalities. The blood, bone-marrow and colon data had been used in other analyses of this study but not for semi-paired estimation.

\subsection*{Design}
Replication units were donors, patients, mice, recording days or neuron types. Every unit was held out in exactly one of three folds (units in a fixed pseudo-random order, the \textit{i}-th in fold \textit{i} mod 3), and in each fold the cells of the other units were training cells. As in development, a fixed pseudo-random rule put 25\% of training cells in a reservoir of paired cells, and the others formed an unpaired pool, each contributing X or Y by a second pseudo-random draw. Features were chosen from each fold's training cells: 200 X features by the binned-dispersion rule of the external test, surface proteins detected in at least 1\% of cells (at most 200, centred log-ratios), 200 accessible peaks or nerve-cord partner groups by the same dispersion rule, and every electrophysiological feature. Cells were centred within sample and condition in droplet data, and within condition across units in Patch-seq and connectome data, in which each unit contributes few cells. Conditions were cell types, transcriptomic families or subclasses, and descending or ascending neurons.

The arms were the proposed estimator, unchanged; the self-tuning variant chosen in development (below); James--Stein shrinkage in marker coordinates, SCOSE, FCOSE and cross-validated low rank (paired cells only), SemiCCA and reference ridge regression, each pooled, refitted per condition and with the proposed map; the proposed shrinkage in bases estimated from the paired cells; and all training cells used as paired cells. Budgets were 25 to 800 paired cells, up to the size of each fold's reservoir, with five draws per budget and fold, and every arm received the same paired cells and the same unpaired pool. Champollion was not included because of its cost (about 9~s per fit). The endpoint was the noise-unbiased recovered fraction of within-population cross-correlation pooled within condition over held-out units and summed over folds and conditions, as in the external test. Because the dependence that can be recovered differs between data sets, targets were set relative to the recovered fraction with every training cell paired: half of it (primary), a quarter and three quarters. Paired cells needed were interpolated as before, and a curve that never reached a target counted as needing the largest budget. Intervals came from 2,000 bootstrap resamples of units, recomputing the targets and the reference. Predictions of every data set were recorded before its held-out statistics were computed.

\subsection*{Hypotheses}
G1 (primary) required the geometric mean over the nine data sets of the saving against the paired-only envelope, at the primary target, to have a 95\% lower bound above one, the interval combining the \textit{r}-th bootstrap resample of every data set; G2 required the same against SemiCCA. G3 compared the self-tuning variant with the fixed estimator (geometric mean of the ratio of paired cells needed, two-sided, without a threshold). G4 tested a structural law: across the nine data sets and the three earlier ones, the saving predicted from the concentration of dependence in the blocks along the unpaired axes, computed from training cells only (the paired reservoir, which gives the signal in each block, and the unpaired pool; never held-out cells), should be positively rank-correlated with the observed saving (one-sided permutation test, 10,000 permutations, $P<0.05$).

\paragraph{Predicted saving.} With the fold-0 reservoir of $N$ paired training cells as the population, the mean per-cell product in the unpaired two-sided eigenbasis gives, for each block $b$, a signal $a_b^2=\max(\|\hat\theta_b\|^2-\tau_b/N,0)$, where $\tau_b$ is the trace of the covariance of the per-cell products. The risk of the oracle linear shrinkage of each block, $\sum_ba_b^2\tau_b/(Ba_b^2+\tau_b)$, and that of a single block of all coefficients, $A\tau/(BA+\tau)$ with $A=\sum_ba_b^2$ and $\tau=\sum_b\tau_b$, fall to half of $A$ at budgets $B_{\mathrm{blocks}}$ and $B_{\mathrm{single}}=\tau/A$. The predicted saving is $B_{\mathrm{single}}/B_{\mathrm{blocks}}$. It is large when the dependence lies in a few blocks whose share of the per-cell noise is small, the condition of Proposition~\ref{prop:bjs}'s corollary. For the three earlier data sets the reservoirs of the bound check were used, and their observed savings were read at half their reference's recovered fraction from their stored curves.

\subsection*{Amendments}
The prespecified targets were fractions of what an unshrunk estimator using every training cell as a paired cell recovered. After three data sets had been scored, this reference turned out to be dominated by noise in conditions with few cells. Its squared norm swamped the noise-unbiased denominator, so its recovered fraction was strongly negative (COVID-19 blood \GenRefStephenson\%, colon \GenRefColon\%) even though it recovered 42--87\% in the large conditions. The targets were then negative, every curve began above them, and the savings equalled one by construction. The same happened later in motor-cortex Patch-seq (\GenRefScala\%) and bone-marrow multiome (\GenRefBmMulti\%). In blood from 8 donors the reference recovered \GenRefHao\%, and neither the proposed estimator nor the paired-only envelope reached half of it within 800 paired cells, so that saving was also one. Before the other six data sets were scored, we therefore recorded an amendment. It set targets at fractions (0.5 primary, 0.25 and 0.75) of the largest recovered fraction that any method reached at any budget, recomputed in every bootstrap resample, and left the savings, censoring, intervals and hypotheses otherwise unchanged (G1$'$--G4$'$). After the fourth data set had been scored, a second amendment declared a data set not informative for savings if the bootstrap lower bound of that level was not above zero, that is, if no method recovered dependence in held-out units; geometric means are reported with and without this rule. The complete amended rule, targets and exclusion together, therefore preceded the scoring of five of the nine data sets. COVID-19 blood, colon and blood from 8 donors were scored before either amendment, and motor-cortex Patch-seq between them; the exclusion removed only motor-cortex Patch-seq, and the other five data sets and the secondary analysis were scored under the complete amended rule. The recorded reference values were later corrected by an erratum; they had been quoted 100 times too small. Both analyses are reported. The amended savings were computed from the same recorded predictions by code that first reproduces every saving of the prespecified rule exactly.

\subsection*{Self-tuning variant}
The fixed estimator uses one block partition and one James--Stein factor per block for every data set. Before the benchmark, three self-tuning variants were compared on the two development screens under a rule written before they were run: a block partition chosen among 42 candidates by Stein's unbiased risk estimate (SURE); preconditioned shrinkage, in which every coefficient of a block is shrunk by $\tau^2/(\tau^2+s_j)$, with $s_j$ its own noise variance and $\tau^2=\max\{0,(\|m\|^2-t)/d\}$ for a block of $d$ coefficients with mean $m$ and noise trace $t$; and both combined, choosing per block the rule with the smaller SURE. With equal noise variances the preconditioned rule equals the fixed estimator's factor $1-t/\|m\|^2$, so it differs only in weighting each coefficient by its own noise. Its SURE, exact for Gaussian block means including the dependence of $\tau^2$ on $m$, is $\sum_jb_j^2m_j^2+2\sum_js_j(1-b_j)+(4/d)\sum_jm_js_j(Sm)_j/(\tau^2+s_j)^2-t$ with $b_j=s_j/(\tau^2+s_j)$; a synthetic check gave mean SURE 2.51 against a mean loss of 2.53. Averaged over budgets, the preconditioned variant recovered \SpSelfCriterion\% against \SpFixedCriterion\% for the fixed estimator in the two screens (at 50 paired cells, \SpSelfFrFifty\% against \SpFixedFrFifty\% in the melanoma screen and \SpSelfPaFifty\% against \SpFixedPaFifty\% in the monocyte screen), whereas choosing the partition by SURE lost recovery at small budgets (\SpSelfPartitionCriterion\%), where SURE is noisy and selection among many partitions overfits. The preconditioned variant was therefore fixed as the benchmark's self-tuning arm.

\subsection*{Results}
Supplementary Table~\ref{tab:gen} gives, for every data set, the best recovery that any method reached and the method that reached it, the paired cells that the proposed estimator and the paired-only envelope needed to reach half of it, and the savings against paired-only estimation and against SemiCCA, together with the saving under the prespecified rule. Across the \GenNInformative{} informative data sets the geometric-mean saving was \GenGOne{} against paired-only estimation (95\% CI, \GenGOneLo--\GenGOneHi; G1$'$ \GenGOneMet) and \GenGTwo{} against SemiCCA (\GenGTwoLo--\GenGTwoHi; G2$'$ \GenGTwoMet). Counting every data set with its censored saving, it was \GenGOneAll{} (\GenGOneAllLo--\GenGOneAllHi) and \GenGTwoAll{} (\GenGTwoAllLo--\GenGTwoAllHi). At a quarter and at three quarters of the best recovery, the savings against paired-only estimation were \GenGOneQuarter{} (\GenGOneQuarterLo--\GenGOneQuarterHi) and \GenGOneThreeQuarters{} (\GenGOneThreeQuartersLo--\GenGOneThreeQuartersHi). Against every comparator in its original form, including reference regression and SemiCCA refitted per condition, the geometric-mean saving was \GenGOriginal{} (\GenGOriginalLo--\GenGOriginalHi). Under the prespecified rule, G1 was \GenFrozenGOneMet{} (\GenFrozenGOne; \GenFrozenGOneLo--\GenFrozenGOneHi) and G2 \GenFrozenGTwoMet{} (\GenFrozenGTwo; \GenFrozenGTwoLo--\GenFrozenGTwoHi); \GenFrozenNTrivial{} of its nine savings were exactly one, because the targets were below every curve (four data sets) or above the proposed and comparator curves (blood, 8 donors). These values of one encode unresolved comparisons, neither measured savings nor lower bounds on them, so the prespecified geometric means show only that the stated criterion was met. The self-tuning variant needed more paired cells than the fixed estimator (ratio of fixed to self-tuned cells, \GenGThree; \GenGThreeLo--\GenGThreeHi; G3$'$). The structural law was \GenLawFrozenMet{} with the prespecified savings ($\rho=\GenLawFrozenRho$, $P=\GenLawFrozenP$; G4) and \GenLawMet{} with the amended ones (Spearman $\rho=\GenLawRho$ over \GenLawN{} data sets, one-sided $P=\GenLawP$; G4$'$). Six of the eleven amended savings were censored. A post hoc concordance analysis, which counts only pairs whose order is known when the smaller saving is exact, gave Harrell's\cite{harrell1982} $C=\GenLawC$ over \GenLawCPairs{} pairs (permutation $P=\GenLawCP$). In the five data sets with exact savings, the rank correlation was \GenLawUncensoredRho{} and the observed saving was \GenLawRatioMin--\GenLawRatioMax{} times the predicted one. This analysis is exploratory.

The prespecified secondary analysis with unpaired neurons of other animals is reported, with its post hoc reanalysis, in Supplementary Note~\ref{sec:deploy}.

\begin{table}[h]
\caption{\textbf{Benchmark.} Best: the largest recovered fraction of within-condition dependence in held-out units that any method reached within the budgets (25 to 800 paired cells, up to each fold's reservoir), and the method that reached it. Paired cells to half of best: interpolated on each curve's running maximum; $>$, not reached within the largest budget, so the saving is a lower bound ($\geq$); $\leq$, reached at the smallest budget, so the saving is again a lower bound (male fly connectome). Savings: cells the comparator envelope needs divided by cells the proposed estimator needs, with 95\% intervals from 2,000 bootstrap resamples of held-out units. n.e., not estimable (no method recovered dependence). Prespecified rule: the saving at half the recovered fraction of the unshrunk fully paired reference (Supplementary Note~\ref{sec:gen}).}
\label{tab:gen}
\centering\scriptsize
\resizebox{\textwidth}{!}{% Generated by extension/generality/make_assets.py. Do not edit.
\begin{tabular}{@{}lrrlrrlll@{}}
\toprule
Data set & Cells & Best (\%) & Reached by & \multicolumn{2}{c}{Paired cells to half of best} & Saving vs & Saving vs & Prespecified\\
\cmidrule(lr){5-6}
 & & & & Proposed & Paired only & paired only & SemiCCA & rule\\
\midrule
\multicolumn{9}{@{}l}{\textit{RNA--protein}}\\
Blood (8 donors) & 72,000 & 35 & ref.\ regr.\ + map & 408 & $>$800 & $\geq$2.0 (1.00--2.1) & $\geq$2.0 (1.00--2.1) & 1.00\\
Blood, COVID-19 (118 patients) & 70,800 & 13 & ref.\ regression & 390 & $>$800 & $\geq$2.0 (1.00--2.6) & 1.02 (0.53--1.3) & 1.00\\
Bone marrow (9 donors) & 36,000 & 46 & proposed & 168 & $>$800 & $\geq$4.7 (1.00--6.0) & 3.7 (1.00--4.6) & 2.5\\
Colon (12 patients) & 15,953 & 25 & proposed & 126 & $>$800 & $\geq$6.4 (3.3--7.7) & $\geq$6.4 (1.9--7.7) & 1.00\\
\multicolumn{9}{@{}l}{\textit{RNA--chromatin}}\\
Bone marrow multiome (10 donors) & 36,450 & 25 & proposed & 285 & $>$800 & $\geq$2.8 (1.7--3.0) & $\geq$2.8 (1.7--3.0) & 1.00\\
\multicolumn{9}{@{}l}{\textit{RNA--electrophysiology (Patch-seq)}}\\
Motor cortex (263 mice) & 1,203 & 0.0 & none (low rank predicts zero) & -- & -- & n.e. & n.e. & 1.00\\
Visual cortex (362 days) & 3,395 & 42 & ref.\ regression & 157 & 219 & 1.4 (1.2--1.6) & 1.3 (1.10--1.5) & 1.4\\
\multicolumn{9}{@{}l}{\textit{Brain--nerve-cord wiring (connectomes)}}\\
Fly CNS, female (913 types) & 2,487 & 68 & proposed & 25 & 90 & 3.6 (3.1--3.8) & 2.3 (2.0--2.4) & 3.7\\
Fly CNS, male (1,010 types) & 3,009 & 69 & ref.\ regr.\ + map & $\leq$25 & 73 & 2.9 (2.6--3.1) & 1.9 (1.7--2.1) & 3.1\\
\bottomrule
\end{tabular}
}
\end{table}

\paragraph{The saving in every test data set.} The main text summarizes the saving over paired-only estimation by its geometric mean over the test data sets in which dependence was recovered: the external test at its prespecified target and the nine benchmark data sets at the amended primary target. Supplementary Table~\ref{tab:savings} gives, for each data set, the target and the saving that enter this summary, whether the saving is a lower bound, and the target and saving under the benchmark's original plan. The summary was computed after all tests had been scored, combines targets that differ between data sets and includes lower bounds; it describes the tested data sets.

\begin{table}[h]
\caption{\textbf{Saving of paired cells in every test data set.} Saving: paired cells needed by paired-only estimation (the best at each budget of James--Stein shrinkage, SCOSE, FCOSE and cross-validated low rank, pooled or per condition) divided by those needed by the estimator, with 95\% intervals from resampling held-out units; $\geq$, a curve did not cross the target within the budgets, so the saving is a lower bound. Rule of the summary: 30\% recovery in the external test (prespecified) and half of the best recovery that any method reached in the benchmark (amended). Original plan: half of the recovered fraction of the unshrunk estimator with every training cell paired; a negative target lies below every curve and an unreached one above both, and either gives a saving of one by construction. The geometric means carry 95\% intervals combining the $r$-th resample of every data set.}
\label{tab:savings}
\centering\scriptsize
\resizebox{\textwidth}{!}{% Generated by extension/synthesis/make_assets.py. Do not edit.
\begin{tabular}{@{}llrlrl@{}}
\toprule
 & & \multicolumn{2}{c}{Rule of the summary} & \multicolumn{2}{c}{Original plan} \\
\cmidrule(lr){3-4}\cmidrule(lr){5-6}
Data set & Units held out & Target (\%) & Saving (95\% CI) & Target (\%) & Saving \\
\midrule
T cells, ORF screen (external test) & ORFs & 30 & 4.9 (3.2--6.0) & 30 & 4.9 \\
Blood, 8 donors & donors & 17.7 & $\geq$2.0 (1.0--2.1) & 35.0 & 1 (neither reached) \\
Blood, COVID-19 & patients & 6.4 & $\geq$2.0 (1.0--2.6) & \ensuremath{-}521.2 & 1 (negative target) \\
Bone marrow & donors & 23.0 & $\geq$4.7 (1.0--6.0) & 32.9 & $\geq$2.5 \\
Colon & patients & 12.7 & $\geq$6.4 (3.3--7.7) & \ensuremath{-}118.6 & 1 (negative target) \\
Bone marrow multiome & donors & 12.5 & $\geq$2.8 (1.7--3.0) & \ensuremath{-}74.0 & 1 (negative target) \\
Motor cortex Patch-seq & mice & 0.0 & no dependence recovered & \ensuremath{-}136.9 & 1 (negative target) \\
Visual cortex Patch-seq & recording days & 20.8 & 1.4 (1.2--1.6) & 21.5 & 1.4 \\
Fly CNS, female & neuron types & 34.1 & 3.6 (3.1--3.8) & 37.9 & 3.7 \\
Fly CNS, male & neuron types & 34.3 & $\geq$2.9 (2.6--3.1) & 38.5 & 3.1 \\
\midrule
\multicolumn{3}{@{}l}{Geometric mean, nine test data sets with dependence recovered} & 3.1 (2.4--3.2) & & \\
\multicolumn{3}{@{}l}{Geometric mean, eight benchmark data sets with dependence recovered} & 2.9 (2.3--3.0) & & \\
\multicolumn{3}{@{}l}{Geometric mean, nine benchmark data sets, original plan (G1)} & & & 1.5 (1.4--1.6) \\
\bottomrule
\end{tabular}
}
\end{table}

\section{Predicting the saving: theory and a prospective test}\label{sec:pred}

Propositions~\ref{prop:bjs} and~\ref{prop:minimax} concern idealized estimators. This note bounds the rule that is implemented (Proposition~\ref{prop:oracle}), derives the saving of paired cells that the unpaired block structure buys (Proposition~\ref{prop:law}), and states what unpaired cells can and cannot reveal about it (Proposition~\ref{prop:reveal}). It then gives the development and the prespecified prospective test of the resulting law.

\subsection*{Setting}
Everything is conditional on the unpaired pool, which fixes the bases $U$ and $V$, the scaling, the population means used for centring and the blocks. Paired cells are independent draws from the population. For a block $b$, the product of an RNA eigen-block $R_b$ and a protein eigen-block $C_b$ with $d_b$ coefficients, let $z_{ib}$ be the coefficients of cell $i$'s product $\tilde x_i\tilde y_i^\top$ in the block ($\tilde x=U^\top x$, $\tilde y=V^\top y$), and
\[
\theta_b=\mathbb{E}z_{ib},\quad \Sigma_b=\operatorname{Cov}z_{ib},\quad \tau_b=\operatorname{tr}\Sigma_b,\quad \ell_b=\lambda_{\max}(\Sigma_b),\quad \kappa_b=\frac{\operatorname{Var}\|z_{ib}-\theta_b\|^2}{\tau_b^2},\quad r_b^2=\|\theta_b\|^2 .
\]
From $B$ cells the implemented rule forms the block mean $m_b$, the unbiased estimate $t_b=\sum_i\|z_{ib}-m_b\|^2/\{B(B-1)\}$ of $\operatorname{tr}\operatorname{Cov}(m_b)=\tau_b/B$, and $\hat\theta_b=(1-t_b/\|m_b\|^2)_+m_b$; paired-only James--Stein is the same rule with one block of all coefficients. For a fixed number $c$, $\mathbb{E}\|cm_b-\theta_b\|^2=c^2\tau_b/B+(1-c)^2r_b^2$ for any distribution with these first two moments, and its minimum over $c$,
\[
\rho(r_b^2,\tau_b;B)=\frac{r_b^2\tau_b}{Br_b^2+\tau_b},
\]
is the oracle-linear risk. It depends on the block only through its signal $r_b^2$ and its per-cell noise trace $\tau_b$.

\begin{proposition}[The implemented rule without Gaussian assumptions]\label{prop:oracle}
If $B\ge2$ and $\mathbb{E}\|z_{ib}\|^4<\infty$, then
\[
\Big|\big(\mathbb{E}\|\hat\theta_b-\theta_b\|^2\big)^{1/2}-\rho(r_b^2,\tau_b;B)^{1/2}\Big|\le\Big(\frac{36\,\ell_b}{B}+\frac{22\,\kappa_b\tau_b}{B^2}\Big)^{1/2}.
\]
\end{proposition}
\noindent\textit{Proof.} Drop $b$. Write $s=\tau/B$, $A=\|m\|^2$, $\mu=\mathbb{E}A=r^2+s$ and $a^*=1-s/\mu\in[0,1)$; the rule $a^*m$ has risk $\rho$. Because $t\ge0$, the implemented factor is $a=\min\{\max(1-t/A,0),1\}$ ($a=0$ if $A=0$), and $\hat\theta=am$. Minkowski's inequality in $L^2$ gives $|(\mathbb{E}\|\hat\theta-\theta\|^2)^{1/2}-\rho^{1/2}|\le(\mathbb{E}[(a-a^*)^2A])^{1/2}=:\Delta^{1/2}$. Clipping to $[0,1]$ is 1-Lipschitz and fixes $a^*$, so $|a-a^*|\le|t/A-s/\mu|$ for $A>0$ and $|a-a^*|\le1$ always. On $\{A\ge\mu/2\}$, $(t/A-s/\mu)^2A=\{(t-s)\mu+s(\mu-A)\}^2/(A\mu^2)\le4(t-s)^2/\mu+4s^2(A-\mu)^2/\mu^3$. On $\{A<\mu/2\}$, $(a-a^*)^2A<\mu/2$, and Chebyshev's inequality gives $\Pr(A<\mu/2)\le4\operatorname{Var}(A)/\mu^2$. As $\mathbb{E}t=s$, $\Delta\le\{4\operatorname{Var}(t)+6\operatorname{Var}(A)\}/\mu$. With $e_i=z_i-\theta$,
\[
\operatorname{Var}(A)=\frac{4\theta^\top\Sigma\theta}{B}+\frac{4\,\mathbb{E}[\theta^\top e(\|e\|^2-\tau)]}{B^2}+\frac{\operatorname{Var}\|e\|^2+2(B-1)\operatorname{tr}\Sigma^2}{B^3},\qquad \operatorname{Var}(t)=\frac{1}{B^2}\Big(\frac{\operatorname{Var}\|e\|^2}{B}+\frac{2\operatorname{tr}\Sigma^2}{B(B-1)}\Big),
\]
the second being the variance of the unbiased trace of a sample covariance, a U-statistic of order two\cite{hoeffding1948}. Using $\theta^\top\Sigma\theta\le\ell r^2$, $\operatorname{tr}\Sigma^2\le\ell\tau$, the Cauchy--Schwarz bound $|\mathbb{E}[\theta^\top e(\|e\|^2-\tau)]|\le(\ell r^2)^{1/2}\kappa^{1/2}\tau$ and $4xy\le2x^2+2y^2$, $\operatorname{Var}(A)\le6\ell r^2/B+2\ell\tau/B^2+3\kappa\tau^2/B^3$ and $\operatorname{Var}(t)\le\kappa\tau^2/B^3+\ell\tau/B^2$. Hence $4\operatorname{Var}(t)+6\operatorname{Var}(A)\le36\ell r^2/B+16\ell\tau/B^2+22\kappa\tau^2/B^3$, and dividing by $\mu=r^2+\tau/B$ gives $\Delta\le36\ell/B+22\kappa\tau/B^2$. $\square$

Dividing the bound by the unshrunk root-noise scale $\sqrt{\tau/B}$ gives $\sqrt{36/d_{\mathrm{eff}}+22\kappa/B}$, where $d_{\mathrm{eff}}=\tau/\ell$ is the block's effective dimension. Thus the absolute root-risk gap is small on this scale for blocks of many effective dimensions and budgets large relative to $\kappa$. This does not imply uniform relative accuracy against the oracle risk, which vanishes at zero signal. The constant is not sharp; in 72 simulated cases with Gaussian, log-normal and Poisson margins the gap never exceeded a tenth of the bound. With one block the statement covers paired-only James--Stein. Given the pool, $\theta_b$ includes the outer product of the errors of the unpaired means used for centring; the root risks against $\theta_b$ and the population cross-covariance differ by at most the norm of that outer product.

\begin{proposition}[The saving is set by where dependence lies relative to where noise lies]\label{prop:law}
With blocks, the oracle-linear risk is $R_K(B)=\sum_b\rho(r_b^2,\tau_b;B)$; with one block of all coefficients it is $R_1(B)=\rho(r^2,\tau;B)$, $r^2=\sum_br_b^2$, $\tau=\sum_b\tau_b$. For a relative risk $\varepsilon\in(0,1)$ let $B_K(\varepsilon)$ and $B_1(\varepsilon)$ solve $R_K(B)=\varepsilon r^2$ and $R_1(B)=\varepsilon r^2$, and let $S(\varepsilon)=B_1(\varepsilon)/B_K(\varepsilon)$. Write $w_b=r_b^2/r^2$ for the dependence shares and $\pi_b=\tau_b/\tau$ for the noise shares, and assume $r^2>0$ and $\tau_b>0$. Then:
\begin{enumerate}
\item[(a)] $S(\varepsilon)=(1-\varepsilon)/\{\varepsilon\beta(\varepsilon)\}$, where $\beta$ solves $\sum_bw_b\pi_b/(\beta w_b+\pi_b)=\varepsilon$; the saving depends on the dependence only through $w$, and $B_1=(\tau/r^2)(1-\varepsilon)/\varepsilon$.
\item[(b)] $S(\varepsilon)\ge1$, with equality, at any and then every $\varepsilon$, if and only if $w=\pi$.
\item[(c)] $S(\varepsilon)\le1/\min\{\pi_b:w_b>0\}$, with equality when all dependence lies in one block.
\item[(d)] $S(\varepsilon)\to1/\sum_{b:w_b>0}\pi_b$ as $\varepsilon\to0$, and $S(\varepsilon)\to\sum_bw_b^2/\pi_b=1+\chi^2(w\,\|\,\pi)$ as $\varepsilon\to1$.
\end{enumerate}
\end{proposition}
\noindent\textit{Proof.} $\rho(x,y;B)=xy/(Bx+y)=\min_c\{c^2y/B+(1-c)^2x\}$ is a minimum of linear functions, hence jointly concave, and it is positively homogeneous of degree one. (a) With $B=\beta\tau/r^2$, $\rho(r_b^2,\tau_b;B)/r^2=w_b\pi_b/(\beta w_b+\pi_b)$ and $R_1(B)/r^2=1/(1+\beta)$; the left side of the equation for $\beta$ decreases strictly from 1 to 0. (b) $R_1(B)=\min_c\sum_b\{c^2\tau_b/B+(1-c)^2r_b^2\}\ge\sum_b\min_{c_b}\{c_b^2\tau_b/B+(1-c_b)^2r_b^2\}=R_K(B)$, with equality if and only if the blockwise minimizers $r_b^2/(r_b^2+\tau_b/B)$ coincide, that is, $w=\pi$. At $B=B_1(\varepsilon)$, $R_K\le\varepsilon r^2$ and $R_K$ decreases in $B$, so $B_K\le B_1$, strictly unless $w=\pi$. (c) Let $\tau_+=\tau\min\{\pi_b:w_b>0\}$. For $w_b>0$, $\rho(r_b^2,\tau_b;B)\ge\rho(r_b^2,\tau_+;B)$, and $x\mapsto\rho(x,\tau_+;B)$ is concave and vanishes at zero, hence subadditive, so $R_K(B)\ge\rho(r^2,\tau_+;B)$; $R_K(B)\le\varepsilon r^2$ then requires $B\ge(\tau_+/r^2)(1-\varepsilon)/\varepsilon$. (d) As $\varepsilon\to0$, $\beta\to\infty$ and the left side is $\beta^{-1}\sum_{w_b>0}\pi_b+O(\beta^{-2})$; as $\varepsilon\to1$, $\beta\to0$ and it is $1-\beta\sum_bw_b^2/\pi_b+O(\beta^2)$. $\square$

\begin{proposition}[What unpaired cells reveal]\label{prop:reveal}
For parts (a) and (b), centre each modality at its population mean.
\begin{enumerate}
\item[(a)] If $x$ and $y$ are independent, with finite fourth moments, then
\[
\Sigma_b=\operatorname{Cov}(\tilde x_{R_b})\otimes\operatorname{Cov}(\tilde y_{C_b}),\qquad
\tau_b=\operatorname{tr}\operatorname{Cov}(\tilde x_{R_b})\operatorname{tr}\operatorname{Cov}(\tilde y_{C_b}).
\]
Thus $\tau_b$ and $\ell_b$ are functions of the two marginal distributions. If $(x,y)$ is jointly Gaussian, $\tau_b$ is the independence value plus $r_b^2$, and $r_b^2$ is at most that independence value.
\item[(b)] Given two Gaussian marginals with positive-definite covariances and any $w$ on the simplex, there are jointly Gaussian laws with these marginals and dependence shares $w$ for every sufficiently small $r^2>0$. As $r^2\to0$, the limiting savings compatible with the unpaired data fill $[1,1/\min_b\pi_b]$, where $\pi$ denotes the independence noise shares, and $B_K(\varepsilon)=\beta(\varepsilon)\tau/r^2$ is unbounded.
\item[(c)] For $m\ge2$ independent paired observations, conditional on the unpaired pool as in the setting above, the raw estimate $\widehat r_b^2=\|m_b\|^2-t_b$ is unbiased and has variance
\[
\operatorname{Var}(\widehat r_b^2)
=\frac{4\theta_b^\top\Sigma_b\theta_b}{m}
+\frac{2\operatorname{tr}(\Sigma_b^2)}{m(m-1)}.
\]
For $r_b^2>0$, its relative standard deviation is at most
\[
\frac{\{\operatorname{Var}(\widehat r_b^2)\}^{1/2}}{r_b^2}
\le\left\{\frac{4\ell_b}{m r_b^2}
+\frac{2\ell_b\tau_b}{m(m-1)r_b^4}\right\}^{1/2}.
\]
\end{enumerate}
\end{proposition}
\noindent\textit{Proof.} (a) For population-centred Gaussian variables, Isserlis' theorem\cite{isserlis1918} gives
\[
\begin{aligned}
\operatorname{Cov}(\tilde x_k\tilde y_l,\tilde x_{k'}\tilde y_{l'})
={}&\operatorname{Cov}(\tilde x_k,\tilde x_{k'})\operatorname{Cov}(\tilde y_l,\tilde y_{l'})\\
&+\operatorname{Cov}(\tilde x_k,\tilde y_{l'})\operatorname{Cov}(\tilde x_{k'},\tilde y_l).
\end{aligned}
\]
Summing the diagonal gives $\tau_b$. Under independence the covariance factorizes for any centred distributions. Summing the inequalities $\operatorname{Cov}(\tilde x_k,\tilde y_l)^2\le\operatorname{Var}(\tilde x_k)\operatorname{Var}(\tilde y_l)$ proves the bound on $r_b^2$.

(b) In the eigenbases, with marginal covariances $\Lambda$ and $M$, take unit-Frobenius matrices $E_b$ supported on the blocks and $\Theta=c\sum_bw_b^{1/2}E_b$. The joint covariance is positive semidefinite when $\|\Lambda^{-1/2}\Theta M^{-1/2}\|_{\mathrm{op}}\le1$, which holds for sufficiently small $c$. Its dependence shares are $w$, its squared signal is $r^2=c^2$, and its noise shares tend to their independence values by (a). At these limiting noise shares, the saving is a continuous function of $w$ on the simplex by Proposition~\ref{prop:law}(a). It equals 1 at $w=\pi$ and $1/\min_b\pi_b$ at the vertex of the least noisy block, so it takes every intermediate value; Proposition~\ref{prop:law}(b,c) excludes values outside that interval. For fixed $w$, $\beta(\varepsilon)>0$ and $\tau$ tends to a positive value, so the required budget diverges as $r^2\to0$.

(c) Drop $b$ and write $e_i=z_i-\theta$. The raw estimate is the U-statistic
\[
\widehat r^2=\frac{1}{m(m-1)}\sum_{i\ne j}z_i^\top z_j,
\qquad
\widehat r^2-r^2=\frac{2}{m}\sum_i\theta^\top e_i
+\frac{1}{m(m-1)}\sum_{i\ne j}e_i^\top e_j.
\]
Both terms have mean zero and are uncorrelated. The first has variance $4\theta^\top\Sigma\theta/m$. In the second, only matching unordered index pairs contribute to the variance, giving $2\operatorname{tr}(\Sigma^2)/\{m(m-1)\}$. The relative bound follows from $\theta^\top\Sigma\theta\le\ell r^2$ and $\operatorname{tr}(\Sigma^2)\le\ell\tau$. $\square$

Estimated centring means leave residual population means. Under independence, let these be $a$ and $c$ in the RNA and protein blocks, with marginal covariances $A$ and $D$. For row-wise vectorization of their product, the exact correction to part (a) is
\[
\Sigma_b=A\otimes D+(aa^\top)\otimes D+A\otimes(cc^\top).
\]
It follows by expanding each variable about its population mean. The additional terms vanish with exact centring; the product-noise covariance with estimated centring remains a function of the two marginal distributions.

Unpaired profiles give the independence noise shares and hence bounds on the oracle saving. The saving itself and the paired budget require the dependence shares, estimated from paired cells. Part (c) quantifies the precision of a pilot's raw block signals. Its second variance term remains positive at zero signal and controls precision for weak signals. For spherical product noise $\Sigma_b=\sigma^2 I_{d_b}$, achieving fixed relative precision in a weak signal can require order $\sigma^2\sqrt{d_b}/r_b^2$ paired cells. All pilot cells count towards the study's paired total. In the implementation, negative raw signals are set to zero before totals and shares are formed; these positive-part estimates and normalized shares are plug-in estimates, not unbiased ones.

\subsection*{Development}
The benchmark's structural test (G4, Supplementary Note~\ref{sec:gen}) compared a predicted saving with the saving against the paired-only envelope, in held-out units and for the mapped estimator. Proposition~\ref{prop:law} concerns the saving of block shrinkage in the unpaired bases over the same shrinkage without them, for the population's own dependence. After the benchmark had been scored, we measured that quantity in its nine data sets (post hoc): 25\% of cells were held out as truth cells, 30\% formed a reservoir of paired cells and 45\% an unpaired pool, and block and single-block James--Stein were fitted to 20 draws of reservoir cells per budget and scored by the noise-unbiased recovered fraction of the truth cells' dependence. With block moments of the reservoir, the law of Proposition~\ref{prop:law} was within \PredDevLawErrMax\% of the observed saving in all \PredDevN{} problems with an estimable saving at a recovered fraction of 0.5 (median error \PredDevLawErr\%; savings \PredDevSaveMin{} to \PredDevSaveMax, including panels of 50 and 800 genes), and a random orthonormal basis pair gave no saving, as the law predicts. The dependence concentrated in the leading blocks: averaged over the \PredProfileN{} data sets with recoverable dependence, RNA eigen-coordinates 1--5 carried \PredProfileTop\% of the signal and 1--20 carried \PredProfileTopTwenty\%. It overpredicted at a recovered fraction of 0.3, which the proposed estimator reaches with 25--50 cells, where its smallest blocks fall short of the oracle-linear risk. Development also fixed the pilot rule and a zero-pair signal profile.

\subsection*{Prospective test: protocol}
The plan, the code and the development records were registered before any of the \PredNData{} data sets below was extracted, the extracted files before any prediction, and all predictions of all data sets were recorded before any statistic of truth cells was computed. Data sets: mouse spleen and lymph node CITE-seq with 111 and 206 antibodies\cite{totalvi2021}; human blood and MALT lymphoma CITE-seq with 14 antibodies (10x Genomics, via ref.~\citenum{totalvi2021}); human bone-marrow CITE-seq with 25 antibodies\cite{stuart2019}; human fetal cortex multiome (RNA--chromatin)\cite{trevino2021}; adult mouse cortex SNARE-seq (RNA--chromatin)\cite{chen2019snare}; the optic-lobe and central-brain wiring of visual projection and centrifugal neurons in the FlyWire\cite{dorkenwald2024,schlegel2024} and BANC\cite{bates2026} connectomes; and human cortical GABAergic Patch-seq (RNA--electrophysiology)\cite{lee2023gaba}. The FlyWire and BANC files had been used before for other neurons. A problem was a data set with an X panel of 100, 200 or 400 features chosen from unpaired pool cells (a panel identical to a smaller one was dropped), giving \PredNProblems{} problems. Roles, estimand, estimators and scoring were those of development. The pilot was the first $\min(200,\,0.2N)$ of the $N$ reservoir cells, and budgets were $25\times2^k$ up to $\min(12{,}800,\,N-\text{pilot})$ with 20 draws. The observed saving was the ratio of cells needed by single-block and block James--Stein, interpolated as elsewhere; it was estimable when both curves crossed the target within the budgets from below. Predictions: the law with block moments of all reservoir cells (primary); the same from the pilot; a zero-pair transfer that combined the development signal-share profile with the problem's unpaired noise shares (Proposition~\ref{prop:reveal}a); and the unpaired bound $1/\min_b\pi_b$. Hypotheses at a recovered fraction of 0.5, evaluated if at least 8 problems were estimable: H1, median $|\log_2(\text{observed}/\text{law})|<\log_2 1.25$; H2, positive Spearman correlation of law and observation (one-sided permutation $P<0.05$); H3, median saving of the random-basis control in $[0.8,1.25]$; H4, pilot median error below $\log_2 1.5$; H5, transfer median error below $\log_2 1.5$; H6, every estimable saving within $[0.8,\,1.25/\min_b\pi_b]$. Two deviations were recorded before any prediction: the plan's cell cap of 3,000 per population was dropped, because it would have discarded most cells of the single-population data sets, and a file was renamed because its planned name clashed with that of a module the registered code uses.

\subsection*{Prospective test: results}
Of the \PredNProblems{} problems, \PredNEstimable, from \PredNEstimableData{} data sets, had an estimable saving at the primary target (Supplementary Table~\ref{tab:pred}). In the \PredNEasy{} CITE-seq data sets without cell-type conditions, block James--Stein exceeded the target with 25 paired cells, the smallest budget, because their pooled cross-covariance includes the differences between cell types; in the two chromatin data sets and the small Patch-seq data set, at least one curve did not reach it within the budgets. All six hypotheses were met: the law's median error was \PredLawErr\% (H1; 95\% interval from resampling data sets, \PredLawErrLo--\PredLawErrHi\%), with \PredLawWithin\% of problems within 25\%; its rank correlation with the observed saving was \PredRho{} (H2, $P\PredRhoRel$); the random-basis control saved \PredRandomObs{} (H3; law, \PredRandomLaw); the pilot's median error was \PredPilotErr\% (H4; \PredPilotErrLo--\PredPilotErrHi\%) and the zero-pair transfer's \PredTransferErr\% (H5; \PredTransferErrLo--\PredTransferErrHi\%, \PredTransferWithin\% of problems within 25\%); and no observed saving left the unpaired bounds (H6). The law overpredicted systematically: observed savings were a median \PredLawRatio{} times the predicted ones. At a recovered fraction of 0.7 the law's median error was \PredLawErrSeven\% (\PredNSeven{} problems), and at 0.3 it was \PredLawErrThree\% (\PredNThree{} problems; observed/predicted \PredLawRatioThree).

Secondary analyses, prespecified without thresholds, locate the bias. The law predicted the cells that single-block James--Stein needed at 0.5 closely (observed/predicted \PredCellsJsRatio) but underpredicted those that block James--Stein needed (\PredCellsBlocksRatio), the finite-sample shortfall of the smallest blocks that Proposition~\ref{prop:oracle} allows. A bootstrap version of the law, which resamples reservoir cells instead of assuming the oracle-linear risk, had median error \PredBootErr\% at 0.5 (observed/predicted \PredBootRatio), \PredBootErrSeven\% at 0.7 and \PredBootErrThree\% at 0.3. Exploratory checks of the censored problems agreed with the law in \PredCensAgree{} of \PredCensTotal{} comparisons of predicted cells with the censoring; the exceptions were two of the three SNARE-seq panels and one Patch-seq panel. In the three SNARE-seq panels block James--Stein needed \PredSnareObsMin--\PredSnareObsMax{} cells where the law predicted \PredSnareLawMin--\PredSnareLawMax, more than twice as many in two of them. Sparse chromatin counts have heavy-tailed products, for which the kurtosis term of Proposition~\ref{prop:oracle} is large.

\begin{table}[h]
\caption{\textbf{Prospective test of the saving law.} Savings at a recovered fraction of 0.5 of each population's dependence: cells needed by single-block James--Stein divided by cells needed by block James--Stein in the unpaired bases. Law: Proposition~\ref{prop:law} with block moments of the reservoir; pilot: the same from the pilot cells; no pairs: the development signal profile with unpaired noise shares; bound: $1/\min_b\pi_b$ from unpaired cells. n.e., not estimable (reason below each data set); $\geq$, single-block James--Stein did not reach the target. Cells: cells of eligible populations in all roles.}
\label{tab:pred}
\centering\scriptsize
\resizebox{\textwidth}{!}{% Generated by extension/predictability/make_assets.py. Do not edit.
\begin{tabular}{@{}llrrlllll@{}}
\toprule
Data set & Pair & Cells & Panel & Observed & Law & Pilot & No pairs & Bound \\
\midrule
Mouse spleen and lymph node, 111 proteins & RNA--protein & 14,577 & 100 & 5.38 & 6.76 & 7.57 & 7.41 & 92 \\
 &  &  & 200 & 8.33 & 9.70 & 8.77 & 12.09 & 151 \\
 &  &  & 400 & 11.57 & 14.55 & 13.40 & 18.85 & 233 \\
Mouse spleen and lymph node, 206 proteins & RNA--protein & 13,617 & 100 & 7.06 & 10.03 & 17.44 & 9.23 & 145 \\
 &  &  & 200 & 11.43 & 14.72 & 29.64 & 15.03 & 235 \\
 &  &  & 400 & $\geq$11.48 & 20.16 & 54.45 & 23.57 & 367 \\
Human blood, 14 proteins & RNA--protein & 6,855 & 100 & n.e. & 2.39 & 2.46 & 1.26 & 21 \\
 &  &  & 200 & n.e. & 3.35 & 3.43 & 1.62 & 28 \\
 &  &  & 400 & n.e. & 4.88 & 5.06 & 2.19 & 37 \\
\multicolumn{9}{@{}l}{\quad\textit{n.e.: both above target at 25 cells}} \\
Human MALT lymphoma, 14 proteins & RNA--protein & 6,838 & 100 & n.e. & 4.99 & 4.68 & 2.07 & 24 \\
 &  &  & 200 & n.e. & 5.90 & 5.34 & 2.46 & 36 \\
 &  &  & 400 & n.e. & 8.10 & 7.17 & 3.37 & 48 \\
\multicolumn{9}{@{}l}{\quad\textit{n.e.: both above target at 25 cells}} \\
Human bone marrow, 25 proteins & RNA--protein & 33,454 & 100 & n.e. & 4.77 & 4.63 & 1.65 & 14 \\
 &  &  & 200 & n.e. & 6.89 & 6.64 & 2.29 & 23 \\
 &  &  & 400 & n.e. & 10.36 & 9.92 & 3.26 & 34 \\
\multicolumn{9}{@{}l}{\quad\textit{n.e.: both above target at 25 cells}} \\
Human fetal cortex, RNA--chromatin & RNA--chromatin & 8,357 & 100 & n.e. & 3.29 & 49.15 & 9.57 & 187 \\
 &  &  & 200 & n.e. & 4.46 & 17.14 & 12.81 & 240 \\
 &  &  & 400 & n.e. & 5.81 & 32.96 & 17.96 & 354 \\
\multicolumn{9}{@{}l}{\quad\textit{n.e.: target not reached}} \\
Mouse cortex, RNA--chromatin & RNA--chromatin & 10,309 & 100 & $\geq$1.18 & 28.25 & 5.37 & 11.03 & 226 \\
 &  &  & 200 & $\geq$1.54 & 40.71 & 2.36 & 17.52 & 362 \\
 &  &  & 400 & $\geq$2.04 & 51.91 & 19.37 & 26.41 & 541 \\
\multicolumn{9}{@{}l}{\quad\textit{n.e.: paired-only not reached}} \\
Fly visual neurons, FlyWire & wiring & 7,101 & 100 & 1.87 & 2.34 & 2.17 & 1.57 & 56 \\
 &  &  & 200 & 2.49 & 2.93 & 2.71 & 1.85 & 22 \\
 &  &  & 302 & 3.09 & 3.48 & 3.17 & 2.13 & 24 \\
Fly visual neurons, BANC & wiring & 4,350 & 100 & 3.26 & 3.54 & 3.15 & 2.23 & 29 \\
 &  &  & 200 & 3.85 & 4.27 & 3.90 & 2.70 & 27 \\
Human cortex Patch-seq & RNA--electrophysiology & 613 & 100 & n.e. & 1.89 & 1.29 & 1.38 & 74 \\
 &  &  & 200 & n.e. & 1.79 & 1.66 & 1.49 & 26 \\
 &  &  & 400 & n.e. & 1.77 & 1.61 & 1.61 & 17 \\
\multicolumn{9}{@{}l}{\quad\textit{n.e.: target not reached}} \\
\bottomrule
\end{tabular}
}
\end{table}

\section{Learning without pairs in blood: protocol, identification analyses and complete results}\label{sec:cross}
The protocol was written and registered
before extraction, fitting and prediction, and before any recipient scoring
cell was read. It is retrospective with respect to the released paper. The
colon cohort had informed method development; the blood recipient cohort
\cite{hao2021} had not been used in this project.

\subsection*{Data and features}
\textbf{Training populations.} Stephenson et al.\ \cite{stephenson2021}
(E-MTAB-10026): samples without
lipopolysaccharide stimulation; cells annotated as platelets or red cells
removed; patients as units (\XsPatients{} patients, 13 with two samples
pooled); at most 2,000 cells per patient, chosen pseudo-randomly
(\XsTrainCells{} cells). Within each patient a second pseudo-random split divided the cells
into an RNA half and a protein half (\XsHalfCells{} RNA-half cells); the
pairing-free estimators saw only RNA moments of the first and protein moments
of the second. Paired comparators used all selected cells.

\textbf{Recipients.} The confirmatory blood cohort was the 3$'$ CITE-seq study of
Hao et al.\ \cite{hao2021} (GEO GSE164378; \XsHaoDonors{} donors, three
vaccination time points pooled per donor, \XsHaoCells{} cells,
\XsHaoScoring{} scoring cells). The colon cohort was the 12 donors of the
large-panel analysis with a new pseudo-random split (\XsColonScoring{} scoring cells).
Scoring halves were written to separate files at split time and read only
after the predictions had been recorded.

\textbf{Features.} Proteins: the \XsProteins{} antibody targets present in all
three panels, matched by specificity (isotype controls and three ambiguous
targets excluded; where the blood recipient cohort measured a target with two
clones, the clone annotated as human-specific, otherwise the first, was used),
as $\log(1+\text{count})$ minus its mean over the matched targets of the
cell. Genes: $\log(1+{}$counts per $10^4)$; the 200 genes with the highest
variance-to-mean ratio over the selected training cells (mean above 0.05),
among symbols shared by all three studies and detected in at least 1\% of the
adaptation cells of both recipient studies, plus six panel genes not already
selected (\XsGenes{} genes). No protein value, pairing, scoring cell or
prediction entered feature selection.

\subsection*{Arms and tuning}
Pairing-free: (1) ridge regression of patient protein means on patient RNA
means in pooled within-patient standard-deviation units, penalty scale
\XsPfScale{} times the trace of the between-patient RNA-mean covariance,
interior to its grid; (2) the Gaussian moment likelihood of
Proposition~\ref{prop:pairing} with $W$ of rank \XsPfRank{}, penalty $10^{-3}$
times that scale, rank chosen from $\{2,4,8,16,32\}$ by held-out patients'
likelihood; the final fit stopped after \XsPfMomentIter{} iterations with
gradient norm $\XsPfMomentGrad$. Paired: the profile-likelihood interaction
(ridge \XsPairedRidge{}, interior), reference regression (penalty
\XsRegLambda{} times the mean diagonal, interior), pooled within-patient
cross-correlations and cross-covariances, and the closed form with identity
marginals. Benchmark: ridge regression on the recipient's own adaptation
pairs, tuned by three folds of adaptation cells. All tuned arms used a
five-value grid, three patient-grouped folds (patients assigned pseudo-randomly) and the
edge rule of the protocol.

\textbf{Amendment.} Before any prediction was computed, the rank search of
the moment estimator was restricted to its five-value grid without the edge
extension, because extending it to full rank (122) would have required three
cross-validation fits of more than an hour each on the available two cores;
the within-colon re-analysis kept the rank of its first version (8). The
information available at the time was the tuning outcome of the other arms and
the fact that the rank search had reached the edge of its grid. The chosen rank
(\XsPfRank) is therefore at the edge. The final fit stopped when the line
search could not decrease the objective further; three restarts from that point
changed the interaction by less than $\XsRestartBChange$ (relative) and no
recipient's endpoint by more than $\XsRestartEChange$.

\subsection*{Identifiability computations}
The effective supported subspace $S$ was spanned by the eigenvectors of the
between-patient RNA-mean matrix $G$ (pooled within-patient standard-deviation
units) whose ridge factor $g/(g+\lambda)$ was at least one half, with
$\lambda$ the fitted pairing-free penalty. $S$ is contained in the span $M$ of
Proposition~\ref{prop:means} and depends on the penalty. The share of the
recipient's standardized RNA variance inside $S$ is
$\operatorname{tr}(P_SR_xP_S)/\operatorname{tr}R_x$ for its adaptation
correlation matrix $R_x$ (the analyses of this note use measured $R_x$;
Supplementary Note~\ref{sec:noise} repeats them with its latent part). The
unsupported share is $\|R_x^{\mathrm{s}}(I-P_S)(R_x^{\mathrm{s}})^{-1}T\|_F/\|T\|_F$,
with $R_x^{\mathrm{s}}$ the shrunk RNA correlation matrix of the scoring cells
and $T$ their cross-correlation. The set of Proposition~\ref{prop:means}(4) was
computed with the adaptation correlation matrices and the fitted $W$ on $S$.
With measured RNA covariance, its whitened block $F$ had largest singular value
between \XsMaxFMin{} and \XsMaxFMax. It exceeded one in \XsMaxFAboveOne{} of
\XsIdUnits{} recipient analyses, including every total analysis (minimum
\XsMaxFTotalMin), so the exact set was empty there. Supplementary
Note~\ref{sec:noise} shows that RNA sampling noise explains most of this and
separates the remaining incompatibility from sampling uncertainty. The radii
below cap the singular values of $F$ at one, which is the same as clipping the
negative eigenvalues of $I-F^\top F$ at zero; they describe a repaired set, not
the identified one, and are kept for comparison with the earlier version. The
per-gene analysis (each gene's row of the cross-correlation matrix) was added
after scoring and is descriptive. In the per-gene analysis, the unsupported
share of each gene's relationships had Spearman correlation
\XsGeneSpearmanUnidHao{} with its prediction error in the confirmatory cohort and
\XsGeneSpearmanUnidColon{} in colon, whereas the share of the gene's own axis
inside $S$ had \XsGeneSpearmanAxisHao{} and \XsGeneSpearmanAxisColon. The
unsupported part of the target was \XsHaoIdFloor{} of its norm (range
\XsHaoIdFloorMin--\XsHaoIdFloorMax) in the confirmatory cohort's total analysis,
\XsHaoWIdFloor{} within its level-1 types and \XsColonIdFloor{} in colon; the
capped radius was \XsHaoIdRadius, \XsHaoWIdRadius{} and \XsColonIdRadius{} of
the norm of the truth, respectively.

\subsection*{Within-colon re-analysis of pairing-free learning}
This re-analysis replaces an earlier version of the within-colon analysis, with the same recipients, genes, proteins, units,
halves and endpoint. Each training biopsy's RNA moments came from one pseudo-random half
of its cells and its protein moments from the other. Folds for tuning were
donor-grouped. The ridge penalty was chosen from $\{10^{-4},\dots,1\}$ times the
trace of the between-biopsy RNA-mean covariance with the edge rule. The moment
estimator kept rank 8 and was run for up to 20,000 L-BFGS-B iterations
(gradient tolerance $10^{-7}$); within this limit, \XsRigConverged{} of 12 fits converged. The
other \XsRigContinued{} were continued post hoc from their final state for up to
20,000 further iterations, after
which \XsRigContinueConverged{} of them had converged; recipient errors changed by at
most \XsRigContinueMaxChange{} and interactions by at most \XsRigContinueMaxWChange{}
(relative Frobenius norm). Supplementary Table~\ref{tab:rigor} reports the errors
at the limit. The permutation control
reassigned protein moments between biopsies of different donors only, 200
times for the primary estimator and 20 times for the moment estimator, with
the same fitting rules. Supplementary Table~\ref{tab:rigor} gives every recipient.

\begin{table}[h]
\centering
\caption{\textbf{Within-colon pairing-free learning, corrected re-analysis.}
Relative cross-covariance error $E$ for each recipient, learned from the other
donors' biopsies without paired cells; $p$ is the donor-aware permutation
$p$-value (200 permutations for means, 20 for rank 8); $\dim M$ is the number
of gene directions the ridge estimator kept by at least one half. An asterisk
marks a rank-8 fit that stopped without meeting the convergence criterion.}
\label{tab:rigor}
\small
% Generated by extension/cross_study/make_assets.py. Do not edit.
\begin{tabular}{@{}lrrrrrrr@{}}
\toprule
Recipient & Biopsies & Penalty & Means $E$ & $p$ & Rank-8 $E$ & $p$ & $\dim M$\\
\midrule
HS1 & 19 & 0.01 & 0.612 & 0.005 & 0.619 & 0.048 & 11\\
HS2 & 19 & 0.01 & 0.662 & 0.005 & 0.701 & 0.048 & 11\\
HS3 & 19 & 0.01 & 0.470 & 0.005 & 0.632 & 0.048 & 12\\
HS4 & 19 & 0.01 & 0.422 & 0.005 & 0.940 & 0.143 & 11\\
HS5 & 19 & 0.01 & 0.399 & 0.005 & 0.547$^*$ & 0.048 & 11\\
HS6 & 19 & 0.01 & 0.588 & 0.005 & 0.624 & 0.048 & 11\\
HS7 & 19 & 0.01 & 0.637 & 0.005 & 0.696$^*$ & 0.048 & 11\\
HS8 & 20 & 0.01 & 0.555 & 0.005 & 0.478$^*$ & 0.048 & 11\\
HS9 & 19 & 0.01 & 0.468 & 0.005 & 0.763 & 0.095 & 11\\
HS10 & 20 & 0.01 & 0.827 & 0.005 & 0.806 & 0.048 & 12\\
HS11 & 19 & 0.01 & 0.579 & 0.005 & 0.617 & 0.048 & 11\\
HS12 & 20 & 0.01 & 0.574 & 0.005 & 0.734 & 0.048 & 12\\
\bottomrule
\end{tabular}

\end{table}

\subsection*{Complete results}
Supplementary Table~\ref{tab:cross} gives every arm's donor-mean error. The null of H1 kept
each training site's patients but paired their RNA and protein moments at
random within site; it therefore retained any site-level batch alignment,
which makes it conservative.

\begin{table}[h]
\centering
\caption{\textbf{Cross-study transfer: donor-mean relative error of the predicted
cross-correlation matrix.} Interactions were learned in the training study and
applied unchanged; the own paired half uses the recipient's adaptation pairs,
which the setting does not permit. Type columns centre every cell by its
cell-type mean in both assays. The last column is the mean correlation between
predicted and held-out entries in colon (all cells).}
\label{tab:cross}
\small
% Generated by extension/cross_study/make_assets.py. Do not edit.
\begin{tabular}{@{}lcccccc@{}}
\toprule
 & \multicolumn{3}{c}{Confirmatory blood cohort (8 donors)} & \multicolumn{3}{c}{Colon (12 donors)}\\
\cmidrule(lr){2-4}\cmidrule(l){5-7}
Arm & All cells & 8 types & 31 types & All cells & 14 types & Entry corr.\\
\midrule
No pairs: means (primary) & 0.292 & 1.082 & 1.484 & 0.988 & 1.274 & 0.40\\
No pairs: moments & 0.389 & 1.038 & 1.435 & 0.907 & 1.227 & 0.46\\
Paired closed form & 0.134 & 0.528 & 0.868 & 0.719 & 1.000 & 0.69\\
Reference regression & 0.522 & 0.680 & 0.842 & 0.859 & 0.983 & 0.52\\
Transferred correlation & 0.630 & 2.212 & 4.044 & 1.132 & 2.225 & 0.42\\
Transferred covariance & 0.951 & 8.960 & 18.953 & 4.651 & 15.642 & 0.23\\
Variances only & 0.976 & 0.901 & 0.915 & 0.971 & 0.983 & 0.28\\
Own paired half & 0.120 & 0.348 & 0.541 & 0.539 & 0.936 & 0.84\\
Independence & 1.000 & 1.000 & 1.000 & 1.000 & 1.000 & 0.00\\
\bottomrule
\end{tabular}

\end{table}

\subsection*{Post hoc analyses}
Two exploratory analyses were appended to the protocol after scoring; their
text and every version of the
protocol are kept with the code. Neither changes any confirmatory result.

\textbf{Between-type and within-type parts (P2).} For each recipient donor, the
scoring cells of the kept level-1 (blood) or coarse (colon) types were used to
compute the between-type part of the cross-covariance (type means weighted by
type proportions) and the within-type part (pooled within-type
cross-covariance); their sum is the total cross-covariance of those cells. Both
parts were divided by the scoring cells' total standard deviations. The prespecified
primary channel $W$ predicted each part from the donor's adaptation cells as
the corresponding RNA covariance part, in the same units, times $W^\top$, with
no saturation, so the two predicted parts also add to the direct-channel
prediction of the total. Supplementary Table~\ref{tab:mixture} reports the relative
Frobenius norm of each part of the truth, the relative error and the entry
correlation of each predicted part, and, as a descriptive addition computed
after P2 and not part of its plan, the share
$\operatorname{tr}(P_MR P_M)/\operatorname{tr}R$ of each RNA covariance part $R$
inside $M$. The direct channel's prediction of the within-type part had
\XsMixHaoNormWithin{} times the norm of the truth in the confirmatory cohort. The
direct channel without saturation is cruder than the closed form: for the total
over the kept types its error was \XsMixHaoErrDirect{} in the confirmatory cohort,
against \XsHaoErrPfMeans{} for the prespecified closed-form prediction of all cells.

\begin{table}[h]
\centering
\caption{\textbf{Between-type and within-type parts of the cross-covariance.}
Donor means. Norm share, Frobenius norm of each part of the scoring cells'
cross-covariance relative to the total (the parts are not orthogonal, so the
shares need not sum to one). Relative error and entry correlation, the prespecified
pairing-free channel's prediction of each part. Share inside $S$, fraction of
the corresponding RNA covariance of the adaptation cells inside the effective
supported subspace (descriptive addition).}
\label{tab:mixture}
\small
% Generated by extension/cross_study/make_assets.py. Do not edit.
\begin{tabular}{@{}lcccccccc@{}}
\toprule
 & \multicolumn{2}{c}{Norm share} & \multicolumn{2}{c}{Relative error} & \multicolumn{2}{c}{Entry correlation} & \multicolumn{2}{c}{Share inside $S$}\\
\cmidrule(lr){2-3}\cmidrule(lr){4-5}\cmidrule(lr){6-7}\cmidrule(l){8-9}
Cohort & between & within & between & within & between & within & between & within\\
\midrule
Confirmatory blood (8 donors) & 0.97 & 0.11 & 0.71 & 1.92 & 0.75 & 0.15 & 0.88 & 0.12\\
Colon (12 donors) & 0.81 & 0.43 & 1.14 & 1.38 & 0.35 & 0.07 & 0.57 & 0.09\\
\bottomrule
\end{tabular}

\end{table}

\textbf{Training on variation within cell types (P1).} Training units were the
training donors' cell types in the training study's initial clustering, with at least 50 cells in each
pseudo-random half (\XsPoneUnits{} units). RNA moments came from the RNA half and
protein moments from the protein half of each unit, and each unit's means were
centred on the cell-weighted mean of its type across donors, leaving only
variation among donors within types. The primary estimator was refitted by its
prespecified rules, with donor-grouped folds. The supported subspace had \XsPoneDim{}
directions, against \XsDimM{} for the prespecified fit, and the unsupported part of
the level-1 within-type truth fell to \XsPoneFloor{} of its norm (from
\XsHaoWIdFloor). The within-type error was \XsPoneErr{} at level 1 and
\XsPoneErrLTwo{} at level 2, against \XsHaoWErrPfMeans{} and
\XsHaoLTwoErrPfMeans{} for the prespecified fit, and the total error was
\XsPoneErrTotal. The training labels are the original authors' annotations of
the training study, which used both assays; the analysis therefore assumes
that each assay carries its own cell-type labels, as annotated atlases do.

\textbf{Gene--protein pairs.} Genes whose protein product was measured were
matched to their antibody targets by gene symbol (21 pairs among the prespecified
features). Their donor-mean held-out and predicted within-type correlations
(level-1 types, confirmatory cohort) are descriptive and were computed after
scoring; for FCGR3A and CD16 they were \XsPairFcgraWTruth{} held out,
\XsPairFcgraWPaired{} for the paired closed form, \XsPairFcgraWBench{} for the
recipient's own paired half and \XsPairFcgraWPfMeans{} for the primary
pairing-free channel; for IGHD and IgD, \XsPairIghdWTruth, \XsPairIghdWPaired,
\XsPairIghdWBench{} and \XsPairIghdWPfMeans; for KLRB1 and CD161,
\XsPairKlrbWTruth, \XsPairKlrbWPaired, \XsPairKlrbWBench{} and
\XsPairKlrbWPfMeans; and for IL7R and CD127, \XsPairIlrWTruth,
\XsPairIlrWPaired, \XsPairIlrWBench{} and \XsPairIlrWPfMeans. Supplementary Fig.~\ref{sfig:entries} shows held-out against
predicted entries of the cross-correlation matrix between and within cell types.

\begin{figure}[htbp]
\centering
\resizebox{\textwidth}{!}{% Generated by extension/cross_study/make_assets.py. Do not edit.
\definecolor{donorcol}{HTML}{5D6670}
\definecolor{haocol}{HTML}{217D91}
\definecolor{coloncol}{HTML}{A64442}
\begin{tikzpicture}[font=\sffamily]
\begin{axis}[name=a,width=.34\textwidth,height=.34\textwidth,xmin=-0.7,xmax=1.0,ymin=-0.7,ymax=1.0,
 xlabel={Held-out correlation},ylabel={Predicted correlation},title={\textbf{a}\enspace All cells, no pairs},
 axis x line*=bottom,axis y line*=left,tick align=outside,xmajorgrids=true,grid style={gray!18},axis line style={gray!55},tick label style={font=\sffamily\scriptsize},label style={font=\sffamily\small},clip=true,title style={font=\sffamily\small,at={(0,1)},anchor=base west,yshift=5pt}]
\draw[black!35] (axis cs:-0.7,-0.7) -- (axis cs:1.0,1.0);
\addplot[only marks,mark=*,mark size=.35pt,color=donorcol!45] coordinates {
 (-0.069,-0.040) (-0.061,-0.072) (-0.045,-0.035) (-0.176,-0.167) (-0.065,-0.023) (0.126,0.139) (0.120,0.106) (-0.062,-0.060) (-0.061,-0.043) (0.021,0.017) (-0.056,-0.037) (0.016,0.022)
 (0.154,0.095) (0.116,0.009) (0.007,0.005) (0.008,-0.008) (0.402,0.389) (0.598,0.487) (0.271,0.208) (0.505,0.439) (0.425,0.393) (0.380,0.399) (-0.484,-0.445) (-0.025,-0.020)
 (-0.092,-0.096) (-0.028,-0.045) (0.245,0.124) (-0.027,-0.026) (0.124,0.095) (0.130,0.019) (-0.029,-0.016) (-0.142,-0.097) (0.338,0.278) (-0.046,-0.040) (0.010,0.010) (-0.063,-0.024)
 (-0.161,-0.215) (0.042,0.051) (0.164,0.135) (-0.155,-0.111) (-0.179,-0.057) (0.296,0.156) (-0.414,-0.410) (0.627,0.505) (-0.218,-0.319) (0.137,0.047) (-0.228,-0.195) (-0.191,-0.176)
 (-0.193,-0.172) (-0.568,-0.495) (-0.381,-0.338) (0.184,0.131) (-0.264,-0.194) (-0.148,-0.141) (-0.012,0.193) (-0.083,-0.145) (0.297,0.246) (-0.051,-0.082) (-0.020,0.042) (0.090,-0.026)
 (-0.004,0.018) (-0.247,-0.200) (-0.234,-0.190) (-0.129,-0.097) (-0.659,-0.543) (-0.206,-0.275) (0.596,0.550) (0.520,0.340) (-0.236,-0.100) (-0.056,-0.054) (-0.013,-0.029) (-0.015,0.000)
 (-0.009,-0.023) (0.007,0.036) (0.112,0.111) (0.016,-0.011) (-0.025,0.010) (-0.037,-0.016) (-0.049,-0.038) (-0.020,0.031) (-0.040,-0.004) (-0.042,-0.036) (-0.051,-0.049) (0.065,0.060)
 (-0.093,-0.033) (-0.236,-0.233) (-0.085,-0.086) (0.724,0.511) (0.122,0.048) (0.366,0.289) (0.195,0.130) (-0.063,-0.060) (-0.179,-0.099) (-0.116,-0.089) (0.726,0.516) (-0.038,-0.020)
 (-0.172,-0.073) (-0.000,0.172) (-0.086,-0.054) (-0.089,-0.050) (-0.091,-0.126) (-0.122,-0.132) (0.203,0.236) (-0.272,-0.271) (0.400,0.391) (-0.109,-0.131) (0.035,0.062) (0.109,0.110)
 (0.402,0.312) (0.063,0.061) (-0.209,-0.167) (0.512,0.341) (-0.295,-0.252) (0.012,0.014) (0.120,0.098) (-0.020,0.235) (-0.104,-0.113) (0.432,0.263) (-0.051,-0.153) (-0.032,0.017)
 (0.027,-0.009) (-0.003,0.027) (-0.028,0.019) (-0.033,-0.010) (-0.023,-0.018) (-0.033,0.003) (-0.034,-0.052) (0.168,0.073) (-0.093,-0.031) (0.328,0.232) (-0.112,-0.122) (-0.023,-0.012)
 (-0.019,-0.031) (-0.025,-0.019) (-0.004,0.045) (0.141,0.181) (-0.000,-0.004) (-0.030,-0.057) (-0.128,-0.080) (-0.346,-0.315) (-0.221,-0.097) (-0.143,-0.095) (-0.145,-0.116) (-0.009,-0.021)
 (-0.109,-0.066) (-0.155,-0.116) (-0.213,-0.131) (0.010,0.069) (-0.113,-0.092) (-0.009,-0.086) (-0.008,-0.007) (-0.227,-0.145) (0.096,0.011) (-0.159,-0.167) (-0.014,-0.028) (-0.077,-0.085)
 (-0.015,-0.026) (0.018,0.020) (-0.014,0.064) (-0.027,-0.051) (-0.061,-0.079) (-0.093,-0.119) (-0.130,-0.146) (0.209,0.229) (-0.278,-0.284) (0.438,0.415) (-0.098,-0.145) (-0.009,0.032)
 (0.143,0.068) (-0.008,0.025) (0.039,0.036) (0.136,0.082) (0.214,0.130) (0.123,0.016) (-0.005,0.058) (0.062,0.043) (0.020,0.016) (-0.031,0.001) (-0.045,-0.069) (0.004,-0.096)
 (0.028,0.079) (-0.055,-0.150) (0.027,0.043) (-0.144,-0.134) (-0.143,-0.156) (-0.060,-0.083) (-0.408,-0.394) (-0.147,-0.134) (0.349,0.342) (0.281,0.321) (-0.134,-0.121) (-0.310,-0.209)
 (0.102,0.078) (-0.218,-0.128) (0.134,0.089) (0.317,0.279) (0.154,0.176) (0.063,0.058) (0.090,0.171) (0.368,0.379) (0.574,0.470) (0.225,0.221) (0.454,0.439) (0.400,0.378)
 (0.356,0.385) (-0.445,-0.430) (0.378,0.377) (0.769,0.637) (-0.099,-0.106) (-0.200,-0.123) (0.486,0.487) (-0.195,-0.136) (0.319,0.353) (-0.280,-0.180) (-0.491,-0.451) (-0.334,-0.277)
 (-0.268,-0.138) (-0.056,-0.039) (0.157,0.092) (-0.033,-0.150) (-0.195,-0.173) (-0.340,-0.283) (-0.069,-0.059) (0.246,0.201) (-0.291,-0.200) (-0.222,-0.139) (-0.207,-0.215) (-0.107,-0.077)
 (-0.132,-0.050) (-0.105,-0.112) (-0.022,-0.001) (-0.080,-0.093) (-0.276,-0.294) (-0.090,-0.114) (-0.024,0.043) (-0.108,-0.133) (-0.049,-0.062) (-0.035,0.027) (-0.118,-0.092) (0.594,0.429)
 (-0.082,-0.214) (-0.049,-0.037) (-0.008,0.015) (-0.021,0.010) (-0.001,-0.017) (-0.010,-0.044) (0.006,-0.003) (-0.005,-0.063) (-0.015,-0.078) (-0.040,-0.013) (0.063,0.165) (-0.085,-0.128)
 (0.557,0.421) (-0.232,-0.184) (0.402,0.326) (-0.244,-0.195) (-0.290,-0.210) (-0.264,-0.177) (-0.219,-0.176) (-0.206,-0.297) (-0.154,-0.090) (0.140,0.084) (0.002,0.073) (-0.140,-0.042)
 (-0.087,-0.058) (-0.126,-0.099) (0.059,0.041) (-0.133,-0.090) (-0.176,-0.103) (-0.025,0.052) (-0.095,-0.089) (-0.001,-0.055) (-0.006,-0.014) (-0.221,-0.129) (0.080,0.003) (-0.166,-0.142)
 (0.237,0.211) (-0.036,-0.061) (0.011,0.008) (-0.053,0.016) (-0.162,-0.203) (0.051,0.039) (0.143,0.124) (-0.054,-0.071) (-0.072,-0.091) (0.105,0.162) (-0.143,-0.158) (0.225,0.235)
 (-0.071,-0.057) (0.003,0.038) (-0.147,-0.123) (-0.083,-0.069) (-0.156,-0.105) (-0.426,-0.371) (-0.202,-0.206) (0.032,0.041) (-0.217,-0.165) (-0.103,-0.082) (-0.115,-0.134) (0.346,0.318)
 (-0.178,-0.042) (-0.319,-0.321) (-0.169,-0.128) (0.164,0.137) (-0.193,-0.162) (-0.090,-0.096) (-0.082,-0.101) (-0.045,-0.052) (-0.251,-0.261) (-0.091,-0.099) (0.205,0.209) (0.180,0.223)
 (-0.085,-0.083) (0.764,0.613) (-0.220,-0.171) (0.354,0.362) (-0.209,-0.183) (-0.298,-0.228) (-0.443,-0.364) (-0.219,-0.179) (-0.167,-0.321) (-0.174,-0.117) (0.156,0.068) (-0.015,0.078)
 (-0.143,-0.062) (-0.121,-0.089) (-0.125,-0.117) (0.115,0.072) (0.350,0.298) (0.681,0.575) (-0.098,-0.055) (-0.282,-0.131) (0.576,0.483) (-0.207,-0.117) (0.296,0.305) (-0.267,-0.280)
 (0.203,0.139) (-0.015,0.060) (-0.034,-0.076) (0.032,0.057) (-0.031,0.001) (-0.103,-0.106) (0.092,0.084) (0.107,0.133) (0.018,0.023) (0.014,0.004) (0.088,0.099) (0.084,0.004)
 (-0.102,-0.063) (0.211,0.366) (0.064,0.097) (-0.073,-0.053) (-0.054,-0.022) (-0.047,-0.038) (-0.086,-0.034) (-0.141,-0.073) (-0.120,-0.150) (0.340,0.268) (-0.056,-0.032) (-0.114,-0.134)
 (0.349,0.311) (-0.123,-0.031) (-0.312,-0.348) (-0.179,-0.149) (0.138,0.121) (-0.210,-0.183) (0.028,0.019) (0.026,0.052) (0.018,0.029) (0.084,0.098) (0.025,0.084) (-0.098,-0.097)
 (-0.039,-0.043) (-0.021,-0.014) (0.570,0.485) (-0.165,-0.132) (0.255,0.277) (-0.161,-0.147) (-0.222,-0.158) (-0.326,-0.264) (-0.169,-0.140) (-0.131,-0.215) (0.450,0.329) (0.620,0.486)
 (0.313,0.115) (0.501,0.348) (0.468,0.379) (0.451,0.369) (-0.506,-0.410) (-0.125,-0.100) (-0.175,-0.109) (0.021,0.055) (-0.089,-0.057) (0.002,-0.049) (-0.011,-0.006) (-0.176,-0.118)
 (0.080,-0.045) (-0.066,0.003) (0.220,0.175) (-0.024,-0.007) (0.003,-0.002) (-0.038,-0.030) (-0.091,-0.065) (0.012,0.009) (0.093,0.065) (-0.192,-0.157) (-0.238,-0.147) (0.403,0.267)
 (-0.544,-0.496) (0.825,0.678) (-0.282,-0.294) (0.184,0.073) (0.118,0.120) (0.371,0.294) (0.070,0.067) (-0.159,-0.128) (0.480,0.330) (-0.249,-0.229) (0.018,0.026) (0.119,0.100)
 (0.047,0.137) (-0.079,-0.065) (-0.113,-0.103) (0.141,0.019) (0.060,0.076) (0.037,-0.071) (0.056,0.034) (0.102,0.061) (0.086,0.062) (0.053,0.043) (0.266,0.202) (0.078,0.115)
 (-0.338,-0.258) (-0.118,-0.052) (-0.065,-0.051) (-0.126,-0.118) (-0.036,-0.028) (0.040,0.007) (-0.070,-0.037) (-0.042,-0.025) (0.114,0.137) (-0.026,-0.016) (-0.052,-0.022) (0.114,0.010)
 (0.266,0.158) (0.144,-0.079) (0.115,-0.024) (0.112,0.077) (0.136,0.046) (-0.161,-0.091) (-0.075,-0.025) (-0.205,-0.206) (-0.085,-0.089) (0.697,0.490) (0.116,0.060) (0.349,0.275)
 (0.222,0.138) (-0.073,-0.077) (-0.062,-0.030) (0.073,0.114) (-0.017,0.002) (0.001,0.030) (-0.027,-0.019) (-0.093,-0.190) (0.003,0.034) (0.066,0.071) (-0.142,-0.103) (0.087,0.109)
 (-0.070,-0.066) (-0.397,-0.299) (0.200,0.108) (-0.246,-0.170) (-0.005,-0.018) (0.106,0.067) (0.165,0.160) (0.030,0.029) (-0.084,-0.037) (0.306,0.155) (-0.118,-0.101) (0.008,0.040)
 (0.059,0.041) (-0.128,-0.131) (0.329,0.279) (-0.068,0.032) (-0.313,-0.335) (-0.185,-0.146) (0.229,0.169) (-0.202,-0.169) (-0.048,-0.027) (-0.046,-0.035) (-0.032,-0.046) (-0.036,-0.036)
 (-0.054,-0.079) (0.265,0.157) (-0.161,-0.024) (0.765,0.420) (-0.058,-0.008) (0.025,0.015) (-0.042,-0.040) (0.044,0.029) (0.062,0.059) (0.041,0.014) (0.011,0.001) (0.012,-0.014)
 (0.406,0.366) (0.564,0.442) (0.229,0.181) (0.467,0.415) (0.437,0.371) (0.401,0.385) (-0.456,-0.423) (-0.179,-0.140) (-0.233,-0.154) (0.063,0.090) (-0.111,-0.098) (-0.062,-0.120)
 (-0.004,-0.020) (-0.293,-0.199) (0.112,0.075) (-0.397,-0.369) (-0.275,-0.249) (-0.211,-0.140) (-0.038,-0.026) (0.128,0.104) (-0.088,-0.063) (-0.159,-0.142) (-0.275,-0.236) (-0.144,-0.126)
 (-0.173,-0.122) (0.273,0.218) (-0.402,-0.384) (0.638,0.559) (-0.211,-0.306) (0.050,0.039) (-0.074,-0.083) (-0.073,-0.029) (-0.076,-0.064) (-0.176,-0.157) (-0.126,-0.104) (0.058,0.099)
 (-0.089,-0.077) (-0.046,-0.045) (0.003,-0.137) (-0.018,0.009) (0.002,0.007) (0.021,-0.003) (0.011,-0.010) (-0.030,0.096) (0.013,0.002) (-0.165,-0.152) (-0.163,-0.166) (-0.076,-0.087)
 (-0.470,-0.433) (-0.170,-0.126) (0.402,0.386) (0.322,0.317) (-0.165,-0.137) (0.130,0.017) (-0.069,-0.040) (0.075,0.059) (-0.056,-0.030) (-0.080,-0.036) (-0.005,0.027) (-0.065,-0.025)
 (-0.075,-0.268) (-0.152,-0.090) (-0.366,-0.325) (-0.259,-0.099) (-0.176,-0.096) (-0.149,-0.125) (-0.017,-0.025) (-0.125,-0.072) (0.288,0.184) (0.494,0.392) (-0.091,-0.053) (-0.138,-0.023)
 (0.386,0.314) (-0.136,-0.053) (0.253,0.242) (-0.209,-0.151) (-0.377,-0.376) (-0.262,-0.236) (-0.213,-0.151) (-0.035,-0.034) (0.120,0.108) (-0.047,-0.010) (-0.148,-0.144) (-0.274,-0.238)
 (0.024,0.044) (0.035,0.026) (0.451,0.291) (-0.202,-0.094) (-0.215,-0.171) (0.292,0.244) (0.242,0.133) (0.074,0.081) (0.321,0.227) (0.040,0.040) (-0.112,-0.073) (0.368,0.202)
 (-0.242,-0.152) (0.022,0.038) (0.146,0.068) (-0.072,0.035) (0.263,0.190) (-0.084,-0.044) (-0.191,-0.247) (-0.107,-0.087) (0.057,0.086) (-0.128,-0.136) (0.036,0.029) (0.034,0.002)
 (0.025,-0.013) (0.108,0.057) (0.030,-0.028) (-0.116,-0.097) (-0.041,0.059) (-0.026,-0.031) (-0.056,-0.068) (0.018,0.036) (-0.015,-0.037) (0.006,0.017) (0.189,0.127) (0.062,0.066)
 (0.001,0.013) (0.015,0.052) (0.462,0.384) (0.601,0.468) (0.270,0.142) (0.523,0.404) (0.469,0.401) (0.468,0.411) (-0.494,-0.437) (0.205,0.051) (0.044,0.066) (-0.014,-0.019)
 (-0.023,-0.050) (-0.085,0.027) (0.018,-0.011) (0.123,-0.011) (-0.038,-0.050) (-0.469,-0.391) (-0.309,-0.240) (0.121,0.132) (-0.074,-0.017) (0.053,0.039) (-0.053,-0.047) (-0.185,-0.138)
 (-0.285,-0.245) (-0.193,-0.165) (0.018,0.031) (0.092,0.099) (-0.568,-0.485) (0.483,0.386) (-0.286,-0.262) (0.016,0.041) (-0.018,-0.003) (-0.020,0.022) (-0.013,-0.011) (-0.017,0.023)
 (-0.055,-0.016) (-0.021,-0.100) (0.063,0.063) (-0.016,-0.022) (-0.140,0.005) (0.207,0.132) (0.356,0.336) (-0.303,-0.404) (-0.197,-0.127) (0.316,0.172) (-0.179,-0.132) (-0.009,0.027)
 (-0.007,0.012) (-0.006,0.004) (0.012,0.016) (-0.004,-0.050) (0.064,0.032) (-0.062,-0.063) (0.174,0.120) (0.535,0.426) (-0.148,-0.125) (0.259,0.258) (-0.153,-0.135) (-0.184,-0.151)
 (-0.309,-0.242) (-0.154,-0.127) (-0.112,-0.227) (-0.166,-0.101) (-0.399,-0.352) (-0.285,-0.115) (-0.194,-0.109) (-0.159,-0.139) (-0.016,-0.026) (-0.126,-0.071) (0.321,0.206) (0.508,0.405)
 (-0.107,-0.050) (-0.139,-0.010) (0.456,0.366) (-0.127,-0.039) (0.346,0.268) (-0.237,-0.226) (-0.478,-0.433) (-0.326,-0.280) (-0.255,-0.148) (-0.047,-0.032) (0.165,0.116) (-0.076,-0.088)
 (-0.190,-0.172) (-0.329,-0.280) (-0.053,-0.053) (0.200,0.155) (-0.225,-0.155) (-0.165,-0.116) (-0.153,-0.158) (-0.090,-0.043) (-0.097,-0.039) (-0.042,-0.018) (-0.048,-0.006) (-0.044,-0.021)
 (-0.062,0.005) (-0.123,-0.023) (-0.064,-0.111) (0.243,0.184) (-0.051,-0.020) (-0.117,-0.133) (0.370,0.316) (-0.109,-0.004) (-0.306,-0.346) (-0.176,-0.156) (0.166,0.136) (-0.203,-0.180)
 (0.005,-0.001) (0.010,0.026) (-0.003,0.002) (0.039,0.054) (0.130,0.002) (-0.021,-0.022) (-0.066,-0.060) (0.080,0.115) (0.056,-0.029) (-0.042,-0.000) (0.068,0.040) (-0.035,0.013)
 (-0.051,0.000) (-0.009,-0.032) (-0.025,0.002) (-0.035,-0.098) (0.315,0.258) (0.403,0.300) (0.177,0.088) (0.343,0.277) (0.330,0.263) (0.312,0.275) (-0.331,-0.287) (-0.061,-0.053)
 (-0.136,-0.139) (-0.054,-0.061) (0.420,0.320) (0.054,0.036) (0.216,0.177) (0.100,0.080) (-0.039,-0.041) (-0.120,-0.077) (-0.074,-0.065) (0.469,0.351) (-0.028,-0.028) (-0.113,-0.046)
 (0.013,0.113) (-0.054,-0.050) (-0.066,-0.043) (0.072,0.085) (0.037,0.035) (0.366,0.259) (0.043,0.011) (-0.337,-0.258) (0.364,0.378) (0.196,0.131) (-0.027,-0.031) (-0.063,-0.029)
 (-0.039,-0.030) (-0.073,0.004) (-0.147,-0.059) (-0.089,-0.111) (0.281,0.225) (-0.041,-0.027) (-0.076,-0.047) (0.226,0.146) (-0.102,-0.078) (-0.198,-0.310) (-0.118,-0.109) (0.062,0.080)
 (-0.137,-0.138) (0.048,0.009) (0.042,0.027) (0.030,0.017) (0.140,0.086) (0.033,0.027) (-0.154,-0.122) (-0.037,0.014) (-0.028,-0.023) (-0.161,-0.160) (-0.025,0.014) (0.012,-0.027)
 (-0.052,-0.014) (-0.028,-0.001) (0.162,0.168) (0.004,0.015) (-0.046,-0.022) (0.431,0.319) (0.540,0.413) (0.265,0.063) (0.484,0.320) (0.419,0.348) (0.443,0.360) (-0.438,-0.383)
 (0.346,0.292) (0.602,0.509) (-0.095,-0.080) (-0.173,-0.078) (0.452,0.410) (-0.153,-0.100) (0.283,0.315) (-0.248,-0.193) (-0.468,-0.423) (-0.316,-0.272) (-0.256,-0.140) (-0.044,-0.032)
 (0.153,0.104) (-0.056,-0.102) (-0.184,-0.162) (-0.323,-0.266) (-0.172,-0.167) (-0.235,-0.195) (0.437,0.357) (-0.531,-0.503) (0.776,0.678) (-0.205,-0.267) (0.093,0.084) (0.138,0.103)
 (0.329,0.263) (0.058,0.057) (-0.163,-0.095) (0.381,0.283) (-0.224,-0.176) (0.030,0.043) (0.098,0.087) (0.027,0.033) (-0.038,0.002) (-0.054,-0.019) (0.039,-0.025) (0.042,0.052)
 (0.023,-0.045) (0.022,0.030) (-0.222,-0.173) (-0.211,-0.167) (-0.117,-0.081) (-0.589,-0.480) (-0.205,-0.248) (0.539,0.495) (0.418,0.282) (-0.205,-0.088) (0.668,0.575) (-0.195,-0.162)
 (0.342,0.327) (-0.179,-0.178) (-0.275,-0.198) (-0.466,-0.346) (-0.204,-0.176) (-0.148,-0.178) (-0.111,-0.063) (0.122,0.091) (0.015,0.065) (-0.096,-0.015) (-0.052,-0.031) (-0.110,-0.061)
 (-0.005,-0.002) (0.228,0.137) (0.403,0.322) (-0.070,-0.049) (-0.110,-0.011) (0.304,0.265) (-0.109,-0.043) (0.188,0.196) (-0.171,-0.089) (-0.119,-0.042) (-0.075,-0.051) (0.357,0.284)
 (-0.028,-0.003) (-0.057,-0.032) (0.036,0.097) (-0.044,-0.025) (-0.048,-0.033) (0.020,0.026) (0.027,0.002) (0.297,0.255) (-0.059,-0.074) (-0.167,-0.115) (0.298,0.299) (0.206,0.143)
 (-0.061,-0.032) (-0.006,0.002) (-0.050,-0.026) (-0.145,-0.112) (-0.062,-0.050) (0.003,0.035) (-0.057,-0.069) (-0.033,-0.026) (-0.045,-0.091) (0.192,0.226) (-0.075,-0.041) (-0.140,-0.223)
 (-0.079,-0.092) (0.054,0.061) (-0.087,-0.109) (-0.059,-0.013) (-0.060,-0.024) (-0.054,-0.031) (-0.061,-0.017) (-0.070,-0.135) (0.300,0.170) (-0.121,-0.084) (0.555,0.418) (-0.108,-0.111)
 (-0.034,-0.022) (0.022,-0.004) (-0.046,-0.027) (-0.024,-0.010) (0.133,0.138) (-0.006,-0.003) (-0.043,-0.027) (-0.139,-0.076) (-0.357,-0.315) (-0.253,-0.077) (-0.172,-0.073) (-0.142,-0.110)
 (-0.008,-0.015) (-0.130,-0.079) (0.362,0.309) (0.691,0.555) (-0.073,-0.068) (-0.200,-0.064) (0.491,0.450) (-0.192,-0.096) (0.363,0.345) (-0.277,-0.271) (-0.431,-0.382) (-0.308,-0.254)
 (-0.248,-0.125) (-0.057,-0.040) (0.153,0.093) (-0.031,-0.091) (-0.174,-0.155) (-0.305,-0.257) (-0.165,-0.164) (-0.216,-0.174) (0.330,0.271) (-0.475,-0.449) (0.744,0.655) (-0.215,-0.305)
 (0.023,0.049) (-0.142,-0.122) (-0.119,-0.102) (-0.132,-0.109) (-0.344,-0.309) (-0.240,-0.223) (0.065,0.092) (-0.170,-0.140) (-0.094,-0.085) (-0.111,-0.159) (0.311,0.255) (-0.145,-0.045)
 (-0.291,-0.370) (-0.166,-0.153) (0.159,0.154) (-0.191,-0.177) (0.184,0.167) (0.183,0.153) (0.089,0.074) (0.429,0.389) (0.162,0.237) (-0.532,-0.455) (-0.316,-0.273) (-0.098,-0.113)
 (0.356,0.324) (-0.102,-0.105) (0.205,0.168) (-0.113,-0.098) (-0.152,-0.107) (-0.244,-0.207) (-0.115,-0.103) (-0.074,-0.140) (-0.134,-0.099) (0.086,0.071) (-0.031,0.038) (-0.125,-0.046)
 (-0.074,-0.059) (-0.095,-0.090) (0.076,0.082) (-0.084,-0.059) (-0.114,-0.066) (0.009,0.030) (-0.041,-0.052) (-0.068,-0.084) (-0.001,-0.021) (-0.111,-0.063) (0.049,-0.010) (-0.111,-0.068)
 (-0.044,-0.064) (0.500,0.356) (-0.031,-0.022) (-0.125,-0.043) (-0.006,0.107) (-0.066,-0.035) (-0.073,-0.037) (0.012,0.033) (0.018,0.019) (0.285,0.199) (-0.124,-0.065) (-0.142,-0.132)
 (0.172,0.093) (0.185,0.101) (0.041,0.046) (0.064,0.060) (0.027,0.026) (0.023,-0.020) (0.160,0.127) (0.015,-0.079) (-0.002,-0.009) (0.048,0.025) (-0.009,-0.004) (0.035,0.047)
 (-0.031,-0.010) (-0.042,-0.076) (-0.001,-0.044) (-0.016,0.048) (-0.017,-0.023) (0.039,0.039) (0.035,0.023) (0.024,0.027) (0.113,0.076) (0.026,0.006) (-0.126,-0.099) (-0.039,0.013)
 (-0.019,-0.011) (0.746,0.615) (-0.235,-0.199) (0.412,0.378) (-0.251,-0.211) (-0.322,-0.243) (-0.442,-0.343) (-0.242,-0.198) (-0.188,-0.267) (0.206,0.207) (0.237,0.219) (0.064,0.126)
 (0.222,0.259) (0.215,0.205) (0.255,0.244) (-0.304,-0.287) (0.393,0.337) (0.743,0.611) (-0.097,-0.079) (-0.223,-0.098) (0.542,0.490) (-0.199,-0.109) (0.372,0.353) (-0.298,-0.268)
 (-0.189,-0.124) (-0.097,-0.083) (0.455,0.359) (-0.034,-0.025) (-0.083,-0.044) (-0.013,0.127) (-0.081,-0.061) (-0.091,-0.070) (-0.037,-0.034) (0.159,0.138) (-0.208,-0.126) (-0.141,-0.101)
 (-0.139,-0.144) (-0.080,-0.005) (-0.092,-0.024) (-0.022,-0.058) (-0.084,-0.063) (-0.067,-0.065) (-0.223,-0.138) (-0.160,-0.115) (0.062,0.062) (-0.081,-0.056) (-0.064,-0.053) (-0.052,0.101)
 (0.295,0.378) (-0.041,-0.108) (-0.108,-0.120) (-0.072,-0.068) (0.105,0.082) (-0.073,-0.061) (-0.022,-0.011) (-0.030,-0.031) (-0.006,-0.010) (-0.042,-0.053) (-0.026,-0.132) (0.008,0.032)
 (0.075,0.078) (-0.084,-0.038) (-0.104,-0.094) (-0.024,-0.018) (0.035,0.000) (-0.045,-0.023) (0.001,-0.008) (0.114,0.110) (-0.017,-0.010) (-0.039,-0.048) (-0.090,-0.059) (0.103,0.051)
 (0.003,0.033) (-0.065,-0.030) (-0.053,-0.039) (-0.070,-0.064) (0.021,0.013) (-0.084,-0.091) (-0.099,-0.074) (0.001,0.037) (-0.060,-0.037) (-0.011,-0.041) (-0.004,0.007) (-0.132,-0.097)
 (0.052,0.047) (-0.175,-0.123) (-0.093,-0.056) (-0.082,-0.096) (-0.028,-0.018) (0.042,0.051) (0.021,-0.038) (-0.063,-0.044) (-0.104,-0.081) (-0.074,-0.077) (-0.090,-0.080) (0.127,0.115)
 (-0.188,-0.158) (0.304,0.293) (-0.117,-0.112) (-0.008,0.019) (-0.208,-0.164) (-0.172,-0.141) (-0.175,-0.136) (-0.462,-0.403) (-0.327,-0.277) (0.171,0.123) (-0.234,-0.137) (-0.134,-0.111)
 (-0.104,-0.100) (0.332,0.341) (-0.132,-0.040) (-0.271,-0.296) (-0.163,-0.132) (0.155,0.118) (-0.181,-0.166) (-0.045,-0.030) (-0.048,-0.032) (-0.032,-0.037) (-0.041,-0.027) (-0.054,-0.080)
 (0.244,0.146) (-0.141,-0.045) (0.542,0.370) (0.512,0.379) (-0.147,-0.093) (0.200,0.240) (-0.129,-0.108) (-0.185,-0.122) (-0.257,-0.203) (-0.147,-0.098) (-0.118,-0.267) (-0.092,-0.065)
 (0.085,0.052) (-0.000,0.048) (-0.081,-0.018) (-0.056,-0.045) (-0.072,-0.063) (0.041,0.039) (-0.062,-0.008) (-0.073,-0.058) (-0.040,-0.009) (-0.007,-0.015) (0.044,0.009) (0.006,0.017)
 (-0.012,-0.014) (0.010,-0.052) (0.084,0.092) (0.059,0.061) (0.028,0.025) (0.024,0.033) (-0.000,-0.030) (-0.001,-0.018) (0.053,0.062) (0.071,0.065) (0.008,0.028) (0.014,0.000)
 (0.314,0.234) (-0.114,-0.085) (-0.111,-0.081) (0.208,0.190) (0.192,0.114) (-0.014,0.001) (0.091,0.088) (0.007,-0.003) (-0.066,-0.042) (0.121,0.098) (-0.166,-0.130) (0.143,0.131)
 (0.017,0.016) (-0.116,-0.146) (0.336,0.293) (-0.138,-0.017) (-0.297,-0.317) (-0.176,-0.143) (0.170,0.135) (-0.200,-0.175) (-0.160,-0.113) (-0.152,-0.122) (-0.079,-0.055) (-0.429,-0.344)
 (-0.153,-0.174) (0.394,0.335) (0.297,0.200) (-0.137,-0.108) (0.615,0.536) (-0.171,-0.132) (0.298,0.303) (-0.150,-0.153) (-0.147,-0.080) (-0.358,-0.283) (-0.193,-0.166) (-0.137,-0.138)
 (0.069,0.087) (0.113,0.083) (0.047,0.067) (0.096,0.128) (0.082,0.085) (0.061,0.061) (-0.073,-0.077) (-0.055,-0.027) (-0.148,-0.158) (-0.058,-0.065) (0.485,0.368) (0.104,0.055)
 (0.246,0.208) (0.151,0.099) (-0.049,-0.057) (-0.242,-0.208) (-0.157,-0.149) (0.322,0.277) (-0.033,-0.062) (-0.057,-0.008) (-0.018,0.157) (-0.099,-0.100) (-0.157,-0.126) (-0.093,-0.076)
 (0.103,0.129) (-0.116,-0.094) (-0.286,-0.206) (0.041,-0.022) (-0.153,-0.144) (-0.069,-0.030) (0.055,0.011) (0.094,0.076) (0.003,0.009) (-0.099,-0.057) (0.110,0.074) (-0.114,-0.092)
 (0.003,0.031) (0.039,0.016) (0.030,-0.060) (-0.062,-0.063) (-0.070,-0.060) (0.155,0.133) (0.044,0.028) (0.047,-0.009) (0.026,0.019) (-0.044,-0.015) (-0.047,-0.022) (-0.038,-0.031)
 (-0.028,-0.001) (-0.053,-0.081) (0.271,0.138) (-0.173,-0.088) (0.562,0.408) (-0.125,-0.104) (0.049,0.032) (-0.098,-0.066) (0.041,0.048) (0.370,0.205) (0.179,0.107) (0.015,0.015)
 (0.039,-0.080) (-0.070,-0.056) (0.026,0.006) (-0.018,0.035) (-0.068,-0.032) (-0.059,-0.041) (-0.053,-0.070) (0.069,0.048) (0.259,0.240) (0.480,0.413) (-0.062,-0.074) (-0.131,-0.077)
 (0.310,0.313) (-0.114,-0.088) (0.214,0.246) (-0.191,-0.122) (-0.124,-0.102) (-0.088,-0.055) (-0.071,-0.024) (-0.017,0.001) (0.062,-0.019) (-0.006,-0.050) (-0.050,-0.021) (-0.088,-0.072)
 (-0.041,-0.020) (-0.062,-0.091) (0.081,0.088) (-0.110,-0.071) (0.167,0.175) (-0.066,0.005) (0.053,0.015) (0.097,0.063) (0.066,0.074) (0.050,0.040) (0.023,0.019) (0.238,0.174)
 (0.028,-0.028) (0.002,0.025) (0.092,0.050) (-0.015,-0.014) (0.108,0.239) (-0.042,-0.067) (-0.052,-0.002) (-0.026,-0.046) (0.039,0.060) (-0.028,-0.043) (0.060,0.054) (0.047,0.046)
 (0.035,0.037) (0.186,0.159) (0.044,0.057) (-0.180,-0.146) (-0.079,-0.058) (0.026,0.051) (0.647,0.555) (-0.187,-0.166) (0.357,0.323) (-0.197,-0.180) (-0.266,-0.217) (-0.419,-0.333)
 (-0.203,-0.174) (-0.142,-0.196) (-0.058,-0.019) (0.019,-0.007) (-0.017,0.040) (-0.071,-0.015) (-0.032,-0.034) (-0.060,-0.041) (0.040,0.032) (-0.055,-0.058) (-0.103,-0.090) (-0.037,-0.005)
 (-0.030,-0.009) (-0.118,-0.127) (0.001,0.006) (-0.091,-0.049) (0.068,0.060) (-0.078,-0.041) (0.159,0.120) (-0.020,-0.018) (0.007,-0.011) (-0.030,0.002) (-0.044,-0.038) (0.012,0.009)
 (0.061,0.051) (-0.125,-0.100) (0.037,0.055) (0.015,0.011) (-0.384,-0.293) (0.234,0.162) (-0.171,-0.162) (-0.021,0.002) (0.027,0.039) (-0.028,-0.025) (0.061,0.056) (0.322,0.241)
 (-0.089,-0.044) (0.029,0.055) (0.027,0.019) (0.031,0.018) (0.039,-0.112) (-0.073,-0.063) (-0.079,-0.084) (0.066,0.115) (0.067,0.014) (0.033,0.020) (0.037,0.024) (-0.080,-0.049)
 (-0.072,-0.057) (-0.044,-0.034) (-0.203,-0.151) (-0.081,0.018) (0.195,0.172) (0.112,0.107) (0.006,0.006) (0.649,0.547) (-0.184,-0.172) (0.333,0.312) (-0.197,-0.172) (-0.260,-0.211)
 (-0.407,-0.335) (-0.193,-0.168) (-0.134,-0.212) (0.188,0.182) (0.229,0.225) (0.123,0.127) (0.205,0.206) (0.185,0.183) (0.186,0.199) (-0.183,-0.195) (0.363,0.344) (0.668,0.581)
 (-0.082,-0.094) (-0.185,-0.117) (0.472,0.461) (-0.165,-0.122) (0.305,0.320) (-0.275,-0.245) (-0.072,-0.063) (0.144,0.132) (-0.022,-0.012) (0.011,0.014) (0.008,-0.004) (-0.038,-0.139)
 (0.018,0.025) (0.083,0.068) (-0.177,-0.144) (-0.208,-0.135) (0.335,0.234) (-0.476,-0.442) (0.713,0.595) (-0.246,-0.250) (0.141,0.055) (0.108,0.030) (0.049,0.023) (0.037,0.024)
 (0.071,0.079) (-0.007,0.048) (0.067,0.047) (-0.006,-0.032) (0.016,0.014) (-0.097,-0.166) (0.293,0.272) (-0.136,-0.033) (-0.256,-0.319) (-0.146,-0.141) (0.116,0.122) (-0.175,-0.168)
 (-0.190,-0.156) (-0.173,-0.131) (-0.107,-0.084) (-0.429,-0.339) (-0.150,-0.249) (0.575,0.456) (0.173,0.137) (0.311,0.267) (-0.037,0.015) (0.017,0.040) (-0.028,-0.020) (0.029,0.056)
 (-0.018,-0.037) (-0.018,-0.091) (0.020,0.026) (0.013,0.006) (-0.048,-0.067) (-0.053,-0.070) (-0.023,-0.133) (-0.051,-0.123) (-0.048,-0.072) (-0.047,-0.076) (0.056,0.057) (0.038,0.073)
 (0.062,0.018) (-0.008,-0.021) (-0.030,-0.022) (0.039,0.008) (-0.031,0.006) (0.042,0.098) (-0.028,0.257) (-0.088,-0.060) (-0.050,-0.065) (0.377,0.289) (-0.015,-0.031) (-0.083,-0.036)
 (0.008,0.114) (-0.040,-0.045) (-0.048,-0.042) (-0.076,-0.030) (-0.083,-0.016) (0.137,0.033) (-0.156,-0.119) (0.247,0.170) (-0.138,-0.052) (0.167,0.039) (-0.033,-0.050) (-0.044,-0.025)
 (-0.034,-0.043) (-0.091,-0.114) (-0.079,-0.066) (0.057,0.033) (-0.051,-0.055) (-0.021,-0.031) (-0.062,0.020) (0.117,0.075) (0.154,0.167) (-0.164,-0.245) (-0.098,-0.071) (0.111,0.073)
 (-0.097,-0.071) (-0.131,-0.115) (-0.111,-0.106) (-0.072,-0.066) (-0.354,-0.295) (-0.084,-0.220) (0.285,0.292) (0.235,0.183) (-0.188,-0.090) (-0.362,-0.303) (0.153,0.112) (-0.250,-0.228)
 (0.124,0.127) (0.455,0.288) (0.389,0.224) (0.126,0.100) (0.127,0.026) (0.317,0.259) (0.407,0.302) (0.175,0.108) (0.347,0.280) (0.330,0.262) (0.299,0.283) (-0.332,-0.289)
 (0.206,0.128) (0.419,0.325) (-0.023,-0.033) (-0.165,-0.064) (0.228,0.261) (-0.139,-0.069) (0.154,0.208) (-0.149,-0.040) (-0.023,0.017) (0.104,0.094) (0.026,0.010) (0.002,-0.030)
 (-0.031,-0.006) (0.008,0.033) (0.017,0.005) (0.027,0.048) (0.031,0.077) (0.094,0.177) (-0.268,-0.251) (0.197,0.140) (-0.196,-0.260) (0.048,0.161) (-0.044,-0.069)};
\draw[haocol!60,line width=.3pt] (axis cs:0.890,0.766) -- (axis cs:0.350,-0.080);
\addplot[only marks,mark=*,mark size=1.6pt,color=haocol] coordinates { (0.890,0.766) };
\node[font=\sffamily\tiny,anchor=west,color=haocol,inner sep=1pt] at (axis cs:0.360,-0.080) {\textit{LYZ}--CD371};
\draw[haocol!60,line width=.3pt] (axis cs:0.869,0.609) -- (axis cs:0.350,-0.165);
\addplot[only marks,mark=*,mark size=1.6pt,color=haocol] coordinates { (0.869,0.609) };
\node[font=\sffamily\tiny,anchor=west,color=haocol,inner sep=1pt] at (axis cs:0.360,-0.165) {\textit{CD79A}--CD19};
\draw[coloncol!60,line width=.3pt] (axis cs:0.819,0.570) -- (axis cs:0.350,-0.250);
\addplot[only marks,mark=*,mark size=1.6pt,color=coloncol] coordinates { (0.819,0.570) };
\node[font=\sffamily\tiny,anchor=west,color=coloncol,inner sep=1pt] at (axis cs:0.360,-0.250) {\textit{IL7R}--CD127};
\draw[coloncol!60,line width=.3pt] (axis cs:0.765,0.420) -- (axis cs:0.350,-0.335);
\addplot[only marks,mark=*,mark size=1.6pt,color=coloncol] coordinates { (0.765,0.420) };
\node[font=\sffamily\tiny,anchor=west,color=coloncol,inner sep=1pt] at (axis cs:0.360,-0.335) {\textit{IGHD}--IgD};
\draw[coloncol!60,line width=.3pt] (axis cs:0.755,0.294) -- (axis cs:0.350,-0.420);
\addplot[only marks,mark=*,mark size=1.6pt,color=coloncol] coordinates { (0.755,0.294) };
\node[font=\sffamily\tiny,anchor=west,color=coloncol,inner sep=1pt] at (axis cs:0.360,-0.420) {\textit{FCGR3A}--CD16};
\draw[coloncol!60,line width=.3pt] (axis cs:0.688,0.259) -- (axis cs:0.350,-0.505);
\addplot[only marks,mark=*,mark size=1.6pt,color=coloncol] coordinates { (0.688,0.259) };
\node[font=\sffamily\tiny,anchor=west,color=coloncol,inner sep=1pt] at (axis cs:0.360,-0.505) {\textit{KLRB1}--CD161};
\end{axis}
\begin{axis}[at={(a.east)},anchor=west,xshift=1.3cm,name=b,width=.34\textwidth,height=.34\textwidth,xmin=-0.6,xmax=0.8,ymin=-0.6,ymax=0.8,
 xlabel={Held-out correlation},ylabel={Predicted correlation},title={\textbf{b}\enspace Within types, no pairs},
 axis x line*=bottom,axis y line*=left,tick align=outside,xmajorgrids=true,grid style={gray!18},axis line style={gray!55},tick label style={font=\sffamily\scriptsize},label style={font=\sffamily\small},clip=true,title style={font=\sffamily\small,at={(0,1)},anchor=base west,yshift=5pt}]
\draw[black!35] (axis cs:-0.6,-0.6) -- (axis cs:0.8,0.8);
\addplot[only marks,mark=*,mark size=.35pt,color=donorcol!45] coordinates {
 (-0.019,0.013) (-0.014,-0.027) (-0.022,-0.012) (-0.054,-0.042) (-0.025,0.021) (0.007,0.016) (0.005,0.004) (-0.008,-0.011) (-0.032,-0.037) (0.001,0.008) (-0.031,-0.020) (-0.001,0.005)
 (-0.025,-0.003) (0.018,-0.042) (0.001,0.004) (-0.006,-0.005) (-0.122,0.004) (0.080,0.030) (-0.024,0.106) (-0.040,0.041) (-0.030,0.009) (-0.155,-0.024) (-0.015,-0.025) (0.014,-0.013)
 (-0.008,-0.005) (0.007,-0.006) (-0.034,-0.072) (-0.095,-0.070) (-0.010,-0.013) (0.071,-0.016) (-0.008,0.017) (-0.224,-0.150) (-0.019,0.020) (-0.008,0.003) (-0.001,0.034) (-0.001,0.002)
 (-0.098,-0.291) (0.002,0.041) (-0.001,0.023) (-0.014,0.031) (-0.010,0.068) (0.025,-0.096) (-0.005,-0.043) (0.026,0.107) (-0.035,-0.159) (0.065,0.002) (-0.029,-0.032) (-0.005,-0.038)
 (0.008,-0.021) (-0.047,-0.015) (-0.023,-0.002) (0.014,-0.002) (0.002,-0.012) (-0.005,-0.023) (0.016,0.152) (-0.019,-0.066) (-0.186,-0.022) (0.002,0.024) (0.014,0.054) (0.052,-0.034)
 (0.006,0.008) (-0.025,0.018) (-0.017,0.005) (-0.021,-0.000) (-0.065,-0.015) (-0.057,-0.108) (0.083,0.037) (0.038,-0.042) (-0.009,0.007) (0.029,0.009) (-0.001,-0.020) (-0.023,0.005)
 (0.015,-0.011) (0.009,0.004) (0.039,0.036) (0.017,-0.004) (-0.006,0.028) (-0.002,0.006) (-0.000,-0.003) (-0.004,0.041) (0.002,0.027) (-0.004,-0.004) (-0.005,-0.017) (-0.001,0.006)
 (0.004,-0.002) (-0.016,-0.014) (0.005,-0.006) (0.078,0.076) (0.044,0.043) (0.021,0.039) (-0.065,-0.017) (-0.003,0.029) (0.006,-0.001) (-0.021,-0.020) (0.080,0.079) (0.001,-0.005)
 (-0.002,0.001) (0.023,0.064) (-0.005,-0.009) (-0.004,0.002) (0.004,-0.049) (-0.008,-0.030) (-0.008,0.089) (-0.006,-0.011) (0.048,0.077) (0.020,-0.013) (-0.036,-0.002) (-0.024,0.014)
 (0.078,0.068) (-0.017,0.020) (-0.459,-0.268) (-0.036,-0.011) (-0.010,-0.061) (-0.009,-0.000) (0.014,0.029) (0.005,0.178) (-0.005,-0.008) (-0.009,0.011) (-0.002,-0.113) (0.001,0.028)
 (0.001,-0.019) (-0.003,0.017) (-0.002,0.043) (-0.009,0.003) (-0.003,-0.001) (-0.023,0.024) (-0.002,-0.017) (0.007,-0.010) (0.003,0.007) (0.002,0.028) (-0.062,-0.081) (-0.002,0.001)
 (-0.060,-0.046) (0.012,-0.006) (0.016,0.032) (0.059,0.117) (0.009,0.000) (0.005,-0.006) (0.008,0.016) (-0.026,-0.059) (0.001,-0.004) (0.003,0.003) (-0.025,-0.014) (0.001,0.022)
 (-0.005,-0.015) (0.006,-0.004) (-0.030,0.005) (0.003,0.019) (-0.038,-0.035) (0.086,-0.052) (-0.004,0.015) (-0.008,-0.006) (-0.002,-0.017) (-0.049,-0.084) (0.011,-0.019) (0.000,0.005)
 (0.002,-0.003) (-0.002,0.004) (-0.014,0.061) (0.017,-0.011) (-0.009,-0.035) (0.009,-0.031) (-0.005,-0.037) (-0.001,0.059) (-0.001,-0.017) (0.058,0.031) (0.043,-0.002) (-0.093,-0.044)
 (0.062,0.023) (-0.069,-0.026) (-0.014,-0.011) (-0.030,-0.043) (-0.006,0.003) (0.171,0.076) (-0.016,0.037) (-0.023,-0.000) (0.001,-0.001) (-0.003,0.027) (0.004,-0.020) (-0.041,-0.164)
 (0.004,0.044) (-0.066,-0.123) (0.008,0.021) (-0.000,-0.003) (-0.008,-0.048) (0.017,-0.029) (-0.060,-0.108) (-0.022,0.001) (0.017,0.046) (-0.006,0.133) (-0.003,-0.038) (-0.037,-0.025)
 (-0.013,-0.002) (-0.016,0.039) (0.038,0.016) (-0.059,-0.022) (-0.234,-0.088) (-0.022,-0.003) (-0.020,0.069) (-0.100,0.017) (0.071,0.044) (-0.048,0.142) (-0.053,0.076) (-0.013,0.010)
 (-0.126,-0.015) (0.001,-0.034) (0.017,0.105) (0.239,0.107) (0.006,-0.044) (0.004,-0.004) (-0.075,0.049) (-0.033,-0.034) (-0.091,0.002) (-0.007,0.160) (-0.018,-0.042) (-0.014,0.006)
 (-0.030,-0.016) (-0.000,-0.001) (-0.023,-0.024) (0.034,-0.119) (-0.002,-0.006) (-0.052,-0.010) (-0.013,-0.018) (0.026,0.049) (-0.019,-0.007) (-0.034,-0.074) (-0.056,-0.114) (0.004,-0.058)
 (-0.008,0.027) (-0.018,-0.052) (-0.009,0.023) (-0.001,-0.022) (-0.030,-0.059) (-0.004,-0.007) (-0.063,0.020) (-0.002,-0.029) (-0.001,-0.019) (0.003,-0.074) (-0.006,0.031) (0.099,0.084)
 (-0.009,-0.150) (0.002,-0.025) (-0.083,0.006) (-0.017,0.006) (0.001,0.006) (0.002,-0.030) (-0.000,-0.007) (-0.037,-0.077) (-0.006,-0.083) (0.010,0.014) (0.018,0.119) (0.003,-0.028)
 (0.100,0.080) (-0.017,-0.007) (0.050,0.057) (0.019,0.009) (-0.014,-0.016) (-0.062,-0.054) (-0.013,-0.017) (-0.012,-0.165) (-0.003,0.008) (0.235,0.060) (0.092,0.078) (0.105,0.070)
 (0.017,0.001) (0.000,-0.005) (-0.110,-0.068) (0.005,0.020) (-0.006,0.018) (-0.020,0.005) (-0.019,-0.033) (0.049,-0.008) (-0.007,0.006) (-0.021,-0.007) (-0.005,-0.013) (-0.314,-0.232)
 (-0.004,-0.006) (0.001,-0.005) (-0.005,0.022) (0.002,0.038) (-0.092,-0.238) (0.002,0.023) (-0.000,0.026) (-0.002,-0.022) (-0.005,-0.016) (-0.005,0.059) (-0.002,-0.015) (0.025,0.035)
 (-0.001,0.012) (-0.031,-0.006) (-0.018,-0.012) (-0.023,0.004) (-0.031,-0.001) (0.004,-0.047) (-0.007,-0.041) (-0.025,-0.045) (-0.035,-0.049) (-0.014,-0.004) (-0.007,-0.038) (0.038,0.065)
 (0.008,-0.016) (-0.077,-0.031) (-0.005,0.018) (-0.010,-0.005) (-0.002,0.006) (0.002,-0.009) (0.006,-0.028) (0.005,-0.015) (-0.016,-0.063) (-0.013,-0.008) (-0.020,0.010) (-0.011,0.084)
 (0.001,-0.008) (0.124,0.100) (-0.022,0.010) (-0.036,0.047) (0.009,0.015) (-0.010,-0.005) (0.029,-0.055) (0.005,0.003) (-0.013,-0.251) (-0.006,-0.019) (0.352,0.094) (0.079,0.105)
 (0.150,0.078) (0.002,-0.023) (0.044,-0.003) (-0.096,-0.045) (0.022,0.026) (0.017,0.060) (0.007,0.019) (-0.013,-0.012) (0.077,0.016) (-0.010,0.001) (-0.038,0.010) (-0.011,-0.038)
 (0.035,-0.000) (-0.002,-0.015) (-0.005,-0.020) (0.001,0.010) (-0.001,-0.001) (-0.001,-0.007) (0.010,0.002) (-0.004,0.013) (-0.009,0.009) (-0.001,-0.006) (0.109,0.063) (0.014,-0.006)
 (0.006,0.023) (0.149,0.303) (0.051,0.058) (-0.008,-0.008) (-0.008,-0.011) (-0.001,-0.006) (-0.076,-0.055) (-0.007,-0.004) (-0.009,-0.026) (0.001,0.020) (-0.001,0.004) (0.022,-0.028)
 (0.048,0.086) (0.013,-0.036) (-0.003,-0.040) (0.006,-0.005) (-0.126,-0.105) (-0.011,-0.016) (-0.002,-0.009) (-0.000,0.027) (-0.003,0.009) (-0.005,0.039) (0.000,0.072) (-0.003,-0.024)
 (-0.008,-0.019) (-0.004,-0.015) (0.054,0.075) (-0.007,0.007) (-0.026,0.025) (-0.001,0.001) (-0.005,0.002) (0.028,0.006) (-0.004,0.004) (-0.013,-0.086) (0.048,0.019) (0.001,0.055)
 (0.023,-0.100) (0.040,-0.026) (0.046,0.059) (0.078,0.033) (-0.031,-0.005) (0.001,-0.014) (-0.030,-0.006) (0.016,0.013) (-0.009,-0.000) (0.056,0.010) (-0.009,0.007) (-0.001,0.008)
 (0.005,-0.095) (-0.011,0.049) (0.000,0.037) (0.002,0.002) (0.001,0.001) (0.001,-0.015) (-0.033,-0.033) (-0.000,-0.013) (0.006,0.003) (-0.002,0.029) (-0.010,0.016) (0.071,0.001)
 (-0.044,-0.045) (0.023,0.100) (-0.047,-0.017) (0.143,0.026) (-0.025,0.021) (0.045,0.050) (-0.019,0.018) (-0.421,-0.254) (-0.025,-0.016) (0.032,-0.037) (-0.010,0.011) (0.004,0.025)
 (-0.004,0.044) (0.003,0.009) (-0.005,0.007) (0.029,-0.039) (-0.003,0.002) (0.058,-0.013) (-0.003,-0.014) (-0.011,-0.028) (-0.015,-0.009) (-0.019,-0.009) (-0.057,-0.018) (-0.018,0.035)
 (-0.020,-0.021) (0.009,0.057) (-0.014,0.001) (-0.027,-0.015) (0.004,0.000) (0.007,0.004) (-0.007,0.002) (-0.023,-0.009) (-0.070,-0.021) (-0.018,-0.008) (0.002,-0.000) (0.111,0.007)
 (0.088,0.019) (0.100,-0.163) (0.102,-0.071) (0.045,0.042) (0.126,0.016) (-0.083,-0.028) (-0.002,-0.018) (-0.011,-0.007) (-0.000,-0.012) (0.075,0.055) (-0.010,0.025) (0.022,0.024)
 (-0.033,-0.018) (-0.007,0.017) (-0.060,-0.005) (-0.072,0.010) (-0.005,-0.004) (-0.004,0.031) (-0.004,-0.005) (-0.045,-0.163) (-0.012,0.023) (0.000,0.021) (-0.025,-0.004) (-0.008,0.023)
 (0.028,0.021) (-0.015,-0.062) (0.022,0.069) (-0.084,0.020) (0.095,0.030) (0.071,0.025) (-0.078,0.013) (-0.002,0.002) (-0.034,-0.008) (0.090,-0.001) (0.091,0.018) (-0.003,0.032)
 (0.011,-0.001) (-0.009,-0.040) (0.044,0.031) (0.004,-0.017) (-0.014,-0.021) (-0.012,-0.008) (0.099,0.055) (-0.007,-0.006) (-0.000,0.004) (0.001,-0.008) (0.005,-0.021) (0.010,-0.016)
 (0.003,-0.018) (-0.042,-0.000) (-0.015,0.091) (0.514,0.101) (0.006,0.022) (-0.000,-0.001) (-0.001,-0.008) (0.021,0.009) (-0.017,-0.021) (-0.034,-0.037) (-0.010,-0.011) (-0.009,-0.042)
 (-0.059,0.005) (0.042,0.026) (-0.043,0.074) (-0.037,0.049) (0.027,0.015) (-0.067,0.001) (-0.005,-0.032) (-0.002,-0.008) (-0.001,0.015) (0.047,0.053) (-0.027,-0.020) (-0.002,-0.081)
 (-0.001,-0.011) (-0.032,-0.059) (-0.005,0.016) (-0.000,-0.008) (-0.008,-0.032) (0.003,-0.025) (0.016,-0.003) (-0.005,0.005) (-0.072,0.010) (0.001,-0.009) (-0.022,-0.004) (-0.006,-0.007)
 (-0.008,-0.002) (-0.049,0.020) (-0.007,-0.014) (0.123,0.214) (-0.039,-0.162) (-0.070,-0.032) (-0.003,-0.024) (-0.003,0.028) (-0.012,-0.012) (-0.025,-0.041) (-0.007,-0.008) (-0.016,0.052)
 (-0.003,-0.001) (-0.005,-0.005) (-0.004,-0.155) (-0.003,0.010) (0.007,-0.017) (0.012,0.024) (-0.001,-0.024) (-0.007,0.130) (0.003,-0.004) (0.001,0.002) (-0.005,-0.043) (0.011,-0.026)
 (-0.054,-0.100) (-0.020,0.044) (-0.001,0.037) (-0.007,0.091) (-0.005,-0.028) (-0.042,-0.071) (-0.011,0.009) (-0.079,-0.021) (0.014,0.024) (0.012,0.024) (0.104,0.049) (-0.000,0.020)
 (-0.014,-0.207) (-0.006,0.004) (-0.024,-0.031) (-0.018,0.011) (-0.040,-0.003) (-0.009,-0.013) (-0.016,-0.006) (-0.010,-0.005) (0.020,-0.033) (0.026,0.070) (-0.021,0.002) (-0.002,0.013)
 (0.050,-0.008) (-0.013,0.013) (-0.000,0.014) (-0.015,-0.009) (-0.005,-0.046) (-0.009,-0.031) (-0.007,-0.028) (0.016,-0.009) (-0.012,0.015) (-0.031,0.069) (0.005,-0.015) (-0.042,-0.025)
 (-0.004,0.019) (0.006,0.011) (0.119,0.068) (-0.006,-0.057) (-0.000,-0.031) (0.047,0.034) (0.027,-0.003) (-0.008,0.021) (0.007,0.026) (0.006,0.009) (-0.005,-0.025) (0.112,-0.005)
 (0.031,0.003) (0.007,0.030) (0.063,0.013) (0.001,0.146) (0.088,0.008) (-0.003,-0.003) (-0.009,-0.070) (-0.000,-0.001) (-0.068,-0.019) (-0.008,-0.037) (-0.003,0.016) (-0.000,-0.016)
 (-0.003,-0.028) (-0.006,-0.017) (-0.004,-0.056) (0.006,-0.032) (0.011,0.085) (-0.003,-0.013) (0.002,-0.022) (-0.001,0.024) (0.015,-0.004) (-0.008,0.001) (0.062,0.012) (-0.045,0.003)
 (-0.007,0.007) (0.004,0.050) (0.022,0.043) (0.093,0.071) (0.006,-0.013) (0.036,0.028) (0.062,0.066) (0.034,0.036) (-0.020,-0.025) (0.067,0.004) (-0.017,0.017) (0.010,-0.016)
 (0.021,0.001) (-0.079,-0.016) (0.030,0.026) (-0.000,-0.041) (-0.015,0.019) (-0.058,-0.061) (-0.027,-0.014) (0.005,0.012) (-0.017,0.029) (-0.011,-0.003) (0.009,-0.048) (-0.005,0.022)
 (-0.030,-0.039) (-0.017,-0.015) (-0.002,0.011) (0.027,0.073) (-0.020,-0.074) (0.074,0.160) (-0.045,-0.007) (0.019,0.040) (-0.016,-0.003) (0.003,0.035) (0.006,0.002) (0.050,0.059)
 (0.008,-0.009) (0.002,-0.055) (-0.032,-0.016) (0.002,-0.004) (-0.022,0.052) (0.060,0.001) (0.042,0.074) (0.029,-0.093) (-0.024,0.001) (0.202,0.053) (-0.011,-0.002) (-0.005,0.033)
 (-0.003,0.015) (-0.001,0.009) (-0.001,0.011) (0.004,-0.032) (0.008,0.006) (-0.001,-0.028) (0.028,0.011) (0.070,0.010) (-0.005,-0.004) (0.006,0.022) (-0.006,-0.004) (0.000,-0.001)
 (-0.009,0.003) (-0.002,-0.001) (-0.004,-0.105) (-0.010,0.005) (-0.047,-0.041) (-0.031,-0.004) (-0.050,-0.016) (-0.008,-0.022) (-0.006,0.010) (-0.005,-0.002) (0.017,-0.046) (-0.124,-0.025)
 (-0.027,0.028) (-0.003,0.000) (0.075,0.002) (0.008,0.028) (0.074,0.001) (-0.013,-0.084) (-0.009,-0.008) (-0.015,-0.023) (-0.011,-0.025) (0.017,-0.000) (0.004,0.008) (-0.050,0.027)
 (-0.001,-0.023) (-0.041,-0.019) (-0.009,-0.016) (0.021,0.026) (-0.007,0.002) (-0.010,-0.060) (-0.011,-0.058) (-0.009,-0.005) (-0.000,0.015) (-0.010,0.011) (-0.001,0.017) (-0.003,0.007)
 (0.021,0.032) (-0.001,0.020) (-0.006,-0.013) (0.035,0.013) (-0.011,0.014) (0.006,-0.035) (0.122,0.130) (-0.017,-0.014) (-0.011,-0.040) (-0.005,-0.041) (-0.055,-0.038) (-0.010,-0.028)
 (-0.002,-0.019) (0.004,0.013) (0.005,0.002) (-0.001,0.023) (0.098,0.038) (-0.058,-0.009) (-0.008,-0.010) (-0.049,0.011) (0.035,-0.040) (-0.008,0.023) (0.010,0.037) (0.009,0.035)
 (-0.010,-0.009) (-0.011,-0.096) (0.003,0.016) (-0.002,-0.064) (0.004,0.015) (0.034,0.025) (-0.003,-0.003) (-0.004,0.029) (0.038,0.020) (0.002,0.006) (-0.008,-0.006) (-0.008,-0.047)
 (-0.004,-0.000) (-0.001,-0.014) (0.008,0.033) (-0.025,0.013) (0.006,0.015) (-0.046,0.004) (-0.001,0.001) (-0.007,-0.010) (-0.007,-0.016) (0.062,0.054) (-0.003,-0.016) (-0.014,0.003)
 (0.018,0.028) (0.000,-0.014) (-0.011,-0.005) (-0.003,0.031) (-0.010,0.016) (0.211,0.079) (0.011,-0.078) (0.001,-0.049) (0.114,0.139) (0.048,-0.002) (0.019,0.006) (-0.017,-0.001)
 (-0.000,-0.003) (-0.044,-0.014) (-0.005,0.007) (-0.025,-0.020) (0.031,0.030) (0.006,0.009) (0.010,0.054) (0.014,-0.073) (-0.003,-0.007) (-0.000,-0.161) (0.002,-0.013) (-0.063,-0.022)
 (-0.000,-0.023) (-0.002,-0.025) (0.001,-0.003) (-0.006,-0.010) (-0.010,-0.028) (-0.009,-0.008) (-0.000,-0.012) (0.022,0.052) (-0.004,0.003) (-0.005,-0.047) (-0.000,0.029) (0.006,-0.019)
 (-0.005,0.003) (-0.007,0.012) (-0.009,0.044) (-0.002,0.012) (-0.008,-0.042) (0.053,0.010) (0.039,0.050) (0.024,-0.114) (0.035,-0.040) (0.047,0.045) (0.083,0.021) (-0.015,-0.013)
 (0.026,0.031) (-0.010,0.038) (0.008,-0.012) (0.002,-0.018) (0.021,0.028) (0.003,-0.019) (-0.033,0.038) (-0.014,0.007) (-0.018,-0.021) (-0.008,-0.014) (-0.022,-0.022) (0.019,0.009)
 (-0.012,0.000) (-0.019,-0.013) (0.000,-0.002) (-0.042,-0.002) (0.007,-0.024) (-0.018,-0.041) (-0.005,0.054) (-0.013,-0.051) (0.141,0.148) (-0.010,-0.056) (-0.106,-0.037) (0.052,0.034)
 (0.017,0.042) (-0.008,0.019) (-0.289,-0.144) (0.011,0.016) (0.061,0.016) (-0.008,0.022) (0.011,0.022) (0.003,0.008) (0.002,0.051) (0.000,0.002) (-0.004,-0.096) (0.006,0.016)
 (0.025,-0.032) (-0.003,0.005) (-0.022,0.018) (-0.017,0.008) (-0.019,0.014) (-0.070,-0.023) (-0.049,-0.105) (0.063,0.042) (0.003,-0.051) (-0.010,-0.005) (0.045,0.097) (-0.017,-0.003)
 (0.022,0.036) (0.011,-0.014) (-0.042,-0.005) (-0.116,-0.067) (-0.011,-0.015) (-0.011,0.020) (-0.003,-0.002) (0.161,0.071) (0.063,0.055) (0.062,0.058) (0.015,0.008) (-0.025,0.004)
 (-0.108,-0.068) (0.013,-0.042) (0.044,0.064) (-0.013,0.001) (-0.003,0.015) (0.032,0.012) (-0.010,0.006) (-0.020,0.012) (-0.016,0.035) (-0.008,0.018) (-0.010,-0.004) (0.037,0.059)
 (-0.004,0.003) (0.005,0.012) (0.047,0.069) (0.007,0.004) (0.006,0.006) (-0.010,0.006) (0.001,-0.012) (0.124,0.123) (0.025,-0.046) (0.001,-0.011) (0.145,0.166) (0.089,0.049)
 (-0.017,-0.010) (0.007,0.018) (-0.007,0.002) (-0.015,-0.002) (-0.011,-0.021) (-0.016,0.006) (-0.005,-0.005) (-0.004,-0.007) (0.010,-0.033) (0.054,0.108) (-0.004,-0.007) (-0.004,-0.094)
 (0.000,-0.023) (-0.023,-0.011) (0.006,-0.029) (-0.009,0.026) (-0.011,0.016) (-0.028,0.007) (-0.010,0.034) (-0.010,-0.105) (0.011,0.016) (0.020,-0.035) (0.085,0.067) (0.013,-0.010)
 (-0.005,-0.007) (0.002,0.004) (0.005,-0.003) (-0.008,-0.003) (-0.011,0.016) (-0.001,0.000) (-0.001,-0.019) (-0.000,0.016) (-0.031,-0.038) (-0.020,0.041) (-0.051,0.023) (-0.011,0.002)
 (-0.012,0.006) (-0.006,-0.001) (0.011,0.018) (0.080,0.050) (-0.005,0.011) (-0.011,0.020) (0.013,0.027) (-0.030,-0.005) (-0.062,0.000) (-0.018,-0.081) (-0.006,0.005) (-0.014,0.000)
 (-0.011,-0.013) (-0.009,-0.021) (-0.001,-0.005) (0.016,0.003) (-0.005,-0.015) (-0.034,-0.015) (0.007,-0.027) (-0.011,-0.025) (-0.046,0.001) (0.010,-0.017) (0.092,0.106) (0.021,-0.079)
 (-0.149,-0.054) (-0.031,-0.015) (0.002,-0.009) (-0.014,-0.015) (-0.008,-0.041) (-0.013,-0.042) (-0.018,0.015) (-0.018,-0.031) (-0.007,-0.013) (0.005,-0.081) (0.006,-0.039) (-0.005,-0.021)
 (-0.014,-0.165) (-0.002,-0.032) (-0.010,0.034) (-0.002,-0.024) (-0.008,0.004) (-0.007,-0.005) (-0.006,-0.007) (-0.031,0.013) (-0.003,0.004) (-0.032,-0.022) (-0.027,-0.051) (-0.005,-0.011)
 (0.014,0.034) (0.002,-0.022) (0.035,0.017) (-0.008,-0.011) (-0.016,-0.009) (-0.040,-0.028) (-0.008,-0.013) (0.008,-0.044) (-0.006,-0.010) (0.091,0.112) (0.003,0.046) (0.055,0.076)
 (0.001,0.010) (0.032,0.018) (-0.006,0.015) (-0.013,-0.006) (-0.016,0.000) (0.019,0.001) (-0.012,0.001) (-0.024,-0.041) (-0.008,-0.011) (-0.006,0.026) (-0.004,-0.026) (0.009,0.012)
 (0.022,-0.009) (0.013,0.019) (-0.003,-0.016) (-0.007,0.004) (-0.001,0.008) (-0.008,0.002) (-0.020,-0.000) (-0.007,0.015) (-0.005,0.005) (0.057,0.059) (0.007,-0.031) (-0.006,-0.060)
 (0.008,-0.069) (0.053,0.019) (-0.008,0.003) (0.033,0.003) (-0.009,-0.000) (-0.110,-0.086) (-0.009,-0.003) (0.025,-0.034) (-0.006,-0.007) (-0.001,-0.003) (0.001,-0.011) (-0.001,-0.013)
 (-0.013,0.012) (0.014,-0.019) (0.008,-0.041) (0.011,0.049) (0.005,-0.000) (-0.001,0.018) (0.002,0.001) (-0.003,0.010) (-0.001,-0.001) (-0.006,-0.016) (-0.005,-0.016) (0.008,0.040)
 (0.003,0.007) (0.020,0.029) (-0.006,-0.013) (0.021,0.036) (-0.019,-0.010) (-0.018,-0.000) (-0.020,-0.005) (-0.011,-0.004) (-0.006,-0.037) (-0.004,0.021) (0.022,0.080) (-0.022,0.092)
 (-0.026,0.093) (0.022,0.028) (0.014,0.047) (-0.006,-0.037) (0.012,0.026) (0.049,0.058) (0.001,0.014) (-0.008,-0.009) (0.003,0.018) (-0.014,0.009) (-0.059,-0.010) (-0.017,-0.045)
 (0.003,-0.003) (-0.002,-0.005) (0.009,0.039) (0.002,-0.012) (0.012,-0.001) (-0.000,0.046) (-0.002,-0.011) (-0.006,-0.005) (0.004,0.001) (-0.010,0.007) (-0.003,-0.011) (0.002,0.004)
 (0.009,0.007) (-0.007,0.036) (-0.001,0.007) (0.054,0.015) (-0.013,0.005) (0.018,-0.004) (0.036,0.054) (-0.006,0.015) (0.016,0.012) (0.023,0.003) (-0.003,0.003) (-0.012,0.159)
 (0.206,0.314) (-0.011,-0.036) (-0.013,-0.036) (-0.015,-0.030) (0.048,0.077) (-0.009,-0.015) (-0.006,0.020) (-0.007,-0.009) (-0.009,-0.014) (-0.023,-0.019) (-0.023,-0.105) (-0.005,-0.005)
 (-0.009,0.004) (-0.004,0.001) (0.003,-0.006) (0.002,-0.001) (0.019,0.005) (0.001,0.001) (-0.011,-0.005) (-0.040,0.001) (-0.013,-0.004) (-0.001,-0.051) (-0.007,-0.005) (0.157,0.058)
 (0.039,0.028) (0.060,0.038) (0.004,0.001) (0.008,-0.011) (-0.068,-0.046) (-0.003,-0.028) (-0.001,0.009) (0.013,0.007) (-0.005,-0.006) (-0.011,-0.014) (-0.002,0.011) (-0.006,-0.009)
 (0.005,0.059) (-0.020,-0.012) (-0.008,0.012) (-0.016,-0.018) (-0.010,-0.005) (-0.011,0.004) (0.040,-0.011) (-0.002,0.001) (-0.017,-0.017) (-0.006,-0.020) (-0.005,-0.008) (-0.034,0.015)
 (0.007,-0.006) (0.038,0.074) (-0.031,0.002) (-0.061,-0.016) (-0.029,0.001) (-0.008,0.000) (-0.013,0.005) (-0.008,-0.029) (-0.013,-0.016) (0.016,0.011) (-0.020,-0.000) (-0.013,0.004)
 (0.004,0.023) (0.072,0.184) (-0.003,-0.003) (-0.009,-0.028) (-0.015,-0.011) (-0.009,-0.028) (-0.006,-0.029) (-0.005,-0.008) (-0.008,-0.010) (-0.002,-0.010) (-0.012,-0.014) (-0.005,-0.029)
 (-0.019,0.014) (-0.004,0.031) (0.129,0.050) (0.068,0.063) (-0.010,0.027) (-0.064,0.029) (0.016,0.026) (0.000,0.008) (0.061,-0.020) (-0.000,0.024) (-0.014,-0.166) (-0.002,-0.010)
 (0.099,0.061) (0.023,0.042) (0.038,0.053) (-0.001,-0.008) (0.009,-0.004) (-0.027,0.005) (-0.008,0.040) (-0.002,-0.009) (-0.016,-0.024) (-0.042,-0.041) (0.026,0.027) (-0.026,0.002)
 (0.051,0.022) (-0.012,-0.022) (-0.044,-0.027) (-0.024,-0.023) (-0.002,0.015) (0.012,0.020) (0.022,-0.002) (-0.006,-0.011) (0.007,0.022) (-0.008,-0.001) (-0.006,0.020) (0.001,-0.005)
 (0.103,0.087) (0.014,-0.037) (0.006,-0.019) (0.057,0.057) (0.061,0.022) (-0.022,-0.015) (-0.005,0.001) (-0.002,-0.010) (-0.068,-0.042) (0.028,0.008) (-0.027,-0.007) (-0.004,0.035)
 (0.004,0.003) (0.003,-0.049) (0.040,0.048) (-0.005,-0.013) (-0.006,-0.005) (-0.009,0.004) (-0.008,-0.038) (-0.007,-0.013) (-0.011,0.023) (-0.007,-0.005) (-0.001,0.008) (-0.059,-0.035)
 (-0.014,-0.042) (0.052,0.035) (0.005,-0.021) (-0.028,-0.049) (0.042,0.086) (-0.017,0.020) (-0.001,0.042) (0.025,0.001) (-0.057,-0.011) (-0.113,-0.090) (-0.011,-0.009) (-0.017,0.074)
 (-0.014,0.008) (0.003,0.015) (-0.002,0.040) (0.004,0.058) (0.002,0.013) (-0.014,-0.015) (-0.004,-0.011) (0.004,-0.008) (-0.007,-0.010) (-0.003,-0.010) (0.035,0.044) (0.042,0.045)
 (0.015,0.027) (-0.016,-0.006) (-0.003,-0.006) (-0.009,-0.016) (-0.006,-0.022) (0.014,0.041) (0.007,-0.051) (-0.006,0.005) (-0.005,0.125) (0.004,-0.019) (-0.017,-0.004) (-0.009,-0.014)
 (-0.002,0.036) (-0.057,-0.024) (0.001,-0.019) (-0.079,-0.082) (-0.029,-0.015) (-0.013,0.017) (0.020,-0.002) (-0.049,-0.009) (-0.001,0.014) (-0.031,0.052) (-0.014,0.002) (0.004,-0.009)
 (-0.007,0.015) (-0.000,0.003) (-0.005,-0.059) (-0.003,-0.027) (-0.009,0.013) (0.102,0.067) (-0.010,-0.005) (0.082,0.014) (-0.016,-0.010) (-0.007,0.006) (-0.010,-0.006) (-0.009,-0.002)
 (-0.017,-0.001) (-0.003,-0.038) (0.005,0.007) (-0.005,-0.014) (0.042,0.044) (-0.010,-0.012) (-0.003,-0.002) (-0.007,0.009) (0.003,0.014) (0.086,0.004) (-0.048,-0.034) (-0.014,-0.007)
 (-0.006,-0.098) (-0.005,-0.010) (0.057,-0.001) (0.008,0.020) (0.019,0.012) (-0.007,-0.003) (0.010,-0.016) (-0.010,-0.007) (0.008,0.020) (0.057,0.050) (0.008,-0.020) (-0.003,-0.014)
 (-0.037,0.006) (0.007,-0.018) (-0.038,0.017) (-0.012,0.058) (-0.005,-0.015) (-0.005,-0.001) (-0.006,0.008) (-0.003,0.009) (-0.001,-0.032) (-0.003,-0.012) (-0.002,0.016) (-0.016,-0.010)
 (-0.002,0.004) (-0.012,-0.037) (0.007,0.015) (-0.008,-0.015) (-0.004,0.025) (-0.015,0.053) (0.032,-0.004) (0.008,0.003) (-0.017,0.007) (-0.009,-0.000) (-0.215,-0.102) (0.006,0.013)
 (0.069,0.020) (-0.008,0.029) (0.007,0.009) (0.001,0.019) (0.059,0.202) (-0.008,-0.002) (-0.011,0.027) (-0.004,-0.025) (0.013,0.052) (0.001,-0.019) (-0.006,0.010) (-0.007,0.006)
 (-0.006,0.003) (-0.019,0.018) (-0.005,-0.003) (0.008,-0.001) (0.022,-0.001) (-0.018,0.017) (0.009,0.038) (-0.001,-0.010) (0.043,0.029) (-0.014,-0.022) (-0.016,-0.002) (-0.046,-0.017)
 (-0.011,-0.015) (0.004,-0.001) (0.000,0.024) (0.019,-0.011) (0.001,0.037) (-0.000,0.025) (0.006,-0.006) (-0.007,0.002) (-0.012,-0.004) (0.006,-0.002) (0.001,0.007) (-0.048,-0.020)
 (0.001,0.013) (0.047,0.007) (0.004,0.008) (-0.010,0.012) (0.002,0.026) (-0.038,-0.026) (0.021,0.015) (-0.005,-0.006) (0.005,-0.006) (-0.004,0.010) (-0.000,-0.013) (0.002,0.002)
 (-0.003,0.004) (-0.009,-0.008) (0.006,0.026) (-0.049,-0.020) (-0.015,-0.027) (-0.055,-0.063) (-0.017,0.014) (-0.021,0.013) (-0.029,-0.014) (0.000,0.010) (-0.002,0.006) (0.145,0.033)
 (0.009,-0.010) (-0.057,-0.025) (0.005,0.004) (0.006,-0.007) (0.004,-0.105) (-0.000,-0.019) (0.002,-0.013) (0.017,0.042) (0.003,-0.035) (0.069,0.037) (0.001,-0.003) (-0.009,0.012)
 (-0.004,-0.000) (-0.006,-0.003) (-0.039,0.002) (-0.014,0.082) (0.009,0.009) (0.002,0.027) (0.015,0.004) (0.061,0.031) (-0.003,-0.031) (0.030,0.023) (-0.020,-0.016) (-0.021,-0.011)
 (-0.038,-0.018) (-0.004,-0.015) (0.007,-0.037) (0.019,0.026) (0.015,0.047) (0.012,0.051) (0.010,0.043) (0.024,0.034) (0.022,0.049) (-0.002,-0.023) (0.012,0.071) (0.033,0.066)
 (0.025,-0.020) (-0.003,-0.034) (-0.033,0.053) (0.004,-0.022) (-0.096,-0.010) (-0.020,-0.029) (-0.100,-0.068) (-0.002,0.010) (0.010,0.014) (0.003,0.017) (0.005,0.006) (-0.011,-0.120)
 (-0.007,0.014) (-0.000,0.013) (-0.012,0.013) (-0.015,0.004) (0.025,-0.010) (-0.006,-0.036) (0.029,0.077) (-0.032,-0.014) (0.063,-0.006) (0.067,-0.007) (0.001,-0.005) (-0.007,-0.005)
 (-0.074,-0.009) (-0.014,0.006) (0.090,0.054) (-0.014,-0.025) (-0.011,-0.013) (0.010,-0.089) (0.025,0.057) (-0.013,-0.005) (-0.006,-0.102) (0.005,-0.022) (-0.047,-0.018) (-0.005,-0.032)
 (-0.018,-0.025) (-0.001,0.005) (-0.012,-0.001) (-0.039,-0.013) (0.002,-0.077) (0.007,0.029) (0.002,-0.007) (0.066,0.042) (0.002,0.013) (-0.002,0.019) (0.001,-0.003) (0.004,0.023)
 (-0.008,-0.009) (-0.014,-0.026) (-0.001,0.007) (-0.004,-0.002) (-0.001,-0.002) (0.004,-0.014) (0.006,-0.106) (0.008,-0.069) (-0.000,-0.015) (-0.000,-0.037) (0.000,0.010) (-0.001,0.032)
 (0.009,-0.037) (0.000,-0.004) (-0.007,-0.030) (-0.003,-0.030) (-0.010,0.011) (0.021,0.060) (-0.003,0.296) (0.007,-0.010) (0.001,-0.014) (0.014,0.042) (0.006,-0.021) (0.008,0.009)
 (0.013,0.043) (0.005,-0.015) (-0.004,-0.009) (-0.020,0.022) (-0.001,0.029) (0.053,-0.030) (-0.006,-0.009) (-0.068,-0.024) (-0.057,0.054) (0.150,0.051) (0.000,-0.018) (-0.003,-0.001)
 (-0.001,-0.012) (-0.010,-0.037) (-0.005,-0.006) (0.011,0.011) (-0.001,-0.021) (0.003,-0.010) (0.006,0.051) (0.016,-0.000) (0.012,0.028) (-0.007,-0.074) (-0.005,-0.006) (-0.006,0.002)
 (-0.006,-0.000) (-0.020,-0.020) (-0.005,-0.011) (-0.013,-0.013) (-0.049,-0.026) (0.004,-0.134) (0.024,0.046) (-0.006,0.025) (-0.007,0.003) (0.001,-0.038) (0.007,-0.007) (0.000,-0.041)
 (-0.027,0.001) (0.068,0.040) (0.105,0.075) (0.004,0.000) (0.003,-0.111) (-0.000,0.002) (0.036,0.013) (-0.006,0.013) (0.003,0.011) (0.030,0.009) (-0.025,0.015) (-0.009,0.002)
 (0.002,-0.054) (0.044,0.040) (0.020,0.007) (-0.004,-0.009) (-0.036,0.051) (-0.012,-0.001) (-0.023,0.060) (-0.010,0.083) (-0.048,-0.008) (0.012,0.023) (-0.014,0.006) (-0.003,-0.028)
 (-0.005,0.006) (0.034,0.035) (-0.005,-0.014) (-0.031,0.002) (-0.009,0.055) (-0.010,0.076) (0.047,-0.055) (-0.034,-0.029) (-0.011,-0.012) (0.072,0.160) (0.087,-0.005)};
\draw[coloncol!60,line width=.3pt] (axis cs:0.484,0.198) -- (axis cs:0.290,-0.080);
\addplot[only marks,mark=*,mark size=1.6pt,color=coloncol] coordinates { (0.484,0.198) };
\node[font=\sffamily\tiny,anchor=west,color=coloncol,inner sep=1pt] at (axis cs:0.300,-0.080) {\textit{IL7R}--CD127};
\draw[coloncol!60,line width=.3pt] (axis cs:0.629,0.124) -- (axis cs:0.290,-0.165);
\addplot[only marks,mark=*,mark size=1.6pt,color=coloncol] coordinates { (0.629,0.124) };
\node[font=\sffamily\tiny,anchor=west,color=coloncol,inner sep=1pt] at (axis cs:0.300,-0.165) {\textit{FCGR3A}--CD16};
\draw[coloncol!60,line width=.3pt] (axis cs:0.514,0.101) -- (axis cs:0.290,-0.250);
\addplot[only marks,mark=*,mark size=1.6pt,color=coloncol] coordinates { (0.514,0.101) };
\node[font=\sffamily\tiny,anchor=west,color=coloncol,inner sep=1pt] at (axis cs:0.300,-0.250) {\textit{IGHD}--IgD};
\draw[coloncol!60,line width=.3pt] (axis cs:0.515,0.066) -- (axis cs:0.290,-0.335);
\addplot[only marks,mark=*,mark size=1.6pt,color=coloncol] coordinates { (0.515,0.066) };
\node[font=\sffamily\tiny,anchor=west,color=coloncol,inner sep=1pt] at (axis cs:0.300,-0.335) {\textit{KLRB1}--CD161};
\end{axis}
\begin{axis}[at={(b.east)},anchor=west,xshift=1.3cm,name=c,width=.34\textwidth,height=.34\textwidth,xmin=-0.6,xmax=0.8,ymin=-0.6,ymax=0.8,
 xlabel={Held-out correlation},ylabel={Predicted correlation},title={\textbf{c}\enspace Within types, paired},
 axis x line*=bottom,axis y line*=left,tick align=outside,xmajorgrids=true,grid style={gray!18},axis line style={gray!55},tick label style={font=\sffamily\scriptsize},label style={font=\sffamily\small},clip=true,title style={font=\sffamily\small,at={(0,1)},anchor=base west,yshift=5pt}]
\draw[black!35] (axis cs:-0.6,-0.6) -- (axis cs:0.8,0.8);
\addplot[only marks,mark=*,mark size=.35pt,color=donorcol!45] coordinates {
 (-0.019,-0.016) (-0.014,-0.013) (-0.022,-0.018) (-0.054,-0.044) (-0.025,-0.013) (0.007,0.033) (0.005,0.009) (-0.008,-0.007) (-0.032,-0.034) (0.001,0.001) (-0.031,-0.036) (-0.001,0.004)
 (-0.025,-0.020) (0.018,0.035) (0.001,0.002) (-0.006,-0.004) (-0.122,-0.069) (0.080,0.120) (-0.024,-0.001) (-0.040,-0.002) (-0.030,-0.015) (-0.155,-0.131) (-0.015,-0.018) (0.014,0.024)
 (-0.008,0.003) (0.007,0.010) (-0.034,-0.046) (-0.095,-0.098) (-0.010,-0.035) (0.071,0.098) (-0.008,-0.006) (-0.224,-0.184) (-0.019,0.010) (-0.008,-0.016) (-0.001,-0.017) (-0.001,0.016)
 (-0.098,-0.074) (0.002,-0.004) (-0.001,-0.018) (-0.014,-0.017) (-0.010,-0.014) (0.025,0.035) (-0.005,-0.054) (0.026,0.094) (-0.035,-0.021) (0.065,0.077) (-0.029,-0.025) (-0.005,-0.018)
 (0.008,-0.006) (-0.047,-0.058) (-0.023,-0.023) (0.014,0.040) (0.002,-0.024) (-0.005,-0.015) (0.016,0.009) (-0.019,-0.003) (-0.186,-0.132) (0.002,0.004) (0.014,-0.002) (0.052,0.068)
 (0.006,0.015) (-0.025,-0.030) (-0.017,-0.030) (-0.021,-0.028) (-0.065,-0.102) (-0.057,-0.034) (0.083,0.096) (0.038,0.031) (-0.009,-0.020) (0.029,0.049) (-0.001,-0.005) (-0.023,-0.023)
 (0.015,0.009) (0.009,0.016) (0.039,0.045) (0.017,0.014) (-0.006,-0.002) (-0.002,-0.003) (-0.000,0.005) (-0.004,-0.005) (0.002,0.000) (-0.004,-0.003) (-0.005,-0.002) (-0.001,-0.002)
 (0.004,-0.015) (-0.016,-0.018) (0.005,-0.003) (0.078,0.103) (0.044,0.008) (0.021,0.073) (-0.065,-0.080) (-0.003,-0.004) (0.006,-0.019) (-0.021,-0.031) (0.080,0.186) (0.001,0.002)
 (-0.002,-0.016) (0.023,0.010) (-0.005,-0.010) (-0.004,0.010) (0.004,-0.009) (-0.008,-0.006) (-0.008,0.008) (-0.006,-0.011) (0.048,0.031) (0.020,0.013) (-0.036,-0.037) (-0.024,-0.024)
 (0.078,0.081) (-0.017,-0.017) (-0.459,-0.360) (-0.036,0.044) (-0.010,-0.021) (-0.009,-0.009) (0.014,0.024) (0.005,0.007) (-0.005,0.007) (-0.009,-0.021) (-0.002,-0.001) (0.001,0.005)
 (0.001,0.010) (-0.003,0.008) (-0.002,0.003) (-0.009,0.002) (-0.003,0.002) (-0.023,0.004) (-0.002,-0.004) (0.007,0.006) (0.003,-0.002) (0.002,0.012) (-0.062,-0.126) (-0.002,0.009)
 (-0.060,-0.072) (0.012,0.015) (0.016,0.040) (0.059,0.086) (0.009,0.010) (0.005,-0.007) (0.008,0.002) (-0.026,-0.026) (0.001,0.003) (0.003,-0.002) (-0.025,-0.032) (0.001,0.032)
 (-0.005,-0.013) (0.006,-0.006) (-0.030,-0.045) (0.003,0.024) (-0.038,-0.030) (0.086,0.062) (-0.004,-0.013) (-0.008,-0.010) (-0.002,-0.010) (-0.049,-0.042) (0.011,0.007) (0.000,0.005)
 (0.002,-0.002) (-0.002,-0.005) (-0.014,-0.027) (0.017,-0.005) (-0.009,-0.008) (0.009,-0.003) (-0.005,-0.009) (-0.001,-0.009) (-0.001,0.001) (0.058,0.060) (0.043,0.025) (-0.093,-0.081)
 (0.062,0.040) (-0.069,-0.043) (-0.014,-0.006) (-0.030,-0.066) (-0.006,0.009) (0.171,0.108) (-0.016,-0.016) (-0.023,0.003) (0.001,0.006) (-0.003,-0.001) (0.004,0.000) (-0.041,-0.027)
 (0.004,0.004) (-0.066,-0.054) (0.008,0.007) (-0.000,-0.007) (-0.008,-0.012) (0.017,0.021) (-0.060,-0.035) (-0.022,-0.018) (0.017,0.030) (-0.006,0.009) (-0.003,-0.016) (-0.037,-0.035)
 (-0.013,-0.017) (-0.016,-0.022) (0.038,0.027) (-0.059,-0.013) (-0.234,-0.148) (-0.022,-0.022) (-0.020,0.000) (-0.100,-0.059) (0.071,0.095) (-0.048,-0.016) (-0.053,-0.014) (-0.013,-0.008)
 (-0.126,-0.094) (0.001,-0.011) (0.017,0.055) (0.239,0.203) (0.006,0.008) (0.004,-0.022) (-0.075,0.010) (-0.033,-0.071) (-0.091,-0.034) (-0.007,-0.017) (-0.018,-0.047) (-0.014,-0.030)
 (-0.030,-0.024) (-0.000,-0.021) (-0.023,-0.013) (0.034,0.030) (-0.002,-0.013) (-0.052,-0.038) (-0.013,-0.013) (0.026,0.080) (-0.019,-0.052) (-0.034,-0.080) (-0.056,-0.084) (0.004,-0.015)
 (-0.008,-0.016) (-0.018,-0.017) (-0.009,-0.002) (-0.001,-0.001) (-0.030,-0.038) (-0.004,-0.004) (-0.063,-0.009) (-0.002,-0.006) (-0.001,0.005) (0.003,0.025) (-0.006,-0.015) (0.099,0.101)
 (-0.009,-0.014) (0.002,0.024) (-0.083,-0.102) (-0.017,-0.002) (0.001,-0.015) (0.002,-0.008) (-0.000,0.011) (-0.037,-0.056) (-0.006,-0.015) (0.010,0.002) (0.018,0.022) (0.003,0.002)
 (0.100,0.145) (-0.017,-0.015) (0.050,0.087) (0.019,0.000) (-0.014,-0.045) (-0.062,-0.067) (-0.013,-0.022) (-0.012,-0.027) (-0.003,-0.012) (0.235,0.179) (0.092,0.116) (0.105,0.088)
 (0.017,0.022) (0.000,-0.013) (-0.110,-0.090) (0.005,0.006) (-0.006,-0.015) (-0.020,-0.000) (-0.019,-0.009) (0.049,0.044) (-0.007,-0.016) (-0.021,-0.008) (-0.005,-0.012) (-0.314,-0.250)
 (-0.004,0.005) (0.001,-0.015) (-0.005,-0.031) (0.002,0.018) (-0.092,-0.059) (0.002,-0.017) (-0.000,-0.040) (-0.002,-0.002) (-0.005,0.001) (-0.005,-0.005) (-0.002,-0.010) (0.025,0.007)
 (-0.001,0.003) (-0.031,-0.033) (-0.018,-0.009) (-0.023,-0.019) (-0.031,-0.024) (0.004,-0.047) (-0.007,-0.019) (-0.025,-0.011) (-0.035,-0.030) (-0.014,-0.020) (-0.007,-0.013) (0.038,0.047)
 (0.008,-0.032) (-0.077,-0.079) (-0.005,-0.012) (-0.010,-0.016) (-0.002,-0.013) (0.002,-0.003) (0.006,-0.002) (0.005,0.013) (-0.016,-0.012) (-0.013,-0.011) (-0.020,-0.007) (-0.011,0.006)
 (0.001,-0.008) (0.124,0.176) (-0.022,-0.026) (-0.036,-0.020) (0.009,-0.003) (-0.010,-0.027) (0.029,0.011) (0.005,-0.015) (-0.013,-0.014) (-0.006,-0.017) (0.352,0.247) (0.079,0.112)
 (0.150,0.134) (0.002,0.002) (0.044,0.051) (-0.096,-0.073) (0.022,0.062) (0.017,0.063) (0.007,-0.007) (-0.013,-0.042) (0.077,0.156) (-0.010,-0.019) (-0.038,0.003) (-0.011,-0.037)
 (0.035,0.054) (-0.002,-0.007) (-0.005,-0.008) (0.001,-0.004) (-0.001,-0.004) (-0.001,-0.010) (0.010,-0.003) (-0.004,0.001) (-0.009,-0.008) (-0.001,-0.016) (0.109,0.109) (0.014,0.026)
 (0.006,-0.007) (0.149,0.135) (0.051,0.054) (-0.008,-0.005) (-0.008,0.004) (-0.001,0.004) (-0.076,-0.064) (-0.007,-0.012) (-0.009,-0.029) (0.001,0.007) (-0.001,0.003) (0.022,-0.004)
 (0.048,0.080) (0.013,-0.043) (-0.003,-0.014) (0.006,-0.016) (-0.126,-0.066) (-0.011,-0.019) (-0.002,0.005) (-0.000,0.003) (-0.003,-0.008) (-0.005,0.012) (0.000,0.001) (-0.003,0.005)
 (-0.008,-0.024) (-0.004,-0.004) (0.054,0.085) (-0.007,-0.008) (-0.026,-0.004) (-0.001,-0.002) (-0.005,-0.018) (0.028,0.006) (-0.004,-0.012) (-0.013,-0.013) (0.048,0.070) (0.001,0.039)
 (0.023,0.006) (0.040,0.027) (0.046,0.086) (0.078,0.123) (-0.031,-0.058) (0.001,-0.008) (-0.030,-0.035) (0.016,0.006) (-0.009,-0.014) (0.056,0.039) (-0.009,-0.013) (-0.001,0.001)
 (0.005,0.001) (-0.011,-0.005) (0.000,0.002) (0.002,-0.004) (0.001,0.009) (0.001,0.000) (-0.033,-0.018) (-0.000,0.006) (0.006,0.009) (-0.002,-0.022) (-0.010,-0.010) (0.071,0.067)
 (-0.044,-0.079) (0.023,0.085) (-0.047,-0.035) (0.143,0.134) (-0.025,-0.018) (0.045,0.059) (-0.019,-0.013) (-0.421,-0.333) (-0.025,0.025) (0.032,0.024) (-0.010,-0.011) (0.004,0.012)
 (-0.004,-0.020) (0.003,-0.004) (-0.005,-0.019) (0.029,0.035) (-0.003,-0.012) (0.058,0.046) (-0.003,-0.008) (-0.011,-0.002) (-0.015,-0.002) (-0.019,-0.009) (-0.057,0.071) (-0.018,-0.010)
 (-0.020,-0.035) (0.009,0.029) (-0.014,-0.014) (-0.027,-0.069) (0.004,0.006) (0.007,0.016) (-0.007,-0.010) (-0.023,-0.023) (-0.070,-0.049) (-0.018,-0.018) (0.002,0.002) (0.111,0.079)
 (0.088,0.043) (0.100,0.065) (0.102,0.061) (0.045,0.055) (0.126,0.111) (-0.083,-0.077) (-0.002,-0.021) (-0.011,-0.012) (-0.000,-0.003) (0.075,0.081) (-0.010,-0.024) (0.022,0.059)
 (-0.033,-0.049) (-0.007,0.005) (-0.060,-0.050) (-0.072,-0.038) (-0.005,-0.003) (-0.004,-0.000) (-0.004,0.000) (-0.045,-0.033) (-0.012,0.004) (0.000,-0.005) (-0.025,-0.030) (-0.008,0.031)
 (0.028,0.034) (-0.015,-0.075) (0.022,0.028) (-0.084,-0.033) (0.095,0.076) (0.071,0.044) (-0.078,-0.041) (-0.002,-0.001) (-0.034,-0.032) (0.090,0.030) (0.091,0.025) (-0.003,-0.007)
 (0.011,-0.000) (-0.009,-0.023) (0.044,0.045) (0.004,0.002) (-0.014,-0.022) (-0.012,-0.024) (0.099,0.082) (-0.007,-0.025) (-0.000,0.004) (0.001,0.005) (0.005,0.007) (0.010,0.026)
 (0.003,0.022) (-0.042,-0.047) (-0.015,-0.026) (0.514,0.361) (0.006,0.003) (-0.000,-0.008) (-0.001,-0.009) (0.021,0.006) (-0.017,0.002) (-0.034,-0.013) (-0.010,-0.011) (-0.009,-0.003)
 (-0.059,-0.016) (0.042,0.078) (-0.043,-0.015) (-0.037,0.004) (0.027,0.030) (-0.067,-0.020) (-0.005,-0.013) (-0.002,0.003) (-0.001,-0.003) (0.047,0.055) (-0.027,-0.029) (-0.002,-0.007)
 (-0.001,-0.011) (-0.032,-0.037) (-0.005,-0.012) (-0.000,-0.006) (-0.008,-0.007) (0.003,-0.022) (0.016,0.015) (-0.005,-0.007) (-0.072,-0.044) (0.001,-0.004) (-0.022,-0.005) (-0.006,-0.015)
 (-0.008,-0.016) (-0.049,-0.001) (-0.007,-0.033) (0.123,0.167) (-0.039,-0.020) (-0.070,-0.057) (-0.003,-0.004) (-0.003,-0.000) (-0.012,-0.004) (-0.025,-0.037) (-0.007,-0.003) (-0.016,0.024)
 (-0.003,-0.004) (-0.005,0.003) (-0.004,-0.002) (-0.003,0.007) (0.007,0.002) (0.012,0.015) (-0.001,0.008) (-0.007,-0.016) (0.003,0.017) (0.001,-0.002) (-0.005,-0.007) (0.011,0.011)
 (-0.054,-0.031) (-0.020,-0.019) (-0.001,0.018) (-0.007,0.009) (-0.005,-0.019) (-0.042,-0.066) (-0.011,0.013) (-0.079,-0.056) (0.014,0.009) (0.012,0.021) (0.104,0.132) (-0.000,0.009)
 (-0.014,-0.008) (-0.006,-0.018) (-0.024,-0.052) (-0.018,-0.022) (-0.040,-0.051) (-0.009,-0.018) (-0.016,-0.033) (-0.010,-0.010) (0.020,0.019) (0.026,0.023) (-0.021,-0.010) (-0.002,-0.011)
 (0.050,0.054) (-0.013,-0.002) (-0.000,-0.000) (-0.015,-0.016) (-0.005,-0.007) (-0.009,-0.010) (-0.007,-0.025) (0.016,0.004) (-0.012,-0.012) (-0.031,-0.027) (0.005,-0.007) (-0.042,-0.013)
 (-0.004,-0.007) (0.006,-0.008) (0.119,0.140) (-0.006,-0.095) (-0.000,-0.011) (0.047,0.047) (0.027,0.033) (-0.008,0.003) (0.007,0.003) (0.006,-0.003) (-0.005,-0.012) (0.112,0.048)
 (0.031,-0.003) (0.007,0.003) (0.063,0.051) (0.001,0.002) (0.088,0.033) (-0.003,-0.022) (-0.009,-0.010) (-0.000,-0.005) (-0.068,-0.018) (-0.008,-0.010) (-0.003,-0.006) (-0.000,-0.008)
 (-0.003,-0.001) (-0.006,-0.035) (-0.004,-0.006) (0.006,-0.004) (0.011,0.017) (-0.003,-0.002) (0.002,0.008) (-0.001,0.008) (0.015,0.015) (-0.008,-0.011) (0.062,0.067) (-0.045,0.010)
 (-0.007,0.002) (0.004,0.002) (0.022,0.048) (0.093,0.115) (0.006,0.024) (0.036,0.063) (0.062,0.069) (0.034,0.091) (-0.020,-0.034) (0.067,0.005) (-0.017,-0.006) (0.010,0.011)
 (0.021,0.011) (-0.079,-0.051) (0.030,-0.001) (-0.000,0.004) (-0.015,0.002) (-0.058,-0.096) (-0.027,-0.038) (0.005,0.064) (-0.017,-0.018) (-0.011,-0.008) (0.009,-0.003) (-0.005,-0.015)
 (-0.030,-0.041) (-0.017,-0.033) (-0.002,0.030) (0.027,0.036) (-0.020,-0.101) (0.074,0.096) (-0.045,-0.026) (0.019,0.017) (-0.016,-0.000) (0.003,0.006) (0.006,0.008) (0.050,0.053)
 (0.008,0.007) (0.002,-0.013) (-0.032,-0.021) (0.002,-0.000) (-0.022,-0.019) (0.060,0.015) (0.042,0.092) (0.029,0.013) (-0.024,-0.018) (0.202,0.152) (-0.011,-0.025) (-0.005,0.003)
 (-0.003,0.000) (-0.001,-0.003) (-0.001,-0.000) (0.004,0.004) (0.008,0.002) (-0.001,0.000) (0.028,0.009) (0.070,0.072) (-0.005,-0.016) (0.006,0.007) (-0.006,-0.010) (0.000,-0.010)
 (-0.009,-0.035) (-0.002,-0.015) (-0.004,-0.015) (-0.010,-0.020) (-0.047,-0.074) (-0.031,-0.034) (-0.050,-0.053) (-0.008,-0.029) (-0.006,-0.019) (-0.005,-0.010) (0.017,0.013) (-0.124,-0.084)
 (-0.027,-0.016) (-0.003,-0.007) (0.075,0.088) (0.008,0.035) (0.074,0.036) (-0.013,-0.008) (-0.009,-0.025) (-0.015,-0.019) (-0.011,-0.024) (0.017,0.006) (0.004,-0.016) (-0.050,-0.047)
 (-0.001,-0.005) (-0.041,-0.021) (-0.009,-0.009) (0.021,0.049) (-0.007,-0.025) (-0.010,-0.047) (-0.011,-0.026) (-0.009,-0.000) (-0.000,-0.005) (-0.010,0.002) (-0.001,0.007) (-0.003,0.005)
 (0.021,0.020) (-0.001,0.004) (-0.006,-0.015) (0.035,0.037) (-0.011,0.008) (0.006,-0.016) (0.122,0.079) (-0.017,-0.031) (-0.011,-0.013) (-0.005,-0.018) (-0.055,-0.003) (-0.010,-0.019)
 (-0.002,0.015) (0.004,0.009) (0.005,-0.005) (-0.001,0.045) (0.098,0.079) (-0.058,-0.035) (-0.008,-0.000) (-0.049,-0.047) (0.035,0.046) (-0.008,-0.005) (0.010,0.020) (0.009,0.008)
 (-0.010,-0.008) (-0.011,-0.010) (0.003,-0.003) (-0.002,-0.006) (0.004,0.028) (0.034,0.056) (-0.003,-0.002) (-0.004,0.016) (0.038,0.051) (0.002,0.035) (-0.008,-0.019) (-0.008,-0.016)
 (-0.004,-0.004) (-0.001,-0.001) (0.008,0.015) (-0.025,-0.037) (0.006,0.019) (-0.046,-0.034) (-0.001,0.001) (-0.007,-0.021) (-0.007,-0.007) (0.062,0.100) (-0.003,0.005) (-0.014,-0.004)
 (0.018,0.005) (0.000,-0.004) (-0.011,0.006) (-0.003,0.003) (-0.010,-0.013) (0.211,0.220) (0.011,-0.041) (0.001,-0.034) (0.114,0.114) (0.048,0.053) (0.019,0.013) (-0.017,-0.008)
 (-0.000,0.008) (-0.044,-0.036) (-0.005,-0.014) (-0.025,-0.053) (0.031,0.028) (0.006,0.002) (0.010,-0.002) (0.014,0.026) (-0.003,-0.017) (-0.000,-0.006) (0.002,-0.002) (-0.063,-0.036)
 (-0.000,-0.015) (-0.002,-0.006) (0.001,-0.001) (-0.006,-0.004) (-0.010,-0.010) (-0.009,-0.004) (-0.000,-0.004) (0.022,0.025) (-0.004,0.001) (-0.005,-0.018) (-0.000,-0.001) (0.006,0.009)
 (-0.005,-0.004) (-0.007,-0.003) (-0.009,0.003) (-0.002,0.004) (-0.008,0.005) (0.053,0.064) (0.039,0.047) (0.024,0.010) (0.035,0.035) (0.047,0.048) (0.083,0.104) (-0.015,-0.034)
 (0.026,0.025) (-0.010,0.022) (0.008,-0.015) (0.002,-0.014) (0.021,0.054) (0.003,0.006) (-0.033,-0.003) (-0.014,-0.012) (-0.018,-0.033) (-0.008,-0.021) (-0.022,-0.023) (0.019,-0.009)
 (-0.012,-0.014) (-0.019,-0.009) (0.000,-0.012) (-0.042,-0.031) (0.007,-0.006) (-0.018,-0.021) (-0.005,0.026) (-0.013,-0.060) (0.141,0.198) (-0.010,-0.026) (-0.106,-0.085) (0.052,0.018)
 (0.017,0.028) (-0.008,-0.009) (-0.289,-0.209) (0.011,0.025) (0.061,0.044) (-0.008,-0.016) (0.011,0.016) (0.003,-0.001) (0.002,-0.003) (0.000,-0.003) (-0.004,-0.005) (0.006,-0.002)
 (0.025,0.013) (-0.003,-0.001) (-0.022,-0.020) (-0.017,-0.020) (-0.019,-0.016) (-0.070,-0.063) (-0.049,-0.035) (0.063,0.076) (0.003,0.002) (-0.010,-0.008) (0.045,0.105) (-0.017,-0.022)
 (0.022,0.039) (0.011,0.008) (-0.042,-0.041) (-0.116,-0.130) (-0.011,-0.021) (-0.011,-0.026) (-0.003,-0.013) (0.161,0.124) (0.063,0.094) (0.062,0.060) (0.015,0.024) (-0.025,-0.032)
 (-0.108,-0.081) (0.013,0.009) (0.044,0.035) (-0.013,-0.006) (-0.003,-0.007) (0.032,0.037) (-0.010,-0.006) (-0.020,-0.007) (-0.016,-0.012) (-0.008,-0.007) (-0.010,0.005) (0.037,0.041)
 (-0.004,0.005) (0.005,0.005) (0.047,0.025) (0.007,0.002) (0.006,0.020) (-0.010,-0.010) (0.001,-0.014) (0.124,0.143) (0.025,-0.016) (0.001,-0.011) (0.145,0.145) (0.089,0.084)
 (-0.017,-0.011) (0.007,0.011) (-0.007,-0.009) (-0.015,-0.025) (-0.011,-0.006) (-0.016,-0.012) (-0.005,-0.004) (-0.004,-0.003) (0.010,-0.005) (0.054,0.042) (-0.004,-0.006) (-0.004,-0.004)
 (0.000,-0.005) (-0.023,-0.013) (0.006,-0.006) (-0.009,0.000) (-0.011,0.001) (-0.028,-0.008) (-0.010,0.014) (-0.010,-0.004) (0.011,0.001) (0.020,-0.010) (0.085,0.119) (0.013,0.003)
 (-0.005,0.001) (0.002,0.010) (0.005,-0.005) (-0.008,-0.007) (-0.011,0.007) (-0.001,-0.001) (-0.001,-0.002) (-0.000,-0.012) (-0.031,-0.058) (-0.020,-0.021) (-0.051,-0.051) (-0.011,-0.019)
 (-0.012,-0.034) (-0.006,-0.007) (0.011,0.024) (0.080,0.048) (-0.005,0.024) (-0.011,-0.011) (0.013,0.017) (-0.030,-0.014) (-0.062,-0.014) (-0.018,-0.029) (-0.006,-0.014) (-0.014,-0.024)
 (-0.011,-0.019) (-0.009,-0.005) (-0.001,-0.013) (0.016,0.004) (-0.005,0.000) (-0.034,-0.024) (0.007,0.001) (-0.011,-0.017) (-0.046,-0.041) (0.010,-0.004) (0.092,0.139) (0.021,0.005)
 (-0.149,-0.126) (-0.031,-0.011) (0.002,-0.006) (-0.014,-0.017) (-0.008,-0.060) (-0.013,-0.024) (-0.018,0.036) (-0.018,-0.016) (-0.007,-0.009) (0.005,-0.015) (0.006,0.044) (-0.005,-0.028)
 (-0.014,-0.012) (-0.002,-0.015) (-0.010,0.017) (-0.002,-0.025) (-0.008,0.018) (-0.007,0.005) (-0.006,-0.012) (-0.031,0.134) (-0.003,0.011) (-0.032,-0.076) (-0.027,-0.053) (-0.005,-0.029)
 (0.014,0.016) (0.002,-0.005) (0.035,0.036) (-0.008,-0.011) (-0.016,-0.015) (-0.040,-0.052) (-0.008,-0.000) (0.008,-0.002) (-0.006,-0.005) (0.091,0.057) (0.003,0.011) (0.055,0.038)
 (0.001,0.014) (0.032,0.061) (-0.006,-0.006) (-0.013,-0.004) (-0.016,-0.011) (0.019,0.011) (-0.012,-0.008) (-0.024,-0.031) (-0.008,-0.010) (-0.006,-0.005) (-0.004,-0.002) (0.009,0.006)
 (0.022,0.007) (0.013,0.062) (-0.003,-0.016) (-0.007,-0.016) (-0.001,-0.004) (-0.008,-0.008) (-0.020,0.005) (-0.007,-0.004) (-0.005,-0.010) (0.057,0.070) (0.007,-0.049) (-0.006,-0.028)
 (0.008,0.023) (0.053,0.037) (-0.008,-0.006) (0.033,0.018) (-0.009,-0.009) (-0.110,-0.097) (-0.009,0.019) (0.025,0.019) (-0.006,-0.004) (-0.001,0.016) (0.001,-0.009) (-0.001,-0.005)
 (-0.013,-0.009) (0.014,0.031) (0.008,0.001) (0.011,0.063) (0.005,-0.003) (-0.001,-0.006) (0.002,-0.006) (-0.003,-0.001) (-0.001,-0.021) (-0.006,-0.005) (-0.005,-0.004) (0.008,0.016)
 (0.003,-0.001) (0.020,0.032) (-0.006,-0.015) (0.021,0.032) (-0.019,-0.020) (-0.018,-0.018) (-0.020,-0.023) (-0.011,-0.017) (-0.006,-0.010) (-0.004,-0.002) (0.022,0.021) (-0.022,-0.024)
 (-0.026,-0.016) (0.022,0.010) (0.014,0.052) (-0.006,-0.018) (0.012,0.046) (0.049,0.054) (0.001,0.009) (-0.008,-0.024) (0.003,0.052) (-0.014,-0.004) (-0.059,-0.003) (-0.017,-0.027)
 (0.003,0.001) (-0.002,-0.001) (0.009,0.051) (0.002,-0.006) (0.012,0.002) (-0.000,0.001) (-0.002,-0.002) (-0.006,-0.003) (0.004,-0.000) (-0.010,-0.014) (-0.003,-0.004) (0.002,-0.024)
 (0.009,-0.006) (-0.007,-0.004) (-0.001,-0.001) (0.054,0.032) (-0.013,-0.013) (0.018,0.004) (0.036,0.029) (-0.006,0.006) (0.016,-0.016) (0.023,0.006) (-0.003,0.001) (-0.012,-0.009)
 (0.206,0.050) (-0.011,-0.013) (-0.013,-0.007) (-0.015,-0.010) (0.048,0.038) (-0.009,-0.003) (-0.006,-0.004) (-0.007,-0.004) (-0.009,-0.013) (-0.023,-0.024) (-0.023,-0.017) (-0.005,-0.004)
 (-0.009,0.009) (-0.004,-0.004) (0.003,-0.019) (0.002,0.004) (0.019,0.016) (0.001,-0.005) (-0.011,-0.001) (-0.040,-0.022) (-0.013,-0.004) (-0.001,-0.003) (-0.007,-0.005) (0.157,0.110)
 (0.039,0.059) (0.060,0.049) (0.004,0.009) (0.008,-0.001) (-0.068,-0.050) (-0.003,0.001) (-0.001,-0.003) (0.013,0.004) (-0.005,-0.010) (-0.011,-0.003) (-0.002,-0.004) (-0.006,-0.014)
 (0.005,0.002) (-0.020,-0.026) (-0.008,-0.010) (-0.016,-0.001) (-0.010,-0.016) (-0.011,-0.001) (0.040,0.034) (-0.002,-0.009) (-0.017,-0.015) (-0.006,-0.011) (-0.005,-0.002) (-0.034,-0.010)
 (0.007,-0.004) (0.038,0.073) (-0.031,-0.019) (-0.061,-0.045) (-0.029,-0.020) (-0.008,-0.006) (-0.013,-0.013) (-0.008,-0.033) (-0.013,-0.017) (0.016,0.004) (-0.020,-0.013) (-0.013,-0.007)
 (0.004,-0.013) (0.072,0.060) (-0.003,-0.030) (-0.009,-0.013) (-0.015,-0.017) (-0.009,0.013) (-0.006,-0.009) (-0.005,-0.002) (-0.008,-0.001) (-0.002,0.002) (-0.012,0.014) (-0.005,0.008)
 (-0.019,-0.015) (-0.004,-0.009) (0.129,0.117) (0.068,0.107) (-0.010,-0.011) (-0.064,-0.039) (0.016,0.012) (0.000,0.002) (0.061,0.061) (-0.000,0.002) (-0.014,-0.013) (-0.002,-0.008)
 (0.099,0.071) (0.023,0.027) (0.038,0.030) (-0.001,0.002) (0.009,0.024) (-0.027,-0.018) (-0.008,0.008) (-0.002,0.003) (-0.016,-0.025) (-0.042,-0.033) (0.026,0.021) (-0.026,-0.019)
 (0.051,0.067) (-0.012,-0.005) (-0.044,-0.027) (-0.024,-0.017) (-0.002,-0.004) (0.012,0.011) (0.022,0.020) (-0.006,0.009) (0.007,0.004) (-0.008,0.007) (-0.006,-0.004) (0.001,-0.017)
 (0.103,0.117) (0.014,-0.033) (0.006,-0.009) (0.057,0.064) (0.061,0.059) (-0.022,-0.015) (-0.005,0.007) (-0.002,-0.005) (-0.068,-0.060) (0.028,0.009) (-0.027,-0.030) (-0.004,0.007)
 (0.004,0.006) (0.003,-0.012) (0.040,0.054) (-0.005,-0.024) (-0.006,-0.009) (-0.009,-0.011) (-0.008,0.013) (-0.007,-0.017) (-0.011,-0.012) (-0.007,-0.013) (-0.001,0.011) (-0.059,-0.041)
 (-0.014,-0.021) (0.052,0.038) (0.005,0.010) (-0.028,-0.034) (0.042,0.122) (-0.017,-0.033) (-0.001,0.013) (0.025,0.007) (-0.057,-0.048) (-0.113,-0.129) (-0.011,-0.028) (-0.017,-0.024)
 (-0.014,-0.004) (0.003,-0.000) (-0.002,-0.003) (0.004,-0.003) (0.002,-0.001) (-0.014,-0.020) (-0.004,0.002) (0.004,-0.005) (-0.007,-0.007) (-0.003,-0.010) (0.035,0.040) (0.042,0.027)
 (0.015,0.036) (-0.016,-0.025) (-0.003,0.002) (-0.009,-0.010) (-0.006,-0.006) (0.014,0.043) (0.007,0.001) (-0.006,-0.010) (-0.005,-0.014) (0.004,-0.001) (-0.017,0.009) (-0.009,-0.001)
 (-0.002,0.015) (-0.057,-0.043) (0.001,-0.038) (-0.079,-0.103) (-0.029,-0.017) (-0.013,-0.016) (0.020,0.006) (-0.049,-0.015) (-0.001,-0.001) (-0.031,-0.018) (-0.014,0.001) (0.004,-0.009)
 (-0.007,-0.007) (-0.000,0.007) (-0.005,-0.006) (-0.003,-0.008) (-0.009,0.001) (0.102,0.069) (-0.010,-0.006) (0.082,0.050) (-0.016,-0.013) (-0.007,0.002) (-0.010,-0.007) (-0.009,0.001)
 (-0.017,0.003) (-0.003,0.005) (0.005,-0.002) (-0.005,-0.008) (0.042,0.042) (-0.010,0.001) (-0.003,-0.005) (-0.007,-0.006) (0.003,-0.005) (0.086,0.084) (-0.048,0.005) (-0.014,-0.008)
 (-0.006,0.002) (-0.005,-0.001) (0.057,0.034) (0.008,0.012) (0.019,0.009) (-0.007,-0.004) (0.010,0.011) (-0.010,-0.012) (0.008,0.030) (0.057,0.067) (0.008,-0.003) (-0.003,-0.013)
 (-0.037,0.006) (0.007,-0.024) (-0.038,-0.013) (-0.012,-0.008) (-0.005,-0.012) (-0.005,-0.006) (-0.006,-0.009) (-0.003,-0.012) (-0.001,0.030) (-0.003,0.019) (-0.002,-0.010) (-0.016,-0.007)
 (-0.002,-0.011) (-0.012,-0.007) (0.007,0.018) (-0.008,-0.024) (-0.004,0.021) (-0.015,-0.012) (0.032,0.038) (0.008,0.006) (-0.017,0.001) (-0.009,-0.009) (-0.215,-0.176) (0.006,0.031)
 (0.069,0.057) (-0.008,-0.011) (0.007,0.026) (0.001,-0.006) (0.059,0.022) (-0.008,-0.005) (-0.011,-0.006) (-0.004,-0.004) (0.013,0.022) (0.001,-0.004) (-0.006,0.000) (-0.007,0.003)
 (-0.006,-0.004) (-0.019,0.043) (-0.005,-0.004) (0.008,-0.004) (0.022,0.024) (-0.018,-0.019) (0.009,0.039) (-0.001,-0.012) (0.043,0.050) (-0.014,-0.021) (-0.016,-0.026) (-0.046,-0.058)
 (-0.011,-0.015) (0.004,0.004) (0.000,0.001) (0.019,0.012) (0.001,0.004) (-0.000,0.000) (0.006,0.004) (-0.007,0.001) (-0.012,-0.007) (0.006,-0.008) (0.001,0.005) (-0.048,-0.036)
 (0.001,-0.004) (0.047,0.051) (0.004,0.005) (-0.010,-0.020) (0.002,0.014) (-0.038,-0.038) (0.021,0.012) (-0.005,-0.004) (0.005,-0.008) (-0.004,0.009) (-0.000,0.010) (0.002,-0.008)
 (-0.003,-0.007) (-0.009,-0.007) (0.006,0.038) (-0.049,-0.037) (-0.015,-0.049) (-0.055,-0.057) (-0.017,-0.022) (-0.021,-0.015) (-0.029,-0.012) (0.000,-0.000) (-0.002,0.006) (0.145,0.108)
 (0.009,-0.007) (-0.057,-0.038) (0.005,0.006) (0.006,0.001) (0.004,-0.013) (-0.000,-0.007) (0.002,0.001) (0.017,0.024) (0.003,-0.005) (0.069,0.066) (0.001,-0.008) (-0.009,-0.005)
 (-0.004,-0.006) (-0.006,-0.004) (-0.039,-0.010) (-0.014,-0.005) (0.009,0.012) (0.002,0.008) (0.015,0.015) (0.061,0.047) (-0.003,-0.012) (0.030,0.023) (-0.020,-0.015) (-0.021,-0.021)
 (-0.038,-0.053) (-0.004,-0.008) (0.007,0.001) (0.019,0.017) (0.015,0.035) (0.012,0.004) (0.010,0.021) (0.024,0.028) (0.022,0.055) (-0.002,-0.006) (0.012,0.023) (0.033,0.077)
 (0.025,0.006) (-0.003,-0.018) (-0.033,0.038) (0.004,-0.014) (-0.096,-0.022) (-0.020,-0.016) (-0.100,-0.071) (-0.002,0.015) (0.010,-0.001) (0.003,-0.015) (0.005,0.017) (-0.011,-0.010)
 (-0.007,-0.012) (-0.000,-0.008) (-0.012,-0.017) (-0.015,-0.022) (0.025,0.030) (-0.006,-0.036) (0.029,0.081) (-0.032,-0.015) (0.063,0.065) (0.067,0.012) (0.001,0.005) (-0.007,-0.005)
 (-0.074,-0.043) (-0.014,-0.015) (0.090,0.064) (-0.014,-0.009) (-0.011,-0.011) (0.010,-0.007) (0.025,0.046) (-0.013,-0.021) (-0.006,-0.014) (0.005,-0.013) (-0.047,-0.021) (-0.005,-0.012)
 (-0.018,-0.005) (-0.001,-0.003) (-0.012,-0.011) (-0.039,-0.019) (0.002,0.014) (0.007,0.024) (0.002,0.005) (0.066,0.077) (0.002,-0.012) (-0.002,0.006) (0.001,0.000) (0.004,0.009)
 (-0.008,-0.004) (-0.014,-0.006) (-0.001,0.004) (-0.004,0.000) (-0.001,0.000) (0.004,-0.003) (0.006,0.004) (0.008,0.002) (-0.000,-0.007) (-0.000,0.006) (0.000,-0.002) (-0.001,0.007)
 (0.009,0.004) (0.000,-0.004) (-0.007,-0.006) (-0.003,-0.004) (-0.010,-0.007) (0.021,0.012) (-0.003,-0.002) (0.007,-0.006) (0.001,0.004) (0.014,0.029) (0.006,0.002) (0.008,-0.002)
 (0.013,0.012) (0.005,-0.001) (-0.004,0.003) (-0.020,-0.018) (-0.001,0.004) (0.053,0.042) (-0.006,-0.017) (-0.068,-0.051) (-0.057,-0.025) (0.150,0.150) (0.000,-0.006) (-0.003,-0.003)
 (-0.001,-0.001) (-0.010,-0.017) (-0.005,-0.006) (0.011,0.018) (-0.001,-0.010) (0.003,-0.005) (0.006,0.003) (0.016,0.015) (0.012,0.038) (-0.007,-0.013) (-0.005,0.001) (-0.006,-0.004)
 (-0.006,-0.008) (-0.020,-0.018) (-0.005,-0.012) (-0.013,-0.019) (-0.049,-0.041) (0.004,0.010) (0.024,0.026) (-0.006,-0.004) (-0.007,-0.009) (0.001,0.019) (0.007,0.016) (0.000,-0.001)
 (-0.027,-0.013) (0.068,0.094) (0.105,0.199) (0.004,0.019) (0.003,0.012) (-0.000,0.009) (0.036,0.032) (-0.006,-0.003) (0.003,0.004) (0.030,0.042) (-0.025,-0.015) (-0.009,-0.008)
 (0.002,0.009) (0.044,0.042) (0.020,0.016) (-0.004,-0.009) (-0.036,-0.011) (-0.012,-0.016) (-0.023,-0.004) (-0.010,-0.006) (-0.048,-0.029) (0.012,0.016) (-0.014,0.007) (-0.003,-0.017)
 (-0.005,-0.002) (0.034,0.023) (-0.005,-0.007) (-0.031,-0.020) (-0.009,-0.014) (-0.010,-0.008) (0.047,0.064) (-0.034,-0.014) (-0.011,-0.008) (0.072,0.069) (0.087,0.086)};
\draw[coloncol!60,line width=.3pt] (axis cs:0.629,0.523) -- (axis cs:0.290,-0.080);
\addplot[only marks,mark=*,mark size=1.6pt,color=coloncol] coordinates { (0.629,0.523) };
\node[font=\sffamily\tiny,anchor=west,color=coloncol,inner sep=1pt] at (axis cs:0.300,-0.080) {\textit{FCGR3A}--CD16};
\draw[coloncol!60,line width=.3pt] (axis cs:0.484,0.387) -- (axis cs:0.290,-0.165);
\addplot[only marks,mark=*,mark size=1.6pt,color=coloncol] coordinates { (0.484,0.387) };
\node[font=\sffamily\tiny,anchor=west,color=coloncol,inner sep=1pt] at (axis cs:0.300,-0.165) {\textit{IL7R}--CD127};
\draw[coloncol!60,line width=.3pt] (axis cs:0.514,0.361) -- (axis cs:0.290,-0.250);
\addplot[only marks,mark=*,mark size=1.6pt,color=coloncol] coordinates { (0.514,0.361) };
\node[font=\sffamily\tiny,anchor=west,color=coloncol,inner sep=1pt] at (axis cs:0.300,-0.250) {\textit{IGHD}--IgD};
\draw[coloncol!60,line width=.3pt] (axis cs:0.515,0.296) -- (axis cs:0.290,-0.335);
\addplot[only marks,mark=*,mark size=1.6pt,color=coloncol] coordinates { (0.515,0.296) };
\node[font=\sffamily\tiny,anchor=west,color=coloncol,inner sep=1pt] at (axis cs:0.300,-0.335) {\textit{KLRB1}--CD161};
\end{axis}
\end{tikzpicture}
}
\caption{\textbf{Gene--protein dependence between and within cell types in the confirmatory cohort.}
Held-out cross-correlations against predictions, averaged over the \XsHaoDonors{} donors, for every sixteenth of the \XsGenes{}~$\times$~\XsProteins{} gene--protein pairs (grey) and for labelled pairs (coloured). \textbf{a}, All cells, pairing-free channel. \textbf{b}, Within level-1 cell types, pairing-free channel. \textbf{c}, Within level-1 cell types, paired closed form. Red, the four genes with the largest held-out within-type correlation with their own protein; blue, two lineage pairs. The diagonal marks perfect prediction. The labelled pairs and averages are descriptive and were chosen after scoring.}
\label{sfig:entries}
\end{figure}


\section{Learning without pairs in a perturbation screen}\label{sec:pert}
This note describes the prespecified test of whether intervention-induced population
variation lets RNA--protein dependence within conditions be recovered without
paired cells.

\subsection*{Data and split}
The Perturb-CITE-seq screen of Frangieh et al.\cite{frangieh2021} was obtained
from the scPerturb release\cite{peidli2024} (Zenodo record 7041849): \PtCells{} patient-derived melanoma cells
with CRISPR knockouts of \PtTargets{} genes, cultured alone, with interferon-$\gamma$
or with autologous tumour-infiltrating lymphocytes, with RNA and 24
antibody-derived tags per cell. The release has no replicate or batch
annotation. Targets were read from the detected guides, because the release's
perturbation column names only the first of several targeted genes and labels
cells with targeting and non-targeting guides as controls. We used the
\PtSingleCells{} cells whose guides all targeted one gene and the \PtNTCells{}
cells whose guides were all non-targeting.

Every fourth target in a fixed pseudo-random order was held out (\PtHeldTargets{} targets), and
every cell whose guides targeted a held-out gene was excluded from training and
tuning. Within each training population (target by condition), cells in a fixed
pseudo-random order formed an RNA half and a protein half. Within each
held-out group and each condition of the non-targeting cells, the first half
formed the adaptation cells and the second half the scoring cells, which
alternated between sub-halves A and B. The gene panel was fixed from training
cells only: the \PtEncodingGenes{} genes that encode proteins targeted by the
antibodies and were detected in at least 1\% of training cells, then the genes of highest dispersion
(detected in at least 5\%, not mitochondrial or ribosomal, log dispersion
z-scored within 20 bins of mean log expression) up to \PtGenes{} genes. Protein
features were the centred log ratios of the \PtProteins{} target antibodies. The
split files were recorded before any statistic of held-out or
non-targeting cells other than library sizes was computed.

\subsection*{Protocol}
The protocol was written after a development analysis on
training targets only, in which every fourth training target played the held-out
role (47 targets, 118 groups). There, the channel fitted to
uncentred population means supported two between-condition directions and
recovered \PtDevPooled\% of within-condition dependence, the condition-centred
channel \PtDevPrimary\%, and the paired closed form \PtDevPaired\%; the
condition-centred channel was chosen as primary, because the target is the dependence within
conditions, and the code was run end to end on this split. The protocol, the
code and the training fits were then registered. Predictions for
eligible groups (at least 40 adaptation cells) were computed from adaptation
cells, each assay separately, and recorded
before the scoring files were opened; the evaluation recomputed every
prediction and checked it against the record. There were no amendments.

Arms were the primary condition-centred pairing-free channel with measured RNA
correlation; its matched controls on pseudo-populations and derangements (ten
draws each); the paired closed form, reference regression and pooled
within-population cross-correlation of the condition, fitted on paired training
cells; independence; and, as secondary arms, the uncentred channel with its own
controls, the primary channel with noise-corrected RNA correlation, with its
regression form and with the pooled training correlation matrices of the
condition instead of the group's, the paired closed form with corrected RNA
correlation, and ridge regression on the group's own adaptation pairs. The
recovered fraction (Methods) removes scoring noise from both the gain over
independence and the size of the target, so that arms are compared on the true
within-group dependence; the mean relative Frobenius error, which scoring noise
inflates towards one for every arm, is reported alongside.

\subsection*{Results}
All \PtGroups{} eligible held-out groups came from \PtHeldEligible{} of the
\PtHeldTargets{} held-out targets. The primary channel supported \PtDimS{}
direction (the uncentred channel \PtDimSPooled), and the direction held a median
of \PtCoverage\% of the groups' corrected RNA variance (\PtCoverageNt\% for
non-targeting cells). H1 (recovery above zero) was \PtHOne, H2 (recovery above
both controls) was \PtHTwo, H3 (recovery in non-targeting cells) was \PtHThree{}
and H4 (at least half of the paired closed form's recovery) was \PtHFour. All arms
are in Supplementary Table~\ref{tab:pert}.

\begin{table}[h]
\caption{\textbf{Recovered fraction of within-condition dependence (\%).} Held-out
targets: \PtGroups{} target-by-condition groups, 95\% intervals from 2,000
resamples of targets; conditions are cultured alone, with interferon-$\gamma$ and
with T cells. Non-targeting cells: three conditions pooled, intervals from 2,000
resamples of guide sets. Error, mean relative Frobenius error over held-out
groups against all scoring cells (independence, 1), inflated by scoring noise.
Controls are averages over ten draws.}
\label{tab:pert}
\centering\scriptsize\setlength{\tabcolsep}{3pt}
% Generated by extension/perturbation/make_assets.py. Do not edit.
\begin{tabular}{lcccccc}
\toprule
& \multicolumn{4}{c}{Held-out targets} & \multicolumn{2}{c}{Non-targeting cells} \\
\cmidrule(lr){2-5}\cmidrule(lr){6-7}
Arm & All (95\% CI) & Alone & IFN$\gamma$ & T cells & All (95\% CI) & Error \\
\midrule
No pairs: perturbations, measured RNA (primary) & 6.8 (4.1 to 9.4) & 11.5 & 16.2 & \ensuremath{-}5.4 & 18.6 (16.1 to 21.1) & 0.992 \\
No pairs: perturbations, corrected RNA & 7.0 (4.4 to 9.5) & 11.6 & 16.1 & \ensuremath{-}4.8 & 17.9 (15.5 to 20.3) & 0.991 \\
No pairs: perturbations, regression form & 9.2 (7.9 to 10.5) & 6.0 & 13.2 & 8.0 & -- & 0.989 \\
No pairs: perturbations, training marginals & 15.1 (13.0 to 17.2) & 15.2 & 24.2 & 6.7 & 18.9 (16.5 to 21.3) & 0.982 \\
No pairs: pseudo-populations (10 draws) & \ensuremath{-}0.1 (\ensuremath{-}0.1 to \ensuremath{-}0.1) & -- & -- & -- & -- & -- \\
No pairs: derangements (10 draws) & \ensuremath{-}0.4 (\ensuremath{-}0.6 to \ensuremath{-}0.3) & -- & -- & -- & -- & -- \\
No pairs: conditions pooled & \ensuremath{-}36.4 (\ensuremath{-}42.0 to \ensuremath{-}31.6) & \ensuremath{-}19.7 & \ensuremath{-}26.6 & \ensuremath{-}58.3 & \ensuremath{-}11.9 (\ensuremath{-}16.1 to \ensuremath{-}7.7) & 1.039 \\
No pairs: conditions pooled, pseudo-populations & \ensuremath{-}46.4 (\ensuremath{-}52.3 to \ensuremath{-}41.3) & -- & -- & -- & -- & -- \\
No pairs: conditions pooled, derangements & \ensuremath{-}75.4 (\ensuremath{-}82.0 to \ensuremath{-}69.8) & -- & -- & -- & -- & -- \\
Paired closed form & 37.3 (31.7 to 42.2) & 23.5 & 42.8 & 43.0 & 94.3 (92.8 to 95.7) & 0.953 \\
Paired closed form, corrected RNA & 39.7 (34.4 to 44.4) & 27.4 & 44.8 & 44.6 & -- & 0.950 \\
Reference regression & 53.4 (49.4 to 57.2) & 43.8 & 54.8 & 59.7 & 91.3 (89.6 to 92.9) & 0.935 \\
Transferred correlation & 96.7 (92.7 to 100.7) & 98.3 & 96.6 & 95.5 & 98.5 (97.0 to 99.8) & 0.884 \\
Own paired adaptation cells & 24.0 (19.8 to 27.7) & 9.5 & 25.3 & 34.1 & 89.7 (88.1 to 91.1) & 0.969 \\
\bottomrule
\end{tabular}

\end{table}

\subsection*{Programs and perturbation choice}
These analyses were
specified after the prespecified evaluation and registered before any of their statistics was computed; the
held-out scoring cells had already been opened for the aggregate endpoint, so the analyses are
exploratory. MSigDB Hallmark gene sets were downloaded from gsea-msigdb.org. The interferon block comprised \PgIfnGenes{}
panel genes and the antibodies to HLA class~I (HLA-A, -B, -C and B2M), PD-L1 and
CD47, whose genes belong to the interferon sets. Prespecified criteria: P1, a
pairing-free share of the paired closed form above one half in the interferon
block, was \PgPOne; P2, the same along the supported direction, was \PgPTwo{}
(share \PgAlongShare\%; \PgAlongShareLo{} to \PgAlongShareHi\%); P3, higher
recovery in the interferon block than in the complement (difference \PgDiff\%;
\PgDiffLo--\PgDiffHi\%), was \PgPThree. The support of the interferon program, the
mean squared correlation of its genes with the RNA score along the supported
direction in adaptation cells, was \PgSupportIfn, against at most
\PgSupportOtherMax{} for the \PgOtherCount{} other programs (Supplementary
Table~\ref{tab:programs}).

For perturbation choice, development on training targets only
showed that channels fitted on 10 to 40 targets, selected at
random or by exposure, retained many spurious directions and were often worse
than independence, whereas removing a few targets from the full set left the
estimator stable; the prespecified analysis therefore removed targets. The four
targets of highest exposure, \DsTopFour, had exposures at least \DsExposureGap{}
times that of any other target. After removal of the 5, 10 and 20 targets of
highest exposure, the channel supported \DsFiveDimExp, \DsTenDimExp{} and
\DsTwentyDimExp{} directions, against a mean of \DsFiveDimRand, \DsTenDimRand{} and
\DsTwentyDimRand{} after random removals, none of which removed all. Recovery over
all genes and proteins after removing the 5 of highest exposure was \DsFiveAllExp\%
(random removals, mean \DsFiveAllRand\%, 5th percentile \DsFiveAllRandLow\%), and
in the interferon block \DsFiveIfnExp\% (random \DsFiveIfnRand\%); after removing
10, \DsTenAllExp\% (random \DsTenAllRand\%) and \DsTenIfnExp\% (random
\DsTenIfnRand\%; \DsTenIfnBelow{} of \DsDraws{} random removals were at or below
it). Removing the 5 and 10 targets of strongest effect (\DsTopStrengthFive{} first)
left interferon-block recovery at \DsFiveIfnStr\% and \DsTenIfnStr\%
(random minus strongest, \DsFiveIfnDiffStr\% and \DsTenIfnDiffStr\%; 95\% CI,
\DsFiveIfnDiffStrLo--\DsFiveIfnDiffStrHi\% and \DsTenIfnDiffStrLo--\DsTenIfnDiffStrHi\%).
Effect strength was a comparator, not a criterion. The design criterion was \DsDOne. Leaving single targets out, the loss of
recovery was largest for \DsLooTop, but its rank correlation with exposure over all
targets was \DsLooSpearman{} (permutation $P=\DsLooP$), because targets that move
the same direction substitute for each other.

\begin{table}[h]
\caption{\textbf{Recovery by program (\%).} Held-out target-by-condition groups of
the melanoma screen and non-targeting cells. Blocks: interferon genes and
antibodies; the RNA score along the supported direction with all antibodies;
interferon genes with all antibodies; all other genes and antibodies; genes of
other Hallmark programs with all antibodies. Arms: no pairs, the primary pairing-free
channel; paired CF, paired closed form; ref.\ regr., reference regression; transf.\
corr., pooled within-condition correlation of paired training cells; own pairs, ridge
regression on the group's paired adaptation cells; replicate, the adaptation half
used directly as the estimate. Rel., split-half reliability of the scoring target
(pooled over groups).}
\label{tab:programs}
\centering\scriptsize\setlength{\tabcolsep}{3pt}
% Generated by extension/perturbation/make_assets.py. Do not edit.
\begin{tabular}{lccccccc}
\toprule
Block & No pairs & Paired CF & Ref.\ regr. & Transf.\ corr. & Own pairs & Replicate & Rel.\ \\
\midrule
Interferon block & 42.5 & 64.2 & 68.8 & 94.0 & 48.3 & \ensuremath{-}24.3 & 0.29 \\
Along $S$ & 27.3 & 69.3 & 76.3 & 87.0 & 47.4 & 6.6 & 0.35 \\
Interferon genes, all proteins & 16.4 & 33.5 & 49.1 & 92.3 & 23.9 & \ensuremath{-}281.3 & 0.12 \\
Complement & 5.7 & 36.2 & 54.7 & 97.9 & 18.7 & \ensuremath{-}308.3 & 0.11 \\
E2F targets & \ensuremath{-}5.1 & \ensuremath{-}18.6 & 4.6 & 100.6 & \ensuremath{-}43.9 & \ensuremath{-}554.5 & 0.07 \\
Epithelial--mesenchymal transition & 7.8 & 56.1 & 68.7 & 96.3 & 38.7 & \ensuremath{-}123.2 & 0.18 \\
G2M checkpoint & \ensuremath{-}7.5 & \ensuremath{-}9.4 & 12.3 & 101.3 & \ensuremath{-}43.8 & \ensuremath{-}525.2 & 0.07 \\
Hypoxia & 3.9 & 32.9 & 50.2 & 94.3 & 24.2 & \ensuremath{-}275.1 & 0.12 \\
TNF-$\alpha$ signalling via NF-$\kappa$B & 5.4 & 41.0 & 57.4 & 99.1 & 29.0 & \ensuremath{-}238.2 & 0.13 \\
\midrule
Non-targeting: interferon block & 61.3 & 92.7 & 90.8 & 97.9 & 89.3 & 94.2 & 0.95 \\
Non-targeting: complement & 12.4 & 95.3 & 92.3 & 98.6 & 89.2 & 87.5 & 0.80 \\
\bottomrule
\end{tabular}

\end{table}

\subsection*{Limitations of this test}
The screen used one patient-derived model without biological replicates, so
intervals reflect variation among targets and guides only. Held-out groups had a
median of \PtAdaptMedian{} adaptation cells, which limits the accuracy of the
within-assay covariance matrices that closed-form transfer uses and favours
arms that do not use them. The panel had 20 surface proteins.

\section{Replication in an independent screen}\label{sec:rep}
This note describes the prespecified replication in the ECCITE-seq screen of Papalexi
et al.\cite{papalexi2021}.

\subsection*{Data, split and protocol}
The scPerturb release\cite{peidli2024} (Zenodo record 7041849) contains 20,729 THP-1 cells stimulated with interferon-$\gamma$,
with knockouts of \RpTargets{} genes and non-targeting guides, three biological
replicates (the 77 cells of a fourth hashtag were not used; \RpCells{} cells
remained), RNA and four surface proteins. Each target-by-replicate group was
split in a fixed pseudo-random order of cells into a training half (RNA and protein
halves) and a test half (adaptation and scoring cells); non-targeting cells were
split into adaptation and scoring cells only. The gene panel was chosen from
training halves by the prespecified rule (\RpGenes{} genes, including the four genes
encoding the antibodies). The protocol and the code were registered after the split and before any statistic of test or
non-targeting cells other than library sizes was computed; the pipeline had been
run on training halves only to check the code. Predictions were computed from
adaptation cells and recorded, and the scoring files were
then opened. There were no amendments.

The prespecified rules of the melanoma screen were reused, with the replicate in the
role of the condition. Every target with an eligible test group was held out in
turn; its channels, controls and paired comparators were fitted on the training
halves of the other targets. For the design criterion, each fold's channel was
also refitted without the five training targets of highest exposure, without the
five of strongest effect and without five random targets (50 draws). The
interferon block comprised \RpIfnGenes{} panel genes of the Hallmark
interferon sets and the antibodies to CD86 and PD-L1.

\subsection*{Results}
Cross-validation over targets chose the smallest penalty scale in every fold,
which retained \RpDimSMedian{} to \RpDimSMax{} supported directions. The five
prespecified criteria were not met: recovery was \RpPrimary\% (95\% CI,
\RpPrimaryLo{} to \RpPrimaryHi\%) over all genes and proteins, lower than both
controls (pseudo-populations \RpAllPseudo\%, derangements \RpAllDerange\%), the
interferon-block share of the paired closed form was \RpIfnShare\%
(\RpIfnShareLo{} to \RpIfnShareHi\%), recovery in non-targeting cells was \RpNtPrimary\%
(\RpNtPrimaryLo{} to \RpNtPrimaryHi\%), and removing the five targets of highest
exposure did not lower recovery relative to random removals (difference
\RpRemDiff\%; \RpRemDiffLo{} to \RpRemDiffHi\%). Paired comparators recovered
\RpPaired\% (closed form) and \RpAllTrans\% (transferred correlation) over all genes
and proteins, and \RpIfnPaired\% and \RpIfnTrans\% in the interferon block. By
replicate, interferon-block recovery of the pairing-free channel was
\RpRepOneIfnPf\%, \RpRepThreeIfnPf\% and \RpRepFourIfnPf\%, against \RpRepOneIfnPaired\%,
\RpRepThreeIfnPaired\% and \RpRepFourIfnPaired\% for the paired closed form.

After the scoring files had been opened, the channel was refitted in every fold with the penalty scale
fixed (exploratory). At 0.1, the value chosen in the
melanoma screen, it retained \RpFixedPointOneDim{} direction and recovered
\RpFixedPointOneAll\% over all genes and proteins and \RpFixedPointOneIfn\% in the
interferon block; at 1, \RpFixedOneDim{} directions, \RpFixedOneAll\% and
\RpFixedOneIfn\%.

\begin{table}[h]
\caption{\textbf{Recovered fraction in the replication screen (\%).} Test targets:
\RpGroups{} target-by-replicate groups from \RpTestTargets{} targets, each held out in
turn. Non-targeting cells: three replicates. Controls and random removals are
averages over their draws.}
\label{tab:rep}
\centering\scriptsize\setlength{\tabcolsep}{3pt}
% Generated by extension/perturbation_replication/make_assets.py. Do not edit.
\begin{tabular}{lccccc}
\toprule
& \multicolumn{3}{c}{Test targets} & \multicolumn{2}{c}{Non-targeting cells} \\
\cmidrule(lr){2-4}\cmidrule(lr){5-6}
Arm & All & Interferon block & Complement & All & Interferon block \\
\midrule
No pairs: perturbations (primary) & \ensuremath{-}212.4 & \ensuremath{-}3.0 & \ensuremath{-}200.5 & \ensuremath{-}89.5 & 20.7 \\
No pairs: corrected RNA & \ensuremath{-}214.4 & \ensuremath{-}6.3 & \ensuremath{-}194.3 & \ensuremath{-}92.7 & 15.4 \\
No pairs: pseudo-populations & \ensuremath{-}10.2 & \ensuremath{-}3.2 & \ensuremath{-}15.6 & \ensuremath{-}15.5 & \ensuremath{-}4.0 \\
No pairs: derangements & \ensuremath{-}45.6 & \ensuremath{-}12.4 & \ensuremath{-}53.6 & \ensuremath{-}25.8 & \ensuremath{-}8.5 \\
No pairs: without 5 highest-exposure targets & \ensuremath{-}172.3 & \ensuremath{-}5.6 & \ensuremath{-}164.0 & \ensuremath{-}132.1 & 13.9 \\
No pairs: without 5 strongest targets & \ensuremath{-}219.3 & \ensuremath{-}22.5 & \ensuremath{-}173.0 & \ensuremath{-}163.7 & \ensuremath{-}7.6 \\
No pairs: without 5 random targets & \ensuremath{-}216.3 & \ensuremath{-}3.7 & \ensuremath{-}226.5 & \ensuremath{-}109.6 & 10.9 \\
Paired closed form & 32.4 & 83.7 & 25.1 & 82.0 & 89.1 \\
Reference regression & 37.1 & 84.3 & 30.6 & 81.3 & 89.4 \\
Transferred correlation & 81.2 & 94.3 & 83.7 & 89.7 & 85.5 \\
Own paired adaptation cells & 21.8 & 60.7 & 3.1 & 74.2 & 71.0 \\
Adaptation half as replicate & \ensuremath{-}204.6 & 18.9 & \ensuremath{-}390.1 & 60.2 & 88.0 \\
\bottomrule
\end{tabular}

\end{table}

\subsection*{Limitations of this test}
The screen had 25 targets, most of them in one pathway, and four proteins, so
the interferon block had only two antibodies. Groups had 40 to 170 adaptation
cells, which limits the within-assay covariance matrices that closed-form transfer
uses.

\section{Follow-up checks and analyses}\label{sec:review}
This note reports follow-up checks and analyses, all specified after every test had been scored. The analyses were registered together, with their code and data, before any of their quantities was computed. None changes a published estimate. Two runs departed from the planned order of execution without any change to code, setting or value.

\subsection*{Scoring of repeated draws}
At each budget a test draws paired cells several times. For predictions $C_1,\dots,C_D$ from $D$ draws and held-out truth $T$, the numerator of the recovered fraction is $G(C)=2\langle C,T\rangle-\|C\|^2$, and
\[
G\Big(\frac1D\sum_dC_d\Big)=\frac1D\sum_dG(C_d)+\frac1D\sum_d\Big\|C_d-\frac1D\sum_{d'}C_{d'}\Big\|^2 ,
\]
so scoring the averaged prediction would credit a study with the paired cells of all its draws. Each curve in this paper is the mean of per-draw scores: the tests store each arm's mean prediction and the mean over draws of $\|C_d\|^2$, from which $2\langle\bar C,T\rangle-\overline{\|C\|^2}=\frac1D\sum_dG(C_d)$. A check reruns the registered prediction code of each test unchanged, observes every draw's prediction as the code accumulates it, scores each on its own and compares (Supplementary Table~\ref{tab:draws}). Across \RvDrawTests{} tests and data sets, the curves averaged over per-draw scores equalled the published curves to within \RvDrawMaxDiff, the regenerated means and squared norms equalled the stored ones, and the rerun of the bone-marrow test reproduced its recorded predictions exactly. Scoring the averaged predictions instead would have raised the curves by a median of \RvDrawInflMedianMin{} to \RvDrawInflMedianMax{} percentage points per test, and by up to \RvDrawInflMax{} points where draws disagreed most. The development comparison and the validation's readout score each draw inside the draw loop. A second check recomputes every curve from the stored per-draw scores without the data.

\begin{table}[h]
\caption{\textbf{Every recovery curve is the mean of per-draw scores.} Per test or data set: held-out units whose predictions were regenerated (of all units; each benchmark data set was rerun for the first of its three folds), draws per budget, the largest absolute difference between the curve averaged over per-draw scores and the published curve (or, for partial reruns, the fold's stored terms), the largest relative difference between the stored and regenerated mean squared norms, and the median and largest amount, in percentage points of recovered fraction, by which scoring the averaged prediction would have exceeded the published curve. External test: every draw's prediction is stored.}
\label{tab:draws}
\centering\footnotesize
% Generated by extension/synthesis/make_assets.py. Do not edit.
\begin{tabular}{@{}lrrrrrr@{}}
\toprule
 & Units & Draws & \multicolumn{2}{c}{Largest difference} & \multicolumn{2}{c}{Averaged prediction} \\
\cmidrule(lr){4-5}\cmidrule(l){6-7}
Test & rerun & & curve & msq & median & max \\
\midrule
External test & 1/1 & 20 & $3\times10^{-16}$ & -- & 9.0 & 61 \\
Bone-marrow test & 12/12 & 5 & $3\times10^{-9}$ & $4\times10^{-16}$ & 28.9 & 537 \\
Validation test & 34/34 & 5 & $4\times10^{-9}$ & $10^{-16}$ & 9.1 & 360 \\
Benchmark: banc & 1/3 & 5 & $4\times10^{-16}$ & $3\times10^{-16}$ & 8.0 & 32 \\
Benchmark: banc\_crossanimal & 1/3 & 5 & $6\times10^{-16}$ & $4\times10^{-16}$ & 27.5 & 123 \\
Benchmark: bmmc\_cite & 1/3 & 5 & $2\times10^{-15}$ & $3\times10^{-16}$ & 8.0 & 86 \\
Benchmark: bmmc\_multiome & 1/3 & 5 & $2\times10^{-16}$ & $4\times10^{-16}$ & 2.3 & 43 \\
Benchmark: colon & 1/3 & 5 & $10^{-16}$ & $4\times10^{-16}$ & 4.3 & 31 \\
Benchmark: gouwens\_visp & 1/3 & 5 & $10^{-16}$ & $3\times10^{-16}$ & 6.2 & 16 \\
Benchmark: hao & 1/3 & 5 & $6\times10^{-16}$ & $4\times10^{-16}$ & 7.1 & 208 \\
Benchmark: malecns & 1/3 & 5 & $4\times10^{-16}$ & $3\times10^{-16}$ & 9.5 & 33 \\
Benchmark: scala\_m1 & 1/3 & 5 & $10^{-16}$ & $4\times10^{-16}$ & 5.5 & 29 \\
Benchmark: stephenson & 1/3 & 5 & $2\times10^{-16}$ & $4\times10^{-16}$ & 5.1 & 39 \\
Law: banc\_vpn & 2 problems & None & $4\times10^{-16}$ & $3\times10^{-16}$ & 11.3 & 19 \\
Law: bmcite & 3 problems & None & $8\times10^{-16}$ & $4\times10^{-16}$ & 2.7 & 25 \\
Law: fafb\_vpn & 3 problems & None & $6\times10^{-16}$ & $4\times10^{-16}$ & 13.6 & 25 \\
Law: fetal\_cortex & 3 problems & None & $10^{-16}$ & $4\times10^{-16}$ & 4.7 & 51 \\
Law: human\_gaba & 3 problems & None & $2\times10^{-16}$ & $4\times10^{-16}$ & 19.5 & 50 \\
Law: malt10k & 3 problems & None & $10\times10^{-16}$ & $5\times10^{-16}$ & 3.4 & 24 \\
Law: pbmc10k & 3 problems & None & $10^{-15}$ & $4\times10^{-16}$ & 4.1 & 27 \\
Law: sln111 & 3 problems & None & $4\times10^{-16}$ & $6\times10^{-16}$ & 10.3 & 67 \\
Law: sln206 & 3 problems & None & $6\times10^{-16}$ & $6\times10^{-16}$ & 8.5 & 33 \\
Law: snare\_cortex & 3 problems & None & $7\times10^{-16}$ & $4\times10^{-16}$ & 5.0 & 82 \\
\bottomrule
\end{tabular}

\end{table}

\subsection*{The noise estimate after contrasts}
Proposition~\ref{prop:oracle} assumes independent product rows, as when paired cells are centred on means estimated from other cells (development, the external test, the benchmark and the prospective test of the law). The validation test and the reanalyses of the bone-marrow and fly tests centre each population by Helmert contrasts. Contrast rows are uncorrelated, but their products are dependent unless the cells are Gaussian, so the noise estimate $t_b$ of Proposition~\ref{prop:oracle} need not be unbiased for them. Its exact expectation is as follows.

Consider one block and populations $p=1,\dots,P$ with $n_p\ge2$ independent cells each, the scaling and the bases held fixed. Population $p$ contributes $m_p=n_p-1$ rows $h_{pj}=\sum_iH^{(p)}_{ji}u_{pi}$, where $H^{(p)}$ has orthonormal rows orthogonal to the vector of ones, and $z_{pj}$ is the block of $h^x_{pj}h^{y\top}_{pj}$. With $N=\sum_pm_p$ rows, $m_b$ is their mean and $t_b=\sum\|z_{pj}-m_b\|^2/\{N(N-1)\}$. For a coefficient $(k,l)$ of the block, let $\sigma^{(p)}_{kl}$ be the population's covariance of $\tilde x_k$ and $\tilde y_l$ and $\kappa^{(p)}_{kl}=\mathbb{E}[\bar x_k^2\bar y_l^2]-\operatorname{Var}\tilde x_k\operatorname{Var}\tilde y_l-2(\sigma^{(p)}_{kl})^2$ their fourth cumulant ($\bar x,\bar y$ centred at the population means). Write $\mu_p$ for the block of $\sigma^{(p)}$, $\bar\mu=\sum_pm_p\mu_p/N$, $\gamma_p=\sum_{(k,l)}\{\operatorname{Var}\tilde x_k\operatorname{Var}\tilde y_l+(\sigma^{(p)}_{kl})^2\}$, $\kappa_p=\sum_{(k,l)}\kappa^{(p)}_{kl}$ and $F_p=\sum_{j,i}(H^{(p)}_{ji})^4$.

\begin{proposition}[The noise estimate after orthonormal contrasts]\label{prop:contrast}
\[
\operatorname{tr}\operatorname{Cov}(m_b)=\frac{1}{N^2}\sum_p\Big(m_p\gamma_p+\frac{m_p^2}{n_p}\kappa_p\Big),\qquad
\mathbb{E}t_b-\operatorname{tr}\operatorname{Cov}(m_b)=\frac{\sum_pm_p\|\mu_p-\bar\mu\|^2-\sum_p\kappa_p\big(m_p^2/n_p-F_p\big)}{N(N-1)} .
\]
For Helmert contrasts $F_p=m_p+3-2\mathcal H_{m_p}-3/n_p$, with $\mathcal H_m$ the $m$-th harmonic number, so that $m_p^2/n_p-F_p=2\mathcal H_{m_p}-4m_p/n_p$, which lies between 0 and $2\ln m_p$.
\end{proposition}
\noindent\textit{Proof.} Contrasts remove the population mean, so take the cells of population $p$ centred, with coordinates $a=\tilde x_k$ and $b=\tilde y_l$ for one coefficient. For independent centred cells, $\mathbb{E}[a_ib_{i'}a_{r}b_{r'}]=\sigma_{ab}^2\delta_{ii'}\delta_{rr'}+\sigma_{aa}\sigma_{bb}\delta_{ir}\delta_{i'r'}+\sigma_{ab}^2\delta_{ir'}\delta_{i'r}+\kappa_{ab}\delta_{ii'rr'}$, where the last term is one when all four indices coincide. Summing against $H_{ji}H_{ji'}H_{kr}H_{kr'}$ and using $\sum_iH_{ji}H_{ki}=\delta_{jk}$ gives $\mathbb{E}z_j=\sigma_{ab}$ and $\operatorname{Cov}(z_j,z_k)=\delta_{jk}(\sigma_{aa}\sigma_{bb}+\sigma_{ab}^2)+\kappa_{ab}\sum_iH_{ji}^2H_{ki}^2$. Rows of different populations are independent. Summing over the coefficients of the block, over $j$ and over $k$, with $\sum_jH_{ji}^2=1-1/n_p$, gives the variance of $m_b$. For rows with means $\mu_j$, $\mathbb{E}\sum_j\|z_j-m_b\|^2=\sum_j\operatorname{tr}\operatorname{Cov}z_j+\sum_j\|\mu_j-\bar\mu\|^2-N\operatorname{tr}\operatorname{Cov}(m_b)$, and dividing by $N(N-1)$ and subtracting leaves the covariances between distinct rows, $\sum_{j\ne k}\sum_iH_{ji}^2H_{ki}^2=m_p^2/n_p-F_p$ within population $p$. For the Helmert row $j$, with $j$ entries $1/\sqrt{j(j+1)}$ and one entry $-j/\sqrt{j(j+1)}$, $\sum_iH_{ji}^4=1+1/j-3/(j+1)$, and summing over $j$ gives $F_p$. Finally $2\mathcal H_m-4m/(m+1)$ is 0 at $m=1$, increases with $m$ because $2/(m+1)\ge4/\{(m+1)(m+2)\}$, and is at most $2\ln m$ because $\mathcal H_m\le1+\ln m$. $\square$

The first term of the bias is non-negative: populations whose cross-covariances differ inflate the estimate, and shrinkage becomes more cautious. The second vanishes for Gaussian cells ($\kappa_p=0$) and otherwise has the sign opposite to the fourth cumulant. Relative to the fourth-cumulant part of $\operatorname{tr}\operatorname{Cov}(m_b)$, it is of order $(\ln m_p)/m_p$ for Helmert contrasts, whose rows are dominated by a single new cell; for a basis with entries of equal size ($F_p=m_p/n_p$) it would miss that part entirely. The referee's example, a population of three cells with $x\in\{-3,0,3\}$ of probabilities $1/18,8/9,1/18$ and $y=x\pm1$, has $\gamma=3$, $\kappa=6$, $\operatorname{Var}(m)=3/2+6/3=3.5$ and $\mathbb{E}t=3.5-6\cdot(1/3)/2=2.5$. The proposition holds the scaling fixed, whereas the implemented pipeline estimates the standard deviations from the same contrasts; the calibration below includes that step.

\paragraph{Calibration.} For every fold of the validation test (atlas axes) and of the bone-marrow test (other-site axes, the reanalysis with contrasts) and each budget, 500 independent resamples of paired cells were drawn with replacement from the fold's reservoir and passed through the complete pipeline: Helmert contrasts in draw order, scaling by the contrasts' standard deviations, rotation and blocks. The calibration ratio is the mean of a noise estimate over resamples divided by the variance of the block mean across resamples, pooled over folds, with 95\% intervals from resampling donors (validation) or held-out batches (bone marrow) and then resamples within folds (Supplementary Tables~\ref{tab:noisecal} and~\ref{tab:noiseblocks}). The implemented estimate was slightly conservative at every budget: its ratio, summed over blocks, was \RvCalValMin{} to \RvCalValMax{} in the validation test and \RvCalBmMin{} to \RvCalBmMax{} in the bone-marrow test, and \RvCalValBlockMin{} to \RvCalValBlockMax{} and \RvCalBmBlockMin{} to \RvCalBmBlockMax{} block by block, highest in the blocks of the leading RNA axes, where the populations' cross-covariances differ most, as the first term of Proposition~\ref{prop:contrast} predicts; the excess grew with the budget. Random orthonormal contrasts gave \RvCalValOrthoMin{} to \RvCalValOrthoMax. Reordering the cells within populations left the ratios unchanged, as it must, since the order changes the estimate's variance but not its expectation; the variance was appreciable at small budgets, where the most affected block's estimate varied over orders with a median coefficient of variation of \RvCalValOrderCvMax\% at 25 paired cells, \RvCalValOrderCvHundred\% at 100 and \RvCalValOrderCvEight\% at 800, and its shrinkage factor by a median range of \RvCalValOrderFactorMax, \RvCalValOrderFactorHundred{} and \RvCalValOrderFactorEight. The jackknife over cells, a dependence-aware alternative, overestimated the noise at small budgets (ratio \RvCalValJackMax{} at 25 paired cells, falling to \RvCalValJackMin{} at 800), and centring on the drawn cells' own averages underestimated it (\RvCalValOwnMin{} at 25), reproducing the failure of the bone-marrow test's original centring. The jackknife's larger estimates shrink more, and on the same resamples they raised the recovered fraction by \RvCalValJackGainMin{} to \RvCalValJackGainMax{} percentage points in the validation test and by \RvCalBmJackGainMin{} to \RvCalBmJackGainMax{} in the bone-marrow test, most at small budgets (Supplementary Table~\ref{tab:noisecal}): at these budgets the positive-part rule with an unbiased noise estimate shrinks less than would be best. The implemented estimate, fixed before the tests, is kept in every result reported here; its calibration after contrasts is empirical, and Proposition~\ref{prop:contrast} gives the form and order of its bias.

\begin{table}[h]
\caption{\textbf{Calibration of the noise estimate after contrasts.} Mean noise estimate over 500 independent resamples of paired cells, divided by the variance of the block mean across the same resamples, summed over the eight blocks and pooled over folds (unrounded to three decimals). For the implemented estimate (Helmert contrasts in draw order): its 95\% interval from resampling donors (validation) or held-out batches (bone marrow) and resamples within folds; the range of its ratio over the eight blocks; and, for the block most affected by the order of cells, the median coefficient of variation of the estimate and the median range of the block's shrinkage factor over 20 random orders of the cells within populations. Alternatives on the same resamples: Helmert contrasts in a random order, random orthonormal contrasts, the jackknife over cells, and centring on the drawn cells' own averages ($\sqrt{n/(n-1)}$ times the deviations, as in the bone-marrow test, compared with the variance of its own block mean). Last rows: fraction of the dependence recovered by the atlas estimate with condition maps (validation) or the other-site estimate (bone marrow) on the same resamples, with the implemented and with the jackknife noise.}
\label{tab:noisecal}
\centering\scriptsize
% Generated by extension/synthesis/make_assets.py. Do not edit.
\begin{tabular}{@{}lrrrrrr@{}}
\toprule
 & \multicolumn{6}{c}{Paired cells} \\
\cmidrule(l){2-7}
Noise estimate & 25 & 50 & 100 & 200 & 400 & 800 \\
\midrule
\multicolumn{7}{@{}l}{\textit{Validation test}} \\
Helmert, draw order (implemented) & 1.005 & 1.008 & 1.011 & 1.014 & 1.016 & 1.016 \\
\quad 95\% interval & 1.003--1.006 & 1.006--1.010 & 1.009--1.012 & 1.012--1.015 & 1.014--1.017 & 1.015--1.018 \\
\quad range over blocks & 1.00--1.05 & 1.00--1.05 & 1.00--1.06 & 1.01--1.06 & 1.01--1.07 & 1.01--1.07 \\
\quad order, CV (largest) & 12.3\% & 10.8\% & 8.5\% & 6.2\% & 4.2\% & 2.8\% \\
\quad order, factor range & 0.223 & 0.184 & 0.103 & 0.057 & 0.032 & 0.016 \\
Helmert, random order & 1.005 & 1.008 & 1.011 & 1.014 & 1.015 & 1.016 \\
Random orthonormal & 1.003 & 1.005 & 1.007 & 1.008 & 1.010 & 1.010 \\
Jackknife over cells & 1.478 & 1.258 & 1.134 & 1.073 & 1.045 & 1.031 \\
Own-mean centring & 0.674 & 0.798 & 0.898 & 0.958 & 0.988 & 1.003 \\
Recovered (\%), implemented & \ensuremath{-}1.3 & 8.8 & 19 & 28 & 33 & 35 \\
Recovered (\%), jackknife & 1.3 & 10 & 22 & 31 & 34 & 35 \\
\addlinespace
\multicolumn{7}{@{}l}{\textit{Bone-marrow test}} \\
Helmert, draw order (implemented) & 1.005 & 1.012 & 1.016 & 1.024 & 1.029 & 1.029 \\
\quad 95\% interval & 1.003--1.008 & 1.008--1.016 & 1.012--1.021 & 1.018--1.030 & 1.021--1.036 & 1.022--1.036 \\
\quad range over blocks & 1.00--1.08 & 1.00--1.16 & 1.00--1.18 & 1.01--1.26 & 1.01--1.28 & 1.01--1.31 \\
\quad order, CV (largest) & 12.4\% & 11.6\% & 9.7\% & 7.8\% & 6.0\% & 4.4\% \\
\quad order, factor range & 0.218 & 0.148 & 0.113 & 0.068 & 0.032 & 0.015 \\
Helmert, random order & 1.005 & 1.012 & 1.016 & 1.024 & 1.028 & 1.029 \\
Random orthonormal & 1.005 & 1.010 & 1.015 & 1.023 & 1.030 & 1.034 \\
Jackknife over cells & 1.538 & 1.328 & 1.193 & 1.117 & 1.075 & 1.051 \\
Own-mean centring & 0.650 & 0.759 & 0.855 & 0.934 & 0.982 & 1.006 \\
Recovered (\%), implemented & \ensuremath{-}4.2 & 3.8 & 10 & 15 & 18 & 17 \\
Recovered (\%), jackknife & 1.7 & 7.8 & 15 & 20 & 20 & 18 \\
\bottomrule
\end{tabular}

\end{table}

\begin{table}[h]
\caption{\textbf{Calibration of the implemented noise estimate block by block.} As Supplementary Table~\ref{tab:noisecal}, for Helmert contrasts in draw order, in each of the eight blocks (RNA axes 1--5, 6--20, 21--60 and the rest, crossed with protein axes 1--3 and the rest), with 95\% intervals.}
\label{tab:noiseblocks}
\centering\scriptsize
% Generated by extension/synthesis/make_assets.py. Do not edit.
\begin{tabular}{@{}lrrrrrr@{}}
\toprule
 & \multicolumn{6}{c}{Paired cells} \\
\cmidrule(l){2-7}
Block (RNA/protein axes) & 25 & 50 & 100 & 200 & 400 & 800 \\
\midrule
\multicolumn{7}{@{}l}{\textit{Validation test}} \\
1--5/1--3 & 1.048 & 1.048 & 1.050 & 1.052 & 1.061 & 1.054 \\
 & (1.032--1.062) & (1.030--1.065) & (1.030--1.070) & (1.036--1.066) & (1.047--1.076) & (1.040--1.068) \\
1--5/rest & 1.038 & 1.051 & 1.059 & 1.063 & 1.066 & 1.069 \\
 & (1.030--1.045) & (1.043--1.059) & (1.052--1.065) & (1.056--1.069) & (1.061--1.071) & (1.064--1.074) \\
6--20/1--3 & 1.011 & 1.009 & 1.015 & 1.022 & 1.025 & 1.030 \\
 & (1.004--1.018) & (1.003--1.015) & (1.009--1.022) & (1.016--1.029) & (1.019--1.031) & (1.024--1.037) \\
6--20/rest & 1.003 & 1.007 & 1.011 & 1.016 & 1.019 & 1.019 \\
 & (1.001--1.005) & (1.005--1.009) & (1.009--1.013) & (1.013--1.018) & (1.016--1.021) & (1.018--1.021) \\
21--60/1--3 & 1.005 & 1.005 & 1.010 & 1.013 & 1.015 & 1.018 \\
 & (1.002--1.009) & (1.001--1.008) & (1.006--1.014) & (1.008--1.017) & (1.011--1.019) & (1.014--1.023) \\
21--60/rest & 1.001 & 1.003 & 1.005 & 1.008 & 1.010 & 1.012 \\
 & (1.000--1.002) & (1.002--1.004) & (1.004--1.007) & (1.007--1.009) & (1.009--1.012) & (1.011--1.014) \\
rest/1--3 & 1.003 & 1.005 & 1.008 & 1.010 & 1.012 & 1.012 \\
 & (1.000--1.005) & (1.003--1.008) & (1.005--1.010) & (1.008--1.013) & (1.009--1.015) & (1.009--1.015) \\
rest/rest & 1.000 & 1.002 & 1.004 & 1.006 & 1.007 & 1.008 \\
 & (1.000--1.001) & (1.001--1.003) & (1.003--1.005) & (1.005--1.007) & (1.006--1.009) & (1.007--1.009) \\
\addlinespace
\multicolumn{7}{@{}l}{\textit{Bone-marrow test}} \\
1--5/1--3 & 1.082 & 1.161 & 1.184 & 1.265 & 1.284 & 1.310 \\
 & (1.041--1.125) & (1.116--1.206) & (1.142--1.224) & (1.203--1.321) & (1.214--1.346) & (1.236--1.375) \\
1--5/rest & 1.017 & 1.034 & 1.040 & 1.059 & 1.071 & 1.074 \\
 & (1.012--1.023) & (1.023--1.046) & (1.030--1.051) & (1.045--1.076) & (1.053--1.090) & (1.056--1.092) \\
6--20/1--3 & 1.031 & 1.049 & 1.077 & 1.099 & 1.113 & 1.113 \\
 & (1.008--1.053) & (1.031--1.069) & (1.053--1.102) & (1.070--1.128) & (1.084--1.139) & (1.087--1.136) \\
6--20/rest & 1.002 & 1.009 & 1.015 & 1.023 & 1.032 & 1.033 \\
 & (1.000--1.005) & (1.006--1.013) & (1.010--1.021) & (1.018--1.029) & (1.023--1.041) & (1.026--1.040) \\
21--60/1--3 & 1.018 & 1.028 & 1.034 & 1.038 & 1.046 & 1.042 \\
 & (1.008--1.028) & (1.019--1.036) & (1.023--1.045) & (1.027--1.048) & (1.034--1.057) & (1.030--1.054) \\
21--60/rest & 1.000 & 1.003 & 1.005 & 1.010 & 1.013 & 1.012 \\
 & (0.998--1.001) & (1.001--1.005) & (1.003--1.007) & (1.007--1.012) & (1.010--1.016) & (1.010--1.015) \\
rest/1--3 & 1.015 & 1.020 & 1.026 & 1.028 & 1.031 & 1.028 \\
 & (1.009--1.022) & (1.013--1.028) & (1.020--1.032) & (1.023--1.034) & (1.024--1.038) & (1.022--1.035) \\
rest/rest & 1.001 & 1.003 & 1.005 & 1.008 & 1.010 & 1.009 \\
 & (0.999--1.002) & (1.001--1.004) & (1.003--1.006) & (1.006--1.010) & (1.008--1.012) & (1.007--1.011) \\
\bottomrule
\end{tabular}

\end{table}

\subsection*{Axes and condition maps}
The estimator takes two things from a reference: axes, along which the paired estimate is shrunk block by block, and each cell type's covariance, through which the condition map turns the estimate pooled over types into one for each type. In the validation test, every draw of paired cells was regenerated with the test's registered code (the regenerated means and squared norms equalled the stored ones to within \RvValMeanDiff{} and \RvValMsqDiff) and scored with four versions for the atlas: axes and maps (the test's arm); axes only, the same pooled estimate given to every cell type; maps only, paired-only single-block James--Stein pooled over types and passed through the atlas's maps; and neither, paired-only single-block James--Stein pooled over types. Axes only was also scored for the other-pool and own-sample references.

Each alone did little (Supplementary Table~\ref{tab:axesmaps}). With 100 paired cells, the axes recovered \RvAxesHundred\% of the dependence within cell types (95\% CI, \RvAxesHundredLo--\RvAxesHundredHi\%) and the maps \RvMapOnlyHundred\% (\RvMapOnlyHundredLo--\RvMapOnlyHundredHi\%), against \RvNeitherHundred\% with neither and \RvBothHundred\% (\RvBothHundredLo--\RvBothHundredHi\%) with both. With 800, the axes recovered \RvAxesEight\%, the maps \RvMapOnlyEight\% and both \RvBothEight\%. Neither alone reached the primary target within 800 paired cells, whereas both together reached it with \RvBothCells{} (95\% CI, \RvBothCellsLo--\RvBothCellsHi; the regenerated draws add the readout's budget of 150 paired cells, which moves the interpolation slightly from the prespecified test's \VlCellsAtlas). The same held for the study's own unpaired cells and its other pools: without maps they recovered \RvOwnAxesEight\% and \RvOtherAxesEight\% with 800 paired cells. An estimate pooled over cell types can recover only the part of the within-type dependence that the types share; the maps make it type-specific, and the axes make the pooled estimate accurate from few cells.

\paragraph{Reweighting the reference.} A condition map passes the pooled paired estimate through the reference's pooled covariance, which weights cell types by the reference's cells, whereas the pooled contrast estimate weights them by their contrast rows (the sum over a type's populations of $n_p-1$). The reweighted map moves the reference's pooled latent covariance $S$ to $S+\sum_c(w_c-w^0_c)S_c$, with $S_c$ the reference's covariance of type $c$ in pooled standard deviations, $w^0_c$ the reference's share of type $c$ and $w_c$ the draw's share of contrast rows, rescales it to the reweighted pooled standard deviations and rebuilds the maps exactly as the unweighted ones are built; with $w=w^0$ it gives the unweighted maps, which the code checks (largest difference 0). It was applied to the full atlas and to the two skewed atlases of Supplementary Table~\ref{tab:robust}, in every draw. Reweighting rescued both skewed atlases: with B cells and CD14 monocytes from 15 patients, the estimator needed \RvFewBMonoRwCells{} paired cells (95\% CI, \RvFewBMonoRwCellsLo--\RvFewBMonoRwCellsHi) with reweighted maps against \RvFewBMonoMappedCells{} with the unweighted ones, and with CD4 and CD8 T cells from 15 patients \RvFewTRwCells{} (\RvFewTRwCellsLo--\RvFewTRwCellsHi) against \RvFewTMappedCells; the full atlas needed \RvFullRwCells{} with reweighted maps and \RvBothCells{} with its own. With 25 paired cells, when some cell types had almost no contrast rows, the reweighted maps were far worse than predicting no dependence for every atlas (\RvFullRwTwentyFive\% for the full atlas).

\begin{table}[h]
\caption{\textbf{Axes and condition maps in the validation test.} Fraction of the within-type dependence recovered (\%), pooled over the \VlNFolds{} held-out samples, and paired cells to the test's primary target (\VlTarget\%) with 95\% intervals from resampling donors. Top: the atlas with its axes and maps (the test's arm), its axes only, its maps only (paired-only single-block James--Stein pooled over types, through the atlas's maps) and neither; the other pools' and the own samples' unpaired cells with and without maps; the best paired-only method at each budget. Bottom: the unchanged atlas and the two atlases with a skewed mix of cell types (Supplementary Table~\ref{tab:robust}), with their unweighted maps, with maps reweighted to each draw's contrast rows per cell type, and without maps.}
\label{tab:axesmaps}
\centering\footnotesize
% Generated by extension/synthesis/make_assets.py. Do not edit.
\begin{tabular}{@{}lrrrrrrrr@{}}
\toprule
 & \multicolumn{7}{c}{Recovered (\%) with paired cells} & Paired cells \\
\cmidrule(lr){2-8}
Reference and version & 25 & 50 & 100 & 150 & 200 & 400 & 800 & to the target (95\% CI) \\
\midrule
\multicolumn{9}{@{}l}{Atlas} \\
\quad Axes and maps & 0.5 & 11 & 23 & 29 & 33 & 42 & 48 & 145 (124--183) \\
\quad Axes only & \ensuremath{-}0.7 & 4.0 & 7.8 & 9.3 & 10 & 13 & 14 & $>$800 \\
\quad Maps only & \ensuremath{-}0.2 & 1.8 & 4.7 & 7.1 & 9.5 & 18 & 28 & $>$800 \\
\quad Neither & 0.0 & 1.1 & 2.4 & 3.4 & 4.4 & 7.1 & 9.9 & $>$800 \\
\addlinespace
\multicolumn{9}{@{}l}{Other pools and donors} \\
\quad Axes and maps & 8.8 & 23 & 34 & 39 & 43 & 50 & 57 & 70 (58--85) \\
\quad Axes only & 1.6 & 7.5 & 10 & 11 & 12 & 13 & 14 & $>$800 \\
\addlinespace
\multicolumn{9}{@{}l}{Own samples} \\
\quad Axes and maps & 5.6 & 19 & 27 & 31 & 34 & 40 & 45 & 109 (82--165) \\
\quad Axes only & 0.5 & 6.5 & 9.2 & 9.7 & 10 & 12 & 13 & $>$800 \\
\addlinespace
\multicolumn{9}{@{}l}{Paired cells only} \\
\quad Best of eight methods & 0.5 & 2.4 & 3.5 & 5.1 & 7.7 & 14 & 13 & $>$800 \\
\midrule
\multicolumn{9}{@{}l}{\textit{Reweighting the atlas}} \\
\multicolumn{9}{@{}l}{All 118 atlas patients} \\
\quad Unweighted maps & 0.5 & 11 & 23 & 29 & 33 & 42 & 48 & 145 (124--183) \\
\quad Reweighted maps & \ensuremath{-}37 & 8.3 & 24 & 31 & 35 & 44 & 50 & 131 (114--162) \\
\quad No maps & \ensuremath{-}0.7 & 4.0 & 7.8 & 9.3 & 10 & 13 & 14 & $>$800 \\
\addlinespace
\multicolumn{9}{@{}l}{B cells and CD14 monocytes from 15} \\
\quad Unweighted maps & \ensuremath{-}13 & \ensuremath{-}6.1 & 2.8 & 5.7 & 10.0 & 17 & 20 & $>$800 \\
\quad Reweighted maps & \ensuremath{-}25 & 10.0 & 23 & 29 & 32 & 40 & 46 & 148 (122--192) \\
\quad No maps & \ensuremath{-}0.5 & 4.2 & 8.0 & 9.3 & 10 & 12 & 14 & $>$800 \\
\addlinespace
\multicolumn{9}{@{}l}{CD4 and CD8 T cells from 15} \\
\quad Unweighted maps & \ensuremath{-}3.3 & 7.2 & 14 & 15 & 14 & 8.8 & \ensuremath{-}2.7 & $>$800 \\
\quad Reweighted maps & \ensuremath{-}35 & 0.6 & 12 & 18 & 22 & 32 & 41 & 311 (266--378) \\
\quad No maps & \ensuremath{-}1.3 & 2.5 & 5.9 & 7.6 & 8.9 & 12 & 13 & $>$800 \\
\bottomrule
\end{tabular}

\end{table}

\subsection*{One study, one draw, one atlas}
The validation's intervals resample held-out donors with predictions averaged over five draws of paired cells and one atlas. A study draws its paired cells once. With each draw index scored on its own, and in 2,000 resamples that drew donors and then one draw index per fold (shared by all methods), the saving over paired-only estimation (V1) and the ratio of the paired cells that reference regression and the estimator needed (K2) were recomputed. In ten bootstrap replicates of the atlas's patients, the atlas context was rebuilt and the estimator and the four non-pooled forms of reference regression refitted on the first draw of every fold; a baseline replicate took every patient once through the same construction. Scored one draw at a time, the estimator needed \RvDrawCellsMin{} to \RvDrawCellsMax{} paired cells, the saving over paired-only estimation was at least \RvDrawVOneMin{} to \RvDrawVOneMax, and K2 was \RvDrawKTwoMin{} to \RvDrawKTwoMax{} (resampled, \RvDrawKTwoMed; 95\% interval \RvDrawKTwoLo--\RvDrawKTwoHi, above one in \RvDrawKTwoShare\% of resamples). With the atlas's patients resampled, the estimator needed \RvAtlasCellsMin{} to \RvAtlasCellsMax{} paired cells and K2 was \RvAtlasKTwoMin{} to \RvAtlasKTwoMax{} over the \RvAtlasReps{} replicates (with donors also resampled, \RvAtlasKTwoLo--\RvAtlasKTwoHi; Supplementary Table~\ref{tab:onedraw}). The estimator's small advantage over reference regression given the atlas therefore depends on the draw and on the atlas's patients; its advantage over paired-only estimation does not.

\begin{table}[h]
\caption{\textbf{One study, one draw, one atlas.} Top: each draw index of the validation test scored on its own (the atlas fixed): paired cells the estimator needed to reach the primary target, the saving over the paired-only envelope (V1; a lower bound when paired-only estimation did not reach the target within 800) and the ratio of the paired cells needed by the best form of reference regression to those needed by the estimator (K2); then 2,000 resamples of donors with one draw index per fold, median and 95\% interval. Bottom: the atlas's patients resampled (ten bootstrap replicates; first draw of paired cells), with the replicate that takes every patient once; then the replicates combined with 200 resamples of donors each. The central 95\% of these nested resamples is a sensitivity summary: ten atlas replicates do not establish nominal coverage over new reference atlases.}
\label{tab:onedraw}
\centering\footnotesize
% Generated by extension/synthesis/make_assets.py. Do not edit.
\begin{tabular}{@{}lrrr@{}}
\toprule
 & Paired cells, estimator & Saving over paired-only (V1) & Regression over estimator (K2) \\
\midrule
\multicolumn{4}{@{}l}{\textit{Each draw of paired cells on its own (atlas fixed)}} \\
\quad Draw 1 & 152 & 5.27 & 1.09 \\
\quad Draw 2 & 142 & 5.65 & 1.17 \\
\quad Draw 3 & 154 & 5.19 & 0.95 \\
\quad Draw 4 & 133 & 6.03 & 1.04 \\
\quad Draw 5 & 151 & 5.30 & 1.08 \\
\quad One draw per fold, donors resampled & & 5.53 (4.21--6.68) & 1.04 (0.85--1.45) \\
\addlinespace
\multicolumn{4}{@{}l}{\textit{Atlas patients resampled (first draw)}} \\
\quad Every patient once & 147 & 5.43 & 1.13 \\
\quad Replicate 1 & 166 & 4.83 & 0.99 \\
\quad Replicate 2 & 160 & 5.00 & 1.05 \\
\quad Replicate 3 & 146 & 5.49 & 1.10 \\
\quad Replicate 4 & 170 & 4.71 & 1.00 \\
\quad Replicate 5 & 151 & 5.30 & 1.27 \\
\quad Replicate 6 & 159 & 5.04 & 1.08 \\
\quad Replicate 7 & 141 & 5.68 & 1.18 \\
\quad Replicate 8 & 158 & 5.05 & 1.06 \\
\quad Replicate 9 & 160 & 5.00 & 1.04 \\
\quad Replicate 10 & 152 & 5.28 & 1.13 \\
\quad Replicates and donors resampled & & 5.04 (3.85--6.34) & 1.07 (0.82--1.40) \\
\bottomrule
\end{tabular}

\end{table}

\subsection*{Named within-type relationships}
Eight relationships between a gene program and a protein were fixed from prior knowledge of blood cell states before any was scored (Supplementary Table~\ref{tab:relationships}). A relationship's value in a held-out sample and cell type is the mean, over the program's genes, of their within-type correlation with the protein; a prediction's value is the same mean of the predicted matrix. Endpoints were the recovered fraction of each relationship over held-out units, the share of units in which the prediction had the sign of the held-out value (among units whose two halves agreed in sign) and, per donor, the difference in absolute error between the estimator and the best paired-only method or the best form of reference regression, both chosen per relationship and budget on the pooled error, which favours the comparators. With 100 paired cells, the estimator recovered more of the relationship than reference regression in \RvRelBetterReg{} of the \RvRelN{} relationships and more than paired-only estimation in \RvRelBetterPo; its absolute error was lower than regression's in \RvRelDonorWinReg\% of the comparisons between donors and lower than paired-only estimation's in \RvRelDonorWinPo\%. Every relationship had the expected sign in the held-out cells, two of them weakly: interferon-stimulated genes with CD169 in CD14 monocytes ($r=\RvRelOneTruth$) and cytotoxic genes with CD56 in NK cells ($r=\RvRelSevenTruth$, negative in \RvRelSevenTruthShare\% of units). The second was too weak for 100 paired cells: the estimator's predictions scattered around zero (mean \RvRelSevenAtlasPred) and recovered less than predicting no association (\RvRelSevenAtlasRf\%), and reference regression's did so too, by less (\RvRelSevenRegRf\%).

\begin{table}[h]
\caption{\textbf{Named within-type relationships with 100 paired cells.} Programs: interferon-stimulated (R1: ISG15, IFI6, IFI44L, MX1, OAS1, OAS2, HERC5, XAF1, EPSTI1, IFITM3, STAT2, IRF9), MHC class II (R2, R3: HLA-DRA, HLA-DRB1, HLA-DRB5, HLA-DQA1, HLA-DQB1, HLA-DPA1, HLA-DPB1, HLA-DMA, HLA-DMB), cytotoxic in CD8 T cells (R4, R5: GZMB, GZMH, PRF1, GNLY, NKG7, FGFBP2, CST7, KLRD1, ADGRG1) and in NK cells (R6, R7: GZMB, PRF1, FGFBP2, SPON2, GNLY, NKG7), and LEF1 (R8); expected signs positive except R5 and R7. Units, held-out samples in which the cell type was scored; held-out $r$, the mean over units of the relationship's value; recovered, the fraction of the relationship recovered by the estimator with the atlas (Est.), reference regression with the atlas (Regr.) and paired-only estimation (Paired), each comparator in its best form for the relationship; correct sign, among units whose two halves agreed in sign; donors, the number of donors in which the estimator's absolute error was lower, of all donors with the cell type scored.}
\label{tab:relationships}
\centering\scriptsize
% Generated by extension/synthesis/make_assets.py. Do not edit.
\begin{tabular}{@{}lrrrrrrrrrr@{}}
\toprule
 & & Held-out & \multicolumn{3}{c}{Recovered (\%)} & \multicolumn{3}{c}{Correct sign (\%)} & \multicolumn{2}{c}{Donors, estimator better} \\
\cmidrule(lr){4-6}\cmidrule(lr){7-9}\cmidrule(l){10-11}
Relationship & Units & $r$ & Est. & Regr. & Paired & Est. & Regr. & Paired & than regr. & than paired \\
\midrule
R1\enspace CD14 Mono: interferon, CD169 & 34 & 0.07 & 27 & 21 & 15 & 72 & 79 & 76 & 14/18 & 15/18 \\
R2\enspace CD14 Mono: MHC II, HLA-DR & 34 & 0.47 & 80 & 46 & 19 & 100 & 100 & 93 & 18/18 & 18/18 \\
R3\enspace B: MHC II, HLA-DR & 34 & 0.15 & 81 & 54 & 39 & 100 & 100 & 100 & 17/18 & 17/18 \\
R4\enspace CD8 T: cytotoxic, CD57 & 34 & 0.50 & 65 & 63 & 57 & 99 & 93 & 94 & 12/18 & 12/18 \\
R5\enspace CD8 T: cytotoxic, CD27 & 34 & \ensuremath{-}0.49 & 69 & 66 & 56 & 99 & 100 & 99 & 10/18 & 16/18 \\
R6\enspace NK: cytotoxic, CD16 & 34 & 0.38 & 53 & 37 & 18 & 98 & 95 & 81 & 18/18 & 18/18 \\
R7\enspace NK: cytotoxic, CD56 & 34 & \ensuremath{-}0.05 & \ensuremath{-}81 & \ensuremath{-}20 & 0.0 & 36 & 46 & 0 & 4/18 & 3/18 \\
R8\enspace CD4 T: LEF1, CD45RA & 34 & 0.33 & 55 & 38 & 29 & 100 & 100 & 96 & 18/18 & 18/18 \\
\bottomrule
\end{tabular}

\end{table}

\subsection*{The association test and the truth it defines}
The readout defines a reproducible association by Fisher's $z$ with $P<0.01$ and the same sign in both halves of the held-out cells, down to 20 cells per half. Permuting protein vectors across the cells of each half, which keeps both marginal distributions and removes the pairing (20 permutations per unit and half), gives the test's null distribution under the actual data. The test was calibrated at the thresholds the readout uses: over all \RvFisherUnits{} units, \RvFisherFive\% of null tests gave $P<0.05$ and \RvFisherOne\% gave $P<0.01$, and in every range of cells per half \RvFisherOneMin\% to \RvFisherOneMax\% gave $P<0.01$ (Supplementary Table~\ref{tab:fisher}); the far tail was somewhat heavier (\RvFisherTenth\% at $P<0.001$). No scored unit had fewer than \RvFisherSmallLo{} cells per half, and in the smallest (\RvFisherSmallRange{} cells) \RvFisherSmallOne\% of null tests gave $P<0.01$. Independently permuted halves passed the readout's rule, $P<0.01$ in both with the same sign, for \RvReplNull{} of every 100,000 pairs, against a nominal \RvReplNominal{} and an observed \RvReplObserved\% of pairs, so false associations make up about \RvReplFalseShare\% of those counted.

The readout was also recomputed with the held-out truth centred within the authors' finer cell states (for example naive and memory B cells) instead of the five broad types, with RNA and protein residualized on the log library sizes of both modalities and the 10x lane, and with both; reproducible associations were redefined under each truth, with Fisher's $z$ scaled by $\sqrt{n-3-k}$ for $k$ covariates (Supplementary Table~\ref{tab:truths}). Removing the technical covariates changed little: there were \RvTechReps{} reproducible associations instead of \RvPrimaryReps, and with 100 paired cells the atlas's gain over paired-only estimation was \RvTechOneMinusPo{} points in direction (95\% CI, \RvTechOneMinusPoLo--\RvTechOneMinusPoHi) and \RvTechTwoMinusPo{} in the top pairs. Within finer states, most associations disappeared (\RvFineReps), and the gain fell to \RvFineOneMinusPo{} points in direction (\RvFineOneMinusPoLo--\RvFineOneMinusPoHi) and \RvFineTwoMinusPo{} in the top pairs; the recovered fraction became negative (\RvFineRfAtlas\%, against \RvFineRfPo\% for paired-only estimation and \RvFineRfReg\% for reference regression), because the estimate, which targets the dependence within a broad type, includes the part that lies between the finer states. Much of the within-type dependence that the readout scores thus lies between cell states within the broad types; estimating it within finer states would require centring the paired cells within them.

\begin{table}[h]
\caption{\textbf{The association test under permutation.} Share of Fisher-$z$ tests at or below each $P$ value when protein vectors were permuted across the cells of each half of a held-out unit (20 permutations), by the number of cells per half; tests are gene--protein pairs.}
\label{tab:fisher}
\centering\footnotesize
% Generated by extension/synthesis/make_assets.py. Do not edit.
\begin{tabular}{@{}lrrrr@{}}
\toprule
 & & \multicolumn{3}{c}{Share of null tests at or below (\%)} \\
\cmidrule(l){3-5}
Cells per half & Tests & 5\% & 1\% & 0.1\% \\
\midrule
40--79 & 888,000 & 4.54 & 1.07 & 0.18 \\
80--159 & 17,760,000 & 4.71 & 1.10 & 0.18 \\
160 or more & 132,312,000 & 4.92 & 1.06 & 0.14 \\
All & 150,960,000 & 4.90 & 1.07 & 0.15 \\
\bottomrule
\end{tabular}

\end{table}

\begin{table}[h]
\caption{\textbf{The readout under other held-out truths.} The atlas estimate minus paired-only estimation (best of eight methods at each budget) and minus reference regression (best of six forms), with 100 paired cells and 95\% intervals from resampling donors, in the share of reproducible associations whose direction was predicted (E1), the share of reproducible associations of the predicted direction among the 100 largest predictions (E2) and the recovered fraction. Truths: centred within the five broad types (the published readout), within the authors' finer cell states, within broad types after residualizing on log library sizes and the 10x lane, and both.}
\label{tab:truths}
\centering\scriptsize
% Generated by extension/synthesis/make_assets.py. Do not edit.
\begin{tabular}{@{}lrrlrrr@{}}
\toprule
 & & & & \multicolumn{3}{c}{Estimator minus comparator (points)} \\
\cmidrule(l){5-7}
Held-out truth & Units & Associations & Comparator & Direction (E1) & Top 100 (E2) & Recovered \\
\midrule
Broad types (published) & 170 & 120,649 & Paired only & 11 (10 to 12) & 19 (18 to 20) & 19 (17 to 20) \\
 &  &  & Regression & \ensuremath{-}0.1 (\ensuremath{-}0.5 to 0.2) & \ensuremath{-}2.8 (\ensuremath{-}3.6 to \ensuremath{-}1.9) & 0.1 (\ensuremath{-}1.3 to 1.6) \\
\addlinespace
Finer states & 170 & 31,509 & Paired only & 2.8 (2.1 to 3.4) & 6.5 (6.1 to 7.0) & \ensuremath{-}71 (\ensuremath{-}87 to \ensuremath{-}58) \\
 &  &  & Regression & \ensuremath{-}2.0 (\ensuremath{-}2.7 to \ensuremath{-}1.3) & \ensuremath{-}1.9 (\ensuremath{-}2.2 to \ensuremath{-}1.6) & \ensuremath{-}70 (\ensuremath{-}85 to \ensuremath{-}57) \\
\addlinespace
Technical covariates & 170 & 98,739 & Paired only & 10 (9.6 to 11) & 18 (17 to 19) & 13 (11 to 16) \\
 &  &  & Regression & 0.1 (\ensuremath{-}0.3 to 0.5) & \ensuremath{-}1.7 (\ensuremath{-}2.3 to \ensuremath{-}1.1) & \ensuremath{-}3.9 (\ensuremath{-}5.8 to \ensuremath{-}2.0) \\
\addlinespace
Both & 170 & 26,750 & Paired only & 3.2 (2.6 to 3.6) & 5.6 (5.1 to 6.1) & \ensuremath{-}89 (\ensuremath{-}107 to \ensuremath{-}73) \\
 &  &  & Regression & \ensuremath{-}0.6 (\ensuremath{-}1.2 to \ensuremath{-}0.0) & \ensuremath{-}1.0 (\ensuremath{-}1.2 to \ensuremath{-}0.7) & \ensuremath{-}87 (\ensuremath{-}104 to \ensuremath{-}71) \\
\bottomrule
\end{tabular}

\end{table}

\subsection*{The law with data sets as units}
The prospective test's rank correlation was tested by permuting problems, but panels of 100, 200 and 400 features of one data set are nested, and the ten problems with a measurable saving came from four data sets of two studies (mouse lymphoid CITE-seq with two antibody panels; fly visual neurons in two connectomes). The law's error was therefore summarized per data set and without each data set in turn, the rank correlation was computed between data-set means with an exact permutation test over data sets and between problems with permutations only within data sets, and the slope of observed on predicted saving was estimated with a fixed effect per data set (Supplementary Table~\ref{tab:lawunits}). At half of the dependence, the law overpredicted in all \RvLawHalfSets{} data sets, with median errors of \RvLawHalfLooMin\% to \RvLawHalfLooMax\% when any one was left out (95\% CI of the median error \RvLawHalfErrLo--\RvLawHalfErrHi\% resampling data sets and \RvLawHalfErrStLo--\RvLawHalfErrStHi\% resampling studies). The data sets' averages were ranked in order ($\rho=\RvLawHalfMeansRho$, exact $P=\RvLawHalfMeansP$, the smallest attainable with four), the problems were ranked within data sets ($P=\RvLawHalfWithinP$ over all \RvLawHalfWithinPerm{} permutations within data sets), and the slope with data-set effects was \RvLawHalfSlope{} (\RvLawHalfSlopeLo--\RvLawHalfSlopeHi). At 30\% of the dependence the law overpredicted in all \RvLawThreeSets{} data sets with a measurable saving (median error \RvLawThreeErr\%); at 70\%, only the two fly data sets had one (\RvLawSevenErr\%). Of the \RvLawHalfNonN{} problems without a measurable saving at half of the dependence, \RvLawHalfBelow{} had both curves above the target with 25 paired cells, \RvLawHalfNotReached{} had a curve that did not reach it within the budgets, and \RvLawHalfOneAbove{} had one curve starting above it.

\begin{table}[h]
\caption{\textbf{The law with data sets as units.} Per data set with a measurable saving: problems, median absolute error of the law ($2^{|\log_2(\text{observed}/\text{law})|}-1$), median ratio of the observed to the predicted saving, and the median error over the other data sets. All: the median error with 95\% intervals from resampling data sets and, in the same parentheses, studies. Then the rank correlation of data-set means with its exact permutation $P$ value, the problem-level rank correlation with permutations within data sets, the slope of $\log$ observed on $\log$ predicted saving with a fixed effect per data set (95\% interval from resampling data sets), and the problems without a measurable saving by reason.}
\label{tab:lawunits}
\centering\footnotesize
% Generated by extension/synthesis/make_assets.py. Do not edit.
\begin{tabular}{@{}lrrrr@{}}
\toprule
Data set & Problems & Median error (\%) & Observed / law & Without this data set (\%) \\
\midrule
\multicolumn{5}{@{}l}{\textit{Half of the dependence recovered (primary)}} \\
\quad Fly visual neurons, BANC & 2 & 10 & 0.91 & 25 \\
\quad Fly visual neurons, FlyWire & 3 & 18 & 0.85 & 26 \\
\quad Mouse lymphoid organs, 111 proteins & 3 & 26 & 0.80 & 18 \\
\quad Mouse lymphoid organs, 206 proteins & 2 & 35 & 0.74 & 17 \\
\quad All & 10 & 21 (11--29; studies 13--26) & & \\
\multicolumn{5}{@{}p{.95\textwidth}}{\quad rank correlation of data-set means 1.00 (exact $P=0.042$, 4 data sets); problem-level $\rho=0.98$, permutations within data sets $P=0.007$ (144); slope with data-set effects 1.08 (0.99--1.26).} \\
\multicolumn{5}{@{}p{.95\textwidth}}{\quad Not measurable: 10 a curve does not reach the target within the budgets; 8 target below both curves at the smallest budget; 1 one curve starts above the target.} \\
\midrule
\multicolumn{5}{@{}l}{\textit{30\% recovered}} \\
\quad Human cortex Patch-seq & 2 & 144 & 0.41 & 61 \\
\quad Mouse lymphoid organs, 111 proteins & 3 & 68 & 0.60 & 81 \\
\quad Mouse lymphoid organs, 206 proteins & 3 & 54 & 0.65 & 99 \\
\quad All & 8 & 74 (54--144; studies 61--144) & & \\
\multicolumn{5}{@{}p{.95\textwidth}}{\quad rank correlation of data-set means 1.00 (exact $P=0.167$, 3 data sets); problem-level $\rho=0.95$, permutations within data sets $P=0.028$ (72); slope with data-set effects 0.98 (-2.39--1.22).} \\
\multicolumn{5}{@{}p{.95\textwidth}}{\quad Not measurable: 9 target below both curves at the smallest budget; 7 a curve does not reach the target within the budgets; 5 one curve starts above the target.} \\
\midrule
\multicolumn{5}{@{}l}{\textit{70\% recovered}} \\
\quad Fly visual neurons, BANC & 2 & 2 & 1.02 & 5 \\
\quad Fly visual neurons, FlyWire & 3 & 5 & 0.95 & 2 \\
\quad All & 5 & 4 (2--5; studies 4--4) & & \\
\multicolumn{5}{@{}p{.95\textwidth}}{\quad problem-level $\rho=1.00$, permutations within data sets $P=0.083$ (12); slope with data-set effects 1.02 (0.99--1.15).} \\
\multicolumn{5}{@{}p{.95\textwidth}}{\quad Not measurable: 15 a curve does not reach the target within the budgets; 9 one curve starts above the target.} \\
\bottomrule
\end{tabular}

\end{table}

\subsection*{A pilot that sets the budget}
For each of the 29 problems of the prospective test, target accuracy (a recovered fraction of 0.3, 0.5 or 0.7) and pilot size (50, 100 or 200 paired cells, at most a fifth of the reservoir), 100 replicates drew a random pilot from the reservoir, estimated each block's signal and noise from it, set the total budget to the paired cells that Proposition~\ref{prop:law} gives for block shrinkage at the target (never fewer than the pilot), drew the remaining cells so that the pilot counted towards the total, fitted block James--Stein to all of them and scored the recovered fraction against the truth cells. A replicate whose pilot showed no positive signal, or whose budget exceeded the reservoir, used the whole reservoir and was flagged. The same fit at the budget chosen from all reservoir cells served as reference (Supplementary Table~\ref{tab:pilotdesign}). At half of the dependence, pilots of 50, 100 and 200 cells chose budgets that reached the target in \RvPilotHalfFiftyReached\%, \RvPilotHalfHundredReached\% and \RvPilotHalfTwoHundredReached\% of replicates, and reached at least the target less 5 percentage points in \RvPilotHalfFiftyWithin\%, \RvPilotHalfHundredWithin\% and \RvPilotHalfTwoHundredWithin\%. Their median totals were \RvPilotHalfFiftyBudget, \RvPilotHalfHundredBudget{} and \RvPilotHalfTwoHundredBudget{} paired cells: from 100 cells on, the pilot alone usually exceeded the budget that the law gave. The budget from all reservoir cells was a median of \RvPilotHalfResRatio{} times the paired cells that block shrinkage needed, and it reached the target in \RvPilotThreeResReached\%, \RvPilotHalfResReached\% and \RvPilotSevenResReached\% of replicates at 30\%, 50\% and 70\%: the law assumes the best factor in every block, and the positive-part rule needs more cells than that factor would. At 70\%, \RvPilotSevenFiftyFlagged\% to \RvPilotSevenTwoHundredFlagged\% of replicates were flagged. By data set, pilots fell more than 5 points short of the target in more than half of the replicates only in human cortical Patch-seq and in the chromatin of fetal cortex at 50\%, and in these and the chromatin of adult mouse cortex at 70\%.

\begin{table}[h]
\caption{\textbf{A pilot that sets the budget.} For each target and pilot size, over the 29 problems of the prospective test (each weighted alike; problems whose reservoir was too small for a pilot size used the largest pilot allowed, a fifth of the reservoir): the median over problems of each problem's median achieved recovered fraction, with the interquartile range over problems; the share of replicates reaching the target, and the share reaching at least the target less 5 percentage points (overshooting included); the median total budget, pilot included; its median ratio to the paired cells that the observed curve needed (problems whose curve crossed the target); and the share of replicates flagged (no positive signal in the pilot, or a budget beyond the reservoir, so that the whole reservoir was used). All reservoir cells: the budget chosen with the block moments of every reservoir cell.}
\label{tab:pilotdesign}
\centering\scriptsize
% Generated by extension/synthesis/make_assets.py. Do not edit.
\begin{tabular}{@{}llrrrrrrr@{}}
\toprule
 & & & Achieved & Reached & Reached target & & Budget / & Flagged \\
Target & Pilot & Problems & (\%; IQR) & (\%) & $-$5 points (\%) & Budget & needed & (\%) \\
\midrule
30\% & 50 cells & 29 & 41 (34--90) & 73 & 78 & 50 & 0.99 & 0 \\
 & 100 cells & 29 & 41 (35--94) & 77 & 82 & 100 & 1.02 & 1 \\
 & 200 cells & 29 & 48 (34--97) & 83 & 87 & 200 & 1.16 & 1 \\
 & All reservoir cells & 29 & 23 (12--26) & 34 & 45 & 32 & 0.57 & 0 \\
\addlinespace
50\% & 50 cells & 29 & 54 (48--90) & 64 & 72 & 84 & 1.16 & 5 \\
 & 100 cells & 29 & 53 (48--94) & 65 & 75 & 100 & 1.12 & 8 \\
 & 200 cells & 29 & 53 (48--97) & 65 & 77 & 200 & 0.87 & 10 \\
 & All reservoir cells & 29 & 46 (40--48) & 26 & 50 & 86 & 0.77 & 0 \\
\addlinespace
70\% & 50 cells & 29 & 68 (60--90) & 46 & 56 & 181 & 0.82 & 22 \\
 & 100 cells & 29 & 69 (61--94) & 48 & 59 & 181 & 0.85 & 24 \\
 & 200 cells & 29 & 70 (61--97) & 56 & 61 & 200 & 0.99 & 25 \\
 & All reservoir cells & 29 & 67 (61--69) & 23 & 49 & 181 & 0.82 & 0 \\
\bottomrule
\end{tabular}

\end{table}

\subsection*{Savings at fixed accuracy}
The savings summary uses each data set's own target and counts curves that never reach it at the largest budget. At fixed recovered fractions from 5\% to 40\%, the paired cells of the estimator and of paired-only estimation were read from the published curves both with and without the running maximum, and each comparison was classed as estimable (both curves cross the level between two budgets), unresolved above (a curve never reaches it) or unresolved below (both start above it); estimable ratios have 95\% intervals from the tests' resamples of held-out units (Supplementary Tables~\ref{tab:fixed} and~\ref{tab:fixeddetail}). Read from the raw curves, the saving could be measured in \RvFixedTenN, \RvFixedTwentyN{} and \RvFixedThirtyN{} of the \RvFixedSets{} data sets at 10\%, 20\% and 30\% recovery (geometric means \RvFixedTenGm, \RvFixedTwentyGm{} and \RvFixedThirtyGm), and the running maximum changed \RvFixedSameRunningMax{} of the counts. At 20\%, the other comparisons were unresolved because paired-only estimation did not reach the level within 800 paired cells (\RvFixedTwentyPoAbove{} data sets), because neither method did (\RvFixedTwentyBothAbove) or because the estimator had already exceeded it with 25 paired cells (\RvFixedTwentyEstBelow). The summary's savings, read at each data set's own target and counting curves that never reach it, therefore rest on comparisons that fixed levels mostly leave unresolved.

\begin{table}[h]
\caption{\textbf{Savings at fixed accuracy.} Over the nine benchmark data sets and the external test, at each recovered fraction: comparisons that could be made (both the estimator and the paired-only envelope cross the level between two budgets), the geometric mean of their savings, and comparisons left unresolved because a curve never reached the level (above) or one or both curves started above it with 25 paired cells (below), read from the raw curves and from their running maximum.}
\label{tab:fixed}
\centering\footnotesize
% Generated by extension/synthesis/make_assets.py. Do not edit.
\begin{tabular}{@{}lrrrrrrrr@{}}
\toprule
 & \multicolumn{4}{c}{Raw curves} & \multicolumn{4}{c}{Running maximum} \\
\cmidrule(lr){2-5}\cmidrule(l){6-9}
Recovered & Estimable & Geometric mean & Above & Below & Estimable & Geometric mean & Above & Below \\
\midrule
5\% & 3 & 4.29 & 4 & 3 & 3 & 4.29 & 4 & 3 \\
10\% & 2 & 1.74 & 6 & 2 & 2 & 1.74 & 6 & 2 \\
15\% & 2 & 2.07 & 6 & 2 & 2 & 2.07 & 6 & 2 \\
20\% & 2 & 2.46 & 6 & 2 & 2 & 2.46 & 6 & 2 \\
30\% & 1 & 4.88 & 7 & 2 & 1 & 4.88 & 7 & 2 \\
40\% & 2 & 3.39 & 8 & 0 & 2 & 3.39 & 8 & 0 \\
\bottomrule
\end{tabular}

\end{table}

\begin{table}[h]
\caption{\textbf{Savings at fixed accuracy per data set.} Ratio of the paired cells needed by paired-only estimation to those needed by the estimator at 10\%, 20\% and 30\% of the dependence recovered, read from the raw curves, with 95\% intervals from the tests' resamples of held-out units; otherwise the reason the comparison is unresolved.}
\label{tab:fixeddetail}
\centering\footnotesize
% Generated by extension/synthesis/make_assets.py. Do not edit.
\begin{tabular}{@{}lrrr@{}}
\toprule
Data set & 10\% recovered & 20\% recovered & 30\% recovered \\
\midrule
External test (OverCITE-seq) & 2.62 (2.06--3.27) & 4.35 (3.21--5.19) & 4.88 (3.22--6.03) \\
Blood (Hao) & paired only not reached & paired only not reached & paired only not reached \\
Blood (Stephenson) & not reached & not reached & not reached \\
Bone marrow CITE-seq & paired only not reached & paired only not reached & paired only not reached \\
Colon & paired only not reached & paired only not reached & not reached \\
Bone marrow RNA--chromatin & paired only not reached & paired only not reached & not reached \\
Patch-seq, visual cortex & 1.15 (0.97--1.33) & 1.40 (1.17--1.59) & paired only not reached \\
Patch-seq, motor cortex & not reached & not reached & not reached \\
Fly, female & both above at 25 & estimator above at 25 & estimator above at 25 \\
Fly, male & both above at 25 & estimator above at 25 & estimator above at 25 \\
\bottomrule
\end{tabular}

\end{table}

\bibliographystyle{plainnat}
\bibliography{references}
\end{document}
